% manuscript-sha256: 37f5bb391cab2093baa8efd20d561372f8b041a9481e317853f56f1538f33dcb
% Options for packages loaded elsewhere
\PassOptionsToPackage{unicode}{hyperref}
\PassOptionsToPackage{hyphens}{url}
\PassOptionsToPackage{dvipsnames,svgnames,x11names}{xcolor}
\documentclass[
  11pt,
]{article}
\usepackage{xcolor}
\usepackage{amsmath,amssymb}
\setcounter{secnumdepth}{-\maxdimen} % remove section numbering
\usepackage{iftex}
\ifPDFTeX
  \usepackage[T1]{fontenc}
  \usepackage[utf8]{inputenc}
  \usepackage{textcomp} % provide euro and other symbols
\else % if luatex or xetex
  \usepackage{unicode-math} % this also loads fontspec
  \defaultfontfeatures{Scale=MatchLowercase}
  \defaultfontfeatures[\rmfamily]{Ligatures=TeX,Scale=1}
\fi
\usepackage{lmodern}
\ifPDFTeX\else
  % xetex/luatex font selection
\fi
% Use upquote if available, for straight quotes in verbatim environments
\IfFileExists{upquote.sty}{\usepackage{upquote}}{}
\IfFileExists{microtype.sty}{% use microtype if available
  \usepackage[]{microtype}
  \UseMicrotypeSet[protrusion]{basicmath} % disable protrusion for tt fonts
}{}
\makeatletter
\@ifundefined{KOMAClassName}{% if non-KOMA class
  \IfFileExists{parskip.sty}{%
    \usepackage{parskip}
  }{% else
    \setlength{\parindent}{0pt}
    \setlength{\parskip}{6pt plus 2pt minus 1pt}}
}{% if KOMA class
  \KOMAoptions{parskip=half}}
\makeatother
\usepackage{color}
\usepackage{fancyvrb}
\newcommand{\VerbBar}{|}
\newcommand{\VERB}{\Verb[commandchars=\\\{\}]}
\DefineVerbatimEnvironment{Highlighting}{Verbatim}{commandchars=\\\{\}}
% Add ',fontsize=\small' for more characters per line
\usepackage{framed}
\definecolor{shadecolor}{RGB}{255,255,255}
\newenvironment{Shaded}{\begin{snugshade}}{\end{snugshade}}
\newcommand{\AlertTok}[1]{\textcolor[rgb]{0.75,0.01,0.01}{\textbf{\colorbox[rgb]{0.97,0.90,0.90}{#1}}}}
\newcommand{\AnnotationTok}[1]{\textcolor[rgb]{0.79,0.38,0.79}{#1}}
\newcommand{\AttributeTok}[1]{\textcolor[rgb]{0.00,0.34,0.68}{#1}}
\newcommand{\BaseNTok}[1]{\textcolor[rgb]{0.69,0.50,0.00}{#1}}
\newcommand{\BuiltInTok}[1]{\textcolor[rgb]{0.39,0.29,0.61}{\textbf{#1}}}
\newcommand{\CharTok}[1]{\textcolor[rgb]{0.57,0.30,0.62}{#1}}
\newcommand{\CommentTok}[1]{\textcolor[rgb]{0.54,0.53,0.53}{#1}}
\newcommand{\CommentVarTok}[1]{\textcolor[rgb]{0.00,0.58,1.00}{#1}}
\newcommand{\ConstantTok}[1]{\textcolor[rgb]{0.67,0.33,0.00}{#1}}
\newcommand{\ControlFlowTok}[1]{\textcolor[rgb]{0.12,0.11,0.11}{\textbf{#1}}}
\newcommand{\DataTypeTok}[1]{\textcolor[rgb]{0.00,0.34,0.68}{#1}}
\newcommand{\DecValTok}[1]{\textcolor[rgb]{0.69,0.50,0.00}{#1}}
\newcommand{\DocumentationTok}[1]{\textcolor[rgb]{0.38,0.47,0.50}{#1}}
\newcommand{\ErrorTok}[1]{\textcolor[rgb]{0.75,0.01,0.01}{\underline{#1}}}
\newcommand{\ExtensionTok}[1]{\textcolor[rgb]{0.00,0.58,1.00}{\textbf{#1}}}
\newcommand{\FloatTok}[1]{\textcolor[rgb]{0.69,0.50,0.00}{#1}}
\newcommand{\FunctionTok}[1]{\textcolor[rgb]{0.39,0.29,0.61}{#1}}
\newcommand{\ImportTok}[1]{\textcolor[rgb]{1.00,0.33,0.00}{#1}}
\newcommand{\InformationTok}[1]{\textcolor[rgb]{0.69,0.50,0.00}{#1}}
\newcommand{\KeywordTok}[1]{\textcolor[rgb]{0.12,0.11,0.11}{\textbf{#1}}}
\newcommand{\NormalTok}[1]{\textcolor[rgb]{0.12,0.11,0.11}{#1}}
\newcommand{\OperatorTok}[1]{\textcolor[rgb]{0.12,0.11,0.11}{#1}}
\newcommand{\OtherTok}[1]{\textcolor[rgb]{0.00,0.43,0.16}{#1}}
\newcommand{\PreprocessorTok}[1]{\textcolor[rgb]{0.00,0.43,0.16}{#1}}
\newcommand{\RegionMarkerTok}[1]{\textcolor[rgb]{0.00,0.34,0.68}{\colorbox[rgb]{0.88,0.91,0.97}{#1}}}
\newcommand{\SpecialCharTok}[1]{\textcolor[rgb]{0.24,0.68,0.91}{#1}}
\newcommand{\SpecialStringTok}[1]{\textcolor[rgb]{1.00,0.33,0.00}{#1}}
\newcommand{\StringTok}[1]{\textcolor[rgb]{0.75,0.01,0.01}{#1}}
\newcommand{\VariableTok}[1]{\textcolor[rgb]{0.00,0.34,0.68}{#1}}
\newcommand{\VerbatimStringTok}[1]{\textcolor[rgb]{0.75,0.01,0.01}{#1}}
\newcommand{\WarningTok}[1]{\textcolor[rgb]{0.75,0.01,0.01}{#1}}
\usepackage{longtable,booktabs,array}
\usepackage{calc} % for calculating minipage widths
% Correct order of tables after \paragraph or \subparagraph
\usepackage{etoolbox}
\makeatletter
\patchcmd\longtable{\par}{\if@noskipsec\mbox{}\fi\par}{}{}
\makeatother
% Allow footnotes in longtable head/foot
\IfFileExists{footnotehyper.sty}{\usepackage{footnotehyper}}{\usepackage{footnote}}
\makesavenoteenv{longtable}
\setlength{\emergencystretch}{3em} % prevent overfull lines
\providecommand{\tightlist}{%
  \setlength{\itemsep}{0pt}\setlength{\parskip}{0pt}}
% House preamble for the primitive-axiomatics paper line.
%
%  Typesetting is not evidence for any claim.
%
% Built 2026-09-08.
%
% xelatex, not pdflatex: the manuscript uses unicode characters that pdflatex
% cannot set without per-character declarations.

\PassOptionsToPackage{protrusion=false}{microtype}
\usepackage{fontspec}
\usepackage{unicode-math}
\setmainfont{Latin Modern Roman}
\setmonofont{Latin Modern Mono}[Scale=0.85]
\setmathfont{Latin Modern Math}
\newfontfamily\wdpsymbolfont{Latin Modern Math}

\usepackage{amsmath,amssymb,amsthm}
\usepackage{newunicodechar}
\newcommand{\um}[1]{\ensuremath{#1}}
\newunicodechar{ℝ}{\um{\mathbb{R}}}
\newunicodechar{⁰}{\um{^0}}
\newunicodechar{₀}{\um{_0}}
\newunicodechar{₁}{\um{_1}}
\newunicodechar{₂}{\um{_2}}
\newunicodechar{¹}{\um{^1}}
\newunicodechar{²}{\um{^2}}
\newunicodechar{∘}{\um{\circ}}
\newunicodechar{δ}{\um{\delta}}
\newunicodechar{⟨}{\um{\langle}}
\newunicodechar{⟩}{\um{\rangle}}
\newunicodechar{φ}{\um{\varphi}}
\newunicodechar{ᵀ}{\um{^\top}}
\newunicodechar{𝛷}{\um{\Phi}}
\newunicodechar{Φ}{\um{\Phi}}
\newunicodechar{Γ}{\um{\Gamma}}
\newunicodechar{α}{\um{\alpha}}
\newunicodechar{β}{\um{\beta}}
\newunicodechar{κ}{\um{\kappa}}
\newunicodechar{∂}{\um{\partial}}
\newunicodechar{∈}{\um{\in}}
\newunicodechar{−}{\um{-}}
\newunicodechar{∞}{\um{\infty}}
\newunicodechar{≠}{\um{\neq}}
\newunicodechar{≤}{\um{\le}}
\newunicodechar{≥}{\um{\ge}}
\newunicodechar{⟺}{\um{\iff}}
\newunicodechar{→}{\um{\to}}
\newunicodechar{─}{-}
\newunicodechar{│}{|}
\newunicodechar{┌}{+}
\newunicodechar{┐}{+}
\newunicodechar{└}{+}
\newunicodechar{┘}{+}
\newunicodechar{┬}{+}
\newunicodechar{┴}{+}
\newunicodechar{▼}{v}
\newunicodechar{ℓ}{\um{\ell}}
\newunicodechar{Ω}{\um{\Omega}}
\newunicodechar{Σ}{\um{\Sigma}}
\newunicodechar{·}{\um{\cdot}}
\newunicodechar{×}{\um{\times}}
\newunicodechar{⊂}{\um{\subset}}
\newunicodechar{ₙ}{\um{_n}}
\newunicodechar{∎}{\um{\blacksquare}}
\newunicodechar{∼}{\um{\sim}}
\newunicodechar{∖}{\um{\setminus}}
\newunicodechar{∩}{\um{\cap}}
\newunicodechar{ε}{\um{\varepsilon}}
\newunicodechar{⁻}{\um{^-}}
\newunicodechar{↦}{\um{\mapsto}}
\newunicodechar{∅}{\um{\emptyset}}
\newunicodechar{⋯}{\um{\cdots}}
\newunicodechar{⊆}{\um{\subseteq}}
\newunicodechar{₅}{\um{_5}}
% Draw the remaining mathematical characters from Latin Modern Math.
% Slot numbers avoid recursion through unicode-math's active characters.
% Modifier letters and the two primes use ordinary super/subscripts because
% Latin Modern Math has no glyph for those code points. Primes are declared
% at begin-document because unicode-math reclaims them.
\newunicodechar{∥}{\ensuremath{\text{\wdpsymbolfont\char"2225}}}
\newunicodechar{σ}{\ensuremath{\text{\wdpsymbolfont\char"03C3}}}
\newunicodechar{ρ}{\ensuremath{\text{\wdpsymbolfont\char"03C1}}}
\newunicodechar{θ}{\ensuremath{\text{\wdpsymbolfont\char"03B8}}}
\newunicodechar{⊗}{\ensuremath{\text{\wdpsymbolfont\char"2297}}}
\newunicodechar{⪰}{\ensuremath{\text{\wdpsymbolfont\char"2AB0}}}
\newunicodechar{ψ}{\ensuremath{\text{\wdpsymbolfont\char"03C8}}}
\newunicodechar{ν}{\ensuremath{\text{\wdpsymbolfont\char"03BD}}}
\newunicodechar{⋃}{\ensuremath{\text{\wdpsymbolfont\char"22C3}}}
\newunicodechar{τ}{\ensuremath{\text{\wdpsymbolfont\char"03C4}}}
\newunicodechar{⇒}{\ensuremath{\text{\wdpsymbolfont\char"21D2}}}
\newunicodechar{η}{\ensuremath{\text{\wdpsymbolfont\char"03B7}}}
\newunicodechar{⪯}{\ensuremath{\text{\wdpsymbolfont\char"2AAF}}}
\newunicodechar{⇐}{\ensuremath{\text{\wdpsymbolfont\char"21D0}}}
\newunicodechar{⊕}{\ensuremath{\text{\wdpsymbolfont\char"2295}}}
\newunicodechar{⇔}{\ensuremath{\text{\wdpsymbolfont\char"21D4}}}
\newunicodechar{∣}{\ensuremath{\text{\wdpsymbolfont\char"2223}}}
\newunicodechar{λ}{\ensuremath{\text{\wdpsymbolfont\char"03BB}}}
\newunicodechar{∪}{\ensuremath{\text{\wdpsymbolfont\char"222A}}}
\newunicodechar{≡}{\ensuremath{\text{\wdpsymbolfont\char"2261}}}
\newunicodechar{π}{\ensuremath{\text{\wdpsymbolfont\char"03C0}}}
\newunicodechar{∫}{\ensuremath{\text{\wdpsymbolfont\char"222B}}}
\newunicodechar{∉}{\ensuremath{\text{\wdpsymbolfont\char"2209}}}
\newunicodechar{ω}{\ensuremath{\text{\wdpsymbolfont\char"03C9}}}
\newunicodechar{⊇}{\ensuremath{\text{\wdpsymbolfont\char"2287}}}
\newunicodechar{∏}{\ensuremath{\text{\wdpsymbolfont\char"220F}}}
\newunicodechar{⌊}{\ensuremath{\text{\wdpsymbolfont\char"230A}}}
\newunicodechar{⌋}{\ensuremath{\text{\wdpsymbolfont\char"230B}}}
\newunicodechar{μ}{\ensuremath{\text{\wdpsymbolfont\char"03BC}}}
\newunicodechar{⁺}{\textsuperscript{+}}
\newunicodechar{ⁿ}{\textsuperscript{n}}
\newunicodechar{ᵖ}{\textsuperscript{p}}
\newunicodechar{ᵐ}{\textsuperscript{m}}
\newunicodechar{ᵉ}{\textsuperscript{e}}
\newunicodechar{ʳ}{\textsuperscript{r}}
\newunicodechar{⁴}{\textsuperscript{4}}
\newunicodechar{⁵}{\textsuperscript{5}}
\newunicodechar{⁶}{\textsuperscript{6}}
\newunicodechar{⁷}{\textsuperscript{7}}
\newunicodechar{₃}{\textsubscript{3}}
\newunicodechar{₄}{\textsubscript{4}}
\newunicodechar{↝}{\ensuremath{\rightsquigarrow}}
\AtBeginDocument{%
  \newunicodechar{′}{\textsuperscript{\ensuremath{\prime}}}%
  \newunicodechar{″}{\textsuperscript{\ensuremath{\prime\prime}}}%
}
\usepackage[margin=1in]{geometry}
\usepackage{longtable,booktabs,array}
\usepackage{pdflscape}
\usepackage[protrusion=false]{microtype}
\microtypesetup{protrusion=false}
\usepackage{enumitem}
\setlist{nosep,leftmargin=1.4em}

\usepackage[dvipsnames]{xcolor}
\usepackage[colorlinks=true,linkcolor=RoyalBlue,urlcolor=RoyalBlue,
            citecolor=RoyalBlue,breaklinks=true]{hyperref}
\usepackage{fvextra}
\usepackage{framed}
\fvset{breaklines=true,breakanywhere=true,fontsize=\small}

\AtBeginEnvironment{longtable}{\small}
\setlength{\tabcolsep}{4pt}
\renewcommand{\arraystretch}{1.15}

\providecommand{\tightlist}{%
  \setlength{\itemsep}{0pt}\setlength{\parskip}{0pt}}

\definecolor{quoterule}{gray}{0.75}
\renewenvironment{quote}
  {\begin{leftbar}\small}
  {\end{leftbar}}

\usepackage[htt]{hyphenat}
\usepackage{xurl}
% Reserve enough TOC space for this monograph's three-digit page numbers.
\makeatletter
\renewcommand{\@pnumwidth}{2.5em}
\renewcommand{\@tocrmarg}{3.5em}
\makeatother
\emergencystretch=3em
\sloppy

\newenvironment{nonclaim}{\begin{quote}\textbf{Non-claim:\ }}{\end{quote}}
\newenvironment{claimblock}{\begin{quote}}{\end{quote}}
\usepackage{bookmark}
\IfFileExists{xurl.sty}{\usepackage{xurl}}{} % add URL line breaks if available
\urlstyle{same}
\hypersetup{
  pdftitle={The Witnessed Difference Programme - From Finite Distinction to Witnessed Measurement},
  pdfauthor={Jeremy H. Carroll},
  colorlinks=true,
  linkcolor={RoyalBlue},
  filecolor={Maroon},
  citecolor={Blue},
  urlcolor={Blue},
  pdfcreator={LaTeX via pandoc}}

\author{}
\date{2026-10-16}

\begin{document}

\begin{center}
{\LARGE\bfseries The Witnessed Difference Programme - From Finite Distinction to Witnessed Measurement \par}
\vspace{0.5em}
{\large Finite obstruction, information loss, and sufficient observation \par}
\vspace{1.0em}
{\textbf{Jeremy H. Carroll}} \par
\vspace{0.4em}
{\small 2026-10-16 \quad $\cdot$ \quad Version 2.4.5} \par
\vspace{0.2em}
{\footnotesize Creative Commons Attribution 4.0 International} \par
\end{center}
\vspace{1.0em}

{
\hypersetup{linkcolor=}
\setcounter{tocdepth}{2}
\tableofcontents
}
\subsection{Abstract}\label{abstract}

The Witnessed Difference Programme studies what must be specified for a
distinction to be retained, compared and reported without losing the question
being asked. Its finite core begins with explicit comparison data, permitted
changes and report maps. A report preserves a question exactly when the answer
is constant on each report fiber. With a declared linear repair map,
discrepancies are compared modulo its image; finite cochain complexes separate
removable discrepancies, closed residual classes and failures of closure.
Exact certificates can check proposed repairs and lower bounds on the remaining
discrepancy.

A central result distinguishes defects that can occur from changes admitted as
repairs. For finite rational linear data it gives sharp counts of additional
rational linear observation rows needed to distinguish realizable defects,
repair classes on the entire defect space, or both. Further results treat transport between presentations, retained
histories and certified progress in finite decision problems. A paired
count-guard family has the homotopy type of a real no-equal arrangement
complement, giving an exact capacity threshold for coherence and a first
nonzero homology obstruction.

The wider programme asks how these constructions can support richer
mathematical and physical theories while making every further assumption
explicit. Later parts develop additive and tensor constructions, conditional
complex representations, finite quantum states and processes, and selected
model-specific limits. Each depends on declared algebraic, operational,
topological or analytic inputs. The present work does not derive those inputs,
or quantum theory, from an unstructured mark. Written proofs, selected Lean
formalizations, executable checks and external interpretations are recorded
separately.

\subsection{Mechanism inventory}\label{mechanism-inventory}

The monograph is one \textbf{Paper} containing the following internal papers.
They share the protected-question discipline, while retaining their separate
input contracts and proof obligations.

\begin{longtable}[]{@{}
  >{\raggedright\arraybackslash}p{(\linewidth - 4\tabcolsep) * \real{0.3333}}
  >{\raggedright\arraybackslash}p{(\linewidth - 4\tabcolsep) * \real{0.3333}}
  >{\raggedright\arraybackslash}p{(\linewidth - 4\tabcolsep) * \real{0.3333}}@{}}
\toprule\noalign{}
\begin{minipage}[b]{\linewidth}\raggedright
Internal paper
\end{minipage} & \begin{minipage}[b]{\linewidth}\raggedright
Location and mechanism
\end{minipage} & \begin{minipage}[b]{\linewidth}\raggedright
Evidence boundary
\end{minipage} \\
\midrule\noalign{}
\endhead
\bottomrule\noalign{}
\endlastfoot
Finite distinction and probability & Chapters 1--9 & Finite carriers, declared weights and cochain identities; optional quantum structure is explicit \\
Observer and measurement & Chapters 10--15 and 23; Appendices T--V, X--Y & Protected questions, calibrated records and finite recovery; authority is a separate input \\
Histories and admissible operations & Chapters 16--17; Appendices I--L, N, P & Declared history relations, operation closure and finite schedule witnesses \\
Conditional quantum interfaces & Chapters 18--19A & Scalar, positivity, readout and composition premises are not consequences of report minimality \\
Limits and analytic extensions & Appendices M, O, W & Analytic hypotheses and external theorems remain separate from finite and Lean checks \\
Relative completeness and proof supplements & Chapters 20--22; Appendices A--H, Q--S & Question-relative theorems, source provenance and bounded mechanization scopes \\
\end{longtable}

\subsection{Reader's mathematical entrance}\label{readers-mathematical-entrance}

The organizing question is simple: when a presentation changes, which
specified questions remain answerable, which differences can the declared
operations repair, and what additional observations suffice when it loses a
needed distinction?

Let \(p:X\to A\) be a report and \(q:X\to Z\) the answer that must be retained.
A decoder \(\bar q:p(X)\to Z\) with \(q=\bar q\circ p\) exists precisely when

\[
p(x)=p(y)\Longrightarrow q(x)=q(y).
\]

This standard factorization criterion gives a constructive decoder when it
holds and a collision pair when it fails. For example, reporting the first bit
of \((a,b)\) cannot answer the question ``what is \(b\)?'' because \((0,0)\) and
\((0,1)\) have the same report but different answers. Appendix P.1 supplies the
full statement and its transport uses.

The first substantial finite theorem refines this question by separating
realizable defects from admitted repairs. Given rational linear maps with
\(R=\operatorname{im}D_s\), \(B=\operatorname{im}D_r\) and existing report
kernel \(K=\ker Q\), unrestricted additional rational linear observation rows
require exactly \(\dim(R\cap K)\) rows to distinguish individual realizable
defects in \(R\), and \(\dim K-\dim(K\cap B)\) rows to distinguish repair classes
\textbf{globally on \(V\)}, meaning \(\ker(Q,E)\subseteq B\). Ensuring both of these
properties requires \(\dim K-\dim(K\cap B)+\dim(R\cap K\cap B)\) rows.
If instead only candidates in \(R\) must be distinguished up to repair, using
their actual augmented reports, the minimum is
\(\dim(R\cap K)-\dim(R\cap K\cap B)\). These are different questions.
Appendix T supplies the matching lower bounds and constructions for the global
theorem WDP-DIM-1; T.8 states and derives the realizable-domain corollary. Restricted detector menus, noise, conditioning
and physical instruments are separate models, not consequences of this count.

A three-edge cycle gives the complementary repair example. A field assigning
one to every oriented edge has circulation three, whereas every vertex-potential
change has circulation zero. It cannot be erased by those permitted repairs;
with the declared \(\ell^1\) norm, the exact residual cost is three. This is a
finite algebraic obstruction, not yet a physical interpretation.

\subsection{Theorem map: when a presentation preserves the question}\label{theorem-map-when-a-presentation-preserves-the-question}

The finite mathematical thread is whether a changed presentation retains
the \textbf{specified question}. Fifteen results carry that thread. Each row names
the assumptions that make its conclusion meaningful; none derives those
assumptions from the initial mark.

Canonical bindings below abbreviate the fixed \texttt{WDP-} prefix before \texttt{CH}.
Register bindings name the bridge register; R12/R13 bindings name the separate
supplement table in the claim ledger. Row IDs are stable identifiers, not
locations, and the binding does not expand the cited evidence scope.

\begin{longtable}[]{@{}
  >{\raggedright\arraybackslash}p{(\linewidth - 6\tabcolsep) * \real{0.2500}}
  >{\raggedright\arraybackslash}p{(\linewidth - 6\tabcolsep) * \real{0.2500}}
  >{\raggedright\arraybackslash}p{(\linewidth - 6\tabcolsep) * \real{0.2500}}
  >{\raggedright\arraybackslash}p{(\linewidth - 6\tabcolsep) * \real{0.2500}}@{}}
\toprule\noalign{}
\begin{minipage}[b]{\linewidth}\raggedright
Result
\end{minipage} & \begin{minipage}[b]{\linewidth}\raggedright
Premises and dependency
\end{minipage} & \begin{minipage}[b]{\linewidth}\raggedright
Exact conclusion and contribution
\end{minipage} & \begin{minipage}[b]{\linewidth}\raggedright
Binding
\end{minipage} \\
\midrule\noalign{}
\endhead
\bottomrule\noalign{}
\endlastfoot
\textbf{R1. Repair obstruction and certificate} (§4) & A declared finite cochain window, admissible repairs, and an exact norm/linear programme. & \texttt{Phi({[}h-d0\ f{]})=Phi({[}h{]})}: changing a representative cannot remove a nonzero quotient obstruction. A matched primal repair and bounded dual circulation certify the optimum. Quotient algebra is standard; the typed repair/certificate interface is this paper's construction. & CH04-P01; CH04-C01 \\
\textbf{R2. Protected-query factorization} (Appendices J.2, P.1) & Actual maps \texttt{f:X-\textgreater{}Y} and protected query \texttt{q:X-\textgreater{}Z}. & A decoder on \texttt{f(X)} exists exactly when \texttt{q} is constant on every \texttt{f}-fiber; a collision pair is a negative witness. The set-theoretic lemma is standard; its use to audit reported loss and future questions is this paper's application. & GEN-WDP-03 \\
\textbf{R3. Behavioral quotient} (CH03-O01, Appendix I.1) & Finite total labelled updates and exact readouts or validated report equivalence. & Partition refinement yields the greatest report-respecting update congruence and its minimal quotient. The algorithmic principle is standard; the declared observer/report interface fixes what must be preserved. & CH03-O01 \\
\textbf{R4. Carrier transport and reflection} (Appendix P.3; admission/transport supplement) & Supplied cochain maps, with a retraction or homotopy only when explicitly given. & Commuting squares transport repair classes; a down-map retraction reflects equality on the old image. A homotopy identifies closed classes but need not identify arbitrary repair classes. Standard cochain algebra is applied to concrete carrier-change and loss boundaries. & GEN-WDP-71 \\
\textbf{R5. Finite local/global test} (CH11-O01, Appendix J.2) & A specified finite context site with coherent restrictions and a declared global report set. & A matching family glues uniquely exactly when its restriction fiber contains one global report. The sheaf criterion is standard; the constrained-assignment countermodel locates this paper's finite obstruction. & CH11-O01 \\
\textbf{R6. Guarded five-edit model histories} (Appendix P.13--P.21) & The actual five-edit raw semantics, tagged events, stable birth names when needed, and declared precedence/guards. & Legal schedules have an exact raw endpoint. Under all-order legality, ideal prices descend by zero diamond curls and swap proofs have a local presentation. P.18 shows that guarded legal schedules can be disconnected despite equal endpoints, and gives the exact componentwise price criterion. P.20 establishes that star-shaped legality makes every legal swap complex connected with trivial fundamental group, so parallel legal swap proofs are joined by legal local 3-witnesses, and that compatibly ordered native precedences give star-shaped legality. P.21 shows that, for at most three literal terms, count total at most four and two successful edits, preserving the joint mark, literal balance record, native offers, prices and continuations forces raw equality (29,771 classes) under either terminal convention, attained by the report of mark, faces and counts capped at two. These are selected finite source theorems, not a global source quotient functor. & GEN-WDP-40--45; GEN-WDP-69; GEN-WDP-73 \\
\textbf{R7. Contextual histories and additive tensor completion} (§17.7; derivation supplements) & An admitted typed generating signature, endpoint-typed protected questions, finite alternative formation, signed completion, and an explicitly required biadditive independent composition. & Two-sided contextual equivalence is the greatest congruence preserving every protected question in every compatible context. Its free additive envelope has the unique homomorphism factorization property of Grothendieck group completion, and the free signed group on generator pairs represents biadditive maps. This derives an integer additive tensor calculus while leaving dagger, positivity, complex scalars, Hilbert completion and Born readout open. & GEN-WDP-58--59; GEN-WDP-72 \\
\textbf{R8. Dagger and observer Gram radical} (§17.8; derivation supplement) & A declared involution on generators, dagger-closed protected questions when descending through compression, and a supplied additive observer into a finite integer coordinate carrier. & Reversal extends uniquely to an involutive additive dagger and acts componentwise on the algebraic tensor. The observer induces a positive semidefinite Gram form whose radical is exactly its kernel, so equal reports are exactly null signed differences. Positivity is nonunique and rank one supplies no integer complex structure. & GEN-WDP-60--61 \\
\textbf{R9. Finite Hermitianization} (§17.9; derivation supplement) & The realized observer carrier is rationalized; a rational-linear map \(J\) is supplied with \(J^2=-I\) and is skew-adjoint for the observer Gram form. & \((a+ib)x=ax+bJx\) makes the carrier a \(\mathbb Q(i)\)-module, and \(h(x,y)=g(x,y)+i g(Jx,y)\) is Hermitian, positive semidefinite and has exactly the real Gram radical. The same positive form admits opposite complex orientations, so this constructs but does not select the complex branch. & GEN-WDP-62--63 \\
\textbf{R10. Certified state/effect readout} (§17.10; derivation supplement) & A nonzero Gaussian-rational Gram factor presents the state; complete effect factors satisfy \(\sum_o C_o^\dagger C_o=I\); the trace pairing is explicitly admitted as the readout. & Frobenius-square certificates give exact nonnegative normalized rational probabilities and replay the trace formula. A depolarized affine family remains positive, normalized and phase invariant while differing on a sharp input, so these conditions do not select the undeformed Born readout. & GEN-WDP-64--65 \\
\textbf{R11. Certified process composition} (§17.11; derivation supplement) & Gram-factor states and a finite Gaussian-rational operator family satisfying \(\sum_oK_o^\dagger K_o=I\) are supplied; sequential or tensor carrier structure is explicitly admitted. & Branch factors give exact nonnegative probabilities whose masses sum to the input mass; sequential and tensor families compose and supplied unitaries are reversible. Equal channels can retain different instruments. Conventional Hadamard/T matrices expose a scalar boundary; R12 realizes the Hadamard ray action over \(\mathbb Q(i)\) while retaining the T obstruction. No gate model capable of arbitrary quantum computation is selected. & GEN-WDP-66--68 \\
\textbf{R12. From separation to transformation} (Chapter 19A) & The R9 carrier has positive definite Hermitian form; R10's pure-state readout is supplied. A bijection of the entire ray space preserves orthogonality in both directions, with \(d\ge3\), or all transition probabilities when \(d=2\). & The map has a projectively unique semilinear similarity representative. Rational phases form an infinite group. Unitarily reachable rays are classified by the Gaussian norm class, while some even-dimensional similarities cross those classes. These standard algebraic ingredients give exact, original-task-bound certificates and sharpen the limits of report preservation. & WDP-R12-O01; WDP-R12-P02--P05; WDP-R12-O02 \\
\textbf{R13. Algebraic generators} (§19A.6) & A \(K\)-linear derivation of \(M_d(K)\) is supplied and obeys the declared adjoint law. & An explicit matrix-unit construction gives \(\delta=-i[H_0,\cdot]\) with \(H_0^\dagger=H_0\), unique modulo a rational scalar. The inner-derivation theorem is standard; the finite certificate checks its actual input table. No time evolution or physical Hamiltonian is selected. & WDP-R13-P01 \\
\textbf{R14. Question-preserving encoding} (§10.8; encoding supplement) & A declared task and protected question, a question-faithful partial parser, a sound checker, and binding of the report to the claimed candidate. & Equal reports with unequal answers exclude a complete report-only decision. Parser fidelity transports checker acceptance to the original question. An exhaustive integer-quadratic experiment compares cochain reconstruction with both enumeration and direct reconstruction; it establishes no cochain-specific efficiency gain. & GEN-WDP-74; QPE-1--QPE-4 \\
\textbf{R15. Sharp linear observation budgets} (Appendix T) & Finite rational spaces, realizable image \(R\), repair image \(B\), existing kernel \(K\), and unrestricted added rational linear rows. & WDP-DIM-1 gives sharp raw-on-\(R\), quotient-on-\(V\), and joint counts. Restricting the repair-class question to actual reports of candidates in \(R\) gives \(\dim(R\cap K)-\dim(R\cap K\cap B)\) rows. Fixed menus and noisy reports require their own models. & WDP-DIM-1; WDP-DIM-R; GEN-WDP-75 \\
\end{longtable}

Section 10.6 connects R1's repair classes to R2's report fibers through
sharp observer completion. The raw-defect question (is a locally accepted
configuration actually defective?) and the quotient question (does it lie in
a nonzero repair class?) have different invisible subspaces and different
measurement budgets. Checked repair retains the resulting certificates as
evaluation evidence, with authorization as a separate operation.

Three forms of completion must be reported separately. A \textbf{written theorem}
is complete when its finite or conditional proof, assumptions and
countermodels are inspectable here. A \textbf{mechanized theorem} is complete
only for the exact Lean statement and its audited axioms; an abstract
model declaration does not verify the source program. \textbf{Concrete
implementation verification} binds the actual Rust/Python inputs,
transition semantics, checker and tested scope; a finite replay does not
prove every run. Appendix D and the claim ledger record those evidence
levels without interchanging them.

A finished \emph{finite edition} can be defined by a declared collection of
carriers, questions and theorems such as R1--R6, with every assertion in
that collection proved or explicitly refused at its stated scope. The
broader GEN-WDP-O01--O05 programme, physical reconstruction, continuum
limits and higher coherence remain separately open until their own
hypotheses and proofs are supplied. The edition boundary does not change
their status or turn a checked finite example into an unrestricted theorem.

\subsubsection{Publication Scope Boundaries and Open Reconstruction Obligations}\label{publication-scope-boundaries-and-open-reconstruction-obligations}

To maintain rigorous claim hygiene and prevent over-extension of formal statements, the following explicit scope boundaries distinguish proven finite conditional theorems from open global reconstruction goals:

\begin{enumerate}
\def\labelenumi{\arabic{enumi}.}
\tightlist
\item
  \textbf{Attainment and rational coefficients (§4.2, §10.2)}: Finite-dimensional real or complex normed repair spaces have compact bounded sublevel sets after the repair kernel is removed. Completeness of an arbitrary scalar field does not supply this compactness. General norms, including polyhedral norms with irrational coefficients, need not attain their infimum on rational repairs. A feasible bounded linear programme with rational data has a rational optimum; matching primal and dual certificates verify a proposed optimum, whereas primal feasibility alone does not.
\item
  \textbf{Graded Subspace Invariance (§16.1)}: Anticommutator zero-locus arguments require an explicitly declared graded domain preserved by \(d_s\) and \(d_h\); general un-graded endomaps do not automatically preserve zero loci.
\item
  \textbf{Specialization Inputs for W1--W6 (§18.1)}: Structural specializations W1--W6 represent admitted conditional inputs (reversible updates, Gaussian-integer linearizations, pairing/quadratic readouts) rather than derivations from primitive distinction alone.
\item
  \textbf{PSD Geometry vs.~LP Polyhedra (§18.5)}: General Positive Semidefinite (PSD) membership must be kept distinct from finite linear programming (LP) polyhedra. For instance, the qubit state slice \(\frac{1}{2}\begin{pmatrix} 1+z & x \\ x & 1-z \end{pmatrix}\) is PSD if and only if \(x^2 + z^2 \le 1\), defining a non-polyhedral disk that requires exact Hermitian/PSD certificates rather than finite LP polyhedra.
\end{enumerate}

\begin{center}\rule{0.5\linewidth}{0.5pt}\end{center}

\subsection{Contents}\label{contents}

\textbf{Part I --- Difference, witness, and finite obstruction:}

\begin{itemize}
\tightlist
\item
  Chapter 1. The admitted cut
\item
  Chapter 2. Residuals and the typed witness taxonomy
\item
  Chapter 3. Symmetry, nonselection, and the price of retention
\item
  Chapter 4. Finite cochains, quotients, and certificates
\item
  Chapter 5. Optional quantum-capable carriers
\end{itemize}

\textbf{Part II --- Order, quantity, and finite probability:}

\begin{itemize}
\tightlist
\item
  Chapter 6. Recurrence, succession, and constitutive order
\item
  Chapter 7. Duration, exactness, periods, and clocks
\item
  Chapter 8. Multiplicity, additive weight, and finite probability
\item
  Chapter 9. Probability scope fences
\end{itemize}

\textbf{Part III --- Observer, observation, and measurement:}

\begin{itemize}
\tightlist
\item
  Chapter 10. A layered observer architecture
\item
  Chapter 11. Intersubjective presentation, descent, and redundancy
\item
  Chapter 12. The observation ladder
\item
  Chapter 13. Measurement as a typed extension of observation
\item
  Chapter 14. Self-reference and the awareness boundary
\item
  Chapter 15. Authority is a separate layer
\end{itemize}

\textbf{Part IV --- Histories and the quantum boundary:}

\begin{itemize}
\tightlist
\item
  Chapter 16. Spatial/history diagnostics and the curved-complex fork
\item
  Chapter 17. Finite weighted histories
\item
  Chapter 18. Quantum reconstruction as an audit
\item
  Chapter 19. The unpaid quantum interfaces
\item
  Chapter 19A. From certified observation to admissible transformation
\end{itemize}

\textbf{Part V --- Relative scope and the open programme:}

\begin{itemize}
\tightlist
\item
  Chapter 20. Relative completeness for observer-indexed presentations
\item
  Chapter 21. Falsifiers, nonclaims, and open obligations
\item
  Chapter 22. Synthesis: what the finite calculus establishes
\item
  Chapter 23. Observer completion on a finite 2D toric-code carrier
\end{itemize}

\textbf{Appendices:} A dependency and status table; B end-to-end example; C source
integration map; D formalization status and evidence architecture; E glossary;
F bibliography and citation policy; G finite decision procedures; H scoped completion results;
I native integration; J expert derivations and native ownership; K
source-bound observer adaptation; L separation of the archived source
derivation from WDP; M Meron characters and finite shadows; N source form
specifications; O presentation closure from finite shadows to completion; P the
earlier component calculi as six reusable constructions, with five-edit model
balance-edit, path, result-mark and guarded-history results; Q transfer of
learned rejection rules; R a plain-language picture of observers and
authority; S certified progress; T sharp completion with distinct images; U
corrected toric instrument mathematics; V locality proofs; W Trotter proofs;
X calibrated reports and observation cost; Y recovery costs from joint report
supports.

\begin{center}\rule{0.5\linewidth}{0.5pt}\end{center}

\subsection{How to read this monograph}\label{how-to-read-this-monograph}

There are three useful routes through the text.

\begin{enumerate}
\def\labelenumi{\arabic{enumi}.}
\tightlist
\item
  A reader interested in the foundational chain can read Chapters 1--3, 6--9,
  10--15, and 20--23. This route foregrounds what is admitted, what is
  constructed, and what is refused.
\item
  A reader interested in the finite mathematics can follow the central chain:
  \textbf{report preservation} (the entrance above and Appendix P.1) → \textbf{repair}
  (Chapter 4) → \textbf{sufficient observation} (§10.6 and Appendix T) →
  \textbf{observation cost} (§§14.6--14.7 and Appendices X--Y) → \textbf{lawful histories}
  (Appendices P.13--P.23). P.22 gives the guarded-schedule topology and
  coherence thresholds; P.23 gives integral chain transport and certificate
  bounds. Chapters 3--8, 16--19 and 19A supply further algebraic and quantum
  constructions with their own admitted inputs.
\item
  A reader interested in observers and measurement can begin with the typed
  glossary below, then read Chapters 10--15, Appendix T and Appendices X--Y.
  The shared-error example in Y.6 shows why retaining correlations can change
  recovery from obstructed to attainable. This route is independent of the
  optional quantum-carrier branch.
\end{enumerate}

No reading route should treat chapter order as a metaphysical order. The order
is a dependency order for the constructions presented here. A later formalism
may admit a different starting package and still be mathematically coherent.

Appendix L follows the actual archived source derivation into WDP. Appendix M
follows the Meron and finite-shadow branch toward continuum models, with
separate proofs for cancellation, observable preservation and conditional
numerical limits. Appendix N follows the source form specifications and proposed
SQL profile;
Appendix O follows metric presentation closure and actual Meron source limits.
The spatial proof distinguishes signed finite generating functions from the
constructed Gaussian random-distribution target.

\subsubsection{Proof supplements and stable locations}\label{proof-supplements-and-stable-locations}

Several full proofs extend beyond the summaries printed in this PDF. They are
part of the accompanying source archive, under these relative paths:

\begin{itemize}
\tightlist
\item
  \texttt{docs/guarded\_schedule\_topology.md}: Appendix P.22, the native event complex,
  common-cover argument and its relationship to arrangement complements.
\item
  \texttt{docs/sparse\_guarded\_chain\_transport.md}: Appendix P.23, the minimum-sweep
  construction and fixed-degree computation and certificate-size bounds.
\item
  \texttt{docs/derivations/README.md}: index of the seven full proof supplements
  for §§17.7--17.11 and Chapter 19A, including additive/tensor completion,
  Hermitianization, certified readout, process composition and preservers.
\item
  \texttt{docs/derivation\_review.md} and the supplements linked from Appendix J:
  detailed admission and transport arguments and their evidence boundaries.
\end{itemize}

The archive preserves these paths; \texttt{MANIFEST.sha256} binds their exact bytes.
The repository snapshot
provides an immutable location for the retained supplements. The present
correction has no assigned version DOI; the older v0.5.0 DOI identifies that
older edition. Appendix F includes the external bibliography in this PDF and
preserves the citation keys used in the text.

\subsubsection{Status grammar}\label{status-grammar}

Every major assertion belongs to one of the following classes.

\begin{itemize}
\tightlist
\item
  \textbf{Definition.} A stipulated mathematical object or predicate. A definition
  may be useful or useless, but is not a theorem just because it is precise.
\item
  \textbf{Finite proposition.} A result proved in the text from stated finite
  hypotheses. A proof sketch is included when the argument is short enough to
  audit locally.
\item
  \textbf{Typing proposition.} A claim about what the declared nominal interfaces do
  or do not coerce. It establishes a type-system fence, not every semantic
  nonimplication between possible interpretations of the types.
\item
  \textbf{Conditional result.} A proposition whose conclusion is proved only after
  named imports or regularity assumptions are supplied.
\item
  \textbf{Cited result.} A result attributed to an external source and used
  only within its declared scope. The primary-source registry records the
  intended source and consultation state.
\item
  \textbf{Mechanized elsewhere.} A repository or external theorem is reported as
  machine checked outside this carrier. That report is not promoted here unless
  the exact source, statement, toolchain, and source reference are available.
\item
  \textbf{Open claim.} A well-typed target without a completed proof.
\item
  \textbf{Open interface.} An open claim whose missing object is a typed adapter or
  interface together with its preservation laws.
\item
  \textbf{Conjecture.} An open claim for which examples or analogy suggest a
  direction.
\item
  \textbf{Refuted draft.} A previously printed implication or classification for
  which a counterexample or typing failure is known.
\item
  \textbf{Interpretive proposal.} A reading of the formalism, not a mathematical
  consequence.
\item
  \textbf{Nonclaim.} A boundary the monograph expressly refuses to cross.
\item
  \textbf{Synthesis conclusion.} A bounded summary of results already established in
  named rows; it introduces no additional theorem.
\end{itemize}

These labels describe \textbf{mathematical status}. The claim ledger records two
additional, independent fields: \textbf{formalization scope} (an exact declaration,
a narrower partial declaration, or no bound declaration) and
\textbf{execution/reproduction status} (a dated paired replay or a pending replay).
The strict paired reproduction gate still requires a successful, source-bound
Lean elaboration and the bound Python check. Missing bindings, unavailable
dependencies, or a check that has not been rerun are reproduction limitations;
they do not make a written proof an open mathematical claim. A declaration
name alone also does not establish a successful current build.

The alphanumeric claim identifiers and implementation names used below are
repository lookup keys, not mathematical standards or additional assumptions.
The definitions and cited source statements, rather than those labels,
determine the content of each result.

Identifiers of the form \texttt{WDP-CHxx-*} are the canonical row identifiers shared
by this monograph and the unified claim ledger. The body supplies hypotheses and
proof context, while the ledger is the authority for final evidence bindings
and replay status.

\subsubsection{Standing scope discipline}\label{standing-scope-discipline}

The base combinatorial and operational families are finite: token and event
sets, graphs, covers, path families at a fixed bound, query families, and
comparison tables. Coefficient and value spaces may be infinite when explicitly
typed---for example \(\mathbb R\), \(\mathbb C\), real intervals, or a
finite-dimensional real vector space. Finite cardinality and finite dimension
are never interchangeable. The finite indexing makes many inspections
exhaustive and avoids infinite path sums; it does not license continuum
conclusions. ``Canonical,'' ``minimal,'' ``unrestricted,'' ``forced,'' ``complete,'' and
``equivalent'' are used only with a stated category, order, optimization problem,
or hypothesis class.

\texttt{WDP-CH00-N01} \textbf{Nonclaim.} The monograph does not derive all structure from
Primitive Difference.

\texttt{WDP-CH00-N02} \textbf{Nonclaim.} The monograph does not identify formal witness with
experience, consciousness, agency, or moral status.

\texttt{WDP-CH00-N03} \textbf{Nonclaim.} No local proof, test run, dashboard, archive, or
status label confers authority.

\begin{center}\rule{0.5\linewidth}{0.5pt}\end{center}

\subsection{Part I --- Difference, witness, and finite obstruction}\label{part-i-difference-witness-and-finite-obstruction}

\subsection{Chapter 1 --- The admitted cut}\label{chapter-1-the-admitted-cut}

\subsubsection{1.1 The starting package}\label{the-starting-package}

The programme begins not with a bare slogan but with an admitted finite
package. Let \(X\) be a finite carrier of exhibited poles or tokens. A \textbf{cut} is
a declared occurrence of apartness between roles in an exhibition. At minimum
we require:

\begin{enumerate}
\def\labelenumi{\arabic{enumi}.}
\tightlist
\item
  a way to distinguish the relevant roles in the present exhibition;
\item
  a structural carrier on which those roles can be compared;
\item
  a rule saying which comparisons count as sameness or repair for the present
  question.
\end{enumerate}

We call the first ingredient \textbf{Primitive Difference} and the remaining
organization \textbf{Structure}. They are co-inputs. The phrase ``from difference'' is
therefore shorthand for ``from an exhibited difference inside a declared finite
structural package.'' We call the resulting finite-carrier setting the
\textbf{finite obstruction calculus}: cochains, admissible repairs, quotient defects
and their certificates over a declared finite carrier.

\texttt{WDP-CH01-D01} \textbf{Definition (admitted cut).} An admitted cut is a tuple
\(\mathfrak C=(X,E,\tau,\mathcal K)\) where \(X\) is a finite set of exhibited
tokens, \(E\) is a finite set of incidences or comparisons, \(\tau\) assigns the
roles needed to type those incidences, and \(\mathcal K\) is the declared family
of admissible comparisons. The tuple is data. No part of it is claimed to be
selected from equality on \(X\) alone.

In the admission-versus-evaluation vocabulary of the earlier component papers
(Root W vs.~Root P), a
\textbf{Root W realization} additionally supplies a relation and its protected
irreflexivity law: \(K = (X, R, \operatorname{Irr}(R))\), with witness
\(d = (a, b, w : R(a,b))\). The general admitted-cut tuple does not itself
entail these extra inputs. Once supplied, they type relational distinction;
orientation still depends on the chosen relation, as in the symmetric and
oriented examples of §1.4. Root P's cochain evaluation and repair quotient
are further structures, not consequences of these initial assumptions.

The exact representation can vary. A graph is often convenient: vertices are
tokens and oriented edges are typed comparisons. A hypergraph, relation, or
small category may be more appropriate when comparisons have higher arity.
What matters is that the representation be named before a theorem depends on
it.

\subsubsection{1.2 Difference and local sameness}\label{difference-and-local-sameness}

Within a fixed comparison family, one may ask whether two presentations agree.
Let \(P\) be a set of presentations and let

\[
\operatorname{Same}_{\mathcal K}:P\times P\longrightarrow\{0,1\}
\]

be a declared predicate. A failed comparison is a residual at the level of
predicates. If the presentations live in an abelian group or vector space, the
failure may instead be encoded by a difference \(p-q\). The algebraic encoding is
stronger: subtraction, scalar structure, and a zero element have already been
admitted.

\texttt{WDP-CH01-P01} \textbf{Finite proposition.} A failed declared sameness comparison
exhibits that the two compared presentations are not interchangeable under that
comparison.

This statement is intentionally local. It does not say the presentations are
different under every possible comparison. It does not make the failed
comparison intrinsic to the tokens. It does not supply an ordering of them.

\subsubsection{1.3 Why Structure cannot be erased}\label{why-structure-cannot-be-erased}

Suppose only a finite set \(X\) with equality is given. For \(|X|\ge 2\), every
permutation of \(X\) is an automorphism. Any construction advertised as
parameter-free and natural under all automorphisms must commute with the full
symmetric group. A distinguished element, a direction, an ordered pair of distinct elements, or a nontrivial
proper subset is generally moved by some permutation.
Therefore it cannot be selected invariantly without additional data.

This elementary observation is the first discipline of the programme. It does
not make structure illegitimate. It requires structure to be visible in the
input ledger.

\subsubsection{1.4 Worked example: a two-pole exhibition}\label{worked-example-a-two-pole-exhibition}

Let \(X=\{a,b\}\). Declaring just \(a\ne b\) exhibits apartness but does not select
which pole is source and which is target. The transposition exchanging \(a\) and
\(b\) preserves equality and inequality. If we additionally declare an oriented
edge \(e:a\to b\), then orientation exists because it was supplied by \(e\).

Three conclusions must remain distinct:

\begin{itemize}
\tightlist
\item
  the poles are apart in the exhibition;
\item
  the edge orients them;
\item
  an observer later may register that orientation.
\end{itemize}

The first does not imply the second, and neither implies the third.

The same separation holds one level down, for the relation that is supposed to
exhibit apartness. A pointed relation \((X,R,a,b)\) records a relation \(R\) on a
finite carrier together with two marked poles.

\texttt{WDP-CH01-P02} \textbf{Finite proposition (pointed relations and apartness).} For a
pointed relation \((X,R,a,b)\) on a finite carrier, \(R(a,b)\) does not imply
\(a\ne b\); \(a\ne b\) does not imply \(R(a,b)\); and \(R(a,b)\) together with
irreflexivity of \(R\) implies \(a\ne b\).

\emph{Proof.} A one-point carrier with \(R=\{(a,a)\}\) and \(b=a\) has \(R(a,b)\) and
\(a=b\). The empty relation on a two-point carrier has \(a\ne b\) without
\(R(a,b)\). If \(R\) is irreflexive and \(R(a,b)\), then \(a=b\) would give \(R(a,a)\).
\(\square\)

Irreflexivity is therefore the declared ingredient that lets a related pair
witness apartness; it belongs to the admitted package and does not follow from
relatedness. Symmetry is a second, independent choice. The symmetric two-point
relation \(\{(a,b),(b,a)\}\) admits the swap as an automorphism, and the oriented
relation \(\{(a,b)\}\) does not. Keeping both is nonselection in the sense of
Chapter 3; choosing one is an entry in the injection ledger of §3.2.

\subsubsection{1.5 Ground and spectral ground state}\label{ground-and-spectral-ground-state}

``Ground'' in this chapter means the admitted starting package. A spectral ground
state is instead a lowest-eigenvalue state of a declared operator on a declared
inner-product space. The latter can be studied in the optional carrier branch,
but it is not the foundation of the former and should not share an unqualified
name.

\texttt{WDP-CH01-N01} \textbf{Nonclaim.} No spectral, physical-vacuum, or cosmological ground
state is derived from an admitted cut.

\subsubsection{1.6 The Priced Grounding Span: From Raw Mention to Witnessed Verification}\label{the-priced-grounding-span-from-raw-mention-to-witnessed-verification}

\texttt{WDP-CH01-D02} \textbf{Definition (Priced Observer Boundary State).}\\
An observer boundary specification is an admitted triple \((X_O,Y_O,\pi_O)\).
For this profile take \(X_O=(S_O,R,\mathrm{TL}_O^{\mathrm{priv}})\), with
\(R\in\mathbb{Q}_{\ge0}\) and \(\mathrm{TL}_O^{\mathrm{priv}}\in\mathbb N\).
Choose an external-state report \(\pi_{\mathrm{ext}}\), a threshold \(c_{\min}\)
and a positive integer \(k\). The public report is
\[
\pi_O(S_O,R,t)=\bigl(\pi_{\mathrm{ext}}(S_O),
\mathbf1_{\{R\ge c_{\min}\}},\lfloor t/k\rfloor\bigr).
\]
The Boolean flag may be labeled SOLVENT or EXHAUSTED. Fibers are inverse images
\(\pi_O^{-1}(y)\), not linear kernels: this mixed discrete/nonlinear carrier has
no declared vector-space dimension. Coarsening a balance can hide its exact
value at one instant. Persistent privacy requires equal reports under all
admitted continuations, as in §20.14. Keeping exact balances in a separate host
ledger is an architectural admission, not a consequence of noninjectivity.

\texttt{WDP-CH01-D03} \textbf{Definition (Five-Gate Grounding Span).}\\
The grounding span for an ingress candidate is the sequence of five gated transitions:
\[
\begin{aligned}
\text{External Supply}&\xrightarrow{g_{\mathrm{sup}}}\text{Token-Erased Admission},\\
\text{Token-Erased Admission}&\xrightarrow{g_{\mathrm{adm}}}\text{Asymmetry-Priced Retention},\\
\text{Asymmetry-Priced Retention}&\xrightarrow{g_{\mathrm{ret}}}\text{Quotient-Defect Evaluation},\\
\text{Quotient-Defect Evaluation}&\xrightarrow{g_{\mathrm{ev}}}\text{External Sovereign Receipt},\\
\text{External Sovereign Receipt}&\xrightarrow{g_{\mathrm{auth}}}\text{Authorized transition}.
\end{aligned}
\]
governed by explicit operational preconditions:
1. \textbf{External Supply (\(g_{\mathrm{sup}}\)):} Ingress of external raw tokens. Operates as an identity transition on internal regime state (\texttt{supply\_no\_op}).
2. \textbf{Token-Erased Admission (\(g_{\mathrm{adm}}\)):} Admitted if and only if form \(f < |forms|\), argument arity matches, and argument data exist (\(x < data\)). Erases token spellings in favor of ordinal indices (\(NS\text{-}14\)), debiting ingress cost \(c_{\mathrm{adm}}\) from ledger \(R\). If \(R < c_{\mathrm{adm}}\), transitions to refusal \texttt{REFUSE\_INSUFFICIENT\_FUNDS}.
3. \textbf{Asymmetry-Priced Retention (\(g_{\mathrm{ret}}\)):} Admits a retention policy across \(\Delta t \le \mathrm{TL}_O\) and a separately chosen fee \(c_{\mathrm{ret}}\Delta t\). \texttt{WDP-CH03-P03} concerns the selected two-item no-form grammar; it does not derive this temporal fee or physical retention law.
4. \textbf{Quotient-Defect Evaluation (\(g_{\mathrm{ev}}\)):} Evaluates discrepancy \(\omega \in C^1\) via quotient distance \(\Phi_1(\omega) = \inf_{\alpha \in C^0} \|\omega - d_0 \alpha\|_1\).
5. \textbf{External Sovereign Receipt (\(g_{\mathrm{auth}}\)):} Emits machine receipt \(\mathrm{Receipt}(\Phi_1, \varphi^*)\). Carries diagnostic evaluation only; confers zero sovereign standing (\texttt{ENC-0195}, \texttt{WDP-CH15-D01}).

\texttt{WDP-CH01-L01} \textbf{Lemma (Token-Erased Admission Obstruction).}\\
Let an unhosted syntactic generator operate with empty boundary presentation (\(B_O = \varnothing\), \(S_O = \varnothing\)). Then gate \(g_{\mathrm{adm}}\) cannot evaluate boundary consistency, returning typed refusal \texttt{REFUSE\_EMPTY\_BOUNDARY}.

\texttt{WDP-CH01-N02} \textbf{Nonclaim.}\\
The five-gate grounding span defines the formal operational requirements of the declared mathematical model; it is not asserted as an empirical model of human cognitive architecture.

\begin{center}\rule{0.5\linewidth}{0.5pt}\end{center}

\subsection{Chapter 2 --- Residuals and the typed witness taxonomy}\label{chapter-2-residuals-and-the-typed-witness-taxonomy}

\subsubsection{2.1 Residual before witness}\label{residual-before-witness}

A \textbf{residual} is the part of a declared target that is not removed by an
admissible repair. Let \(V\) be a finite-dimensional real vector space of defects,
\(A\) a finite-dimensional space of repairs, \(d:A\to V\) a linear repair map, and
\(\delta\in V\) an exhibited defect. For a declared seminorm or cost functional
\(J:V\to[0,\infty)\), define

\[
\Phi_J([\delta])=\inf_{\alpha\in A}J(\delta-d\alpha).
\]

Here \([\delta]\) denotes the coset of \(\delta\) modulo \(\operatorname{im}d\) when
that quotient is meaningful. The expression is not primitive: it imports a
linear carrier, a repair map, and a cost.

In this lexicon, the exhibited defect \(\delta\),
the repair kernel \(\ker d\), and the quotient defect \(\Phi_J([\delta])\) are
classified strictly as \textbf{Root P}: the evaluation kernel and quotient defect
structure on cochains. Root P computes obstruction to global exactness and
repair over an admitted complex. It presupposes Root W as underlying carrier and
never substitutes for it.

\texttt{WDP-CH02-D01} \textbf{Definition (residual locus witness).} A residual-locus witness
is a typed locus \(L\) together with evidence that the restriction or component of
the residual at \(L\) is nonzero according to the declared test.

The locus is a witness \emph{to a proposition}. It is not thereby a memory, an
observer, an agent, or a measurement apparatus.

\subsubsection{2.2 Seven witness roles}\label{seven-witness-roles}

The single word \emph{witness} is insufficient. We use the following typed roles.

\begin{enumerate}
\def\labelenumi{\arabic{enumi}.}
\tightlist
\item
  \texttt{ResidualLocusWitness{[}P{]}}: a locus at which residual proposition \(P\) is
  exhibited.
\item
  \texttt{RetentionWitness{[}P{]}}: a mark or state component enabling later re-access to
  evidence for \(P\).
\item
  \texttt{RepairWitness{[}P{]}}: an admissible \(\alpha\) exhibiting that \(P\) can be
  repaired, often by satisfying \(d\alpha=\delta\).
\item
  \texttt{ContextWitness{[}P{]}}: a local assignment or section witnessing \(P\) inside a
  declared context.
\item
  \texttt{DualCertificateWitness{[}P{]}}: a functional separating a defect from the repair
  image.
\item
  \texttt{DescentReceipt{[}P{]}}: evidence that compatible local presentations glue or
  transport according to a declared descent rule.
\item
  \texttt{CounterexampleWitness{[}P{]}}: a finite inhabitant refuting a universally
  quantified \(P\).
\end{enumerate}

The general notation

\[
\operatorname{Witness}_{k}[P](x,e)
\]

means that evidence \(e\) makes subject \(x\) a witness of kind \(k\) for proposition
\(P\). The proposition and kind are part of the type.

\texttt{WDP-CH02-P01} \textbf{Typing proposition.} In the nominal typed interface, no one of
the seven role tags coerces to another unless an explicit adapter and its
hypotheses are declared.

This is a syntactic fence, not a proof of all 42 directed semantic
nonimplications. Simple models do separate several important cases: a nonzero
residual component need not be stored; a stored mark need not repair anything;
a repair element need not define a separating functional; a local section need
not glue; and a counterexample need not be re-accessible.

\subsubsection{2.3 A triangle example}\label{a-triangle-example}

Let \(G\) be an oriented triangle with vertices \(0,1,2\) and edges \(01,12,20\).
Let \(C^0(G;\mathbb R)=\mathbb R^3\) and
\(C^1(G;\mathbb R)=\mathbb R^3\). The coboundary is

\[
(d_0 f)_{ij}=f_j-f_i.
\]

Take \(\delta=(1,1,1)\) in the edge order \((01,12,20)\). Its circulation is
\(1+1+1=3\). Every coboundary has zero circulation because the vertex terms
cancel. Thus the cycle is a \texttt{DualCertificateWitness} against exact repair: the
linear functional \(\ell(x)=x_{01}+x_{12}+x_{20}\) vanishes on
\(\operatorname{im}d_0\) and satisfies \(\ell(\delta)=3\).

Each edge is also a possible residual locus if the local test reads its nonzero
component. But no edge stores the result just by being nonzero, and the
functional \(\ell\) is not an observer.

\subsubsection{2.4 Witness and evidence}\label{witness-and-evidence}

Evidence is role-relative. A proof term can witness a proposition in a formal
system. A data record can witness that a test returned a value. A residual can
witness failure of a repair condition. These are not interchangeable. Whenever
a chapter moves between them, it must provide an adapter---for example, a theorem
connecting an executable output to a mathematical predicate, or a provenance
verification record connecting a stored artifact to the source that generated it.

\texttt{WDP-CH02-N01} \textbf{Nonclaim.} A typed witness does not acquire epistemic,
phenomenal, legal, or institutional authority just by being called evidence.

\begin{center}\rule{0.5\linewidth}{0.5pt}\end{center}

\subsection{Chapter 3 --- Symmetry, nonselection, and the price of retention}\label{chapter-3-symmetry-nonselection-and-the-price-of-retention}

\subsubsection{3.1 The finite symmetry obstruction}\label{the-finite-symmetry-obstruction}

Let \(X\) be a finite set, \(|X|=n\ge2\), carrying no structure beyond equality.
The group \(S_X\) acts transitively on \(X\). A parameter-free choice of a point
would be an \(S_X\)-equivariant map from the one-point trivial \(S_X\)-set to \(X\).
Such a map selects \(x\in X\) fixed by every permutation, but no point is fixed by
all permutations when \(n\ge2\).

\texttt{WDP-CH03-P01} \textbf{Finite proposition (no equivariant point selection).} There is
no \(S_X\)-equivariant selection of a distinguished element of an unstructured
finite set \(X\) with \(|X|\ge2\).

The same proof applies to an oriented edge or nontrivial proper subset whenever
the full symmetry group acts without a global fixed target of the requested
type.

This is a nonselection theorem, not an impossibility of choice. A label, base
point, order, random seed, boundary condition, prior state, or external input
can select. The theorem requires the asymmetry source to be named.

\subsubsection{3.2 Injection ledger}\label{injection-ledger}

For any proposed selection, record the first input that breaks the relevant
symmetry. Typical entries are:

\begin{itemize}
\tightlist
\item
  a named element or base point;
\item
  an orientation;
\item
  an ordering or clock;
\item
  a measure or weighting;
\item
  a boundary condition;
\item
  a representation or coordinate system;
\item
  an observer profile;
\item
  a stochastic device and its distribution;
\item
  a physical law not present in the root package.
\end{itemize}

\texttt{WDP-CH03-D01} \textbf{Definition (asymmetry-price frontier).} Relative to a specified
implementation class and dependency preorder, the asymmetry-price frontier is
the family of inclusion-minimal declared input sets whose removal restores a
symmetry that moves the output. A scalar ``price'' additionally requires a stated
cost function; several frontier elements may tie or remain incomparable.

Minimality here is relative to the declared implementation class and order. No
uniqueness or global-minimality theorem is claimed for the frontier in general;
CH03-P03 below settles one declared instance, and a tie survives even there.

A small search makes the point concrete. A \textbf{declare-then-assert grammar} has a
finite set \(D\) of data, a list of declared forms with positive arities, a set of
asserted instances (a form with a tuple of data of its arity), and a set of forms
carrying a rule that asserts every tuple of that form. A \textbf{context} is a declared
form whose slots hold admitted data or holes, at least one of them a hole; a datum
\(x\) satisfies the context when filling every hole with \(x\) gives an asserted
instance or a tuple of a ruled form. Two data are \textbf{separated} when some context
is satisfied by exactly one of them. A \textbf{response} is a list of items, each of
which adds a datum, declares a form of positive arity, asserts one instance of an
already declared form on admitted data, or admits the rule of an already declared
form. An isomorphism of results is a pair of bijections, on data and on forms,
preserving arities, instances, rules, and the required pair as an unordered pair.
Fix a base with \(a\ne b\) in \(D\), no forms, no instances, no rules, and the single
required pair \(\{a,b\}\). A response is valid when every item applies and the
result separates \(a\) from \(b\). The first cost is the item count; the second
refines it lexicographically by the sum of the arities of the newly declared
forms.

\texttt{WDP-CH03-P03} \textbf{Finite proposition (separation price and the surviving tie).}
Over this base:

\begin{enumerate}
\def\labelenumi{\arabic{enumi}.}
\tightlist
\item
  no response of at most one item is valid, and declaring a unary form \(F\) and
  asserting \(F(a)\) is valid;
\item
  for each \(k\ge1\), declaring a \(k\)-ary form and asserting it on
  \((a,\dots,a)\) is valid, and distinct \(k\) give non-isomorphic results, so the
  item count leaves infinitely many isomorphism classes of minimal responses;
\item
  under the lexicographic refinement, up to renaming the newly declared unary
  form, exactly two marked responses remain: the unary marks on \(a\) and on
  \(b\). They form one isomorphism class, and the
  transposition \((a\,b)\), an automorphism of the base, exchanges them, so no
  \((a\,b)\)-equivariant rule selects one of them.
\end{enumerate}

\emph{Proof.} (1) The base declares no form, so it has no context, and an instance
or rule item cannot apply to it. A single new datum adds no context. A single
new form adds contexts but asserts nothing, so every context fails at both \(a\)
and \(b\). For the unary mark, the context \(F(\_)\) holds at \(a\) and fails at \(b\).
(2) The all-hole context \(F(\_,\dots,\_)\) holds at \(a\) and fails at \(b\). The
result has exactly one form, and an isomorphism preserves its arity. (3) By (1)
a minimal response has two items. Its first item declares a form \(F\): an
instance or rule cannot apply to the base, and after a first datum the single
remaining item cannot both declare and assert. Its second item is neither a
datum nor a second form, which assert nothing, nor the rule of \(F\), which makes
every context hold at both \(a\) and \(b\). It therefore asserts one instance of
\(F\), and the arity sum equals the arity of \(F\), whose least value is \(1\). A
unary form has the single context \(F(\_)\), which holds exactly at the asserted
datum \(x\), so separation forces \(x\in\{a,b\}\). The transposition \((a\,b)\) fixes
the base with its required pair and carries the mark on \(a\) to the mark on \(b\).
The two results are therefore isomorphic, and \(H=\{1,(a\,b)\}\) acts on the two
surviving responses without a fixed point. By CH03-P02 (stated in Chapter 5),
no \(H\)-equivariant choice exists. \(\square\)

\emph{Worked count.} For a binary form, an asserted pair \((x,y)\) separates \(a\) from
\(b\) exactly when \(x\) or \(y\) lies in \(\{a,b\}\): the contexts \(F(\_,\_)\),
\(F(c,\_)\), and \(F(\_,c)\) test the diagonal, the second slot, and the first slot.
Up to isomorphism the separating pairs are \((a,a)\) and \((a,b)\) and, when \(D\) has
a third datum \(c\), also \((a,c)\) and \((c,a)\). With at most two items and arity at
most two, the minimal responses therefore fall into \(3\) classes when
\(D=\{a,b\}\) and \(5\) when \(|D|\ge3\), and into \(1\) class under the refinement.
The bounded counts grow with the arity bound, as clause 2 requires.

The proposition is relative to this grammar, these two costs, and this
isomorphism relation. It asserts no minimality over other formal languages, and
the value two is a property of the declare-then-assert grammar. Choosing \(a\)
over \(b\) is a further declared input of the kind listed in the injection ledger,
not an output of the search. The absorbed archived working-paper source proved clause 1
for every base with no forms and no instances, the unary witnesses, clause 2 for
two data, and the uniqueness half of clause 3 for two and three data in Lean at
its recorded commit. Those modules are copies of an external derivation study,
which the atlas rules exclude. They were not ported, and this row takes no Lean
credit.

\subsubsection{3.3 Retention is additional structure}\label{retention-is-additional-structure}

A difference at one exhibition does not by itself remain available in another.
Retention requires a carrier that persists, an update rule, and a criterion of
re-identification. A generic retention packet is

\[
\mathcal R=(M,u,r,\approx),
\]

where \(M\) is a memory carrier, \(u:M\times I\to M\) is an update under inputs \(I\),
\(r:M\to O\) is a report map, and \(\approx\) is the declared comparison used to
recognize retained content.

This packet is a definition, not a derivation. Earlier drafts sought an order-free minimal retention constructor and component
independence results. The declared finite behavioral classification below is
now proved; Appendix J also classifies typed component choices and separates additional regularity laws from typing.

\texttt{WDP-CH03-O01} \textbf{Finite proposition (resolved finite target).} For a finite nonempty carrier and finite input alphabet with total labelled deterministic updates and exact readouts or a validated report equivalence, descending partition refinement computes the greatest report-respecting update congruence, its minimal quotient, and behavioral classification by quotient isomorphism, without a global chronology.
The proof, native ownership and scope are given in Appendix I.

\texttt{WDP-CH03-O02} \textbf{Finite proposition (typed component independence).} For fixed finite nonempty carrier \(M\), input alphabet \(I\) and nonempty report alphabet \(O\), all typed retention extensions form the Cartesian product \(M^{I\times M}\times O^M\times\mathcal P(O\times O)\). Named compatibility laws may couple its factors. For arbitrary report comparison the all-word relation is the greatest report-respecting forward-compatible relation; quotient and literal-readout descent require their separate stated conditions. Appendix J establishes this result and the complete declared predicate classifications.

\subsubsection{3.4 Worked countermodel: witness without retention}\label{worked-countermodel-witness-without-retention}

At a single instant, let a detector bit be \(1\) exactly while an input contact is
closed. If the bit resets to \(0\) when contact opens and no other state persists,
then the bit is a \texttt{ResidualLocusWitness} for current contact but is not a
\texttt{RetentionWitness} for past contact. Adding a latch creates retention. The latch
is not contained in the original event; it is a new mechanism.

\subsubsection{3.5 Relation to automata}\label{relation-to-automata}

Myhill--Nerode equivalence is a standard downstream framework for quotienting
word histories by future distinguishability {[}REF-AUT-002{]}, but its inputs are
words and therefore already ordered. The native finite construction in Appendix I establishes its own minimal-quotient
property for the declared total update packet; no global chronology or free
selection of update structure is inferred. Bisimulation
likewise presupposes a labelled transition system {[}REF-AUT-005{]}. Channel theory
presupposes classifications and their connecting maps {[}REF-AUT-004{]}.
These frameworks clarify what retention can mean after their input structures
are admitted.

\subsubsection{3.6 Component independence: retention countermodels CM1--CM5}\label{component-independence-retention-countermodels-cm1cm5}

To probe the independence structure of \(\mathcal R=(M,u,r,\approx)\) under
Claim CH03-O02, and to give a candidate branching-time / non-total-order behavioral
relation for Claim CH03-O01, we give five concrete finite countermodels. They separate components in these instances. Appendix J supplies the full typed extension-space theorem and complete named predicate classifications; the finite behavioral classification is proved in Appendix I:

\begin{itemize}
\tightlist
\item
  \textbf{CM1 (Carrier Independence / Static Memory):} Let \(M = \{m_0, m_1\}\) with
  \(|M| \ge 2\), input alphabet \(I = \{\text{tick}\}\), identity update
  \(u(m, \text{tick}) = m\), injective report \(r(m_i) = r_i\), and diagonal
  comparison \(\approx = \{(r_0, r_0), (r_1, r_1)\}\). The memory carrier is
  non-trivial and observable, but updates perform no state transitions. Having
  a memory carrier \(M\) does not force active recording dynamics \(u\).
\item
  \textbf{CM2 (Update Independence / Readout Blindness):} Let \(M = \{m_0, m_1\}\),
  \(I = \{\text{toggle}\}\), update \(u(m_0, \text{toggle}) = m_1\) and
  \(u(m_1, \text{toggle}) = m_0\), but constant report \(r(m) = o_0\) for all \(m\).
  The internal state dynamics are invertible and non-trivial, but the readout
  is completely blind. Active updating \(u\) does not induce an informative
  report \(r\) without an explicit observability axiom.
\item
  \textbf{CM3 (Report Independence / Coarse-Graining):} Let \(M = \{m_0, m_1, m_2\}\),
  \(I = \{\text{cycle}\}\), cyclic update \(u(m_i, \text{cycle}) = m_{(i+1)\bmod 3}\),
  and non-injective report \(r(m_0) = r(m_2) = \text{even}\), \(r(m_1) = \text{odd}\).
  The dynamics are bijective, but report fibers merge states, showing that
  report equivalence \(\ker(r)\) is strictly coarser than state equality, so a
  single report cannot reconstruct \(M\). (The depth-1 relation below does separate
  \(m_0\) from \(m_2\): their successors \(m_1\) and \(m_0\) report differently.)
\item
  \textbf{CM4 (Comparison Independence / Non-Trivial Tolerance):} Let \(M = \{m_0, m_1\}\),
  \(I=\{\text{tick}\}\), \(u(m,\text{tick})=m\), \(O=M\), and
  \(r=\operatorname{id}_M\). Set \(\approx=M\times M\), a non-trivial tolerance
  that is also an equivalence relation. The states and reports are
  distinct, yet declared recognition identifies them. The comparison relation
  \(\approx\) is logically independent of \((M, u, r)\).
\item
  \textbf{CM5 (Lossy Update / Irreversible Reset):} Let \(M = \{m_0, m_1\}\),
  \(r = \operatorname{id}_M\), and \(I = \{\text{reset}, \text{keep}\}\) with
  \(u(m,\text{reset})=m_0\) and \(u(m,\text{keep})=m\) for all \(m\).
  The readout is exactly faithful, but
  the reset input collapses all history. Faithful observation of current state
  does not ensure retention of past interactions.
\end{itemize}

\textbf{Branching-Time / Non-Total-Order Behavioral Equivalence (CH03-O01):} Standard
automata equivalence uses words in \(I^*\) to compose admitted updates. Each word
has a local order; these words do not supply a global chronology. The same
behavior can be presented
by taking input events organized along branching computation trees of depth \(k\) (where
each branch carries local transition depth, but no global linear time axis is postulated
across branches): two states \(s \in M_1, s' \in M_2\) are \(k\)-equivalent if \(r_1(s) \approx r_2(s')\)
and for every declared input event \(i \in I\), the stepped states \(u_1(s, i)\) and
\(u_2(s', i)\) are \((k-1)\)-equivalent, with \(0\)-equivalence given by
\(r_1(s)\approx r_2(s')\) alone. This gives a bounded-depth behavioral relation on
finite branching structures without presupposing a global total order or linear
clock; each branch is still a finite sequence of events. It is an equivalence
relation when \(\approx\) is one, as it is in CM4. For a reflexive symmetric
comparison that is not transitive, the behavioral relation need not be
transitive: use three states with identity reports and updates, with
\(a\approx b\) and \(b\approx c\) but \(a\not\approx c\).

\texttt{WDP-CH03-N01} \textbf{Nonclaim.} The symmetry obstruction does not establish that
retention, dynamics, or observers are obstructed or unrealizable. It establishes that their
selection is not free.

\subsubsection{3.7 Presentation relabeling and invariant exports}\label{presentation-relabeling-and-invariant-exports}

\texttt{WDP-CH03-D02} \textbf{Definition (Presentation Relabeling and Invariant Passage Charge).}\\
Let \(\mathcal{P}_1\) and \(\mathcal{P}_2\) be presentations related by a declared
invertible alphabet map \(\sigma\). A protected export \(E\) is relabeling-invariant
when \(E(\sigma x)=E(x)\) on its declared domain. A passage cost, cochain action
and induced quotient action are additional data. Zero cost does not by itself
imply the identity action on \(Q^1=C^1/\operatorname{im}(d_0)\) or invariance of E.

\texttt{WDP-CH03-P04} \textbf{Proposition (limit of an invariant export).} If
\(E(\sigma x)=E(x)\) but a protected answer satisfies \(q(\sigma x)\ne q(x)\),
no decoder using only E recovers q on the whole domain. The proof is the
same-report contradiction in §14.8. It requires an actual target-separating pair;
a relabeling that preserves both E and q presents no such obstruction.

\texttt{WDP-CH03-N02} \textbf{Nonclaim.}\\
Proposition \texttt{WDP-CH03-P04} assumes invariance of a specified export. It does not
establish invariance of arbitrary internal verdicts or natural-language
tokenizers, and does not prohibit computing a defect from a supplied table.

\begin{center}\rule{0.5\linewidth}{0.5pt}\end{center}

\subsection{Chapter 4 --- Finite cochains, quotients, and certificates}\label{chapter-4-finite-cochains-quotients-and-certificates}

\subsubsection{4.1 The graph-cochain host}\label{the-graph-cochain-host}

Let \(G=(V,E)\) be a finite oriented graph. For a normed coefficient space \(A\),
define

\[
C^0(G;A)=A^V,\qquad C^1(G;A)=A^E,
\]

with coboundary

\[
(d_0f)(u,v)=f(v)-f(u).
\]

An edge field \(\delta\in C^1(G;A)\) is \textbf{exact} when
\(\delta=d_0f\). The quotient

\[
Q^1(G;A)=C^1(G;A)/\operatorname{im}d_0
\]

records edge fields modulo vertex-potential repair. It is a quotient relative
to the declared graph, coefficient space, and repair map. It is not the first
cohomology group unless a next differential \(d_1\) is supplied and
\(d_1d_0=0\).

\subsubsection{4.2 Quotient defect}\label{quotient-defect}

For a norm \(J\) on \(C^1(G;A)\), define

\[
\Phi_J([\delta])=\inf_{f\in C^0(G;A)}J(\delta-d_0f).
\]

\texttt{WDP-CH04-P01} \textbf{Finite proposition (well-defined quotient defect).}
\(\Phi_J\) is independent of the representative of \([\delta]\).

\emph{Proof.} If \(\delta'=\delta+d_0g\), then

\[
\inf_f J(\delta'-d_0f)
=\inf_fJ(\delta-d_0(f-g))
=\inf_hJ(\delta-d_0h).
\]

The translation \(h=f-g\) is a bijection on \(C^0(G;A)\). \(\square\)

\texttt{WDP-CH04-P02} \textbf{Conditional result (attainment).} If the repair
space is finite-dimensional over \(\mathbb{R}\) or \(\mathbb{C}\),
and \(J\) is a norm, the infimum is attained after quotienting the kernel of \(d_0\).

\emph{Proof sketch.} Restrict to a complement of \(\ker d_0\). On this complement,
\(f\mapsto J(\delta-d_0f)\) is continuous and coercive because \(d_0\) is injective
and all finite-dimensional real norms are equivalent (view a complex space as
a real space). Over \(\mathbb{R}\),
a continuous coercive function attains its minimum on compact sublevel sets. \(\square\)

The finite-dimensional, real/complex scalar, and norm hypotheses matter. Finite dimension
and a norm do \textbf{not} ensure attainment over general rational carriers \(\mathbb{Q}\).
For an explicit counterexample, take two parallel equally oriented edges
between two vertices, fix the first potential to zero, and write
\(d_0 f = (x,x) \in \mathbb{Q}^2\) for the second potential \(x \in \mathbb{Q}\).
Let the edge field be \(\delta = (1,0)\).
Equip \(C^1(G; \mathbb{R})\) with the strictly convex weighted \(\ell^4\) norm
\(J(u_1, u_2) = (u_1^4 + 2 u_2^4)^{1/4}\). The residual is \(\delta - d_0 f = (1-x, -x)\), with
\(J(\delta - d_0 f)^4 = (1-x)^4 + 2x^4\). The unique real minimizer is
\(x^* = 1/(1 + \sqrt[3]{2}) \notin \mathbb{Q}\). For all rational potentials
\(f \in C^0(G; \mathbb{Q})\), \(J(\delta - d_0 f) > J(1-x^*, -x^*)\), so the infimum over
rational potentials is not attained. For scalar rational coefficients, positive
rational weights and rational repair data, the weighted \(L^1\) problem of §4.3
is a feasible bounded rational linear programme and has a rational optimizer.
More generally the same conclusion holds for a norm given by finitely many
rational linear forms and rational input data. Polyhedrality alone is insufficient:
\(J(u,v)=|u+\sqrt{2}v|+|v|\) is a norm with a polyhedral real unit ball, but
on the same repairs its unique real minimizer is
\(x^*=1/(1+\sqrt{2})\). Rational repairs approach that value without attaining it.

\subsubsection{\texorpdfstring{4.3 Weighted \(L^1\) geometry}{4.3 Weighted L\^{}1 geometry}}\label{weighted-l1-geometry}

Choose positive edge weights \(w_e>0\) and set

\[
\|x\|_{1,w}=\sum_{e\in E}w_e\|x_e\|_A.
\]

This norm is especially appropriate when disjoint positive witness masses must
add exactly. Its dual is a weighted \(L^\infty\) constraint. In the real scalar case,
finite-dimensional convex duality {[}REF-MATH-005{]} gives a circulation certificate
of the form

\[
\Phi_{1,w}([\delta])
=\max\left\{\langle y,\delta\rangle:
d_0^{\!*}y=0,\ |y_e|\le w_e\right\},
\]

under the standard linear-programming hypotheses. The constraint
\(d_0^{\!*}y=0\) says that \(y\) is a circulation annihilating coboundaries.

\texttt{WDP-CH04-C01} \textbf{Conditional result.} The displayed primal--dual identity is an
imported finite-dimensional convex-duality result specialized to the declared
weighted graph complex. It is not proved here for arbitrary coefficient
spaces.

\subsubsection{4.4 Closure requires another degree}\label{closure-requires-another-degree}

Suppose a finite \(2\)-cell structure supplies
\(d_1:C^1\to C^2\). A closure diagnostic may be defined by

\[
\Gamma(\delta)=J_2(d_1\delta).
\]

If \(J_2\) is a norm, this descends to \(Q^1\) if and only if \(d_1d_0=0\)
on the admissible repairs. Necessity follows by comparing a repair with the
zero representative; sufficiency follows from linearity. For a seminorm the
exact condition is \(\operatorname{im}(d_1d_0)\subseteq\ker J_2\). The failure

\[
K=d_1d_0
\]

is itself a typed curvature or composition defect. Calling \(Q^1\) ``cohomology''
while \(K\ne0\) is an error. Section 16.7 computes \(K\) for declared restriction
maps on a finite poset, where it is exactly the failure of the composition law.

\subsubsection{4.5 Worked example: exact path and obstructed cycle}\label{worked-example-exact-path-and-obstructed-cycle}

On a path \(0\to1\to2\), let \(\delta=(2,3)\). Taking
\(f=(0,2,5)\) gives \(d_0f=\delta\), so \(\Phi([\delta])=0\) for every norm.

The path is one case of a general fact. On a finite forest,
\(\dim\operatorname{im}d_0=(|V|-c)\dim A=|E|\dim A\) for a finite-dimensional coefficient space \(A\), where \(c\) is the number of components. For any vector coefficient space,
\(d_0\) is onto \(C^1\) and every edge cochain is exact. A graph without cycles
therefore has zero quotient for every coefficient choice, and the edgeless
graph has no edge cochains at all. Neither carries a nontrivial obstruction.
For the same reason every \(U(1)\) connection of §5.2 on a forest is
gauge-equivalent to the trivial connection: fix the gauge at one root per
component and propagate along the unique paths. Its connection Laplacian is
then unitarily equivalent to the ordinary weighted graph Laplacian.

On the oriented triangle, take \(\delta=(1,1,1)\). As in Chapter 2, circulation is
\(3\), whereas every coboundary has circulation \(0\). Hence \(\delta\) is not exact.
For uniform weights and scalar coefficients, symmetry shows that the best
\(L^1\) repair leaves total residual at least \(3\) under the chosen orientation;
the dual circulation \(y=(1,1,1)\) certifies the bound. The example illustrates a
quotient obstruction. It says nothing about quantum contextuality until a
measurement cover and empirical model are also supplied.

\subsubsection{4.6 Replayed vector-valued three-cycle}\label{replayed-vector-valued-three-cycle}

The source corpus also supplied one numerical Python sanity check. On the
oriented three-cycle it used the vector-valued edge defect

\[
\delta=((1,0),(0,1),(0,0)).
\]

Its circulation is \(c=(1,1)\). Every repaired residual
\(r=\delta-d_0\alpha\) has the same sum \(c\), so the triangle inequality gives

\[
\sum_e\|r_e\|_1\ge\|c\|_1=2,
\qquad
\sum_e\|r_e\|_2\ge\|c\|_2=\sqrt2.
\]

Both bounds are attained by
\(\alpha=((0,0),(1,0),(1,1))\), which leaves residual
\(((0,0),(0,0),(1,1))\). The executable replacement in
\texttt{code/witnessed\_difference/vector\_obstruction.py} therefore computes the exact
certificate rather than depending on a continuous numerical optimizer. It also
reproduces the historical display values \texttt{2.000000} and \texttt{1.414214}. This is a
finite example under declared norms, not a general optimization oracle or a
physical observer theorem.

\subsubsection{4.7 Evaluation preorders and the degree-1 window}\label{evaluation-preorders-and-the-degree-1-window}

Sections 4.1--4.6 declared the cochain host directly. The primitive-axiomatics
source (Appendix C) builds the same host from local comparisons, and its route
is worth recording because every step is a choice.

Let \(V\) be a finite set of evaluation sites and \(R\subseteq V\times V\) a
declared relation of direct comparisons. Its reflexive--transitive closure
\(R^{*}\) is a preorder: reflexivity is added by the closure, and concatenating
two finite comparison chains gives a finite chain. Directed cycles are
admitted. A cycle does not contradict the construction; it is a place where a
residual can survive repair.

The preorder is not yet a cochain complex. A declared route

\[
P\longrightarrow N(P)\longrightarrow \mathbb F[N(P)]\longrightarrow
\text{normalized chains}\longrightarrow\text{dual cochains}\longrightarrow
\text{2-truncation}
\]

supplies \(C^0=\mathbb F^{V}\), \(C^1=\mathbb F^{E}\), and \(C^2=\mathbb F^{S_2}\)
with simplicial coboundaries, where \(\mathbb F\) is a declared coefficient field,
\(E\) is a declared set of retained nondegenerate 1-simplices of the nerve (pairs
\(u\ne v\) with \((u,v)\in R^{*}\), oriented by the comparison), and \(S_2\) is a
declared set of retained nondegenerate 2-simplices, each of whose faces lies in
\(E\) or is degenerate; degenerate faces vanish in normalized cochains. When
\(R^{*}\) has directed cycles, the nerve has nondegenerate simplices in every
degree, and the 2-truncation keeps only the declared \(E\) and \(S_2\). The nerve \(N(P)\) is a
simplicial set, not a cochain complex and not itself a Dold--Kan object. A
Dold--Kan comparison, if wanted, applies to the degree-reversed chain complex
\(A_2=C^0\), \(A_1=C^1\), \(A_0=C^2\) and represents the declared truncation, not
the nerve.

\texttt{WDP-CH04-D01} \textbf{Definition (evaluation preorder and degree-1 window).} The
evaluation preorder of a declared finite direct-comparison relation \(R\) on \(V\)
is its reflexive--transitive closure \(R^{*}\). A degree-1 diagnostic window is a
declared tuple \(\mathcal W=(C^0,C^1,C^2,d_0,d_1)\) of finite-dimensional
coefficient spaces and linear maps together with a certificate that
\(d_1d_0=0\). The certificate is automatic on the simplicial route above and must
be verified for any hand-declared pair \((d_0,d_1)\).

When the certificate fails, \(K=d_1d_0\) is the composition defect of §4.4 and
none of §§4.8--4.11 applies.

\emph{Orientation convention.} Sections 2.3 and 4.5 orient the triangle cyclically,
\((01,12,20)\). The simplicial route orients it by vertex order, \((01,12,02)\).
Reversing the edge \(20\) to \(02\) negates that component, so the cyclic field
\(\delta=(1,1,1)\) is the simplicial field \((1,1,-1)\); both have circulation \(3\).
Every quantity below is invariant under reorientation once the cochain is
transported with its edges.

\subsubsection{4.8 Closure descent and the degree-1 exact sequence}\label{closure-descent-and-the-degree-1-exact-sequence}

\texttt{WDP-CH04-D02} \textbf{Definition (quotient residual, closure charge, and
signature).} In a window \(\mathcal W\) with a declared norm \(J\) on \(C^1\), set
\(Q^1=C^1/\operatorname{im}d_0\); define the descended closure map
\(\kappa^1:Q^1\to C^2\) by \(\kappa^1([h])=d_1h\); call \(H^1:=\ker\kappa^1\) the
closed sector; and define the quotient residual mass
\(\Phi_1([h])=\Phi_J([h])\) of §4.2, the closure mass
\(\Gamma_2([h])=\|d_1h\|_1\), and, when \(\Phi_1([h])>0\), the closure ratio
\(R_{cl}([h])=\Gamma_2([h])/\Phi_1([h])\). The diagnostic class of \([h]\) is
\texttt{ExactClass} when \([h]=0\), \texttt{ClosedResidual} when \([h]\ne0\) and
\(\Gamma_2([h])=0\), and \texttt{ActiveClosureFailure} when \(\Gamma_2([h])>0\). A window
whose carrier, norm, closure data, or \(d_1d_0\) certificate is missing or ill
typed receives a typed refusal (\texttt{NoRule}), not a class.

Because \(d_1d_0=0\) is certified, a class lies in \(\ker\kappa^1\) exactly when
one (equivalently, every) representative is closed. Thus \(H^1\) is the usual
\(Z^1/B^1\) of the window. The name ``cohomology'' is earned by the certificate, in
keeping with the migration rule of §21.12.

\emph{Norm hypothesis.} Sections 4.9--4.11 take \(\mathbb F=\mathbb R\) or
\(\mathbb C\), or \(\mathbb F=\mathbb Q\) with \(J\) the restriction of a norm on the
real extension \(C^1\otimes\mathbb R\), such as a weighted \(L^1\) or \(L^2\) norm
with rational weights. The restriction matters. The function
\(J(x)=|x_{01}+\sqrt2\,x_{12}+\sqrt3\,x_{02}|\) is a norm on \(\mathbb Q^3\),
because \(1,\sqrt2,\sqrt3\) are linearly independent over \(\mathbb Q\); yet on the
filled triangle of §4.11 the values \(J(\delta-d_0f)\) for \(\delta=(1,1,0)\)
accumulate at \(0\), so \(\Phi_1([\delta])=0\) although \(\Gamma_2([\delta])=2\).
Under the hypothesis, \(\Phi_1\) of a rational class equals \(\Phi_1\) of the same
class over \(\mathbb R\), since rational repairs are dense among real ones and \(J\)
is continuous; and a rational field lies in the real image of \(d_0\) only if it
lies in the rational image.

\texttt{WDP-CH04-P03} \textbf{Finite proposition (closure descent and exactness).} In a
window with \(d_1d_0=0\), the map \(\kappa^1\) is well defined on \(Q^1\) and

\[
0\longrightarrow H^1\longrightarrow Q^1\xrightarrow{\;\kappa^1\;}
\operatorname{im}d_1\longrightarrow0
\]

is exact. The sequence is an extension. In finite dimension it splits, but a
splitting \(Q^1\cong H^1\oplus\operatorname{im}d_1\) requires a chosen section
and is not canonical.

\emph{Proof.} If \(h'=h+d_0\alpha\), then \(d_1h'=d_1h+d_1d_0\alpha=d_1h\), so
\(\kappa^1\) is constant on repair classes. The inclusion of
\(H^1=\ker\kappa^1\) is injective with image \(\ker\kappa^1\), which is exactness at
\(H^1\) and at \(Q^1\). Every element \(d_1h\) of \(\operatorname{im}d_1\) equals
\(\kappa^1([h])\), which is exactness at \(\operatorname{im}d_1\). Closedness of a
class is representative independent because \(\kappa^1\) is. \(\square\)

The Lean package states this on the literal quotient of \(C^1\) by the repair
relation \(h\sim h+d_0\alpha\) in an abstract additive setting; §D.1 records the
declarations.

\subsubsection{4.9 The obstruction signature}\label{the-obstruction-signature}

\texttt{WDP-CH04-P04} \textbf{Finite proposition (diagnostic trichotomy).} In a window with
\(d_1d_0=0\) and a norm \(J\) on finite-dimensional \(C^1\) satisfying the norm
hypothesis, every class
\([h]\in Q^1\) has exactly one diagnostic class, \(\Phi_1([h])=0\) holds exactly for
\texttt{ExactClass}, and every \texttt{ActiveClosureFailure} class has \(\Phi_1([h])>0\).

\emph{Proof.} The quotient seminorm is faithful: over \(\mathbb R\) or \(\mathbb C\)
because \(\operatorname{im}d_0\) is a closed subspace of a finite-dimensional
space, and over \(\mathbb Q\) by the comparison with the real extension under the
norm hypothesis. So \(\Phi_1([h])=0\) iff \([h]=0\). If \([h]=0\) then
\(\Gamma_2([h])=\|d_1\,0\|_1=0\). The three defining conditions---\([h]=0\);
\([h]\ne0\) with \(\Gamma_2=0\); \(\Gamma_2>0\)---are therefore pairwise disjoint and
exhaustive. If \(\Gamma_2([h])>0\) then \([h]\ne0\), hence \(\Phi_1([h])>0\).
\(\square\)

\texttt{WDP-CH04-P05} \textbf{Finite proposition (signature bounds).} Under the hypotheses
of CH04-P04, write
\(\|d_1\|_{J\to1}=\sup\{\|d_1x\|_1:J(x)\le1\}\), a maximum over \(\mathbb R\) or
\(\mathbb C\) and finite over \(\mathbb Q\) by the norm hypothesis. Then

\begin{enumerate}
\def\labelenumi{\arabic{enumi}.}
\tightlist
\item
  \(\Phi_1([h])\ge0\) and \(\Gamma_2([h])\ge0\);
\item
  \(\Phi_1([h])=0\) if and only if \([h]=0\);
\item
  \(\Gamma_2([h])=0\) if and only if \([h]\in H^1\);
\item
  \(\Gamma_2([h])\le\|d_1\|_{J\to1}\,\Phi_1([h])\);
\item
  if \(H^1\ne0\), no constant \(c>0\) satisfies \(\Gamma_2\ge c\,\Phi_1\) on all of
  \(Q^1\);
\item
  on \(Q^1/H^1\), \(\Gamma_2\) is a norm equivalent to the quotient norm.
\end{enumerate}

\emph{Proof.} Items 1--3 follow from the definitions and CH04-P04. For item 4,
\(d_1(h+d_0\alpha)=d_1h\) for every repair \(\alpha\), so
\(\Gamma_2([h])=\|d_1(h+d_0\alpha)\|_1\le\|d_1\|_{J\to1}J(h+d_0\alpha)\); take the
infimum over \(\alpha\). For item 5, a nonzero \([h]\in H^1\) has \(\Phi_1>0\) and
\(\Gamma_2=0\). For item 6, \(\kappa^1\) induces an injective linear map
\(Q^1/H^1\to C^2\), and all norms on a finite-dimensional real or complex space
are equivalent; over \(\mathbb Q\), apply this on the real extension and restrict.
\(\square\)

For the unweighted \(L^1\) norm, \(\|d_1\|_{1\to1}\) is the largest absolute column
sum of the matrix of \(d_1\).

\texttt{WDP-CH04-P06} \textbf{Finite proposition (closure-ratio bound).} Under the
hypotheses of CH04-P04, if \(\Phi_1([h])>0\), then \(0\le R_{cl}([h])\le\|d_1\|_{J\to1}\).

\emph{Proof.} Divide item 4 of CH04-P05 by the positive number
\(\Phi_1([h])\). \(\square\)

The normalized ratio \(r_{cl}=R_{cl}/\|d_1\|_{J\to1}\) is defined only when
\(d_1\ne0\). If \(d_1=0\), every class is closed and \(r_{cl}\) is the undefined
quotient \(0/0\), not zero. Readings such as ``\(R_{cl}\) near
\(\|d_1\|_{J\to1}\) means the residual is dominated by closure failure'' are
interpretive summaries of the bound, not further theorems.

The unnormalized ratio depends on the declared norms. Rescaling \(J\) to \(cJ\) and
the closure norm to \(\lambda\|\cdot\|_1\) multiplies \(R_{cl}\) by
\(\lambda/c\), while \(r_{cl}\) is unchanged, because \(\|d_1\|_{J\to1}\) rescales in
the same way.

The source also proposed a repair policy that accepts a transition only when
\(\Phi_1\) decreases and \(R_{cl}\) stays bounded, with a termination argument for
such runs. A policy of that kind belongs to a separately checked routing layer
and is not absorbed here; see CH04-O01.

\subsubsection{4.10 Dual certificates without importing strong duality}\label{dual-certificates-without-importing-strong-duality}

\texttt{WDP-CH04-P07} \textbf{Finite proposition (dual-certificate soundness).} Over
\(\mathbb F=\mathbb Q\), let \(J=\|\cdot\|_{1,w}\) with positive rational weights
\(w_e\). If \(y\in C^1\) satisfies
\(d_0^{*}y=0\) and \(|y_e|\le w_e\) for every edge, then
\(\langle y,\delta\rangle\le\Phi_{1,w}([\delta])\). Consequently, a repair \(f\)
and such a \(y\) with \(\|\delta-d_0f\|_{1,w}=\langle y,\delta\rangle\) certify
\(\Phi_{1,w}([\delta])=\langle y,\delta\rangle\) and show that \(f\) attains it.

\emph{Proof.} For every \(f\),

\[
\langle y,\delta\rangle
=\langle y,\delta-d_0f\rangle+\langle d_0^{*}y,f\rangle
=\langle y,\delta-d_0f\rangle
\le\sum_e|y_e|\,|(\delta-d_0f)_e|
\le\|\delta-d_0f\|_{1,w}.
\]

Taking the infimum over \(f\) gives the lower bound. A matching pair meets its
own lower bound. \(\square\)

The same argument holds verbatim over \(\mathbb R\) with positive real weights;
it uses the order of the field, so it does not apply over \(\mathbb C\) or
\(\mathbb F_2\).

Here \(d_0^{*}\) is the transpose of \(d_0\) for the standard pairing:
\((d_0^{*}y)(v)=\sum_{e\to v}y_e-\sum_{v\to e}y_e\), the net inflow at \(v\). The
condition \(d_0^{*}y=0\) says that \(y\) is a circulation.

CH04-C01 asserts more: for every \(\delta\) a matching pair exists. That
existence statement is imported finite-dimensional convex duality and remains a
conditional result. CH04-P07 needs no import. The executable package
computes \((f,y)\) with an exact two-phase simplex method using Bland's rule
{[}REF-MATH-010{]} and then checks the pair independently against
CH04-P07; a candidate pair that fails any check is refused, not reported.
Soundness rests on that check, not on the solver.

\texttt{WDP-CH04-P10} \textbf{Finite proposition (repairable mass and certificate pairs).} Under
the hypotheses of CH04-P07,

\begin{enumerate}
\def\labelenumi{\arabic{enumi}.}
\tightlist
\item
  \(0\le\Phi_{1,w}([\delta])\le\|\delta\|_{1,w}\), so the repairable mass
  \(R_{\rm rep}(\delta)=\|\delta\|_{1,w}-\Phi_{1,w}([\delta])\) is nonnegative;
  \(\|\delta\|_{1,w}\) and \(R_{\rm rep}\) depend on the representative \(\delta\),
  while \(\Phi_{1,w}\) depends only on its class;
\item
  for a repair \(f\) with residual \(h=\delta-d_0f\) and a circulation \(y\) with
  \(|y_e|\le w_e\), the equality \(\langle y,\delta\rangle=\|h\|_{1,w}\) holds
  exactly when \(y_e=w_e\operatorname{sign}(h_e)\) wherever \(h_e\ne0\); in
  particular, if \(f\) attains \(\Phi_{1,w}([\delta])\) and
  \(\langle y,\delta\rangle=\Phi_{1,w}([\delta])\), then \(y\) saturates the weights
  on the support of \(h\);
\item
  neither the minimizing residual nor the optimal circulation is unique in
  general.
\end{enumerate}

\emph{Proof.} The zero repair gives the upper bound in item 1, and CH04-P01 gives
class dependence. For item 2, the chain in the proof of CH04-P07 reads
\(\langle y,\delta\rangle=\sum_ey_eh_e\le\sum_e|y_e||h_e|\le\sum_ew_e|h_e|\), and
equality holds throughout exactly when \(y_eh_e=w_e|h_e|\) for every \(e\). For
item 3, on the triangle of §4.5 with edges \((01,12,20)\), unit weights, and
\(\delta=(1,0,0)\), the repairs \(f=(0,0,0)\), \((0,1,0)\), and
\((0,\tfrac23,\tfrac13)\) leave the residuals \((1,0,0)\), \((0,1,0)\), and
\((\tfrac13,\tfrac13,\tfrac13)\), each of mass \(1\) and each met by \(y=(1,1,1)\).
On the theta graph (two vertices joined by three parallel edges \(0\to1\)) with
\(\delta=(1,0,0)\), the repair \(f=0\) is met by each of \(y=(1,-1,0)\),
\((1,0,-1)\), and \((1,-\tfrac12,-\tfrac12)\); so \(\Phi_1=1\) with three optimal
circulations. \(\square\)

On the path \(a\to b\to c\), \(\delta=(1,0)\) is exact, so \(\Phi_1=0\) and
\(R_{\rm rep}=1\), while the zero representative of the same class has
\(R_{\rm rep}=0\). On the lollipop, a triangle \(0\to1\to2\to0\) with a pendant
edge \(2\to3\) and \(\delta=(1,0,0,1)\), the repair \(f=(0,0,0,1)\) and the
circulation \(y=(1,1,1,0)\) meet at \(1\), so \(\|\delta\|_1=2\), \(\Phi_1=1\), and
\(R_{\rm rep}=1\): the pendant unit is repairable and the cycle unit is not. Only
the scalars are determined by the input. Two correct solvers may return
different pairs; each pair is checked against CH04-P07, not against the other.

\emph{Declared selections.} Over \(\mathbb R\) the optimal residuals form a nonempty
bounded convex polyhedron, and so do the optimal circulations once CH04-C01
supplies one. A declared strictly convex tie-break, such as least Euclidean
norm in the declared edge coordinates, selects exactly one of each. The
selection is noncanonical without those declared data. Once the coordinates
and a unique tie-break are fixed, the selected residual and dual are
class-invariant: replacing \(\delta\) by \(\delta+d_0g\) leaves
\(\{\delta-d_0f:f\in C^0\}\) unchanged, and every feasible dual annihilates
\(d_0g\), so its objective and optimal face are unchanged. The site potential
itself need not be class-invariant. On the triangle, both \(\delta=(1,0,0)\)
and \(\delta'=(2,-1,0)=\delta+d_0(0,1,0)\) select the same residual
\((\tfrac13,\tfrac13,\tfrac13)\) and dual
\((1,1,1)\); on the theta graph it selects \((1,0,0)\) and
\((1,-\tfrac12,-\tfrac12)\). The absorbed archived working-paper source 008 checked the
lollipop pair and the selected triangle and theta pairs with exact rational
code at its recorded commit. The current
\href{code/unit_tests/test_observer_domain_and_selection.py}{selection regression}
independently checks the two triangle representatives, their exact primal/dual
witnesses and the squared-norm identity selecting the barycentre; it does not
port the historical general selection code. On the triangle's optimal
residual face, \(r_i\ge0\) and \(\sum_i r_i=1\), hence
\(\sum_i r_i^2-\tfrac13=\sum_i(r_i-\tfrac13)^2\ge0\),
with equality only at the stated barycentre. This establishes its unique
least-Euclidean-norm selection on the entire real optimal face.

Two routes between the same endpoints give a circulation, so a failed
commuting diagram is a dual certificate.

\texttt{WDP-CH04-P11} \textbf{Finite proposition (path discrepancy bounds the quotient
defect).} Under the hypotheses of CH04-P07, let \(p\) and \(q\) be directed paths
from \(s\) to \(t\), read as edge-count vectors in \(\mathbb Q^E\), and let
\(\gamma=p-q\). Then \(d_0^{*}\gamma=0\), the discrepancy
\(\Delta=\sum_{e\in p}\delta_e-\sum_{e\in q}\delta_e\) equals
\(\langle\gamma,\delta\rangle\) and does not change when \(\delta\) is replaced by
\(\delta-d_0f\), and if \(\gamma\ne0\),

\[
|\Delta|\le c_\gamma\,\Phi_{1,w}([\delta]),\qquad
c_\gamma=\max_{e}\frac{|\gamma_e|}{w_e}.
\]

In particular \(\Delta\ne0\) certifies \(\Phi_{1,w}([\delta])>0\).

\emph{Proof.} Traversing a directed path adds \(1\) to the net inflow at its end, \(-1\)
at its start, and \(0\) at every interior visit, so
\(d_0^{*}p=\mathbf 1_t-\mathbf 1_s=d_0^{*}q\) and \(d_0^{*}\gamma=0\). Hence
\(\Delta=\langle p,\delta\rangle-\langle q,\delta\rangle=\langle\gamma,\delta\rangle\)
and \(\langle\gamma,d_0f\rangle=\langle d_0^{*}\gamma,f\rangle=0\), which is the
invariance of CH05-P04, item 2, for the 1-cycle \(\gamma\). Both
\(\pm\gamma/c_\gamma\) are circulations with \(|y_e|\le w_e\), so CH04-P07 gives
\(\pm\langle\gamma,\delta\rangle\le c_\gamma\Phi_{1,w}([\delta])\). \(\square\)

For simple paths \(\gamma_e\in\{-1,0,1\}\), so with unit weights the bound is
\(|\Delta|\le\Phi_1\). On the triangle \(A\to B\), \(B\to C\), \(A\to C\) with
\(\delta=(\tfrac1{10},\tfrac15,\tfrac25)\), the routes \(ABC\) and \(AC\) give
\(\Delta=-\tfrac1{10}\) and \(\gamma=(1,1,-1)\). The repair \(f=(0,\tfrac1{10},\tfrac3{10})\)
leaves \(\tfrac1{10}\) on \(AC\) and meets \(y=(-1,-1,1)\), so
\(\Phi_1=\tfrac1{10}=|\Delta|\), while \(\|\delta\|_1=\tfrac7{10}\) and
\(R_{\rm rep}=\tfrac35\): the edge-mass bound \(|\Delta|\le\sum_e|\delta_e|\) is
seven times weaker. With \(\delta=(\tfrac1{10},\tfrac15,-\tfrac25)\) instead,
\(\Delta=\Phi_1=\|\delta\|_1=\tfrac7{10}\). With weights \((\tfrac12,1,1)\) and
\(\delta=(1,0,0)\), \(c_\gamma=2\) and \(f=0\) meets \(y=(\tfrac12,\tfrac12,-\tfrac12)\),
so \(\Phi_{1,w}=\tfrac12\) and \(|\Delta|=1=c_\gamma\Phi_{1,w}\).

The discrepancy locates repair-invariant mass on the cycle \(\gamma\); it does
not say which edge is wrong. A version with declared transports into a common
comparison space and a declared gauge per edge reduces to the triangle
inequality under its hypotheses and is not carried; with identity restrictions a
section-induced field \(\delta=d_0s\) is exact and every discrepancy vanishes.
The absorbed source that stated the path bound had no Lean and no executable
check of it.

\subsubsection{4.11 Worked example: a filled triangle with one non-closing edge field}\label{worked-example-a-filled-triangle-with-one-non-closing-edge-field}

\texttt{WDP-CH04-P08} \textbf{Finite proposition (filled-triangle signature).} Take vertices
\(\{0,1,2\}\), edges \((01,12,02)\), and one 2-simplex \((0,1,2)\), with

\[
d_0f=(f_1-f_0,\;f_2-f_1,\;f_2-f_0),\qquad d_1x=x_{01}+x_{12}-x_{02}.
\]

For \(\delta=(1,1,0)\) over \(\mathbb Q\) with unit weights: \(d_1d_0=0\);
\(\Phi_1([\delta])=2\), attained at \(f=(0,1,0)\) and certified by the circulation
\(y=(1,1,-1)\); \(\Gamma_2([\delta])=2\); \(\|d_1\|_{1\to1}=1\); and
\(R_{cl}([\delta])=1\), which saturates CH04-P06.

\emph{Proof.} \(d_1d_0f=(f_1-f_0)+(f_2-f_1)-(f_2-f_0)=0\). For \(f=(0,1,0)\),
\(d_0f=(1,-1,0)\) and \(\delta-d_0f=(0,2,0)\), which has \(L^1\) mass \(2\). The vector
\(y=(1,1,-1)\) has \(|y_e|\le1\), and its net inflow is \(0-(1-1)=0\) at vertex \(0\),
\(1-1=0\) at vertex \(1\), and \((1-1)-0=0\) at vertex \(2\); equivalently
\(\langle y,d_0f\rangle=(f_1-f_0)+(f_2-f_1)-(f_2-f_0)=0\) for every \(f\), so
\(d_0^{*}y=0\). Since
\(\langle y,\delta\rangle=2\), CH04-P07 gives \(\Phi_1=2\). Next,
\(d_1\delta=2\), so \(\Gamma_2=2\). The matrix of \(d_1\) is the row \((1,1,-1)\), whose
largest absolute column sum is \(1\). Finally \(R_{cl}=2/2=1\). \(\square\)

The class is \texttt{ActiveClosureFailure}. Compare the hollow triangle of §4.5: with
no 2-simplex, \(C^2=0\), so every class is closed; \(\delta=(1,1,1)\) has
\(\Phi_1=3\), certified by \(y=(1,1,1)\), and \(\Gamma_2=0\). That class is
\texttt{ClosedResidual} and witnesses item 5 of CH04-P05: filling the triangle
converts a closed residual into closure charge.

Growing a complex is a natural candidate for a repair mechanism. It is not one.

\texttt{WDP-CH04-P09} \textbf{Finite proposition (growth does not lower \(\Phi_1\)).} Let \(K\)
be a subcomplex of a finite complex \(K'\), with weights on \(K\) restricted from
\(K'\), and let \(h\) be the restriction of an edge field \(h'\) on \(K'\). Then
\(\Phi_1(h)\le\Phi_1(h')\). Attaching 2-simplices leaves \(\Phi_1\) unchanged, and
can raise \(\Gamma_2\). Subdividing an edge through a new vertex of degree two,
with the old value split into two summands and both new edges carrying the old
weight, leaves \(\Phi_1\) unchanged.

\emph{Proof.} For any repair \(\alpha'\) on \(K'\), dropping the edges outside \(K\) from
the weighted sum gives
\(\|h'-d_0\alpha'\|_{K'}\ge\|h-d_0(\alpha'|_K)\|_K\ge\Phi_1(h)\); take the
infimum over \(\alpha'\). Attaching 2-simplices changes neither \(C^0\), \(C^1\), nor
\(d_0\). For a subdivided edge \(u\to m\to v\) with values \(x_1+x_2=x\) and weight
\(w\),
\(w|x_1-(f_m-f_u)|+w|x_2-(f_v-f_m)|\ge w|x-(f_v-f_u)|\) by the triangle
inequality, with equality at \(f_m=f_u+x_1\); so every repair of the subdivided
complex has a repair of the original of no greater cost, and conversely. \(\square\)

On the hollow triangle with the cyclic field of §4.5, \(\Phi_1=3\) and
\(\Gamma_2=0\); after filling, \(\Phi_1=3\) and \(\Gamma_2=3\). Splitting an edge
value as \((\tfrac12,\tfrac12)\) or \((5,-4)\) leaves \(\Phi_1=3\). None of these
growth moves lowers \(\Phi_1\). A decrease requires deleting cells, changing
declared weights (including unequal weights on a subdivided edge), or changing
existing edge values: deleting edge \(20\) leaves a tree with \(\Phi_1=0\), and
subdividing edge \(01\) with new weights \((\tfrac12,\tfrac12)\) lowers \(\Phi_1\) to
\(\tfrac32\).

\emph{Restricted repair sets.} The same argument covers a declared set
\(A\subseteq C^0(K)\) of admissible repairs, with
\(\Phi_A(h)=\inf_{\alpha\in A}\|h-d_0\alpha\|_{1,w}\). If \(K'\) arises from \(K\) by
the growth moves of CH04-P09, \(h'\) extends \(h\) as there, and the declared
repair set \(A'\subseteq C^0(K')\) restricts into \(A\) on the old vertices, then
\(\Phi_{A'}(h')\ge\Phi_A(h)\): for \(\alpha'\in A'\) the proof of CH04-P09 gives
\(\|h'-d_0\alpha'\|_{K'}\ge\|h-d_0(\alpha'|_K)\|_K\ge\Phi_A(h)\). A decrease
needs a repair set whose restrictions leave \(A\). On one edge \(0\to1\) with value
\(1\) and \(A=\{0\}\), \(\Phi_A=1\). Subdividing through \(m\) with values \((1,0)\) and
\(A'=\{\alpha':\alpha'(0)=\alpha'(1)=0\}\) keeps
\(\Phi_{A'}=\min_t(|1-t|+|t|)=1\). Taking \(A'=C^0(K')\) gives \(0\), but only
because its restrictions include repairs with \(\alpha(1)\ne\alpha(0)\). That is
a change of declared repair data, not growth.

\subsubsection{4.12 What the signature does not do}\label{what-the-signature-does-not-do}

\texttt{WDP-CH04-N01} \textbf{Nonclaim.} A diagnostic class, signature value, or
certificate is not a routing terminal, repair action, stabilization verification record, or
authority. In particular, \texttt{ExactClass} is not \texttt{ExactRepair}, and
\texttt{ClosedResidual} is not \texttt{StabilizedSurvivor}.

Routing a class to a terminal requires a separately checked policy that consumes
the operation, witness, budget, and stabilization evidence. The source's
terminal enumerations, lift specifications, guarded repair policy, and reflective
completion of contextual systems are not automatically carried by a diagnostic.
Chapter 21 and Appendix J give the declared finite routing interface and its limits as CH04-O01.

\texttt{WDP-CH04-N02} \textbf{Nonclaim.} Growth of the complex does not repair a quotient
residual: passing to a larger complex, face attachment, and equal-weight
subdivision never lower \(\Phi_1\); a decrease requires deleting cells, changing
declared weights (including unequal subdivision weights), or changing existing
edge values; and survival of a class to the limit page of a filtration spectral
sequence records cohomology of the final complex rather than positivity of
\(\Phi_1\).

Projection is not repair either. On the filled triangle with \(\delta=(1,1,0)\),
orthogonal projection onto \(\ker d_1\) subtracts \((2/3,2/3,-2/3)\in\operatorname{im}d_1^{*}\)
and returns \((1/3,1/3,2/3)\), which is closed. But
\(\operatorname{im}d_1^{*}\) meets \(\operatorname{im}d_0\) only in \(0\), so the
difference is not a repair and \(\Phi_1\) of the original class is still \(2\). The
projection replaces the datum; it does not diagnose it.

The window supports one sharp falsifier. Two implementations given the same
declared window, coefficient field, weights, and cochain must return the same
signature. A divergence refutes an implementation or exposes an undeclared
assumption. Certificate pairs are not compared this way: by CH04-P10 two
correct implementations may return different residuals and circulations, and
each pair is checked against CH04-P07 on its own.

\textbf{Interpretive remark (two opposite misreadings of the quotient).} The quotient
\(C^1/\operatorname{im}d_0\) blocks two errors. The first treats every visible
difference as an obstruction, although it may be a coordinate choice, a unit
mismatch, a calibration offset or a naming convention, all of which are
coboundaries once the repair data declare them. The second treats every
difference as removable redescription; the quotient refutes it whenever a class
survives all declared repairs. Under a declared repair space, the coboundaries
are the differences that redescription removes, and the nonzero class is what
survives it. A difference in formatting between two generated answers is of the
first kind under a formatting repair. A lost boundary condition, proof
obligation or database constraint is of the second. Whether a difference can be
removed is a question about the declared repair data, never about the
appearance of the difference.

\begin{center}\rule{0.5\linewidth}{0.5pt}\end{center}

\subsection{Chapter 5 --- Optional quantum-capable carriers}\label{chapter-5-optional-quantum-capable-carriers}

\subsubsection{5.1 Branch declaration}\label{branch-declaration}

This chapter begins a branch, not the trunk. We now \textbf{declare} complex
coefficients, phases, inner products, and matrix representations. Those imports
are recorded because later quantum comparisons require them. Nothing in
Chapters 1--4 selects them uniquely.

\texttt{WDP-CH05-D01} \textbf{Definition (quantum-capable carrier package).} A finite
quantum-capable carrier is

\[
\mathfrak Q=(G,w,\mathbb C^N,U,\nabla,\langle-,-\rangle),
\]

where \(G\) is a finite graph or cell complex, \(w\) is a positive weighting,
\(\mathbb C^N\) is a declared complex fiber, \(U\) assigns unitary transport to
oriented edges, \(\nabla\) is the associated covariant difference, and
\(\langle-,-\rangle\) is a declared Hermitian inner product.

The adjective ``quantum-capable'' means only that familiar finite quantum
structures can be represented. It does not mean that a physical quantum theory
has been derived.

\subsubsection{\texorpdfstring{5.2 \(U(1)\) connection and holonomy}{5.2 U(1) connection and holonomy}}\label{u1-connection-and-holonomy}

For a line fiber, assign \(U_{uv}\in U(1)\) with
\(U_{vu}=U_{uv}^{-1}\). For a vertex field \(\psi:V\to\mathbb C\), set

\[
(\nabla_U\psi)_{uv}=\psi(v)-U_{uv}\psi(u).
\]

A gauge transformation \(g:V\to U(1)\) acts by

\[
\psi(u)\mapsto g(u)\psi(u),\qquad
U_{uv}\mapsto g(v)U_{uv}g(u)^{-1}.
\]

\texttt{WDP-CH05-P01} \textbf{Finite proposition (gauge covariance).} Under the displayed
action,

\[
(\nabla_{U^g}\psi^g)_{uv}=g(v)(\nabla_U\psi)_{uv}.
\]

The proof is a direct substitution. Consequently, any norm invariant under
vertexwise phase multiplication gives a gauge-invariant defect magnitude.

For an oriented loop \(\gamma=e_1\cdots e_k\), the holonomy is

\[
\operatorname{Hol}_U(\gamma)=U_{e_k}\cdots U_{e_1}.
\]

In the abelian \(U(1)\) case it is gauge invariant. In a nonabelian unitary group
it transforms by conjugation, so conjugacy-invariant functions such as trace or
spectrum are gauge invariant.

The same product is defined on any walk that may traverse an edge against its
declared orientation: such a step contributes \(U_{vu}=U_{uv}^{-1}\), which is the
complex conjugate in \(U(1)\). \textbf{Finite proposition (walk transport).} For a walk
\(\gamma\) from \(u\) to \(v\), the transport \(T_U(\gamma)\) satisfies
\(T_{U^g}(\gamma)=g(v)\,T_U(\gamma)\,g(u)^{-1}\). The gauge factors of each
interior vertex cancel in adjacent steps. On a closed walk (\(u=v\)) the \(U(1)\)
transport is therefore gauge invariant. On an open walk it is only covariant at
its two endpoints. An open-path phase is not an invariant observable of the
connection until its endpoint frames are fixed by a further declaration.

\subsubsection{5.3 Connection Laplacian}\label{connection-laplacian}

Let \(W\) be the positive diagonal edge-weight operator. The finite connection
Laplacian is

\[
L_U=\nabla_U^{*}W\nabla_U.
\]

\texttt{WDP-CH05-P02} \textbf{Finite proposition.} \(L_U\) is Hermitian and positive
semidefinite.

\emph{Proof.} Hermiticity follows from the adjoint construction. For every \(\psi\),

\[
\langle\psi,L_U\psi\rangle
=\langle\nabla_U\psi,W\nabla_U\psi\rangle\ge0
\]

because \(W\) is positive. \(\square\)

The kernel of such an operator is exactly the joint kernel of its pieces, but
only when the weights are definite.

\texttt{WDP-CH05-P07} \textbf{Finite proposition (kernel of a weighted sum of forms).} For
finitely many linear maps \(A_a\) out of a finite-dimensional inner-product space
over \(\mathbb Q\), \(\mathbb Q(i)\), \(\mathbb R\), or \(\mathbb C\), and
positive-definite weights \(W_a\), the operator \(H=\sum_aA_a^{*}W_aA_a\) is
Hermitian and positive semidefinite, and \(H\psi=0\) exactly when \(A_a\psi=0\) for
every \(a\). The equivalence can fail when some \(W_a\) is only positive
semidefinite.

\emph{Proof.} \(\langle\psi,H\psi\rangle=\sum_a\langle A_a\psi,W_aA_a\psi\rangle\ge0\).
If \(H\psi=0\), every summand vanishes, and definiteness gives \(A_a\psi=0\); the
converse is immediate. With \(W=0\) and \(A\ne0\), \(H=0\) while \(A\psi\ne0\).
\(\square\)

The connection Laplacian of CH05-P02 is the case of a single covariant
difference, so its kernel is exactly the set of covariantly constant fields.

Gauge-related connections yield unitarily conjugate Laplacians, hence identical
spectra. This is a representation theorem inside the declared package, not a
selection theorem for the package.

\subsubsection{5.4 Real cochains and phases}\label{real-cochains-and-phases}

Given a real edge cochain \(a\) and a declared scale \(\alpha\in\mathbb R\), one may
set

\[
U_e=\exp(i\alpha a_e).
\]

If \(a\) changes by a real coboundary \(d_0f\), then \(U\) changes by the \(U(1)\) gauge
transformation \(g(v)=\exp(i\alpha f(v))\). Thus exponentiation descends from real
cochain classes to gauge classes once \(\alpha\) is fixed.

The phase scale is not selected by the real quotient. Distinct real cochains can
also exponentiate to the same phase because of \(2\pi\) periodicity. The map is
therefore structured and potentially lossy.

The scale must also be fixed before the cochain is seen. Choosing
\(\alpha=1/I_1(a)\), so that each cochain is normalized by its own raw \(\ell^1\)
mass, does not define a phase convention. It is not linear in \(a\). It is not
constant on cochain classes either, because \(I_1\) depends on the representative
and gauge-equivalent representatives can then receive different scales. A fixed
\(\alpha\), declared as observer data, is what makes the descent statement above
meaningful. Neither \(I_1\) nor \(\Phi_1\) selects it.

\subsubsection{5.5 Matrix fibers}\label{matrix-fibers}

For \(N\ge2\), a generalized Gell-Mann basis gives a real basis of traceless
Hermitian \(N\times N\) matrices. Together with the identity, it coordinatizes
Hermitian matrices. The basis is useful for calculations but depends on a
declared complex fiber and basis convention. It does not derive \(SU(N)\) from
finite difference.

\textbf{Finite proposition (generalized Gell-Mann basis).} With matrix units
\(E_{jk}\), take \(E_{jk}+E_{kj}\) and \(-iE_{jk}+iE_{kj}\) for \(j<k\), and, for
\(1\le\ell\le N-1\),
\(\sqrt{2/(\ell(\ell+1))}\,\operatorname{diag}(1,\dots,1,-\ell,0,\dots,0)\) with
\(\ell\) leading ones. These \(N^2-1\) matrices \(\lambda_a\) are Hermitian, traceless and
satisfy \(\operatorname{Tr}(\lambda_a\lambda_b)=2\delta_{ab}\), so they form a real
basis of the traceless Hermitian matrices.

\emph{Proof.} There are \(N(N-1)/2\) symmetric, the same number of antisymmetric and
\(N-1\) diagonal elements, \(N^2-1\) in all. Each is Hermitian and traceless by
inspection. Off-diagonal elements with different supports have zero trace
product. The symmetric and antisymmetric elements on one support have trace
product \(\operatorname{Tr}((E_{jk}+E_{kj})(-iE_{jk}+iE_{kj}))=-i+i=0\). The
diagonal elements are orthogonal by direct summation, and the displayed scale
makes every squared Frobenius norm \(2\). Orthogonal nonzero elements are linearly
independent, and their number equals the real dimension \(N^2-1\). \(\square\)

With \(T_a=\lambda_a/2\), the commutators define structure constants by
\([T_a,T_b]=i\sum_cf_{abc}T_c\). They are real, because \(-i[T_a,T_b]\) is Hermitian
and traceless and the \(T_c\) are a real basis. They are totally antisymmetric,
because \(f_{abc}=-2i\operatorname{Tr}([T_a,T_b]T_c)\) and the trace is cyclic. None
of this structure comes from a graph. The label space and its basis may be
attached to vertices or edges only by a separate modeling choice. The exact
companion checks of these identities, together with walk transport, the
connection Laplacian and cochain exponentiation, are the
\texttt{omega\_quantum\_physics} integration tests
\texttt{tests/foundations/quantum\_carrier\_primitive\_difference.rs}.

\subsubsection{5.6 Worked example: phase around a triangle}\label{worked-example-phase-around-a-triangle}

Let each oriented edge of a triangle carry the same real value \(a_e=\theta\) and
take \(\alpha=1\). The loop holonomy is \(e^{3i\theta}\). If
\(\theta=2\pi/3\), the holonomy is \(1\) even though each edge phase is nontrivial.
If \(\theta=\pi/3\), the holonomy is \(-1\). The example shows why local phase and
global loop residue differ. It does not establish a particle, an amplitude
rule, or an experiment.

\texttt{WDP-CH05-N01} \textbf{Nonclaim.} This chapter does not derive Hilbert space,
unitarity as dynamics, tensor composition, entanglement, observables, the Born
rule, collapse, or measurement.

\texttt{WDP-CH05-N02} \textbf{Nonclaim.} The existence of a connection Laplacian does not
identify its lowest eigenvector with a physical vacuum or foundational ground.

\subsubsection{5.7 Nonabelian transport and its exact scope}\label{nonabelian-transport-and-its-exact-scope}

For a declared fiber \(\mathbb C^N\), edge transport can take values in \(U(N)\).
With the convention

\[
(\nabla_U\psi)_{uv}=\psi(v)-U_{uv}\psi(u),
\]

a vertexwise gauge \(g:V\to U(N)\) acts by

\[
\psi(v)\mapsto g(v)\psi(v),\qquad
U_{uv}\mapsto g(v)U_{uv}g(u)^{-1}.
\]

The same substitution as in the line-bundle case gives covariance at the target
vertex. Loop holonomy based at \(v\) transforms as
\(H_\gamma\mapsto g(v)H_\gamma g(v)^{-1}\). Thus its conjugacy class, trace, and
spectrum are invariant. A chosen matrix representative is not invariant.

This distinction is important when a source speaks of ``the holonomy.'' The
gauge-invariant object is generally a conjugacy class or class function, not a
specific matrix. Likewise, a basis of \(\mathfrak{su}(N)\) is computational data;
the Lie algebra does not privilege one generalized Gell-Mann basis.

\subsubsection{5.8 Density matrices are another import}\label{density-matrices-are-another-import}

Given the complex Hilbert carrier, one may declare a density operator
\(\rho\succeq0\) with \(\operatorname{tr}\rho=1\). Unitary conjugation preserves
positivity and trace. A completely positive trace-preserving map can then be
studied, and trace norm has useful contractivity properties under appropriate
hypotheses. Every noun in this sentence---density operator, positivity cone,
trace, complete positivity, tensor extension---is additional structure.

The monograph therefore uses the following import ladder:

\[
\begin{array}{c}
\text{real graph cochains}\\
\downarrow\ \text{declare phase scale and complex line fiber}\\
U(1)\text{ transport}\\
\downarrow\ \text{declare }\mathbb C^N\text{ and Hermitian product}\\
U(N)\text{ transport and connection Laplacian}\\
\downarrow\ \text{declare state/effect cones and normalization}\\
\text{finite operational quantum carrier}\\
\downarrow\ \text{declare instruments/composition/dynamics}\\
\text{candidate quantum theory.}
\end{array}
\]

Each arrow is an interface. None is licensed by similarity of notation.

\subsubsection{5.9 Stability under perturbation}\label{stability-under-perturbation}

Finite-dimensionality gives a useful but limited stability statement. If two
connection matrices differ slightly in an operator norm, their Laplacians also
differ continuously; standard finite-matrix perturbation bounds can then
control eigenvalue changes. To use such a bound one must specify the norm,
weights, graph, and convention for comparing connections. Gauge-equivalent
connections need no perturbation estimate because they are exactly unitarily
conjugate.

\texttt{WDP-CH05-O01} \textbf{Conditional result (resolved finite target).} For two unweighted rational U(1) connection Laplacians on the same ordered finite carrier, the exact maximum absolute-real-plus-imaginary row sum of their difference bounds every ordered eigenvalue displacement; a separately valid gap bound greater than twice this value ensures gap survival.
The proof, native ownership and scope are given in Appendix I.

\subsubsection{5.10 Higher-degree windows on finite simplicial complexes}\label{higher-degree-windows-on-finite-simplicial-complexes}

Sections 5.10 and 5.12 use real or rational cochains, and §5.11 compares
rational with mod-2 coefficients; they sit in this branch because they extend the degree-1 window of §§4.7--4.11 rather than the
trunk. Let \(K\) be a finite abstract simplicial complex on a totally ordered
vertex set, and let \(C^k(K;\mathbb F)=\mathbb F^{\Sigma_k}\) with

\[
(\delta^k\varphi)([v_0,\dots,v_{k+1}])
=\sum_{i=0}^{k+1}(-1)^i\varphi([v_0,\dots,\widehat{v_i},\dots,v_{k+1}]).
\]

\texttt{WDP-CH05-P03} \textbf{Finite proposition (degree-\(k\) windows).} On a finite
oriented simplicial complex, \(\delta^{k+1}\delta^k=0\) for every \(k\). Hence each
triple \((C^{k-1},C^k,C^{k+1})\) with a norm \(J\) on \(C^k\) satisfying the norm
hypothesis of §4.8 is a window in the sense
of CH04-D01: \(\kappa^k([h])=\delta^kh\) is well defined on
\(Q^k=C^k/\operatorname{im}\delta^{k-1}\), the sequence
\(0\to H^k\to Q^k\to\operatorname{im}\delta^k\to0\) is exact, and
\(\Gamma_{k+1}\le\|\delta^k\|_{J\to1}\Phi_k\).

\emph{Proof.} In \(\delta^{k+1}\delta^k\varphi\), each \((k)\)-face of a
\((k+2)\)-simplex arises by deleting two vertices in either order, with opposite
signs {[}REF-MATH-002{]}. The remaining statements are
CH04-P03 and CH04-P05 applied to the degree-\(k\) window. \(\square\)

This is a finite simplicial extension. It says nothing about composition-level
gluing of evaluator defects, cup products, or continuum differential forms.

With the standard inner products, the combinatorial Laplacian
\(\Delta_k=\delta^{k*}\delta^k+\delta^{k-1}\delta^{k-1*}\) has kernel isomorphic to
\(H^k\) over \(\mathbb Q\) or \(\mathbb R\). This is elementary finite linear
algebra: \(\ker\Delta_k=\ker\delta^k\cap\ker\delta^{k-1*}\), the orthogonal
complement of \(\operatorname{im}\delta^{k-1}\) inside \(\ker\delta^k\).
{[}REF-MATH-001{]} develops the real Whitney setting in which it is usually met. The
obstruction signature does not use it.

\subsubsection{5.11 Coefficients matter: a torsion witness}\label{coefficients-matter-a-torsion-witness}

The coefficient field is part of the declared window, not a detail.

\texttt{WDP-CH05-P06} \textbf{Finite proposition (a rational solver misses 2-torsion).} For
the six-vertex, ten-triangle triangulation \(K\) of the real projective plane,
\(H^2(K;\mathbb Q)=0\) while \(H^2(K;\mathbb F_2)\cong\mathbb F_2\).

\emph{Proof.} \(K\) has \(6\) vertices, \(15\) edges, and \(10\) triangles. Exact row
reduction gives \(\operatorname{rank}_{\mathbb Q}\delta^1=10\) and
\(\operatorname{rank}_{\mathbb F_2}\delta^1=9\). Since \(C^3=0\),
\(\dim H^2=10-\operatorname{rank}\delta^1\), which is \(0\) over \(\mathbb Q\) and \(1\)
over \(\mathbb F_2\). This agrees with \(H^2(\mathbb{RP}^2;\mathbb Z)\cong
\mathbb Z/2\) and the algebraic homology coefficient theorem {[}REF-MATH-002{]}. \(\square\)

A window declared over \(\mathbb Q\) therefore reports no degree-2 class where a
mod-2 window exposes one. Neither answer is wrong; they answer different
declared questions. The nontrivial mod-2 class is represented by the indicator
of a single triangle; the all-ones 2-cochain is a mod-2 coboundary, because the
ten triangles pair off. Over \(\mathbb Q\) the single-triangle indicator is exact
(\(\Phi_2=0\)), and every rational repair of it has non-integral values: that is
the concrete sense in which a rational solver misses the torsion.

The algebra behind it is short: multiplication by \(2\) on \(\mathbb Z\) has
cokernel \(\mathbb Z/2\), but becomes invertible after tensoring with \(\mathbb Q\).
Rational coefficients are flat and see no torsion.

\subsubsection{5.12 Balance identities on finite graphs}\label{balance-identities-on-finite-graphs}

\texttt{WDP-CH05-P04} \textbf{Finite proposition (combinatorial balance and invariant
pairings).} Items 1 and 2 hold on every finite oriented graph; item 3 concerns
a finite simplicial complex:

\begin{enumerate}
\def\labelenumi{\arabic{enumi}.}
\tightlist
\item
  \(\sum_v(d_0^{*}x)(v)=0\) for every \(x\in C^1\), so \(d_0^{*}x=\rho\) is
  solvable only if \(\sum_v\rho(v)=0\);
\item
  a functional \(F\) on \(C^1\) with \(F(A+d_0\lambda)=F(A)\) for every \(\lambda\)
  descends uniquely to \(Q^1\); in particular, the pairing \(\langle A,\gamma\rangle\)
  with a 1-cycle \(\gamma\) (\(\partial\gamma=0\)) is invariant;
\item
  on a finite simplicial complex with 3-simplices, \(\delta^2\delta^1=0\).
\end{enumerate}

\emph{Proof.} (1) \(\sum_v(d_0^{*}x)(v)=\langle\mathbf 1,d_0^{*}x\rangle
=\langle d_0\mathbf 1,x\rangle=0\) because \(d_0\) kills constants; equivalently,
each edge contributes \(+x_e\) at its target and \(-x_e\) at its source. (2)
\(\langle A+d_0\lambda,\gamma\rangle=\langle A,\gamma\rangle
+\langle\lambda,\partial\gamma\rangle\). (3) is CH05-P03 at \(k=1\).
\(\square\)

Item 1 needs no closedness, connectedness, or boundary condition. The
condition \(\sum_v\rho(v)=0\) is necessary, not sufficient: on a disconnected
graph each component must be neutral separately. The source stated item 1 only
``on a closed complex'' and attributed a boundary leak to graphs with sinks. On a finite graph there is no such leak; it appears only if some
edge is allowed an endpoint outside the vertex set, which is not a graph in
this sense.

\texttt{WDP-CH05-N03} \textbf{Nonclaim.} Reading \(d_0^{*}\) as divergence, \(d_1A\) as flux or
curvature, item 1 as charge conservation, \(\delta^2\delta^1=0\) as a Bianchi
identity or an absence of monopoles, or averaging over a finite symmetry group
as a Noether law is interpretive vocabulary. No Maxwell equation, Lorentz
structure, physical charge, monopole physics, or physical conservation law is
derived.

\subsubsection{5.13 A projective phase obstruction}\label{a-projective-phase-obstruction}

Declare a finite group \(G\) and unitary matrices \(U(g)\) with
\(U(g)U(h)=\omega(g,h)U(gh)\) for phases \(\omega(g,h)\in U(1)\). Associativity of
matrix multiplication forces the 2-cocycle identity
\(\omega(g,h)\,\omega(gh,k)=\omega(g,hk)\,\omega(h,k)\). Rephasing
\(U'(g)=\lambda(g)U(g)\) multiplies \(\omega\) by the coboundary
\(\lambda(g)\lambda(h)/\lambda(gh)\), so only the class
\([\omega]\in H^2(G;U(1))\) is invariant {[}REF-MATH-011{]}.

\texttt{WDP-CH05-P05} \textbf{Finite proposition (Klein-four Pauli obstruction).} Let
\(G=\mathbb Z_2\times\mathbb Z_2=\{e,x,z,xz\}\) with \(U(e)=I\), \(U(x)=\sigma_x\),
\(U(z)=\sigma_z\), and \(U(xz)=\sigma_y\). Then \(U\) is a projective
representation, its phase function satisfies the 2-cocycle identity, and the
ratio \(\omega(x,z)/\omega(z,x)=-1\) is unchanged by every rephasing. Hence
\([\omega]\ne0\): no rephasing makes \(U\) a linear representation.

\emph{Proof.} \(\sigma_x\sigma_z=-i\sigma_y=-i\,U(xz)\) and
\(\sigma_z\sigma_x=i\sigma_y=i\,U(xz)\), so \(\omega(x,z)=-i\), \(\omega(z,x)=i\), and
the ratio is \(-1\). Under rephasing, the coboundary factors at \((x,z)\) and
\((z,x)\) are equal because \(G\) is abelian, so they cancel in the ratio. A linear
representation would have \(\omega\equiv1\) and ratio \(1\). \(\square\)

The source wrote \(\sigma_x\sigma_z=-U(xz)\); the correct factor is \(-i\). Its
conclusion survives because only the ratio is invariant. The result is finite
group cohomology of a declared complex representation. It derives no spin,
qubit, or quantum theory from difference; CH05-N01 still applies.

The Klein-four class is the case \(N=2\) of a family, and the family has an
additive counterpart whose nonzero form no finite-dimensional carrier over a
characteristic-zero field supports. Positive-characteristic exceptions are
part of CH05-P08 below.

\texttt{WDP-CH05-P08} \textbf{Finite proposition (trace obstruction to a finite commutation
relation).} Let \(R\) be a commutative ring, \(N\ge1\), and \(c\in R\) with \(Nc\ne0\).
No \(A,B\in M_N(R)\) satisfy \(AB-BA=cI_N\). Over a field \(K\), some
\(A,B\in M_N(K)\) satisfy \(AB-BA=cI_N\) for some \(c\ne0\) exactly when \(K\) has
characteristic \(p>0\) and \(p\mid N\). The hypothesis \(Nc\ne0\) cannot be weakened to
``\(R\) has characteristic zero'': \(R=\mathbb Z\times\mathbb F_2\) has
characteristic zero, \(c=(0,1)\) is nonzero with \(2c=0\), and \(A=(0,D)\),
\(B=(0,X)\) with

\[
D=\begin{pmatrix}0&1\\0&0\end{pmatrix},\qquad
X=\begin{pmatrix}0&0\\1&0\end{pmatrix}
\]

satisfy \(AB-BA=cI_2\).

\emph{Proof.} In a commutative ring
\(\operatorname{tr}(AB)=\sum_{i,j}A_{ij}B_{ji}=\operatorname{tr}(BA)\), so
\(\operatorname{tr}(AB-BA)=0\), while \(\operatorname{tr}(cI_N)=Nc\ne0\). Over a
field, \(Nc=0\) with \(c\ne0\) forces \(N=0\) in \(K\), that is, characteristic \(p>0\)
with \(p\mid N\). Conversely, on \(\mathbb F_p[x]/(x^p)\) with basis
\(1,x,\dots,x^{p-1}\) let \(D\) differentiate and \(X\) multiply by \(x\). Then
\([D,X]x^k=(k+1)x^k-kx^k=x^k\) for \(k<p-1\) and
\([D,X]x^{p-1}=-(p-1)x^{p-1}=x^{p-1}\), so \([D,X]=I_p\); for \(p\mid N\),
block-diagonal copies of \(cD\) and \(X\) give \(cI_N\). For \(p=2\) these are the
displayed matrices, and \(M_2(\mathbb Z\times\mathbb F_2)=M_2(\mathbb Z)\times
M_2(\mathbb F_2)\) componentwise. \(\square\)

In particular no complex carrier, and no exact carrier over \(\mathbb Q(i)\) or
\(\mathbb Q(\zeta_N)\), supports \(AB-BA=cI\) with \(c\ne0\). The absorbed
archived working-paper source 004 stated the obstruction for every commutative ring of
characteristic zero; the ring above refutes that form. Its cited external Lean
establishes the case of integer matrices only; it was read, not ported, and gives no
Lean credit. Nothing here represents a position, a momentum, or Planck's
constant; CH05-N01 applies.

\texttt{WDP-CH05-P09} \textbf{Finite proposition (the clock-and-shift cocycle).} Let \(N\ge2\),
let \(\omega\) be a primitive \(N\)-th root of unity, and on \(\mathbb C^N\) with
indices modulo \(N\) set \(Xe_j=e_{j+1}\) and \(Ze_j=\omega^je_j\). Then \(X\) and \(Z\)
are unitary with entries in \(\mathbb Z[\omega]\), \(X^N=Z^N=I\), and
\(XZ=\omega^{-1}ZX\). On \(G=\mathbb Z_N\times\mathbb Z_N\) the assignment
\(U(a,b)=X^aZ^b\) is a projective representation with
\(U(a,b)U(a',b')=\omega^{ba'}U(a+a',b+b')\). For \(g=(a,b)\) and \(h=(a',b')\) its
commutator ratio \(c(g,h)/c(h,g)=\omega^{ba'-b'a}\) is unchanged by every
rephasing and equals \(\omega^{-1}\) at \(g=(1,0)\), \(h=(0,1)\), so the class
\([c]\in H^2(G;U(1))\) has order exactly \(N\). For \(N=2\), \(U\) is the assignment of
CH05-P05 up to the rephasing \(U(1,1)=XZ=-i\sigma_y\).

\emph{Proof.} \(XZe_j=\omega^je_{j+1}\) and \(ZXe_j=\omega^{j+1}e_{j+1}\), so
\(Z^bX^{a'}=\omega^{ba'}X^{a'}Z^b\), which gives the product rule; exponents may
be read modulo \(N\) because \(X^N=Z^N=I\). As earlier in this section, a rephasing
multiplies \(c\) by \(\lambda(g)\lambda(h)/\lambda(g+h)\), which is symmetric in \(g\)
and \(h\) because \(G\) is abelian, so the ratio is unchanged. The cocycle \(c^k\) has
ratio \(\omega^{-k}\ne1\) at \((1,0),(0,1)\) for \(0<k<N\), so it is not a coboundary,
while \(c^N\equiv1\). \(\square\)

The two propositions answer one question in two ways. The additive relation
\(AB-BA=cI\) has no finite solution when \(Nc\ne0\); the multiplicative relation
\(XZ=\omega^{-1}ZX\) holds in every dimension as a projective class of order \(N\).
The absorbed source proposed the second as the finite substitute for the first
and did not prove it. Neither derives quantum kinematics; CH05-N01 applies.

\subsubsection{5.14 Trace-norm certificates and channel contractivity}\label{trace-norm-certificates-and-channel-contractivity}

With the imports of §5.8 in place, two standard finite-dimensional results are
used inside this branch and nowhere upstream.

\texttt{WDP-CH05-C01} \textbf{Conditional result (cited result).} On
\(M_d(\mathbb C)\) with the pairing \(\operatorname{Re}\operatorname{tr}(W^\dagger
X)\), for a subspace \(V\) and an operator \(L\),

\[
\min_{X\in V}\|L-X\|_1=\max\{\operatorname{Re}\operatorname{tr}(W^\dagger L):
W\in V^\perp,\ \|W\|_\infty\le1\},
\]

and every completely positive trace-preserving map \(\mathcal E\) satisfies
\(\|\mathcal E(\rho)-\mathcal E(\sigma)\|_1\le\|\rho-\sigma\|_1\) for density
operators \(\rho,\sigma\) {[}REF-QM-020; REF-QM-021; REF-MATH-005{]}.

The first identity is the Schatten-norm analogue of §4.10; its weak direction is
the same one-line argument with the trace pairing. The second is the operator
analogue of observer passivity (§10.5). Neither makes a channel an observer or
derives measurement. The source checked these results numerically with
floating-point eigensolvers and a declared tolerance; this monograph imports
them as theorems and does not treat a numerical check as a certificate.

Two lifts of real edge fields into this branch show what the first identity
does and does not add to §4.10.

\emph{Diagonal lift.} For real \(x\in\mathbb R^E\),
\(\|\operatorname{diag}(x)\|_1=\|x\|_1\), because the singular values of a real
diagonal matrix are the \(|x_e|\). Take \(L=\operatorname{diag}(\delta)\) and the
real subspace \(V=\{\operatorname{diag}(d_0f):f\in\mathbb R^V\}\). The primal
side is then the unit-weight \(\Phi_1([\delta])\). On the dual side only the real
parts \(y_e=\operatorname{Re}W_{ee}\) of the diagonal of \(W\) enter the pairing
with \(L\) and the condition \(W\in V^\perp\); they form a circulation, and
\(|y_e|\le|W_{ee}|\le\|W\|_\infty\le1\). Conversely \(\operatorname{diag}(y)\) is
feasible for every such circulation. So the §4.10 problem is the diagonal
instance of the first identity, whose weak direction there is CH04-P07. For
\(h=(1,-2,\tfrac12)\), \(\|\operatorname{diag}(h)\|_1=\tfrac72=\|h\|_1\). The lift
represents a chosen field: \(\operatorname{diag}(d_0f)\ne0\) whenever
\(d_0f\ne0\), so it does not quotient repairs.

\emph{Density lift.} A nonzero real field \(h\) also has a declared rank-one lift,
\(\psi(h)_e=\operatorname{sign}(h_e)\sqrt{|h_e|/\|h\|_1}\) and
\(\rho(h)=\psi(h)\psi(h)^{\top}\), a symmetric projector of trace one. The real,
phase-free section is a choice: phases \(e^{i\theta_e}\) give other lifts with
the same moduli. The lift remembers exactly the line through \(h\):
\(\rho(h)=\rho(k)\) if and only if \(k=ch\) for some real \(c\ne0\), because
\(\rho(h)\) determines \(\psi(h)\) up to sign and
\(h_e/\|h\|_1=\operatorname{sign}(\psi_e)\psi_e^2\). For unit vectors
\(\psi,\psi'\) the difference \(D=\psi\psi^{\top}-\psi'\psi'^{\top}\) has trace
\(0\), rank at most \(2\), and \(\operatorname{tr}D^2=2-2\langle\psi,\psi'\rangle^2\), so
two lifts have trace-norm separation \(2\sqrt{1-\langle\psi,\psi'\rangle^2}\). Every
functional of lifts is blind to magnitude. Replacing \(\delta\) by
\(\lambda\delta\) with \(\lambda\ne0\) multiplies \(\Phi_1\) by \(|\lambda|\) and the
least-Euclidean-norm residual and dual selections of §4.10 by \(\lambda\) and \(\operatorname{sign}\lambda\),
which leaves their lifts unchanged. This scaling statement does not extend
to an arbitrary strictly convex tie-break, such as a norm centered at a fixed
nonzero target. On the theta graph the selected residual
\((1,0,0)\) and circulation \((1,-\tfrac12,-\tfrac12)\) have
\(\langle\psi,\psi'\rangle^2=\tfrac12\), so their lifts have trace-norm separation
\(\sqrt2\). On the triangle the selections \((\tfrac13,\tfrac13,\tfrac13)\) and
\((1,1,1)\) have equal lifts, at distance \(0\), although \(\Phi_1=1\). A distance
between lifts is not an obstruction magnitude and is not used as one. The lift
is a density operator on the declared Hilbert carrier. It is not thereby
identified with a physical state; unit trace alone supplies no empirical
probability interpretation (CH05-N01).

\subsubsection{5.15 Two finite reconstruction frames}\label{two-finite-reconstruction-frames}

With the imports of §5.8, an operator on \(\mathbb C^d\) can be recovered exactly
from finitely many numbers, and for the two frames below those numbers lie in a
cyclotomic field when the operator's entries do. Both frames are declared
structure on the complex carrier; the calculus selects neither.

Two bases are \emph{unbiased} when
\(|\langle u,v\rangle|^2=\langle u,u\rangle\langle v,v\rangle/d\) for every \(u\) in
one and \(v\) in the other. In dimension \(4\) take the standard basis and

\begin{longtable}[]{@{}ll@{}}
\toprule\noalign{}
Basis & Vectors in \(\mathbb Z[i]^4\) \\
\midrule\noalign{}
\endhead
\bottomrule\noalign{}
\endlastfoot
\(B_1\) & \((1,-1,1,-1)\), \((1,1,-1,-1)\), \((1,-1,-1,1)\), \((1,1,1,1)\) \\
\(B_2\) & \((1,-i,i,1)\), \((1,i,-i,1)\), \((1,-i,-i,-1)\), \((1,i,i,-1)\) \\
\(B_3\) & \((1,-1,-i,-i)\), \((1,1,-i,i)\), \((1,1,i,-i)\), \((1,-1,i,i)\) \\
\(B_4\) & \((1,-i,-1,-i)\), \((1,i,-1,i)\), \((1,i,1,-i)\), \((1,-i,1,i)\) \\
\end{longtable}

\texttt{WDP-CH05-P10} \textbf{Finite proposition (complete unbiased bases in dimensions 3 and
4).} In dimension \(3\), with \(\omega\) a primitive cube root of unity, the
standard basis and the three bases \(\{v_{r,k}\}_{k=0,1,2}\) with
\(v_{r,k}=(\omega^{2rj^2+jk})_{j=0,1,2}\), \(r=0,1,2\), are pairwise unbiased. In
dimension \(4\), the standard basis and \(B_1,\dots,B_4\) are pairwise unbiased. For
the \(d(d+1)\) rank-one projectors \(\Pi=uu^{*}/\langle u,u\rangle\) of either
family, every complex \(d\times d\) matrix \(\rho\) satisfies

\[
\sum_{\Pi}\operatorname{tr}(\rho\Pi)\,\Pi=\rho+\operatorname{tr}(\rho)\,I_d .
\]

With any one of the \(d+1\) bases removed, the identity fails.

\emph{Proof.} Unbiasedness is direct computation over \(\mathbb Z[\omega]\) and
\(\mathbb Z[i]\): in dimension \(3\), \(\langle v_{r,k},v_{r,k'}\rangle=3\delta_{kk'}\),
\(|\langle e_j,v_{r,k}\rangle|^2=1\), and
\(|\langle v_{r,k},v_{r',k'}\rangle|^2=3\) for \(r\ne r'\); in dimension \(4\), each
\(B_m\) is orthogonal with \(\langle u,u\rangle=4\), \(|\langle e_j,u\rangle|^2=1\), and
\(|\langle u,v\rangle|^2=4\) between different \(B_m\). For the identity, let \(V_m\) be
the real span of the traceless Hermitian matrices \(\Pi-I/d\) for \(\Pi\) in basis
\(m\); it has dimension \(d-1\). Across bases
\(\operatorname{tr}(\Pi\Pi')=1/d\), so
\(\operatorname{tr}\big((\Pi-\tfrac Id)(\Pi'-\tfrac Id)\big)=0\): the \(d+1\) spaces
are orthogonal for the positive-definite trace form, and since
\((d+1)(d-1)=d^2-1\) they fill the traceless Hermitian matrices. Let
\(T(\rho)=\sum_\Pi\operatorname{tr}(\rho\Pi)\Pi\). An element
\(H=\sum_jh_j\Pi_j^{(n)}\) of \(V_n\) has \(\operatorname{tr}(H\Pi_j^{(n)})=h_j\) and
\(\operatorname{tr}(H\Pi)=\operatorname{tr}(H)/d=0\) for \(\Pi\) in another basis, so
\(T(H)=H\); and \(T(I)=(d+1)I\). Hence \(T(\rho)=\rho+\operatorname{tr}(\rho)I\) on
Hermitian matrices, and on all matrices by linearity. Without basis \(n\), a
nonzero \(H\in V_n\) has \(T(H)=0\ne H\). \(\square\)

When \(\rho\) is a density operator the numbers \(\operatorname{tr}(\rho\Pi)\) are the
\(d(d+1)\) outcome probabilities of the \(d+1\) bases, and in the language of
CH14-P02 the map \(\rho\mapsto(\operatorname{tr}(\rho\Pi))_\Pi\) is a linear
presentation that is injective on all of \(M_d(\mathbb C)\), with the explicit
inverse above. The projectors have entries in \(\mathbb Q(\omega)\) or
\(\mathbb Q(i)\), so for \(\rho\) with entries there the numbers and the
reconstruction stay there. The absorbed archived working-paper source 004 normalized the
dimension-3 vectors by \(1/\sqrt3\), which does not lie in
\(\mathbb Q(\omega)=\mathbb Q(\sqrt{-3})\), and did not list its dimension-4
bases. It cited an external Lean check of orthogonality and unbiasedness for the
unnormalized vectors used here; that check was read, not built or ported, and
gives no Lean credit, and the reconstruction identity had only an executable
test there. Informational completeness is a property of the declared bases; it
derives no measurement postulate (CH05-N01, CH19-O04).

Let \(N\) be odd, let \(\omega\), \(X\), and \(Z\) be as in CH05-P09, and let
\(Pe_j=e_{-j}\) with indices modulo \(N\). Then \(P^2=I\) and
\(\operatorname{tr}P=1\), because \(j\equiv-j\) only for \(j=0\) when \(N\) is odd;
for even \(N\) the trace is \(2\).

\texttt{WDP-CH05-P11} \textbf{Finite proposition (the phase-point frame).} For odd \(N\) and
\((q,k)\in\mathbb Z_N^2\) put \(A(q,k)=X^qZ^kPZ^{-k}X^{-q}\). Then
\(A(q,k)e_j=\omega^{2k(q-j)}e_{2q-j}\); each \(A(q,k)\) is Hermitian with
\(\operatorname{tr}A(q,k)=1\); \(\sum_{q,k}A(q,k)=NI_N\); and
\(\operatorname{tr}\big(A(q,k)A(q',k')\big)=N\delta_{qq'}\delta_{kk'}\).
Consequently, for every complex \(N\times N\) matrix \(\rho\) the function
\(W_\rho(q,k)=\operatorname{tr}(\rho A(q,k))/N\) satisfies

\[
\rho=\sum_{q,k}W_\rho(q,k)\,A(q,k),\qquad \sum_{q,k}W_\rho(q,k)=\operatorname{tr}\rho ,
\]

\(W_\rho\) is real when \(\rho\) is Hermitian, and \(W_\rho\) takes values in
\(\mathbb Q(\omega)\) when the entries of \(\rho\) do.

\emph{Proof.} \(X^{-q}e_j=e_{j-q}\); \(Z^{-k}\) multiplies it by \(\omega^{-k(j-q)}\); \(P\)
sends it to \(e_{q-j}\); \(Z^k\) multiplies by \(\omega^{k(q-j)}\); and \(X^q\) sends it
to \(e_{2q-j}\). This gives the formula. A diagonal entry at \(j\) needs
\(2(j-q)\equiv0\), that is \(j=q\) since \(N\) is odd, and equals \(1\). The entry at
\((j,2q-j)\) is the conjugate of the entry at \((2q-j,j)\), so \(A\) is Hermitian.
For fixed \(q\) and \(j\), \(\sum_k\omega^{2k(q-j)}\) is \(N\) when \(j=q\) and \(0\)
otherwise, because then \(\omega^{2(q-j)}\ne1\) is an \(N\)-th root of unity; hence
\(\sum_{q,k}A(q,k)=NI\). The product \(A(q,k)A(q',k')\) sends \(e_j\) to a multiple of
\(e_{2q-2q'+j}\), so its trace vanishes unless \(q=q'\), and then it is
\(\sum_j\omega^{2(k-k')(j-q)}=N\delta_{kk'}\). The \(N^2\) Hermitian operators are
therefore orthogonal for the pairing \(\operatorname{tr}(A^{*}B)\) and form a
basis of \(M_N(\mathbb C)\); the coefficient of \(A(q,k)\) in \(\rho\) is
\(\operatorname{tr}(\rho A(q,k))/N\), and summing the coefficients gives
\(\operatorname{tr}\rho\). For Hermitian \(\rho\) and \(A\),
\(\operatorname{tr}(\rho A)\) is real. \(\square\)

This frame is the standard odd-dimensional discrete phase space; it is proved
here directly, and no external theorem about it is used. The phase factor that
the absorbed archived working-paper source 005 placed in its displacement operators
cancels in \(A(q,k)\). That source established the frame for \(N\in\{3,5\}\) only by
executed tests in an external repository; the Lean it cites establishes \(P^2=I\) and
\(\operatorname{tr}P=1\) for \(N\in\{3,5\}\). It was read, not ported, and gives no
Lean credit.

\texttt{WDP-CH05-P12} \textbf{Finite proposition (a signed phase-point valuation).} Let \(N\ge3\)
be odd and \(J\) the all-ones \(N\times N\) matrix. The density operators
\(\rho_0=e_0e_0^{\mathsf T}\) and \(\rho_+=J/N\) have
\(W_{\rho_0}(q,k)=\delta_{q,0}/N\) and \(W_{\rho_+}(q,k)=\delta_{k,0}/N\), both
nonnegative. For \(m=e_1-e_{N-1}\), the density operator
\(\rho_M=\tfrac12mm^{\mathsf T}\), whose entries are \(0\) and \(\pm\tfrac12\), has
\(W_{\rho_M}(0,0)=-1/N\).

\emph{Proof.} \(W_{\rho_0}(q,k)=A(q,k)_{00}/N\), and the \((0,0)\) entry of \(A(q,k)\) is
nonzero only when \(2q\equiv0\), that is \(q=0\), where it is \(1\). Next,
\(W_{\rho_+}(q,k)=N^{-2}\sum_{i,j}A(q,k)_{ij}=N^{-2}\sum_j\omega^{2k(q-j)}
=\delta_{k,0}/N\). At the origin \(A(0,0)=P\), and \(Pm=e_{N-1}-e_1=-m\), so
\(\operatorname{tr}(\rho_MP)=\tfrac12m^{\mathsf T}Pm=-\tfrac12m^{\mathsf T}m=-1\).
\(\square\)

For \(N=3\) the other eight values of \(W_{\rho_M}\) equal \(\tfrac16\), and
\(-\tfrac13+8\cdot\tfrac16=1\). For \(N=5\) the values lie in \(\mathbb Q(\sqrt5)\):
on the line \(q=0\) they are \(-\tfrac15\), \(\tfrac{1+\sqrt5}{20}\),
\(\tfrac{1-\sqrt5}{20}\), \(\tfrac{1-\sqrt5}{20}\), \(\tfrac{1+\sqrt5}{20}\) for
\(k=0,\dots,4\), on the lines \(q=1\) and \(q=4\) they are \(\tfrac1{10}\), and on
\(q=2,3\) they are \(0\); three are negative and they sum to \(1\). So the frame
represents \(\rho_M\) by real weights on the finite set \(\mathbb Z_N^2\) with
total \(1\) and a negative entry: the kind of signed valuation that the remark
after the representation theorem of §18.2 separates from a probability state.
The Lean that the absorbed source cites establishes \(Pm=-m\),
\(\langle m,m\rangle=2\), and \(\langle m,Pm\rangle=-2\) for \(N\in\{3,5\}\) and
defines the value \(-1/N\) as a constant, as that source recorded; the
connecting step is the computation above. It was read, not ported, and gives no
Lean credit.

The two frames of this section place the negativity differently. The numbers
\(\operatorname{tr}(\rho\Pi)\) of CH05-P10 are nonnegative for every density
operator, but that reconstruction subtracts \(\operatorname{tr}(\rho)I\); the
phase-point reconstruction is a plain sum, but its weights can be negative.
This compares two declared frames and is not a theorem about all frames. A
negative phase-point value is not claimed to certify nonclassicality,
contextuality, or a computational resource; CH05-N01 and CH19-N02 apply.

\subsubsection{Technical note for Chapters 1--5: morphisms and verification records}\label{technical-note-for-chapters-15-morphisms-and-verification-records}

The foundational constructions become composable only after morphisms are
fixed. For admitted cuts, a morphism may consist of maps on tokens and
incidences preserving role types and the selected comparison family. For graph
cochains, a graph map can induce pullback or pushforward only under the usual
variance and coefficient hypotheses. For observer-facing use, the morphism
must also preserve the proposition witnessed; a map that forgets the relevant
edge can destroy a residual witness.

A minimal verification record for any construction in Part I should contain:

\begin{enumerate}
\def\labelenumi{\arabic{enumi}.}
\tightlist
\item
  the finite carriers and their explicit enumeration or source hash;
\item
  the role and comparison declarations;
\item
  the construction performed;
\item
  the proposition claimed of its output;
\item
  the verifier or proof and its assumptions;
\item
  the morphisms under which the proposition is invariant;
\item
  the known countermodel class;
\item
  an explicit statement of what the verification record does not establish.
\end{enumerate}

This verification record schema prevents a common migration error. Suppose a cycle
certificate is valid on graph \(G\), and a later document replaces \(G\) by a
quotient graph \(G'\). The old certificate does not automatically migrate. One
must show that the relevant circulation functional descends, that the target
cochain is transported, and that the repair image maps into the new repair
image. A matching filename or informal diagram is not a proof of that transport.

A useful, non-exhaustive checklist names what a presentation change can alter:
the repair space, the closure map, the carrier and its maps, the coefficient
field, the declared norms, and the differential itself (a perturbation, as in
§16.5). Each needs its own transport statement; a change on one axis does not
license a claim on another.

The same point sharpens the symmetry theorem. For an \(H\)-set \(Y\), an
\(H\)-equivariant map from the trivial one-point \(H\)-set to \(Y\) exists exactly
when \(Y\) has an \(H\)-fixed point. Thus every equivariant selection problem can be
reduced to a fixed-point question in the target action. The full symmetric group
on an unstructured carrier is one instance. If the input includes a partition,
the automorphism group shrinks and fixed targets may appear. The new output is
then relative to the partition.

\texttt{WDP-CH03-P02} \textbf{Finite proposition (fixed-target criterion).} Let a finite group
\(H\) act on a finite set \(Y\). A parameter-free \(H\)-equivariant choice of a point
of \(Y\) exists if and only if \(Y^H\), the set of global fixed points, is nonempty.

This proposition gives an exact audit recipe: identify the automorphism group of
the admitted input, identify its action on the proposed outputs, and compute the
fixed locus. Claims that a target is ``forced'' should specify this action and show
that the fixed locus has the required uniqueness property.

Typed witnesses also need morphisms. An adapter from a residual-locus witness to
a retention witness consists of a write map, a persistence law, and a readback
proof. An adapter from a repair witness to a dual statement may not exist: a
repair certifies membership in an image, whereas a separating functional establishes
nonmembership. They face opposite directions. An adapter from a data record to
a proposition requires a semantics and a validity predicate. These obligations
are small enough to state explicitly and large enough to block invalid prose
inferences.

Finally, the quotient defect admits a stability estimate. By the triangle
inequality,

\[
|\Phi_J([\delta])-\Phi_J([\delta'])|
\le J(\delta-\delta')
\]

whenever \(J\) is a norm and both defects live in the same quotient problem.
Indeed, any repair candidate for one defect supplies an upper bound for the
other after adding the norm of their difference; exchange the defects for the
reverse inequality. This one-Lipschitz property is useful when Chapter 13 adds
measurement uncertainty. It does not compare different graphs, repair spaces,
or gauges without transport data.

\begin{center}\rule{0.5\linewidth}{0.5pt}\end{center}

\subsection{Part II --- Order, quantity, and finite probability}\label{part-ii-order-quantity-and-finite-probability}

\subsection{Chapter 6 --- Recurrence, succession, and constitutive order}\label{chapter-6-recurrence-succession-and-constitutive-order}

\subsubsection{6.1 Recurrence is not yet succession}\label{recurrence-is-not-yet-succession}

If a retained mark can be accessed in two distinguishable exhibitions, one has
a recurrence pattern. Recurrence alone does not say which exhibition is
earlier. A pair of records can be compared symmetrically; order appears only
when an asymmetric relation is declared or constructed from additional data.

Let \(S\) be a finite set of retained states or events. A \textbf{succession candidate}
is a binary relation \(\prec\subseteq S\times S\). To function as a strict partial
order it must be irreflexive and transitive. Local transition data may instead
give a directed relation \(\leadsto\) that is not transitive; its transitive
closure can define \(\prec\) when cycles are absent.

\texttt{WDP-CH06-D01} \textbf{Definition (constitutive succession).} A constitutive
succession packet is

\[
\mathfrak S=(S,\leadsto,\prec,\pi),
\]

where \(\leadsto\) is an exhibited immediate-transition relation, \(\prec\) is a
declared or proved strict partial order containing \(\leadsto\), and \(\pi\)
records the assumptions under which that containment and order are valid.

This definition forces an often-hidden choice into view: whether order is
primitive, generated by acyclic reachability, or inferred by an observer from
records. Those are different models.

\subsubsection{6.2 Immediate succession and cycles}\label{immediate-succession-and-cycles}

If \(G=(S,E)\) is a finite directed acyclic graph, define \(x\prec y\) when there is
a nonempty directed path from \(x\) to \(y\). Then \(\prec\) is irreflexive and
transitive. Conversely, if \(G\) contains a directed cycle, reachability is not
irreflexive.

\texttt{WDP-CH06-P01} \textbf{Finite proposition.} Nonempty reachability on a finite
directed graph is a strict partial order exactly when the graph is acyclic.

\emph{Proof sketch.} Transitivity follows by concatenating paths. A directed cycle
gives \(x\prec x\). If \(x\prec x\), the witnessing nonempty path contains a
directed cycle. \(\square\)

The proposition is graph theory, not a derivation of physical time. It says
what an admitted transition graph must satisfy to support a strict order.

\subsubsection{6.3 No external clock in the definition}\label{no-external-clock-in-the-definition}

A transition record can be indexed for storage without treating the index as
physical duration. If a computation enumerates states \(s_0,s_1,\ldots\), the
enumeration order is part of the implementation. A foundational claim must not
silently transfer that order to the represented system.

The programme therefore distinguishes:

\begin{itemize}
\tightlist
\item
  \textbf{execution order}, the order in which an algorithm processes data;
\item
  \textbf{record order}, an order stored in a trace;
\item
  \textbf{constitutive order}, the relation asserted of the modeled events;
\item
  \textbf{coordinate time}, a numeric label assigned after a clock convention;
\item
  \textbf{duration}, an additive or otherwise composed magnitude on intervals.
\end{itemize}

Any bridge among them is a theorem or a declaration, not a synonym.

\subsubsection{6.4 Worked example: fork without total order}\label{worked-example-fork-without-total-order}

Let \(S=\{a,b,c\}\) with edges \(a\leadsto b\) and \(a\leadsto c\). Reachability gives
\(a\prec b\) and \(a\prec c\), while \(b\) and \(c\) are incomparable. The packet has
succession but no total order. A file format may list \(b\) before \(c\); that list
does not change the constitutive relation.

Adding an edge \(b\leadsto c\) makes one branch later than the other. Adding both
\(b\leadsto c\) and \(c\leadsto b\) destroys strictness by creating a cycle.

\subsubsection{6.5 The finite kernel and its boundary}\label{the-finite-kernel-and-its-boundary}

The finite order facts above are elementary and amenable to proof assistants.
The stronger programme targets are not automatic:

\begin{itemize}
\tightlist
\item
  constructing the event carrier without presupposing ordered inputs;
\item
  proving that a retained-difference mechanism selects one transition relation;
\item
  connecting the order to physical causal structure;
\item
  recovering continuum topology or metric;
\item
  explaining an experienced passage of time.
\end{itemize}

\texttt{WDP-CH06-O01} \textbf{Finite proposition (relational retention and succession).} A declared finite nonempty labelled transition system with exact reports or a validated report equivalence admits greatest strong bisimilarity and its minimal bounded quotient. Its raw directed relation has a strict transitive edge extension exactly when acyclic. Cycles obstruct such an extension; SCC condensation changes the carrier. A nonempty strict order cannot be chosen invariantly from a fully symmetric carrier. Directional transition or symmetry-breaking order data must therefore be declared. See Appendix J.

\emph{Equivariance Analysis and Mechanized Separation.}

\begin{enumerate}
\def\labelenumi{\arabic{enumi}.}
\tightlist
\item
  \textbf{Order Data Necessity via \(\operatorname{Sym}(M)\)-Equivariance:} Let \(M\)
  be a bare finite memory carrier with \(|M| \ge 2\), constant report, and no distinguished relational or ordering data.
  Its symmetry group is the full symmetric group \(G = \operatorname{Sym}(M)\).
  Any strict partial order \(\prec \ \subset M \times M\) must be irreflexive and
  asymmetric (\(a \prec b \implies \neg(b \prec a)\)). Any parameter-free
  construction from order-free retention must be \(G\)-equivariant. On the set of
  distinct pairs \(M \times M \setminus \Delta\), the action of \(G\) is transitive,
  because \(G\) acts 2-transitively on \(M\).
  The only \(G\)-invariant subsets of distinct pairs are \(\emptyset\) and
  \(M \times M \setminus \Delta\). The empty relation \(\emptyset\) is formally a
  trivial strict partial order, but contains no pairs and therefore provides no
  succession, while the full set of pairs is symmetric, violating asymmetry.
  Hence no nonempty (or nontrivial) \(G\)-equivariant strict partial order exists
  on an order-free carrier \(M\): to obtain any actual succession, additional order
  data (or symmetry-breaking selection) is mathematically necessary. While this
  symmetry argument establishes the necessity of order data in the parameter-free
  \(G\)-equivariant setting, Appendix J also classifies declared finite relational retention and strict succession extensions. It does not select an undeclared transition relation.
\item
  \textbf{Constitutive succession and immediate transitions:} On the three-cut
  carrier \texttt{Chain3} with immediate edges
  \(\alpha\leadsto\beta\leadsto\gamma\), nonempty reachability relates
  \(\alpha\) to \(\gamma\), although there is no immediate edge
  \(\alpha\leadsto\gamma\). Thus reachability and immediate transition differ
  on the same ordered pair \((\alpha,\gamma)\). Appendix D records the
  separate formalization evidence for
  \texttt{lean4/WitnessedDifference/Succession.lean}.
\end{enumerate}

\texttt{WDP-CH06-N01} \textbf{Nonclaim.} A finite acyclic relation is not a continuum
spacetime, a causal-set model of our universe {[}REF-TIME-001; REF-TIME-002{]}, or
an account of temporal experience.

\begin{center}\rule{0.5\linewidth}{0.5pt}\end{center}

\subsection{Chapter 7 --- Duration, exactness, periods, and clocks}\label{chapter-7-duration-exactness-periods-and-clocks}

This chapter merges the two duration drafts. Succession and duration are kept
separate: order answers \emph{which intervals compose}; duration assigns a magnitude
to those intervals.

\subsubsection{7.1 Edge valuation}\label{edge-valuation}

Let \(G=(S,E)\) be a finite directed acyclic graph or, more generally, a directed
graph on which a chosen path family is defined. A duration gauge is a map

\[
\lambda:E\longrightarrow M
\]

into a commutative additive monoid \((M,+,0)\). For a path
\(p=e_1\cdots e_k\), define

\[
D_\lambda(p)=\sum_{j=1}^{k}\lambda(e_j).
\]

\texttt{WDP-CH07-D01} \textbf{Definition (finite duration gauge).} The pair of a path system
and an additive edge valuation is a finite duration gauge. Positivity,
Archimedean structure, units, or real values are additional hypotheses.

\texttt{WDP-CH07-P01} \textbf{Finite proposition (concatenation additivity).} Whenever path
concatenation \(pq\) is defined,

\[
D_\lambda(pq)=D_\lambda(p)+D_\lambda(q).
\]

The proof is reassociation of a finite sum. This is the strongest unconditional
duration theorem available at this level.

\subsubsection{7.2 Exact, closed nonexact, active failure, undetermined}\label{exact-closed-nonexact-active-failure-undetermined}

When edge durations lie in an abelian group, ask whether
\(\lambda=d_0t\) for a vertex potential \(t\). Four statuses must not be collapsed.

\begin{enumerate}
\def\labelenumi{\arabic{enumi}.}
\tightlist
\item
  \textbf{Exact.} A displayed \(t\) satisfies \(d_0t=\lambda\).
\item
  \textbf{Closed nonexact.} A declared next differential gives
  \(d_1\lambda=0\), while a period or dual certificate establishes
  \(\lambda\notin\operatorname{im}d_0\).
\item
  \textbf{Active closure failure.} \(d_1\lambda\ne0\) on an exhibited \(2\)-cell.
\item
  \textbf{Exactness undetermined.} The available structure or search did not decide
  the preceding cases.
\end{enumerate}

``Closed nonexact'' is unavailable when no valid \(d_1\) exists. On a graph without
declared \(2\)-cells, loop periods can still refute exactness, but the word
\emph{closed} should be used only relative to a named complex.

\subsubsection{7.3 Path independence}\label{path-independence}

\texttt{WDP-CH07-P02} \textbf{Finite proposition.} If \(\lambda=d_0t\), then all directed
paths from \(a\) to \(b\) have the same duration \(t(b)-t(a)\).

\emph{Proof.} The edge differences telescope along each path. \(\square\)

Conversely, suppose a finite directed graph has a base vertex \(o\) from which every vertex is reachable, and sums agree along all directed paths from \(o\) to each vertex. Define \(t(v)\) by this common sum. Appending any edge \(u\to v\) to an \(o\)-to-\(u\) path gives \(t(v)=t(u)+\lambda(u,v)\), so \(d_0t=\lambda\). The converse needs this reachability and closure under edge extension; connectedness of the underlying undirected graph or an arbitrary selected path family is insufficient.

\subsubsection{7.4 Worked examples}\label{worked-examples}

\textbf{Chain.} On \(a\to b\to c\) let \(\lambda=(2,5)\). With
\(t(a)=0,t(b)=2,t(c)=7\), the valuation is exact and
\(D(a\to b\to c)=7\).

\textbf{Diamond.} Let \(a\to b\to d\) and \(a\to c\to d\). Give the upper path weights
\(2,5\) and lower path weights \(3,4\). Both total \(7\); taking
\(t=(0,2,3,7)\) shows exactness. Changing the last lower edge from \(4\) to \(6\)
makes the lower path total \(9\). The discrepancy \(2\) is a period witness around
the undirected cycle and refutes exactness. With unit weights it also bounds
the quotient defect from below (CH04-P11), and here \(\Phi_1=2\) exactly: the
potential \(t=(0,2,3,7)\) leaves the single residual \(2\) on the last lower edge.

\textbf{Unit tally.} Assign every edge value \(1\). Path duration becomes edge count.
This is a valid declared gauge, but two refinements of the same modeled process
can change the count. Unit tally therefore requires a granularity convention.
It is additive under concatenation, yet it is not exact on a graph with a
directed cycle: a directed cycle of length \(n\ge1\) has unit-tally circulation
\(n\ne0\), so no vertex potential produces it.

\textbf{Commuting square.} Let two routes of two steps each join the same endpoints,
one carrying values \((1,1)\) and the other \((1,0)\). The valuation is additive
along each route, but the circulation around the square is \(1\), so it is not
exact. Counting steps cannot tell the routes apart, since both have two items.

\subsubsection{7.5 Clock packets}\label{clock-packets}

A clock is not just a duration function. It is a comparator between a target
process and a repeatable reference process.

\texttt{WDP-CH07-D02} \textbf{Definition (finite clock packet).} A clock packet is

\[
\mathfrak T=(C,\leadsto_C,\lambda_C,r_C,\kappa,\varepsilon),
\]

where \(C\) is a reference-state carrier, \(\leadsto_C\) its transition relation,
\(\lambda_C\) its duration gauge, \(r_C\) a readout, \(\kappa\) a declared coupling or
comparison with target intervals, and \(\varepsilon\) an error or tolerance
model.

A clock is \textbf{adequate for a target family} only relative to stated criteria:
coverage, stability, resolution, monotonicity, and error bounds. No clock is
adequate just because it ticks.

\subsubsection{7.6 Calibration joins duration to measurement}\label{calibration-joins-duration-to-measurement}

Calibration supplies a relation between an instrument readout and a reference
quantity. In this monograph, duration becomes a measurement quantity only when
a clock packet is attached to the measurement architecture of Chapter 13. The
edge valuation alone is a mathematical magnitude; the calibrated result
requires a measurand, procedure, unit, and uncertainty statement.

\texttt{WDP-CH07-O01} \textbf{Finite proposition (resolved finite target).} For a positive calibrated finite occurrence chain whose edge residuals are bounded by nonnegative declared budgets, every interval residual is bounded by the sum of those budgets, and endpoint-preserving block-sum refinement preserves the bound.
The proof, native ownership and scope are given in Appendix I.

\texttt{WDP-CH07-N01} \textbf{Nonclaim.} Additive finite duration does not select a global unrestricted
unit, a continuum parameter, Lorentzian geometry, or a physical arrow of time.

\begin{center}\rule{0.5\linewidth}{0.5pt}\end{center}

\subsection{Chapter 8 --- Multiplicity, additive weight, and finite probability}\label{chapter-8-multiplicity-additive-weight-and-finite-probability}

\subsubsection{8.1 Matching before number}\label{matching-before-number}

Let \(A\) and \(B\) be finite exhibited families. A \textbf{matching} is a finite set of
pairing incidences satisfying:

\begin{enumerate}
\def\labelenumi{\arabic{enumi}.}
\tightlist
\item
  no endpoint participates in two pairings;
\item
  every endpoint in both families participates in a pairing;
\item
  the pairing preserves the comparisons declared relevant.
\end{enumerate}

These conditions describe a bijective correspondence without taking a number
object as primitive. They still presuppose the exhibited families and a frozen
comparison scope.

\texttt{WDP-CH08-D01} \textbf{Definition (comparison-relative congruence).} \(A\) and \(B\) are
congruent relative to \(\mathcal K\) when a matching preserving \(\mathcal K\) is
exhibited.

This construction does not show that plurality itself is derived from one cut.
It describes how finite multiplicity can be compared after plural exhibited
tokens are available.

\subsubsection{8.2 Witness mass}\label{witness-mass}

Let \(W=\{w_1,\ldots,w_n\}\) be a finite family of typed witnesses and let
\(m:W\to\mathbb R_{\ge0}\) be a declared mass. For \(A\subseteq W\), define

\[
\mu(A)=\sum_{w\in A}m(w).
\]

For disjoint \(A,B\),

\[
\mu(A\sqcup B)=\mu(A)+\mu(B).
\]

This is finite additivity by construction. If \(m(w)=1\), \(\mu(A)\) is a tally.
If witnesses are indistinguishable under a declared transitive symmetry and
\(m\) is required to be equivariant, equal weights follow on that orbit. The
equivariance requirement is the load-bearing assumption; indifference is not
obtained from ignorance alone.

\subsubsection{8.3 Normalized valuation}\label{normalized-valuation}

Let \(\mathcal P=\{A_1,\ldots,A_k\}\) be a declared finite partition of \(W\) and
assume \(\mu(W)>0\). Define

\[
p(A_i)=\frac{\mu(A_i)}{\mu(W)}.
\]

\texttt{WDP-CH08-P01} \textbf{Finite proposition.} The values \(p(A_i)\) are nonnegative and
sum to \(1\).

The proof is finite arithmetic. The result is a normalized additive valuation
on a declared finite partition. Calling it probability is mathematically
standard in a finite sample model, but its physical interpretation remains
open.

\subsubsection{8.4 Refinement}\label{refinement}

Suppose a cell \(A_i\) is refined into disjoint cells
\(B_{i1},\ldots,B_{ir}\). Finite additivity gives

\[
p(A_i)=\sum_jp(B_{ij}).
\]

If the mass assigned to a witness changes under refinement, this equality can
fail because a new model has been introduced. A refinement theorem therefore
requires a transport rule preserving total mass.

\texttt{WDP-CH08-O01} \textbf{Finite proposition (resolved finite target).} Every nonnegative additive valuation on a finite Boolean event algebra has unique algebra-atom masses and measurable-map pushforward; nonnegative masses on individual witnesses extending it are unique exactly when every nonsingleton algebra atom has mass zero.
The proof, native ownership and scope are given in Appendix I.

This conditional result is load-bearing for any assertion that \(L^1\) aggregation is
forced rather than just chosen.

\subsubsection{\texorpdfstring{8.5 Why \(L^1\) appears---and what it does not exclude}{8.5 Why L\^{}1 appears---and what it does not exclude}}\label{why-l1-appearsand-what-it-does-not-exclude}

For nonnegative vectors with disjoint support,

\[
\|x+y\|_1=\|x\|_1+\|y\|_1.
\]

For \(p>1\), \(\|x+y\|_p\) is generally not the sum of the norms. Thus \(L^1\) is the
natural norm when the magnitude itself must be degree-one homogeneous and
exactly additive over disjoint positive sectors.

This does not ban \(L^2\). Squared \(L^2\) magnitude is additive on orthogonal
components:

\[
\|x\oplus y\|_2^2=\|x\|_2^2+\|y\|_2^2.
\]

The two expressions answer different questions. Moving from witness mass to
amplitude geometry changes the object and its rules.

The primitive-axiomatics source isolates the hypotheses under which an
aggregator of edge magnitudes must be \(L^1\). Let \(\mathcal A\) send finite
families \((a_e)_{e\in E}\) of nonnegative edge magnitudes to \([0,\infty)\).

With additivity over disjoint edge sets assumed outright, the conclusion is
finite. If \(\mathcal A\) is edge-decomposable, additive over disjoint edge sets,
positively homogeneous on each edge, invariant under relabeling edges, and
faithful (\(\mathcal A_{\{e\}}(1)>0\)), then additivity over the singleton
decomposition gives \(\mathcal A((a_e))=\sum_e\mathcal A_{\{e\}}(a_e)\),
homogeneity gives \(\mathcal A_{\{e\}}(a_e)=a_e\mathcal A_{\{e\}}(1)\),
relabeling invariance makes \(\mathcal A_{\{e\}}(1)=c\) independent of \(e\), and
faithfulness gives \(c>0\). So \(\mathcal A((a_e))=c\sum_ea_e\). The conditional
part is the replacement of additivity by the weaker replica extensivity.

\texttt{WDP-CH08-C01} \textbf{Conditional result (\(L^1\) aggregation from replica
extensivity).} For every finite list \(a\) of nonnegative reals, let
\(\mathcal A(a)\ge0\). Assume \(\mathcal A(\varnothing)=0\),
\(\mathcal A((t))=ct\) for a fixed \(c>0\), and a continuous associative operation
\(F:[0,\infty)^2\to[0,\infty)\) with identity \(0\) such that
\[
\mathcal A(a\sqcup b)=F(\mathcal A(a),\mathcal A(b)).
\]
Assume replica extensivity: \(\mathcal A(a^{\sqcup N})=N\mathcal A(a)\) for
positive integers \(N\). Then \(F(x,y)=x+y\) and
\(\mathcal A((a_e))=c\sum_e a_e\).

\emph{Proof.} For \(t\ge0\), a singleton of value \(t/c\) has aggregate \(t\).
Replicas and the composition law give \(F(mt,nt)=(m+n)t\) for nonnegative
integers \(m,n\); the cases with a zero argument use the identity. For any
nonnegative rational \(x,y\), choose a common denominator \(N\) and put
\(t=1/N\) to obtain \(F(x,y)=x+y\). Continuity and density of rational pairs
extend this equality to all nonnegative real \(x,y\). The composition law
then gives the claimed sum by induction on list length. \(\square\)

The singleton condition follows from positive homogeneity, relabeling
invariance and faithfulness on singleton inputs. The explicit composition
law is the meaning of edge decomposability used here. This finite argument
does not require Aczél's representation theorem {[}REF-MATH-006{]}.

The hypotheses are not idle. Bottleneck (max), density (mean), saturating
(capped-sum), and \(L^p\) aggregators with \(p>1\) fail replica extensivity, as do
typical aggregators with cross-edge coupling. The sum of squares
\(\sum_ea_e^2\) fails only homogeneity; a weighted sum \(\sum_ec_ea_e\) with
unequal positive \(c_e\) fails only relabeling invariance; and the zero
aggregator fails only faithfulness. Edge-decomposability is the format in which
the others are stated. The result selects \(L^1\) for evaluators that require independent
defects to add. It does not make \(L^1\) globally unrestricted, and the classification of
gauges by edge symmetry monoids belongs to a separate source that is not
absorbed here.

Identifying this nonnegative-magnitude aggregator with a norm on signed cochains additionally assumes dependence only on coordinate magnitudes (independent sign invariance), or additivity on every disjoint-support pair. Additivity restricted to the positive cone alone does not classify signed norms.

\subsubsection{8.6 Worked example: a declared partition}\label{worked-example-a-declared-partition}

Let \(W=\{w_1,w_2,w_3,w_4\}\) with masses \((1,1,2,4)\). Partition it into
\(A=\{w_1,w_2\}\) and \(B=\{w_3,w_4\}\). Then

\[
\mu(W)=8,\qquad p(A)=\frac14,\qquad p(B)=\frac34.
\]

Refining \(B\) gives probabilities \(1/4\) and \(1/2\), which sum to \(3/4\).
Nothing in this calculation says the witnesses were random trials or that the
numbers predict frequencies. Those readings require a sampling or physical
model.

\begin{center}\rule{0.5\linewidth}{0.5pt}\end{center}

\subsection{Chapter 9 --- Probability scope fences}\label{chapter-9-probability-scope-fences}

\subsubsection{9.1 Five distinct notions}\label{five-distinct-notions}

The following must not be collapsed:

\begin{enumerate}
\def\labelenumi{\arabic{enumi}.}
\tightlist
\item
  a finite normalized valuation on an exhibited partition;
\item
  a limiting frequency in a repeatable trial model;
\item
  a propensity or physical chance;
\item
  a Bayesian degree of belief;
\item
  a countably additive measure on a sigma-algebra.
\end{enumerate}

Chapter 8 constructs only the first. A finite valuation can be embedded into
some of the others after additional data are supplied, but there is no unique
embedding in general.

\subsubsection{9.2 Normalization does not select the Born rule}\label{normalization-does-not-select-the-born-rule}

Let a nonzero finite amplitude vector be
\(z=(z_1,\ldots,z_n)\). For any \(q>0\), define

\[
p_i^{(q)}=\frac{|z_i|^q}{\sum_j|z_j|^q}.
\]

Then \(p_i^{(q)}\ge0\) and \(\sum_i p_i^{(q)}=1\). Normalization alone therefore
does not select \(q=2\).

\texttt{WDP-CH09-P01} \textbf{Finite proposition (normalization nonselection).} On any
amplitude vector with at least two unequal nonzero magnitudes, the normalized
vectors for exponents \(q,r>0\) with \(q\ne r\) are distinct, while each is a
probability vector.

\emph{Proof.} For two nonzero coordinates of unequal magnitude, their normalized
coordinate ratio is \((|z_i|/|z_j|)^q\). Because the positive base differs from
\(1\), equality of the ratios forces \(q=r\). \(\square\)

For example, for magnitudes \((1,2)\), \(q=1\) gives \((1/3,2/3)\) and \(q=2\) gives
\((1/5,4/5)\). Any derivation of the Born rule must use structure beyond positivity
and normalization---for example, noncontextuality and Hilbert-space assumptions
of a Gleason-type theorem {[}REF-PROB-002{]}, or operational axioms in a
reconstruction {[}REF-QM-003; REF-QM-004; REF-QM-006{]}.

The reverse direction loses information too. \textbf{Finite proposition (mass-to-amplitude
embedding is many-to-one).} Given nonnegative masses \(\mu_j\), not all zero, and
separately declared phases \(\varphi_j\), put \(p_j=\mu_j/\sum_k\mu_k\) and
\(\psi_j=\sqrt{p_j}\,e^{i\varphi_j}\). Then \(\sum_j|\psi_j|^2=1\). The map
\((\mu,\varphi)\mapsto\psi\) forgets the overall scale of \(\mu\), and the phases
are inputs rather than outputs. If the masses come from a signed cochain through
absolute values, the signs are lost as well. \emph{Proof.} Unit norm is the
normalization; \(\mu\) and \(c\mu\) for \(c>0\) give the same \(\psi\), and changing
\(\varphi\) changes \(\psi\) without changing \(\mu\). \(\square\) Reading \(|\psi_j|^2\)
back as a probability returns the declared \(p_j\) by construction. That is a
consistency check of the embedding, not a derivation of the Born rule.

\subsubsection{9.3 Finite additivity is not countable additivity}\label{finite-additivity-is-not-countable-additivity}

On a finite algebra, finite and countable additivity coincide trivially because
only finitely many disjoint nonempty cells exist. This does not provide a
sigma-algebra over an infinite outcome space, regularity, convergence, or a
continuum measure. Those are separate imports.

\subsubsection{9.4 Probability and observers}\label{probability-and-observers}

A partition may be selected by an observer profile: which distinctions the
observer can report, which outcomes are grouped, and which events are outside
scope. The resulting valuation is then profile-relative. This does not make it
subjective in a psychological sense; it makes the dependence on a partition
map explicit.

Later chapters will require an observer to publish the partition and value map
used in a measurement. Without that declaration, two normalized tables may
disagree because they answer different questions.

\texttt{WDP-CH09-N01} \textbf{Nonclaim.} The finite witness valuation is not a derivation of
chance, frequency convergence, exchangeability, a continuum probability space,
or the Born rule.

\texttt{WDP-CH09-N02} \textbf{Nonclaim.} Equal weights under a declared symmetry do not
establish an unrestricted principle of indifference.

\subsubsection{9.5 Order-preserving maps and temporal coarse-graining}\label{order-preserving-maps-and-temporal-coarse-graining}

Two succession packets can describe the same finite events at different
resolutions. Let \((S,\prec_S)\) and \((T,\prec_T)\) be strict partial orders. A map
\(f:S\to T\) is monotone when \(x\prec_S y\) implies either
\(f(x)=f(y)\) or \(f(x)\prec_T f(y)\), depending on whether the target convention
permits collapsed intervals. If equality is forbidden, use strict monotonicity.

Collapsing a chain segment can preserve causal order while changing unit-tally
duration. For example, \(a\prec b\prec c\) may map to
\(f(a)=u,f(b)=v,f(c)=v\). The image remembers that \(u\) precedes \(v\) but erases the
second step. A duration transport theorem must therefore say whether collapsed
edges contribute zero, are summed into a coarse edge, or make the map
inadmissible.

\texttt{WDP-CH06-D02} \textbf{Definition (duration-respecting coarse-graining).} Given
duration packets \((S,\lambda_S)\) and \((T,\lambda_T)\), a monotone map \(f\) is
duration respecting relative to an edge/path assignment \(F\) when every source
path \(p\) satisfies

\[
D_{\lambda_T}(F(p))=D_{\lambda_S}(p)
\]

or lies within a separately stated error bound. Such preservation is extra
data; order preservation alone does not imply it.

\subsubsection{9.6 Duration uncertainty}\label{duration-uncertainty}

Suppose edge durations are not exact values but intervals
\(I_e=[\ell_e,u_e]\). For a finite path \(p\), interval addition yields

\[
I_p=\left[\sum_{e\in p}\ell_e,\sum_{e\in p}u_e\right].
\]

This is a conservative bound when edge values can vary independently inside
their stated hard intervals. Correlation constraints can only preserve or
narrow that Minkowski bound because they select a subset of the Cartesian
product. A wider set would require changing the interval semantics---for example,
treating the inputs as marginal confidence sets rather than hard bounds---and
must be represented as a different model. If two paths from \(a\) to \(b\) have disjoint duration
intervals, exact path independence is refuted robustly. If their intervals
overlap, exactness is not established; the data are just inconclusive at that
resolution.

A clock comparison adds its own uncertainty. If a target interval maps to \(k\)
reference ticks, and each tick has interval \(I\), naive interval addition gives
\(kI\). Drift, shared calibration error, and quantization may be correlated and
cannot always be added as independent contributions. Chapter 13's uncertainty
packet is the correct place to state those rules.

\texttt{WDP-CH07-P03} \textbf{Finite proposition (interval accumulation).} Under the model
that each edge value can be selected independently from its declared interval,
the set of possible finite path sums is contained in the Minkowski sum of the
edge intervals, which for real intervals is the displayed endpoint sum.

\subsubsection{9.7 Orbit weights}\label{orbit-weights}

Let a finite group \(H\) act on a witness family \(W\). If a mass function
\(m:W\to\mathbb R_{\ge0}\) is \(H\)-invariant, then it is constant on each orbit.
The converse also holds: assigning one nonnegative mass to each orbit defines an
invariant mass.

\texttt{WDP-CH08-P02} \textbf{Finite proposition (invariant mass classification).}
\(H\)-invariant nonnegative mass functions on \(W\) are in bijection with
nonnegative functions on the orbit set \(W/H\).

\emph{Proof.} Invariance forces equal values on orbit-related witnesses. Conversely,
pulling a function back along the orbit projection is invariant. \(\square\)

This theorem states exactly what symmetry buys. It does not select the mass of
one orbit relative to another. If the action is transitive and total mass is
fixed, equal atomic mass follows uniquely. If there are several orbits, their
relative weights remain free. A statement of ``equal probability by symmetry''
must therefore name the group, action, orbit structure, and total-mass
normalization.

\subsubsection{9.8 Coarse partitions and sufficient reports}\label{coarse-partitions-and-sufficient-reports}

An observer often sees only a partition of \(W\). Let
\(q:W\to P\) assign each witness to a report class. The pushforward mass is

\[
(q_*\mu)(p)=\sum_{w:q(w)=p}m(w).
\]

Normalization commutes with this finite pushforward when total mass is
preserved. If two fine-grained models have the same pushforward, no evaluator
using only \(P\) can distinguish them. This is the probability analogue of the
presentation blind spot in Chapter 14.

A report is sufficient for a particular query when the query factors through
\(q\). Sufficiency is query-relative. A partition sufficient for predicting one
outcome can erase information needed for another.

\subsubsection{9.9 Frequencies, exchangeability, and updating}\label{frequencies-exchangeability-and-updating}

Suppose a protocol repeats a finite measurement \(N\) times and records outcome
counts \(n_i\). The empirical frequency \(n_i/N\) is a finite normalized valuation
on the observed trial multiset. It is a fact about those records. A law of large
numbers relates such frequencies to a probability parameter only after a
stochastic model---often independence or a weaker condition---is supplied.

Exchangeability is another substantive assumption: the joint assignment is
invariant under permutation of trial labels. Finite exchangeability by itself
does not provide an infinite de Finetti representation {[}REF-PROB-005{]}, and an
observer's choice to ignore order does not prove the physical process
exchangeable.

Bayesian updating likewise begins with a prior and likelihood model. The finite
obstruction calculus can record those objects and audit their transformations;
it does not select a uniquely rational prior from symmetry alone. A maximum-
entropy rule {[}REF-PROB-006{]} adds a constraint set, reference measure, and
optimization principle.

\subsubsection{9.10 Probability-to-measurement interface checklist}\label{probability-to-measurement-interface-checklist}

Before a finite valuation is used as a measurement probability, record:

\begin{enumerate}
\def\labelenumi{\arabic{enumi}.}
\tightlist
\item
  the trial or instrument protocol;
\item
  the outcome partition and its version;
\item
  whether weights are counts, modeled chances, likelihoods, or decision
  weights;
\item
  the sampling or exchangeability assumptions;
\item
  missing-data and stopping rules;
\item
  calibration and detection-efficiency treatment;
\item
  the uncertainty statement;
\item
  the inference rule and its scope.
\end{enumerate}

Without these fields, a normalized table remains valid arithmetic but has no
unique empirical reading.

\begin{center}\rule{0.5\linewidth}{0.5pt}\end{center}

\subsection{Part III --- Observer, observation, and measurement}\label{part-iii-observer-observation-and-measurement}

\subsection{Chapter 10 --- A layered observer architecture}\label{chapter-10-a-layered-observer-architecture}

\subsubsection{10.1 Why one tuple is not enough}\label{why-one-tuple-is-not-enough}

The source papers contain a quotient observer, a nine-component evaluator, a
six-component source packet, a five-component operational packet, and a
self-witnessing extension. These objects overlap, but they are not definitionally
equal. Attempting to compress them into one long tuple hides which components
are mathematical presentation, temporal memory, evaluation, policy, or
institutional authority.

We instead use four layers.

\texttt{WDP-CH10-D01} \textbf{Definition (observer presentation).} An observer presentation
is

\[
\mathsf{Pres}_{\mathcal O}=(X_{\mathcal O},P_{\mathcal O},q_{\mathcal O},
\mathcal C_{\mathcal O}),
\]

where \(X_{\mathcal O}\) is the admitted input carrier,
\(P_{\mathcal O}\) is the presentation or report carrier,
\(q_{\mathcal O}:X_{\mathcal O}\to P_{\mathcal O}\) is a presentation map, and
\(\mathcal C_{\mathcal O}\) is a declared context/profile family.

\texttt{WDP-CH10-D02} \textbf{Definition (observer memory).} An observer memory is a
retention packet

\[
\mathsf{Mem}_{\mathcal O}=(M_{\mathcal O},u_{\mathcal O},r_{\mathcal O},
\approx_{\mathcal O})
\]

as in Chapter 3, with a typed interface from presentations into memory and from
memory into reports.

\texttt{WDP-CH10-D03} \textbf{Definition (observer evaluator).} An observer evaluator is

\[
\mathsf{Eval}_{\mathcal O}=(D_{\mathcal O},A_{\mathcal O},d_{\mathcal O},
J_{\mathcal O},\Gamma_{\mathcal O},T_{\mathcal O}),
\]

consisting of a defect carrier, admissible repair carrier and map, declared
cost, any well-typed closure tests, and a terminalizer returning a status in a
declared finite status set.

\texttt{WDP-CH10-D04} \textbf{Definition (observer policy).} An observer policy specifies
what the system does with evaluator outputs: store, request more data, emit a
report, refuse an ill-typed query, or take a domain-specific action. Policy is
not authority.

The layered observer is

\[
\mathcal O=(\mathsf{Pres}_{\mathcal O},\mathsf{Mem}_{\mathcal O},
\mathsf{Eval}_{\mathcal O},\mathsf{Policy}_{\mathcal O})
\]

together with explicit adapters among layers.

\textbf{Interpretive remark (what an observer is, and what a perspective is).} In
this monograph an observer is not a person, a sensor or a conscious state. It
is the declared packet that makes a claim evaluable: what is compared, what
counts as defect, which repairs are admissible, how the residual is measured,
and which terminal statuses exist. Two such packets applied to the same objects
can return different verdicts without contradiction, because they evaluate
different declared questions. A \emph{perspective}, in the same vocabulary, is
available only once such a packet is declared. A viewpoint, camera angle or
opinion becomes a perspective when it can be terminalized by a stated
evaluator. A human or a test harness can operate the packet; neither is the
packet. One reading of this picture must be corrected. It is sometimes said
that the model proposes and the evaluator authorizes. Under CH10-D04 and
CH12-P01 the evaluator returns a status, and a policy acts on it. Authorization
is a separate record supplied by an authority packet (Chapter 15), so the
correct slogan is that the model proposes, the evaluator evaluates, and an
independently declared regime may authorize.

\subsubsection{10.2 Quotient evaluators}\label{quotient-evaluators}

The quotient construction of Chapter 4 gives a useful evaluator. Fix its
repair image \(B_{\mathcal O}=\operatorname{im}d_{\mathcal O}\) and let the
presentation expose \([\delta]\in Q^1_{\mathcal O}=C^1/B_{\mathcal O}\).
The evaluator computes

\[
\Phi_{\mathcal O}([\delta])
=\inf_{\alpha\in A_{\mathcal O}}
J_{\mathcal O}(\delta-d_{\mathcal O}\alpha).
\]

If the presentation instead retains an existing quotient \(C^1/B\), this
formula descends to that quotient exactly when \(B\subseteq B_{\mathcal O}\)
in the faithful finite-dimensional norm setting. Sufficiency follows by
translating the repair variable; necessity follows by comparing the zero
representative with each \(b\in B\), since zero distance to the closed repair
image means membership. Equality of repair images retains exactly the same
zero classes. A strict inclusion gives a seminorm on the old quotient,
with kernel \(B_{\mathcal O}/B\); those extra zero classes must be declared.

For example, take \(C^1=\mathbb R^2\), \(B=\operatorname{span}(e_0)\) and
\(B_{\mathcal O}=\operatorname{span}(e_1)\) with the \(\ell^1\) norm. The old
quotient identifies \(0\) and \(e_0\), but the new formula gives values 0 and 1.
This mismatched profile is refused before numerical evaluation; it does not
define an evaluator on the old classes. The
\href{code/witnessed_difference/quotient_observer.py}{bound profile evaluator}
checks the exact repair-image inclusion before invoking the rational solver.

Under the finite-dimensional real or complex normed-space hypotheses of
Chapter 4, the infimum is attained. Rational input data and a rational
polyhedral norm likewise give a rational optimizer through a feasible bounded
rational linear programme, as for the weighted \(L^1\) problem of §4.3.
Attainment over rational repairs can fail for other norms, including
polyhedral norms with irrational coefficients (§4.2). A terminalizer may
return \texttt{Accept} when the value is below a declared tolerance, \texttt{Reject} when a certified
lower bound exceeds it, and \texttt{Undetermined} otherwise. The tolerance, repair family,
and cost are observer-profile data.

This is evaluation, not perception in the phenomenological sense. It is also
not proposal generation: the evaluator compares a supplied candidate against a
declared profile.

\texttt{WDP-CH10-N01} \textbf{Nonclaim.} A tolerance status belongs to the evaluator, not
to the class. A class whose quotient residual lies below a declared tolerance is
not thereby exact, absent, infinite, or continuous, and changing the tolerance
changes the status while leaving the class, its residual, and its certificate
unchanged.

For the cyclic field \(\delta=(2,2,2)\) on the hollow triangle, \(\Phi_1=6\) with
repair residual \((0,0,6)\) and circulation \(y=(1,1,1)\). A terminalizer with
tolerance \(10\) returns \texttt{Accept}. With tolerance \(1\) the certified lower bound \(6\)
exceeds the tolerance and it returns \texttt{Reject}. An exact class is accepted at
every positive tolerance, so reading ``below tolerance'' as ``infinite'' or
``continuous'' would give that name to every exact class. In the terms of §21.9,
moving the tolerance is a profile change, not an evidential update.

\subsubsection{10.3 Adapters for legacy packets}\label{adapters-for-legacy-packets}

An adapter from a legacy packet \(L\) to the layered observer is a partial map

\[
\operatorname{adapt}_L:L\rightharpoonup\mathcal O
\]

with a domain predicate listing the required fields and laws. Typical mappings
are:

\begin{itemize}
\tightlist
\item
  a quotient observer supplies \texttt{ObserverPresentation}; an evaluator requires its own test or certificate;
\item
  a report map and context family map into \texttt{ObserverPresentation};
\item
  update and re-access maps map into \texttt{ObserverMemory};
\item
  a mathematical terminal status belongs to \texttt{ObserverEvaluator}; a separately declared status-to-action map supplies \texttt{ObserverPolicy};
\item
  an authorization field does \textbf{not} map into any of these layers and must be
  carried separately.
\end{itemize}

\texttt{WDP-CH10-O01} \textbf{Conditional result (source-bound partial adapter completeness).}
Every identified observer family in the selected source corpus admits an
explicit typed partial adapter preserving its declared maps, relations,
undefined fields and separate grammar/terminal interfaces. For surjective
sort maps, the target interpretation exists uniquely exactly when mapped
operation results and predicate truth are constant on input fibers. Missing
semantics are recoverable exactly when their own visible-fiber criterion
holds. Appendix K supplies the complete criterion, source-specific adapters
and independent necessary-hypothesis countermodels. All 44 original texts
and all 529 absorbed source occurrences are verified. Arbitrary history
functions require their actual declared semantics; finite transducer and
rational certificate verification retain their explicit executable domains.

\subsubsection{10.4 Worked observer}\label{worked-observer}

Let \(X=\mathbb R^3\) represent three sensor readings. Let the presentation be
\(q(x_1,x_2,x_3)=(x_1-x_2,x_2-x_3)\). Let memory retain the previous presented
pair. Write the current presented pair as \((u,v)\). With diagonal repair
\(d(a)=(a,a)\) and \(L^1\) cost, the evaluator computes
\(\inf_a(|u-a|+|v-a|)=|u-v|\).
The value is zero exactly when the two reported differences are equal.
Subtracting a common offset cannot change their discrepancy. A policy may
emit \texttt{stable}, \texttt{unstable}, or \texttt{insufficient-history}.

Every component is visible. Nothing about this construction authorizes a
decision outside the policy's declared domain, and nothing makes the system
conscious.

\subsubsection{10.5 The quotient observer as an adapter}\label{the-quotient-observer-as-an-adapter}

The primitive-axiomatics source defines an observer as a repair-compatible
quotient of \(Q^1\). That legacy packet fits the layered architecture through one
adapter, and the fit is exact enough to state.

\texttt{WDP-CH10-D06} \textbf{Definition (quotient-observer adapter).} A quotient observer
on a window (§4.7) is a linear map \(\mathcal O:Q^1\to Q^1_{\mathcal O}\),
equivalently a linear report map \(R:C^1\to\mathbb F^k\) with \(Rd_0=0\). It enters
the layered observer only as a presentation report map and as an evaluator
cost, the induced norm on the \textbf{realized report space} \(\operatorname{im}R\),
\(\Phi_{\mathcal O}(y)=\inf\{\Phi_1(r):r\in Q^1,\ \mathcal O(r)=y\}\). It
supplies no memory, policy, or authority layer, and requests for those layers
are refused.

The ambient codomain may contain unrealizable report values. For \(R=0\) and
\(y=1\) the fiber is empty, so the displayed expression is not a finite norm
value there. The evaluator refuses reports outside \(\operatorname{im}R\)
instead of issuing an ordinary numerical acceptance. On the realized
space, the real-extension norm hypothesis gives the stated norm and
passivity properties. Attainment holds over the finite-dimensional real or
complex carrier, and for rational queries with rational polyhedral norms;
a general rational norm restriction need not attain its infimum (§4.2).

Because \(\Phi_1(r)\) is itself an infimum over representatives,
\(\Phi_{\mathcal O}(y)=\inf\{J(h):h\in C^1,\ Rh=y\}\). For weighted \(L^1\) this is
a finite linear program, which the executable package solves exactly.

\texttt{WDP-CH10-P01} \textbf{Finite proposition (observer passivity).} For a quotient
observer, \(\Phi_{\mathcal O}(\mathcal O(q))\le\Phi_1(q)\) for every \(q\in Q^1\).
Hence an observer cannot report positive obstruction for an exact class, and
states identified by the core profile remain identified by every quotient
observer.

\emph{Proof.} The class \(q\) lies in the fiber over \(\mathcal O(q)\), and an infimum is
at most any member of its set. If \(\Phi_1(q)=0\), then
\(0\le\Phi_{\mathcal O}(\mathcal O(q))\le0\). If a core profile \(p\) has
\(p(x)=p(y)\), then \(\mathcal O(p(x))=\mathcal O(p(y))\) because \(\mathcal O\)
is a function; linearity is not needed. \(\square\)

Two consequences are worth naming. If an observer sends an exposed context
residual (§11.5) to zero, it carries no context witness for that anomaly. Two
reports \(q_A,q_B\) from observers \(\mathcal O_A,\mathcal O_B\) share a core
preimage exactly when \(\mathcal O_A^{-1}(q_A)\cap\mathcal O_B^{-1}(q_B)\) is
nonempty; an empty fiber means the reports cannot describe one core class.
This adapter covers the repair-compatible quotient family. Appendix K supplies
the other source-specific partial adapters and completes CH10-O01.

\texttt{WDP-CH10-P02} \textbf{Finite proposition (declared observer families).} For a
declared finite family of report maps, set \(\mathcal O\preceq\mathcal P\) when
\(\ker\mathcal O\subseteq\ker\mathcal P\), that is, when \(\mathcal P\) identifies
whatever \(\mathcal O\) identifies. This relation is a preorder; modulo equal
kernels it is a partial order; the identity observer, when present, is finest,
and the zero observer, when present, is coarsest.

\emph{Proof.} Kernel inclusion is reflexive and transitive, and mutual inclusion is
equality of kernels. Every kernel contains \(\{0\}\) and is contained in the
whole space. \(\square\)

Over an infinite field the set of all quotient observers on a nonzero \(Q^1\) is
infinite, and when \(\dim Q^1\ge2\) even their kernels form an infinite family.
Finiteness belongs to the declared family.

\texttt{WDP-CH10-P03} \textbf{Finite proposition (grammar resolution bound).} A
deterministic report map \(r\) from a declared finite set of core classes into a
finite grammar \(\mathcal G\) exposes at most \(|\operatorname{im}r|\le|\mathcal G|\)
visible classes. If the core set has more than \(|\mathcal G|\) elements, two
distinct core classes share a report.

\emph{Proof.} Visible classes are the fibers of \(r\), and there is one fiber per
element of \(\operatorname{im}r\subseteq\mathcal G\). The pigeonhole principle
gives the collision. \(\square\)

For reports of length at most \(L\) over an alphabet \(\Sigma\),
\(|\mathcal G|\le\sum_{\ell=0}^{L}|\Sigma|^\ell\). Grammar design is therefore a
resolution constraint on the observer, not a presentation detail.

Two consequences follow. A quotient observer separates all of \(Q^1\) exactly
when \(\ker R=\operatorname{im}d_0\), which needs at least \(\dim Q^1\) report
coordinates. And a reporter that reads windows of radius \(d\) over \(k\) letters
has a grammar of at most \(k^{2d+1}\) reports, so a longer run of distinct local
situations must repeat a report.

\subsubsection{10.6 Sharp observer completion: raw defects and repair classes}\label{sharp-observer-completion-raw-defects-and-repair-classes}

\textbf{GEN-WDP-46 (finite rational observer-completion theorem).} Supply rational
matrices \(D:\mathbb R^p\to V\) and \(Q:V\to W\), interpreted over their real
extensions, and put
\[
B=\operatorname{im}D,\qquad K=\ker Q,\qquad H=B\cap K,
\qquad k=\dim K,\quad r=\dim H.
\]
Here \(B\) is both the image of realizable raw defects and the space of allowed
repair changes. This is an explicit hypothesis: distinct realization and
repair operators require distinct subspaces. Give \(V,W\), and each augmented
report space faithful rational polyhedral norms. For an arbitrary cochain
\(v\in V\), define
\[
I(b)=\|b\|_V\quad(b\in B),\qquad
\Phi(v)=\operatorname{dist}_V(v,B),\qquad
\Phi_Q(Qv)=\operatorname{dist}_W(Qv,Q(B)).
\]

\begin{longtable}[]{@{}
  >{\raggedright\arraybackslash}p{(\linewidth - 4\tabcolsep) * \real{0.3000}}
  >{\raggedright\arraybackslash}p{(\linewidth - 4\tabcolsep) * \real{0.3000}}
  >{\raggedleft\arraybackslash}p{(\linewidth - 4\tabcolsep) * \real{0.4000}}@{}}
\toprule\noalign{}
\begin{minipage}[b]{\linewidth}\raggedright
Required property
\end{minipage} & \begin{minipage}[b]{\linewidth}\raggedright
Exact criterion
\end{minipage} & \begin{minipage}[b]{\linewidth}\raggedleft
Minimum additional rational scalar measurements
\end{minipage} \\
\midrule\noalign{}
\endhead
\bottomrule\noalign{}
\endlastfoot
Raw: \(Qb=0\Rightarrow b=0\) for every \(b\in B\) & \(H=0\) & \(r\) \\
Quotient: \(\Phi_Q(Qv)=0\Rightarrow\Phi(v)=0\) for every \(v\in V\) & \(K\subseteq B\) & \(k-r\) \\
Both & \(K=0\) & \(k\) \\
\end{longtable}

These are unrestricted linear measurement counts. A restriction to local
sensors, coordinate reads, finite precision, or a fixed sensor menu requires
its own attainability theorem. The augmented norm is supplied separately;
minimum row count alone does not minimize the recovery constant.

\emph{Proof.} Raw invisibility is precisely \(H\). For the quotient,
\(Q^{-1}(Q(B))=B+K\), so observed zero implies original zero exactly when
\(K\subseteq B\). The two criteria together require \(K=0\). The canonical
exact sequence
\[
0\longrightarrow H\longrightarrow K\longrightarrow(K+B)/B
\longrightarrow0
\]
separates the two invisible dimensions. For an augmentation \(A=(Q,R)\),
raw soundness requires \(R|_H\) injective, quotient soundness requires
\(\ker(R|_K)\subseteq H\), and both require \(R|_K\) injective. Rank-nullity
gives the three lower bounds. Extend a rational basis of \(H\) first to \(K\),
then by the remaining directions of \(B\), and finally to \(V\). Its dual
coordinates on the first \(r\), the next \(k-r\), or all first \(k\) directions
attain the respective bounds. \(\square\)

For a fixed sound augmentation \(A\), the least recovery constants in
\[
\|b\|_V\le C_{\rm raw}\|Ab\|_U,\qquad
\Phi(v)\le C_{\rm quo}\operatorname{dist}_U(Av,A(B))
\]
are finite and rational under the declared polyhedral norms. Exhaustive
rational vertex enumeration and matching primal/dual distance certificates
prove sharpness. A nonzero invisible vector gives an infinite constant
for the corresponding failed property. The constant is zero when its
recovery domain is zero: \(B=0\) for raw recovery, or \(V/B=0\) for quotient
recovery. A zero report space with a nonzero recovery domain gives infinity.
The \href{code/observer_completion/sharp_observer_completion_theorem.md}{full derivation and certificate specification}
give the norm hypotheses, complete finite constructions, counterexamples,
and deterministic error bounds. No efficiency claim is made.

The source three-patch carrier has \((\dim B,k,r)=(15,6,3)\), giving minimum
counts \((3,3,6)\). Its realized family satisfies \(I(Ds)=6|\lambda|\) and
\(QDs=0\), while \(\Phi(Ds)=0\) because every section defect is exact. A separate
invisible nonexact cochain witnesses quotient failure. Three cycle reads
restore quotient recovery with sharp \(\ell^1\) constant 1; six missing
coordinate reads restore both properties with constants 1. The coupled
carriers have \((\dim B,k,r)=(2,2,0)\), giving counts \((0,2,2)\): raw
soundness already holds, but quotient soundness fails. Their source report
coordinate conventions differ and are checked separately. The original
coupled raw constant is 2 for stacked \(\ell^1\) reports and 4 for the maximum
individual score.

This theorem joins the raw false-acceptance test (a locally accepted but
actually defective configuration) to WDP's quotient-observer
test without conflating them. In §10.5's convention \(QD=0\), quotient
soundness specializes to \(K=B\); raw soundness is deliberately lost when
\(B\ne0\). Projection Obstruction's physical-residual-to-cochain adapter must
still be supplied separately. A checked-repair evaluator may evaluate the retained exact
certificates; neither a repaired observer nor a recovery certificate grants
authority to install it or execute a policy. The
\href{code/observer_completion/README.md}{reference package} records written
proofs, source-pinned replay, 15 adversarial controls and 27 scoped Python
tests. It supplies no new Lean formalization or verification of the native
Rust owners. This bridge claim supplements the canonical claim table.

\begin{center}\rule{0.5\linewidth}{0.5pt}\end{center}

\subsubsection{10.7 Sufficient observation under physical locality}\label{sufficient-observation-under-physical-locality}

The finite completion theorem of §10.6 permits arbitrary declared linear
observation rows. A physical instrument must additionally realize its rows
with the declared spatial support, communication and precision constraints.
These are separate obligations; algebraic completion does not supply an
instrument.

For density matrices rho and sigma with equal \textbf{joint} marginals on a
region R, let O\_R be supported there and let \textbar\textbar O-O\_R\textbar\textbar{} \textless= beta. Then

\[
 |\operatorname{Tr}((\rho-\sigma)O)|
 \le \|\rho-\sigma\|_1\,\beta\le2\beta.
\]

Subtract the zero local readout Tr((rho-sigma)O\_R), then apply trace/operator
norm duality. A supplied Lieb--Robinson estimate and a specified conditional
expectation can provide beta after the spatial tail is summed. Equal single-site
marginals alone do not establish the joint-marginal premise.

For a physical-to-target map F defined on \textbf{states}, the separation is between
F(rho) and F(sigma), not F(rho-sigma), which is generally outside its domain.
A reconstruction map on readouts with Lipschitz constant kappa and endpoint
errors at most alpha satisfies delta \textless= 2 alpha + kappa eta, where eta is
the actual readout separation. When eta \textless= 2 beta and delta \textgreater{} 2 alpha,
stable reconstruction requires the corresponding conditioning price. This
does not turn an error upper bound into a device error floor.

Readout noise balls overlap if the separation is at most twice their radius.
At a common report one reconstruction error is at least half the target
separation. No exact waiting threshold T \textgreater= r/v follows from an exponential
upper envelope without its constants, desired precision and lower evidence.

The \href{docs/meron_field_limit_derivations/unified_lieb_robinson_observer_completion_derivation.md}{corrected derivation}
supplies the full finite recovery/completion theorem, locality and conditioning
proofs, actual CAR exchange signs and executable counterexamples. GEN-WDP-56
in §20.11 binds its written, formal and executable scopes separately. This
section supersedes the duplicate locality bridge imported at remote-main
commit 9b55e6c1; the original §10.6 and its canonical claims are retained.

\begin{center}\rule{0.5\linewidth}{0.5pt}\end{center}

\subsubsection{10.8 Question-preserving encodings and a direct checker}\label{question-preserving-encodings-and-a-direct-checker}

The \href{docs/question_preserving_encoding.md}{encoding bridge and experiment}
apply the existing report-fiber criterion to a bounded symbolic task. For a
fixed question \(q\) and encoding \(E\), equal reports with unequal answers
exclude a decoder that correctly decides both. A checker may still soundly
refuse both; the obstruction does not require false acceptance.

For the positive bridge, let \(P\) be an explicit partial parser and assume
\(P(E(c))\) is defined with \(q_Z(P(E(c)))=q_C(c)\). If a checker \(K\) obeys
\(K(z,w)=\mathrm{accept}\Rightarrow q_Z(z)=\mathrm{true}\), then accepting
\(K(P(E(c)),w)\) implies \(q_C(c)=\mathrm{true}\). This follows by substitution
of the two premises. Checker soundness and the binding of a received report
to an original candidate remain explicit obligations; cochain closure does
not supply either premise. This specializes the faithful-representation
soundness interface of §20.6 to an explicit parser and bounded experiment.

The finite implementation represents integer quadratics with coefficients
in \([-2,2]\). Its protected question is an exact derivative identity together
with a specified value at zero. The edge differences around a triangle lose
the common coefficient offset: \((0,0,0)\) and \((1,1,1)\) have the same report
but different answers to the zero task. Adding one coefficient as an anchor
gives an explicit inverse parser. The parser calls the canonical closure
functional and rejects malformed reports; the direct polynomial checker
then evaluates the original question.

Across 225 tasks and 125 candidates, enumeration uses 20,375 checker calls.
Both direct algebraic reconstruction and anchored cochain reconstruction
use 125. The latter adds a scalar and a closure test for each successful
proposal. Thus this declared experiment demonstrates no cochain-specific
search or diagnostic efficiency gain. All 28,000 incorrect task/candidate
pairs nevertheless have closed encoded cochains. Correctness rests on the
checker; representation consistency, source/world interpretation and action
authorization remain distinct. The supplement records full proofs, cost
scope, failure controls and native replay bindings.

\subsubsection{\texorpdfstring{10.9 The Priced Observer Boundary State and Referential Binding Protocol \(\mathcal{G}_{\mathrm{bind}}\)}{10.9 The Priced Observer Boundary State and Referential Binding Protocol \textbackslash mathcal\{G\}\_\{\textbackslash mathrm\{bind\}\}}}\label{the-priced-observer-boundary-state-and-referential-binding-protocol-mathcalg_mathrmbind}

\texttt{WDP-CH10-D07} \textbf{Definition (Refined Observer Boundary State).}\\
The observer boundary state is formally specified as the projection pair:
\[\pi_O: X_O \to Y_O, \quad X_O = (S_O, R, \mathrm{TL}_O^{\mathrm{priv}}), \quad Y_O = (B_O, R_{\mathrm{pub}}, \mathrm{TL}_O^{\mathrm{pub}})\]
where \(R_{\mathrm{pub}} = \mathbf{1}_{\{ R \ge c_{\min} \}}\) projects a coarsened operational solvency predicate, while the exact quantitative ledger \(R \in \mathbb{Q}_{\ge 0}\) is arbitrated out-of-band by the external host runtime (\texttt{ENC-0195}) and contained within the private projection fiber \(\operatorname{Fib}_O := \ker \pi_O\) satisfying \(\dim(\operatorname{Fib}_O) > 0\).

\texttt{WDP-CH10-P04} \textbf{Proposition (Private Fiber Invariance in Boundary Observation).}\\
For all internal perturbations \(v \in \operatorname{Fib}_O\) and all external observation projections \(\pi_O\), \(\pi_O(x + v) = \pi_O(x)\). Public observation is invariant under internal fiber shifts.

\texttt{WDP-CH10-D08} \textbf{Definition (Referential Binding Protocol \(\mathcal{G}_{\mathrm{bind}}\)).}\\
The protocol \(\mathcal{G}_{\mathrm{bind}}\) coordinates reference across the observer boundary via four discrete steps:
1. \textbf{Ingress:} Receiving candidate token stream under Gate 1;
2. \textbf{Admission:} Validating form and data preconditions under Gate 2, debiting ledger \(R\);
3. \textbf{Retention:} Evaluating structural asymmetry cost under Gate 3;
4. \textbf{Attestation:} Computing quotient defect \(\Phi_1\) under Gate 4 and emitting receipt under Gate 5.

\subsection{Chapter 11 --- Intersubjective presentation, descent, and redundancy}\label{chapter-11-intersubjective-presentation-descent-and-redundancy}

\subsubsection{11.1 Many observers, one carrier}\label{many-observers-one-carrier}

Let observers \(\mathcal O_i\) have presentation maps
\(q_i:X\to P_i\). A comparison map \(r_{ij}:P_i\rightharpoonup P_j\) may allow one
presentation to be transported into another. These maps are data and may be
partial.

For a family of overlapping contexts \(U_i\), local reports \(s_i\) are compatible
when their restrictions agree after declared transport on overlaps. A descent
problem asks whether there exists a global report \(s\) whose restrictions recover
the \(s_i\).

\texttt{WDP-CH11-D01} \textbf{Definition (observer descent receipt).} A descent receipt is a
tuple containing the local reports, overlap maps, compatibility checks, a global
candidate when one exists, and verification that its restrictions agree.

The definition resembles sheaf descent, but a sheaf is not obtained until a
site or cover, restriction maps, and gluing laws are supplied.

\subsubsection{11.2 Local satisfiability is not global satisfiability}\label{local-satisfiability-is-not-global-satisfiability}

Every proper subfamily of a finite constraint family can be jointly satisfiable
while the whole family is inconsistent. For example, each context in a
three-cycle can admit assignments satisfying its own parity constraint, and
each proper pair of contexts can admit a compatible assignment, although no
single assignment satisfies all three parity constraints. This is not a
compatible family of fixed local sections: if fixed sections actually agree on
all overlaps, the ordinary gluing problem has already been solved at the level
of variable assignments. Constraint-cycle examples motivate obstruction
diagnostics, but contextuality in the quantum-foundational sense additionally
requires a measurement scenario and empirical model.

The same gap has an exact quotient form. Let \(G\) be the oriented \(m\)-cycle with
vertices \(0,\dots,m-1\), edges \(e_k=(k,k+1)\) with indices mod \(m\), unit weights,
and an edge field \(\delta\) over \(\mathbb Q\) with circulation
\(c=\sum_k\delta_{e_k}\). For each \(j\) let \(P_j=E\setminus\{e_j\}\) be the path
patch that omits \(e_j\), and set \(J_{\max}(x)=\max_j\sum_{e\in P_j}|x_e|\).

\texttt{WDP-CH11-P02} \textbf{Finite proposition (three readings of one cycle class).} For
\(m\ge3\): (1) the restriction of \(\delta\) to each patch \(P_j\) is exact, so each
patch reports \(\Phi_1=0\); (2) the unit \(L^1\) quotient defect is
\(\Phi_1([\delta])=|c|\); (3) \(J_{\max}\) is a norm on \(C^1(G;\mathbb Q)\) and
\(\Phi_{J_{\max}}([\delta])=\tfrac{m-1}{m}|c|\), attained by a rational repair.

\emph{Proof.} (1) A path is a tree: fix the potential at one end and add the edge
values along the path. (2) With the cyclic orientation every coboundary sums to
\(0\) around the cycle, so every residual \(r=\delta-d_0f\) has \(\sum_e r_e=c\) and
\(\sum_e|r_e|\ge|c|\). The repair built along \(P_j\) leaves \(c\) on \(e_j\) alone, and
\(y=\pm(1,\dots,1)\) is a circulation with \(|y_e|\le1\) and \(\langle y,\delta\rangle=|c|\),
so CH04-P07 certifies the value. (3) Each edge lies in \(m-1\ge2\) patches, so
\(J_{\max}(x)=0\) forces \(x=0\), and a maximum of seminorms is a seminorm. For
every repair,

\[
\max_j\sum_{e\in P_j}|r_e|
\;\ge\;\frac1m\sum_j\sum_{e\in P_j}|r_e|
\;=\;\frac{m-1}{m}\sum_e|r_e|
\;\ge\;\frac{m-1}{m}|c|.
\]

The image of \(d_0\) is the sum-zero hyperplane, because the cycle is connected
and \(d_0\) has rank \(m-1\). So the residual with every entry \(c/m\) equals
\(\delta-d_0f\) for a rational \(f\), and then every patch reads
\(\tfrac{m-1}{m}|c|\). \(\square\)

For the cyclic field \(\delta=(2,2,2)\) on the triangle the three readings are
\(0\), \(6\), and \(4\). Each patch alone is repaired exactly. The union of edges
carries the whole circulation: \(f=(0,2,4)\) leaves \((0,0,6)\), and \(y=(1,1,1)\)
certifies \(6\). One shared repair read patch by patch through the maximum leaves
\(4\), attained at \(f=0\) with residual \((2,2,2)\). On the triangle each patch
repair \(f^{(j)}\) of \(P_j\) is unique up to an additive constant, since each patch
meets all three vertices. Two patch repairs differ by a constant on the
endpoints of their shared edge, and the three constants
\(f^{(0)}-f^{(1)}\), \(f^{(1)}-f^{(2)}\), \(f^{(2)}-f^{(0)}\) always sum to \(c\),
whatever constants are chosen. So the local repairs glue exactly when \(c=0\):
the family is compatible up to gauge, and the gauge constants carry the
obstruction. The three numbers are values of three declared evaluator costs on
one class. None of them is the size of the class, and the proposition concerns
a single cycle with unit weights.

Provenance: the absorbed archived working-paper source 011 proved in Lean, over the
integers at its recorded commit, the triangle case of (1) and the non-exactness
of \((2,2,2)\). It checked the three numbers by exact rational Python. Its Lean
evaluates one candidate for (3) and does not prove the minimum. Not ported; no
Lean credit.

\texttt{WDP-CH11-O01} \textbf{Finite proposition (explicit observer descent).} On the finite inclusion site of intersection-closed contexts with union covers, checked report sets and coherent restriction maps form a presheaf. It is a sheaf exactly when every matching-family fiber contains one global report. Constrained assignment reports admit a glue exactly when their pointwise union is admitted; the equality-constraint countermodel refutes unrestricted local-to-global gluing. Appendix J defines the category and coverage and establishes these criteria.

The same gap appears for composed procedures. If a step \(\gamma_A\) passes a
local evaluator \(E_A\) and a step \(\gamma_B\) passes \(E_B\), the composite
\(\gamma_B\circ\gamma_A\) need not pass a global evaluator \(E_G\). CH11-P02 is an
instance: every proper patch passes and the cycle fails. Two local passes
compose to a global pass only when the boundary between them preserves what
\(E_G\) inspects. A tool or subagent that passes its own check is therefore
evidence-bearing only for the question that check declares.

\subsubsection{11.3 Redundancy}\label{redundancy}

Suppose several conditionally independent channels carry presentations of one
source. Redundancy can improve robustness: a report recoverable from any
\(k\) of \(n\) channels survives some channel failures. A finite redundancy theorem
must state the channel model, failure family, decoder, and target predicate.

Quantum Darwinism uses environmental fragments, decoherence, and preferred
observables to study a related but stronger physical setting
{[}REF-OBS-001; REF-OBS-002{]}. The monograph may borrow the structural
question---when do independent channels support stable agreement?---but not the
physical mechanism or its conclusions.

\subsubsection{11.4 Agreement and authority}\label{agreement-and-authority}

Even unanimous observer reports do not create authority. Agreement can be a
premise in an external decision rule, but the rule and its jurisdiction must be
declared. Conversely, an authorized institution can issue a decision that is
not an observer-consensus theorem. Keeping these types separate prevents a
descent result from becoming a governance claim.

\subsubsection{11.5 Contextual congruence and witness rank}\label{contextual-congruence-and-witness-rank}

The primitive-axiomatics source supplies a finite notion of context-relative
failure that complements descent. A finite contextual system \(S\) consists of a
finite state set \(X\), a finite context family \(C\) acting by maps
\(\rho(c):X\to X\), and a profile map \(p:X\to Q^1\) into the quotient of a window
(§4.8).

\texttt{WDP-CH11-D02} \textbf{Definition (contextual congruence and context witness).} Set
\(x\sim y\) when \(p(x)=p(y)\). The relation \(\sim\) is a contextual congruence when
\(x\sim y\) implies \(\rho(c)x\sim\rho(c)y\) for every context \(c\). A critical
context witness is a triple \((x,y,c)\) with \(p(x)=p(y)\) and
\(p(\rho(c)x)\ne p(\rho(c)y)\). Its exposed residual is
\(\nu(x,y,c)=p(\rho(c)x)-p(\rho(c)y)\in Q^1\). Such a triple is a
\texttt{ContextWitness} in the sense of Chapter 2.

The source's slogan, ``obstruction alone does not force refinement;
context-visible non-congruence does,'' is a design principle for when to enlarge
a presentation. It is not a theorem about which enlargement to choose.

\texttt{WDP-CH11-P01} \textbf{Finite proposition (witness-rank invariance).} The witness
rank \(\operatorname{wrank}(S)=\dim\operatorname{span}\{\nu(w)\}\), taken over all
critical context witnesses \(w\), is invariant under invertible linear coordinate
changes of \(Q^1\) and under isomorphisms of contextual systems, meaning
bijections of states and contexts that intertwine the actions, composed with a
linear isomorphism of profile spaces.

\emph{Proof.} A linear isomorphism preserves linear independence, hence the
dimension of a span. An isomorphism of systems carries critical witnesses
bijectively to critical witnesses and transforms their exposed residuals by the
profile isomorphism. \(\square\)

Paths of contexts refine the picture. For context paths \(a,b\) with common
endpoints that are declared homotopic, the path defect
\(\Theta(x;a,b)=p(\rho(a)x)-p(\rho(b)x)\) lies in \(Q^1\). It is closed and
non-exact when \(\Theta\in H^1\setminus\{0\}\) and open when
\(\kappa^1(\Theta)\ne0\). Homotopy fillers are added only by declaration; none
arises from the nerve of the context family alone.

The source also stated a reflective completion of contextual systems, a left
adjoint to the inclusion of congruence-complete systems, as a conjecture that
depends on an unproved initiality lemma for minimal refinements. It is not
absorbed here; see CH04-O01.

\subsubsection{\texorpdfstring{11.6 The Inter-Observer Descent Operad \(\mathcal{O}_{\mathrm{Bridge}}\), Relational Correspondence Span, and Represented Bridge Quotient Class \(Q^1_{\mathcal{B}_{12}}\)}{11.6 The Inter-Observer Descent Operad \textbackslash mathcal\{O\}\_\{\textbackslash mathrm\{Bridge\}\}, Relational Correspondence Span, and Represented Bridge Quotient Class Q\^{}1\_\{\textbackslash mathcal\{B\}\_\{12\}\}}}\label{the-inter-observer-descent-operad-mathcalo_mathrmbridge-relational-correspondence-span-and-represented-bridge-quotient-class-q1_mathcalb_12}

\texttt{WDP-CH11-D03} \textbf{Definition (Coloured Planar Operad \(\mathcal{O}_{\mathrm{Bridge}}\)).}\\
The inter-observer descent operad is a coloured planar operad \(\mathcal{O}_{\mathrm{Bridge}} = (\mathcal{C}, \Sigma, \circ, \mathcal{I})\) with colour palette:
\[\mathcal{C} = \{\mathsf{World}_{\mathcal{O}_i}, \mathsf{Fact}_{\mathcal{O}_i}, \mathsf{Export}_{\mathcal{O}_i}, \mathsf{Bridge}_{\mathcal{O}_1 \to \mathcal{O}_2}, \mathsf{Defect}_1, \mathsf{Receipt}_{\mathrm{joint}}\}\]
where colour \(\mathsf{Fact}_{\mathcal{O}_i} := (R_i, \sim_i, G_i, \rho_i)\) encapsulates the quotient reference carrier \(R_i\) together with the admissible internal gauge group action representation \((G_i, \rho_i)\). Generator slot arities are strictly linearized:
- \(g_{\mathrm{fact}, i} \in \mathcal{O}(\mathsf{World}_{\mathcal{O}_i}; \mathsf{Fact}_{\mathcal{O}_i})\), arity 1;
- \(g_{\mathrm{defect}} \in \mathcal{O}(\mathsf{Fact}_{\mathcal{O}_1}, \mathsf{Bridge}_{\mathcal{O}_1 \to \mathcal{O}_2}, \mathsf{Fact}_{\mathcal{O}_2}; \mathsf{Defect}_1)\), arity 3, evaluating \(\alpha_{\mathrm{designation}}\) over encapsulated \((G_1, G_2)\) representations without auxiliary wire inputs;
- \(g_{\mathrm{receipt}} \in \mathcal{O}(\mathsf{Defect}_1; \mathsf{Receipt}_{\mathrm{joint}})\), arity 1;
yielding overall planar composition tree \(\tau = g_{\mathrm{receipt}} \circ_1 (g_{\mathrm{defect}} \circ (g_{\mathrm{fact}, 1}, 1_{\mathsf{Bridge}}, g_{\mathrm{fact}, 2}))\) of arity 3 with zero wire crossings and fan-out 1.

\texttt{WDP-CH11-D04} \textbf{Definition (Relational Correspondence Span).}\\
Given two observers \(\mathcal{O}_1, \mathcal{O}_2\) with respective report quotient spaces \(R_1, R_2\), the inter-observer comparison is formalized by the relational correspondence span:
\[R_1 \xleftarrow{\pi_1^B} K_{12} \xrightarrow{\pi_2^B} R_2\]
where \(K_{12} \subseteq R_1 \times R_2\) is a closed sub-relation replacing the unconstructed categorical pullback apex \(\mathcal{O}_1 \wedge \mathcal{O}_2\).

\texttt{WDP-CH11-D05} \textbf{Definition (Four Congruence Annihilators).}\\
Admissibility of inter-observer transport requires satisfaction of four congruence annihilators evaluated by \(g_{\mathrm{defect}}\):
1. \(\alpha_{\mathrm{core}}\): verifies support compatibility on span apex \(K_{12}\);
2. \(\alpha_{\mathrm{equiv}}\): verifies quotient descent preservation (\(k \sim_1 k' \iff k \sim_2 k'\));
3. \(\alpha_{\mathrm{designation}}\): verifies gauge invariance under internal presentation relabeling across encapsulated gauge orbits:
\[\forall (g_1, g_2) \in G_1 \times G_2,\; \forall k \in K_{12},\; \pi_1^B(\rho_1(g_1) \cdot k) \sim_1 \pi_1^B(k) \land \pi_2^B(\rho_2(g_2) \cdot k) \sim_2 \pi_2^B(k);\]
4. \(\alpha_{\mathrm{defect}}\): verifies coboundary exactness \(\Phi_1(\omega_{12}) = \inf_\alpha \|\omega_{12} - d_0^{12} \alpha\|_1 = 0\).

\texttt{WDP-CH11-T02} \textbf{Theorem (Represented Bridge Quotient Class and LP Dual Circulation Witness).}\\
Let \(\mathcal{B}_{12}\) be an inter-observer bridge complex terminating at degree 1 (\(C^2_{\mathcal{B}_{12}} = \varnothing\)). Then:
1. The 1-cocycle condition is vacuous (\(Z^1_{\mathcal{B}_{12}} = C^1_{\mathcal{B}_{12}}\)), and the degree-1 obstruction class is represented by the quotient class:
\[q_{\mathcal{B}_{12}} = [\omega_{12}] \in Q^1_{\mathcal{B}_{12}} := C^1_{\mathcal{B}_{12}} / \operatorname{im}(d_0^{12})\]
2. The dual space is canonically isomorphic to the adjoint kernel \((Q^1_{\mathcal{B}_{12}})^* \cong \ker((d_0^{12})^*)\);
3. By finite-dimensional linear programming duality, there exists an optimal dual witness \(\varphi^* \in \ker((d_0^{12})^*)\) with \(\|\varphi^*\|_\infty \le 1\) such that:
\[\langle \varphi^*, \omega_{12} \rangle = \|q_{\mathcal{B}_{12}}\|_{Q^1_{\mathcal{B}_{12}}}\]
attaining zero rational duality gap.

\texttt{WDP-CH11-R01} \textbf{Remark (Connection to Trace-Zero Monograph Theorem 3.1).}\\
When potential repairs \(\alpha \in C^0_{\mathcal{B}_{12}}\) are unrestricted, the dual witness condition is \((d_0^{12})^* \varphi^* = 0\) (exact vanishing divergence, matching Problem (8)). When potential repairs are constrained by trace-zero or mean-zero conditions across observer frames, the dual witness condition relaxes to \((d_0^{12})^* \varphi^* \in \mathbb{R}I\) (scalar divergence, matching Problem (7)).

\begin{center}\rule{0.5\linewidth}{0.5pt}\end{center}

\subsection{Chapter 12 --- The observation ladder}\label{chapter-12-the-observation-ladder}

\subsubsection{12.1 Five levels}\label{five-levels}

The word \emph{observation} often compresses several events. We separate them.

\begin{enumerate}
\def\labelenumi{\arabic{enumi}.}
\tightlist
\item
  \textbf{Interaction.} A system and target participate in a declared coupling.
\item
  \textbf{Registration.} The coupling changes a record carrier.
\item
  \textbf{Typed witness.} The record is linked by a declared evidence relation to a
  proposition about the interaction.
\item
  \textbf{Re-comparison.} The record can be accessed and compared under the same or
  a transported profile.
\item
  \textbf{Evaluator acceptance.} A declared evaluator returns an admissible status
  for the typed proposition.
\end{enumerate}

\texttt{WDP-CH12-D01} \textbf{Definition (observation record and validated observation).}
An observation record is

\[
\mathsf{Obs}=(i,g,w,m,e)
\]

whose fields record interaction \(i\), registration \(g\), typed witness \(w\),
memory/re-comparison \(m\), and evaluator result \(e\). A \textbf{validated
observation} is a well-typed observation record for which the first four
level predicates verify and \(e=\mathsf{Accept}\).

A receipt for a failed registration, invalid witness, missing re-access,
refusal, rejection, or undetermined result records an observation attempt but
does not inhabit the validated observation type.

\subsubsection{12.2 Failure modes}\label{failure-modes}

The ladder makes failure localizable.

\begin{itemize}
\tightlist
\item
  \textbf{No coupling:} the proposed interaction did not occur under the model.
\item
  \textbf{No registration:} coupling occurred but left no admissible record.
\item
  \textbf{Misregistration:} the record does not satisfy the witness relation.
\item
  \textbf{No retention:} the record existed but cannot be re-accessed.
\item
  \textbf{Profile drift:} the later comparison uses an incompatible context.
\item
  \textbf{Evaluator refusal:} the query is ill typed or outside the evaluator's
  domain.
\item
  \textbf{Evaluator rejection:} a well-typed claim fails its declared criterion.
\item
  \textbf{Undetermined:} available evidence does not settle the criterion.
\end{itemize}

These statuses should not be replaced by a single Boolean unless the loss of
information is deliberate.

\subsubsection{12.3 Observation is not authorization}\label{observation-is-not-authorization}

\texttt{WDP-CH12-P01} \textbf{Typing proposition.} \texttt{ObservationAccept} does not imply
\texttt{AuthStanding} or \texttt{Promote} in the layered architecture, because those predicates
belong to a separate authority object and no coercion between the types is
defined.

This proposition is almost syntactic, but it repairs a consequential ambiguity
in the measurement source. An evaluator may accept that a thermometer read
\(20.1\,^{\circ}\mathrm C\) under its procedure. That does not authorize opening a
facility, publishing a medical conclusion, or changing a controlled standard.

\subsubsection{12.4 Worked ladder}\label{worked-ladder}

A finite switch closes, changing a voltage at a sensor. An analog-to-digital
converter stores the integer \texttt{813}. A calibration table associates the code
with a voltage interval. Memory stores the code, calibration version, and time
index. An evaluator checks checksum, range, and calibration validity.

\begin{itemize}
\tightlist
\item
  Interaction: switch/sensor coupling.
\item
  Registration: stored code \texttt{813}.
\item
  Witness: code and calibration receipt support a voltage proposition.
\item
  Re-comparison: the stored record can be checked later.
\item
  Evaluator acceptance: validity predicates pass.
\end{itemize}

This is already a validated observation. It becomes a measurement only after
the separate measurand specification and result carrier, value map, unit,
calibration, uncertainty/error semantics, metrological-traceability evidence
when claimed, and provenance/re-access fields of Chapter 13 are supplied.

\begin{center}\rule{0.5\linewidth}{0.5pt}\end{center}

\subsection{Chapter 13 --- Measurement as a typed extension of observation}\label{chapter-13-measurement-as-a-typed-extension-of-observation}

\subsubsection{13.1 Metrological repair}\label{metrological-repair}

The earlier measurement draft correctly recognized that measurement requires a
scale or partition and re-access. It was incomplete in four ways: it used
authority where evaluator standing was required; it did not distinguish
measurand from indication; it lacked calibration and uncertainty; and it lacked
traceability and instrument-response semantics. We repair those omissions here
without claiming a full physical metrology theory.

\texttt{WDP-CH13-D01} \textbf{Definition (finite measurement packet).} A finite measurement
packet in this WDP-specific interface is

\[
\mathfrak M=(\mathsf{Obs},Q,Y,h,u,\mathcal S,\mathsf{Cal},
\mathsf{Unc},\mathsf{Err},\mathsf{Tr}_{\mathrm{met}},\mathsf{Prov}),
\]

where:

\begin{itemize}
\tightlist
\item
  \(\mathsf{Obs}\) is a validated observation record, a stronger WDP-specific
  prerequisite rather than part of the VIM definition of measurement
  {[}REF-MET-001{]};
\item
  \(Q\) specifies the measurand: the quantity intended to be measured, including
  the relevant system, quantity kind, and stated conditions;
\item
  \(Y\) is the distinct carrier for reported quantity values or results;
\item
  \(h\) maps admissible indications/records to reported values;
\item
  \(u\) is a unit or scale convention when the value type requires one;
\item
  \(\mathcal S\) is the permitted value/partition structure;
\item
  \(\mathsf{Cal}\) is a versioned calibration record relating reference quantity
  values and their uncertainties, indications and their uncertainties, and the
  conditions and relation used to obtain a result;
\item
  \(\mathsf{Unc}\) is the measurement-uncertainty model and budget, when that
  metrological claim is made;
\item
  \(\mathsf{Err}\) separately records any correction, measurement-error estimate,
  worst-case error bound, confidence set, or result enclosure, with its own
  semantics;
\item
  \(\mathsf{Tr}_{\mathrm{met}}\) records any claim of metrological traceability as
  a property of the measurement result established through a documented
  unbroken calibration chain, each calibration contributing to the stated
  measurement uncertainty;
\item
  \(\mathsf{Prov}\) is the separate provenance and re-access record linking the
  result to the procedure, calibration version, evaluator, and raw record.
\end{itemize}

The packet is \textbf{valid relative to a profile} only when each component is typed,
the evaluator accepts its declared checks, and the result lies within the
procedure's domain.

\subsubsection{13.2 Measurand, indication, result}\label{measurand-indication-result}

The measurand specification \(Q\) is not a value carrier and neither is the raw
record. An instrument produces an indication \(z\); a value map and calibration
yield an estimate \(\hat y=h_{\mathsf{Cal}}(z)\) in \(Y\). A measurement result
should pair that estimate with
an uncertainty description rather than present an exact numeral when the
procedure does not justify one.

A finite interval
\(E(z)=[\hat y-\epsilon_-,\hat y+\epsilon_+]\) can be declared as a worst-case
result enclosure or error bound, and a confidence set can be declared under a
specified sampling model. Neither object is automatically a VIM/GUM measurement
uncertainty {[}REF-MET-001; REF-MET-002; REF-MET-003{]}. A claim of measurement uncertainty must state its uncertainty
components, evaluation and propagation method, coverage convention when used,
and relationship to the reported result. Probabilistic uncertainty requires a
probability model and must not be inferred just from finite witness weights.

When serialized for machine audit and custody ledgers, a validated observation or measurement packet could be recorded in a 12-field evidential receipt such as:

\[
\begin{aligned}
R = (&\text{claim\_id}, \text{theory\_profile}, \text{inputs}, \text{operation}, \text{result}, \text{witness\_role},\\
&\text{assumptions}, \text{provenance}, \text{uncertainty}, \text{falsifier}, \text{tool\_version}, \text{dependencies}).
\end{aligned}
\]

Such a record would keep raw indication, calibration lineage, and uncertainty
model separate. The schema is a proposal of this text, not a package interface:
the executable \texttt{Receipt} in \texttt{code/witnessed\_difference/receipts.py} has six
fields (subject, input and output digests, claim IDs, implementation version,
and receipt digest).

\subsubsection{13.3 Calibration}\label{calibration}

Let reference inputs \(y_1,\ldots,y_k\) have accepted reference quantity values
with stated uncertainties and conditions, and let the instrument produce
indications \(z_1,\ldots,z_k\) with their indication uncertainties. A calibration
record relates those standards and indications and records how that relation is
used to obtain a result and its uncertainty. A fitted affine calibration curve,

\[
\hat y=az+b,
\]

is one component rather than the whole calibration. The parameters \(a,b\) are
fitted or otherwise certified from the reference data; reference-value,
indication, model, and fit-residual contributions enter the appropriate
uncertainty or error accounts. Extrapolation outside the calibrated range must
be separately flagged.

\texttt{WDP-CH13-C01} \textbf{Conditional result.} Write \(I\prec J\) for compact indication
intervals with \(\sup I<\inf J\). Let \(F\) be the complete declared interval map,
including calibration and every uncertainty/error enlargement used in the
report. If \(F\) is order-preserving on the admitted range in the exact sense
\(I\prec J\Rightarrow F(I)\prec F(J)\), then ordered separated indication
intervals yield ordered separated reported-value intervals. A monotone point
calibration alone is insufficient: the interval hypothesis must be verified,
for example by an explicit gap-versus-enlargement bound.

\subsubsection{13.4 Metrological traceability, provenance, and re-access}\label{metrological-traceability-provenance-and-re-access}

Metrological traceability is a property of a measurement result, not a synonym
for a hash chain or data lineage. A claim of it requires a documented unbroken
chain of calibrations to a stated reference, with every calibration contributing
to the measurement uncertainty. The chain may terminate in an external
standard. This monograph can type and verify supplied finite records and their
internal links; it does not certify an external institution, reference, or
calibration event.

Provenance records where the data came from and which transformations were
applied. Re-access is the stronger WDP-specific ability of a later evaluator to
retrieve the raw indication, calibration version, transformation, uncertainty
inputs, and profile used for the report. Neither property alone establishes
metrological traceability. A number without the provenance/re-access packet is
an assertion, not a replayable measurement result in this architecture.

\subsubsection{13.5 Worked measurement}\label{worked-measurement}

Consider a four-level finite sensor with raw codes
\(Z=\{0,1,2,3\}\). A calibration table maps them to intervals in volts:

\[
0\mapsto[0.0,0.4],\quad
1\mapsto[0.4,0.9],\quad
2\mapsto[0.9,1.5],\quad
3\mapsto[1.5,2.0].
\]

An interaction registers code \(2\). The checksum and range tests pass; memory
stores the raw code and calibration hash. The measurand is terminal voltage
under the declared loading condition. The reported result may be the interval
\([0.9,1.5]\,\mathrm V\), or a representative value with an explicitly derived
uncertainty. It is not justified to report ``\(1.2\) V exactly.''

If a new calibration changes code \(2\) to \([1.0,1.6]\) V, the old result remains
interpretable because its provenance/re-access packet identifies the old
version. A policy may request re-evaluation, but the new table does not
retroactively alter the stored raw code. Any claim that either result is
metrologically traceable must separately verify its calibration hierarchy and
uncertainty contributions.

\subsubsection{13.6 Measurement and probability}\label{measurement-and-probability}

Repeated measurement can produce a finite table of outcomes. Chapter 8 can
normalize their declared weights. To interpret that table as sampling
frequencies or a stochastic instrument, one must add a trial protocol,
exchangeability or independence conditions as appropriate, and a probability
model. A single measurement packet does not contain those assumptions.

\texttt{WDP-CH13-O01} \textbf{Conditional result (resolved finite target).} For sequential affine calibration stages with declared gain, offset and residual error bounds, the radius recurrence (abs(a)+alpha)r+alpha abs(x)+beta+eta contains every admissible realization without independence assumptions; nominal calibration updates have exact interval sup error abs(da x+db)+abs(da)r.
The proof, native ownership and scope are given in Appendix I.

\texttt{WDP-CH13-N01} \textbf{Nonclaim.} Measurement validity does not imply authority,
causal explanation, physical completeness, or observer independence.

\begin{center}\rule{0.5\linewidth}{0.5pt}\end{center}

\subsection{Chapter 14 --- Self-reference and the awareness boundary}\label{chapter-14-self-reference-and-the-awareness-boundary}

\subsubsection{14.1 Self-presentation}\label{self-presentation}

An observer can receive as input a representation of some of its own state.
Let \(\sigma:\mathcal O\to X_{\mathcal O}\) encode the accessible self-state and
let \(q_{\mathcal O}\sigma\) be its presentation. A self-report is then an ordinary
report with a reflexive provenance path.

\texttt{WDP-CH14-D01} \textbf{Definition (self-witnessing event).} A self-witnessing event is
a validated observation whose target proposition concerns a typed component of the same
observer packet and whose provenance/re-access record exhibits that reflexive
target.

The definition requires no phenomenological vocabulary. A program checking its
own version hash is a simple instance. It witnesses a proposition about an
accessible representation, not every fact about itself.

\subsubsection{14.2 The refuted implication ladder}\label{the-refuted-implication-ladder}

An earlier draft printed a hierarchy roughly of the form

\[
O_3\Rightarrow O_2\Rightarrow O_1\Rightarrow O_0.
\]

As stated, the levels mixed incompatible context actions and made \(O_0\) and
\(O_1\) contrary in some models. The full carrier ladder is therefore a
\texttt{RefutedDraft}, not a theorem.

The repair is to use independent predicates and explicit bridges:

\begin{itemize}
\tightlist
\item
  \(R_1\): a self-report is linked to an accessible self-state;
\item
  \(R_2\): the report remains stable under a declared context transport;
\item
  \(R_3\): an evaluator can compare the report with later behavior or state.
\end{itemize}

Appendix J supplies the complete finite classifications under explicit bridge data. Comparison availability and calibration success are separate predicates; partial maps cannot be completed by defaults.

\subsubsection{14.3 Blind spots}\label{blind-spots}

If the presentation map \(q_{\mathcal O}\) is not injective, there are distinct
input states with the same presentation. No evaluator operating only on that
presentation can distinguish them.

\texttt{WDP-CH14-P01} \textbf{Finite proposition (presentation blind spot).} If
\(x\ne y\) but \(q_{\mathcal O}(x)=q_{\mathcal O}(y)\), then every function
\(F:P_{\mathcal O}\to Z\) assigns the same output to the presentations of \(x\) and
\(y\).

This trivial factorization fact is a precise limit on self-knowledge in the
model. It is not the philosophical hard problem and does not imply that a human
or machine lacks awareness.

\subsubsection{14.4 Capacity and diagonal questions}\label{capacity-and-diagonal-questions}

A finite observer with a finite report alphabet cannot injectively encode an
arbitrarily large state space. Pigeonhole arguments yield blind spots relative
to the selected encoding. Stronger diagonal or self-reference results require
a language capable of naming its own reports, an interpretation map, and a
logic. Those structures must be listed before invoking Gödelian language.

Frauchiger--Renner-style scenarios can serve as adversarial tests for context
mixing in a quantum host {[}REF-OBS-009{]}. They are not evidence for a general
awareness theorem, and no protocol or proof from that work is imported here.

Linear presentations make the capacity bound quantitative. Let
\(q:V\to\mathbb F^M\) be a linear presentation map on a finite-dimensional space
and let \(Y\subseteq V\) be a declared target subspace. Here \(Y\) is a declared
target, not a defect: in Chapter 4 the image of \(d_0\) is the repair image, and a
quotient observer is built to be blind to it (CH10-D06).

\texttt{WDP-CH14-P02} \textbf{Finite proposition (injectivity on a declared subspace).} A
linear presentation map \(q\) on a finite-dimensional space is injective on a
declared subspace \(Y\) exactly when \(Y\) meets \(\ker q\) only in \(0\). This
requires at least \(\dim Y\) report coordinates, and with fewer, exact elimination
over the coefficient field exhibits a nonzero \(y\in Y\) with \(q(y)=0\).

\emph{Proof.} Injectivity on \(Y\) means \(q(y)=q(y')\) implies \(y=y'\) for
\(y,y'\in Y\), that is, \(q(y-y')=0\) implies \(y-y'=0\). By rank--nullity,
\(\dim(Y\cap\ker q)\ge\dim Y-M\), and Gaussian elimination on the matrix of \(q\)
restricted to a basis of \(Y\) produces a nonzero kernel vector when
\(M<\dim Y\). \(\square\)

\texttt{WDP-CH14-P03} \textbf{Finite proposition (recovery bound).} If a linear
presentation map \(q\) is injective on a finite-dimensional declared subspace \(Y\),
and norms are declared on \(Y\) and on reports, over \(\mathbb R\) or \(\mathbb C\) or,
over \(\mathbb Q\), as restrictions of norms on the real extensions, then some
\(c>0\) satisfies
\(\|y\|\le c\|q(y)\|\) for every \(y\in Y\), and a report perturbed by \(\eta\) with
\(\|\eta\|\le\varepsilon\) still gives
\(\|y\|\le c(\|q(y)+\eta\|+\varepsilon)\).

\emph{Proof.} The restriction of \(q\) to \(Y\) is an injective linear map between
finite-dimensional spaces, so it has a bounded inverse on its image, because
all norms on a finite-dimensional real or complex space are equivalent; over
\(\mathbb Q\) the constant for the real extension serves. The hypothesis is
needed: with \(Y=\mathbb Q^2\), \(q\) the identity, the \(L^1\) norm on \(Y\), and the
report norm \(|r_1+\sqrt2\,r_2|\), no constant exists. The noisy bound follows from
\(\|q(y)\|\le\|q(y)+\eta\|+\|\eta\|\). \(\square\)

For rational matrices and rational polyhedral norms, §10.6 and its
\href{code/observer_completion/sharp_observer_completion_theorem.md}{full certificate construction}
replace this existential bound by the \textbf{least} raw or quotient recovery
constant. Complete vertex coverage supplies a global upper bound and a
saturating witness supplies the matching lower bound. Measurement units,
the report norm and the noise bound are part of the instance. Rescaling a
measurement row while leaving its physical noise or units unstated is not
an informational improvement: a change of units must transport the norm
and the error bound as well. The integer coupled fixture, for example,
has sharp raw constant 2 for stacked \(\ell^1\) reports and 4 for the maximum
individual score.

For Euclidean norms on a nonzero real space \(Y\), assuming \(q\) is injective
on \(Y\), let \(V\) be a basis matrix with Gram matrix \(G=V^TV\). Then
the best constant is \(c=1/\sqrt t\) for the largest \(t\) with
\(V^Tq^TqV-tG\succeq0\). An exact \(LDL^T\) factorization certifies a rational lower
bound on \(t\). The exact Rayleigh quotient of a rational vector bounds \(t\) only
from above, so it bounds the best constant from below and cannot certify a
recovery constant. For \(Y=\{0\}\) the optimal nonnegative constant is \(0\)
and there is no largest finite \(t\). Over complex coordinates, replace every
transpose here by the conjugate transpose. In the coupled example
\(Y=\{(s_0,s_1,s_0,s_1)\}\subseteq\mathbb Q^4\) with \(q(y)=(y_0,y_3)\), each of the
two coordinate reads has a nontrivial kernel on its own edge, yet
\(\|q(y)\|_2^2=\tfrac12\|y\|_2^2\) on \(Y\), so \(c=\sqrt2\). If every probe contributes
one report coordinate at a declared positive cost, a suite that is injective on
\(Y\) costs at least the sum of the \(\dim Y\) cheapest probes.

\texttt{WDP-CH14-N02} \textbf{Nonclaim.} Injectivity of a presentation map on a declared
subspace does not detect inputs outside that subspace, does not certify
detection of unmodeled or adversarial behavior, and does not select, price, or
authorize a sensor suite.

\subsubsection{14.5 The awareness firewall}\label{the-awareness-firewall}

The monograph reserves \textbf{awareness} for a bridge not supplied by the present
finite architecture. Functional self-report, global availability, integration,
or higher-order representation may be compared with theories of consciousness
{[}REF-OBS-004; REF-OBS-005; REF-OBS-006; REF-OBS-007; REF-OBS-008{]}, but no
equivalence is asserted.

\texttt{WDP-CH14-N01} \textbf{Nonclaim.} Self-witnessing, self-modeling, evaluator closure,
or report integration does not establish phenomenal consciousness.

\texttt{WDP-CH14-O01} \textbf{Finite proposition (declared functional bridges).} For independently declared finite presentation, report, context and calibration maps, factorization and stability have the complete implication classification of Appendix J. Explicit accessible self-state, reflexive provenance, later-state and partial comparison maps separately classify linkage, stability, comparison availability and comparison success; later-target and successful linked-report existence are exact fiber-constancy questions. No phenomenology premise occurs.

\textbf{Interpretive remark (reports, not experience).} The firewall concerns which
part of an experience report can enter a shared evaluation, not whether the
experience exists. Two kinds of report differ here. An \emph{adjectival} report
attaches a quality to an object (``the apple appears red''). It has external
anchors: the object, other reporters, instruments and a shared colour space.
An \emph{adverbial} report describes a manner of undergoing (``I am remembering
vividly''). Its anchors are temporal profile, salience and later reports by the
same subject, so it is harder to place in a declared packet. Neither report is
nothing, and neither governs a shared verdict by being felt. Each enters a
shared evaluation only through a declared report, probe and comparison
structure. No claim about consciousness follows.

\begin{center}\rule{0.5\linewidth}{0.5pt}\end{center}

\subsubsection{14.6 Minimum-cost sufficient observation}\label{minimum-cost-sufficient-observation}

The unrestricted rational row counts in §10.6 answer a dimension question.
A supplied instrument catalogue with prices asks a different finite
optimization question. Let \texttt{X} and \texttt{I} be finite, with existing report
\texttt{p:X-\textgreater{}P}, protected question \texttt{q:X-\textgreater{}A}, measurements \texttt{m\_i:X-\textgreater{}M\_i} and positive
rational prices \texttt{c\_i}. Define the unresolved unordered pairs

\[U=\{\{x,y\}:p(x)=p(y),\ q(x)\ne q(y)\},\qquad
E_u=\{i:m_i(x)\ne m_i(y)\}\quad(u=\{x,y\}).\]

\textbf{Finite priced-observation theorem (GEN-WDP-47).} A fixed suite \texttt{S} is
sufficient exactly when it meets every \texttt{E\_u}. Equivalently its distinction
sets cover \texttt{U}. Its minimum cost is the exact weighted hitting-set optimum,
or infinity if some \texttt{E\_u} is empty. A feasible suite gives an upper bound;
optimality additionally requires a replayed lower bound covering every
cheaper alternative. An LP bound alone may be strictly smaller: for the
seven-point Fano incidence instance, the fractional optimum is \texttt{7/3},
while the integer optimum is \texttt{3}. The \href{docs/minimum_cost_sufficient_observation.md}{full finite proof and
certificates} give the explicit
table, primal/dual fractional witnesses and integer counting argument.

For adaptive observation, let \texttt{C} be the states remaining after the free
report and previous measurements. Put \texttt{V(C)=0} if \texttt{q} is constant on \texttt{C},
including the empty set. Otherwise, with a minimum over splitting
measurements only,

\[V(C)=\min_{i:\,|m_i(C)|>1}
 \left(c_i+\max_{a\in m_i(C)}V(C\cap m_i^{-1}(a))\right).\]

An empty minimum is infinity. Each child has smaller cardinality, so exact
dynamic programming terminates. Induction establishes the Bellman lower bound
for every strategy; choosing a minimizing measurement at each finite-valued
node realizes it. Since \texttt{p} is already observed, the initial worst-case
price is \texttt{max\_\{z\ in\ p(X)\}\ V(p\^{}\{-1\}(z))}, with zero for empty \texttt{X}, not \texttt{V(X)}
in general. It is at most the optimal fixed-suite price. Asking the full
shelf establishes feasibility precisely when there is no differently answered
pair indistinguishable under \texttt{p} and every measurement. Such a pair defeats
every adaptive strategy because both states follow the same branch at
every step.

For a concrete strict gap, use four states with distinct answers and
constant \texttt{p}. Let \texttt{A=(0,0,1,1)}, \texttt{B=(0,1,0,0)}, \texttt{C=(0,0,0,1)} with prices
\texttt{1,3,3}. The pairs \texttt{(1,3)}, \texttt{(1,2)} and \texttt{(3,4)} force \texttt{A}, \texttt{B} and \texttt{C},
respectively, so every sufficient fixed suite costs \texttt{7}. Asking \texttt{A} first
and then \texttt{B} or \texttt{C} costs \texttt{4}; the path through state \texttt{1} must ask both
\texttt{A} and \texttt{B}, proving that lower bound. A refinement of the free report can
only decrease either optimum; coarsening can only increase it.

The supplied derivation is preserved with its \href{code/finite_observation/source_provenance.json}{correction and source
record}. In particular its
original eight-state adaptive table was infeasible, its report-monotonicity
implication was reversed, and a cover witness was incorrectly listed as a
lower bound. The corrected finite checker binds the original instance and
replays submitted optimality data. This is exact bounded computation, not
a polynomial-time algorithm, an empirical speedup, instrument physics or
permission to buy or use the selected observations. No canonical \texttt{CH14}
identifier is reused or silently strengthened.

\begin{center}\rule{0.5\linewidth}{0.5pt}\end{center}

\subsubsection{14.7 Realizable observation and the cost of preserving future operations}\label{realizable-observation-and-the-cost-of-preserving-future-operations}

The measurements of §14.6 leave the measured state unchanged. A measurement
that changes state needs an additional preservation test. Let \(k:X\to K\)
be a finite quotient preserving the protected answers, target truth, enabled
future actions, their prices and their successor-block sets. Supply actual
measurement transitions, post-state reports, admissibility conditions and
positive rational costs.

Track the nonempty relation \(z\subseteq X\times X\) of possible original
and current states. For a measurement \(i\) with successors \(T_i(y)\)
and post-state report \(\lambda_i\), the observed update is

\[
z_{i,o}=\{(x,y'):\exists y\ ((x,y)\in z,\ y'\in T_i(y),\
\lambda_i(y')=o)\}.
\]

\textbf{Finite observation-realization theorem.} A completed measurement procedure
identifies the original future class and preserves that class at handoff
exactly when every terminal relation satisfies
\(\exists C\ \forall(x,y)\in z:\ k(x)=C=k(y)\).
The least worst-case cost is the finite reachability/minimax optimum of
this relation arena. An initial-containing loss trap certifies
nonrealizability; a terminating policy and all-action lower inequalities
certify an attained optimum. Initial free-report fibers are treated separately.
The full argument, constructive decoder and certificate conditions are in
the \href{docs/derivations/realizable_observation.md}{derivation}.

This is more than a refinement of a present-answer partition. In the
specified eight-state example, the least fixed nondisturbing measurement
family costs \(6\), adaptive nondisturbing observation costs \(4\), and
temporary disturbance followed by restoration costs \(15/4\). Only
identifying the original class costs \(7/2\) but leaves the wrong future
operation enabled. Removing restoration makes the selected disturbing
measurement family insufficient. If preservation is required at every
prefix, the disturbing action is excluded and the optimum is \(4\).

The \href{code/realizable_observation/finite_replay_v3.json}{finite replay}
contains five instances and twelve hostile controls, checked by a separate
construction verifier and the existing finite-game certificate checker.
Five core-only Lean declarations check relation composition, retention of
the original coordinate, decoder uniqueness and the current-state collision;
they use no axioms. These are distinct from the complete written cost proof
and from physical realization of a measurement apparatus.

\subsubsection{14.8 Reasoning under a declared information boundary}\label{reasoning-under-a-declared-information-boundary}

\texttt{WDP-CH14-T01} \textbf{Theorem (target-separating report obstruction).} Let
\(E:W\to R\) encode all information available at the decision, and let
\(q:W\to Q\) be the protected answer. If there are admitted worlds \(w,w'\)
with \(E(w)=E(w')\) and \(q(w)\ne q(w')\), there is no function \(f:R\to Q\)
such that \(f(E(x))=q(x)\) for every \(x\in W\).

\textbf{Proof.} Such an \(f\) would give
\(q(w)=f(E(w))=f(E(w'))=q(w')\), a contradiction. Noninjectivity of \(E\)
alone is insufficient: a constant question is recoverable from a constant report.

\texttt{WDP-CH14-C01} \textbf{Corollary (post-processing preserves the collision).} For
any deterministic computation \(c\) on the available report,
\(c(E(w))=c(E(w'))\). Thus additional internal computation alone cannot settle
this target on both worlds. For interaction, the premise must cover all admitted
observation/action histories. A one-step blind spot can disappear after an
intervention; swapping two coordinates exposes a previously hidden coordinate.

\texttt{WDP-CH14-N03} \textbf{Nonclaim.} The theorem concerns a specified information map
and target, not every autonomous system or language model. Internal computation
can prove consequences of supplied tables and assumptions. Physical actuators
or external verifiers change the available information and require a new check
of the theorem's premises; they do not supply an automatic exemption or guarantee.
An invariant export cannot distinguish inputs that it identifies, whereas an
arbitrary internal verdict need not be invariant. Zero passage charge alone
does not establish the required collision.

This replaces the earlier three-axis formulation. The unsourced Lemma D and
retired NB-9 do not support it. Formal self-reference limits require their own
coding, expressiveness and soundness hypotheses; an institutional grant still
requires a separately declared authority rule. The corrected argument agrees
with \href{docs/universal_epistemic_no_go_theorem.md}{the epistemic boundary analysis}.
The complete inference calculus, source-task interpretation and formal evidence
are given in §20.14 and its derivation supplement.

\subsection{Chapter 15 --- Authority is a separate layer}\label{chapter-15-authority-is-a-separate-layer}

\subsubsection{15.1 Typed authority}\label{typed-authority}

Authority concerns who or what is permitted to issue, accept, promote, or act on
a statement within a jurisdiction. It is not an observer component.

\texttt{WDP-CH15-D01} \textbf{Definition (authority packet).} An authority packet is

\[
\mathfrak A=(R,J,\mathsf{Perm},\mathsf{Proc},\mathsf{Appeal}),
\]

where \(R\) is a set of roles, \(J\) a jurisdiction, \(\mathsf{Perm}\) a typed
permission relation, \(\mathsf{Proc}\) a procedure for exercising permissions,
and \(\mathsf{Appeal}\) any review or revocation mechanism.

The packet is external to the layered observer. A bridge from measurement or
proof status into an authority procedure is a policy choice in a jurisdiction,
not a theorem of obstruction calculus.

\subsubsection{15.2 Four separations}\label{four-separations}

The following implications are invalid without explicit bridge rules:

\[
\begin{aligned}
\text{witness exists} &\not\Rightarrow \text{observer accepts},\\
\text{observer accepts} &\not\Rightarrow \text{measurement valid},\\
\text{measurement valid} &\not\Rightarrow \text{authority standing},\\
\text{authority standing} &\not\Rightarrow \text{claim true}.
\end{aligned}
\]

An authorized role can make an incorrect statement; a correct measurement can
be outside a role's jurisdiction; an evaluator can accept an internally valid
record that lacks metrological fields; and a residual can exist without being
registered.

\subsubsection{15.3 Repository status}\label{repository-status}

A green test run says that specified commands returned expected results in a
specified environment. A formal proof says that a term inhabits a proposition
under the proof assistant's logic and imports. A release tag identifies bytes.
None of these confers institutional, scientific, or moral authority. They may
be inputs to an authority process, but only through a declared rule.

\texttt{WDP-CH15-N01} \textbf{Nonclaim.} This monograph assigns no authorization and changes
no jurisdiction. It carries no operational authority or institutional standing.

\subsubsection{15.4 Observer morphisms}\label{observer-morphisms}

To compare observers, define the morphisms rather than comparing tuple lengths.
A morphism \(F:\mathcal O\to\mathcal O'\) can contain maps

\[
F_X:X_{\mathcal O}\to X_{\mathcal O'},\quad
F_P:P_{\mathcal O}\to P_{\mathcal O'},\quad
F_M:M_{\mathcal O}\to M_{\mathcal O'}
\]

that commute with selected presentation, update, and report maps. Evaluator
transport additionally requires a map on defects and repairs preserving the
repair action and controlling the cost. A strong isometry condition may require
exact equality of costs; a simulation condition may require only a one-sided
bound.

There is no single correct morphism class for every purpose. If the question is
whether one observer can reproduce another's reports, behavioral simulation is
appropriate. If the question is whether obstruction magnitudes agree, gauge
and cost preservation are needed. If the question is whether authority
transfers, no observer morphism is enough: an authority-packet morphism must be
defined separately.

\texttt{WDP-CH10-D05} \textbf{Definition (report simulation).} Observer \(\mathcal O'\) report-
simulates \(\mathcal O\) on input family \(D\) when there is an input adapter and a
report decoder making the declared report diagram commute on \(D\). This does not
imply memory-state isomorphism or evaluator equivalence.

\subsubsection{15.5 A finite descent example}\label{a-finite-descent-example}

Take three contexts
\(U_{12}=\{x_1,x_2\}\), \(U_{23}=\{x_2,x_3\}\), and
\(U_{31}=\{x_3,x_1\}\) with Boolean values. Impose local parity constraints

\[
x_1\oplus x_2=0,\qquad
x_2\oplus x_3=0,\qquad
x_3\oplus x_1=1.
\]

Each context has local assignments. Any two constraints can be satisfied
together, but all three cannot: adding the left sides modulo two gives \(0\),
while adding the right sides gives \(1\). The summed parity is a finite
counterexample witness to global assignment.

This example has the \emph{shape} of a contextual obstruction. To call it an
empirical contextuality model, one must additionally say what the measurements
and outcomes are, provide compatible distributions on contexts as required,
and connect the logical support to those distributions. The arithmetic
contradiction alone does not supply an experiment.

\subsubsection{15.6 Observation receipt invariants}\label{observation-receipt-invariants}

An observation receipt should survive benign representation changes. Suppose a
raw record is re-encoded by a bijection \(c:R\to R'\). Registration and readback
maps can be transported through \(c\); the witnessed proposition should remain
the same if its semantics is transported as well. A checksum over raw bytes,
however, changes unless recomputed and provenance-linked. Thus semantic
invariance and byte identity are different predicates.

For each observation, retain at least:

\begin{itemize}
\tightlist
\item
  interaction identifier and target role;
\item
  registration source and raw record digest;
\item
  witness kind and proposition identifier;
\item
  context/profile identifier;
\item
  memory/re-access pointer;
\item
  evaluator version, inputs, output, and any certificate;
\item
  transformations applied to the raw record;
\item
  explicit exclusions.
\end{itemize}

This schema is deliberately more demanding than storing the final status. It
allows later reviewers to distinguish a changed evaluator from changed data.

\subsubsection{15.7 Measurement equations and uncertainty propagation}\label{measurement-equations-and-uncertainty-propagation}

Let a reported scalar be \(y=h(z,c)\), where \(z\) is an indication and \(c\) is a
calibration parameter vector. A first-order uncertainty approximation around
\((z_0,c_0)\) is

\[
\operatorname{Var}(y)\approx
\nabla h^{\mathsf T}\Sigma\nabla h,
\]

where \(\Sigma\) is the joint covariance of indication and calibration inputs.
This familiar formula is conditional on a probabilistic model, differentiability,
and the validity of local linearization. It is not part of the finite root.

For purely bounded errors, one can instead use interval or Lipschitz transport.
If \(h\) is \(L\)-Lipschitz on the admitted domain and
\(\|x-x_0\|\le\epsilon\), then

\[
|h(x)-h(x_0)|\le L\epsilon.
\]

This deterministic statement often matches the finite calculus better. A
measurement packet should declare which uncertainty semantics it uses. Mixing a
worst-case bound with a standard deviation without an adapter is a type error.

\subsubsection{15.8 Repeated instruments and state change}\label{repeated-instruments-and-state-change}

Repeatability has at least three meanings:

\begin{enumerate}
\def\labelenumi{\arabic{enumi}.}
\tightlist
\item
  repeated readout returns the same code;
\item
  repeated measurement returns compatible value intervals;
\item
  the first measurement prepares a state that, under declared assumptions, ensures the second result.
\end{enumerate}

The third belongs to an operational state-transition theory and, in quantum
models, to instrument structure. The finite measurement packet by itself
supports only the first two when suitable stability predicates are provided.
Davies--Lewis and Ozawa supply primary quantum-instrument comparisons
{[}REF-MET-006; REF-MET-007{]}; their results should not be back-projected into the
general definition.

\subsubsection{15.9 Self-model calibration}\label{self-model-calibration}

A self-report can be calibrated against externally accessible behavior without
being exhaustive. Let an observer emit a finite report \(r\in R\) and let a test
procedure observe a behavior \(b\in B\). A calibration relation
\(K\subseteq R\times B\) records which pairs count as agreement. Repeated tests
can reveal systematic mismatches. The result concerns \(K\) and the sampled
behavior family, not the total internal state.

This gives a precise repaired hierarchy:

\begin{itemize}
\tightlist
\item
  \textbf{self-presentation:} a report is generated from an encoded self-state;
\item
  \textbf{self-retention:} the report is available later;
\item
  \textbf{self-comparison:} a later procedure compares it with a typed target;
\item
  \textbf{self-calibration:} the comparison satisfies a declared relation on a test
  family;
\item
  \textbf{cross-context robustness:} the relation is preserved under named context
  transports.
\end{itemize}

No implication should be asserted in reverse, and forward implications require
the adapter laws. A stable report can be consistently wrong; a calibrated
report can cover only a narrow test family; a broad functional self-model still
does not settle phenomenology.

\subsubsection{15.10 Authority bridge example}\label{authority-bridge-example}

Suppose a laboratory procedure says: ``A calibration-valid measurement below
threshold \(\theta\), signed by role \texttt{Operator}, permits action \(A\).'' This rule is
a bridge

\[
\mathsf{MeasurementReceipt}\times\mathsf{RoleCredential}
\longrightarrow\mathsf{PermissionDecision}.
\]

The measurement supplies a value and uncertainty; the credential supplies
standing; the procedure supplies the jurisdictional inference. If any input is
missing, the bridge refuses. If a later calibration changes the value interval,
the procedure determines whether prior decisions are reviewed. The obstruction
calculus can audit the inputs and rule. It does not create the rule's legitimacy.

\subsubsection{15.11 Part III interface theorem schema}\label{part-iii-interface-theorem-schema}

The most useful theorem schema for this part is compositional and conditional.
If:

\begin{enumerate}
\def\labelenumi{\arabic{enumi}.}
\tightlist
\item
  an interaction-to-record map satisfies a declared fidelity relation;
\item
  the record-to-proposition semantics is valid on a stated domain;
\item
  memory preserves the required fields;
\item
  the evaluator is sound for its status predicates;
\item
  calibration and uncertainty transport are valid;
\item
  provenance links all transformations and any metrological-traceability claim
  satisfies the documented calibration-chain and uncertainty requirements;
\end{enumerate}

then the final measurement result is valid \textbf{relative to that package}. Each
premise has its own falsifier. The conclusion does not include authorization or
truth outside the measurand/profile domain.

\texttt{WDP-CH13-O02} \textbf{Finite proposition (resolved finite target).} The declared finite registration/witness/measurement instance replays the retained indication, composed affine map, correction and uncertainty enclosure against separate re-access data; altered records, understated uncertainty and rehashed false enclosures are rejected.
The proof, native ownership and scope are given in Appendix I.

\subsubsection{15.12 Causal representations are proposals}\label{causal-representations-are-proposals}

Structural causal models {[}REF-CAUSE-001; REF-CAUSE-002{]} express interventional
and counterfactual queries. This section places such queries in the
separations above. It adds no theorem about causal inference itself.

\textbf{Definition (causal proposal).} A causal proposal is a represented
interventional or counterfactual claim in a declared causal representation
\(M\), for example ``under \(\mathrm{do}(A=a)\) the outcome \(Y\) changes by
\(\Delta\)'' or ``had \(A\) not occurred, \(Y\) would have differed''.

\textbf{Definition (causal overclaim).} A causal overclaim passes from ``\(M\) can
represent the intervention or counterfactual'' to a standing such as ``the system
is trustworthy'' or ``acting on the answer is authorized'' without the evaluation
and authority packets that standing requires.

\textbf{Finite proposition (a model and query do not fix a verdict).} Let a domain of
packets \((M,q,\mathcal O)\) carry a protected status \(s(M,q,\mathcal O)\). If two
packets share \((M,q)\) and differ in \(s\), then no function of \((M,q)\) alone
equals \(s\) on that domain.

\emph{Proof.} This is the fiber criterion of P.1 and L.2 with report \((M,q)\). A
function of the report takes one value on the shared fiber, and \(s\) takes
two. \(\square\)

The hypothesis is met by a two-variable instance. Let \(M\) be the binary model
\(Y:=X\) with \(X\in\{0,1\}\), and let \(q\) state that
\(\mathbb E[Y\mid\mathrm{do}(X=1)]-\mathbb E[Y\mid\mathrm{do}(X=0)]=1\), which is
true in \(M\). An evaluator \(\mathcal O_{\mathrm{in}}\) declares the target
population to be the one \(M\) represents and supplies the support repair, so the
declared support defect is exact and \(\Phi=0\). An evaluator
\(\mathcal O_{\mathrm{out}}\) declares a target outside the represented support
and admits no transport repair; its one-coordinate defect has \(\Phi=1\). With
the same threshold rule (accept when \(\Phi\le\tfrac12\)), the first returns
accept and the second refuse, although \(M\) and \(q\) are unchanged. The defect
values here are declared by the two evaluators, not computed from data; the
instance shows only that the hypothesis is satisfiable. Even the accept status
is an evaluation, not an authorization. By §15.2 and P.19, a regime must
still supply the grant.

An earlier informal statement of this point said that there is no
observer-independent map from queries to authority. As stated it is false: a
constant assignment is such a map. The fiber statement above is the corrected
form, and it concerns agreement with a declared status on a declared domain.
Likewise, descent through an evaluator yields evaluable evidence. It does not
yield authority.

An interventional proposal enters an evaluation only after the following are
declared, each with its own falsifier: the unit or population carrier; the
admissible variables; the operational meaning of the intervention; the
comparison made; the identification assumptions; the data support in the
intervention region; any remaining confounding, selection or transport defect;
the residual after admissible repair; and the terminal statuses. A
counterfactual proposal additionally needs the factual state, the structural
equations or update rule, the evidence (abduction) regime, the prediction
semantics, and a ledger of where the answer extrapolates beyond support.
Without these, counterfactual language is a richer proposal syntax. Because the
evaluator and the regime remain separate, this section asserts no causal
claim about any system.

\subsubsection{\texorpdfstring{15.13 Separation V: Authority Isolation from Prose Assertions (\texttt{AuthorityEffect:\ None} and Executable Witness Dominance)}{15.13 Separation V: Authority Isolation from Prose Assertions (AuthorityEffect: None and Executable Witness Dominance)}}\label{separation-v-authority-isolation-from-prose-assertions-authorityeffect-none-and-executable-witness-dominance}

\texttt{WDP-CH15-D02} \textbf{Definition (Executable Machine Receipt).}\\
An executable machine receipt is a typed record \(\mathrm{Receipt} = (c, \Delta, \varphi^*, \mathrm{status})\) emitted by a sovereign test harness or compiled kernel, where \(c\) is the command, \(\Delta\) the diagnostic vector, \(\varphi^*\) the rational dual witness, and \(\mathrm{status} \in \{\mathsf{PASS}, \mathsf{FAIL}, \mathsf{REFUSE}\}\).

\texttt{WDP-CH15-P04} \textbf{Proposition (Prose Authority Impotence).}\\
Markdown declarations, natural language assertions, and cited rules carry \(\operatorname{AuthorityEffect} = \text{None}\). An autonomous agent citing prose rules does not constitute verification. Compliance is established solely by the generation of a passing machine receipt.

\texttt{WDP-CH15-N02} \textbf{Nonclaim.}\\
Proposition \texttt{WDP-CH15-P04} defines repository verification governance; it does not model legal or social authority structures.

\begin{center}\rule{0.5\linewidth}{0.5pt}\end{center}

\subsection{Part IV --- Histories and the quantum boundary}\label{part-iv-histories-and-the-quantum-boundary}

\subsection{Chapter 16 --- Spatial/history diagnostics and the curved-complex fork}\label{chapter-16-spatialhistory-diagnostics-and-the-curved-complex-fork}

\subsubsection{16.1 Two directions}\label{two-directions}

Suppose a model has a spatial or contextual degree \(p\) and a historical degree
\(q\). Let

\[
C^{p,q}
\]

be a finite bigraded family with candidate maps

\[
d_s:C^{p,q}\to C^{p+1,q},\qquad
d_h:C^{p,q}\to C^{p,q+1}.
\]

A genuine anticommuting double complex requires

\[
d_s^2=0,\qquad d_h^2=0,\qquad d_sd_h+d_hd_s=0.
\]

Earlier drafts treated these identities as likely constructions while also
recording defects that can make them false. Here the identities are gates, not
labels.

\texttt{WDP-CH16-D01} \textbf{Definition (curvature defects).} Define

\[
K_s=d_s^2,\qquad K_h=d_h^2,\qquad
\Lambda=d_sd_h+d_hd_s.
\]

The vanishing locus (or flat locus) \(Z\) is the subfamily on which all three
defects vanish. With this anticommuting convention the unsigned total operator
\(D=d_s+d_h\) satisfies \(D^2=K_s+\Lambda+K_h\), so \(D^2 x = 0\) for each \(x \in Z\).
However, the vanishing locus \(Z\) is \textbf{not} automatically a subcomplex: forming a
genuine chain complex \((Z, D|_Z)\) requires the additional hypothesis that \(Z\) is
preserved by the differentials (\(d_s(Z) \subseteq Z\) and \(d_h(Z) \subseteq Z\)).

A finite countermodel demonstrates that differential preservation does not follow
from pointwise vanishing: let \(C^{0,0} = \operatorname{span}(x)\),
\(C^{1,0} = \operatorname{span}(y)\), \(C^{0,1} = \operatorname{span}(u)\),
\(C^{1,1} = \operatorname{span}(z)\), and \(C^{1,2} = \operatorname{span}(w)\), with all
other components zero. Set \(d_s(x) = y\), \(d_s(u) = -z\), and \(d_s = 0\) elsewhere;
set \(d_h(x) = u\), \(d_h(y) = z\), \(d_h(z) = w\), and \(d_h = 0\) elsewhere. On the
element \(x \in C^{0,0}\), all three defects vanish:
\(K_s(x) = d_s(y) = 0\), \(K_h(x) = d_h(u) = 0\), and
\(\Lambda(x) = d_s(u) + d_h(y) = -z + z = 0\), so \(x \in Z\).
Yet for its differential image \(d_s(x) = y\), the historical curvature is
\(K_h(d_s x) = d_h^2(y) = d_h(z) = w \ne 0\).
Thus \(d_s(x) \notin Z\), so \(d_s(Z) \not\subseteq Z\). Without higher Bianchi
coherence ensuring differential stability, the zero-defect locus is just a subset,
not a chain subcomplex. The alternative signed convention \(D=d_s+(-1)^p d_h\)
instead pairs with commuting bidegree maps; the two conventions must not be mixed.

\subsubsection{16.2 The three outcomes}\label{the-three-outcomes}

There are three honest outcomes.

\begin{enumerate}
\def\labelenumi{\arabic{enumi}.}
\tightlist
\item
  \textbf{Flat complex.} Prove all square-zero and mixed identities on a declared
  domain. Cohomology and spectral-sequence language then becomes available
  {[}REF-HOM-001{]}.
\item
  \textbf{Curved system.} Retain the nonzero defects as typed curvature and use a
  curved differential formalism whose axioms are separately checked
  {[}REF-HOM-002{]}.
\item
  \textbf{Higher coherent repair.} Supply explicit homotopies or higher operations
  satisfying the relevant coherence identities {[}REF-HOM-003; REF-HOM-004{]}.
\end{enumerate}

Just renaming a defect ``higher structure'' is not Outcome 3. Higher cells are
new data and their identities are new obligations.

\subsubsection{16.3 Projection-induced memory}\label{projection-induced-memory}

Let a full state evolve by \(T:X\to X\), while an observer sees a surjective
presentation \(q:X\to P\) (or replace \(P\) by \(\operatorname{im}q\)). If a section
\(i:P\to X\) is chosen, a reduced update can be written
\(S=qTi\). Descent of the full dynamics through \(q\) instead requires

\[
q(x)=q(y)\Longrightarrow q(Tx)=q(Ty).
\]

Without descent, the next presentation can depend on information discarded by
\(q\). To retain a single-valued deterministic update on the unchanged
presentation carrier, one must add enough state or history to distinguish the
relevant fibers. Alternatives include refining \(P\) or adopting relation-valued,
set-valued, or stochastic reduced dynamics.

\texttt{WDP-CH16-P01} \textbf{Finite proposition (descent criterion).} For a surjective
\(q:X\to P\), a function \(\bar T:P\to P\) with \(\bar Tq=qT\) exists exactly when
\(q(x)=q(y)\) implies \(q(Tx)=q(Ty)\).

\emph{Proof.} Necessity follows by applying \(\bar T\) to equal presentations.
For sufficiency, define \(\bar T(p)=q(Tx)\) for any \(x\) with \(q(x)=p\); the
condition makes this well defined. \(\square\)

Selecting \(S=qTi\) without descent gives one reduced dynamics, but it depends on
the section and does not establish quotient dynamics.

The same criterion governs any target, not only dynamics. For a surjective
presentation \(q:X\to P\) and a map \(F\) on \(X\), call the pairs with equal
presentation and unequal target the obstruction set of \(F\).

\texttt{WDP-CH16-P02} \textbf{Finite proposition (presentation descent and sections).} For a
surjective presentation \(q\) from \(X\) onto \(P\) and a target \(F\) on \(X\), the
obstruction set is empty exactly when \(F\) factors through \(q\); the factor is then
unique; and \(F=F\circ i\circ q\) for one section \(i\) of \(q\) exactly when the
obstruction set is empty, in which case the equality holds for every section.

\emph{Proof.} Factorization is the fiber-constancy criterion behind CH16-P01,
applied to \(F\) in place of \(q\circ T\): \(F\) factors through the surjection \(q\)
exactly when \(F\) is constant on the fibers of \(q\), which is emptiness of the
obstruction set. Uniqueness uses
surjectivity: two factors agree on \(q(x)\) for every \(x\). If
\(F=F\circ i\circ q\), then \(q(x)=q(y)\) gives \(F(x)=F(i(q(x)))=F(i(q(y)))=F(y)\).
Conversely, an empty obstruction set gives \(F(i(q(x)))=F(x)\) for every section,
since \(q(i(q(x)))=q(x)\). \(\square\)

Without surjectivity, uniqueness of the factor off the image of \(q\) is not
guaranteed, and no section exists at all. For a singleton answer target, there
is still exactly one factor on the entire report target; with two answer values,
an unconstrained report point can instead admit distinct extensions.

\subsubsection{16.4 Worked example: forgotten bit}\label{worked-example-forgotten-bit}

Let \(X=\{0,1\}^2\), \(q(x_1,x_2)=x_1\), and
\(T(x_1,x_2)=(x_2,x_1)\). States \((0,0)\) and \((0,1)\) have the same presentation,
but their next presentations are \(0\) and \(1\). Hence \(T\) does not descend through
\(q\). Remembering the discarded second bit repairs prediction. The example gives
a precise sense in which projection can induce memory; it is not by itself a
Mori--Zwanzig kernel theorem {[}REF-HOM-005{]}.

\texttt{WDP-CH16-N01} \textbf{Nonclaim.} No cohomology, spectral sequence, or homotopy-
coherent descent is available until its defining identities and data are
supplied.

Even when a filtration spectral sequence is available, it is built from the
filtered complex alone and does not see \(\Phi_1\) (compare CH04-N02).

\subsubsection{16.5 Chain contractions and perturbation}\label{chain-contractions-and-perturbation}

Presentation changes between cochain complexes are often contractions: a smaller
complex that retains the cohomology, with explicit comparison maps. This section
fixes the convention and records what those maps do and do not preserve.

\texttt{WDP-CH16-D02} \textbf{Definition (chain contraction and retraction defect).} A
contraction of a finite cochain complex \(C\) onto \(D\) over a declared field
consists of chain maps \(i:D\to C\) and \(r:C\to D\) with \(ri=1_D\), together with a
map \(h\) of degree \(-1\) satisfying \(1_C-ir=dh+hd\). It is normalized when \(hi=0\),
\(rh=0\), and \(hh=0\). Its retraction defect is \(\varepsilon=1-ir\).

\texttt{WDP-CH16-P03} \textbf{Finite proposition (retraction defect).} For chain maps \(i\) and
\(r\) with \(ri=1\) on a finite complex over a field, \(ir\) is idempotent,
\(\varepsilon=1-ir\) commutes with \(d\) and so is closed in the degree-zero window of
the endomorphism complex, a homotopy \(h\) with \(\varepsilon=dh+hd\) exists exactly
when the class of \(\varepsilon\) vanishes, which holds exactly when \(\ker r\) is
acyclic, and any such \(h\) can be modified to satisfy the normalization
conditions; the retraction class is therefore \texttt{ExactClass} or \texttt{ClosedResidual}
and never \texttt{ActiveClosureFailure}.

\emph{Proof.} \((ir)(ir)=i(ri)r=ir\). Since \(i\) and \(r\) are chain maps,
\(d\varepsilon=d-idr=\varepsilon d\). In the endomorphism complex, the differential
of a degree-zero map \(f\) is \(df-fd\), so \(\varepsilon\) is a cocycle, and it is a
coboundary exactly when some \(h\) gives \(\varepsilon=dh+hd\). Over a field
\(C=\operatorname{im}i\oplus\ker r\) as complexes, and \(\varepsilon\) is the
projection onto \(\ker r\); it is nullhomotopic exactly when \(\ker r\) is acyclic.
Given any \(h\), set \(\pi=1-ir\), \(h'=\pi h\pi\), and \(h''=h'dh'\); these satisfy the
homotopy equation together with the three normalization conditions.
\(\square\)

\texttt{WDP-CH16-C01} \textbf{Conditional result (perturbation lemma, imported standard).}
With the convention \(1-ir=dh+hd\), if \((i,r,h)\) is a normalized contraction,
\(\delta\) has degree one with \((d+\delta)^2=0\), and \(1+\delta h\) is invertible, then
with \(A=(1+\delta h)^{-1}\delta\) the maps \(d_D+rAi\), \(i-hAi\), \(r-rAh\), and
\(h-hAh\) form a normalized contraction of \((C,d+\delta)\) onto \(D\)
{[}REF-HOM-007{]}.

The signs depend on the convention. The source that proposed this layer used
the formulas with \((1-\delta h)^{-1}\) that belong to the opposite convention
\(ir-1=dh+hd\); with its own convention they fail to give a differential on \(D\).
The executable package checks every perturbed contraction it produces rather
than trusting the formula. The same source listed side conditions for a
transfer step that omitted the homotopy equation itself; with \(C=\mathbb Q^2\),
\(D=\mathbb Q\), \(d=0\), and \(h=0\) every listed condition holds while the
cohomology dimensions differ.

\subsubsection{16.6 Truncation leakage}\label{truncation-leakage}

Restricting an operator to a subspace is a presentation change of a different
kind. Let \(A\) act on a finite-dimensional inner-product space and let \(P\) be an
orthogonal projector.

\texttt{WDP-CH16-P04} \textbf{Finite proposition (truncation leakage).} For a linear operator
\(A\) on a finite-dimensional inner-product space over \(\mathbb Q\), \(\mathbb Q(i)\),
\(\mathbb R\), or \(\mathbb C\) and an orthogonal projector \(P\),
\(PA^{*}AP=(PAP)^{*}(PAP)+L_P\) with \(L_P=((1-P)AP)^{*}((1-P)AP)\) positive
semidefinite; \(L_P=0\) exactly when the range of \(P\) is \(A\)-invariant; and a
vector \(\psi\) in the range of \(P\) with \(PAP\psi=0\) satisfies \(A\psi=0\) exactly
when \(\langle\psi,L_P\psi\rangle=0\), so \(L_P=0\) is sufficient but not necessary
for a reduced zero to be genuine.

\emph{Proof.} Insert \(1=P+(1-P)\) between \(A^{*}\) and \(A\) and use \(P^2=P=P^{*}\) and
\((1-P)^2=1-P\). The range of \(P\) is invariant exactly when \((1-P)AP=0\), which is
exactly when \(L_P=0\). For \(\psi=P\psi\) with \(PAP\psi=0\),
\(A\psi=(1-P)AP\psi\), whose squared norm is \(\langle\psi,L_P\psi\rangle\).
\(\square\)

On \(\mathbb Q(i)^3\) with \(P=\operatorname{diag}(1,1,0)\) and \(A=e_3e_2^{T}\),
\(L_P=\operatorname{diag}(0,1,0)\ne0\), yet \(e_1\) is a genuine zero; \(e_2\) is a
reduced zero that is not genuine. Leakage-free is also not the same as descent
through \(P\): for nonsymmetric \(A\), \((1-P)AP=0\) does not force \(PA(1-P)=0\).

\texttt{WDP-CH16-N02} \textbf{Nonclaim.} The leakage term does not classify vectors: a
vector in the range of \(P\) can satisfy \(A\psi=0\) while \(L_P\) has arbitrarily
large norm, and no Hamiltonian, energy, zero-mode, probability-leak, or
superspace interpretation is derived from the identity.

\subsubsection{16.7 Composition curvature of a restriction system}\label{composition-curvature-of-a-restriction-system}

Restriction data often come without a composition law. Let \((P,\le)\) be a
finite poset. For each \(x\in P\) declare a vector space \(F(x)\) over a declared
field, and for each strict comparable pair \(a<c\) a linear map
\(r_{c,a}:F(c)\to F(a)\). Nothing is assumed about how these maps compose. Index
cochains by strict chains, valued in the fiber at the bottom:

\[
C^0=\prod_{x}F(x),\quad C^1_\le=\prod_{a<c}F(a),\quad
C^2_\le=\prod_{a<b<c}F(a),\quad C^3_\le=\prod_{a<b<c<d}F(a),
\]

with

\[
(d_0s)_{c,a}=r_{c,a}s(c)-s(a),\qquad
(d_1g)_{c,b,a}=r_{b,a}g_{c,b}-g_{c,a}+g_{b,a},
\]

\[
(d_2h)_{d,c,b,a}=r_{b,a}h_{d,c,b}-h_{d,c,a}+h_{d,b,a}-h_{c,b,a}.
\]

\texttt{WDP-CH16-P05} \textbf{Finite proposition (composition curvature).} Put
\(K_{c,b,a}=r_{b,a}r_{c,b}-r_{c,a}\) for every strict triple \(a<b<c\). For every
\(s\in C^0\), every \(g\in C^1_\le\), and every strict chain \(a<b<c<d\),

\[
(d_1d_0s)_{c,b,a}=K_{c,b,a}\,s(c),\qquad
(d_2d_1g)_{d,c,b,a}=K_{c,b,a}\,g_{d,c}.
\]

Hence \(d_1d_0=0\) exactly when \(K=0\), which holds exactly when
\(r_{b,a}r_{c,b}=r_{c,a}\) on every strict triple; and then \(d_2d_1=0\) as well.

\emph{Proof.} Using linearity of \(r_{b,a}\),
\((d_1d_0s)_{c,b,a}=r_{b,a}(r_{c,b}s(c)-s(b))-(r_{c,a}s(c)-s(a))+(r_{b,a}s(b)-s(a))
=(r_{b,a}r_{c,b}-r_{c,a})s(c)\). Since \(C^0\) is the full product, \(s(c)\) ranges
over all of \(F(c)\), so \(d_1d_0=0\) forces every \(K_{c,b,a}=0\). For the second
identity, expand \(d_2\) on the four components of \(d_1g\): the terms in
\(g_{d,b}\), \(g_{d,a}\), \(g_{c,b}\), \(g_{c,a}\), and \(g_{b,a}\) cancel in pairs and
leave \((r_{b,a}r_{c,b}-r_{c,a})g_{d,c}\). \(\square\)

On a chain \(a<b<c\) with one-dimensional fibers, \(r_{c,b}=2\), \(r_{b,a}=3\), and
\(r_{c,a}=5\), the curvature is \(K=3\cdot2-5=1\). The section
\((s(a),s(b),s(c))=(0,0,1)\) has \((d_0s)_{c,b}=2\), \((d_0s)_{b,a}=0\), and
\((d_0s)_{c,a}=5\), so \((d_1d_0s)_{c,b,a}=3\cdot2-5+0=1=K\,s(c)\). Declaring
\(r_{c,a}=6\) instead makes the triple flat. A diagnostic that reads only the
cover pairs \(a\prec b\prec c\) never sees \(r_{c,a}\) and cannot detect \(K\).

This is Outcome 2 of §16.2 in its smallest form: \(K\) is typed curvature, and
\(d_1d_0=K\) is exact bookkeeping, not a cohomology class. On the flat locus,
with identities adjoined by definition (\(r_{a,a}=1\)), the declared maps satisfy
the identity and composition laws, so they form a functor on
\((P,\le)^{\rm op}\) {[}REF-MATH-004{]}, and the differentials square to zero
through degree three. The identities use only additivity of \(r_{b,a}\), so
abelian groups may replace vector spaces. No norm is used, and no \(H^2\) class,
no descent of \(d_1\) to a quotient of \(C^1_\le\), and no higher-coherent repair
is claimed; CH16-N01 still applies. The absorbed archived working-paper source 009
checked the first identity and the flat equivalence in Lean for a single triple
of maps on one carrier type, and the case \(g=d_0s\) of the second identity for a
single quadruple; that module is not ported and earns no Lean binding here.

\begin{center}\rule{0.5\linewidth}{0.5pt}\end{center}

\subsection{Chapter 17 --- Finite weighted histories}\label{chapter-17-finite-weighted-histories}

\subsubsection{17.1 Path carrier}\label{path-carrier}

Let \(G=(V,E)\) be a finite directed graph with a declared admissibility predicate
on paths. For vertices \(a,b\), let \(\Omega_{a,b}^{(n)}\) be the finite set of
admissible paths of length \(n\) from \(a\) to \(b\). A history is a path, not an
instantaneous state.

A length cutoff \(n\le N\) creates a finite set, but the union of length-bounded
paths is not generally closed under concatenation: two paths of length at most
\(N\) can concatenate to a path longer than \(N\). Thus it is not a category unless
the composition law is restricted or the graph is suitably acyclic.

\subsubsection{17.2 Semiring weights}\label{semiring-weights}

Let \((R,\oplus,\otimes,0,1)\) be a declared semiring and assign edge weights
\(w:E\to R\). The weight of a path
\(p=e_1\cdots e_n\) is

\[
w(p)=w(e_1)\otimes\cdots\otimes w(e_n),
\]

with the empty path assigned \(1\). The finite length-\(n\) history aggregate is

\[
Z_{a,b}^{(n)}=\bigoplus_{p\in\Omega_{a,b}^{(n)}}w(p).
\]

\texttt{WDP-CH17-P01} \textbf{Finite proposition (matrix path sum).} For the full family of
directed graph paths of length \(n\) in the stated graph, with weighted adjacency
matrix \(M\) and conventional ordered multiplication matching path composition,

\[
(M^n)_{a,b}=Z_{a,b}^{(n)}.
\]

\emph{Proof sketch.} Induct on \(n\). The step expands
\((M^{n+1})_{a,b}=\bigoplus_c(M^n)_{a,c}\otimes M_{c,b}\) and distributes over the
last edge of every length-\((n+1)\) path. \(\square\)

For an admissibility predicate compiled into a finite augmented-state graph,
let \(\lambda_a\) be the declared row of admissible augmented start states over
original vertex \(a\) and \(\rho_b\) the declared column of accepting augmented end
states over original vertex \(b\), including any start/end weights. If augmented
paths are in weight-preserving bijection with admissible original paths, then
the correct aggregate is \(\lambda_a M^n\rho_b\), generally a weighted block sum,
not a single entry indexed by the original vertices. An arbitrary path-global
predicate is not counted by a matrix power unless such finite-state compilation
and endpoint maps are supplied.

The theorem is finite. An all-length sum
\(I\oplus M\oplus M^2\oplus\cdots\) requires nilpotence or a finite acyclic
truncation, a complete or \(*\)-continuous semiring, or an analytic convergence
structure. A bare Conway/Kleene star operation can define an algebraic \(M^*\)
without identifying it with that countable indexed sum, and a bare semiring
does not evaluate the infinite series.

\subsubsection{17.3 Families of meaning}\label{families-of-meaning}

The same path formula has different interpretations for different semirings.

\begin{itemize}
\tightlist
\item
  Boolean: reachability.
\item
  Natural numbers: path count.
\item
  Nonnegative reals: positive path weight.
\item
  Tropical min-plus: optimal cost.
\item
  Complex numbers: amplitude-like cancellation, after complex structure is
  declared.
\end{itemize}

Changing the semiring changes the theory. A proof over nonnegative reals cannot
be advertised as a theorem about complex interference without a transport
argument.

\subsubsection{17.4 Cancellation and zerosumfree carriers}\label{cancellation-and-zerosumfree-carriers}

A semiring is zerosumfree when \(a\oplus b=0\) implies \(a=b=0\). In such a carrier,
nonzero path contributions cannot cancel to zero just through addition.

\texttt{WDP-CH17-P02} \textbf{Conditional result.} In a zerosumfree semiring, a
finite sum is zero only if every summand is zero. The binary property extends by
induction on the number of summands, using associativity.

Complex numbers are not zerosumfree: \(1+(-1)=0\). That cancellation is a
necessary ingredient for some amplitude phenomena, but it is not sufficient
for quantum theory. Quadratic readout, tensor composition, contextual structure,
and dynamics remain separate.

\subsubsection{17.5 Obstruction action and history dependence}\label{obstruction-action-and-history-dependence}

Suppose a local action assigns \(A(e)\in\mathbb R\) and define

\[
A(p)=\sum_{e\in p}A(e),\qquad
w_\beta(p)=e^{-\beta A(p)}.
\]

This is a declared positive-real weighting. If \(A=d_0f\) on edges, then
\(A(p)=f(b)-f(a)\) depends only on endpoints, so all \(a\)-to-\(b\) paths have equal
weight. Nonzero periods can make the weighting genuinely history-dependent.

The converse requires care: equal weights at one chosen \(\beta\) may hide action
differences if the exponential is noninjective in another coefficient system.
Over real exponentials with \(\beta\ne0\), equality of weights implies equality of
actions.

\subsubsection{17.6 Worked semiring example}\label{worked-semiring-example}

Take two paths from \(a\) to \(b\): \(p\) with edge weights \(2,3\) and \(q\) with one
edge weight \(5\).

\begin{itemize}
\tightlist
\item
  Over \((\mathbb N,+,\cdot)\), the total is \(2\cdot3+5=11\).
\item
  Over min-plus, interpreting the listed numbers as costs, the total is
  \(\min(2+3,5)=5\).
\item
  Over complex numbers, choosing path weights \(1\) and \(-1\) gives total \(0\).
\end{itemize}

The graph is unchanged; the aggregate's meaning is not.

\subsubsection{17.7 Contextual histories, free differences, and algebraic tensor composition}\label{contextual-histories-free-differences-and-algebraic-tensor-composition}

The path carrier above presupposes a graph and a weight semiring. The more
primitive construction begins with an admitted typed generating signature
\(\Sigma=(O,E,s,t)\) and freely forms well-typed histories. Given an
endpoint-indexed family of protected questions, define two parallel histories
to be equivalent when every compatible prefix and suffix gives the same answer
to every protected question. This is the greatest two-sided congruence with
that preservation property. Equal current reports alone are insufficient: an
allowed continuation can distinguish them. The full construction and the
counterexample are proved in
\href{docs/derivations/generative_histories_contextual_compression_derivation.md}{\texttt{generative\_histories\_contextual\_compression\_derivation.md}}.

For a declared generator carrier \(X\), finite alternative formation is the free
commutative monoid \(\mathbb N^{(X)}\). Its signed completion is the free
abelian group \(\mathbb Z^{(X)}\). The canonical map
\(\mathbb N^{(X)}\to\mathbb Z^{(X)}\) has the Grothendieck extension property:
every additive map from positive alternatives into an abelian group extends
uniquely to signed differences. Applying this construction to every
contextual-history hom-set and extending composition bilinearly gives the free
preadditive envelope of the contextual history category.

\texttt{GEN-WDP-58} \textbf{Finite algebraic proposition (free signed history
completion).} For every admitted generator type \(X\), the alternative carrier
and signed carrier satisfy the free commutative-monoid and free abelian-group
initial factorisation properties. Every signed class is a difference of two positive
finite alternatives, and every alternative-monoid map into an abelian group
has a unique group-homomorphic extension.

For two free signed carriers, explicitly requiring independent composition to
be additive in each argument yields

\[
\mathbb Z^{(X)}\otimes_{\mathbb Z}\mathbb Z^{(Y)}
\cong \mathbb Z^{(X\times Y)}.
\]

\texttt{GEN-WDP-59} \textbf{Finite algebraic proposition (biadditive tensor
factorization).} Every biadditive map
\(B:\mathbb Z^{(X)}\times\mathbb Z^{(Y)}\to G\) into an abelian group factors
through a unique group homomorphism
\(\overline B:\mathbb Z^{(X\times Y)}\to G\), with
\(\overline B(a\otimes b)=B(a,b)\). The proof defines the factor on pair
generators and applies the free-group uniqueness theorem successively in the
right and left inputs.

These results derive the algebraic tensor object only after biadditivity has
been declared. They do not derive a Hilbert tensor product or a physical
independence principle. A determinant-one \(2\times2\) integer tensor is already
non-pure, but algebraic non-factorability is not yet physical entanglement.
Complex scalars, a dagger, positivity, norm completion, normalized states and
Born readout remain separate interfaces. Full proofs, exact information-loss
boundaries, core-only Lean, and replay evidence are in
\href{docs/derivations/additive_tensor_completion_derivation.md}{\texttt{additive\_tensor\_completion\_derivation.md}}.

\subsubsection{17.8 Dagger extension and the observer Gram radical}\label{dagger-extension-and-the-observer-gram-radical}

A dagger is additional reversal data. On a typed generating signature it
sends every generator \(e:x\to y\) to an involutive reverse
\(e^\dagger:y\to x\) and reverses the order of a history. It descends through
the protected contextual quotient if and only if its equivalence relation is
invariant under reversal. Closure of the protected questions under reversal
is a sufficient condition for this invariance. Once supplied on a generator carrier, free-group
uniqueness extends it to an additive involution on signed differences.

\texttt{GEN-WDP-60} \textbf{Finite algebraic proposition (dagger extension).} A declared
generator involution has a unique additive extension \(D\) to the free signed
carrier, with \(D^2=I\). For componentwise involution on generator pairs,

\[
(a\otimes b)^\dagger=a^\dagger\otimes b^\dagger.
\]

The tensor equation follows because both sides are biadditive and agree on
every pair of generator units. It does not assert that a native guarded edit
has an executable reverse.

Now supply an additive observer
\(L:\mathbb Z^{(X)}\to\mathbb Z^m\). Its coordinate Gram form is

\[
g_L(a,b)=L(a)^T L(b).
\]

\texttt{GEN-WDP-61} \textbf{Finite algebraic proposition (observer Gram radical).} The form
\(g_L\) is symmetric, biadditive and positive semidefinite, and

\[
\operatorname{Rad}(g_L)=\ker L,
\qquad
L(u)=L(v)\iff g_L(u-v,u-v)=0.
\]

The proof is the exact sum-of-squares identity
\(g_L(a,a)=\sum_iL(a)_i^2\). Quotienting by the radical produces the realized
image geometry and removes exactly the observer-blind signed differences.
This is relative to the supplied observer. Distinct positive forms already
exist in rank one, while rank one over the integers admits no endomorphism
squaring to \(-I\). Thus dagger plus positivity does not select complex
structure, Hilbert completion or probability readout. The full derivation,
Lean scope and replay are in
\href{docs/derivations/dagger_observer_pairing_derivation.md}{\texttt{dagger\_observer\_pairing\_derivation.md}}.

\subsubsection{17.9 Finite Hermitianization from a checked complex structure}\label{finite-hermitianization-from-a-checked-complex-structure}

Let the realized observer carrier be rationalized to a finite rational vector
space \(V\) with Gram form \(g\). Complex structure is a new admitted datum: a
rational-linear map \(J:V\to V\) satisfying

\[
J^2=-I,
\qquad
g(Jx,y)=-g(x,Jy).
\]

The second law is equivalent to \(J\) preserving \(g\). In coordinates the checker
recomputes the observer Gram matrix \(G=L^TL\) and verifies \(J^2=-I\),
\(J^TG=-GJ\), and \(J^TGJ=G\).

\texttt{GEN-WDP-62} \textbf{Finite conditional theorem (Gaussian-rational carrier).} Define

\[
(a+ib)x=ax+bJx.
\]

Then \(V\) is a vector space over \(\mathbb Q(i)\). Scalar associativity follows
exactly from rational linearity and \(J^2=-I\); no analytic completion is used.

\texttt{GEN-WDP-63} \textbf{Finite conditional theorem (sharp Hermitianization).} Define

\[
h(x,y)=g(x,y)+i g(Jx,y).
\]

With the first slot conjugate linear, \(h\) is a positive-semidefinite Hermitian
form and

\[
\operatorname{Rad}(h)=\operatorname{Rad}(g).
\]

For an observer Gram form, \(h(x,x)=0\) exactly when \(Lx=0\). Thus the complex
extension preserves the observer's distinction boundary. On the rational
plane, \(J(x_0,x_1)=(-x_1,x_0)\) gives the standard Gaussian-rational pairing.
The equally admissible orientation \(-J\) has the same positive \(g\) but conjugates
the off-diagonal Hermitian values, proving that positivity does not select
\(J\). The full proof, refusal example, Lean formalization and source-bound
replay are in
\href{docs/derivations/finite_hermitianization_derivation.md}{\texttt{finite\_hermitianization\_derivation.md}}.

\subsubsection{17.10 Certified finite states, effects and readout}\label{certified-finite-states-effects-and-readout}

Hermitianization supplies a carrier and pairing, but no normalized state or
effect cone. Assume a supplied standard Hermitian coordinate system over \(\mathbb Q(i)\);
this need not follow from Hermitianization over that same field. For the
explicit rectangular Gram-factor class, let \(B\ne0\) present

\[
\rho_B=\frac{B^\dagger B}{\lVert B\rVert_F^2},
\]

and let a binary effect be presented by factors \(C,D\) satisfying
\(C^\dagger C+D^\dagger D=I\).

\texttt{GEN-WDP-64} \textbf{Finite conditional theorem (certified state/effect readout).}
The admitted trace readout has the exact factor form

\[
\operatorname{tr}(\rho_B C^\dagger C)
=\frac{\lVert CB^\dagger\rVert_F^2}{\lVert B\rVert_F^2}.
\]

The accepted and rejected values are nonnegative rationals summing to one.
The proof is finite trace cyclicity plus the submitted completeness equation;
the checker reconstructs every Gram matrix and refuses a zero state or
incomplete effect. The certificate language covers matrices with supplied
\(\mathbb Q(i)\) Gram factors, not every positive matrix over a larger scalar
field.

\texttt{GEN-WDP-65} \textbf{Finite countermodel (readout nonselection).} For
\(0\le\lambda\le1\),

\[
p_\lambda(\rho,E)=\lambda\operatorname{tr}(\rho E)
 +(1-\lambda)\frac{\operatorname{tr}E}{n}
\]

is positive, normalized on every finite POVM, phase invariant and affine for
\(n\ge1\). For \(n\ge2\), a sharp state/effect pair gives
\(p_\lambda=\lambda+(1-\lambda)/n\), which differs from the undeformed value \(1\)
when \(\lambda\ne1\). At \(n=1\), every normalized state is \([1]\) and an effect
\([e]\) gives \(p_\lambda=e\) for every \(\lambda\). Thus the theorem constructs an exact Born-form interface after
that readout is admitted; positivity and normalization do not uniquely select
it. Full proof, Lean scope, refusal controls and replay are in
\href{docs/derivations/finite_state_effect_readout_derivation.md}{\texttt{finite\_state\_effect\_readout\_derivation.md}}.

\subsubsection{17.11 Certified finite processes and composition}\label{certified-finite-processes-and-composition}

For an ensemble of state factors \(B_s\), let

\[
T=\sum_s\lVert B_s\rVert_F^2>0,
\qquad
\rho=T^{-1}\sum_sB_s^\dagger B_s.
\]

A finite process certificate is a labelled family \((K_o)\) satisfying
\(\sum_oK_o^\dagger K_o=I\). Its branch factors are
\(B'_{s,o}=B_sK_o^\dagger\).

\texttt{GEN-WDP-66} \textbf{Finite conditional theorem (certified process composition).}
Every branch mass

\[
T_o=\sum_s\lVert B_sK_o^\dagger\rVert_F^2
\]

is nonnegative and \(\sum_oT_o=T\). The unconditioned output is
\(\sum_oK_o\rho K_o^\dagger\). Applying \((K_o)\) and then \((L_r)\) has composite
factors \(L_rK_o\); complete families on admitted independent carriers tensor;
and a supplied unitary factor preserves mass and is reversed by its adjoint.
These are exact algebraic statements over \(\mathbb Q(i)\).

\texttt{GEN-WDP-67} \textbf{Finite countermodel (channel/instrument nonselection).} The
one-outcome instrument \(\{I\}\) and the two-outcome instrument
\(\{(3/5)I,(4/5)I\}\) induce the same identity channel. Their outcome reports
are respectively \((1)\) and \((9/25,16/25)\). A presentation retaining only the
unconditioned channel therefore loses future conditioning and branch questions.

\texttt{GEN-WDP-68} \textbf{Exact scalar obstruction.} No rational \(r\) satisfies
\(r^2=1/2\). The standard Hadamard matrix and the usual T phase consequently do
not lie in the current exact \(\mathbb Q(i)\) matrix language. This does not rule
out approximate synthesis or a different rational gate set; those require a
topology, error metric, alphabet and compilation theorem. Full proofs, exact
replay, refusal controls and Lean scope are in
\href{docs/derivations/finite_process_composition_derivation.md}{\texttt{finite\_process\_composition\_derivation.md}}.
Chapter 19A distinguishes a conventional matrix from its ray action:
the Gaussian Hadamard there realizes that ray action exactly, while the
T eigenvalue ratio and the order-three Fourier requirements remain genuine
field obstructions.

\texttt{WDP-CH17-O01} \textbf{Finite proposition (resolved finite target).} For a finite rational row-source transition matrix and a nonempty block partition, representative-independent endpoint-block path sums at every finite length exist exactly when outgoing block sums are constant within each source block, equivalently M H = H B; B is unique and finite partition enumeration is exhaustive.
The proof, native ownership and scope are given in Appendix I.

\emph{Partial Characterization and Mechanized Path Algebra.}

\begin{enumerate}
\def\labelenumi{\arabic{enumi}.}
\tightlist
\item
  \textbf{Semiring Lumpability Condition:} Let \(G = (V, E, W)\) be a finite directed graph
  with transition matrix \(M \in \operatorname{Mat}_{|V| \times |V|}(R)\) over a
  semiring \((R, \oplus, \otimes, 0, 1)\). Let \(\mathcal P = \{V_1, \dots, V_k\}\)
  be a partition of \(V\) with block projection matrix \(P \in \operatorname{Mat}_{k \times |V|}(\{0, 1\})\)
  (where \(P_{I, u} = 1 \iff u \in V_I\)) and transposed membership matrix
  \(P^T \in \operatorname{Mat}_{|V| \times k}(\{0, 1\})\). The representative-independent
  outgoing-block condition requires that for each source vertex \(u \in V_I\) and target
  block \(V_J\), the aggregate transition weight
  \(\bigoplus_{v \in V_J} M_{u, v} = \bar M_{I, J}\) is independent of the choice of
  representative \(u \in V_I\). In matrix form, evaluating the \((u, J)\)-entry shows
  this corresponds to the intertwining relation \(M P^T = P^T \bar M\). (The dual
  relation \(P M = \bar M P\) corresponds to pushforward of distribution column vectors
  \(\bar x = P x\) under the stated matrix convention, commuting via \(P(M x) = \bar M(P x)\).)
  Zerosumfreeness prevents nontrivial
  cancellation to zero (\(a \oplus b = 0 \implies a = 0 \land b = 0\)), precluding
  paths with nonzero weights from destructively canceling to the zero element,
  though it does not in general force equality between arbitrary sums to be termwise
  rigid without additional cancellative or idempotence assumptions. A comprehensive
  classification of coarse-grainings preserving arbitrary path aggregates across
  general semirings remains a broader research problem. Claim CH17-O01
  selects the rational endpoint-block aggregate, whose complete finite
  criterion and partition classification are proved in Appendix I.
\item
  \textbf{Mechanization of Finite Path Distributivity and Coboundary Holonomy:} In
  \texttt{lean4/WitnessedDifference/SemiringTransport.lean}, we formally verify under
  \texttt{FWSemiring} that finite sums distribute on the right:
  \texttt{WitnessedDifference.lsum\_mul\_right} establishes \(\operatorname{mul}(\operatorname{lsum} xs, y) = \operatorname{lsum}(xs.\operatorname{map}(\lambda x \implies \operatorname{mul} x y))\).
  Furthermore, \texttt{coboundary\_trivial\_holonomy} demonstrates that coboundary connection
  potentials telescope along walks and evaluate to 0 around any closed walk.
  Neither declaration involves a partition or a coarse-graining; the lumpability
  condition of item 1 is not mechanized.
\end{enumerate}

\texttt{WDP-CH17-N01} \textbf{Nonclaim.} A finite semiring path sum is not a continuum path
integral, a measure over all trajectories, or the Feynman amplitude rule
{[}REF-AUT-007{]}.

\begin{center}\rule{0.5\linewidth}{0.5pt}\end{center}

\subsection{Chapter 18 --- Quantum reconstruction as an audit}\label{chapter-18-quantum-reconstruction-as-an-audit}

\subsubsection{18.1 The benchmark question and candidate finite operational axiom package}\label{the-benchmark-question-and-candidate-finite-operational-axiom-package}

The correct question is not whether the notation resembles quantum mechanics.
It is whether a stated set of operational principles singles out complex
finite-dimensional quantum theory from a comparison class that includes
classical and post-quantum models {[}REF-QM-003; REF-QM-004; REF-QM-005;
REF-QM-006{]}.

To make this audit mathematically precise without smuggled assumptions, we
formulate a candidate finite operational axiom package. Axioms W1--W3 abstract
the finite witness relation, presented report carrier, and rational mass structures
introduced in Chapters 1--13, while Axioms W4--W6 are additional specialization
hypotheses introduced to probe the boundary with complex quantum mechanics.
Earlier packets in Chapters 1--13 do not require reversible dynamics (W4) or
quadratic readout (W6); in general, observer updates may be non-injective
descents and readouts may be affine or report-based.

\begin{enumerate}
\def\labelenumi{\arabic{enumi}.}
\tightlist
\item
  \textbf{Axiom W1 (Primitive Difference Carrier):} An admitted finite irreflexive
  witness relation \(K = (X, R, \operatorname{Irr} R)\) with witness
  \(d = (a, b, w : R a b)\).
\item
  \textbf{Axiom W2 (Presented Report Carrier):} A finite carrier presentation
  \(C = (X, Q, \operatorname{pres} : X \to Q)\) where two points share a report
  \(\operatorname{sameReport}(x, y)\) if and only if
  \(\operatorname{pres}(x) = \operatorname{pres}(y)\).
\item
  \textbf{Axiom W3 (Integer/Rational Weights):} A mass assignment
  \(w : Q \to \mathbb{Z}\) (or \(\mathbb{Q}\)) with strictly positive total mass
  \(W = \sum_{q \in Q} w(q) > 0\).
\item
  \textbf{Axiom W4 (Reversible Discrete Dynamics, Specialization):} State updates are bijections
  \(\sigma \in S_Q\) on the report carrier preserving total mass.
\item
  \textbf{Axiom W5 (Free Linearization \& Admitted Sesquilinear Form, Specialization):} The free
  module \(\operatorname{Vec}_K(Q)\) over \(\mathbb{Z}[i]\) (or \(\mathbb{Z}\))
  equipped with the standard positive-definite inner pairing
  \(\langle u, v \rangle = \sum_{q \in Q} \bar{u}_q v_q\).
\item
  \textbf{Axiom W6 (Phase-Blind Quadratic Readout, Specialization):} The coordinatewise readout map
  \(\pi : \operatorname{Vec}_K(Q) \to \mathbb{Z}_{\ge 0}^Q\) given by
  \(\pi(v)_q = |v_q|^2\).
\end{enumerate}

\subsubsection{18.2 The classical-simplex fork and classical representation theorem}\label{the-classical-simplex-fork-and-classical-representation-theorem}

Finite nonnegative normalized witness weights form a simplex when all such
weights and their convex mixtures are admitted. Normalization alone does not
imply nonnegativity: W3 also permits the signed weights \((2,-1)\). A simplex has unique decompositions into extreme points
under the standard finite-dimensional model. The full density-operator state
space of complex quantum theory is nonsimplicial in Hilbert dimension at least
two; the one-dimensional case is a singleton simplex, and commutative classical
subtheories can also be simplicial.

Nonsimpliciality is not a quantum signature by itself.

\texttt{WDP-CH18-P01} \textbf{Finite proposition (a finite-vertex non-simplex polytope).} The convex polytope
with four extreme points \(v_1=(1,0)\), \(v_2=(0,1)\), \(v_3=(-1,0)\), and \(v_4=(0,-1)\) (whose convex hull is a continuum of states generated by finitely many vertices) is
not a simplex: its centre equals both \(\tfrac12v_1+\tfrac12v_3\) and
\(\tfrac12v_2+\tfrac12v_4\). A nonsimplicial finite-dimensional state polytope therefore exists
outside quantum theory, and nonsimpliciality alone is not a quantum
discriminator.

\emph{Proof.} Both mixtures equal \((0,0)\), and the two decompositions use disjoint
sets of extreme points. The square, with affine effects, is the state space of
the gbit of generalized probabilistic theories {[}REF-QM-005{]}, which is not a
quantum state space. \(\square\)

The centre has a one-parameter family of decompositions,
\(t(v_1+v_3)+(\tfrac12-t)(v_2+v_4)\), while points on an edge still decompose
uniquely.

Under Axioms W1--W6, the reachable representation theorem establishes that finite
reversible discrete dynamics admits a faithful representation within an integer
permutation-isometry envelope:

\textbf{Theorem (Representation in Permutation-Isometry Envelope).}
Let \((X, Q, \operatorname{pres})\) be a finite presented carrier satisfying
Axioms W1--W4 with integer weights \(w : Q \to \mathbb{Z}\) and total mass
\(W = \sum_{q \in Q} w(q) > 0\). Every reversible report update \(\sigma \in S_Q\)
admits a faithful representation as an isometric automorphism \(P_\sigma \in
\operatorname{GL}_{|Q|}(\mathbb{Z})\) preserving the standard Euclidean form
\(\langle u, v \rangle = u^T v\) with two-sided inverse \(P_{\sigma^{-1}} = P_\sigma^T\),
and identically preserving total mass \(\sum_{q \in Q} (P_\sigma w)_q = W\).

\emph{Proof.} Let \(|Q| = n\) with standard basis vectors \(e_q\) for \(q \in Q\). For any
\(\sigma \in S_Q\), define the linear operator \(P_\sigma\) by \(P_\sigma e_q = e_{\sigma(q)}\).
Since \(\sigma\) is a bijection, \(P_\sigma\) is a permutation matrix with entries
in \(\{0, 1\} \subset \mathbb{Z}\). For any vectors \(u, v \in \mathbb{Z}^n\),
\(\langle P_\sigma u, P_\sigma v \rangle = \sum_{q} u_{\sigma(q)} v_{\sigma(q)} = \sum_{q'} u_{q'} v_{q'} = \langle u, v \rangle\),
so \(P_\sigma\) is an exact isometry with \(P_\sigma^T P_\sigma = I_n\). Furthermore,
\(\sum_{q \in Q} (P_\sigma w)_q = \sum_{q \in Q} w_{\sigma^{-1}(q)} = \sum_{q' \in Q} w_{q'} = W\).
Distinct permutations act differently on some basis vector, so the action of
\(S_Q\) on the report carrier \(Q\) embeds faithfully into the integer
permutation-isometry envelope, preserving total mass. This does not make the
presentation \(\operatorname{pres}:X\to Q\) injective. \(\square\)

\emph{Remark on Stochastic vs.~Signed Valuations.} Axiom W3 permits signed integer
weights with positive total (\(W > 0\)) without requiring pointwise nonnegativity.
For example, the weight vector \(w = (2, -1)\) satisfies \(\sum_q w(q) = 1 > 0\) but
contains a negative weight, so it is a signed discrete valuation rather than a
probability state on a simplex. Consequently, the theorem establishes a
permutation-isometry representation over \(\mathbb{Z}\), not a classical stochastic
or probabilistic realization; a stochastic realization requires the additional
standing hypothesis of pointwise nonnegativity \(w(q) \ge 0\).

\textbf{The classical corner.} The theorem above assumes Axiom W4, that reversible
updates are bijections of \(Q\). The next four propositions derive that assumption
from different premises, show which premise does the work, and locate where the
package stops being classical. They are mechanized in
\texttt{lean4/WitnessedDifference/ClassicalCorner.lean}, a Mathlib-backed module.
Other modules also import Mathlib; a standalone core-only replay therefore
does not establish compilation of this result or of the full package.

\texttt{WDP-CH18-P02} \textbf{Finite proposition (nonnegative weights force permutations).}
Let \(Q\) be a finite report carrier and \(W > 0\) a rational number, and call
\(w : Q \to \mathbb{Q}\) a \emph{state of mass \(W\)} when \(w(q) \ge 0\) for every \(q\) and
\(\sum_q w(q) = W\). A linear update \(T\) sends every state of mass \(W\) to a state
of mass \(W\) if and only if \(T\) is column-stochastic. It has an inverse \(S\) with
the same property (\(ST = I\)) if and only if \(T = P_\sigma\) for some
\(\sigma \in S_Q\). Once W3 carries pointwise nonnegativity, the content of W4 is
a theorem for linear updates that are reversible within the states.

\emph{Proof.} Applying \(T\) to \(W e_j\) shows that column \(j\) of \(T\) is nonnegative with
sum \(1\); conversely, a column-stochastic matrix preserves nonnegativity and total
mass. Let \(ST = I\) with both matrices column-stochastic. For \(i \ne j\) the entry
\((ST)_{ij} = \sum_k S_{ik} T_{kj} = 0\) is a sum of nonnegative terms, so
\(S_{ik} T_{kj} = 0\). Column \(j\) of \(T\) has an entry \(T_{kj} > 0\); column \(k\) of
\(S\) then vanishes off row \(j\) and, summing to \(1\), equals \(e_j\). Since \(T\) is
square, \(ST = I\) gives \(TS = I\), and reading \(TS = I\) at column \(k\) gives
\(T e_j = e_k\). Two columns of \(T\) cannot use the
same \(k\), since column \(k\) of \(S\) would then equal two different basis vectors,
so \(j \mapsto k\) is an injection of the finite set \(Q\) into itself, hence a
permutation. Every permutation matrix is column-stochastic with a
column-stochastic inverse. \(\square\)

The same module restates the preservation clauses of the theorem above for
every finite carrier: a permutation matrix over \(\mathbb{Z}\) preserves the
standard form, the total mass, and the \(\ell^1\) norm. \texttt{QuantumImports.lean}
establishes them only for \(Q = \mathrm{Bool}\).

\texttt{WDP-CH18-P03} \textbf{Finite proposition (signed weights do not force permutations).}
With the signed weights that W3 admits, total mass and invertibility do not force
a permutation. The matrices
\(T = \begin{pmatrix} 2 & 1 \\ -1 & 0 \end{pmatrix}\) and
\(T^{-1} = \begin{pmatrix} 0 & -1 \\ 1 & 2 \end{pmatrix}\) both preserve
\(\sum_q w(q)\), \(T\) is not a permutation matrix, and \(T\) sends the basis weight
\((1, 0)\) to the signed weight \((2, -1)\) of the remark above.

\emph{Proof.} Both matrices have unit column sums, \(T^{-1} T = I\), and
\(T_{00} = 2 \notin \{0, 1\}\). \(\square\)

\texttt{WDP-CH18-P04} \textbf{Finite proposition (unitaries over \(\mathbb{Z}[i]\) are
monomial).} Let \(U\) be a \(Q \times Q\) matrix over the Gaussian integers with
\(U^{*} U = I\). Then there are \(\sigma \in S_Q\) and units
\(u_q \in \{\pm 1, \pm i\}\) with \(U e_q = u_q e_{\sigma(q)}\). Consequently the
phase-blind readout of W6 moves by the same permutation,
\(\pi(U v)_{\sigma(q)} = \pi(v)_q\), and no such \(U\) separates two vectors with
equal readout.

\emph{Proof.} The diagonal of \(U^{*} U = I\) gives \(\sum_k N(U_{kj}) = 1\), where
\(N(a + bi) = a^2 + b^2\) is a nonnegative integer. Exactly one entry of column \(j\)
therefore has norm \(1\), which makes it a unit, and the others vanish. If two
columns \(j \ne j'\) had their units in the same row \(k\), the entry
\((U^{*} U)_{j j'} = \overline{U_{kj}}\, U_{kj'}\) would be a product of units and
hence nonzero. So the rows are distinct and \(j \mapsto k\) is a permutation. Then
\((U v)_{\sigma(q)} = u_q v_q\), and \(N\) is multiplicative with \(N(u_q) = 1\).
\(\square\)

Along any finite run of such updates, the readout is the initial readout
relabelled by one permutation (Lean: \texttt{gaussian\_run\_readout}; not a ledger row).
The package is therefore classical under either of two strengthened readings: W3
with pointwise nonnegativity and linear updates reversible within the states
(CH18-P02), or dynamics that preserve the form of W5 over \(\mathbb{Z}[i]\), read
through W6 (CH18-P04). Under either reading every readout the package predicts
comes from a deterministic reversible model on \(Q\), and the phase fiber of
principle P2 in §18.3 is never observed. For W3 as stated, with signed weights,
CH18-P03 shows that the first route fails. Over \(\mathbb{Z}\) the same argument,
with \(N(a) = a^2\), gives the signed permutation matrices (not mechanized).

\texttt{WDP-CH18-P05} \textbf{Finite proposition (where the classical corner ends).} The
unnormalized Hadamard matrix
\(H = \begin{pmatrix} 1 & 1 \\ 1 & -1 \end{pmatrix}\) used for interference in
§18.3 satisfies \(H^{*} H = 2I\). It separates the readout fiber of \((1, 1)\) and
\((1, -1)\), which CH18-P04 shows no unitary over \(\mathbb{Z}[i]\) can do, so it is
not a unitary over \(\mathbb{Z}[i]\). Unitaries that separate readout fibers exist
over finite localizations of the scalars. The rotation
\(R = \tfrac15 \begin{pmatrix} 3 & -4 \\ 4 & 3 \end{pmatrix}\) over
\(\mathbb{Z}[\tfrac15]\) sends the readouts of \((1, 1)\) and \((1, -1)\), both
\((1, 1)\), to \((\tfrac{1}{25}, \tfrac{49}{25})\) and
\((\tfrac{49}{25}, \tfrac{1}{25})\). The matrix
\(M = \tfrac12 \begin{pmatrix} 1+i & 1-i \\ 1-i & 1+i \end{pmatrix}\) over
\(\mathbb{Z}[i][\tfrac12]\) sends the readouts of \((1, 1)\) and \((1, i)\), both
\((1, 1)\), to \((1, 1)\) and \((2, 0)\), and \(M^2\) is the toggle
\(\begin{pmatrix} 0 & 1 \\ 1 & 0 \end{pmatrix}\). Extending the scalars extends the
readout of W6 to rational values.

\emph{Proof.} Direct computation over \(\mathbb{Z}[i]\), \(\mathbb{Q}\), and
\(\mathbb{Q}(i)\). \(\square\)

So unitary interference needs neither the continuum field (import 1) nor the
circle group (import 2). It needs a non-monomial unitary: none exists over
\(\mathbb{Z}[i]\) (CH18-P04), and a finite localization of \(\mathbb{Z}\) or
\(\mathbb{Z}[i]\) already has one. The similitude \(H\) leaves the classical corner
by another route, since it is not an isometry. The square root \(M\) of the toggle
also shows that the missing half-step of principle P11 is missing only among
report permutations. This does not reconstruct quantum theory: CH18-N01 and the
eight interfaces of §18.5 stand.

\textbf{The strengthened package S.} Call S the package W1, W2, W3⁺, W4′, W6′ and W7,
where W3⁺ is W3 with pointwise nonnegative weights, W4′ says that an update is a
linear map that sends states to states (no reversibility is required), W6′ says
that a measurement is a \emph{classical instrument}, a family \((T_o)_{o\in O}\) of
nonnegative updates whose sum is column-stochastic, and W7 says that the
empirical model of the measurement cover has a nonnegative global section, which
by CH19-P05 means it is the image of a table in which every context reads the
same runs. S is not a strengthening of W1--W6: it omits W5, replaces W6 by W6′, and drops
the reversibility of W4. It is a package of hypotheses, not a theorem about the
packets of Chapters 10--13: with CH19-P04's free tuple-formation and repetition
rule, those admit tables with contextual
images, so W7 is a hypothesis on the recorded runs and does not follow from the
packet types.

\texttt{WDP-CH18-P07} \textbf{Finite proposition (stochastic updates).} Over a finite
report carrier, the updates of W4′ are the column-stochastic matrices. They are
closed under composition and contain the identity. A stochastic update whose
entries are all \(0\) or \(1\) is exactly the update \(M_f\) induced by a function
\(f\) of reports, and so is every presentation map. Every stochastic update is a
mixture of function updates,
\(T=\sum_fw_fM_f\) with \(w_f=\prod_jT_{f(j)j}\ge0\) and \(\sum_fw_f=1\). Its units are
the permutations (CH18-P02).

\emph{Proof.} A product of column-stochastic matrices is column-stochastic. A
column-stochastic \(\{0,1\}\)-matrix has exactly one \(1\) in each column, which is
the matrix of a function. Expanding \(\prod_j\sum_iT_{ij}=1\) gives
\(\sum_fw_f=1\), and fixing \(f(k)=i\) in the same expansion gives
\(\sum_{f:f(k)=i}w_f=T_{ik}\) because the other columns still sum to one; this is
\(T=\sum_fw_fM_f\). \(\square\)

\texttt{WDP-CH18-P08} \textbf{Finite proposition (classical facts about S, and their
sharpness).} Let \((T_o)_{o\in O}\) be a classical instrument on \(\mathbb
Q^Q\) and \(p\) a state. The law of an outcome word \(o_1\dots o_n\) is the mass of
\(T_{o_n}\cdots T_{o_1}p\). It is nonnegative; the laws of the words of length \(n\)
add up to the mass of \(p\); and the law of a word is the sum of the laws of its
one-letter extensions. Under W7 the empirical model of the measurement cover has
a nonnegative global section by definition, and every such model on the CHSH
cover has \(\mathrm{CHSH}\le2\). Both hypotheses matter to these conclusions.
Without W3⁺ the signed shear of CH18-P03 has a negative entry, so it is not
stochastic; it preserves mass and has a mass-preserving inverse but is not a
permutation. Without W7 the PR box of CH19-P05 is an empirical model with
\(\mathrm{CHSH}=4\) and no nonnegative global section.

\emph{Proof.} Nonnegativity follows by induction on the word, for a state \(p\ge0\)
and finite \(O\). Because \(\sum_oT_o\) is column-stochastic,
\(\sum_o\sum_i(T_oq)_i=\sum_iq_i\) for every \(q\), which gives consistency and, by
induction on \(n\), normalization. The CHSH bound and the PR-box values are those
of CH19-P05; its outcome-availability and free table-formation hypotheses are
needed only for the equivalence with global tables. \(\square\)

Read in the vocabulary of §19.10, under S a state is a distribution on \(Q\), an
update is a mixture of functions (CH18-P07), an instrument outcome law is a
hidden-Markov law, and a noncontextual model has a hidden-variable
representation. This reading is an interpretation of CH18-P07 and CH18-P08, and
no unrestricted structure-preserving representation of the packets of Chapters 10--13 into a classical joint-run model is asserted here; Appendix J constructs the exact admitted report-instrument quotients.

\texttt{WDP-CH18-O01} \textbf{Conditional result (reconstruction classification and obstruction).} Signed report weights have a classical state representation exactly when every observable block mass is nonnegative; signed instrument updates admit a classical report-instrument quotient exactly under the stated fiber-intertwining and stochastic-quotient conditions. The unrestricted witness/observer/measurement class admits accepted PR tables, refuting reconstruction that preserves their declared probabilities together with classical joint runs or bipartite local Hilbert measurement composition. Classical and quantum alternatives require the explicit additional packages in Appendix J; neither is forced by the original packets.

The original two alternatives were:

\begin{itemize}
\tightlist
\item
  a representation theorem into classical finite stochastic theory; or
\item
  a reconstruction theorem with all additional quantum assumptions exposed.
\end{itemize}

Outcome (a) is proved for two strengthened readings of the package W1--W6 of
§18.1: W3 with pointwise nonnegativity and linear updates reversible within the
states (CH18-P02), and dynamics that preserve the form of W5 over
\(\mathbb{Z}[i]\) (CH18-P04). CH18-P07 and CH18-P08 prove classical facts for the package S,
which drops reversibility and adds classical instruments and joint runs; without
W3⁺ or W7 those conclusions fail. CH18-P03 shows that mass and invertibility
alone do not force permutations, and CH18-P04 shows that the preserved form over
\(\mathbb{Z}[i]\) does. Appendix J resolves the unrestricted classical reading negatively: full signed indicators can be negative, and accepted PR tables admit no nonnegative joint representation. They also violate the imported Hilbert CHSH bound, so the original packet class cannot universally reconstruct standard quantum theory while preserving fixed local factor measurements. The constructive report-partition and instrument-factorization criteria identify the precise classical subclasses. A separately declared Hilbert/matrix-unit package gives a conditional quantum composition theorem, with the probability readout still exposed as an assumption. CH18-P05 remains an earlier algebraic milestone.

\subsubsection{18.3 Discriminator suite and audit table}\label{discriminator-suite-and-audit-table}

\texttt{WDP-CH18-D01} \textbf{Definition (reconstruction audit).} A reconstruction audit is
a table whose rows are candidate principles and whose columns record:

\begin{enumerate}
\def\labelenumi{\arabic{enumi}.}
\tightlist
\item
  the candidate principle;
\item
  its status in the finite obstruction calculus;
\item
  its formal statement and proved theorem or defect locator;
\item
  classical models satisfying it;
\item
  quantum models satisfying it;
\item
  post-quantum countermodels and conclusion.
\end{enumerate}

\emph{Note on claim hierarchy and internal anchors:} The principles P1--P12 and
internal identifiers (\texttt{H006}, \texttt{H013}, \texttt{H021}, etc.) refer to exploratory
computational and formal test fixtures in \texttt{code/witnessed\_difference/quantum.py}
and \texttt{lean4/WitnessedDifference/QuantumImports.lean}. They do not constitute
standalone canonical claim rows in the canonical ledger. The canonical rows
governing this chapter are CH18-D01, CH18-P01 to CH18-P08, the conditional reconstruction classification
CH18-O01, and the nonclaim CH18-N01.

The table below audits the 12 candidate principles across six questions:
nonclassical geometry, contextuality, composition, probability readout,
continuity and reversibility, and purification or dilation. A row labelled
``derived'' means a consequence of its stated specialization assumptions,
including W4--W6 or an additional operator when that row requires one. It does
not mean that the primitive distinction alone supplies those assumptions.
``Refuted'' applies to the specified implication or host model, not to quantum
mechanics or to every possible reconstruction route.

\begin{longtable}[]{@{}
  >{\raggedright\arraybackslash}p{(\linewidth - 10\tabcolsep) * \real{0.1667}}
  >{\raggedright\arraybackslash}p{(\linewidth - 10\tabcolsep) * \real{0.1667}}
  >{\raggedright\arraybackslash}p{(\linewidth - 10\tabcolsep) * \real{0.2500}}
  >{\raggedright\arraybackslash}p{(\linewidth - 10\tabcolsep) * \real{0.1389}}
  >{\raggedright\arraybackslash}p{(\linewidth - 10\tabcolsep) * \real{0.1389}}
  >{\raggedright\arraybackslash}p{(\linewidth - 10\tabcolsep) * \real{0.1389}}@{}}
\toprule\noalign{}
\begin{minipage}[b]{\linewidth}\raggedright
Candidate Principle
\end{minipage} & \begin{minipage}[b]{\linewidth}\raggedright
Status in the finite obstruction calculus
\end{minipage} & \begin{minipage}[b]{\linewidth}\raggedright
Proved Theorem / Defect Locator
\end{minipage} & \begin{minipage}[b]{\linewidth}\raggedright
Classical Model
\end{minipage} & \begin{minipage}[b]{\linewidth}\raggedright
Quantum Model
\end{minipage} & \begin{minipage}[b]{\linewidth}\raggedright
Post-Quantum / Conclusion
\end{minipage} \\
\midrule\noalign{}
\endhead
\bottomrule\noalign{}
\endlastfoot
\textbf{P1: Permutation Invariance} & \textbf{DERIVED} & Form-preserving isometry with two-sided inverse (\texttt{H006.permutation\_isometry}) & Symmetric group \(S_n\) & Permutation matrices \(\subset U(n)\) & Isometry envelope over \(\mathbb{Z}[i]^n\) is the monomial matrices with entries in \(\{\pm 1, \pm i\}\) (CH18-P04 gives one inclusion; the other is immediate); strictly smaller than \(U(n)\). \\
\textbf{P2: Coordinate-readout fiber} & \textbf{CON-SEQUENCE OF W5--W6} & For the admitted quadratic readout \(\pi(v)_q=\vert v_q\vert^2\), equal reports can leave sign or phase freedom (\texttt{H009}). & A declared readout can identify distinct underlying representatives. & On normalized complex vectors with fixed coordinate probabilities and support of size \(r\), the fiber modulo common phase is \((S^1)^{r-1}\); \(\mathbb CP^{n-1}\) is the full pure-state ray space, not this single fiber. & A separately supplied mixer can distinguish some equal-readout vectors. This does not select complex scalars or the quadratic readout. \\
\textbf{P3: Interference Non-Descent} & \textbf{DERIVED (only with \(H\) admitted)} & Hadamard mixing separates equal-weight fiber states; no descent through fiber (\texttt{H010.int\_no\_descent}). \(H^{*}H = 2I\), so \(H\) lies outside W4, and no unitary over \(\mathbb{Z}[i]\) separates a readout fiber (CH18-P04) & Trivial descent (\(f(\pi) = \pi'\)) & Unitary basis rotation \(U \psi\) & Runs on \(\mathbb{Z}\) with \(\pm 1\) once \(H\) is admitted; an isometric separator is unrealizable over \(\mathbb{Z}[i]\) (CH18-P04) and exists over \(\mathbb{Z}[\tfrac15]\) (CH18-P05); complex field \(\mathbb{C}\) is NOT selected. \\
\textbf{P4: Phase Group Selection} & \textbf{REFUTED (derived); IMPORTED} & Sign group \(\{\pm 1\}\) suffices; circle group \(U(1)\) is an explicit import (\texttt{H011Circle}) & Trivial phase \(\{1\}\) & Continuous circle \(U(1)\) & Interference requires only signs; \(U(1)\) requires explicit continuity import. \\
\textbf{P5: Quadratic Born Scaling} & \textbf{DERIVED (discrete; uses \(H\))} & On uniform superpositions \(v = (1, 1)\), Hadamard norm ratio \(\Vert H v\Vert_p^p / \Vert v\Vert_p^p = 2\) holds identically iff \(p=2\); vanishing third-order interference \(I_3 \equiv 0\) follows in the declared additive-amplitude/quadratic-measure model (Sorkin 1994 {[}REF-QM-018{]}; \texttt{H021}) & Linear \(L^1\) norm & \(L^2\) Hilbert norm \(\langle \psi \vert \psi \rangle\) & \(p=2\) uniquely preserves Hadamard scaling on uniform superposition, while \(I_3 \equiv 0\) is a structural consequence of additive amplitudes and quadratic probabilities. \\
\textbf{P6: Gleason Uniqueness} & \textbf{REFUTED (host); IMPORTED (\(n \ge 3\))} & For state \((1,0)\), the standard frame and \(((3,4)/5,(-4,3)/5)\) give fourth-power sums \(1\) and \(337/625\), while both squared-coordinate sums are \(1\) (\texttt{H013Contract}). This counterexample is not Gleason uniqueness {[}REF-QM-017{]} & Simplex valuation & Spectral measure \(P_i\) on continuum lattice & Gleason's theorem applies to normalized weight-one frame functions on the complete projection structure {[}REF-QM-017{]}. These finite real-ray reports do not determine a complex Hermitian operator or certify its positivity. \\
\textbf{P7: Algebraic Non-Factorable Witness} & \textbf{DERIVED (algebraic)} & Vector \(\Phi = (1, 0, 0, 1)^T\) on \(Q \times Q\) has nonfactorable coefficient matrix (\(\det = 1 \ne 0\)), proving tensor rank \(\ge 2\) in \(\mathbb{Z}[i]^{Q \times Q}\) (\texttt{H014\_bell\_not\_product}) & Factored distributions & Entangled bipartite states \(\sum c_{ij} \vert i j \rangle\) & Algebraic non-product vectors exist in free modules over Cartesian products; physical quantum entanglement requires composite state tensor structure and LOCC constraints, which remain unpaid (§19.2). \\
\textbf{P8: Tensor Uniqueness} & \textbf{REFUTED (forced); IMPORTED} & Weaker combining interface admits \texttt{dumpRight} \(\ne\) \texttt{tensorComb} (\texttt{H014Interface}) & Cartesian product \(X \times Y\) & Tensor product \(\mathcal{H}_A \otimes \mathcal{H}_B\) & Injective composite tagging excludes degenerate maps like \texttt{dumpRight}, but does not uniquely force the Hilbert tensor product or its categorical factorisation property. \\
\textbf{P9: Non-Contextuality (KS)} & \textbf{REFUTED (classical); OBSTRUCTED} & Cabello 18-ray in \(\{-1, 0, 1\}^4\) with 9 tetrads yields mod-2 parity contradiction \(9 \not\equiv 2k\) (CH18-P06; {[}REF-QM-014{]}; \texttt{H020} checks the supplied rays, projective distinctness, all 54 orthogonality products and exact two-incidence before parity) & Non-contextual \(\{0, 1\}\) valuations & Kochen-Specker state space & Refutes non-contextual binary valuations on this ray configuration. No classification of its fractional frame-polytope vertices is supplied. \\
\textbf{P10: Reversible Dilation} & \textbf{DERIVED (census)} & Census on 2-point carrier with 1-bit prepared ancilla (\(d_A=2, d_E=2\)): exactly \(2\cdot 2=4\) reversible permutations dilate collapse: the two prepared states must go to the two states with system value \(0\), and the other two to the rest (\texttt{H015CollapseDilations} enumerates every permutation of the supplied composite carrier) & Markov embedding & Stinespring dilation (\(d_E \le d_A d_B\)) {[}REF-QM-016{]} & Collapse has no same-basis isometry (\texttt{H015\_collapse\_no\_isometry} checks its Gram matrix); while finite channels admit finite ancillary dilations (\(r \le d_A d_B\) {[}REF-QM-016{]}), quantum dilation requires importing CPTP maps, complex Hilbert spaces, and unitary rotations. \\
\textbf{P11: Continuous Dynamics} & \textbf{REFUTED (host); REFUSED} & Toggle permutation on \(\mathrm{Bool}\) has no report square root (\texttt{import7\_unforced} in \texttt{QuantumImports.lean}) & Continuous Markov flow & Stone's theorem \(e^{-i H t}\) {[}REF-QM-019{]} & In topology, \(\mathrm{SO}(n, \mathbb{Q})\) is dense in \(\mathrm{SO}(n, \mathbb{R})\), not discrete; operational discreteness consists in restricting evolution to discrete integer time steps \(\tau \in \mathbb{Z}\) without Lie infinitesimal generators. The missing half-step is missing only among report permutations: the unitary \(M\) of CH18-P05 over \(\mathbb{Z}[i][\tfrac12]\) squares to the toggle. \\
\textbf{P12: Finite Matrix Recovery} & \textbf{DERIVED (\(\mathbb{Z}[i]\))} & \(\operatorname{End}_{\mathbb Z[i]}(\mathbb Z[i]^Q) \cong M_n(\mathbb{Z}[i])\) with conjugate-transpose adjoint (standard linear-algebra theorem; \texttt{H017Fin} checks supplied finite basis images, matrix, adjoint and composition) & Stochastic matrices & Matrix algebra \(M_n(\mathbb{C})\) & The code checks declared Gaussian-integer linear instances; it does not establish the unrestricted isomorphism by sampling or recover continuous \(M_n(\mathbb{C})\). \\
\end{longtable}

\subsubsection{18.4 Contextuality in exact arithmetic}\label{contextuality-in-exact-arithmetic}

The Mermin--Peres square {[}REF-QM-015{]}, the Peres 33-ray configuration in \(\mathbb{Z}[\zeta_8]^3\),
and the Cabello 18-ray configuration in \(\{-1, 0, 1\}^4\) {[}REF-QM-014{]} establish contextuality
obstructions in exact arithmetic without continuous angles. The eighteen-ray case is a
finite parity count on the following table of nine contexts.

\begin{longtable}[]{@{}ll@{}}
\toprule\noalign{}
Context & Vectors in \(\mathbb Q^4\) \\
\midrule\noalign{}
\endhead
\bottomrule\noalign{}
\endlastfoot
\(C_1\) & \((0,0,0,1)\), \((0,0,1,0)\), \((1,1,0,0)\), \((1,-1,0,0)\) \\
\(C_2\) & \((0,0,0,1)\), \((0,1,0,0)\), \((1,0,1,0)\), \((1,0,-1,0)\) \\
\(C_3\) & \((1,-1,1,-1)\), \((1,-1,-1,1)\), \((1,1,0,0)\), \((0,0,1,1)\) \\
\(C_4\) & \((1,-1,1,-1)\), \((1,1,1,1)\), \((1,0,-1,0)\), \((0,1,0,-1)\) \\
\(C_5\) & \((0,0,1,0)\), \((0,1,0,0)\), \((1,0,0,1)\), \((1,0,0,-1)\) \\
\(C_6\) & \((1,-1,-1,1)\), \((1,1,1,1)\), \((1,0,0,-1)\), \((0,1,-1,0)\) \\
\(C_7\) & \((1,1,-1,1)\), \((1,1,1,-1)\), \((1,-1,0,0)\), \((0,0,1,1)\) \\
\(C_8\) & \((1,1,-1,1)\), \((1,-1,-1,-1)\), \((1,0,1,0)\), \((0,1,0,-1)\) \\
\(C_9\) & \((1,1,1,-1)\), \((1,-1,-1,-1)\), \((1,0,0,1)\), \((0,1,-1,0)\) \\
\end{longtable}

\texttt{WDP-CH18-P06} \textbf{Finite proposition (an eighteen-ray parity obstruction).} In the
table, each context consists of four pairwise orthogonal vectors of \(\mathbb Q^4\),
no two of the eighteen vectors are parallel, and every vector lies in exactly two
contexts. No map \(v\) from the eighteen vectors to \(\mathbb F_2\) satisfies
\(\sum_{r\in C_k}v(r)=1\) for \(k=1,\dots,9\); in particular no \(\{0,1\}\)-valuation
puts exactly one \(1\) in every context.

\emph{Proof.} The \(54\) inner products within contexts are \(0\) by direct integer
computation, no two vectors are proportional, and each vector occurs in exactly
two rows. Adding the nine equations gives
\(\sum_k\sum_{r\in C_k}v(r)=\sum_r2v(r)=0\) in \(\mathbb F_2\) on the left and
\(9=1\) on the right. A solution with values in \(\{0,1\}\) or in \(\mathbb Z\) would
reduce modulo \(2\) to one over \(\mathbb F_2\). \(\square\)

The proposition concerns a declared finite configuration. It becomes the
Kochen--Specker obstruction for \(\mathbb C^4\) {[}REF-QM-008; REF-QM-014{]} only after
the vectors are read as rank-one projectors of the Hilbert carrier of §5.8 and a
noncontextual valuation is defined on them, and it yields no Bell violation, no
correlation bound, and no reconstruction (CH18-N01, CH19-N02). CH19-P05 applies
the same count to the four CHSH contexts. The absorbed archived working-paper source 004
printed a different eighteen-vector table in which seven of the nine contexts
contain non-orthogonal pairs, and read the contradiction as a nonzero class in
the first cohomology of the context hypergraph with \(\mathbb Z_2\) coefficients.
That reading is not carried: no cochain complex was declared, and the falsifier
of §21.6 applies. The source cited an external Lean check of a correct table,
equivalent to this one under a signed permutation of coordinates; it was read,
not built or ported, and gives no Lean credit.

The finite frame equations constrain ray probabilities. Such a probability
vector is not a matrix and has no eigenvalues without a specified reconstruction.
The listed rays are real, so their expectations do not determine a complex
Hermitian matrix. For example,
\[
H_t=\tfrac14 I_4+t\,i(E_{12}-E_{21})
\]
has expectation \(1/4\) on every real unit ray. Its eigenvalues are
\(1/4+t,1/4-t,1/4,1/4\); for \(t>1/4\) it is not positive semidefinite, yet
its ray report is also the report of \(I_4/4\). Thus these finite tests do not
certify positivity of a submitted operator. This example does not prove that
the frame polytope has a nonquantum vertex. The image of the density-matrix
set under the ray-readout map is a spectrahedral shadow; neither equality
nor strict inclusion in this particular frame polytope is established here.

Gleason's theorem {[}REF-QM-017{]}, on a real or complex separable Hilbert space
of dimension at least three, concerns the complete projection structure:
a nonnegative, normalized weight-one frame function is represented by a
unique density operator. A finite list of rays does not supply those hypotheses.

\texttt{WDP-CH18-N01} \textbf{Nonclaim.} The primitive finite calculus alone is not shown
to select the quantum set of correlations, a complex Hilbert representation,
or quantum dynamics. Under separately supplied finite Hilbert-space, tensor
and trace assumptions, CH19-O02 establishes the CHSH bound \(S^2\le8\). That
conditional result does not derive those assumptions from the primitive
calculus.

\subsubsection{18.5 The eight external interfaces exposed by this reconstruction route}\label{the-eight-external-interfaces-exposed-by-this-reconstruction-route}

To reach textbook continuous quantum mechanics \(M_n(\mathbb{C})\), this reconstruction
route exposes eight external mathematical interfaces along its boundary. We describe
them as interfaces exposed by this specific constructive progression, rather than
asserting that all eight are logically independent necessities for every conceivable
axiomatic route:

\begin{enumerate}
\def\labelenumi{\arabic{enumi}.}
\item
  \textbf{Import 1: The Continuous Field \(\mathbb{C}\).} Cauchy completeness, archimedean
  topology, and algebraic closure (refused on discrete host; host uses \(\mathbb{Z}[i]\)).
\item
  \textbf{Import 2: The Continuous Circle Group \(U(1)\).} Connected Lie group phase
  symmetries (refused on host; discrete interference selects only \(\{\pm 1\}\)).
\item
  \textbf{Import 3: The Spectral Theorem for Self-Adjoint Operators.} Orthonormal
  eigenbases (refused on host; Hermitian matrices in \(M_n(\mathbb{Z}[i])\) need not have eigenvalues in \(\mathbb{Z}[i]\)).
\item
  \textbf{Import 4: The complete projection structure in dimension at least three.}
  Gleason's theorem represents normalized weight-one frame functions by unique
  density operators under its full real or complex Hilbert-space hypotheses.
  The finite real-ray example above does not provide those hypotheses or a
  classification of all probability vectors satisfying its frame equations.
\item
  \textbf{Import 5: Unique Hilbert Tensor Product.} Postulating the Hilbert tensor
  product \(\mathcal{H}_A \otimes \mathcal{H}_B\) and its initial bilinear factorisation property;
  injective composite report tagging excludes degenerate combining maps such as
  \texttt{dumpRight}, but does not uniquely force the Hilbert tensor product structure.
\item
  \textbf{Import 6: CPTP Channels \& Unitary Stinespring Dilation.} Completely positive
  trace-preserving (CPTP) channel structure and complex unitary rotations on an enlarged
  system (finite-dimensional quantum channels admit finite ancillary constructions with
  environment dimension \(r \le d_A d_B\) {[}REF-QM-016{]}; the missing structures are CPTP condition,
  complex Hilbert inner products, and continuous unitary groups, whereas the discrete carrier
  provides only a finite permutation census).
\item
  \textbf{Import 7: Continuous 1-Parameter Unitary Groups / Stone's Theorem.} Continuous
  time evolution \(t \in \mathbb{R}\) generated by self-adjoint Hamiltonians \(e^{-i H t}\)
  {[}REF-QM-019{]} (refused on discrete host; \(\mathrm{SO}(n, \mathbb{Q})\) is topologically
  dense in \(\mathrm{SO}(n, \mathbb{R})\), but discrete dynamics is restricted to integer
  steps \(\tau \in \mathbb{Z}\)).
\item
  \textbf{Import 8: The positive-semidefinite state cone.} Require state matrices
  to satisfy \(\rho\succeq0\). A fixed finite collection of ray inequalities
  gives polyhedral constraints and, in general, does not characterize this
  entire cone. For a supplied exact Hermitian matrix, positivity can instead
  be checked by valid principal-minor, Gram-factor or congruence certificates;
  semidefinite programming is another method for problems with matrix
  positivity constraints. These methods check the supplied state condition;
  they do not derive its operational interpretation.
\end{enumerate}

The route also has steps that are not among these interfaces. CH18-P05 shows
that a finite localization of the scalars, such as \(\mathbb{Z}[\tfrac15]\) or
\(\mathbb{Z}[i][\tfrac12]\), already admits unitaries that separate readout fibers.

\subsubsection{18.6 Audit verdict}\label{audit-verdict}

The defensible verdict is: the finite obstruction calculus, under the stated
specialization hypotheses, establishes an
integer permutation-isometry envelope over \(\mathbb{Z}\), a phase-blind
superposition fiber, an interference non-descent obstruction over \(\mathbb{Z}\)
once the similitude \(H\) is admitted, an algebraic non-factorable tensor-rank
witness, and the finite matrix algebra
\(M_n(\mathbb{Z}[i])\). For the continuous circle group, Hilbert tensor
uniqueness, Gleason uniqueness at dimension \(\ge 3\), CPTP dilations, Lie
algebras, and continuous unitary dynamics it gives one small finite witness each
of why the discrete host does not supply them; the reconstruction route in
§18.5 exposes eight admitted interfaces. This is not a theorem that every
reconstruction must use those interfaces, or that they are independent. Under the two strengthened readings of
CH18-P02 and CH18-P04 the package is classical, which is outcome (a) of CH18-O01
for those readings; CH18-P07 and CH18-P08 add classical facts for the package S,
which has stochastic updates, classical instruments and joint runs.

This verdict can be stated with exact claim discipline across the four-tier foundational ambition taxonomy:

\begin{enumerate}
\def\labelenumi{\arabic{enumi}.}
\tightlist
\item
  \textbf{Tier 1 (Syntactic Representation):} Established for the declared finite encodings. Parse/print and token encode/decode round-trip laws concern their specified carriers; selected implementation replays have the separate scopes recorded in Appendix D.
\item
  \textbf{Tier 2 (Faithful Operational Representation):} Established for the stated finite models and adapter laws. Integer permutation actions and cochain transports preserve the specified operations and questions when the corresponding commutation conditions hold. No unrestricted theory-profile adapter follows from an encoding alone (§20.4).
\item
  \textbf{Tier 3 (Reconstruction Under Explicit Principles):} Characterized with open interfaces. This reconstruction route exposes eight admitted interfaces; under the strengthened readings of CH18-P02 and CH18-P04 the package stays inside the monomial \(\mathbb{Z}[i]\) envelope, which is classical. The similitude \(H\) (P3) and a finite localization of the scalars (CH18-P05) each leave that envelope, and neither is one of the eight interfaces.
\item
  \textbf{Tier 4 (Unique Selection From Primitive Kernel):} Explicitly refused. The monograph does not claim that textbook continuous quantum mechanics \(M_n(\mathbb{C})\) is forced or uniquely selected by primitive difference, and it establishes no theorem that it cannot be (§18.5, CH19-N01).
\end{enumerate}

\begin{center}\rule{0.5\linewidth}{0.5pt}\end{center}

\subsection{Chapter 19 --- The unpaid quantum interfaces}\label{chapter-19-the-unpaid-quantum-interfaces}

\subsubsection{19.1 Context and empirical content}\label{context-and-empirical-content}

A family of contexts is not selected by the existence of cochains. To reach an
empirical contextuality model one needs measurements, outcome sets, compatible
probability distributions, and a cover. Chapter 13 provides a general finite
measurement packet, but a quantum scenario requires a specialized adapter into
effects and instruments.

\texttt{WDP-CH19-O01} \textbf{Conditional result (arbitrary-cover repair and sampling).} For any declared finite cover and nonempty rational run tables, the minimax context \(L^1\) distance to the compatible polytope is the optimum of the exact standard-form LP in Appendix J, with independently verified primal/dual certificates. Under declared within-context iid categorical trials and compatible population marginals, the trial counts give explicit finite-sample pair tolerances and a union-bound failure probability. The empirical image retains only the previously stated outcome information.

CH19-P02 to CH19-P05 below construct this map for tables whose frequencies
agree exactly on every overlap of contexts, and state what survives it.
CH19-P06 and CH19-P07 add a tolerant adapter for tables whose context
frequencies agree only up to a declared tolerance. Appendix J completes arbitrary-context exact repair and derives finite-sample tolerance under the explicitly declared trial protocol and probability model of §13.6. The specialization to effects and
instruments is CH19-O04.

The adapter uses the packet and a declared profile. A declared cover consists
of a finite set \(X\) of measurements, a finite outcome set \(O\), and contexts
\(C_i \subseteq X\). A profile assigns each measurement \(x\) a measurand and a
partition of calibrated results into outcomes. The partition stands in for the
structure \(\mathcal S\) of CH13-D01; it is declared per measurement rather than
carried by the packet. A packet may fill slot \(x\) only when its measurand is
the one declared for \(x\), and it reports the bin of its calibrated value. A run
in context \(C_i\) holds one such packet for each measurement of \(C_i\), and a
table lists the runs recorded in each context. The frequency \(e_i(s)\) of a
section \(s : C_i \to O\) is the normalized multiplicity (Chapter 8) of the runs
that report \(s\). Weights are rational throughout. An empirical model
{[}REF-QM-001{]} is a probability distribution with rational weights on the sections
of each context whose marginals agree on every overlap \(C_i \cap C_j\). It is
noncontextual when a nonnegative rational distribution on the global
assignments \(X \to O\) restricts to it, and strongly contextual when no global
assignment restricts, in every context, to a section of positive weight. For a
model with rational weights a real nonnegative global section yields a rational
one, because the constraints are rational linear inequalities; that step is not
mechanized.

\texttt{WDP-CH19-P02} \textbf{Finite proposition (the packet adapter).} On a declared
cover, the adapter sends a table of measurement packets to its frequencies. It
refuses the table when some context has no runs, naming that context, and when
the frequencies of two contexts disagree on their overlap, naming the two
contexts and the overlap section at which they differ. Otherwise it returns an
empirical model. It returns the model \(m\) exactly when no context is empty and
\(m\) is the table of frequencies.

\emph{Proof.} Frequencies are nonnegative, and in a context with at least one run
they sum to one. Agreement on overlaps is the compatibility condition of an
empirical model. \(\square\)

\texttt{WDP-CH19-P03} \textbf{Finite proposition (what survives the adapter).} A section
has positive weight in the image exactly when a recorded run reports it. The
image depends only on the list of reported outcomes in each context, and two
nonempty tables have the same frequencies in a context exactly when their
outcome counts there are proportional. So a packet field reaches the image only
through the outcome it reports. The world, the presentation map, the indication
and the calibrated value reach it that way. The trace, the evaluator, the
calibration chain and its uncertainty contributions, the unit and scale, the
provenance record, the order of runs, and the absolute counts do not affect it.
Replacing the calibration chain of every packet, for example, changes the
documented traceability and leaves the image unchanged. Two packets that differ
only in their evaluators' verdicts on reports that were not recorded have the
same measurand and readout, so exchanging them leaves the image unchanged.
Every packet of a run does carry the declared measurand, an accepted evaluator
status, a report that the observer's memory returns, and a provenance record
that replays the indication and the calibrated value. It carries them by its
type (CH13-D01), in every table, and none of them can be read back from the
image. The answer to the second half of CH19-O01 is therefore that no observer
or evaluator field survives the map except through the outcome it reports.
Acceptance survives only as the condition for a run to be counted.

\emph{Proof.} A count is positive exactly when some run reports the section. Counts
are functions of the outcome list, and \(c/n = c'/n'\) exactly when \(c n' = c' n\).
The readout reads the report, the indication, the calibrated value, and the
profile's bin, and nothing else. \(\square\)

\texttt{WDP-CH19-P04} \textbf{Finite proposition (the image of freely formed tables is every
rational model).} Suppose every outcome of every measurement is reported by
some admissible packet, context runs may be formed from arbitrary tuples of
those packets, and each run may be repeated any finite number of times. Then
every empirical model on the cover with rational weights is the image of a
finite table. Under this table-formation rule, the packet type alone neither
forces nor excludes contextuality.

\emph{Proof.} For each section \(s\) of context \(i\), choose a packet reporting each
component of \(s\) and form their tuple. Let \(N\) be a common denominator of the
weights in that context, and record \(N e_i(s)\) copies of the tuple. Their
frequencies are exactly \(e_i\). \(\square\)

This is a formal table-realization theorem. Individual outcome availability
does not imply joint physical availability: an operational theory may forbid
some tuples or constrain which runs can coexist. Such restrictions must be
declared separately; the construction does not realize a PR box as a physical
experiment.

\texttt{WDP-CH19-P05} \textbf{Finite proposition (contextuality as a missing global
table).} Call a table global when each of its runs reads every measurement at
once and every context reads the same runs. Assume the freely formed,
repeatable tables of CH19-P04, with the same tuple-formation rule for all
measurements at once. An empirical model with rational weights on any declared cover is noncontextual exactly when some
global table has it as its image. On the CHSH cover, with Alice's settings \(a_0, a_1\), Bob's settings
\(b_0, b_1\), contexts \(\{a_x, b_y\}\) and outcomes in \(\{0, 1\}\), let \(E_{xy}\) be
the expected value of \(+1\) for agreeing outcomes and \(-1\) for differing ones.
Every noncontextual model, and so every global table, has
\(E_{00} + E_{01} + E_{10} - E_{11} \le 2\). The PR box {[}REF-QM-009{]} gives
probability \(\tfrac12\) to each outcome pair with \(a_x \oplus b_y = xy\). It is an
empirical model, its CHSH value is \(4\), and it is strongly contextual. It is the
image of a finite table of accepted measurement packets, and no global table
has it as its image. It has the signed global section
\(\tfrac12(\delta_{0000} - \delta_{0010} + \delta_{0110} + \delta_{1011})\) on
\((a_0, a_1, b_0, b_1)\), and no nonnegative one.

\emph{Proof.} A nonnegative global section with rational weights is the frequency
table of a finite list of global assignments, by freely forming the global
tuples and clearing denominators as in CH19-P04. If a rational empirical model
has a real nonnegative global section, the finite rational linear system
specifying that section has a rational feasible point: restrict to the support
of a real solution, use rational Gaussian elimination and approximate its free
parameters closely enough to keep all positive coordinates positive. Thus the
rational-weight restriction loses no noncontextual rational models.
Conversely, the frequencies of the global assignments in a global
table are a nonnegative global section. One global assignment has CHSH value at
most \(2\), and so does any mixture of them. In the PR box the four context
conditions add up to \(1\) modulo \(2\), while each measurement occurs in two
contexts, so no global assignment meets all four. This is the mod-2 count of
CH18-P06 (§18.4), applied to the four CHSH contexts instead of the nine contexts
there.
The signed section is checked by direct computation. \(\square\)

So, under the stated free table-formation rule, a rational model is contextual exactly when no
single table of global runs reproduces it. The packet type alone imposes no Tsirelson bound, since the PR box is
an image, and CH18-N01 stands.

\texttt{WDP-CH19-P06} \textbf{Finite proposition (the tolerant adapter).} For contexts \(i,j\)
let \(D_{ij}\) be the \(\ell^1\) distance between the marginals of the frequency
tables \(e_i,e_j\) on the overlap \(C_i\cap C_j\). For \(\varepsilon\ge0\) the
tolerant adapter accepts a table exactly when no context is empty and every
\(D_{ij}\le\varepsilon\). It returns the frequencies together with a certificate
that every \(D_{ij}\le\varepsilon\), and otherwise names the empty context or a pair with \(D_{ij}>\varepsilon\). At
\(\varepsilon=0\) it accepts exactly the tables of CH19-P02 and returns the same
frequencies, and acceptance is monotone in \(\varepsilon\). If every \(e_i\) lies
within \(\varepsilon\) in \(\ell^1\) of the \(i\)-th distribution of a compatible
model, then every \(D_{ij}\le2\varepsilon\).

\emph{Proof.} At \(\varepsilon=0\) the condition is agreement of the overlap
marginals, which is compatibility. Marginalization does not increase \(\ell^1\)
distance and the compatible model has equal overlap marginals, so the triangle
inequality gives \(D_{ij}\le\varepsilon+\varepsilon\). \(\square\)

What the tolerant adapter returns is not an empirical model, since its overlap
marginals may differ by up to \(\varepsilon\); it is a table with a certified
signalling defect.

\texttt{WDP-CH19-P07} \textbf{Finite proposition (exact repair on two contexts).} On a
cover with two contexts and a nonempty outcome set, let \(D\) be the defect of the
two frequency tables. Replace each overlap marginal by the average \(\nu\) of the
two, rescale the conditional weights over each overlap section, and, where a
context has no weight over a section with \(\nu>0\), place \(\nu\) on one fixed
extension of that section. The result is an empirical model, and it lies at
\(\ell^1\) distance exactly \(D/2\) from each table (the distance is taken context
by context). Hence a compatible model within \(\varepsilon\) of both tables exists
exactly when \(D\le2\varepsilon\).

\emph{Proof.} Over an overlap section \(t\) with positive weight \(\mu_i(t)\),
rescaling moves \(\mu_i(t)\) to \(\nu(t)\) and changes the weights over \(t\) by
\(|\mu_i(t)-\nu(t)|\) in total. Where \(\mu_i(t)=0\) the added mass is
\(\nu(t)=|\mu_i(t)-\nu(t)|\). Summing over \(t\) gives \(\sum_t|\mu_i(t)-\nu(t)|=D/2\),
because \(\nu-\mu_i=\pm(\mu_1-\mu_2)/2\). The converse is CH19-P06. \(\square\)

The closeness hypothesis of CH19-P06 is a declaration. Appendix J derives a sufficient finite-sample tolerance from actual trial counts under explicit within-context iid and compatible-population assumptions, and solves the exact compatible-polytope distance for arbitrary finite covers. The three-point-mass example shows why half the pairwise defect is insufficient for three contexts. These assumptions cannot be inferred from the table (§9.9).

\subsubsection{19.2 Composition}\label{composition}

Two systems require a rule for states, effects, and transformations of the
composite. Cartesian product, tensor product, minimal and maximal ordered tensor
products, and constrained subspaces yield different theories. Primitive
difference does not choose among them.

\texttt{WDP-CH19-O02} \textbf{Conditional result (paid quantum composition).} Declaring positive finite Hilbert density states, local Hermitian dichotomic observables and trace pairing on the tensor composite gives the exact CHSH bound \(S^2\le8\), excluding PR. Commuting complete local matrix units yield the tensor-factor representation up to a multiplicity space, removed by the stated dimension condition. The native instrument and observer adapters exhibit consistency with the finite architecture. These assumptions are inputs; CP alone does not supply them.

CH19-P08 to CH19-P10 state a composition principle that needs no import and
prove what its clauses do and do not decide. The principle CP has three clauses:
(T) injective tagging, so that the composite report determines both local
reports; (NS) no-signalling, so that a local instrument on one side never
changes the marginal of the other, which on the CHSH product cover is the
compatibility of an empirical model on overlaps; and (B) bilinear positivity, so
that a composite state is a bilinear form on product effects that is nonnegative
and normalized. Clause (B) is formulated here for classical carriers only.

\texttt{WDP-CH19-P08} \textbf{Finite proposition (injective tagging).} For nonempty finite
carriers \(A,B\) and a composite tagging \(c:A\times B\to C\), there are decoders
\(d_A,d_B\) with \(d_A(c(a,b))=a\) and \(d_B(c(a,b))=b\) if and only if \(c\) is
injective. The tagging \(\mathrm{dumpRight}(a,b)=b\) is not injective when \(A\)
has two elements, so it has no decoder for the left factor.

\emph{Proof.} If decoders exist, \(c(a,b)=c(a',b')\) gives \(a=a'\) and \(b=b'\). If \(c\) is
injective, decode with any left inverse. \(\square\)

\texttt{WDP-CH19-P09} \textbf{Finite proposition (no-signalling and presentation maps).} On the
Cartesian composite \(A\times B\) a joint state is a nonnegative weight \(p(a,b)\),
and a classical instrument \((K_o)\) on \(B\) acts as \(1\otimes K_o\). For every \(a\),
\(\sum_o\sum_{b'}\big((1\otimes K_o)p\big)(a,b')=\sum_bp(a,b)\), and the same holds
with the factors exchanged. So the marginal of one factor is the same for any two
instruments on the other. A presentation map \(\pi:A\to Q\) lifts to
\(\pi\circ\mathrm{fst}\) on \(A\times B\), and pushing forward the \(A\)-marginal gives
the \(Q\)-marginal of the lifted map.

\emph{Proof.} Summing over the outcomes and the post-states uses only that
\(\sum_oK_o\) is column-stochastic. The lift is a reordering of a finite sum.
\(\square\)

\texttt{WDP-CH19-P10} \textbf{Finite proposition (classical composites are Cartesian).} Let
\(A,B\) be finite carriers and \(\omega\) a bilinear form on effect vectors of \(A\)
and \(B\) that is nonnegative on every pair of effects (values in \([0,1]\)) and is
\(1\) on the pair of constant effects. Then there is a unique joint distribution
\(p\) on \(A\times B\) with \(\omega(e,f)=\sum_{a,b}e_af_bp(a,b)\), namely
\(p(a,b)=\omega(1_a,1_b)\); conversely every joint distribution defines such a
form. On the CHSH product cover, whose contexts \(\{a_x,b_y\}\) are the products of
two two-setting scenarios, the PR box of CH19-P05 is compatible on every overlap,
so it satisfies (NS), and yet it is not noncontextual and has \(\mathrm{CHSH}=4\).

\emph{Proof.} Evaluating \(\omega\) on indicator effects gives \(p\) and its
nonnegativity, and bilinearity gives the formula. The second part is CH19-P05.
\(\square\)

By its clauses CP excludes taggings that lose a local report and composites
that signal. For classical carriers (B) forces every composite state to be a
joint distribution on \(A\times B\); it does not force every joint distribution
to be a state. Clauses (T) and (NS) do not exclude the PR box, and (B) is not
formulated here for its nonclassical factors, so its status under CP is
undecided. On the Cartesian composite of classical carriers, local
presentations and instruments act on one factor, do not signal, and the
composite report determines both local reports (CH19-P08, CH19-P09); this is the
sense in which CP is consistent with the finite observer architecture here, and
no more is proved. Nothing here shows that CP selects the Hilbert tensor
product. Appendix J adds a sufficient finite Hilbert composition package that excludes PR and conditionally yields tensor factors. Its matrix algebra, positivity and trace assumptions pay import 5 of §18.5 explicitly; they are not consequences of CP.

\subsubsection{19.3 Continuity and reversible dynamics}\label{continuity-and-reversible-dynamics}

Finite-dimensional real or complex vector spaces can support continuous
parameter groups once their topology and transformations are specified.
Finite dimension is different from finite cardinality. Every finite Hausdorff
space is discrete, although finite spaces with other topologies exist.
Consequently, a continuous action of the connected group \(\mathbb R\) on a
finite discrete state set is trivial. Nontrivial continuous reversible dynamics
requires additional structure beyond finitely many witness types.

\subsubsection{19.4 Phase and interference}\label{phase-and-interference}

The complex semiring permits cancellation. A \(U(1)\) connection permits loop
phase. Neither supplies the operational rule connecting amplitude to observed
frequency. Conversely, positive finite probability does not supply phase. A
valid bridge must specify:

\begin{enumerate}
\def\labelenumi{\arabic{enumi}.}
\tightlist
\item
  the amplitude carrier;
\item
  the composition/addition of alternatives;
\item
  the transformation law;
\item
  the measurement/effect pairing;
\item
  the probability readout;
\item
  the empirical interpretation.
\end{enumerate}

\texttt{WDP-CH19-O03} \textbf{Conditional result (phase and operational readout bridge).} A declared continuous character \(\chi_\theta(t)=e^{i\theta t}\) maps real cochains into complex path amplitudes and descends under gauge to phase-blind operational readouts; its cycle-period lattice and anchor dependence are explicit. An independent exact rational-turn/cyclotomic branch is implemented. Declared effect/instrument pairing supplies a conditional operational model, while normalized quadratic and quartic readouts prove nonselection by phase invariance. No Born rule or faithful real phase embedding is inferred.

CH19-P11 and CH19-P12 declare the first three slots of this bridge and prove
what the last three can and cannot depend on. The amplitude carrier is \(\mathbb
Q(i)\). Alternatives add in \(\mathbb Q(i)\). The transformation law multiplies by
the unit-norm phase \(u^{a(v,w)}\) along each step \(v\to w\) of a walk, for an
antisymmetric integer edge cochain \(a\) and the fixed rational phase
\(u=(3+4i)/5\). The effect pairing and the probability readout stay parameters.

\texttt{WDP-CH19-P11} \textbf{Finite proposition (gauge descent of phase-blind readouts).}
Let \(z_1\ne0\). Two pairs of amplitudes \((z_1,z_2)\) and \((w_1,w_2)\) in \(\mathbb
Q(i)\) differ by one common unit-norm phase exactly when \(N(z_1)=N(w_1)\),
\(N(z_2)=N(w_2)\) and \(z_1\bar z_2=w_1\bar w_2\). So a readout of two alternatives
that a common phase does not change is a function of \(|z_1|^2\), \(|z_2|^2\) and
the cross term \(z_1\bar z_2\). For two walks with common endpoints and magnitudes
\(m_1,m_2\), \(z_1\bar z_2=m_1m_2\) times the ratio of the two walk amplitudes,
which for an antisymmetric cochain is the holonomy of the loop formed by the
first walk and the reverse of the second (the mechanized statement is about the
ratio and does not use antisymmetry). If the cochain
changes by a coboundary \(d_0f\), each walk amplitude is multiplied by the same
phase \(u^{f(\text{end})-f(\text{start})}\), so the holonomy and every phase-blind
readout are unchanged: they depend on the cochain only through its class modulo
coboundaries. The quadratic readout \(N(z_1+z_2)\) and the quartic readout
\(N(z_1+z_2)^2\) are both phase-blind and differ (\(4\) against \(16\) for
\(z_1=z_2=1\)), so gauge invariance does not select the quadratic readout.

\emph{Proof.} If \(w=gz\) with \(N(g)=1\), the three invariants agree because \(N\) is
multiplicative and \(g\bar g=1\). Conversely put \(g=w_1\bar z_1/N(z_1)\). Then
\(gz_1=w_1\) and \(N(g)=1\), and the equality \(z_2\bar z_1=w_2\bar w_1\), the
conjugate of the third condition, gives \(gz_2=w_2\). By induction along a walk the
amplitude of \(a+d_0f\) is that of \(a\) times \(u^{f(\text{end})-f(\text{start})}\);
the two walks share this factor, which cancels in their holonomy. \(\square\)

\texttt{WDP-CH19-P12} \textbf{Finite proposition (the exponentiation obstruction).} The
phase \(u=(3+4i)/5\) has no square root among the unit-norm elements of \(\mathbb
Q(i)\), so no group homomorphism from \((\mathbb Q,+)\) to those elements sends \(1\)
to \(u\). Integer exponents compose, \(u^{m+n}=u^mu^n\), and \(u^n=1\) only for
\(n=0\). The map \(a\mapsto u^a\) sends coboundaries to gauge transformations
(CH19-P11). So the exponentiation \(U_e=\exp(i\alpha a_e)\) of §5.4 is not
available on rational cochains with this scale inside \(\mathbb Q(i)\); it needs
an extension of the scalars that contains roots of \(u\) of every order (an
algebraic closure suffices, and import 1 of §18.5 supplies one), or a
restriction to integer cochains. The obstruction is to a homomorphism with
values in the unit-norm elements of \(\mathbb Q(i)\), not a claim that no
exponentiation exists.

\emph{Proof.} A square root \(x+yi\) would satisfy \(x^2+y^2=1\) and \(x^2-y^2=3/5\), hence
\(x^2=4/5\), and \(20=(5x)^2\) is not the square of a rational number. A homomorphism
\(\varphi\) with \(\varphi(1)=u\) would give \(\varphi(\tfrac12)^2=u\). For the order
of \(u\), write \(u=(2+i)/(2-i)\). If \(u^n=1\) with \(n>0\), then \((2+i)^n=(2-i)^n\), so
the imaginary part \(b_n\) of \((2+i)^n=a_n+b_ni\) vanishes. Modulo \(5\),
\(a_n+3b_n\equiv0\) for \(n\ge1\) and \(a_n-3b_n\equiv(-1)^n\), so \(b_n=0\) gives
\(0\equiv\pm1\). Negative \(n\) follow by inversion. \(\square\)

Appendix J supplies the real-character bridge, its period lattice and the exact rational-turn alternative. Probability readout and empirical interpretation are explicitly declared operational parameters. CH19-P11 and the quadratic/quartic countermodel establish that gauge invariance cannot select them (CH19-N01).

\subsubsection{19.5 Born uniqueness}\label{born-uniqueness}

Chapter 9 establishes only nonselection by normalization. Gleason-type theorems
{[}REF-PROB-002{]} begin with a Hilbert space and noncontextual probability assignments of a specified
kind. Operational reconstructions begin with other strong axioms. Those are
legitimate routes when their assumptions are visible. They cannot be
retroactively described as consequences of finite witness mass.

\texttt{WDP-CH19-N01} \textbf{Nonclaim.} Born-quadratic uniqueness is not proved.

\texttt{WDP-CH19-N02} \textbf{Nonclaim.} Complex phases, contextual obstructions, or
destructive cancellation separately do not characterize quantum theory.

\subsubsection{19.6 A small curved-bicomplex example}\label{a-small-curved-bicomplex-example}

Let \(C^{0,0}=C^{1,0}=C^{0,1}=C^{1,1}=\mathbb R\) and define scalar maps

\[
d_s^{0,0}(x)=ax,\qquad d_h^{0,0}(x)=bx,
\]

with the maps into \(C^{1,1}\) given by multiplication by \(c\) from \(C^{1,0}\)
and by \(d\) from \(C^{0,1}\). The mixed defect on \(x\) is

\[
\Lambda(x)=(db+ca)x
\]

for the chosen sign convention. Even in this one-dimensional toy model,
anticommutation is a nontrivial algebraic condition \(db+ca=0\). If it fails,
calling the square a double complex is false.

One can add a homotopy \(h\) only after specifying its degree and an identity such
as \(\Lambda=d h+h d\) in a surrounding complex. The existence of the scalar
\(db+ca\) does not manufacture that homotopy. This example is deliberately
simple: the typing obstruction is already visible before geometry or physics
enters.

\subsubsection{19.7 Coarse-graining histories}\label{coarse-graining-histories}

Let \(\pi:\Omega\to\bar\Omega\) group fine histories into coarse classes. A
natural coarse weight is

\[
\bar w(\bar p)=\bigoplus_{p:\pi(p)=\bar p}w(p).
\]

This definition preserves the total finite aggregate by associativity and
commutativity of \(\oplus\) when the fibers are finite. It may fail to preserve
composition: two fine histories in the same coarse class can compose with a
third history into different coarse classes. For a compositional coarse model,
\(\pi\) must be a congruence relative to the path product, or the coarse carrier
must retain additional memory.

\texttt{WDP-CH17-P03} \textbf{Finite proposition (aggregate-preserving pushforward).} For a
finite history family and a commutative additive operation,

\[
\bigoplus_{\bar p\in\bar\Omega}\bar w(\bar p)
=\bigoplus_{p\in\Omega}w(p).
\]

The theorem is a finite reindexing. It does not prove Markovianity,
decoherence, or consistent histories.

\subsubsection{19.8 Destructive cancellation versus operational interference}\label{destructive-cancellation-versus-operational-interference}

Suppose two complex-weighted paths contribute \(z_1\) and \(z_2\). Their aggregate
amplitude is \(z_1+z_2\), and a separately declared quadratic readout gives

\[
|z_1+z_2|^2
=|z_1|^2+|z_2|^2+2\operatorname{Re}(z_1\overline{z_2}).
\]

The cross term is absent if one first converts each path to a positive weight
and then adds. This elementary identity explains why positive witness mass and
complex amplitude are different layers. The identity itself does not say which
alternatives are physically coherent, which phases occur, or why the quadratic
readout is a probability.

The history chapter should therefore use ``cancellation'' for an algebraic sum
that vanishes and reserve ``interference'' for a model containing alternatives,
phase-sensitive addition, and an operational readout. This vocabulary repair
prevents the zerosumfree theorem from masquerading as a quantum theorem.

The higher-order version of the same bookkeeping is an identity about declared
readouts. For alternatives with amplitudes \(a_j\) and a subset \(S\), let
\(A(S)=\sum_{j\in S}a_j\), and define the alternating sum
\(I_k(S)=\sum_{T\subseteq S}(-1)^{|S|-|T|}P(T)\) for \(|S|=k\).

\texttt{WDP-CH19-P01} \textbf{Finite proposition (Sorkin sums under a quadratic readout).} If
the amplitudes of finitely many alternatives add and the readout is
\(P(S)=B(A(S),A(S))\) for a declared biadditive \(B\), then the alternating sum
\(I_k\) vanishes for every set of \(k\ge3\) alternatives, while \(I_2\) equals
\(B(a,b)+B(b,a)\) and can be nonzero. Vanishing higher-order interference is
therefore a consequence of the declared quadratic readout and not a derivation
of it.

\emph{Proof.} Expanding by biadditivity, \(P(T)=\sum_{i,j\in T}B(a_i,a_j)\). The
coefficient of \(B(a_i,a_j)\) in \(I_k(S)\) is
\(\sum_{\{i,j\}\subseteq T\subseteq S}(-1)^{|S|-|T|}\), which is zero whenever
\(S\setminus\{i,j\}\) is nonempty, that is, whenever \(k\ge3\). For \(k=2\) only
the cross terms survive. \(\square\)

With \(a=1\), \(b=\tfrac35+\tfrac45i\), and \(B(x,y)=x\overline y\), the readouts are
\(P(a)=P(b)=1\) and \(P(a+b)=\tfrac{16}5\), so \(I_2=\tfrac65\); the value exceeds \(1\)
because the readout is unnormalized intensity, not probability. A quartic
readout \(|A|^4\) gives \(I_3=36\) when all three amplitudes equal \(1\), and an additive readout
gives \(I_2=0\). Commutativity of the coefficients is not used: the identity holds
for matrix amplitudes with the conjugate-transpose readout as well. The
condition \(I_3=0\) does not distinguish quantum from classical models, since the
additive readout satisfies it too {[}REF-QM-018{]}.

\subsubsection{19.9 Reconstruction audit matrix}\label{reconstruction-audit-matrix}

The following matrix states the present programme boundary. \texttt{Declared} means the
structure can be added consistently in an optional host; it does not mean the
root selects it.

\begin{longtable}[]{@{}
  >{\raggedright\arraybackslash}p{(\linewidth - 6\tabcolsep) * \real{0.2500}}
  >{\raggedright\arraybackslash}p{(\linewidth - 6\tabcolsep) * \real{0.2500}}
  >{\raggedright\arraybackslash}p{(\linewidth - 6\tabcolsep) * \real{0.2500}}
  >{\raggedright\arraybackslash}p{(\linewidth - 6\tabcolsep) * \real{0.2500}}@{}}
\toprule\noalign{}
\begin{minipage}[b]{\linewidth}\raggedright
Feature
\end{minipage} & \begin{minipage}[b]{\linewidth}\raggedright
Root finite calculus
\end{minipage} & \begin{minipage}[b]{\linewidth}\raggedright
Optional carrier
\end{minipage} & \begin{minipage}[b]{\linewidth}\raggedright
Reconstruction obligation
\end{minipage} \\
\midrule\noalign{}
\endhead
\bottomrule\noalign{}
\endlastfoot
finite state/evidence sets & available & available & specify operational interpretation \\
convex mixtures & not intrinsic & declarable & justify mixing and normalization \\
non-simplex state space & absent & declarable & prove from operational principles \\
effects/instruments & absent & declarable & connect to measurement packets (one finite qubit adapter: CH19-P13, CH19-P14) \\
complex Hilbert space & absent & declared in Ch. 5 & show why real/classical alternatives fail \\
\(U(1)\) phase & absent & declared & connect phase to operational alternatives (integer cochains and rational phases: CH19-P11, CH19-P12) \\
tensor composition & absent & not fixed & state product rule and locality axiom (CP of §19.2 forces classical composite states to be joint distributions) \\
contextual empirical model & absent & possible & supply cover, outcomes, distributions \\
continuous reversible dynamics & absent & possible & supply topology/group and derive relevance \\
purification & absent & possible & state exact axiom and countermodels \\
Born rule & absent & not selected & prove under named assumptions \\
CHSH Tsirelson bound & not derived from the primitive calculus & \(S^2\le8\) under CH19-O02's supplied finite Hilbert, local-observable, tensor and trace assumptions & deriving or selecting those assumptions remains a reconstruction obligation \\
\end{longtable}

This table makes a classical no-go verdict meaningful. If every structure
actually constructed before the optional branch has a faithful simplex model,
then no quantum conclusion follows at that stage. The verdict is not that
quantum theory is false; it is that the current premises underdetermine it.

Under the four-tier ambition taxonomy, this audit matrix establishes that:

\begin{itemize}
\tightlist
\item
  \textbf{Tiers 1 \& 2 (Syntactic \& Faithful Operational Representation)} are satisfied for discrete carrier models (such as permutation isometries and integer cochain lattices over \(\mathbb{Z}\)).
\item
  \textbf{Tier 3 (Reconstruction Under Explicit Principles)} is the exact function of the table above: cataloging which physical structures remain unforced and must be adjoined as explicit operational axioms.
\item
  \textbf{Tier 4 (Unique Selection From Primitive Kernel)} is refused as an ambition for quantum theory: the monograph does not claim that primitive difference forces complex Hilbert spaces, Gleason frame uniqueness, or the Born rule, and the finite import witnesses of Appendix D are not proofs that no derivation could (CH19-N01).
\end{itemize}

\subsubsection{19.10 A classical comparison model}\label{a-classical-comparison-model}

Take a finite set \(\Lambda\) of hidden states. Let normalized states be
probability distributions \(p\in\Delta(\Lambda)\), effects be functions
\(e:\Lambda\to[0,1]\), and transformations be stochastic matrices. Measurement
probability is

\[
\Pr(e\mid p)=\sum_{\lambda\in\Lambda}e(\lambda)p(\lambda).
\]

This model supports finite probability, observers with coarse presentation
maps, memory, stochastic instruments, and measurement verification records. Its state space
is a simplex. Any axiom package satisfied by this model cannot by itself
distinguish quantum theory from classical theory.

Thus a proposed quantum principle should be tested on this model first. If it
passes, either the principle is too weak or it contributes only to a shared
operational foundation. That shared contribution may still be valuable, but it
should not be labelled a quantum discriminator.

No-signalling is one such shared principle. With the Cartesian composite outcome
set and product effects, summing any local partition of unity on one side
returns the other side's marginal statistics. Because the PR box of §19.11 is
also nonsignalling, no-signalling alone is not a discriminator either.

\subsubsection{19.11 A post-quantum comparison}\label{a-post-quantum-comparison}

A nonsignalling correlation such as a PR-box fixture can exceed classical and
quantum CHSH bounds while respecting nonsignalling {[}REF-QM-009; REF-QM-010{]}.
Therefore nonlocality or
contextuality beyond a classical model does not suffice to identify quantum
theory. A reconstruction needs principles that exclude both classical and
post-quantum alternatives in the intended domain. A rational violation of the
classical bound is cheap: with \(P(a,b\mid x,y)=\tfrac14(1+ab\,E_{xy})\) on
outcomes \(\pm1\) and \((E_{00},E_{01},E_{10},E_{11})=(\tfrac45,\tfrac45,\tfrac45,-\tfrac25)\)
the model is nonsignalling with CHSH value \(\tfrac{14}5\), and it is exactly
\(\tfrac35\) of the PR box plus \(\tfrac15\) of the shared-bit model \(a=b\) plus
\(\tfrac15\) of uniform noise. A value below \(2\sqrt2\) does not by itself place a
model in the quantum set.

The fixture should be used operationally: list inputs, outputs, conditional
probabilities, normalization, nonsignalling equalities, and the correlation
value. Do not refer to ``strong contextuality'' without specifying which formal
definition and support structure are being used. CH19-P05 is such a fixture:
it lists the four contexts, the conditional probabilities, normalization,
agreement on overlaps, the CHSH value \(4\), and strong contextuality in the
sense of §19.1. It also establishes the classical bound: every noncontextual model
on the CHSH cover has CHSH value at most \(2\). CH19-O02 and Appendix J.4 prove
the quantum bound under their separately supplied finite Hilbert-space,
local-observable, tensor-product and trace assumptions. CH18-N01 retains
the distinction between that conditional theorem and reconstruction from
the primitive calculus.

\subsubsection{19.12 Quantum measurement adapter}\label{quantum-measurement-adapter}

If the optional Hilbert carrier is adopted, a finite quantum instrument can be
represented by completely positive maps indexed by outcomes whose sum is trace
preserving {[}REF-MET-006; REF-MET-007{]}. A measurement packet can then be
populated as follows:

\begin{itemize}
\tightlist
\item
  interaction: system--apparatus channel or dilation, when supplied;
\item
  registration: outcome label and apparatus record;
\item
  witness: instrument outcome linked to the declared effect;
\item
  value map: outcome-to-reported-value map;
\item
  state change: normalized post-measurement state when defined;
\item
  calibration/uncertainty: empirical characterization of the instrument;
\item
  metrological-traceability evidence: the calibration hierarchy and uncertainty
  contributions, when such traceability is claimed;
\item
  provenance/re-access: implementation, raw-record, transformation, evaluator,
  and version records.
\end{itemize}

This is an adapter specification, not a derivation of instruments. It also
shows why a POVM alone is not the whole measurement packet: a POVM gives outcome
probabilities but not in general the conditional state change, calibration
history, uncertainty budget, or authority rule.

\texttt{WDP-CH19-O04} \textbf{Finite proposition (complete finite instrument replay).} Grouped finite Gaussian-rational Kraus instruments preserve their declared probabilities, post-states and canonical dilation. On a common input/output carrier, outcome repeatability is equivalent to \(K_{jb}K_{oa}=0\) for every \(j\ne o,a,b\). Independent replay checks the complete registration, witness, accepted observation, indication, calibration enclosure and declared documentary hierarchy against actual inputs. Documentary traceability does not authenticate a physical standard or select the Born readout.

CH19-P13 and CH19-P14 give one finite adapter, for a qubit over \(\mathbb Q(i)\)
in Kraus form, with the least structure that states it. A state is a finite
mixture \(\rho=\sum_jw_j|v_j\rangle\langle v_j|\) with rational \(w_j\ge0\) and
\(\sum_jw_j\langle v_j,v_j\rangle=1\); positivity is part of the declaration, so no
spectral theorem (import 3) is used. An instrument is a family of \(2\times2\)
matrices \(K_o\) over \(\mathbb Q(i)\) with \(\sum_oK_o^\dagger K_o=1\), the
trace-preserving condition; no Stinespring dilation and no Choi matrix is used.

\texttt{WDP-CH19-P13} \textbf{Finite proposition (probabilities and state change of the qubit
adapter).} The adapter declares the outcome probabilities
\(p_o:=\sum_jw_j\langle K_ov_j,K_ov_j\rangle\) for a state given as a list of
weights and vectors. Then \(p_o=\sum_jw_j\langle v_j,E_ov_j\rangle\), the trace
\(\operatorname{tr}(E_o\rho)\) of the effect \(E_o=K_o^\dagger K_o\) against the
mixture, the \(p_o\) are nonnegative, and \(\sum_op_o=1\). For
\(p_o>0\) the post-measurement state \(K_o\rho K_o^\dagger/p_o\), the mixture
\(\{(w_j/p_o,K_ov_j)\}\), is again a state of the same form. The probabilities are
an empirical model on the one-measurement cover, so by CH19-P04 the table of
declared probabilities is the image of a finite table of accepted measurement
packets whenever every outcome can be reported. Appendix J now supplies independent native replay through the full declared measurement packet; the historical Lean qubit theorem below does not formalize that additional replay. The rule
\(p_o=\operatorname{tr}(E_o\rho)\) is the declared readout of the adapter, not a
derived one (CH19-N01). That \(\rho\) does not depend on the decomposition of the
mixture is not used.

\emph{Proof.} The Gram identity \(\langle Kv,Kv\rangle=\langle v,K^\dagger Kv\rangle\)
and completeness give \(\sum_o\langle K_ov,K_ov\rangle=\langle v,v\rangle\), so
\(\sum_op_o=\sum_jw_j\langle v_j,v_j\rangle=1\). Nonnegativity is that of
\(\langle x,x\rangle\) over \(\mathbb Q(i)\), and the trace of the post-state is
\(\sum_jw_j\langle K_ov_j,K_ov_j\rangle/p_o=1\). \(\square\)

\texttt{WDP-CH19-P14} \textbf{Finite proposition (a POVM does not determine the
instrument).} The Lüders instrument \(K_o=|o\rangle\langle o|\) and the
measure-and-prepare instrument \(M_o=|w\rangle\langle o|\) with \(w=(3/5,4/5)\) are
both complete and have the same effects \(E_o=|o\rangle\langle o|\), so they give
the same outcome probability on every state. On the input \((1,0)\) the outcome
\(0\) leaves \((1,0)\) under one instrument and \((3/5,4/5)\) under the other, and the
two vectors are not parallel.

\emph{Proof.} Direct computation over \(\mathbb Q(i)\): \(M_o^\dagger M_o=\langle
w,w\rangle|o\rangle\langle o|=|o\rangle\langle o|\). If \((1,0)=c(3/5,4/5)\) then
\(c\cdot4/5=0\), so \(c=0\) and \(1=0\). \(\square\)

So a POVM is not the whole measurement packet already at this size, as the
adapter specification above says. Calibration, uncertainty and authority are
carried by neither instrument, and CH19-P03 shows that the calibration chain of
a packet does not reach the image. Appendix I extends the native adapter to
arbitrary finite input/output dimensions and grouped Kraus maps, with a
canonical outcome-resolved isometric dilation and independent verification record replay.
Appendix J.6 supplies the repeatability criterion and describes replay of
the complete declared packet, including calibration arithmetic and documentary
traceability. These are the finite scopes recorded in CH19-O04; documentary
checks do not authenticate physical standards or identify an actual apparatus.
Unitary completion, general dilation-independence and selection of the Born
rule remain outside the stated result (CH19-N01, imports 4 and 6).

\begin{center}\rule{0.5\linewidth}{0.5pt}\end{center}

\subsection{Chapter 19A --- From certified observation to admissible transformation}\label{chapter-19a-from-certified-observation-to-admissible-transformation}

\subsubsection{19A.1 What observation determines}\label{a.1-what-observation-determines}

This chapter starts after R8--R10. Declare \(K=\mathbb Q(i)\), a carrier
\(V=K^d\), \(d\ge1\), and \(h(x,y)=x^*Hy\) with \(H=H^*>0\).
The preceding observer construction supplies positive semidefiniteness;
requiring zero radical then supplies positive definiteness. Nondegeneracy
without that positivity premise would be insufficient. The field and form
are admitted data, and the probability rule remains R10's declared readout.

For \(f,y\ne0\), write \(r_f=h(f,f)\) and
\[
\tau_f([y])=\frac{N(h(f,y))}{r_f r_y},\qquad N(z)=\bar z z.
\]
For \(d\ge2\), \(\mathbb P(K^d)\) is countably infinite; for \(d=1\),
it has one ray. On the infinite carrier, a finite example cannot certify
an arbitrary black-box permutation.

\textbf{WDP-R12-P01 --- Finite-dimensional proposition (pure-ray separation).}
The table is independent of both representatives, lies in
\(\mathbb Q\cap[0,1]\), and \(\tau_f=\tau_g\) exactly when \([f]=[g]\).
Indeed, with \(\alpha=h(f,y)/r_f\),
\[
r_f h(y-\alpha f,y-\alpha f)=r_f r_y-N(h(f,y)).
\]
Positivity gives the bounds and makes equality at one equivalent to
collinearity. Equal tables agree at the ray \([f]\); representative
independence gives the converse. The \(d^2\) probes
\(e_j,e_j+e_k,e_j+i e_k\), \(j<k\), also suffice by polarization of
\(y^*Hff^*Hy/r_f\). This is a sufficient tomography menu, not a minimal
pure-state measurement claim.

\subsubsection{19A.2 Preservers, branches and projective uniqueness}\label{a.2-preservers-branches-and-projective-uniqueness}

\textbf{WDP-R12-O01 --- Conditional result (global preserver classification).}
For \(d\ge3\), a bijection of the entire ray space preserves
orthogonality in both directions exactly when it has the form
\[
T[f]=[M\sigma(f)],\qquad M^*HM=c\,\sigma(H),\qquad
c\in\mathbb Q_{>0},\quad \sigma\in\{\mathrm{id},\mathrm{conj}\}.
\]
The branch and matrix modulo nonzero scalar are unique, and the resulting
map preserves all transition probabilities.

The proof reconstructs projective lines by double perpendiculars, then
normalizes a coordinate frame. Preservation of affine lines forces a
single coordinate map to preserve addition and multiplication. A field
automorphism fixes \(\mathbb Q\) and sends \(i\) to \(i\) or \(-i\).
Equality of the zero kernels of the transported and original forms
forces one positive rational multiplier. Basis, sum and mixed rays
prove uniqueness. The full coordinate proof is in the
\href{docs/derivations/certified_preserver_derivation.md\#3-separation-preservers-and-their-exact-scope}{preserver supplement, §3}.
The related algebraic literature is recorded as REF-PRES-001; the
written proof here does not rely on an unread classification theorem.

The projective linear-similarity group sits in a split extension by
\(C_2\). A splitting is obtained after choosing an \(h\)-orthogonal
\(K\)-basis, where the Gram matrix is rational diagonal. Bare coordinate
conjugation in an arbitrary basis need not preserve \(H\).
The conjugate branch is an antilinear map on states; it is not the
operator adjoint. Nor is a similarity automatically an R11 unitary.

\subsubsection{19A.3 Phases and the exact reachability obstruction}\label{a.3-phases-and-the-exact-reachability-obstruction}

\textbf{WDP-R12-P02 --- Finite-dimensional proposition (corrected phase group).}
Replacing \(M\) by \(aM\) changes its multiplier to \(N(a)c\).
Equal multipliers leave
\[
K^1=\{a:N(a)=1\}=
\{-1\}\cup\left\{\frac{1-t^2+2it}{1+t^2}:t\in\mathbb Q\right\}.
\]
The inverse parameter is \(\operatorname{Im}a/(1+\operatorname{Re}a)\).
This group is infinite: \((3+4i)/5\) is already outside the four
Gaussian-integer units. Only the torsion subgroup is \(\mu_4\).
Thus a chosen unitary lift can have cocycles in \(K^1\) outside \(\mu_4\);
Chapter 5's cubic phases retain their separate cyclotomic carrier,
and Chapter 18's four-unit theorem retains its integral hypothesis.

\textbf{WDP-R12-P03 --- Conditional arithmetic result (norm realization).}
The integer two-squares theorem {[}REF-NORM-001, Theorem 9.10{]}, applied
to a reduced numerator and denominator, gives
\[
\begin{gathered}
r\in N(K^\times)\iff
v_p(r)\in2\mathbb Z\ \text{for every prime }p\equiv3\pmod4,\\
\mathbb Q_{>0}/N(K^\times)\simeq
\bigoplus_{p\equiv3\ (4)}\mathbb Z/2.
\end{gathered}
\]
Denominator valuations are included. A positive witness exhibits two
rational coordinates; a negative witness exhibits a prime and its exact
odd valuation. A resource limit is neither witness.
The Gaussian Hadamard
\[
B=\frac{1+i}{2}\begin{pmatrix}1&1\\1&-1\end{pmatrix},
\qquad B^*B=I,\quad B^2=iI
\]
implements the conventional Hadamard ray action. This corrects any
reading of §17.11 that rules out that action just because the real
normalization \(1/\sqrt2\) is absent. Order-three Fourier normalization
fails the norm test; its cubic root of unity is a second, distinct
missing scalar. The conventional T eigenvalue ratio also remains outside \(K\).

\textbf{WDP-R12-O02 --- Conditional result (complete unitary ray criterion).}
A supplied positive form admits a unitary taking \([f]\) to \([g]\)
exactly when \(r_f/r_g\in N(K^\times)\). Antiunitaries do not enlarge
these orbits. For sufficiency choose \(a\) with that norm, put \(v=ag\),
and, unless \(f=v\), set
\[
w=f-v,\qquad \beta=h(w,f),\qquad Ux=x-w\,h(w,x)/\beta.
\]
Here \(h(w,w)=\beta+\bar\beta\); positivity ensures \(\beta\ne0\).
Expansion establishes \(h(Ux,Uy)=h(x,y)\) and \(Uf=v\).
The inverse replaces \(\beta\) by \(\bar\beta\). This construction
uses no square root or assumed matching of complementary Gram frames.

A similarity normalizes to a unitary precisely when its multiplier is
a Gaussian norm. Its determinant gives \(N(\det M)=c^d\), so this
holds in odd dimension. In dimension four the supplement's exact
matrix satisfies \(M^*M=3I\), preserves every probability, and crosses
unitary norm classes. The whole similarity group therefore cannot be
identified with the reversible unitary operations. Classifying the
multiplier subgroup for every supplied form is a separate question.

\subsubsection{19A.4 Postselection and the dimension-two boundary}\label{a.4-postselection-and-the-dimension-two-boundary}

\textbf{WDP-R12-P04 --- Finite-dimensional proposition (the specified Lüders branch).}
The operator \(\Pi_g=gg^*H/r_g\) is an \(h\)-selfadjoint idempotent.
The instrument \(\{\Pi_g,I-\Pi_g\}\) is complete, and
\[
\Pr(\Pi_g\mid f)=h(\Pi_g f,\Pi_g f)/r_f=\tau_f([g]).
\]
Its conditional ray is \([g]\) only at positive probability; a zero
branch has no normalized conditional state. This probability is not an
optimal operating cost. For \(f=e_1,g=(0,1,1,1)\), the direct branch
vanishes and the norm class obstructs unitarity, yet projections through
\(f+g\) reach \(g\) with probability \(3/16\). Admissible operations matter.

\textbf{CM-W1.} In dimension two, orthogonality is a perfect matching.
Swapping the axes and exchanging the orthogonal pair
\(\{[1,1],[1,-1]\}\) with \(\{[1,2],[2,-1]\}\) in both directions,
while fixing every other ray, preserves orthogonality. It changes a
transition probability from \(1/2\) to \(4/5\) and is not a semisimilarity.
The six-ray pocket is closed under orthogonal complement, which is what
justifies this extension to the entire carrier.

\textbf{WDP-R12-P05 --- Finite-dimensional proposition (full-probability qubits).}
In dimension two the full transition table does force a projectively
unique semisimilarity. Normalize the images of two orthogonal axes
and the plus ray. The imaginary plus ray then has image with coefficient
\(i\) or \(-i\). After fixing that branch, the four reference readouts
recover the norm, real part and imaginary part of every ray coordinate.
The supplement gives the argument for every positive Hermitian form.
This settles the proposed qubit conjecture at its declared field.

\subsubsection{19A.5 Stabilizers and the price of sufficient observation}\label{a.5-stabilizers-and-the-price-of-sufficient-observation}

\textbf{WDP-R12-O03 --- Conditional result (three distinct consolidation laws).}
For named readout functions \(E\) and a declared transformation group \(G\),
let \(G_E\) fix every named function pointwise and let \(\operatorname{Fix}(G)\)
be the functions fixed by every group element. Then
\[
E\subseteq\operatorname{Fix}(G)\iff G\subseteq G_E.
\]
This antitone Galois pair concerns the actual labelled observations.
For a unital \(h\)-dagger-closed \(K\)-operator algebra \(A\), one also has
\(A''=A\) and the commutant of the unitaries of \(A'\) equals \(A\).
Invariant orthogonal projections on the amplified carrier prove the
bicommutant identity; rational Cayley transforms recover the algebra
from its unitaries. These are algebraic statements without topological closure.

Finally, for an explicitly supplied finite rational linear defect space
\(R\) and linear observation map \(Q_E\), arbitrary additional rational
rows achieve injectivity with the sharp count
\[
k=\dim_{\mathbb Q}(R\cap\ker Q_E).
\]
Each row removes at most one kernel dimension; extended duals of a kernel
basis attain the bound. This is the injectivity case of Appendix T.
Restricted instrument shelves, original-carrier realization, nonlinear
transformation sets and operation costs retain their separate obligations.

For \(d\ge2\), an injective menu has trivial projective pointwise stabilizer;
the converse is false. In dimension one, the identity and conjugate
semilinear classes act identically on the single ray, so faithfulness
requires quotienting this action kernel. The supplement supplies two exact
three-dimensional effects with trivial stabilizer and two distinct rays
reporting \(1/2\) for both. A residual conjugation branch is not automatic:
an imaginary cross-term probe detects it.

\subsubsection{19A.6 From a supplied derivation to its generator}\label{a.6-from-a-supplied-derivation-to-its-generator}

\textbf{WDP-R13-P01 --- Finite-dimensional proposition (inner compatible generators).}
Let \(\delta:M_d(K)\to M_d(K)\) be \(K\)-linear and satisfy Leibniz.
For matrix units define
\[
S=\sum_j\delta(E_{j1})E_{1j}.
\]
Differentiating \(E_{kl}E_{j1}=\delta_{lj}E_{k1}\) gives
\([S,E_{kl}]=\delta(E_{kl})\), hence \(\delta=[S,\cdot]\).
If also \(\delta(A^\dagger)=\delta(A)^\dagger\), then \(S+S^\dagger\)
commutes with all matrix units and is a rational scalar matrix.
Replace \(S\) by \((S-S^\dagger)/2\), and put \(H_0=iS\). Thus
\[
\delta=-i[H_0,\cdot],\qquad H_0^\dagger=H_0,
\]
with \(H_0\) unique modulo a rational scalar. The supplied finite
matrix-unit table can be checked independently for linear extension,
every Leibniz identity, the dagger law, and the submitted generator.
This is an inner-derivation result, not Stone's theorem. No time
parameter, exponential flow, physical Hamiltonian or authorization follows.

\subsubsection{19A.7 Dependencies and evidence scope}\label{a.7-dependencies-and-evidence-scope}

R12 uses the admitted R8 involution, R9 positive form and R10 readout.
It compares its preservers with R11's operations. Chapter 5's phases and
Chapter 18's integral carrier are comparisons with different coefficients,
not hypotheses that turn \(K^1\) into \(\mu_4\). Only the linear
completion clause uses Appendix T. R13 uses a declared matrix algebra,
adjoint and derivation; it does not assume a selected physical dynamics
or need the preserver classification as an algebraic premise.

The \href{docs/derivations/certified_preserver_derivation.md}{full proof supplement},
\href{witnessed_difference_programme_claim_ledger.md\#preserver-and-generator-supplement-r12r13}{supplement evidence rows},
and \href{docs/derivation_evidence.md\#certified-preservers-and-algebraic-generators}{evidence guide}
separate written proofs, the Lean model, concrete Rust certificates, and
Python exploration. These supplement IDs do not alter the historical
154-row canonical inventory or its meaning of status coverage.
Physical reconstruction, branch selection, global synthesis, source
refinement and continuum completion retain their existing open status.

\begin{center}\rule{0.5\linewidth}{0.5pt}\end{center}

\subsection{Part V --- Relative scope and the open programme}\label{part-v-relative-scope-and-the-open-programme}

\subsection{Chapter 20 --- Relative completeness for observer-indexed presentations}\label{chapter-20-relative-completeness-for-observer-indexed-presentations}

\subsubsection{20.1 From ``cosmogenesis'' to a defensible target}\label{from-cosmogenesis-to-a-defensible-target}

The earlier title ``witness-relative cosmogenesis'' promised more than its formal
content supplied. The actual construction concerns the status of a query
relative to a finite observer profile. We therefore treat the chapter as a
calculus of \textbf{relative completeness for observer-indexed presentations}.

Let an observer profile \(\Pi\) contain a grammar of admissible queries, an input
domain, a presentation map, an evaluator, a resource bound, and a policy. For a
query \(q\), define a status only after checking that \(q\) is well typed and within
profile.

\subsubsection{20.2 Status algebra}\label{status-algebra}

A useful finite status set is:

\[
\mathsf{Status}_{\Pi}=
\{\mathsf{Proved},\mathsf{Refuted},\mathsf{Undetermined},
\mathsf{OutsideProfile},\mathsf{Refused}\}.
\]

The terms require disjoint typing rules.

\begin{itemize}
\tightlist
\item
  \texttt{Proved}: a declared proof/evidence predicate is satisfied.
\item
  \texttt{Refuted}: a declared counterevidence predicate is satisfied.
\item
  \texttt{Undetermined}: the query is in profile and well typed, but neither terminal
  predicate is established within the declared resources.
\item
  \texttt{OutsideProfile}: the query is meaningful in the ambient language but is not
  represented by this profile.
\item
  \texttt{Refused}: processing is prohibited or unsafe under a separate policy despite
  the query being syntactically processable.
\end{itemize}

An ill-typed string should be handled by a parse/type error before this status
set, rather than simultaneously being \texttt{OutsideProfile} and \texttt{Refused}.

\texttt{WDP-CH20-D01} \textbf{Definition (profile-relative completeness).} An evaluator is
complete for a finite query family \(Q\) relative to profile \(\Pi\) when every
\(q\in Q\) receives exactly one terminal status according to the declared status
rules, and the evidence attached to \texttt{Proved} and \texttt{Refuted} passes its verifier.

This is a finite coverage property. It says nothing about all mathematical
truths, all meaningful questions, or reality independent of the profile.

\subsubsection{20.3 Exclusivity and coverage}\label{exclusivity-and-coverage}

To prove that statuses are exclusive, the rules must exclude overlap. In
particular, consistency conditions should prevent both \texttt{Proved} and \texttt{Refuted}
for one query under one fixed logic and evidence base. To prove coverage, every
in-profile query must either terminate or be assigned \texttt{Undetermined} under a
bounded procedure.

\texttt{WDP-CH20-O01} \textbf{Finite proposition (resolved finite target).} A declared finite Boolean semantic model with admission, policy and bounded independent evidence replay has five exclusive and exhaustive statuses; injective value-preserving profile adapters transport verified evidence and rerun target evaluation rather than assuming status invariance under changed policy or resources.
The proof, native ownership and scope are given in Appendix I.

\subsubsection{20.4 Observer-relative does not mean arbitrary}\label{observer-relative-does-not-mean-arbitrary}

Profiles can be compared. If \(\Pi_1\) embeds faithfully into \(\Pi_2\), one can ask
whether proofs transport, whether counterexamples remain valid, and whether an
\texttt{OutsideProfile} query becomes expressible. These are naturality questions. A
status is relative to a profile, but constrained by its evidence and transport
laws.

Formally, a theory profile transport specification from a target theory \(T\) to a host substrate \(G\) uses an encoder \(E_T\) and a left-inverse decoder \(D_T\), satisfying four core commutation properties. The decoder is inverse on the encoded image; no surjectivity onto all host states is required:

\begin{enumerate}
\def\labelenumi{\arabic{enumi}.}
\tightlist
\item
  \textbf{Source round trip:}
  \[
  D_T(E_T(s)) = s \quad \text{for all } s \in \operatorname{State}_T.
  \]
\item
  \textbf{Checker fidelity:}
  \[
  \operatorname{Check}_T(\pi, \varphi) \iff \operatorname{Check}_{G,T}(E_T \pi, E_T \varphi).
  \]
\item
  \textbf{Operational step commutation:}
  \[
  E_T(\operatorname{step}_T(s, a)) \simeq \operatorname{step}_{G,T}(E_T s, E_T a).
  \]
\item
  \textbf{Observational commutation:}
  \[
  \operatorname{Obs}_G(E_T(s)) = E_O(\operatorname{Obs}_T(s)).
  \]
\end{enumerate}

Across profile boundaries, the following cumulative specification ladder adds
obligations at eight levels:

\[
\begin{aligned}
&\text{Syntactic} \longrightarrow \text{Judgmental} \longrightarrow \text{Operational} \longrightarrow \text{Observational}\\
&\longrightarrow \text{Probabilistic} \longrightarrow \text{Causal} \longrightarrow \text{Higher Coherence} \longrightarrow \text{Semantic Adequacy}.
\end{aligned}
\]

Each level retains the earlier obligations and adds its own commutation
conditions and falsifiers. The arrows describe that cumulative specification,
not implications between independently stated properties. Proving round-trip
syntactic encoding (Level 1) does not imply operational step commutation
(Level 3) or semantic truth transport (Level 8).

\subsubsection{20.5 No root observer}\label{no-root-observer}

The calculus contains no theorem selecting a privileged observer whose profile
captures every other profile. A candidate root observer would require a
initial factorisation property in a specified category and a proof that the necessary
morphisms exist uniquely in the relevant sense. Neither is present.

\texttt{WDP-CH20-N01} \textbf{Nonclaim.} Relative completeness is not absolute
omniscience, a truth-decider, a Gödel escape, world creation, or cosmogenesis.

\subsubsection{20.6 Faithful representation and certificate completeness}\label{faithful-representation-and-certificate-completeness}

\textbf{GEN-WDP-49 (conditional theorem).} Fix an effectively encoded original
problem P, sound initial witness decoder, and constructive coverage of its
solutions. Each admitted representation change supplies a backward map
preserving accepted witnesses, a forward map preserving solution coverage,
and equality of their decoded original answers. Both laws compose along
every admitted finite history. Therefore every accepted terminal witness
decodes to a solution of P, and

\[
\operatorname{Sol}(P)\ne\varnothing
\iff\exists\tau:\operatorname{Check}(\operatorname{Encode}(P),\tau)
=\mathsf{Solved}.
\]

Here \texttt{Solved} is a verified positive existential claim, represented by
\texttt{Proved}. Soundness, coverage and the existence of enough admitted resources
are explicit hypotheses. The theorem does not reinterpret the status-coverage
definition of §20.2. In particular, a sound backward adapter can discard
every source solution; only forward coverage permits negative-result transfer.
The \href{docs/faithful_representation_and_relative_completeness.md}{complete derivation}
establishes composition, original-answer soundness, certificate existence and fair
search separately. The new core-only Lean module establishes those laws without
assuming the combined end-to-end conclusion.

\subsubsection{20.7 A constructive finite constraint foundation}\label{a-constructive-finite-constraint-foundation}

\textbf{GEN-WDP-50 (finite theorem).} Given finite variable domains and finite
tables of allowed local assignments, add full-domain unary tables at uncovered
variables. Choose one permitted row from each table. Compare occurrences of
each variable using integer indicator differences against its first occurrence.
Zero raw defect is equivalent to agreement on every overlap. Gluing a
compatible family gives a unique original assignment satisfying every table;
restriction gives the inverse map. Thus original solutions and compatible
local families correspond bijectively, including the empty-domain, empty-table,
nullary and empty-variable cases.

The encoder reads supplied local tables and domains; it does not consume
global solution bits. Bijective renaming, exact joins, checked lossless
decomposition and total deterministic auxiliary variables supply constructive
forward/backward witness maps. The local defect domain is a constrained finite
set, not automatically a vector space. A compressed zero detector must be
adequate on this actual set. Choosing the whole ambient space as formal
repairs cannot manufacture an admissible solving endpoint.

The \href{code/faithful_representation/evidence/summary.json}{finite replay} compares
all 4,096 Boolean triangle relation systems with direct original-assignment
enumeration, covering 32,768 local families. Both the consistent and globally
inconsistent cycles have nonempty local relations. These finite checks,
the written all-finite proof, and the new Lean gluing/indicator/product
theorems retain their separate evidence scopes.

\subsubsection{20.8 Recursive preservation, search and stopping}\label{recursive-preservation-search-and-stopping}

\textbf{GEN-WDP-51 (conditional finite-history theorem).} Re-entry retaining the
original problem, assumptions, authorization requirements and witness
interpretation preserves that anchor under finite composition. If each
counted operation consumes one unit of a shared fixed ceiling B, a history
of n operations has spent count s\_0+n at most B. Re-entry cannot reset that
count. Explicit new budget grants are separate from re-entry; records of
requirements or grants do not establish their institutional legitimacy.

A fair enumeration of a declared effective certificate language, with
terminating checks, eventually finds an admitted finite solving witness when
one exists and admission/resources persist. For a finite local product of
size M, canonical enumeration with at most one optional proposal before each
canonical candidate needs at most 2M candidate checks, plus the admitted
transport checks. This is a search-unit bound, not a wall-clock bound.
An exhaustive negative result requires retained coverage; budget exhaustion
and interruption remain \texttt{Undetermined}. Rejected proposals do not change an
accepted mathematical state. Completeness does not follow just from strict
residual descent or an arbitrary continuation policy.

The \href{code/faithful_representation/core.py}{reference checker} independently
replays supplied witnesses, transport parents, enumeration and fuel. General
source refinement, future-operation congruence for new compressed states,
guarded-history coherence and physical reconstruction remain separate
obligations. The full statement, constructive proofs and replay instructions
are in the \href{docs/faithful_representation_and_relative_completeness.md}{supplement}.

\begin{center}\rule{0.5\linewidth}{0.5pt}\end{center}

\subsubsection{20.9 Certified progress, dynamic compression and resource claims}\label{certified-progress-dynamic-compression-and-resource-claims}

\textbf{GEN-WDP-52 (finite theorem).} From a declared finite transition/report model,
compute reachable nonempty beliefs and their universally enabled actions.
The target attractor is exactly the winning region. A natural-ranked strategy
certifies guaranteed finite success; an initial-containing, target-disjoint
region with an actual internal successor for every enabled action certifies
loss. Its dead ends remain in that region. With positive rational prices,
finite-horizon minimax constructs an exact optimum and its own strategy rank.
Independent upper/lower certificate replay establishes that optimum without
rerunning the producer. Initial membership, full winning partition, original
model targets and zero terminal values are required boundaries.

\textbf{GEN-WDP-53 (finite theorem).} A partition of this belief arena preserving
target truth, enabled actions, exact costs and successor-block sets preserves
winning and every achievable finite budget under its two-sided quotient map.
Iterative signature splitting computes the coarsest such refinement of a
supplied partition. Congruence is sufficient; accidental agreement of a
single current answer does not establish it. Quotient strategies lift, and
class-constant fine strategies descend. The mathematical quotient does not
create a physical observer realizing its distinctions.

\textbf{GEN-WDP-54 (conditional resource and guard result).} Optional operations
preserving old outcomes and prices cannot worsen the optimum. A restrictive
guard can raise finite cost or destroy winning, but a failed rejection-rule transfer
alone establishes neither loss nor original-problem impossibility. An optimum above
budget establishes budget exhaustion; an expensive feasible strategy is only an
upper bound. Reports and certificate refusal preserve their distinct meanings.
Original-answer soundness and negative-result coverage still require §20.6's
actual maps; neither a loss verification record nor a supplied authorization flag supplies
them.

The full written proofs and retained counterexamples are in Appendix S below;
the \href{docs/certified_progress_proof_and_evidence.md}{proof/evidence map} binds
39 new selected Lean declarations and independent Python replay to their
separate scopes. The local-acceptance analysis supplies the concrete local/global problem and overlap
question; Projection Obstruction/WDP identifies the information lost by a
report; checked repair verifies the proposed witness and its bound task; the
admission and continuation protocol and the discrete deformation steps must
still supply admissible realized changes and continuation. A change of report
cannot itself repair a nonzero original obstruction. This connects the owners
through preserved questions and checked maps, without equating their different
diagnostics or granting authority.

\subsubsection{20.10 Exact progress quotients and checked task binding}\label{exact-progress-quotients-and-checked-task-binding}

\textbf{GEN-WDP-55 (finite structural theorem).} An exact operation-preserving
quotient of the belief arena exists if and only if presentation fibers agree
in terminal truth and retained reports, enabled action sets, positive rational
prices, and successor class sets. Necessity is for this structural definition;
accidental equality of one task's optimum is weaker. An explicit representative
section constructs the quotient. All achievable budgets, winning and optimal
cost are preserved and reflected. Retained resolution IDs need actual common
task-witness binding; equal target booleans alone do not supply it.

The \href{docs/sharp_certified_progress_theorem.md}{self-contained sharp theorem}
reproves the attractor, certificates, finite optimal-cost horizon and coarsest
refinement, separating standard game theory from the WDP evidence specifications.
The \href{code/certified_progress/calculus.py}{bounded calculus adapter} replays the
actual archived antiderivative through its supplied task checker. It resolves
at operating cost three; a forever-unhelpful proposal-only menu loses without
making the integral unrealizable. No missing proposal-protocol stage or native review binding
is asserted. Source-sound rejection-rule transfer, budget cutoff and optional shortcut
setup costs retain their separate rejection tests.

\subsubsection{20.11 Sufficient observation and conditional physical locality}\label{sufficient-observation-and-conditional-physical-locality}

\textbf{GEN-WDP-56 (finite and conditional theorem family).} For admitted binary
errors E and repair subspace B in a syndrome kernel, syndrome-only recovery
exists exactly when each admitted syndrome fiber lies in one B-class. Added
linear observations must separate every forbidden same-syndrome difference;
the minimum is a finite hitting-suite problem. If \(E\) is a linear admitted-error
subspace, \(B\subseteq E\cap\ker Q\), and arbitrary additional linear rows are
available, the sharp count is \(\dim(\ker(Q|_E)/B)\). A supplied physical quasi-local readout
with tail norm beta distinguishes equal joint-marginal states by at most
2 beta. Reconstruction error alpha and conditioning kappa obey
delta \textless= 2 alpha + kappa eta, where eta is the actual readout separation.

The \href{docs/meron_field_limit_derivations/unified_lieb_robinson_observer_completion_derivation.md}{complete corrected locality derivation}
establishes these statements and retains Bell-marginal, gauge-cycle, CAR exchange
and false error-floor controls. A locality upper bound does not give a physical
error floor. A finite recovery theorem does not manufacture a device adapter.
The analytic Lieb--Robinson hypothesis is imported separately from the new
finite observation/conditioning deductions.

\subsubsection{20.12 Trotter smearing, time consistency and cochain descent}\label{trotter-smearing-time-consistency-and-cochain-descent}

\textbf{GEN-WDP-57 (finite-carrier time theorem and exact adapter theorem).} For
the supplied full sixteen-state, four-site CAR trace, response matrices have
size 8M and retain their declared 1/(L(2M)\^{}2) quadratic normalization.
Adjacent temporal pairing bounds the staggered average by D/(4M).
The all-M rank-one compressed contribution tends to zero under a source-derived
linear growth bound. The complete quadratic has a conservative explicit A/M
error envelope against the fixed four-site continuous-time covariance.
That covariance is positive for the declared real sources by its finite
thermal spectral representation; a strictly positive supplied margin gives
eventual strict positivity with floor(A/m)+1, not an unqualified ceiling.

The \href{docs/meron_field_limit_derivations/trotter_cochain_extinction_derivation.md}{complete corrected Trotter/cochain derivation}
provides the analytic proof, a source-derived constant-profile margin, four
independently checked finite inertia witnesses, and exact rational schedules.
A cochain adapter descends iff it kills the declared repair subspace. A low-pass
projection alone does not ensure this; weighted subdivision requires its
explicit norm convention. The staggered cochain class is not automatically
the physical CAR repair relation. Neither four spectra nor a fixed spatial
carrier's time limit closes all-M inertia or the spatial continuum programme.

Written analytic proofs, core-only Lean certificate consequences and executable
exact witnesses are bound separately in the
\href{docs/derivation_evidence.md}{derivation evidence map}. The new bridge IDs leave
the 154 canonical rows and GEN-WDP-O01--O05 at their existing scopes.

\subsubsection{20.13 Bridge Operad Equivariance and Relative Completeness Across Observers}\label{bridge-operad-equivariance-and-relative-completeness-across-observers}

\texttt{WDP-CH20-T02} \textbf{Theorem (Bridge Operad Descent Equivariance).}\\
Let \(\mathcal{O}_1, \mathcal{O}_2\) be two bounded observers with relational correspondence span \(K_{12}\). If the four congruence annihilators \((\alpha_{\mathrm{core}}, \alpha_{\mathrm{equiv}}, \alpha_{\mathrm{designation}}, \alpha_{\mathrm{defect}})\) evaluate to zero on candidate cochain \(\omega_{12}\), then the induced quotient representations descend equivariantly:
\[q_1 \sim_{\mathcal{B}_{12}} q_2 \iff [R_1] \cong_{K_{12}} [R_2]\]
without requiring the existence of a global root observer.

\texttt{WDP-CH20-C02} \textbf{Corollary (Multi-Observer Consistency Without Root Observer).}\\
Intersubjective agreement across a network of observers is certified pair-by-pair through bridge operad compositions \(\tau_{ij}\), establishing profile-relative consistency while preserving \texttt{WDP-CH20-N01} (refutation of global root observer).

\subsubsection{20.14 From an admitted distinction to witnessed reasoning}\label{from-an-admitted-distinction-to-witnessed-reasoning}

The presentation \texttt{\textless{}\textgreater{}} acquires inferential meaning through an explicitly admitted
carrier, facts, rules and interpretation. With no facts or rules, the everywhere
false model excludes every derivation. Admitting an atom as a fact changes that
result. The glyph therefore does not uniquely select the resulting logic.

For the first reasoning profile, choose a finite definite Horn theory
\(T=(A,\Gamma,R)\). A rule has a finite list of premises and one conclusion.
The judgment \(T\vdash a\) is generated by facts in \(\Gamma\) and rule firing
from already derived premises. A model makes the facts true and preserves truth
under every admitted rule. The following claims have separate bindings in the
ledger's reasoning supplement.

\begin{itemize}
\tightlist
\item
  \texttt{WDP-RSN-01}: every derivation is true in every model, by induction on the proof.
\item
  \texttt{WDP-RSN-02}: the derivable atoms themselves form the least model, hence
  derivability equals truth in all models. A countermodel proves non-entailment.
\item
  \texttt{WDP-RSN-03}: atom substitutions preserving facts and rule consequences
  transport proofs; adding an already derived fact is conservative by applying
  soundness to the old derivability model. This is the reuse/cut rule.
\item
  \texttt{WDP-RSN-04}: every checked trace entry has a derivation. Finite saturation
  inserts at most \(|A|\) atoms; it returns a proof or a closed-set countermodel.
  A smaller exhausted budget yields no negative conclusion.
\item
  \texttt{WDP-RSN-05}: encoding a directed edge \(u\to v\) as \(u\Rightarrow v\) and
  the start vertex as a fact makes proofs correspond to paths. Induction on
  proofs appends edges; induction on paths fires rules. A separate graph checker
  checks the decoded original-task witness.
\item
  \texttt{WDP-RSN-06}: a same-report, different-answer pair prevents a decoder,
  even after deterministic post-processing, as in §14.8.
\item
  \texttt{WDP-RSN-07}: for total deterministic dynamics the same argument applies
  to the whole function from finite action words to reports. An instantaneous
  collision alone does not establish this persistent limit.
\item
  \texttt{WDP-RSN-08}: if the fact and rule semantics are valid in every admitted
  original world, every checked derivation has a true decoded interpretation.
  World adequacy is an explicit premise, not supplied by the inference checker.
\end{itemize}

The \href{docs/derivations/witnessed_reasoning_calculus.md}{full derivation} gives
definitions, proofs, certificate formats, source bindings and falsifiers.
The first fifteen root-library Lean declarations verify the abstract inference
and report arguments. Section 20.15 adds executable checker refinement and
formal graph transport. Finite saturation retains its written termination proof;
whole-program Python/Rust refinement remains open. The historical monograph-wide
Lean census is unchanged.

This establishes a conditional route from admitted distinctions to checkable
reasoning and its information limits. It neither derives empirical premises nor
authorizes action. Probabilistic or defeasible reasoning, full executable
refinement, general source adapters and physical/institutional interpretation
retain their named residual obligations. Finite Horn entailment is decidable;
stronger self-reference limits cannot be inferred for it without changing the
formal system and stating the additional hypotheses.

\subsubsection{20.15 Executable refinement and kernel replay}\label{executable-refinement-and-kernel-replay}

The abstract trace now has a concrete executable realization in the root Lean
library. A finite theory contains explicit atom names, fact indices and ordered
rules. A certificate contains its source theory, goal and steps with indices
into earlier conclusions. The executable checker validates the carrier, source
and query, matches every referenced premise in order, and requires the final
conclusion to equal the requested goal.

\texttt{WDP-RSN-09} \textbf{Theorem (Concrete proof refinement).} For every finite input,
an accepted concrete step is abstractly justified. Successful replay refines
the abstract checked trace, reversing the chronological list to match its
prepend constructor. Hence an accepted proof certificate yields
\(T\vdash a\). The proof follows the executable recursion and establishes its
invariant at every step, including empty and repeated premise lists.

\texttt{WDP-RSN-10} \textbf{Theorem (Concrete countermodel refinement).} A model
certificate checks every fact, every rule, vector length, source/query equality
and a false value for the goal. These checks construct a model of \(T\) omitting
\(a\), so soundness gives \(T\nvdash a\). Positive and negative certificates cannot
both be accepted for the same task and query. A rejected or exhausted replay
provides neither judgment.

\texttt{WDP-RSN-11} \textbf{Theorem (Original graph transport).} For the edge-to-rule
encoding of §20.14, derivability is equivalent to the independent inductive
relation of directed reachability, including the zero-edge path. The two
inductions are now formalized for arbitrary finite edge lists. Composing this
equivalence with the concrete checkers turns accepted graph certificates into
reachability or non-reachability theorems for the supplied graph.

\texttt{WDP-RSN-12} \textbf{Construction (Foreign certificate replay).} Python and Rust
export concrete task and certificate data to closed Lean proof obligations.
The kernel checks each instance by reduction of the verified Boolean checker
using \texttt{by\ decide}; an externally supplied success flag supplies no premise.
Unicode scalar constructors keep names as data. The live Python adapter returns
a logical judgment only after kernel acceptance. Its source hash and query bind
the reported result to the generated obligation.

Seventeen new Lean declarations establish the executable and graph results.
The scoped replay checks 2,048 two-atom Horn queries, 64 two-vertex graph
queries, malformed Python obligations and freshly emitted native Rust cases.
The \href{docs/derivations/executable_reasoning_refinement.md}{full derivation and verification boundary}
separates the general Lean theorems from these bounded interoperability checks.
Whole-program foreign-language equivalence, serialization correctness, compiler
trust, production deployment, empirical adequacy and formal saturation bounds
retain explicit obligations. Neither a closed proof nor a countermodel grants
institutional authorization or supplies missing observations.

\subsection{Chapter 21 --- Falsifiers, nonclaims, and open obligations}\label{chapter-21-falsifiers-nonclaims-and-open-obligations}

\subsubsection{21.1 Why falsifiers belong in the monograph}\label{why-falsifiers-belong-in-the-monograph}

A finite programme can be unusually explicit about failure. Each theorem family
should name a small countermodel that would refute it, and each interface should
name the missing type that prevents composition. This is more informative than
a list of aspirations.

\subsubsection{21.2 Foundational falsifiers}\label{foundational-falsifiers}

\begin{itemize}
\tightlist
\item
  The full-symmetry nonselection theorem fails if a parameter-free equivariant
  selector exists for a moved target under the stated group action.
\item
  A minimality or uniqueness claim for a separating response fails if its
  grammar, cost, and isomorphism relation are not all stated: under item count
  the grammar of CH03-P03 has infinitely many minimal classes, and the
  refinement that leaves one class still leaves a swap pair with no
  equivariant choice.
\item
  A retention theorem fails if its proposed memory packet cannot distinguish a
  displayed pair of histories that the theorem says it retains.
\item
  A ``difference alone'' theorem fails if its construction uses an unrecorded
  label, order, measure, topology, or observer profile.
\end{itemize}

\subsubsection{21.3 Cochain falsifiers}\label{cochain-falsifiers}

\begin{itemize}
\tightlist
\item
  Quotient well-definedness fails if changing a representative changes the
  claimed value.
\item
  Cohomology language fails if \(d_1d_0\ne0\) on the declared domain.
\item
  A dual certificate fails if it does not annihilate every admissible repair or
  violates its norm constraint.
\item
  Gauge invariance fails if gauge-related inputs give nonconjugate or otherwise
  inequivalent outputs under the stated invariant.
\item
  A reported \(\Phi_1\) value for a weighted \(L^1\) norm is uncertified unless a
  repair and a bounded circulation with equal values accompany it (CH04-P07).
\item
  A path-discrepancy bound fails if two directed paths with common
  endpoints give \(|\Delta|>c_\gamma\Phi_{1,w}([\delta])\) (CH04-P11).
\item
  The signature bound fails if a class is exhibited with
  \(\Gamma_2>\|d_1\|_{J\to1}\Phi_1\) on a window whose \(d_1d_0=0\) certificate
  verifies.
\item
  Two implementations given one declared window, coefficient field, weights,
  and cochain must return one signature; a divergence refutes an implementation
  or exposes an undeclared assumption.
\item
  A coefficient-free cohomology claim fails when changing the declared field
  changes the answer, as on the projective plane (CH05-P06).
\item
  A claimed repair fails when it subtracts an element outside
  \(\operatorname{im}d_0\), such as an orthogonal projection onto \(\ker d_1\); the
  result lies in a different class, and the original class and its \(\Phi_1\)
  are not repaired.
\item
  A claim that passing to a larger complex, attaching faces, or equal-weight
  subdivision lowers \(\Phi_1\) fails (CH04-P09); only deleting cells, changing
  declared weights, or changing existing edge values can lower it.
\item
  A restriction system reported flat fails if some strict triple has
  \(r_{b,a}r_{c,b}\ne r_{c,a}\), equivalently if \(d_1d_0s\ne0\) for some section
  (CH16-P05).
\end{itemize}

\subsubsection{21.4 Time and probability falsifiers}\label{time-and-probability-falsifiers}

\begin{itemize}
\tightlist
\item
  Succession as strict order fails on a directed cycle.
\item
  Exact duration fails when two paths with the same endpoints have unequal
  declared sums.
\item
  Clock adequacy fails under an admitted target interval that has no comparable
  reference readout or exceeds the stated error bound.
\item
  Refinement invariance fails if transported atomic mass is not conserved.
\item
  Born selection from normalization is refuted by the positive-power family in
  Chapter 9.
\end{itemize}

\subsubsection{21.5 Observer and measurement falsifiers}\label{observer-and-measurement-falsifiers}

\begin{itemize}
\tightlist
\item
  A legacy observer adapter fails when it silently drops a field used by a
  theorem or invents memory not present in the source packet.
\item
  Observation validity fails when a record cannot be connected to its target
  proposition.
\item
  Measurement validity fails when the measurand specification, result carrier
  and value map, calibration, declared uncertainty/error semantics, or required
  provenance/re-access record is absent. If metrological traceability is
  claimed, absence of the documented calibration hierarchy or any calibration's
  uncertainty contribution refutes that claim.
\item
  Any implication from evaluator acceptance to authority is ill typed unless an
  authority bridge is declared.
\item
  A self-knowledge completeness claim fails whenever the presentation map has a
  displayed nontrivial fiber.
\end{itemize}

\subsubsection{21.6 History and quantum falsifiers}\label{history-and-quantum-falsifiers}

\begin{itemize}
\tightlist
\item
  A length-bounded history category fails when two admissible paths compose to
  a path beyond the bound.
\item
  A star-sum theorem over a bare semiring fails when the infinite sum has no
  declared semantics.
\item
  A zerosumfree noncancellation theorem fails in a carrier with nonzero
  additive inverses.
\item
  A quantum characterization fails on any nonquantum GPT satisfying its proposed
  discriminator set.
\item
  A contextuality/cohomology claim fails when only non-coboundary, not cocycle,
  has been shown.
\item
  An eighteen-ray parity claim fails when some context contains two
  non-orthogonal vectors or some vector does not lie in exactly two contexts;
  an absorbed source's printed table failed the first test.
\item
  A trace obstruction stated for every ring of characteristic zero fails over
  \(\mathbb Z\times\mathbb F_2\); the hypothesis is \(Nc\ne0\).
\end{itemize}

\subsubsection{21.7 Consolidated nonclaims}\label{consolidated-nonclaims}

The monograph does \textbf{not} claim:

\begin{enumerate}
\def\labelenumi{\arabic{enumi}.}
\tightlist
\item
  that Primitive Difference alone produces structure;
\item
  that witness implies memory, observer, or experience;
\item
  that retention uniquely produces time;
\item
  that finite duration is physical or continuum time;
\item
  that finite normalized weight is physical probability;
\item
  that normalization selects the Born rule;
\item
  that a quotient evaluator is a conscious observer;
\item
  that observer acceptance confers authority;
\item
  that a measurement packet is observer-independent;
\item
  that finite path sums are continuum or quantum path integrals;
\item
  that a cochain obstruction alone is quantum contextuality;
\item
  that declared complex carriers are derived from difference;
\item
  that contextuality is sufficient for quantum theory;
\item
  that the calculus reconstructs quantum mechanics;
\item
  that self-reference establishes awareness;
\item
  that profile-relative completeness is absolute completeness;
\item
  that any chapter creates worlds or establishes cosmogenesis;
\item
  that repository tests, formal elaboration, or publication status decides the
  interpretation of the mathematics;
\item
  that an obstruction signature or certificate selects a repair, stabilizes a
  residual, or routes a claim;
\item
  that combinatorial balance identities on finite complexes are physical
  conservation laws;
\item
  that a residual below an evaluator's tolerance is exact, infinite, or
  continuous, or that the tolerance status belongs to the class;
\item
  that a negative value of a declared phase-point frame, or the failure of
  \(AB-BA=cI\) on a finite carrier, is a quantum signature, a computational
  resource, or a statement about physical position and momentum.
\item
  that a positive product-projection residual (Paper 000) represents a computational or logical failure; such residuals are expected and necessary during entangling gate operations, and a raw residual bound does not constitute an obstruction modulo repairs without a declared adapter and repair space;
\item
  that nonzero physical realization error by itself prevents quantum-computing scalability; fault-tolerance threshold theorems show that scalable computation is achievable under local architectures whenever physical error rates fall below threshold.
\end{enumerate}

\subsubsection{21.8 Completed derivations and formal coverage}\label{completed-derivations-and-formal-coverage}

The nineteen registered open targets now have finite or conditional
resolutions in Appendices I--K. Their original identifiers are retained so that
the proof or counterexample can be followed from the earlier target. The
canonical registry contains no open claim or open interface.

This statement concerns the 154 canonical rows at their declared finite or
conditional scopes. The enlarged source-derivation obligations in Appendix L and the
analytic and physical-model obligations in Appendix M remain open. Their
separate register prevents a resolved finite target from being read as a
completed derivation from the primitive mark to a field theory.

The results cover retention and succession classification, clock adequacy,
finite valuations, observer adapters and descent, uncertainty and measurement
replay, self-reference, semiring coarse-graining, reconstruction countermodels,
paid composition and phase bridges, and profile-relative evaluation. Each
construction states the maps, policies or physical assumptions it uses.

Further formalization can extend the partial Lean bindings for the complete
prose results. The real-character, matrix-analysis, probability and optimization
arguments are not all mechanized. Physical calibration, Born selection,
unrestricted quantum reconstruction and external source programmes retain
their explicit nonclaim or conditional scope; the finite adapter result does
not turn them into conclusions of Primitive Difference.

\texttt{WDP-CH04-O01} \textbf{Conditional result (routing and contextual reflection).} Under a declared terminal, lift and guarded-change policy, independently verified quotient/closure diagnostics feed checked finite routing verification records. Context-equivariant profile-preserving systems admit the least same-carrier congruent coarsening, initial in the specified comma category; a profile-reflecting completion of a noncongruent source has no solution under stated axioms. Appendix J supplies the construction and refutation. The diagnostic does not select its policy.

The primitive-axiomatics source records an unconstrained initiality target and
its failure. Appendix J resolves the declared finite profile-preserving
coarsening category and refutes the incompatible profile-reflecting demand;
it does not identify these different categories.

One consequence bounds any such layer in advance. Growth by vertices,
edges, faces, or equal-weight subdivision, with a repair set that restricts
into the old one, never lowers \(\Phi_1\) (CH04-P09 and the remark after
it). A lift-admission policy over those moves therefore never admits a
candidate, and refusal is its only outcome. A policy with content must
declare a change of weights, a deletion, an edit of existing values, or an
enlarged repair set.

\subsubsection{21.9 Evidence monotonicity and revision}\label{evidence-monotonicity-and-revision}

Statuses can change when evidence or profiles change, but not every change is a
monotone upgrade. Adding a proof can move an in-profile query from
\texttt{Undetermined} to \texttt{Proved}. Adding a counterexample can move it to \texttt{Refuted}.
Changing definitions may instead create a different query. Expanding a profile
can move a query from \texttt{OutsideProfile} to an evaluable status. A policy change
can add or remove \texttt{Refused} without altering mathematical evidence.

Every status transition should therefore record one of four causes:

\begin{enumerate}
\def\labelenumi{\arabic{enumi}.}
\tightlist
\item
  \textbf{evidence addition or correction} under the same proposition and profile;
\item
  \textbf{profile change} altering admissible inputs or resources;
\item
  \textbf{statement change} producing a new proposition version;
\item
  \textbf{policy change} altering processing permission.
\end{enumerate}

Only the first is ordinarily an evidential update. If the proposition changes,
the old status must remain attached to the old proposition rather than being
silently overwritten.

A status of \texttt{Undetermined} can hide two causes that this list separates. If a
named requirement is absent---an input, a calibration record, a proof
obligation---the status can move by evidence addition. If a declared resource
bound was consumed, only a profile change that raises the bound can move it.
Recording which of the two holds says which cause could change the status; the
bare label does not.

\subsubsection{21.10 Claim identity}\label{claim-identity}

A claim identifier should bind to normalized statement text, scope, hypotheses,
and version. Reusing one ID after strengthening a quantifier, deleting a
hypothesis, or changing the carrier produces false continuity. Conversely,
typographic repair that preserves the formal statement need not mint a new
claim, but the changelog should say why the identity was preserved.

A useful claim record has fields:

\begin{Shaded}
\begin{Highlighting}[]
\NormalTok{claim\_id}
\NormalTok{statement}
\NormalTok{status}
\NormalTok{hypotheses}
\NormalTok{scope}
\NormalTok{dependencies}
\NormalTok{proof\_or\_counterexample}
\NormalTok{formal\_anchor}
\NormalTok{executable\_anchor}
\NormalTok{source\_revision}
\NormalTok{nonclaims}
\NormalTok{last\_checked}
\end{Highlighting}
\end{Shaded}

The narrative anchor IDs in this monograph are intentionally lighter. They
locate prose; they do not replace these records.

\subsubsection{21.11 Dependency-aware status}\label{dependency-aware-status}

If claim \(C\) uses claims \(A\) and \(B\), its status cannot exceed the weakest
load-bearing dependency unless \(C\) carries an independent proof. An imported
standard can support a conditional result only when the source theorem's
hypotheses are matched. A mechanized special case cannot support an unrestricted
prose theorem. A historical verification record cannot support a new revision without a
transport or rerun.

This yields a simple audit algorithm:

\begin{enumerate}
\def\labelenumi{\arabic{enumi}.}
\tightlist
\item
  parse the claim's explicit dependencies;
\item
  verify each dependency's exact statement and current evidence;
\item
  check hypothesis discharge at the use site;
\item
  search for hidden imports in definitions and examples;
\item
  downgrade or split the claim when its supported scope is narrower;
\item
  preserve the stronger target as an open row rather than deleting it.
\end{enumerate}

\subsubsection{21.12 Canonical terminology migrations}\label{canonical-terminology-migrations}

The following migrations should be applied consistently in future source
consolidation:

\begin{longtable}[]{@{}
  >{\raggedright\arraybackslash}p{(\linewidth - 4\tabcolsep) * \real{0.3333}}
  >{\raggedright\arraybackslash}p{(\linewidth - 4\tabcolsep) * \real{0.3333}}
  >{\raggedright\arraybackslash}p{(\linewidth - 4\tabcolsep) * \real{0.3333}}@{}}
\toprule\noalign{}
\begin{minipage}[b]{\linewidth}\raggedright
Legacy phrase
\end{minipage} & \begin{minipage}[b]{\linewidth}\raggedright
Canonical phrase
\end{minipage} & \begin{minipage}[b]{\linewidth}\raggedright
Reason
\end{minipage} \\
\midrule\noalign{}
\endhead
\bottomrule\noalign{}
\endlastfoot
``witness'' without modifier & typed witness kind & removes seven-way polysemy \\
``observer tuple'' & named observer layer or full layered observer & avoids false equivalence of tuples \\
``observer standing'' & evaluator standing & separates evaluation from authority \\
``authority-standing observation'' & validated observation with evaluator acceptance & repairs the measurement conflict \\
``probability from difference'' & finite normalized valuation after mass and partition inputs & exposes imports \\
``ground state'' at root & admitted ground/cut & avoids spectral homonym \\
``history integral'' & finite semiring history aggregate & avoids continuum implication \\
``interference'' for any cancellation & algebraic cancellation unless operational readout is present & preserves quantum boundary \\
``cosmogenesis'' & profile-relative completeness & matches actual content \\
``cohomology class'' without \(d_1\) & non-coboundary or quotient class & respects typing \\
``curvature'' for \(\operatorname{im}d_1\) & closure charge & no connection or geometry is declared \\
``physical gauge orbit space'' for \(Q^1\) & quotient of edge fields by vertex repairs & removes a physics import \\
``magnetic flux mass,'' ``absence of monopoles'' & closure mass \(\Gamma_2\); \(\delta^2\delta^1=0\) & interpretive vocabulary (CH05-N03) \\
\texttt{ExactClass} as \texttt{ExactRepair} & diagnostic class, routed only by a checked policy & separates diagnosis from routing (CH04-N01) \\
``observer triple \((\mathcal A,\mathcal M,\mathcal J)\)'' & evaluator cost with declared aggregation and symmetry data & an evaluator layer, not a full observer \\
``observer'' for a bare linear map & presentation map \(q_{\mathcal O}\) & a single layer (CH10-D01) \\
``sound observer'' & injective on a declared subspace & avoids collision with evaluator soundness \\
``fake zero-mode'' & reduced zero not certified by the leakage test & a typed truncation blind spot \\
``Hodge repair'' & orthogonal projection (not a repair) & projection can change the repair class \\
``finite/infinite regime'' of a residual against a floor & evaluator tolerance status (\texttt{Accept}, \texttt{Reject}, \texttt{Undetermined}) & the status belongs to the evaluator, not the class (CH10-N01) \\
\end{longtable}

\subsubsection{21.13 Editorial tests for the monograph}\label{editorial-tests-for-the-monograph}

The manuscript itself should pass finite editorial checks:

\begin{itemize}
\tightlist
\item
  every restricted word is locally scoped;
\item
  every theorem-like sentence has a status;
\item
  every external result has a registry key and consultation state;
\item
  no internal file path is treated as an external citation;
\item
  no optional quantum object appears upstream of its branch declaration;
\item
  no observation or measurement statement implies authority;
\item
  every use of ``witness'' is typed by context;
\item
  duration is defined once, with the former duplicate mapped to provenance;
\item
  open claims remain visible rather than being narrated as completed results;
\item
  nonclaims survive summary and conclusion sections.
\end{itemize}

These checks are necessary for a reliable release but do not prove the
mathematics.

\subsubsection{21.14 Worked status transition}\label{worked-status-transition}

Suppose the claim ``all accepted readings lie within \(0.1\) V of the reference'' is
\texttt{Undetermined} because only three of four sensor codes have calibration data. A
fourth reference record is added and a finite exhaustive check passes. Under
the unchanged claim, calibration model, and profile, the status can move to
\texttt{Proved} relative to those four codes. It must not become ``proved for every
voltage'' because the domain is finite and declared.

If instead the calibration model changes from interval lookup to affine
interpolation, that is a statement/profile change. The old proof remains valid
for the old model; the new claim needs a new check. If a safety policy refuses
code \texttt{3} regardless of calibration, the mathematical bound and the policy
status coexist: the reading can be calibration-valid and policy-refused.

\subsubsection{21.15 Research claims versus release claims}\label{research-claims-versus-release-claims}

A research claim answers ``what has been established under these hypotheses?'' A
release claim answers ``what exact artifact contains the statement and evidence,
and can another reader recover it?'' Both are needed for publication, but they
are not the same. A manuscript can contain a sound proof with poor artifact
custody; a well-hashed archive can contain a false theorem. The canonical
package should report both dimensions without substituting one for the other.

\begin{center}\rule{0.5\linewidth}{0.5pt}\end{center}

\subsubsection{21.16 The sharp open research program on observer completion and quantum viability}\label{the-sharp-open-research-program-on-observer-completion-and-quantum-viability}

The integration across FSP, Majorana, MUB6, Paper 000, and WDP turns the question of physical realizability into a sharply answerable research program:

\begin{quote}
\textbf{Can the available observations and recovery dynamics resolve the error classes produced by the noise, within the available time and resources, while preserving the logical quantum state?}
\end{quote}

This formulation unites the four research directions under the finite obstruction calculus:

\begin{itemize}
\item
  \textbf{FSP:} Does the chosen fermionic representation remove sign holonomy on configuration space while preserving the requested physical observables?
\item
  \textbf{Majorana:} Does the syndrome observer \(Q\) distinguish the admitted error set \(\mathcal{E}\) up to stabilizer-trivial changes (\((\mathcal{E} - \mathcal{E}) \cap \ker Q \subseteq B\)), preserving the logical Pauli/error class quotient \(H^1 = \ker Q / B\) on the binary Pauli carrier (rather than the complex Hilbert codespace itself)?
\item
  \textbf{MUB6:} Does the regional cocycle circulation obstruction \(\Phi_1 > 0\) extend to an exhaustive coverage certificate across the full complex Hadamard moduli space?

  The proposed Matolcsi--Matszangosz--Varga--Weiner (2025) Conjecture 2 conditions
  are \(G_H(\pi(\gamma_1))=0\) and
  \(G_H(\pi(\gamma_2))G_{H^*}(\pi(\gamma_2))=0\) for the twenty sign
  permutations. The retained evidence establishes the integer parity identity
  \(1+(-1)^{6m-9}=0\) for \(m\in\{0,1,2,3\}\) and checks explicitly declared
  tensor/two-circulant families. It does not establish the proposed conditions
  for every MUB transition or classify the full complex Hadamard moduli space.
  In particular, the
  \href{../../mub6/artifacts/mub_subgap_r6_closure_receipt.json}{R6 record}
  is explicitly marked \texttt{NOT\_A\_CERTIFICATE}: its Lean moduli reduction decides
  an admitted row-count table, not a theorem about all matrices. The former
  promotion from that lookup to confinement in \(F(x,y)\cup B_6(\theta)\) is
  withdrawn; a matrix-level necessity and classification proof would be needed.
  The Tao matrix supplies a separate exact negative control:
  \(G_{S_6}(\pi(\gamma_1))=9\ne0\) for all twenty permutations, as recorded in
  the \href{../../mub6/artifacts/mub_target_g1_matolcsi_conjecture2_receipt.json}{G1 artifact}.
  Thus the single-matrix Hadamard equations alone do not imply that vanishing
  condition. This says neither that Tao is a counterexample to the conjecture's
  MUB premise nor that the remaining matrices belong to two named families.
  Connecting local circulation or Positivstellensatz certificates to an
  exhaustive, hypothesis-correct MUB6 cover remains open.
\item
  \textbf{Paper 000 / WDP:} Does the discarded correlation residual lie in an unrepairable class modulo stabilizers under local Lieb--Robinson dynamics and bounded readout precision?
\end{itemize}

\begin{center}\rule{0.5\linewidth}{0.5pt}\end{center}

\subsection{Chapter 22 --- Synthesis: what the finite calculus establishes}\label{chapter-22-synthesis-what-the-finite-calculus-establishes}

\subsubsection{22.1 The trunk}\label{the-trunk}

The defensible trunk of the programme is modest and coherent.

An admitted cut contains exhibited difference and the structure needed to type
it. A declared comparison can fail, leaving a residual. Typed witnesses can
locate, retain, repair, separate, or refute claims about that residual, but those
roles require explicit adapters. Full symmetry blocks a deterministic
equivariant point choice of a moved target. It does not block invariant
probability laws; realizing a point from such a law adds a random outcome or
seed. Every successful deterministic point selection must depend on a declared
symmetry-breaking input. When the candidate input family is finite and removing
the full declared dependency set restores a symmetry that moves the output,
that dependency set contains at least one inclusion-minimal frontier element;
a particular implementation may also carry redundant inputs. On a declared
finite graph, cochains and
quotient defects provide exact mathematical tools;
finite-dimensional optimization yields bounded propositions and certificates.

Retention plus an admitted transition relation supports finite succession.
Succession plus an additive valuation supports duration. Exhibited finite
witness families plus declared masses and partitions support normalized finite
probability. Each layer has a clear input ledger and none supplies its usual
physical interpretation for free.

\subsubsection{22.2 The observer/measurement waist}\label{the-observermeasurement-waist}

The central architectural result of the synthesis is the separation

\[
\text{witness}
\longrightarrow
\text{observer presentation/memory/evaluation}
\longrightarrow
\text{observation}
\longrightarrow
\text{measurement},
\]

with authority outside the chain.

A witness is typed evidence for a proposition. An observer packages a
presentation boundary, memory, evaluator, and policy. An observation is a
registered, witnessed, re-accessible interaction accepted by an evaluator. In
the WDP interface, a measurement adds a measurand specification, a distinct
reported-value carrier and map, scale/unit, calibration, uncertainty, any
claimed metrological-traceability evidence, and separate provenance/re-access.
This ``waist'' lets duration and probability meet operational records without
claiming that either was derived as physical law.

\subsubsection{22.3 The branches}\label{the-branches}

Histories form one branch. Projection can fail to descend; if a single-valued
deterministic update on the unchanged presentation carrier is retained, extra
state or history is then required. Refinement or relation-valued, set-valued, or
stochastic dynamics are alternatives. Finite path aggregates over semirings are
well defined at fixed length and have a matrix representation under the stated
admissibility hypothesis. Infinite sums and quantum readings need more.

Quantum-capable carriers form another branch. Declared complex fibers, unitary
connections, holonomy, and connection Laplacians support rigorous finite
geometry. They remain assumption-bearing. A reconstruction audit shows the
additional assumptions: empirical contexts, composition, continuity, purification or
an alternative discriminator, and a probability rule stronger than
normalization.

Relative completeness forms a third branch. A finite evaluator can cover a
declared query family with exclusive statuses if its rules are proved correct.
That is an engineering and logical property of a profile, not an origin story
for a universe.

\subsubsection{22.4 The four-tier ambition hierarchy}\label{the-four-tier-ambition-hierarchy}

The entire programme is organized across four distinct levels of foundational ambition:

\begin{enumerate}
\def\labelenumi{\arabic{enumi}.}
\tightlist
\item
  \textbf{Tier 1 --- Syntactic Representation:} Demonstrated for the declared finite encodings. The round-trip claims concern their specified source and image carriers; selected implementation checks retain their recorded scope. Encoding a packet does not establish its operational interpretation.
\item
  \textbf{Tier 2 --- Faithful Operational Representation:} Demonstrated for the stated constructions and adapter laws. Finite transition quotients preserve their declared reports and updates; permutation actions and cochain maps preserve their specified algebra. Formation, checker fidelity, operational steps and observational answers remain separate obligations for each theory-profile adapter (§20.4).
\item
  \textbf{Tier 3 --- Reconstruction Under Explicit Principles:} Characterized as an audit framework. Domain structures such as empirical contextuality covers, operator-algebraic states, and continuous measurement uncertainties are reconstructed only by adjoining explicit operational axioms and cataloging their external interfaces.
\item
  \textbf{Tier 4 --- Unique Selection From Primitive Kernel:} Explicitly refused. The programme does not claim that textbook continuous spacetime, the complex field \(\mathbb{C}\), the continuous unitary circle group \(U(1)\), Gleason frame uniqueness, or the Born probability rule is forced or uniquely selected by primitive difference. It establishes no theorem that they cannot be; the nonclaims and the finite witnesses of §18.5 and Appendix D mark where the programme stops.
\end{enumerate}

Maintaining strict boundaries between these four tiers prevents confusing a verified syntactic embedding (Tier 1) or finite operational model (Tier 2) with unrestricted ontological necessity (Tier 4).

\subsubsection{22.5 Strongest defensible conclusion}\label{strongest-defensible-conclusion}

\texttt{WDP-CH22-C01} \textbf{Synthesis conclusion.} Finite witnessed-difference systems can
be organized into an explicit obstruction calculus in which every transition
from residual to retention, order, magnitude, probability, observer,
observation, measurement, history, or quantum-capable representation is exposed
as a typed construction with a visible assumption cost.

This conclusion is architectural. It does not claim that all listed open
constructions exist or are unique. Its value is that failed derivations become
local: one can point to the missing order, map, law, calibration, or theorem
rather than covering it with vocabulary.

\subsubsection{22.6 Closing boundary}\label{closing-boundary}

The monograph ends where its evidence ends. It supplies finite definitions,
written proofs of the scoped finite and conditional theorem families, selected
formalizations, bounded computations, and an explicit map of remaining
obligations. Each kind of support retains its stated scope and reproduction
status. It does not turn those materials into a
physical theory by declaration. A future result may close one of the interfaces;
if so, the new hypothesis and proof should be attached at that interface, and
the nonclaims elsewhere should remain intact.

\begin{center}\rule{0.5\linewidth}{0.5pt}\end{center}

\subsection{Chapter 23 --- Observer completion on a finite 2D toric-code carrier}\label{chapter-23-observer-completion-on-a-finite-2d-toric-code-carrier}

\textbf{Active exposition status, 2026-10-07.} The active exposition of this chapter develops the corrected mathematics of finite toric observer completion. The historical draft text originally imported with withdrawn global-kernel and restricted-error detector-rank/adaptivity arguments has been retired to \href{docs/source_archive/toric_code_chapter_23_import.md}{\texttt{docs/source\_archive/toric\_code\_chapter\_23\_import.md}}. The original draft identifiers (\texttt{WDP-CH23-D01} through \texttt{WDP-CH23-P20}) are mapped to their corrected statements (\texttt{WDP-CH23-P01C} through \texttt{WDP-CH23-P18C}) in the crosswalk register of §23.8 and Appendix U.

\textbf{What this chapter does.} It instantiates the operational question of the previous session --- \emph{given what the observer already sees, which further observations are necessary and sufficient, and what is the cheapest way to obtain them?} --- on a declared finite carrier with actual instruments and prices. The finite procedures return a sufficient observation set with a checked lower bound on its cost. Adaptivity depends on the declared candidate set and instruments: P17C gives a worst-case adaptive cost of four against a fixed-suite cost of six. The chapter is a \textbf{conditional result}: it presupposes a declared noise model, a declared schedule, and a declared precision, and it says nothing about any physical device.

\begin{center}\rule{0.5\linewidth}{0.5pt}\end{center}

\subsubsection{23.1 The declaration packet}\label{the-declaration-packet}

\texttt{WDP-CH23-D01} \textbf{Definition (toric-code completion packet).} A \emph{completion packet} is a tuple

\[
\mathfrak T=\bigl(\Lambda,\ \mathcal{S},\ \mathcal{D},\ \mathcal{T},\ \mathcal{P},\ \mathcal{E},\ q,\ B\bigr)
\]

with the following fields, all finite and all declared:

\begin{enumerate}
\def\labelenumi{\arabic{enumi}.}
\tightlist
\item
  \textbf{Carrier.} \(\Lambda=\mathbb Z_L\times\mathbb Z_L\) with \(L\ge 3\), horizontal qubits \(h_{i,j}\) joining \((i,j)\) to \((i,j+1)\) and vertical qubits \(v_{i,j}\) joining \((i,j)\) to \((i+1,j)\) (indices modulo \(L\)), \(n=2L^{2}\) data qubits, and the Pauli carrier \(\mathcal P=\mathbb F_2^{\,n}\oplus\mathbb F_2^{\,n}\) with \(e=(e_X,e_Z)\) and weight convention
  \[
  \mathrm{wt}(e)=|\mathrm{supp}e_X|+|\mathrm{supp}e_Z|,
  \]
  so a single \(Y\) contributes \(2\).
\item
  \textbf{Stabilizers.} \(S=\langle\sigma^\star_v,\ \sigma^\square_f\rangle\), \(\sigma^\star_v\) the \(X\)-star of site \(v\), \(\sigma^\square_f\) the \(Z\)-plaquette of face \(f\).
\item
  \textbf{Detector locations.} A \emph{detector} is a pair \(d=(\sigma,k)\) with \(\sigma\) a stabilizer generator and \(k\) a round index. A schedule \(D\) is a finite subset of \(\mathcal S\times[R]\); the \emph{dark windows} of \(D\) are the maximal open intervals of \([R]\) containing no round whose full stabilizer set is in \(D\).
\item
  \textbf{Classical communication.} Per round, a measurement duration \(\tau_m\) and a link latency \(\tau_c\); the syndrome of round \(k\) is available to the policy at time \(t_k=k(\tau_m+\tau_c)+\tau_c\).
\item
  \textbf{Precision.} Each detector returns one bit (exact readout), or --- in the noisy branch of §23.6 --- one bit subject to a declared independent flip with rational parameter \(\varepsilon\), or no value at all (that is, the location is dark).
\item
  \textbf{Admissible errors.} \(\mathcal E\subseteq\mathcal P\), declared; the standard instance is \(\mathcal E_t=\{e:\mathrm{wt}(e)\le t\}\).
\item
  \textbf{Protected question.} A group homomorphism \(q:\mathcal P\to A\) with \(B=\ker q\) and \(S\subseteq B\). The standard instance is \(q(e)=[e]\) in \(\mathcal P/S\), with \(B=S\): \emph{the recovery left a stabilizer, so the logical state is the prepared one.}
\item
  \textbf{Recovery.} A policy \(\rho:\mathrm{im}\,Q_D\to\mathcal P\) applied to the accumulated detector record.
\end{enumerate}

\texttt{WDP-CH23-D02} \textbf{Definition (measured history record).} For a schedule
\(D\subseteq\mathcal S\times[R]\), put
\(D_{\le k}=\{(\sigma,j)\in D:j\le k\}\). Its record is
\[
Q_{D,k}(e^{(1)},\ldots,e^{(k)})
 =\bigl(\operatorname{syn}_{\sigma}(e^{(j)})\bigr)_{(\sigma,j)\in D_{\le k}}
 \in\mathbb F_2^{D_{\le k}}.
\]
Each coordinate records the error at its actual measurement time. A previous
measurement does not reveal a later syndrome. With no measurement error,
one complete round observes
\[
Q(x\mid z)=(Ax,A_vz),
\]
where \(A\in\mathbb F_2^{L^2\times n}\) is the face--qubit incidence matrix
and \(A_v\in\mathbb F_2^{L^2\times n}\) is the site--qubit incidence matrix.
The rank, separation and observation-cost results below concern this
single-round map, or repeated tests of one fixed error. They do not solve
recovery for an evolving error history.

\begin{center}\rule{0.5\linewidth}{0.5pt}\end{center}

\subsubsection{\texorpdfstring{23.2 Exact structure of the commutation kernel: \(S^\perp\), not \(S\)}{23.2 Exact structure of the commutation kernel: S\^{}\textbackslash perp, not S}}\label{exact-structure-of-the-commutation-kernel-sperp-not-s}

Let \(A\) be the \(L^2\times n\) face-incidence matrix and \(A_v\) the star-incidence matrix over \(\mathbb F_2\). Each edge lies in two faces, so the sum of all rows of \(A\) is zero. The dual graph is the connected torus grid, so the left kernel is the constant vectors and \(\operatorname{rank} A=L^2-1\). Similarly, \(\operatorname{rank} A_v=L^2-1\).

On the Pauli space of dimension \(4L^2\), the block commutation syndrome \(Q(x\mid z)=(Ax,A_v z)\) has rank
\[
\operatorname{rank} Q=(L^2-1)+(L^2-1)=2L^2-2,
\]
so its kernel has dimension
\[
\dim\ker Q=4L^2-(2L^2-2)=2L^2+2.
\]
The stabilizer space \(S\) is spanned by the star rows in the \(X\) block and face rows in the \(Z\) block. Commutation relations give \(AA_v^{\mathsf T}=0\), so \(S\subseteq\ker Q\). Since the two blocks are independent, \(\dim S=2L^2-2\). Hence
\[
\dim(\ker Q/S)=(2L^2+2)-(2L^2-2)=4.
\]

\texttt{WDP-CH23-P01C} \textbf{Finite proposition (commutation kernel structure).} On this carrier, with this weight convention, the full stabilizer commutation syndrome has kernel \(S^\perp\neq S\), and \(\dim(S^\perp/S)=4\). The kernel of \(Q\) is the set of Pauli vectors commuting with every generator of \(S\) (the symplectic orthogonal \(S^\perp\)), which properly contains \(S\). An exhibiting vector is the weight-\(L\) support \(Z_1\) (horizontal edges \(h_{i,0}\) for \(i=0,\ldots,L-1\)) placed in the \(X\) block with \(Z\) block zero: it lies in \(\ker A\), outside the star span, and outside \(S\).

\begin{center}\rule{0.5\linewidth}{0.5pt}\end{center}

\subsubsection{23.3 Distance under declared weight and coarser-question thresholds}\label{distance-under-declared-weight-and-coarser-question-thresholds}

Let \(x\in\ker A\). Let \(\alpha_i(x)\) be the sum of horizontal edges in lattice row \(i\), and \(\beta_j(x)\) the sum of vertical edges in column \(j\). Summing the face constraints along a row gives \(\alpha_i=\alpha_{i+1}\); summing them along a column gives \(\beta_j=\beta_{j+1}\). Both fluxes vanish on the star span, so \((\alpha,\beta)\) descends to an isomorphism \(\ker A/\operatorname{rowspan}(A_v)\cong\mathbb F_2^2\), as proved in Appendix U.2.

If \(\alpha=1\), each row contributes at least one edge, giving weight \(\ge L\). The same bound holds in each sector.

\texttt{WDP-CH23-P02C} \textbf{Finite proposition (code distance).} \(\min\{\operatorname{wt}(e): e\in\ker Q\setminus S\}=L\), and a weight-\(L\) vector exists in each sector.

A class nontrivial in both sectors has weight at least \(2L\) under the declared additive weight.

\texttt{WDP-CH23-P07C} \textbf{Finite proposition (diagonal distance and coarser-question threshold).} Under \(\operatorname{wt}(e)=|\operatorname{supp}e_X|+|\operatorname{supp}e_Z|\), the diagonal class has distance \(2L\), not \(2L-1\). However, this distance does not establish the fault-tolerance threshold \(t\le L-1\) for the coarser question \(B=S+\langle u+v\rangle\). The relevant separation distance is
\[
d_B=\min\{\operatorname{wt}(z):z\in\ker Q\setminus B\},
\]
measuring forbidden differences. Because the diagonal class \(u+v\) belongs to the allowed subspace \(B\), its weight of \(2L\) does not govern \(d_B\). The forbidden difference set \(\ker Q\setminus B\) still contains the pure single-sector logical classes \(u\) and \(v\), each of weight \(L\). Thus \(d_B=L\), retaining the single-sector threshold \(2t<L\). The inference that \(2t<2L\) tolerates \(t\le L-1\) fails: at \(L=3, t=2\), the errors \(e_1=h_{0,0}\) and \(e_2=\{h_{1,0},h_{2,0}\}\) have weights one and two, identical full syndromes, and difference \(u\notin B\), producing an undetectable collision.

\texttt{WDP-CH23-P08C} \textbf{Withdrawal of halving slogans.} The imported slogans that diagonal coarsening halves the required lattice and is factorially cheaper are withdrawn. No cost or lattice-size identity of that form is proved. The coarser question retains the threshold requirement \(2t<L\).

\begin{center}\rule{0.5\linewidth}{0.5pt}\end{center}

\subsubsection{23.4 The restricted fiber criterion and failure certificates}\label{the-restricted-fiber-criterion-and-failure-certificates}

Let \(E_{\mathrm{adm}}\) be a declared finite set of Pauli vectors, \(Q_D\) a declared detector map, and \(q\) a protected question with \(B=\ker q\).

\texttt{WDP-CH23-P03C} \textbf{Finite proposition (restricted fiber criterion).} For every admitted pair \(e,e'\in E_{\mathrm{adm}}\),
\[
Q_D(e)=Q_D(e')\implies q(e)=q(e')
\]
if and only if
\[
(E_{\mathrm{adm}}-E_{\mathrm{adm}})\cap\ker Q_D\subseteq B.
\]
\texttt{WDP-CH23-P04C} \textbf{Corollary (primal recovery table).} When the inclusion holds, choosing one admitted representative \(e_s\) in each nonempty syndrome fiber gives an exact recovery policy \(\rho(s)=e_s\) such that \(e+\rho(Q_D e)\in B\) for all \(e\in E_{\mathrm{adm}}\). Both construction and verification terminate by finite enumeration.

\texttt{WDP-CH23-P05C} \textbf{Corollary (dual failure certificate).} When the inclusion fails, finite search returns an admitted pair \(e_1,e_2\in E_{\mathrm{adm}}\) such that \(Q_D e_1=Q_D e_2\) but \(e_1-e_2\notin B\), providing an independently verifiable failure witness.

\texttt{WDP-CH23-P10C} \textbf{Finite proposition (global rank vs.~admitted collisions).} A global rank defect is not an admitted collision. Dropping one independent check raises \(\dim(\ker Q_D/S)\) from \(4\) to \(5\). At \(t=0\), \(E_{\mathrm{adm}}=\{0\}\), so \((E_{\mathrm{adm}}-E_{\mathrm{adm}})=\{0\}\subseteq B\). The empty purchase is sound at \(t=0\). Whether an enlarged kernel contains an admitted collision is determined by searching \(E_{\mathrm{adm}}-E_{\mathrm{adm}}\), not by quoting the global quotient dimension.

\begin{center}\rule{0.5\linewidth}{0.5pt}\end{center}

\subsubsection{\texorpdfstring{23.5 Full rank sufficiency, \(t=0\) purchase, and the matroid boundary}{23.5 Full rank sufficiency, t=0 purchase, and the matroid boundary}}\label{full-rank-sufficiency-t0-purchase-and-the-matroid-boundary}

Let \(\mathcal E_t=\{e\in\mathcal P:\operatorname{wt}(e)\le t\}\) and \(B=S\).

\texttt{WDP-CH23-P06C} \textbf{Finite proposition (sufficiency of full rank).} For \(1\le t\) and \(2t<L\) with \(B=S\), a passive stabilizer-check set of rank \(2L^2-2\) separates \(\mathcal E_t\) up to \(S\). Sufficiency follows from \(d=L\): any vector of weight \(\le 2t<L\) in \(\ker Q=S^\perp\) already belongs to \(S\). Full rank is not necessary for separation: at \(L=4\), twenty-nine independent checks separate \(\mathcal E_1\) (\(t=1, 2t=2<4\)), dropping one check without producing an admitted collision.

\texttt{WDP-CH23-P11C} \textbf{Finite proposition (uniform cost minimum and zero purchase).} At \(t=0\), the empty detector set is sufficient and has cost \(0\). For \(1\le t\) and \(2t<L\) with \(B=S\), rank \(2L^2-2\) is sufficient but need not be the minimum independent count. The strict saving stated here is the \(L=4,t=1\) example; it is not asserted for every \(L,t\).

\texttt{WDP-CH23-P12C} \textbf{Restriction of the matroid shortcut.} Greedy selection in the stabilizer-check matroid returns a minimum-cost basis only when sufficiency is proved equivalent to full rank. For general \(E_{\mathrm{adm}}\), sufficiency is the hitting condition of P03C, requiring exact enumeration under positive rational costs.

\texttt{WDP-CH23-P13C} \textbf{Finite proposition (uniform-cost check bounds).} The uniform-cost minimum on a declared finite shelf is returned by exhaustive enumeration. Every detector in a minimal sufficient suite is essential for that suite, but need not span the full check space.

\begin{center}\rule{0.5\linewidth}{0.5pt}\end{center}

\subsubsection{23.6 Adaptivity, schedules, and spacetime consistency}\label{adaptivity-schedules-and-spacetime-consistency}

\texttt{WDP-CH23-P17C} \textbf{Finite proposition (refutation of unrestricted adaptive lower bounds).} The imported unrestricted claim that adaptivity buys nothing is false. On a four-hypothesis toy model with tests of cost 2, nonadaptive separation requires three tests (cost 6), while an adaptive policy requires two tests (cost 4). The claimed lower bound of 6 is refuted.

\texttt{WDP-CH23-P18C} \textbf{Refutation of global purchase bounds.} Every unrestricted detector-purchase bound requiring full check rank along every branch is withdrawn. The \(t=0\) family and the \(L=4\) rank-29 suite are explicit counterexamples.

\texttt{WDP-CH23-P15C} \textbf{Conditional status of causal-cone bounds.} Comparing the cone radius \(\lceil v_{\mathrm{LR}}\Delta t\rceil\) with \(\lfloor L/2\rfloor\) is a geometric comparison, not a proof of schedule failure. A failure certificate requires an explicit dual pair in \((E^R-E^R)\cap\ker Q_{D,R}\setminus B^R\).

\texttt{WDP-CH23-P16C} \textbf{Spacetime consistency and passive readout.} The spacetime chain complex, Bernoulli readout model, and repetition threshold \(R\ge L\) remain conditional on an explicit joint error model and state-evolution map. Detectors are passive queries; a classical bit table does not prove physical non-disturbance.

\begin{center}\rule{0.5\linewidth}{0.5pt}\end{center}

\subsubsection{23.7 Terminating observer-completion procedure}\label{terminating-observer-completion-procedure}

\texttt{WDP-CH23-D16} \textbf{Definition (\texttt{ToricObserverCompletion}).} Given \((\Lambda,\mathcal S,D,\mathcal E,B)\), the procedure:
1. enumerates Pauli coordinates of the declared carrier;
2. constructs \(H_D\);
3. computes \(\ker H_D\);
4. enumerates \(\mathcal E-\mathcal E\);
5. tests \((\mathcal E-\mathcal E)\cap\ker H_D\subseteq B\);
6. if true, outputs primal recovery dictionary \(\rho(\sigma)=e_\sigma\);
7. if false, returns dual failure pair \(e_1,e_2\).

\texttt{WDP-CH23-P19} \textbf{Finite proposition (termination and correctness).} The procedure terminates for every finite instance and outputs either an exact primal recovery dictionary or an exact dual failure pair.

\texttt{WDP-CH23-P20} \textbf{Finite proposition (complexity).} For \(d=2L^2\), \(r=|D|\), and \(m=|\mathcal E|\), Gaussian elimination uses \(O(d^3+d^2 r)\) operations over \(\mathbb F_2\). Testing membership in \(B\) is polynomial in \(d\), and difference enumeration is polynomial in \(m\).

\begin{center}\rule{0.5\linewidth}{0.5pt}\end{center}

\subsubsection{23.8 Identifier crosswalk and archival register}\label{identifier-crosswalk-and-archival-register}

The original draft text of Chapter 23 is archived in \href{docs/source_archive/toric_code_chapter_23_import.md}{\texttt{docs/source\_archive/toric\_code\_chapter\_23\_import.md}}. The crosswalk below maps each imported draft identifier to its canonical disposition and corrected identifier:

\begin{longtable}[]{@{}
  >{\raggedright\arraybackslash}p{(\linewidth - 8\tabcolsep) * \real{0.2000}}
  >{\raggedright\arraybackslash}p{(\linewidth - 8\tabcolsep) * \real{0.2000}}
  >{\raggedright\arraybackslash}p{(\linewidth - 8\tabcolsep) * \real{0.2000}}
  >{\raggedright\arraybackslash}p{(\linewidth - 8\tabcolsep) * \real{0.2000}}
  >{\raggedright\arraybackslash}p{(\linewidth - 8\tabcolsep) * \real{0.2000}}@{}}
\toprule\noalign{}
\begin{minipage}[b]{\linewidth}\raggedright
Draft ID
\end{minipage} & \begin{minipage}[b]{\linewidth}\raggedright
Topic
\end{minipage} & \begin{minipage}[b]{\linewidth}\raggedright
Draft Claim
\end{minipage} & \begin{minipage}[b]{\linewidth}\raggedright
Canonical Disposition
\end{minipage} & \begin{minipage}[b]{\linewidth}\raggedright
Corrected ID
\end{minipage} \\
\midrule\noalign{}
\endhead
\bottomrule\noalign{}
\endlastfoot
\texttt{D01} & Packet & Toric completion packet & Proved definition & \texttt{WDP-CH23-D01} \\
\texttt{D02} & Detection & Cumulative syndrome map & Proved definition & \texttt{WDP-CH23-D02} \\
\texttt{P01} & Kernel & \(\ker Q = S\), quotient dim 2 & Corrected: \(\ker Q = S^\perp \ne S\), quotient dim 4 & \texttt{WDP-CH23-P01C} \\
\texttt{P02} & Distance & Distance \(L\) & Corrected: distance \(L\) via winding fluxes & \texttt{WDP-CH23-P02C} \\
\texttt{P03} & Fiber & \((E-E)\cap\ker Q\subseteq B\) & Corrected: sound equivalence proof & \texttt{WDP-CH23-P03C} \\
\texttt{P04} & Recovery & Primal recovery dictionary & Proved & \texttt{WDP-CH23-P04C} \\
\texttt{P05} & Dual & Dual failure pair & Proved with explicit search & \texttt{WDP-CH23-P05C} \\
\texttt{P06} & Sufficiency & \(2t<L\) iff full rank & Corrected: full rank sufficient, not necessary & \texttt{WDP-CH23-P06C} \\
\texttt{P07} & Diagonal & Diagonal distance \(2L-1\), \(t\le L-1\) & Corrected: distance \(2L\), but \(d_B=L\) retaining \(2t<L\) & \texttt{WDP-CH23-P07C} \\
\texttt{P08} & Slogans & Halves lattice, factorially cheaper & Refuted and withdrawn & \texttt{WDP-CH23-P08C} \\
\texttt{P09} & Terminal & Terminal error convention & Proved for stated convention & \texttt{WDP-CH23-P09} \\
\texttt{P10} & Omission & Dropped check gives low-weight collision & Corrected: rank defect is not admitted collision & \texttt{WDP-CH23-P10C} \\
\texttt{P11} & Check count & Full rank necessary for \(2t<L\) & Corrected: false at \(t=0\) and for \(L=4\) rank-29 & \texttt{WDP-CH23-P11C} \\
\texttt{P12} & Matroid & Greedy matroid basis & Corrected: restricted to verified full-rank cases & \texttt{WDP-CH23-P12C} \\
\texttt{P13} & Uniform cost & Minimum count \(2L^2-2\) & Corrected: enumeration minimum, not rank bound & \texttt{WDP-CH23-P13C} \\
\texttt{P14} & Dichotomy & Primal/dual certificate dichotomy & Proved & \texttt{WDP-CH23-P14} \\
\texttt{P15} & Causal cone & Causal cone failure threshold & Conditional on error propagation model & \texttt{WDP-CH23-P15C} \\
\texttt{P16} & Spacetime & Spacetime consistency chain & Conditional on joint error model & \texttt{WDP-CH23-P16C} \\
\texttt{P17} & Adaptivity & Worst-case adaptivity buys nothing & Refuted as unrestricted claim; toy countermodel & \texttt{WDP-CH23-P17C} \\
\texttt{P18} & Cost & Adaptive cost \(\ge 2L^2-2\) & Refuted as unrestricted claim & \texttt{WDP-CH23-P18C} \\
\texttt{P19} & Procedure & Algorithm termination & Proved & \texttt{WDP-CH23-P19} \\
\texttt{P20} & Complexity & Bit-operation bounds & Proved & \texttt{WDP-CH23-P20} \\
\end{longtable}

\begin{center}\rule{0.5\linewidth}{0.5pt}\end{center}

\subsubsection{23.9 Claim hygiene and boundaries}\label{claim-hygiene-and-boundaries}

\subsubsection{1. What this chapter supports}\label{what-this-chapter-supports}

\begin{itemize}
\tightlist
\item
  The full commutation syndrome has kernel \(S^\perp\ne S\) with logical quotient dimension 4.
\item
  The exact observer completion condition is \((E_{\mathrm{adm}}-E_{\mathrm{adm}})\cap\ker Q_D\subseteq B\).
\item
  For \(\mathcal E_t\), \(B=S\), and \(2t<L\), full independent stabilizer rank \(2L^2-2\) is sufficient, but not necessary.
\item
  For coarser diagonal questions \(B=S+\langle u+v\rangle\), the forbidden distance is \(d_B=\min\{\operatorname{wt}(z):z\in\ker Q\setminus B\}=L\), retaining the single-sector threshold \(2t<L\).
\item
  At \(t=0\), the empty detector set is sufficient and costs 0.
\item
  Every finite instance admits either an exact primal recovery dictionary or an exact dual failure pair.
\item
  Unrestricted claims that adaptivity never reduces cost or that full rank is required along every branch are refuted.
\end{itemize}

\subsubsection{2. What this chapter does not support}\label{what-this-chapter-does-not-support}

\begin{itemize}
\tightlist
\item
  A physical toric-code quantum device, its hardware, calibration, or noise spectra.
\item
  Fault tolerance under undeclared stochastic, correlated, or adversarial noise models.
\item
  Increase of fault tolerance beyond \(2t<L\) for diagonal-stabilizer quotients.
\item
  Universality of matroid greedy selection for general error-separation tasks.
\item
  Non-disturbance of quantum states by passive classical readouts.
\item
  Computational scalability for unbounded lattice sizes or spacetime histories.
\item
  Any claim about human cognition, general intelligence, or consciousness.
\end{itemize}

\subsubsection{3. Narrowest defensible claim}\label{narrowest-defensible-claim}

Under the declared finite 2D toric carrier and Pauli error family, observer completion is an unresolved quotient-homology problem governed by \((E_{\mathrm{adm}}-E_{\mathrm{adm}})\cap\ker Q_D\subseteq B\). For the CSS commutation syndrome, the kernel is \(S^\perp\) with quotient dimension 4; full rank \(2L^2-2\) is sufficient for \(1\le t, 2t<L, B=S\), but not necessary; coarser diagonal questions retain \(d_B=L\) and threshold \(2t<L\); and every finite instance terminates with an exact primal recovery table or dual failure witness.

\subsubsection{4. Next test needed}\label{next-test-needed}

Calibrated execution pairing measured lattice error statistics with the finite kernel solver, verifying that the predicted primal dictionary or dual failure pair matches observed error-correction performance on physical benchmark data.

\begin{center}\rule{0.5\linewidth}{0.5pt}\end{center}

\subsection{Appendix A --- Dependency and status table}\label{appendix-a-dependency-and-status-table}

\begin{longtable}[]{@{}
  >{\raggedright\arraybackslash}p{(\linewidth - 6\tabcolsep) * \real{0.2500}}
  >{\raggedright\arraybackslash}p{(\linewidth - 6\tabcolsep) * \real{0.2500}}
  >{\raggedright\arraybackslash}p{(\linewidth - 6\tabcolsep) * \real{0.2500}}
  >{\raggedright\arraybackslash}p{(\linewidth - 6\tabcolsep) * \real{0.2500}}@{}}
\toprule\noalign{}
\begin{minipage}[b]{\linewidth}\raggedright
Construction
\end{minipage} & \begin{minipage}[b]{\linewidth}\raggedright
Required inputs
\end{minipage} & \begin{minipage}[b]{\linewidth}\raggedright
Output
\end{minipage} & \begin{minipage}[b]{\linewidth}\raggedright
Present status
\end{minipage} \\
\midrule\noalign{}
\endhead
\bottomrule\noalign{}
\endlastfoot
Admitted cut & finite exhibition + Structure & typed apartness & Definition \\
Residual locus & comparison + failed repair & local residual witness & Definition / finite checks \\
Symmetry no-selection & full group action + equivariance & obstruction to invariant choice & Finite proposition \\
Separation price & declare-then-assert grammar + item cost + isomorphism & minimal responses and the surviving tie & Finite proposition \\
Retention & memory carrier + update + report + comparison & re-accessible mark & Finite deterministic quotient; broader relational classification open \\
Cochain quotient & graph + coefficient group + \(d_0\) & defect modulo repair & Definition; bounded propositions \\
Closure/cohomology & \(C^2\), \(d_1\), \(d_1d_0=0\) & cocycles/classes & Conditional; often absent \\
Evaluation preorder and window & comparison relation + nerve route or certified \((d_0,d_1)\) & degree-1 window & Definition \\
Obstruction signature & window + norm \(J\) + closure declaration & \((\Phi_1,\Gamma_2,R_{cl})\) and diagnostic class & Finite propositions; routing interface open \\
Exact certificate & weighted \(L^1\) + repair + bounded circulation & certified \(\Phi_1\) & Finite propositions (soundness; repairable mass and non-unique pairs); existence imported \\
Path discrepancy & two same-endpoint paths + weighted \(L^1\) & lower bound on \(\Phi_1\) & Finite proposition \\
Quotient observer & window + report map with \(Rd_0=0\) & presentation map + evaluator cost & Definition; passivity proved \\
Context witness & states + contexts + profile into \(Q^1\) & exposed residuals and witness rank & Definition; rank invariance proved \\
Patch cover of a cycle & oriented cycle + path patches + declared aggregation & independent, shared-maximum, and union readings of one class & Finite proposition \\
Degree-\(k\) and torsion windows & finite simplicial complex + coefficient field & \(Q^k\), \(H^k\), coefficient-dependent classes & Finite propositions (branch) \\
Projective phase class & finite group + declared unitary matrices & \([\omega]\in H^2(G;U(1))\) & Finite proposition (branch) \\
Finite commutation & matrices over a commutative ring; clock and shift over \(\mathbb Z[\zeta_N]\) & trace obstruction; Weyl class of order \(N\) & Finite propositions (branch) \\
Finite reconstruction frames & complex carrier + complete unbiased bases (\(d=3,4\)) or phase-point operators (odd \(N\)) & exact linear reconstruction of \(\rho\); signed valuation \(W_\rho\) & Finite propositions (branch) \\
Complex growth & larger complex, face attachment, or equal-weight subdivision & comparison of \(\Phi_1\) & Finite proposition (growth never lowers \(\Phi_1\)) \\
Presentation descent & surjective presentation + target & factor or obstruction set & Finite proposition \\
Chain contraction & chain maps + homotopy + convention & retraction defect; perturbed contraction & Definition; proposition; imported lemma \\
Restriction curvature & finite poset + declared comparable-pair maps & \(K\), with \(d_1d_0=K\) and \(d_2d_1=K\) on the top edge & Finite proposition \\
Truncation leakage & operator + orthogonal projector & leakage form \(L_P\) & Finite proposition \\
Linear recovery & presentation map + declared subspace + norms & recovery constant or kernel witness & Finite propositions \\
Non-simplex state space & extreme points + affine effects & two decompositions of the centre & Finite proposition (branch) \\
Sorkin sums & additive amplitudes + declared biadditive readout & vanishing \(I_{k\ge3}\) & Finite proposition (branch) \\
Succession & events + directed acyclic relation & strict partial order & Finite proposition after inputs \\
Duration & path system + additive valuation & path magnitude & Definition; exactness propositions \\
Clock & reference process + coupling + error rule & calibrated comparison & Definition; adequacy open \\
Finite probability & witness family + mass + partition & normalized valuation & Finite proposition; atomicity open \\
Observer & presentation + memory + evaluator + policy & profile-relative reports/status & Definition; adapters open \\
Observation & interaction + registration + witness + re-access + evaluation & verification record; validated observation when status is \texttt{Accept} & Definition \\
Measurement & accepted observation + measurand specification + result carrier, map, unit, and calibration + uncertainty and error semantics + optional metrological-traceability evidence + provenance/re-access & measurement result & Repaired WDP definition; calculus open \\
Authority & role + jurisdiction + permission + procedure & authorized act/status & Separate declared layer \\
History sum & path carrier + semiring weights & finite aggregate & Finite proposition at fixed length \\
Quantum carrier & complex/unitary imports & quantum-capable geometry & Conditional branch \\
Certified preserver (R12) & R8 involution, R9 positive definite form, R10 readout, global ray bijection and dimension hypothesis & unique projective semisimilarity; rational phase group; typed norm obstruction & Native written proof; selected Lean theorems; concrete submitted-witness checks \\
Observation stabilization (R12-O03) & named readouts, declared group/algebra, linear defect carrier; Appendix T for the count & Galois closure, dagger bicommutant, sharp linear observation count & Conditional algebra; no nonlinear or restricted-menu optimality inferred \\
Compatible generator (R13) & matrix algebra, \(K\)-linear derivation and declared adjoint law & \(\delta=-i[H_0,\cdot]\), unique modulo rational scalars & Written finite-dimensional theorem; Python basis-table checker; no supplied flow \\
Kochen--Specker parity & declared rays in \(\mathbb Q^4\) + orthogonal contexts & no \(\mathbb F_2\) valuation with sum 1 per context & Finite proposition (branch) \\
Classical package S & nonnegative weights + stochastic updates + classical instruments + joint runs & mixtures of function updates; instrument laws; CHSH bound; sharpness & Finite propositions (classical corner extended); no packet representation \\
Tolerant packet adapter & frequency tables + declared tolerance & table with certified signalling defect; exact repair on two contexts & Finite propositions and finite-cover/i.i.d. concentration scope; general non-i.i.d. sampling open \\
Composition principle CP & injective tagging + no-signalling + bilinear positivity (classical carriers) & composite states of classical carriers are joint distributions; PR box not excluded & Finite propositions; tensor product open \\
Amplitude bridge & integer cochain + rational phase \(u\) + two walks & gauge-class dependence of phase-blind readouts; no rational exponentiation & Finite propositions; readout open \\
Qubit instrument adapter & Kraus family over \(\mathbb Q(i)\) + mixture state & probabilities, post-state, POVM underdetermination & Finite propositions and finite-dimensional instrument laws (Appendix J); infinite-dimensional case open \\
Quantum reconstruction & operational axioms + comparison class & representation/no-go & Open audit \\
Relative completeness & finite grammar, evaluator and status rules; sound decoding and constructive coverage for solving claims & finite status coverage; original-goal witness and search preservation under declared adapters & Finite coverage and §§20.6--20.8 solving theorems; arbitrary profile fidelity remains open \\
Toric-code observer completion & finite carrier \(\Lambda = \mathbb{Z}_L \times \mathbb{Z}_L\) + stabilizers + schedule \(D\) + admissible errors \(\mathcal{E}\) + protected subgroup \(B\) & exact primal recovery table \(\rho\) or dual failure pair \((e_1, e_2)\) & Finite admitted-fiber criterion and corrected Chapter 23/Appendix U bounds; joint spacetime/readout matrices outside supplied family open \\
\end{longtable}

\begin{center}\rule{0.5\linewidth}{0.5pt}\end{center}

\subsection{Appendix B --- A compact end-to-end finite example}\label{appendix-b-a-compact-end-to-end-finite-example}

This example illustrates composition without claiming physical completeness.

Let a device have states \(S=\{s_0,s_1,s_2\}\) and transitions
\(s_0\to s_1\to s_2\). A binary contact changes from open at \(s_0\) to closed at
\(s_1\) and remains closed at \(s_2\).

\begin{enumerate}
\def\labelenumi{\arabic{enumi}.}
\tightlist
\item
  \textbf{Cut.} Open/closed roles are declared distinct.
\item
  \textbf{Residual witness.} At \(s_1\), a sensor record differs from the open
  baseline.
\item
  \textbf{Retention.} A latch stores the change bit through \(s_2\).
\item
  \textbf{Succession.} The transition graph is acyclic, so reachability is a strict
  order.
\item
  \textbf{Duration.} Edge weights \((2,3)\) give duration \(5\) from \(s_0\) to \(s_2\) in
  a declared unit.
\item
  \textbf{Observer presentation.} The observer exposes \texttt{(contact,\ latch)} and hides
  other device state.
\item
  \textbf{Observation.} Interaction, ADC registration, typed contact witness,
  retained record, and an evaluator range check that returns \texttt{Accept} form an
  validated observation.
\item
  \textbf{Measurement.} The measurand \(Q\) is terminal voltage under the declared
  loading condition; \(Y\) is the ordered carrier of closed voltage intervals.
  A calibration table maps ADC code into \(Y\), attaches volts, a separately
  typed worst-case error bound, and a calibration-version verification record. No
  metrological-traceability claim is made without an external calibration
  chain and uncertainty contributions.
\item
  \textbf{Probability.} Repeated, separately declared trials may be tallied and
  normalized; the resulting finite table has no stochastic interpretation
  until a trial model is added.
\item
  \textbf{History.} Positive-real edge weights can be multiplied and alternative
  path weights summed. With only one path, no interference question arises.
\item
  \textbf{Authority.} Whether the result permits an external action is determined
  by a separate jurisdictional rule.
\end{enumerate}

At each step, deleting the new input removes the corresponding conclusion. The
example is therefore a dependency demonstration, not a derivation from one bit
of difference.

\begin{center}\rule{0.5\linewidth}{0.5pt}\end{center}

\subsection{Appendix C --- Source integration map}\label{appendix-c-source-integration-map}

The canonical chapters consolidate the source folders as follows.

\begin{longtable}[]{@{}
  >{\raggedright\arraybackslash}p{(\linewidth - 2\tabcolsep) * \real{0.5000}}
  >{\raggedright\arraybackslash}p{(\linewidth - 2\tabcolsep) * \real{0.5000}}@{}}
\toprule\noalign{}
\begin{minipage}[b]{\linewidth}\raggedright
Source folder
\end{minipage} & \begin{minipage}[b]{\linewidth}\raggedright
Canonical destination
\end{minipage} \\
\midrule\noalign{}
\endhead
\bottomrule\noalign{}
\endlastfoot
\texttt{022\_ground\_state} & Chapters 1--2; spectral appendix separated into Chapter 5 scope \\
\texttt{001\_from\_primitive\_difference\_to\_cochain\_systems} & Chapters 3, 10--11, 16, and open programme \\
\texttt{000\_primitive\_difference\_for\_quantum\_carriers} & Chapters 4--5 \\
\texttt{003\_witnessed\_succession\_and\_duration} & Chapters 6--7 \\
\texttt{024\_duration} & Merged into Chapter 7; no parallel live duration theory \\
\texttt{023\_probability} & Chapters 8--9 \\
\texttt{004\_calculus\_of\_observers\_and\_observation} & Chapters 10--12 and 15 \\
\texttt{025\_measurement} & Repaired and expanded in Chapter 13 \\
\texttt{002\_self\_witnessing\_and\_the\_awareness\_boundary} & Chapter 14 with refuted ladder fenced \\
\texttt{005\_finite\_witness\_history\_integral} & Chapter 17, dependent on Chapter 16 \\
\texttt{006\_quantum\_reconstruction\_or\_classical\_no\_go} & Chapters 18--19 \\
\texttt{007\_witness\_relative\_cosmogenesis} & Retitled and narrowed into Chapter 20 \\
\texttt{../006\_primitive\_axiomatics} (v2.10.0) & §§4.7--4.12, §§5.10--5.14, §8.5, §10.5, §11.5, and \texttt{CH04-O01} in §21.8 \\
\texttt{../../000\_genesis\_alpha\_program/012\_primitive\_difference\_residual\_quotients\_and\_simplex\_growth} & CH04-P09 and CH04-N02, which replace its growth rule \\
\texttt{../../000\_genesis\_alpha\_program/015\_primitive\_projection\_obstruction\_descent\_and\_retraction\_calculus} & §§16.3 and 16.5--16.6; the §4.12 projection remark; the Chapter 1--5 technical note \\
\texttt{../../000\_genesis\_alpha\_program/003\_operational\_probabilistic\_state\_spaces} & CH18-P01, CH19-P01, and the §19.10 no-signalling sentence \\
\texttt{../../000\_genesis\_alpha\_program/035\_cost\_of\_closing\_observer\_blind\_spots} & CH14-P02, CH14-P03, and CH14-N02 \\
\texttt{../../000\_genesis\_alpha\_program/020\_finite\_hilbert\_lift\_adequacy} & CH05-P07 and CH16-P04 \\
\texttt{../../000\_genesis\_alpha\_program/007\_what\_the\_basis\_forces\_minima\_difference\_cost} & CH01-P02 and merged examples in §3.2 and §7.4 \\
\texttt{../../000\_genesis\_alpha\_program/004\_nonlocality\_contextuality\_finite\_commutator\_obstructions} & CH18-P06 (§18.4); CH05-P08 and CH05-P09 (§5.13); CH05-P10 (§5.15); the PR-box mixture sentence in §19.11 \\
\texttt{../../000\_genesis\_alpha\_program/005\_discrete\_phase\_spaces\_stabilizer\_grounded\_semantics} & CH05-P11 and CH05-P12 (§5.15) \\
\texttt{../../000\_genesis\_alpha\_program/006\_no\_free\_conclusions\_minimal\_calculus} & CH03-P03 and its worked count in §3.2; a §21.2 falsifier; an Appendix A row \\
\texttt{../../000\_genesis\_alpha\_program/008\_l1\_coboundary\_certificates\_and\_lift\_refusal} & CH04-P10 and CH04-P11 in §4.10; the restricted-repair remark after CH04-P09; sentences in §4.12 and §7.4; a §21.3 falsifier; the lift-admission paragraph after CH04-O01 in §21.8 \\
\texttt{../../000\_genesis\_alpha\_program/009\_representation\_lifts\_and\_coherence\_defects} & CH16-P05 in the new §16.7 with a pointer in §4.4; the diagonal-lift and density-lift remarks in §5.14; the Euclidean selections after CH04-P10; a §21.3 falsifier \\
\texttt{../../000\_genesis\_alpha\_program/010\_reasoning\_and\_intelligence\_as\_containers} & Not carried: an interpretive essay whose separations are already §10.2, CH12-P01, §§15.2--15.3, §20.2, §20.5, and CH00-N03 \\
\texttt{../../000\_genesis\_alpha\_program/011\_finite\_and\_infinite\_as\_observer\_relative\_regimes\_of\_obstruction} & CH11-P02 in §11.2 and CH10-N01 in §10.2; §21.7 item 21, a §21.12 migration, and an Appendix A row \\
\texttt{../../000\_genesis\_alpha\_program/022\_finite\_presentation\_calculus\_for\_computation} & No rows: its descent criterion is CH16-P02; the awaiting/exhausted split of \texttt{Undetermined} is merged into §21.9 \\
\end{longtable}

Merged-source papers under the observer folder remain provenance sources. They
do not become additional canonical chapters just because they contain useful
definitions or examples.

The primitive-axiomatics source is a sibling package in this series with its own
deposit record. It was absorbed on 2026-09-29 as a source input: its bytes are
bound in the \texttt{absorbed\_snapshots} block of \texttt{source\_provenance.json}, and its
claim IDs appear in the historical crosswalk of the ledger without inheriting
status. Its trace-instrumentation bridge, guarded repair policy, terminal and
lift enumerations, reflective completion, adjudication protocol, and
edge-symmetry gauge classification were not carried. Six corrections to the
source were made on absorption: the balance identity needs no closedness
hypothesis (§5.12); the Pauli product carries a factor \(-i\) (§5.13); the mod-2
class on the projective plane is represented by a single-triangle indicator,
not the all-ones cochain (§5.11); the source's brute-force \(\Phi_1\), which
searched a small box of integer potentials, and its Bland's-rule test were
replaced in the executable package; and a mis-cited identifier for the
trace-distance contraction reference was replaced (REF-QM-020). Its Lean
anchors for the exact sequence, ratio bound, passivity, and triangle were
restated over \(\mathbb Q\) and the literal quotient (Appendix D).

Six archived working papers were absorbed on the same date, each bound by
its own \texttt{absorbed\_snapshots} block. Audit found much of their stated content
false, packaged, or out of scope, and none of it enters here without a new
proof or a correction. Not carried: 012's claim that complex growth repairs
residuals (it cannot; CH04-P09), its spectral-sequence reading of \(\Phi_1\), and
its physics appendix; 015's two-category claim, routing alphabet, hash tape,
and superspace counterexample, and its perturbation formulas as printed (sign
error; CH16-C01 gives the corrected form); 003's cyclotomic census and
mechanization claims; 035's greedy optimality, its floating-point ``dual
witness,'' its \(\ell^1\) modulus bound (false), and its physics and audit readings;
020's claim that vanishing leakage is necessary for a genuine reduced zero
(false; CH16-P04) and its Born-rule reading; and 007's
framing that depends on an external derivation study, and its vendored source code, which the atlas rules exclude. Each paper's
claim IDs appear in the ledger's historical crosswalk without inherited status.

The remaining eight archived working papers (004, 005, 006, 008, 009, 010, 011,
and 022) were absorbed on 2026-09-30 under the same discipline, each bound by
its own \texttt{absorbed\_snapshots} block. Every surviving statement was corrected,
generalized, or reproved before it entered. Not carried: 004's rational CHSH
model (its correlators are not those of its vectors, and the model is a PR-box
mixture), its printed ray table (seven of nine contexts not orthogonal) and its
\(H^1(\mathbb Z_2)\) reading, its characteristic-zero form of the trace
obstruction (false for rings), its microcausality reading, and its \(1/\sqrt3\)
normalization; 005's stabilizer codes (correct, outside scope), its Heyting
semantics, empirical-reference definition, panel result, and firewall, its
truth-maintenance system and the source-withdrawn AGM claim, and its pin to an
external commit other than the required tag; 006's regime
framing that depends on an external derivation study, and its vendored source code, which the atlas rules exclude, its
number-theory catalogue (correct, but outside scope), and its two-fermion
exchange bridge (a physics reading whose cited Lean names are absent from its
package); 008's lift-admission and refusal calculus, net-utility and rollback
rules, and typed terminals (routing layer, CH04-O01), its LP complexity and
coordinate-descent material, its SAT and belief-revision readings, and its
claim that only a repair column equal to the residual lowers \(\Phi_1\) (false);
009's tie-break functional \(\Xi\) with its evaluator hierarchy and transition
gate, its quarantined Hilbert-lift conjecture, and its sheaf predictive-coding
note (circular where not withdrawn); 010 entirely, an interpretive essay whose
measurements and verification records refer to paths absent from this repository; 011's
finite/infinite vocabulary, which labels every exact class infinite, its
reading of the quotient class as a continuum, its Riemann-hypothesis
springboard, the quantum bound it cites, and its refusal flags, which are
policy booleans; and 022's machine proposal and open theorem queue, and its
productivity remark, which fails for runs from one stage of a deterministic
machine with finitely many events. No absorbed Lean or code was ported, and no
new row takes Lean or executable credit. Every absorbed archived working-paper directory
has now been removed from the tree; the bytes remain at the commits recorded in
\texttt{source\_provenance.json}.

\begin{center}\rule{0.5\linewidth}{0.5pt}\end{center}

\subsection{Appendix D --- Formalization status and executable evidence architecture}\label{appendix-d-formalization-status-and-executable-evidence-architecture}

This appendix records the evidence architecture and mechanization status of the
monograph. Selected finite discrete statements are formalized in Lean 4
(\texttt{lean4/WitnessedDifference/}) and backed by executable Python test suites
(\texttt{code/witnessed\_difference/}). This appendix documents the evidence binding
discipline, distinguishing proved formal theorems, executable tests, and open
interfaces without letting one artifact stand for another.

\subsubsection{D.1 Mechanized finite core}\label{d.1-mechanized-finite-core}

The source inventory at remote main \texttt{096551f4f98d9e801d95acfa84fe85deaf309b70},
counted on 2026-10-11, is \textbf{102 \texttt{.lean} files under \texttt{lean4/}}: the lakefile,
the historical root module, and \textbf{100 modules} in
\texttt{lean4/WitnessedDifference/}. This is a source-file count, excluding build
outputs and the separate reference sources under \texttt{code/}. It is not a count
of compiled modules, audited declarations, or fully formalized claims.

The dated evidence surfaces are distinct:

\begin{longtable}[]{@{}
  >{\raggedright\arraybackslash}p{(\linewidth - 6\tabcolsep) * \real{0.2308}}
  >{\raggedright\arraybackslash}p{(\linewidth - 6\tabcolsep) * \real{0.2308}}
  >{\raggedleft\arraybackslash}p{(\linewidth - 6\tabcolsep) * \real{0.3077}}
  >{\raggedright\arraybackslash}p{(\linewidth - 6\tabcolsep) * \real{0.2308}}@{}}
\toprule\noalign{}
\begin{minipage}[b]{\linewidth}\raggedright
Audit date
\end{minipage} & \begin{minipage}[b]{\linewidth}\raggedright
Recorded source scope
\end{minipage} & \begin{minipage}[b]{\linewidth}\raggedleft
Reports
\end{minipage} & \begin{minipage}[b]{\linewidth}\raggedright
Interpretation
\end{minipage} \\
\midrule\noalign{}
\endhead
\bottomrule\noalign{}
\endlastfoot
2026-09-30 & Historical Mathlib-backed root: then 29 source files and 341 theorem/lemma declarations & 202 & Historical root audit, retained in \texttt{lean4/VERIFICATION.md}; not a current full-package build. \\
2026-10-03 & Standalone core spine: 29 supplement modules & 408 & Extends the earlier 28-module, 393-report tensor checkpoint with \texttt{DyadicPresentation}. \\
2026-10-08 & Standalone core spine: 50 modules listed in \texttt{scripts/verify\_finite\_derivations.py} & 691 & Later audit of the expanded spine, with source-bound reports in \texttt{lean4/*\_axioms.txt}; supersedes the 29/408 scope for this entry point. \\
2026-10-09 & \texttt{GuardedSchedule}\allowbreak\texttt{Topology} & 44 & 37 theorems, 6 definitions and 1 inductive declaration; the arrangement homotopy theorem remains a written proof. \\
2026-10-09 & \texttt{GuardedChainTransport} and \texttt{SparseChainTransportBounds} & 143 & 78 and 65 explicit declarations, respectively, including 99 theorems; conditional certificate and arithmetic results, not a formalization of the entire cellular construction. \\
\end{longtable}

The last two entries have their own source hashes and scope exclusions in
\texttt{manifests/guarded\_schedule\_topology\_lean.json},
\texttt{manifests/guarded\_chain\_transport\_lean.json} and
\texttt{manifests/sparse\_chain\_transport\_lean.json}. They do not silently extend the
50-module entry point. The further module compilations recorded in the
2026-10-08 review overlap existing sets and must not be summed as disjoint
coverage. Dates above identify the original stored audits; their receipts bind
source revisions by SHA-256. The 2026-10-11 editorial pass additionally replayed
the unchanged 50-module, 691-declaration core spine. That replay does not renew
the unrelated audits or the paired canonical reproduction gate.

The standalone replay uses Lean 4.33.1 without Mathlib and checks only
\texttt{\{propext,\ Classical.choice,\ Quot.sound\}}, with no admitted proofs:
\texttt{scripts/verify\_finite\_derivations.py\ -\/-lean\ /path/to/lean-4.33.1/bin/lean}.
The default \texttt{lake\ build} and historical \texttt{AxiomReport.lean} have a separate
Mathlib-backed import graph, pinned in \texttt{lean4/lakefile.lean}. Neither source
counts nor selected declaration counts measure whole-row formalization.
For the canonical claims, the following exact bindings and the ledger's
separate formalization/reproduction fields state the relevant scope.

In accordance with the canonical claim-to-declaration map (\texttt{lean4/claim\_map.md}), the formal
package establishes exact bindings for twenty-nine canonical theorems and one definition:

\begin{itemize}
\tightlist
\item
  \textbf{Exact theorem bindings:}

  \begin{itemize}
  \tightlist
  \item
    CH03-P01: No point of a finite set with at least two elements is fixed by every permutation (\texttt{WitnessedDifference.Symmetry.no\_fully\_invariant\_point});
  \item
    CH07-P01: Exact path duration additivity under concatenation (\texttt{WitnessedDifference.FinitePath.duration\_append});
  \item
    CH07-P02: Path independence of exact duration potentials (\texttt{WitnessedDifference.FinitePath.exact\_duration\_path\_independent}, with telescoping lemma \texttt{duration\_potential});
  \item
    CH14-P01: Presentation blind spots under noninjective observation maps (\texttt{WitnessedDifference.presentation\_blind\_spot});
  \item
    CH16-P01: A map descends through a surjective presentation exactly when it is constant on fibers (\texttt{WitnessedDifference.Descent.function\_descends\_iff\_fiber\_constant});
  \item
    CH04-P07: Dual-certificate soundness over \(\mathbb Q\) on a finite edge list (\texttt{dual\_certificate\_lower\_bound}, with \texttt{matching\_certificate\_is\_optimal}, in \texttt{WitnessedDifference.Obstruction});
  \item
    CH04-P08: The filled-triangle signature over \(\mathbb Q\) (\texttt{WitnessedDifference.Obstruction.filled\_triangle\_signature}), whose lower bound is derived from the certificate lemma;
  \item
    CH10-P02: The kernel order on declared observer families (\texttt{declared\_family\_order}, in \texttt{WitnessedDifference.ObserverResolution});
  \item
    CH10-P03: The grammar resolution bound (\texttt{grammar\_resolution\_bound}, in the same namespace);
  \item
    CH01-P02: Pointed relations and apartness (\texttt{pointed\_relation\_independence\_fin} in \texttt{WitnessedDifference.PointedRelation});
  \item
    CH16-P02: Presentation descent and sections (\texttt{presentation\_descent\_signature} in \texttt{WitnessedDifference.Descent});
  \item
    CH18-P02: Nonnegative weights force permutations (\texttt{nonnegative\_weights\_force\_permutations} in \texttt{WitnessedDifference.ClassicalCorner});
  \item
    CH18-P03: Signed weights do not (\texttt{signed\_weights\_break\_the\_corner}, same namespace);
  \item
    CH18-P04: Unitaries over \(\mathbb{Z}[i]\) are monomial with entries \(\pm 1\) or \(\pm i\) and move the readout by the same permutation (\texttt{unitary\_package\_is\_classical}, same namespace);
  \item
    CH18-P05: Where the classical corner ends (\texttt{classical\_corner\_ends}, same namespace);
  \item
    CH19-P02: The packet adapter (\texttt{measurement\_packet\_adapter} in \texttt{WitnessedDifference.PacketAdapter});
  \item
    CH19-P03: What survives the adapter (\texttt{what\_survives\_the\_adapter}, same namespace);
  \item
    CH19-P04: Every rational model is an image (\texttt{every\_model\_is\_a\_packet\_image}, same namespace);
  \item
    CH19-P05: Contextuality as a missing global table (\texttt{contextuality\_is\_a\_missing\_global\_table}, same namespace), with the general equivalence \texttt{noncontextual\_iff\_global\_table} in \texttt{WitnessedDifference.Empirical};
  \item
    CH18-P07: Stochastic updates are mixtures of function updates (\texttt{stochastic\_updates\_are\_mixtures\_of\_functions} in \texttt{WitnessedDifference.ClassicalPackage});
  \item
    CH18-P08: Classical instruments and the sharpness of the package S (\texttt{classical\_package\_and\_sharpness}, same namespace);
  \item
    CH19-P06: The tolerant adapter (\texttt{tolerant\_adapter\_signature} in \texttt{WitnessedDifference.ToleranceAdapter});
  \item
    CH19-P07: Exact repair on two contexts (\texttt{two\_context\_tolerance\_iff}, same namespace, with \texttt{exists\_repair});
  \item
    CH19-P08: Injective tagging (\texttt{locallyRecoverable\_iff\_injective} and \texttt{dumpRight\_not\_recoverable} in \texttt{WitnessedDifference.Composition});
  \item
    CH19-P09: No-signalling of local instruments and observer compatibility (\texttt{noSignalling\_right}, \texttt{noSignalling\_left}, \texttt{marginal\_independent\_of\_choice}, \texttt{presentation\_commutes\_with\_marginal}, same namespace);
  \item
    CH19-P10: Classical composites are Cartesian, and the PR box on the product cover (\texttt{composite\_state\_iff\_joint\_distribution}, \texttt{joint\_unique}, \texttt{no\_signalling\_does\_not\_force\_the\_local\_composite}, same namespace);
  \item
    CH19-P11: Gauge descent of phase-blind readouts (\texttt{common\_phase\_iff}, \texttt{phaseBlind\_readout\_factors}, \texttt{holonomy\_gauge\_invariant}, \texttt{cross\_eq\_holonomy}, \texttt{readout\_depends\_on\_cochain\_class}, \texttt{gauge\_invariance\_does\_not\_select\_the\_readout} in \texttt{WitnessedDifference.AmplitudeBridge});
  \item
    CH19-P12: The exponentiation obstruction (\texttt{exponentiation\_obstruction}, with \texttt{phaseU\_no\_square\_root}, \texttt{no\_rational\_exponentiation} and \texttt{phaseU\_zpow\_ne\_one}, same namespace);
  \item
    CH19-P14: A POVM does not determine the instrument (\texttt{povm\_does\_not\_determine\_the\_instrument} in \texttt{WitnessedDifference.InstrumentAdapter}).
  \end{itemize}
\item
  \textbf{Exact definition binding:}

  \begin{itemize}
  \tightlist
  \item
    CH07-D01: Finite path structure and duration valuation (\texttt{WitnessedDifference.FinitePath}, \texttt{FinitePath.duration}).
  \end{itemize}
\end{itemize}

In addition to these thirty exact bindings, auxiliary modules mechanize focused finite milestones:

\begin{itemize}
\tightlist
\item
  constitutive order vs.~immediate succession separation on \(M_3\) posets (\texttt{Succession.lean});
\item
  transitivity of constitutive order, and irreflexivity and asymmetry under the \(M_3\) hypotheses (\texttt{Succession.lean}); the reachability criterion of CH06-P01 is prose only;
\item
  nonzero natural multiplicity normalization (\texttt{Probability.lean});
\item
  typed observation verification records, status refusal, and authority non-entailment (\texttt{Observer.lean});
\item
  nonempty calibration hierarchies and non-negative uncertainty fields (\texttt{Measurement.lean});
\item
  fixed-length semiring path-sum right-distributivity and coboundary holonomy (\texttt{SemiringTransport.lean});
\item
  one small finite witness for each of the eight external imports of §18.5, such as the absence of an integer square root of \(-1\) and of a Boolean square root of the toggle; they illustrate why the discrete host does not supply each import and are not independence proofs (\texttt{QuantumImports.lean});
\item
  well-definedness and set-level exactness of the closure map on the quotient by the repair relation for CH04-P03, the closure-ratio bound for \(\mathbb Q\)-valued costs for CH04-P06, the operator bound of item 4 of CH04-P05 from the infimum property, and a packaged form of observer passivity whose fiber lower bound is a hypothesis (\texttt{Obstruction.lean}; each narrower than its row);
\item
  the probabilities, the trace formula and the post-measurement state of the qubit adapter of CH19-P13, and their form as an empirical model on the one-measurement cover (\texttt{sum\_prob}, \texttt{prob\_eq\_effect}, \texttt{prob\_nonneg}, \texttt{post}, \texttt{instrumentModel} in \texttt{InstrumentAdapter.lean}); the application of CH19-P04 to that model and the replay of the verification records through the packet of CH13-D01 are prose, so CH19-P13 is narrower than its row;
\item
  restriction and one-edge subdivision for CH04-P09 from the infimum halves (\texttt{ComplexGrowth.lean}), the Sorkin identity for two and three alternatives under any biadditive readout for CH19-P01 (\texttt{Interference.lean}), and the commuting part of the retraction defect for CH16-P03 (\texttt{Descent.lean}), each narrower than its row.
\end{itemize}

The Lean anchors inherited from the primitive-axiomatics source were restated
rather than copied. The source's exact-sequence file proved kernel/range
exactness for an arbitrary function, its ratio bound used natural-number floor
division, its passivity file assumed the inequality it named, and its triangle
replay ranged over integer potentials; \texttt{lean4/PROVENANCE.md} records each
replacement.

The remaining 124 canonical claim rows have partial or absent whole-row Lean
bindings, including definitions and nonclaims, as recorded in
\texttt{lean4/claim\_map.md}. A prose result can be proved without being mechanized.
The nineteen historical targets now have bounded finite or conditional
resolutions in Appendices I--K; that status does not upgrade their partial Lean
bindings or resolve the five separate source-derivation obligations of Appendix L.

\subsubsection{D.2 Executable test suites and property testing}\label{d.2-executable-test-suites-and-property-testing}

The Python test suite across \texttt{code/unit\_tests/} and \texttt{code/tests/python/} (\texttt{integration\_tests/}, \texttt{injection\_tests/}, and \texttt{e2e\_tests/}) provides executable evidence across tests
covering:

\begin{itemize}
\tightlist
\item
  enumeration of the full symmetric groups on small sets, confirming the full-group case of the fixed-point criterion (no invariant point, oriented edge, or strict total order);
\item
  enumeration of small graphs testing exactness, loop periods, and duration statuses;
\item
  observer and measurement packets with deliberately missing fields asserting typed refusal;
\item
  verification that calibration-version updates preserve raw-record custody and provenance;
\item
  bounded path enumeration over exact rationals with cancellation reports,
  together with signed rational coarse-graining and length-one intertwining
  tests; the general-semiring implication has its separately scoped Lean proof;
\item
  classical permutation-isometry representation and quantum audit fixtures (P1--P12). The current \texttt{H015\_collapse\_no\_isometry} computes the supplied basis-map Gram matrix; \texttt{H015CollapseDilations} enumerates dilations on at most eight composite states; \texttt{H016\_no\_report\_half\_step} enumerates endomap roots on at most six report states; \texttt{H017Fin} checks a supplied Gaussian-integer linear instance; and \texttt{H020} verifies the full supplied eighteen-ray incidence and orthogonality witness before using parity. Malformed inputs and census limits refuse. The \href{code/unit_tests/test_quantum_input_evidence.py}{input mutation tests} establish these executable scopes, separately from the written carrier-wide results and historical Lean audits;
\item
  an exact two-phase simplex method with Bland's rule, whose every reported \(\Phi_1\) is
  re-checked by the dual-certificate test of CH04-P07, with tamper tests on each
  certificate field and Beale's cycling example started from its classical basis;
\item
  degree-\(k\) simplicial windows, the six-vertex projective-plane torsion witness,
  balance identities, the Klein-four projective phase class, quotient-observer
  passivity by exact linear programming, observer-family kernel order, grammar
  collisions, and witness-rank invariance;
\item
  exhaustive pointed-relation censuses, growth verification records for restriction, face attachment,
  and subdivision, exact congruence certificates for positive semidefiniteness and
  recovery constants, truncation leakage over \(\mathbb Q(i)\), presentation descent,
  chain contractions with the sign-corrected perturbation lemma re-validated on every
  output, the square state space, and Sorkin sums under declared readouts.
\end{itemize}

\subsubsection{D.3 Evidence separation}\label{d.3-evidence-separation}

For every formalized claim, distinct pointers are preserved for:

\begin{enumerate}
\def\labelenumi{\arabic{enumi}.}
\tightlist
\item
  prose statement;
\item
  claim-ledger row;
\item
  implementation;
\item
  tests;
\item
  formal theorem;
\item
  axiom/import report;
\item
  source revision and hash;
\item
  release artifact, if any.
\end{enumerate}

A missing pointer is \texttt{unknown} or \texttt{not\ present}, not green by inheritance.

\begin{center}\rule{0.5\linewidth}{0.5pt}\end{center}

\subsection{Appendix E --- Glossary}\label{appendix-e-glossary}

\textbf{Notation K.} Appendix T uses \(K=\ker Q\); CH01-D01 uses \(\mathcal K\) for the admitted comparison family; Chapter 16 uses \(K_{c,b,a}=r_{b,a}r_{c,b}-r_{c,a}\) for the restriction-composition defect.

\textbf{Admitted cut.} A finite exhibition of difference together with the structure
needed to type and compare it.

\textbf{Asymmetry price.} The declared input whose removal restores a symmetry that
moves the construction's output.

\textbf{Authority.} A role/jurisdiction/permission relation external to observer and
measurement structure.

\textbf{Calibration.} A versioned relation between reference quantity values and
their uncertainties and instrument indications and their uncertainties,
including the conditions and use under which a result and its uncertainty are
obtained.

\textbf{Closure defect.} A nonzero value of a proposed next differential, or the
failure of a differential composition to vanish.

\textbf{Descent.} The existence of a well-defined global or quotient-level object
compatible with declared local/restriction data.

\textbf{Duration gauge.} An additive valuation on composable finite intervals or
edges.

\textbf{Observation verification record.} A record of an interaction, registration, typed
witness, re-comparison state, and evaluator result; a refusal or rejection
verification record is not an accepted observation.

\textbf{validated observation.} An observation verification record whose first four ladder predicates
verify and whose evaluator result is \texttt{Accept}.

\textbf{Evaluator.} A declared defect/repair/cost/status mechanism. It does not
propose truth or confer authority.

\textbf{Measurement.} In the WDP interface, a validated observation extended by a
measurand specification, a distinct result/value carrier and value map,
scale/unit, calibration, separate uncertainty and error semantics, any claimed
metrological-traceability evidence, and separate provenance/re-access. The
validated observation prerequisite is WDP-specific.

\textbf{Metrological traceability.} A property of a measurement result supported by
a documented unbroken calibration chain to a stated reference, each calibration
contributing to measurement uncertainty; it is not ordinary data provenance.

\textbf{Observer.} A layered presentation, memory, evaluator, and policy with typed
adapters.

\textbf{Primitive Difference.} The admitted role of apartness in a finite
exhibition, not a complete structure generator.

\textbf{Profile.} The contexts, grammar, inputs, comparisons, resources, and
terminal rules under which an observer operates.

\textbf{Residual.} The portion of a declared target not removed by admissible repair.

\textbf{Retention.} A mechanism supporting later re-access and comparison of a
typed mark.

\textbf{Provenance/re-access.} A replayable chain linking a result to its raw record,
procedure, calibration version, evaluator, and transformations. It does not by
itself establish metrological traceability.

\textbf{Typed witness.} Evidence of a declared kind for a declared proposition;
never the unqualified source of observer, measurement, or authority structure.

\begin{center}\rule{0.5\linewidth}{0.5pt}\end{center}

\subsection{Appendix F --- Bibliography and citation policy}\label{appendix-f-bibliography-and-citation-policy}

All external references are registered in \href{docs/references.md}{references.md}. A
registry entry records the support role and affected chapters. An entry marked
\texttt{UNCONSULTED} or \texttt{UNRESOLVED} may identify a research target but may not support
a theorem statement. Secondary expositions can orient the reader, but claims
about an original theorem should be anchored to a primary source when one is
available.

Internal source manuscripts are provenance inputs, not external validation.
Their release labels, claim counts, and artifact references require independent
custody checks. The unified claim ledger, not this narrative, is the definitive
location for evidence status.

\subsubsection{F.1 Bibliography}\label{f.1-bibliography}

The following records resolve every \texttt{REF-*} citation in the main text.
They reproduce bibliographic records already held in the supporting registry;
inclusion does not extend the recorded consultation or theorem-support scope.
Panov and Tril {[}REF-MATH-013{]}, Construction 4.5 and Theorem 4.6, supply the
prior task-complex/arrangement connection used in Appendix P.22.

\textbf{{[}REF-AUT-002{]}} A. Nerode, ``Linear Automaton Transformations,'' \emph{Proceedings of the American Mathematical Society} 9 (4), 541--544 (1958). \href{https://doi.org/10.1090/S0002-9939-1958-0135681-9}{DOI 10.1090/S0002-9939-1958-0135681-9}.

\textbf{{[}REF-AUT-004{]}} J. Barwise and J. Seligman, \emph{Information Flow: The Logic of Distributed Systems}, Cambridge University Press (1997). \href{https://doi.org/10.1017/CBO9780511895968}{DOI 10.1017/CBO9780511895968}.

\textbf{{[}REF-AUT-005{]}} D. Park, ``Concurrency and Automata on Infinite Sequences,'' in \emph{Theoretical Computer Science}, LNCS 104, 167--183 (1981), \href{https://doi.org/10.1007/BFb0017309}{DOI 10.1007/BFb0017309}; R. Milner, \emph{Communication and Concurrency}, Prentice Hall (1989), ISBN 0-13-115007-3.

\textbf{{[}REF-AUT-007{]}} R. P. Feynman, ``Space-Time Approach to Non-Relativistic Quantum Mechanics,'' \emph{Reviews of Modern Physics} 20, 367--387 (1948). \href{https://doi.org/10.1103/RevModPhys.20.367}{DOI 10.1103/RevModPhys.20.367}.

\textbf{{[}REF-CAUSE-001{]}} J. Pearl, \emph{Causality: Models, Reasoning, and Inference}, 2nd ed., Cambridge University Press, 2009. ISBN 9780521895606.

\textbf{{[}REF-CAUSE-002{]}} E. Bareinboim, J. D. Correa, D. Ibeling and T. Icard, ``On Pearl's hierarchy and the foundations of causal inference,'' in \emph{Probabilistic and Causal Inference: The Works of Judea Pearl}, ACM Books, 2022, pp.~507--556.

\textbf{{[}REF-HOM-001{]}} C. A. Weibel, \emph{An Introduction to Homological Algebra}, Cambridge University Press (1994), \href{https://doi.org/10.1017/CBO9781139644136}{DOI 10.1017/CBO9781139644136}.

\textbf{{[}REF-HOM-002{]}} L. Positselski, \emph{Two Kinds of Derived Categories, Koszul Duality, and Comodule-Contramodule Correspondence}, Memoirs of the American Mathematical Society 212 (996) (2011), \href{https://doi.org/10.1090/S0065-9266-2010-00631-8}{DOI 10.1090/S0065-9266-2010-00631-8}; \href{https://arxiv.org/abs/0905.2621}{arXiv:0905.2621}.

\textbf{{[}REF-HOM-003{]}} B. Keller, ``Introduction to \(A_\infty\) Algebras and Modules,'' \emph{Homology, Homotopy and Applications} 3 (1), 1--35 (2001), \href{https://doi.org/10.4310/HHA.2001.v3.n1.a1}{DOI 10.4310/HHA.2001.v3.n1.a1}; \href{https://arxiv.org/abs/math/9910179}{arXiv:math/9910179}.

\textbf{{[}REF-HOM-004{]}} R. M. Vogt, ``Homotopy Limits and Colimits,'' \emph{Mathematische Zeitschrift} 134, 11--52 (1973), \href{https://doi.org/10.1007/BF01219090}{DOI 10.1007/BF01219090}; J.-M. Cordier and T. Porter, ``Homotopy Coherent Category Theory,'' \emph{Transactions of the American Mathematical Society} 349 (1), 1--54 (1997), \href{https://doi.org/10.1090/S0002-9947-97-01752-2}{DOI 10.1090/S0002-9947-97-01752-2}.

\textbf{{[}REF-HOM-005{]}} R. Zwanzig, ``Memory Effects in Irreversible Thermodynamics,'' \emph{Physical Review} 124, 983--992 (1961), \href{https://doi.org/10.1103/PhysRev.124.983}{DOI 10.1103/PhysRev.124.983}; H. Mori, ``Transport, Collective Motion, and Brownian Motion,'' \emph{Progress of Theoretical Physics} 33, 423--455 (1965), \href{https://doi.org/10.1143/PTP.33.423}{DOI 10.1143/PTP.33.423}; A. J. Chorin, O. H. Hald, and R. Kupferman, ``Optimal Prediction with Memory,'' \emph{Physica D} 166, 239--257 (2002), \href{https://doi.org/10.1016/S0167-2789(02)00446-3}{DOI 10.1016/S0167-2789(02)00446-3}.

\textbf{{[}REF-HOM-007{]}} V. K. A. M. Gugenheim, ``On the chain-complex of a fibration,'' \emph{Illinois Journal of Mathematics} 16(3) (1972), \href{https://doi.org/10.1215/ijm/1256065766}{DOI 10.1215/ijm/1256065766}; D. W. Barnes and L. A. Lambe, ``A fixed point approach to homological perturbation theory,'' \emph{Proceedings of the American Mathematical Society} 112(3), 881--892 (1991), \href{https://doi.org/10.1090/S0002-9939-1991-1057939-0}{DOI 10.1090/S0002-9939-1991-1057939-0}, with its correction, \emph{Proc. AMS} 129(3), 941 (2001; online 2000), \href{https://doi.org/10.1090/S0002-9939-00-06018-4}{DOI 10.1090/S0002-9939-00-06018-4}.

\textbf{{[}REF-MATH-001{]}} J. Dodziuk, ``Finite-Difference Approach to the Hodge Theory of Harmonic Forms,'' \emph{American Journal of Mathematics} 98 (1), 79--104 (1976). \href{https://doi.org/10.2307/2373615}{DOI 10.2307/2373615}.

\textbf{{[}REF-MATH-002{]}} A. Hatcher, \emph{Algebraic Topology}, Cambridge University Press (2002). \href{https://pi.math.cornell.edu/~hatcher/AT/AT.pdf}{Author-hosted text}.

\textbf{{[}REF-MATH-004{]}} S. Mac Lane, \emph{Categories for the Working Mathematician}, 2nd ed., Springer (1998). \href{https://doi.org/10.1007/978-1-4757-4721-8}{DOI 10.1007/978-1-4757-4721-8}.

\textbf{{[}REF-MATH-005{]}} R. T. Rockafellar, \emph{Convex Analysis}, Princeton University Press (1970). \href{https://press.princeton.edu/books/paperback/9780691015866/convex-analysis}{Publisher record}.

\textbf{{[}REF-MATH-006{]}} J. Aczél, \emph{Lectures on Functional Equations and Their Applications}, Mathematics in Science and Engineering 19, Academic Press (1966), 510 pp.~\href{https://books.google.com/books?id=1Va4QgAACAAJ}{Google Books record}.

\textbf{{[}REF-MATH-010{]}} R. G. Bland, ``New finite pivoting rules for the simplex method,'' \emph{Mathematics of Operations Research} 2(2), 103--107 (1977). \href{https://doi.org/10.1287/moor.2.2.103}{DOI 10.1287/moor.2.2.103}.

\textbf{{[}REF-MATH-011{]}} K. S. Brown, \emph{Cohomology of Groups}, Graduate Texts in Mathematics 87, Springer (1982). \href{https://doi.org/10.1007/978-1-4684-9327-6}{DOI 10.1007/978-1-4684-9327-6}.

\textbf{{[}REF-MATH-013{]}} Taras Panov and Vsevolod Tril, ``Permutohedral complex and complements of diagonal subspace arrangements'' (2025). \href{https://arxiv.org/abs/2504.04183v1}{arXiv:2504.04183v1}.

\textbf{{[}REF-MET-001{]}} Joint Committee for Guides in Metrology, \emph{International Vocabulary of Metrology---Basic and General Concepts and Associated Terms}, 3rd ed., JCGM 200:2012. \href{https://doi.org/10.59161/JCGM200-2012}{DOI 10.59161/JCGM200-2012}.

\textbf{{[}REF-MET-002{]}} Joint Committee for Guides in Metrology, \emph{Evaluation of Measurement Data---Guide to the Expression of Uncertainty in Measurement}, JCGM 100:2008, \href{https://doi.org/10.59161/JCGM100-2008E}{DOI 10.59161/JCGM100-2008E}, with Amendment 1:2026, \href{https://doi.org/10.59161/PPDI3267}{DOI 10.59161/PPDI3267}.

\textbf{{[}REF-MET-003{]}} Joint Committee for Guides in Metrology, \emph{Guide to the Expression of Uncertainty in Measurement---Part 1: Introduction}, JCGM GUM-1:2023. \href{https://doi.org/10.59161/JCGMGUM-1-2023}{DOI 10.59161/JCGMGUM-1-2023}.

\textbf{{[}REF-MET-006{]}} E. B. Davies and J. T. Lewis, ``An Operational Approach to Quantum Probability,'' \emph{Communications in Mathematical Physics} 17, 239--260 (1970). \href{https://doi.org/10.1007/BF01647093}{DOI 10.1007/BF01647093}.

\textbf{{[}REF-MET-007{]}} M. Ozawa, ``Quantum Measuring Processes of Continuous Observables,'' \emph{Journal of Mathematical Physics} 25, 79--87 (1984). \href{https://doi.org/10.1063/1.526000}{DOI 10.1063/1.526000}.

\textbf{{[}REF-NORM-001{]}} Keith Conrad, \href{https://kconrad.math.uconn.edu/blurbs/ugradnumthy/Zinotes.pdf}{\emph{The Gaussian Integers}}, Theorems 9.6 and 9.10, pp.~32--34.

\textbf{{[}REF-OBS-001{]}} H. Ollivier, D. Poulin, and W. H. Zurek, ``Objective Properties from Subjective Quantum States: Environment as a Witness,'' \emph{Physical Review Letters} 93, 220401 (2004). \href{https://arxiv.org/abs/quant-ph/0307229}{arXiv:quant-ph/0307229}; \href{https://doi.org/10.1103/PhysRevLett.93.220401}{DOI 10.1103/PhysRevLett.93.220401}.

\textbf{{[}REF-OBS-002{]}} H. Ollivier, D. Poulin, and W. H. Zurek, ``Environment as a Witness: Selective Proliferation of Information and Emergence of Objectivity in a Quantum Universe,'' \emph{Physical Review A} 72, 042113 (2005). \href{https://arxiv.org/abs/quant-ph/0408125}{arXiv:quant-ph/0408125}; \href{https://doi.org/10.1103/PhysRevA.72.042113}{DOI 10.1103/PhysRevA.72.042113}.

\textbf{{[}REF-OBS-004{]}} M. Oizumi, L. Albantakis, and G. Tononi, ``From the Phenomenology to the Mechanisms of Consciousness: Integrated Information Theory 3.0,'' \emph{PLOS Computational Biology} 10 (5), e1003588 (2014). \href{https://doi.org/10.1371/journal.pcbi.1003588}{DOI 10.1371/journal.pcbi.1003588}.

\textbf{{[}REF-OBS-005{]}} D. M. Rosenthal, \emph{Consciousness and Mind}, Oxford University Press (2005), \href{https://doi.org/10.1093/oso/9780198236979.001.0001}{DOI 10.1093/oso/9780198236979.001.0001}.

\textbf{{[}REF-OBS-006{]}} U. Kriegel, \emph{Subjective Consciousness: A Self-Representational Theory}, Oxford University Press (2009), \href{https://doi.org/10.1093/acprof:oso/9780199570355.001.0001}{DOI 10.1093/acprof:oso/9780199570355.001.0001}.

\textbf{{[}REF-OBS-007{]}} B. J. Baars, ``Global Workspace Theory of Consciousness: Toward a Cognitive Neuroscience of Human Experience,'' \emph{Progress in Brain Research} 150, 45--53 (2005), \href{https://doi.org/10.1016/S0079-6123(05)50004-9}{DOI 10.1016/S0079-6123(05)50004-9}; S. Dehaene and J.-P. Changeux, ``Experimental and Theoretical Approaches to Conscious Processing,'' \emph{Neuron} 70, 200--227 (2011), \href{https://doi.org/10.1016/j.neuron.2011.03.018}{DOI 10.1016/j.neuron.2011.03.018}.

\textbf{{[}REF-OBS-008{]}} D. Chalmers, ``The Meta-Problem of Consciousness'' (2018), as named in the source corpus; \emph{Journal of Consciousness Studies} 25 (9--10), 6--61 (2018). \href{http://hdl.handle.net/1885/310975}{ANU repository record}.

\textbf{{[}REF-OBS-009{]}} D. Frauchiger and R. Renner, ``Quantum Theory Cannot Consistently Describe the Use of Itself,'' \emph{Nature Communications} 9, 3711 (2018). \href{https://arxiv.org/abs/1604.07422}{arXiv:1604.07422}; \href{https://doi.org/10.1038/s41467-018-05739-8}{DOI 10.1038/s41467-018-05739-8}.

\textbf{{[}REF-PRES-001{]}} Leiba Rodman and Peter Šemrl, ``Orthogonality preserving bijective maps on finite dimensional projective spaces over division rings,'' \emph{Linear and Multilinear Algebra} 56(6), 647--664 (2008), \href{https://www.tandfonline.com/doi/abs/10.1080/03081080701395772}{DOI 10.1080/03081080701395772}.

\textbf{{[}REF-PROB-002{]}} A. M. Gleason, ``Measures on the Closed Subspaces of a Hilbert Space,'' \emph{Journal of Mathematics and Mechanics} 6, 885--893 (1957). \href{https://doi.org/10.1512/iumj.1957.6.56050}{DOI 10.1512/iumj.1957.6.56050}.

\textbf{{[}REF-PROB-005{]}} B. de Finetti, ``La prévision: ses lois logiques, ses sources subjectives,'' \emph{Annales de l'Institut Henri Poincaré} 7 (1), 1--68 (1937). \href{https://www.numdam.org/item/AIHP_1937__7_1_1_0/}{Official Numdam record and scan}.

\textbf{{[}REF-PROB-006{]}} P.-S. Laplace, \emph{Théorie analytique des probabilités} (1812), as a historical indifference source; E. T. Jaynes, ``Information Theory and Statistical Mechanics,'' \emph{Physical Review} 106, 620--630 (1957), \href{https://doi.org/10.1103/PhysRev.106.620}{DOI 10.1103/PhysRev.106.620}.

\textbf{{[}REF-QM-001{]}} S. Abramsky and A. Brandenburger, ``The Sheaf-Theoretic Structure of Non-Locality and Contextuality,'' \emph{New Journal of Physics} 13, 113036 (2011). \href{https://arxiv.org/abs/1102.0264}{arXiv:1102.0264}; \href{https://doi.org/10.1088/1367-2630/13/11/113036}{DOI 10.1088/1367-2630/13/11/113036}.

\textbf{{[}REF-QM-003{]}} L. Hardy, ``Quantum Theory From Five Reasonable Axioms'' (2001). \href{https://arxiv.org/abs/quant-ph/0101012}{arXiv:quant-ph/0101012}.

\textbf{{[}REF-QM-004{]}} G. Chiribella, G. M. D'Ariano, and P. Perinotti, ``Informational Derivation of Quantum Theory,'' \emph{Physical Review A} 84, 012311 (2011). \href{https://arxiv.org/abs/1011.6451}{arXiv:1011.6451}; \href{https://doi.org/10.1103/PhysRevA.84.012311}{DOI 10.1103/PhysRevA.84.012311}.

\textbf{{[}REF-QM-005{]}} J. Barrett, ``Information Processing in Generalized Probabilistic Theories,'' \emph{Physical Review A} 75, 032304 (2007). \href{https://doi.org/10.1103/PhysRevA.75.032304}{DOI 10.1103/PhysRevA.75.032304}.

\textbf{{[}REF-QM-006{]}} L. Masanes and M. P. Müller, ``A Derivation of Quantum Theory from Physical Requirements,'' \emph{New Journal of Physics} 13, 063001 (2011). \href{https://arxiv.org/abs/1004.1483}{arXiv:1004.1483}; \href{https://doi.org/10.1088/1367-2630/13/6/063001}{DOI 10.1088/1367-2630/13/6/063001}.

\textbf{{[}REF-QM-008{]}} S. Kochen and E. P. Specker, ``The Problem of Hidden Variables in Quantum Mechanics,'' \emph{Journal of Mathematics and Mechanics} 17, 59--87 (official record 1968; commonly cited 1967). \href{https://doi.org/10.1512/iumj.1968.17.17004}{DOI 10.1512/iumj.1968.17.17004}.

\textbf{{[}REF-QM-009{]}} S. Popescu and D. Rohrlich, ``Quantum Nonlocality as an Axiom,'' \emph{Foundations of Physics} 24, 379--385 (1994). \href{https://doi.org/10.1007/BF02058098}{DOI 10.1007/BF02058098}.

\textbf{{[}REF-QM-010{]}} B. S. Cirel'son, ``Quantum Generalizations of Bell's Inequality,'' \emph{Letters in Mathematical Physics} 4, 93--100 (1980). \href{https://doi.org/10.1007/BF00417500}{DOI 10.1007/BF00417500}.

\textbf{{[}REF-QM-014{]}} A. Cabello, J. M. Estebaranz, and G. García-Alcaine, ``Bell-Kochen-Specker theorem: A proof with 18 vectors,'' \emph{Physics Letters A} 212, 183--188 (1996); \href{https://arxiv.org/abs/quant-ph/9706009}{arXiv:quant-ph/9706009}.

\textbf{{[}REF-QM-015{]}} N. D. Mermin, ``Simple unified form for the major no-hidden-variables theorems,'' \emph{Physical Review Letters} 65, 3373--3376 (1990); A. Peres, ``Incompatible results of quantum measurements,'' \emph{Physics Letters A} 151, 107--108 (1990).

\textbf{{[}REF-QM-016{]}} W. F. Stinespring, ``Positive functions on C\emph{-algebras,'' }Proceedings of the American Mathematical Society* 6, 211--216 (1955); see also finite-dimensional environment bounds (e.g.~arXiv:2301.05488v3).

\textbf{{[}REF-QM-017{]}} A. M. Gleason, ``Measures on the Closed Subspaces of a Hilbert Space,'' \emph{Journal of Mathematics and Mechanics} 6, 885--893 (1957).

\textbf{{[}REF-QM-018{]}} R. D. Sorkin, ``Quantum mechanics as quantum measure theory,'' \emph{Modern Physics Letters A} 9, 3119--3127 (1994); \href{https://arxiv.org/abs/gr-qc/9401003}{arXiv:gr-qc/9401003}.

\textbf{{[}REF-QM-019{]}} M. H. Stone, ``On one-parameter unitary groups in Hilbert Space,'' \emph{Annals of Mathematics} 33, 643--648 (1932).

\textbf{{[}REF-QM-020{]}} M. B. Ruskai, ``Beyond strong subadditivity? Improved bounds on the contraction of generalized relative entropy,'' \emph{Reviews in Mathematical Physics} 6(5a), 1147--1161 (1994). \href{https://doi.org/10.1142/S0129055X94000407}{DOI 10.1142/S0129055X94000407}.

\textbf{{[}REF-QM-021{]}} J. Watrous, \emph{The Theory of Quantum Information}, Cambridge University Press (2018). \href{https://doi.org/10.1017/9781316848142}{DOI 10.1017/9781316848142}.

\textbf{{[}REF-REL-001{]}} A. Einstein, ``Die Grundlage der allgemeinen Relativitätstheorie,'' \emph{Annalen der Physik} 49(7), 1916, 769--822.

\textbf{{[}REF-REL-002{]}} E. Kretschmann, ``Über den physikalischen Sinn der Relativitätspostulate, A. Einsteins neue und seine ursprüngliche Relativitätstheorie,'' \emph{Annalen der Physik} 53, 1917, 575--614.

\textbf{{[}REF-REL-003{]}} A. Komar, ``Construction of a complete set of independent observables in the general theory of relativity,'' \emph{Physical Review} 111, 1958, 1182. \href{https://doi.org/10.1103/PhysRev.111.1182}{DOI 10.1103/PhysRev.111.1182}.

\textbf{{[}REF-REL-004{]}} C. Rovelli, ``What is observable in classical and quantum gravity?'' \emph{Classical and Quantum Gravity} 8(2), 1991, 297--316. \href{https://doi.org/10.1088/0264-9381/8/2/011}{DOI 10.1088/0264-9381/8/2/011}.

\textbf{{[}REF-REL-005{]}} J. T. François and L. Ravera, ``Mechanics as a general-relativistic gauge field theory, and relational quantization,'' arXiv:2510.19845 {[}physics.gen-ph{]}, 2025. \href{https://arxiv.org/abs/2510.19845}{arXiv}.

\textbf{{[}REF-TIME-001{]}} L. Bombelli, J. Lee, D. Meyer, and R. D. Sorkin, ``Space-Time as a Causal Set,'' \emph{Physical Review Letters} 59, 521--524 (1987). \href{https://doi.org/10.1103/PhysRevLett.59.521}{DOI 10.1103/PhysRevLett.59.521}.

\textbf{{[}REF-TIME-002{]}} J. Henson, ``The Causal Set Approach to Quantum Gravity'' (2006). \href{https://arxiv.org/abs/gr-qc/0601121}{arXiv:gr-qc/0601121}.

\begin{center}\rule{0.5\linewidth}{0.5pt}\end{center}

\subsection{Appendix G --- Finite decision procedures and proof obligations}\label{appendix-g-finite-decision-procedures-and-proof-obligations}

The programme repeatedly appeals to finiteness. This appendix states what that
finiteness actually buys. A finite carrier permits exhaustive algorithms for
many local questions, but decidability of an encoded instance is not the same
as a conceptual derivation, an efficient algorithm, or a theorem about an
unbounded family.

\subsection{G.1 Symmetry audit}\label{g.1-symmetry-audit}

For a finite input structure \(A\) and finite target set \(Y\):

\begin{enumerate}
\def\labelenumi{\arabic{enumi}.}
\tightlist
\item
  compute or enumerate the automorphism group \(\operatorname{Aut}(A)\);
\item
  compute its induced action on \(Y\);
\item
  intersect the fixed sets of the generators;
\item
  report the global fixed locus \(Y^{\operatorname{Aut}(A)}\);
\item
  if a construction returns \(y\notin Y^{\operatorname{Aut}(A)}\), identify the
  input field or implementation convention that broke symmetry.
\end{enumerate}

If the fixed locus is empty, invariant point selection is unrealizable. If it has
one element, that element is unique relative to the stated action. If it has
several elements, symmetry alone underdetermines the choice. This procedure
does not decide which input structure is metaphysically correct; it audits one
finite presentation.

\subsection{G.2 Exactness and certificate audit}\label{g.2-exactness-and-certificate-audit}

Choose bases for finite-dimensional \(C^0\) and \(C^1\), represent \(d_0\) by a
matrix \(D\), and represent the defect by a vector \(\delta\). Exactness is the
linear feasibility problem

\[
D f=\delta.
\]

Gaussian elimination over an exact field can either return \(f\) or show
inconsistency. A separating row vector \(y\) satisfying

\[
yD=0,\qquad y\delta\ne0
\]

is a dual nonmembership certificate. Over integers, modular rings, or floating-
point approximations, the algorithm and certificate semantics differ; the
coefficient domain must be recorded.

For weighted \(L^1\) residual minimization, introduce nonnegative slack variables
and solve a finite linear programme. A primal repair and dual feasible
circulation with matching objective values certify optimality under exact
arithmetic. A floating-point solver output is evidence only after tolerances,
rounding, and certificate checking are specified.

\subsection{G.3 Duration classifier}\label{g.3-duration-classifier}

Given a finite directed graph and exact rational edge valuation:

\begin{enumerate}
\def\labelenumi{\arabic{enumi}.}
\tightlist
\item
  test the declared \(2\)-cell closure equations, if a valid \(d_1\) exists;
\item
  solve \(d_0t=\lambda\);
\item
  if solvable, return \texttt{Exact} with potential \(t\);
\item
  if not solvable, seek a cycle-period or dual certificate;
\item
  return \texttt{ClosedNonExact} only when closure and nonexactness are both
  certified;
\item
  return \texttt{ActiveClosureFailure} with an exhibited cell when closure fails;
\item
  return \texttt{ExactnessUndetermined} when the structural or arithmetic inputs do
  not justify one of the prior statuses.
\end{enumerate}

The four statuses are evidence-bearing outputs, not labels inferred from solver
failure alone. A timeout or missing \(2\)-cell structure yields \texttt{Undetermined},
not \texttt{ClosedNonExact}.

\subsection{G.4 Observer-packet validator}\label{g.4-observer-packet-validator}

A finite schema validator can check whether an observer packet contains the
required fields and whether table-valued maps are total on their declared
domains. Deeper checks include:

\begin{itemize}
\tightlist
\item
  presentation commutation under a proposed adapter;
\item
  update/report consistency for every finite memory state and input;
\item
  representative independence of quotient evaluation;
\item
  soundness of terminal statuses on an exhaustive fixture family;
\item
  preservation of raw-record and calibration identifiers through policy output.
\end{itemize}

Exhaustive success establishes only the finite instance or bounded family actually
enumerated. A generic theorem requires symbolic proof or a verified reduction
from arbitrary instances to the enumeration.

\subsection{G.5 Observation and measurement validator}\label{g.5-observation-and-measurement-validator}

For a candidate measurement verification record, validate in dependency order:

\begin{enumerate}
\def\labelenumi{\arabic{enumi}.}
\tightlist
\item
  interaction and registration identifiers are present;
\item
  the raw record hash matches the stored bytes;
\item
  the witness semantics names a target proposition and context;
\item
  re-access recovers the same record and profile version;
\item
  the first four observation-ladder predicates reproduce, the evaluator inputs
  and result reproduce, and the result is exactly \texttt{Accept}; otherwise the
  object is an attempt verification record rather than the validated observation required
  by a measurement packet;
\item
  the measurand specification is distinct from the result carrier, and the
  indication and result types match the value map;
\item
  unit and scale are valid for the reported value type;
\item
  calibration version covers the indication range;
\item
  uncertainty and error/enclosure procedures are separately named and their
  inputs are present;
\item
  the provenance/re-access chain has no missing internal link;
\item
  when metrological traceability is claimed, the documented unbroken
  calibration hierarchy, stated reference, and every calibration's uncertainty
  contribution verify;
\item
  any authority decision is stored separately with its role and procedure.
\end{enumerate}

The validator should return a structured list of failures. Collapsing missing
calibration, rejected range, and absent authority into one \texttt{false} value makes
repair and audit harder.

\subsection{G.6 Finite probability audit}\label{g.6-finite-probability-audit}

For a witness family \(W\), mass table \(m\), and partition \(\mathcal P\):

\begin{enumerate}
\def\labelenumi{\arabic{enumi}.}
\tightlist
\item
  verify that every mass is in the declared nonnegative carrier;
\item
  verify that partition cells are disjoint and cover \(W\);
\item
  compute total mass and reject normalization when it is zero;
\item
  compute normalized cell values using exact arithmetic when possible;
\item
  verify that they sum to one in the chosen carrier;
\item
  if symmetry is claimed, enumerate group generators and check mass
  invariance;
\item
  for a refinement, verify the transport rule and mass conservation.
\end{enumerate}

This procedure certifies a finite valuation. A stochastic interpretation needs
the separate protocol fields listed in Chapter 9.

\subsection{G.7 History dynamic programme}\label{g.7-history-dynamic-programme}

For a finite graph and semiring with executable operations, the aggregate over
all graph paths of a fixed length can be computed iteratively:

\[
Z^{(0)}=I,\qquad Z^{(n+1)}=Z^{(n)}M.
\]

The recurrence avoids enumerating every path explicitly, though complexity
still depends on graph size, length bound, and semiring operations. For a
finite-state compilation of admissibility, the reported aggregate must use the
declared endpoint selectors \(\lambda_a Z^{(n)}\rho_b\) and a weight-preserving
bijection to the original admissible paths. To attach a proof verification record, record
the orientation/product convention and verify the matrix recurrence on small
explicitly enumerated fixtures. An all-length algorithm must state why it
terminates or how the star is evaluated.

Coarse-graining can be tested by computing fine and pushed-forward aggregates
and comparing them. Composition preservation requires a separate congruence
check; equality of one total aggregate is not enough.

\subsection{G.8 Reconstruction countermodel loop}\label{g.8-reconstruction-countermodel-loop}

For each proposed quantum principle:

\begin{enumerate}
\def\labelenumi{\arabic{enumi}.}
\tightlist
\item
  formalize it independently of quantum terminology where possible;
\item
  test it on finite classical simplex models;
\item
  test it on the intended finite quantum models;
\item
  test it on post-quantum nonsignalling fixtures when well typed;
\item
  record which model families it excludes;
\item
  search for conjunctions whose discriminating power exceeds the individual
  principles;
\item
  only then attempt a representation theorem.
\end{enumerate}

A principle satisfied by all comparison families may be foundational but is not
a discriminator. A principle satisfied by the quantum examples but also by an
unexamined post-quantum family is not yet sufficient. Countermodel search is not
a proof of completeness, but it is an efficient guard against premature
reconstruction claims.

\subsection{G.9 Profile-status verifier}\label{g.9-profile-status-verifier}

For finite query family \(Q\), enumerate each query and check:

\begin{itemize}
\tightlist
\item
  parsing and typing;
\item
  profile membership;
\item
  policy refusal predicates;
\item
  proof and refutation predicates;
\item
  exclusivity of terminal conditions;
\item
  evidence attachment and verifier success.
\end{itemize}

Coverage holds for \(Q\) when each member receives exactly one allowed status.
If both proof and refutation predicates hold, the evaluator is inconsistent on
that query or the predicates refer to different proposition versions. The
verifier should expose the conflict rather than choose by priority.

\subsection{G.10 What algorithms do not settle}\label{g.10-what-algorithms-do-not-settle}

All procedures above consume a representation. They do not select the
representation, the relevant automorphism group, the semantics of a record, the
physical measurand, the observer profile, the semiring, or the reconstruction
axioms. Those choices remain in the assumption ledger. Finiteness turns many
well-typed questions into checkable ones; it does not erase the need to type the
question.

This is the operational form of the monograph's conclusion: every verified
output should be accompanied by the finite input, declared semantics, algorithm
or proof, and scope boundary that made verification possible.

\subsection{Appendix H. Scoped completion results and obstruction review}\label{appendix-h.-scoped-completion-results-and-obstruction-review}

This supplement adds finite results under explicit hypotheses. It does not
replace the original canonical obligations by broader conclusions. The
executable witnesses use exact rational arithmetic; finite examples validate
implementations, while the arguments below state the general finite results.

\subsubsection{H.1 Chain-clock adequacy and refinement}\label{h.1-chain-clock-adequacy-and-refinement}

Declare a finite occurrence chain with positive edge durations \(d_i\), clock
readings \(r_i\), calibration \(k>0\), and nonnegative error budgets \(b_i\). Assume
\(|k(r_{i+1}-r_i)-d_i|\le b_i\) on every edge. For every interval \([a,c]\),

\[\left|k(r_c-r_a)-\sum_{i=a}^{c-1}d_i\right|
\le\sum_{i=a}^{c-1}b_i.\]

Proof: sum the edge residuals, telescope the readings, and apply the triangle
inequality. Restricting the readings to partition endpoints and summing the
durations and budgets inside each block preserves these inequalities.
Conversely, every endpoint-preserving refinement whose fine block sums recover
the declared coarse durations and budgets induces the same coarse bounds.
This is adequacy relative to the declared reference duration and coupling, not
physical calibration evidence. An endpoint reading alone does not supply the
edgewise hypotheses. The implementation is \texttt{ChainClock} in
\texttt{code/witnessed\_difference/finite\_adequacy.py}. The native owner is
\texttt{omega\_measurement::finite\_chain\_clock}; WDP independently replays named
clock adequacy and selected block/refinement verification records. Given the admitted
durations, calibration and starting reading, the exact construction
\(r_{i+1}=r_i+d_i/k\) has zero residual budgets and commutes with block sums.
This conditional construction does not select a physical clock.

\subsubsection{H.2 Atomic representation and stable pushforward}\label{h.2-atomic-representation-and-stable-pushforward}

For a finite set \(X\) with its full Boolean event algebra, a nonnegative,
finitely additive valuation \(v\), with \(v(\varnothing)=0\), has the unique
representation \(v(A)=\sum_{x\in A}m_x\), where \(m_x=v(\{x\})\). Proof: each event
is the disjoint union of its singleton atoms. Conversely, any nonnegative
singleton masses define such a valuation. On a proper Boolean subalgebra,
the same argument gives masses of the algebra's atoms; individual witnesses
inside one atom are not uniquely distinguishable.

A declared map \(p:X\to Y\) transports masses by
\(m'_y=\sum_{p(x)=y}m_x\), so \(v'(B)=v(p^{-1}B)\). For two mass assignments, put
\(\delta_x=m_x-n_x\). Their exact event distance is

\[\max_{A\subseteq X}|v(A)-w(A)|
=\max\left(\sum_{\delta_x>0}\delta_x,
           -\sum_{\delta_x<0}\delta_x\right).\]

The upper bound follows by separating positive and negative terms; the sets
of positive or negative terms attain it. Pushforward is contractive because
coarse events pull back to a subset of the fine events. The half-\(\ell^1\)
formula additionally requires equal total mass. No probability readout or
Born rule is selected by this theorem. See \texttt{AtomicValuation} and the exact
valuation routines in \texttt{finite\_adequacy.py}.

\subsubsection{H.3 A specified finite function sheaf}\label{h.3-a-specified-finite-function-sheaf}

On the inclusion category of subsets of a finite set \(X\), use union covers
and reports \(F(U)=Y^U\). Restriction is selection of the values on a subset;
identity and composition laws follow from function restriction. Compatible
local reports glue by assigning to each point the value of any patch
containing it. Compatibility makes this well defined; coverage establishes both
existence and uniqueness, including the empty carrier/empty-cover case.
The checked Rust adapter is \texttt{omega\_wdp::presheaf}, including restriction,
cover validation and independent replay of the candidate global section.

This theorem does not hold for arbitrary report presheaves. On
\(X=\{a,b\}\), allow arbitrary Boolean singleton reports but require global
reports to have equal values. Restrictions form a presheaf, yet singleton
reports \(a\mapsto0\), \(b\mapsto1\) agree on the empty overlap and have no global
section. Observer categories and their report constraints must therefore be
specified before a general descent claim is made.

\subsubsection{H.4 Conditional factorization and the reflective-completion boundary}\label{h.4-conditional-factorization-and-the-reflective-completion-boundary}

For a surjection \(\eta:A\to B\), a map \(f:A\to M\) factors uniquely through
\(\eta\) exactly when \(\eta(x)=\eta(y)\) implies \(f(x)=f(y)\). Necessity follows
from applying the factor map to an equality. For sufficiency, choose one
preimage of each \(b\) and define \(g(b)=f(x_b)\); kernel compatibility makes the
choice immaterial, and surjectivity establishes uniqueness. An explicit section
makes the construction independent of classical choice. Collapsing the two
Boolean points cannot factor their identity map.

For contextual systems, suppose a morphism \(\eta\) preserves and reflects the
profile relation, and intertwines every context action. If its target profile
is a congruence, its source profile is already a congruence: transport a
profile-related pair to the target, apply the target context congruence,
intertwine, then reflect the resulting relation. These assumptions therefore
cannot complete a genuinely noncongruent source. A one-step refinement also
cannot be placed in the category of congruent systems until membership in
that category is established. This is a set-level obstruction, not a HoTT
identification or a general reflector theorem.

All seven supporting theorems compile independently with Lean 4.33.1 in
\texttt{lean4/WitnessedDifference/RefinementFactorization.lean}. Their separate axiom
report is \texttt{lean4/refinement\_factorization\_axioms.txt}; its union is
\texttt{Classical.choice} and \texttt{Quot.sound}, with no admitted proofs. This supplement
is not registered in the historical 202-declaration axiom census.

\subsubsection{H.5 Exact finite quantum diagnostics}\label{h.5-exact-finite-quantum-diagnostics}

For real rational qubit density matrices, their difference is
\(\Delta=\begin{pmatrix}a&b\\b&-a\end{pmatrix}\), and
\(\Delta^2=(a^2+b^2)I\). Thus trace distance satisfies \(D^2=a^2+b^2\) exactly,
even when \(D\) is irrational. The diagonal-shadow total variation is \(|a|\),
so its squared contraction slack is exactly \(b^2\). Returning a rounded square
root as an exact rational distance is invalid. The Rust adapter now retains
the squared certificate and refuses an irrational rational-root request.

For two declared unweighted finite \(U(1)\) connection Laplacians on the same
ordered carrier, let \(H=L'-L\) and
\(B=\max_i\sum_j(|\operatorname{Re}H_{ij}|+|\operatorname{Im}H_{ij}|)\).
Hermiticity implies \(\|H\|_2\le B\), because the absolute row and column sums
coincide. Weyl's inequality then bounds each ordered eigenvalue displacement
by \(B\), and each adjacent gap shrinkage by \(2B\). A valid independently supplied
gap lower bound \(g>2B\) certifies survival of that gap. The computation does not
certify a caller's spectral-gap premise. \texttt{omega\_wdp::connection} delegates the
actual Laplacians to \texttt{omega\_holonomy}; no numerical eigensolver is substituted
for that premise. General weighted vector-bundle uncertainty remains separate.

\subsubsection{H.6 A finite profile evaluator}\label{h.6-a-finite-profile-evaluator}

Declare a finite Boolean semantic table, an admitted query subset, a refusal
subset, a replay budget and candidate evidence. Recheck candidates against the
table independently of their claimed result. On admitted, allowed queries,
the verified true and false predicates are disjoint because each query has one
Boolean value. Their negation gives Undetermined; the two remaining cases
are OutsideProfile and Refused. The five guards are exclusive and exhaustive.
This is a conditional finite model, not an oracle for arbitrary statements.

An injective, total, table-value-preserving query encoding transports verified
evidence when the source and target both replay it successfully. A change in
admission, policy or budget can change the status, even with preserved semantic
values. The implementation in \texttt{code/finite\_profile\_evaluator.py} checks both
model preservation and replay; it does not equate successful transport with
status invariance under arbitrary resource changes.

\subsubsection{H.7 Outstanding requirements and native decomposition}\label{h.7-outstanding-requirements-and-native-decomposition}

The full row-by-row review is in \texttt{docs/derivation\_review.md}. No unconditional
quantum reconstruction follows: contextual tables need not have joint runs,
no-signalling permits a PR box, gauge invariance permits distinct readouts,
and Gaussian rational phases do not admit the required rational-cochain
extension. Trial-derived statistical tolerances, physical traceability,
nonclassical composition axioms and a well-typed refinement category remain
prerequisites, rather than conclusions of the finite witnesses above.

Native ownership stays with the existing math crates: cochain boundaries and
Betti numbers in \texttt{omega\_homology}, exact image/rank and PSD checks in
\texttt{omega\_linear\_algebra}, quotient defects and primal-dual certificates in
\texttt{omega\_functional}, and transport/Laplacian calculations in \texttt{omega\_holonomy}.
WDP retains packet declarations, adapters and independently checked verification records.
An abstract simplex extension and candidate-face membership check are not a
Kan completion of a simplicial set; the historical adapter names are explicitly
scoped accordingly.

Validation of this supplement: 42 targeted Python tests passed and the seven
new Lean theorems compiled. Six function-sheaf module tests passed in an
isolated source harness. During the first pass, production Rust integration was blocked by a reviewed
specification pin mismatch. Appendix I records its explicit current-state review
and the subsequent production validation; the verification code was preserved. The first pass also ran five quantum/transport tests in an actual-source
harness. The subsequent integration validates the production dependencies.

\subsection{Appendix I. Native integration and discharged finite targets}\label{appendix-i.-native-integration-and-discharged-finite-targets}

The original target identifiers are retained for custody. Eight former open
targets acquired precisely stated finite propositions or conditional results;
the eleven targets left at that stage are resolved in Appendices J--K. The local research report on the source derivation and its primitives inform
the packet vocabulary. Access to the external source repository is not a
prerequisite for these derivations.

\subsubsection{I.1 Greatest behavioral congruence and minimal quotient}\label{i.1-greatest-behavioral-congruence-and-minimal-quotient}

For a finite nonempty carrier \(M\), total updates \(u_i:M\to M\) and exact
readout \(r\), put \(R_0=\ker r\) and

\[R_{k+1}=R_k\cap\bigcap_i (u_i\times u_i)^{-1}R_k.\]

Every \(R_k\) is an equivalence relation by induction. A strict refinement
increases the number of classes, so stabilization occurs after at most
\(|M|-|M/R_0|\) strict refinements. The fixed point \(R_*\) is an update
congruence. If \(B\) is any report-respecting invariant relation, induction gives
\(B\subseteq R_k\) for every \(k\); conversely \(R_*\) itself has those properties.
It is therefore the greatest such relation. Induction also identifies \(R_k\)
with agreement after all generated input words of length at most \(k\).
Words here are finite compositions of the admitted updates, not an external
clock or global ordering of events.

Reports and updates descend to \(M/R_*\). For every surjective report/update
morphism \(f:M\to N\), its kernel lies in \(R_*\), and
\(g(f(m))=[m]\) is well defined and unique. Hence \(|M/R_*|\le |N|\).
Two packets with the same declared input labels have the same behavior set
exactly when every behavior class in their disjoint union meets both carriers;
then their minimal quotients are label/readout-preserving isomorphic.
A declared report equivalence is handled by readout into its equivalence
classes. Nontransitive tolerances are rejected rather than silently closed.

The reusable algorithm and checked refinement/factorization verification records belong
to \texttt{omega\_category\_theory::finite\_moore}; WDP translates named packets into
that owner. Seventeen new Lean theorems establish the semantic congruence,
quotient and comparison properties. The finite termination bound above is
proved in prose and implemented in the checked native refinement routine.

\subsubsection{I.2 Boolean atoms and individual-witness uniqueness}\label{i.2-boolean-atoms-and-individual-witness-uniqueness}

For any finite Boolean event algebra, points with identical membership in all
events form an atom. Each atom is the finite intersection of its signature's
events and complements, hence is measurable. The atoms partition the carrier,
and every measurable event is their disjoint union. Disjoint additivity gives
its unique sum of atom masses. This establishes the representation on a coarse
algebra as well as on a full powerset.

For nonnegative extensions to individual witnesses, a zero-mass atom forces
all its point masses to vanish. A singleton atom fixes its point's mass.
A positive nonsingleton atom admits distinct first-point and last-point
concentrations with identical measurable restrictions. Thus singleton-level
extension is unique exactly when every nonsingleton atom has zero mass.
For a declared measurable map \(f\), pushforward is
\(v'(B)=v(f^{-1}B)\), preserving disjoint additivity and total mass. Coarsening
restricts the event algebra. None of these operations selects a canonical
point inside a positive indistinguishable atom.

\texttt{omega\_finite\_probability::finite\_valuation} owns the checked algebra and
exact mass operations. WDP retains named witness/event declarations.

\subsubsection{I.3 Coarse-graining classified by intertwining}\label{i.3-coarse-graining-classified-by-intertwining}

Select rational weights, allowing signed values, with row-source matrix \(M\).
For a nonempty block partition, let \(H_{x,j}=1\) when \(x\) lies in block \(j\).
A quotient matrix \(B\) preserves representative-independent endpoint-block
path sums at every finite length exactly when each outgoing sum
\(\sum_{y\in B_j}M_{xy}\) is constant for \(x\) in a fixed source block.
Those constants uniquely define \(B\), and the criterion is \(MH=HB\).
Necessity follows at length one. Sufficiency follows by induction:

\[M^0H=H=HB^0,\qquad
M^{n+1}H=M(M^nH)=M(HB^n)=(MH)B^n=HB^{n+1}.\]

The argument uses semiring laws, not positivity or absence of cancellation.
Finite enumeration of all set partitions gives a complete classification for
the selected input matrix. Equal unrestricted totals do not imply this
criterion, and endpoint-block preservation does not imply preservation of
individual histories; signed countermodels distinguish those predicates.
The generic matrix routine lives in the linear-algebra owner; WDP supplies
the rational adapter. The new Lean module establishes the power-intertwining
implication and its length-one converse under explicit associative/unital laws.

\subsubsection{I.4 Sequential calibration and replayed measurement}\label{i.4-sequential-calibration-and-replayed-measurement}

Assume an input enclosure \(|t-x|\le r\), nominal calibration \(at+b\), gain error
\(|\delta a|\le\alpha\), offset error \(|\delta b|\le\beta\), and residual
\(|e|\le\eta\). Subtracting the nominal center gives
\(a(t-x)+\delta a\,x+\delta a(t-x)+\delta b+e\), so the next radius is

\[r'=(|a|+\alpha)r+\alpha|x|+\beta+\eta.\]

Induction propagates every admissible chained realization, including
adversarially correlated errors. Exact nominal composition follows
\((A,B)\mapsto(aA,aB+b)\). A nominal update \((a,b)\mapsto(a',b')\) changes the
output over the input interval by exact supremum
\(|(a'-a)x+(b'-b)|+|a'-a|r\), attained at an endpoint.
This is deterministic uncertainty calculus under declared coefficient bounds,
not a statistical confidence interval or authentic calibration traceability.

The paper's end-to-end instance uses its actual registration, witness,
observation and measurement APIs. It rechecks a separately re-accessed mark,
the composed nominal map, exact calibration correction and declared enclosure.
False records and freshly rehashed false enclosures fail algebraic replay.
For the explicit instance, the retained indication has center \(2\) and radius
\(1/10\). Stage one has \((a,b,\alpha,\beta,\eta)=(2,1,1/10,1/20,0)\),
so its center is \(5\) and radius \(23/50\). Stage two has
\((-3,2,0,0,1/50)\), giving center \(-13\) and radius \(7/5\).
The composed nominal map is \(t\mapsto-6t-1\). The declared correction \(+1\)
and documentary additive hierarchy bound \(1/10\) give corrected center \(-12\)
and radius \(3/2\). Replay binds these values to the actual registration,
witness and measurement packet; it does not infer physical traceability from
the documentary fields.
The native measurement owner uses checked kernel intervals; WDP binds named
calibration stages, input and recomputed output in a replayable verification record.

\subsubsection{I.5 Finite quantum instruments and canonical dilation}\label{i.5-finite-quantum-instruments-and-canonical-dilation}

For grouped rectangular Gaussian-rational Kraus operators \(K_{oi}\) with
\(\sum_{o,i}K_{oi}^*K_{oi}=I\), define
\(\mathcal M_o(\rho)=\sum_i K_{oi}\rho K_{oi}^*\) and
\(E_o=\sum_i K_{oi}^*K_{oi}\). Congruence preserves positivity, and cyclicity of
trace gives \(p_o=\operatorname{tr}(E_o\rho)\), \(p_o\ge0\) and \(\sum_o p_o=1\).
When \(p_o>0\), the normalized post-state is \(\mathcal M_o(\rho)/p_o\) and stays
Gaussian rational. Zero-probability post-state normalization is refused.

The declared dilation \(V\psi=\sum_{o,i}K_{oi}\psi\otimes|o,i\rangle\) obeys
\(V^*V=I\). Projecting onto an outcome's environment coordinates and tracing out
the environment recovers each \(\mathcal M_o(\rho)\), not only their sum.
The exact channel and Hermitian owners supply this algebra. The new instrument
owner supplies grouping and normalized post-states; WDP independently replays
verification records against the actual input and ordered instrument presentation.
A shared POVM need not determine the post-state, as the dimension-three
countermodel verifies. Unitary completion, Born-rule selection, apparatus
identification and metrological traceability remain outside this result.

\subsubsection{I.6 Native profile evaluator and integration repairs}\label{i.6-native-profile-evaluator-and-integration-repairs}

The finite Boolean evaluator of Appendix H now has a checked native adapter.
It replays an ordered evidence prefix under the declared budget, respects
admission/refusal guards, binds a verification record to the profile and candidates, and
recomputes it during verification. Profile transport checks a total injective
value-preserving encoding and evaluates transported evidence in the target's
own policy and resource model. It never copies a status as evidence.

Production integration required two narrow repairs. The quotient-transport
projection received a new explicit review against the actual committed
current source/projection, with the prior row archived verbatim. Verification
code and specification laws were unchanged. The quantum owner also had two
incompatible \texttt{PauliString} definitions; the phase-carrying type was renamed
\texttt{PhasedPauliString}, preserving the consumed legacy API and both test families.
These repairs enabled tests against the real production dependencies.

The native ownership map and exact validation results are recorded in
\texttt{docs/derivation\_review.md}. The historical package axiom census remains
unchanged; each new core-only Lean module has its own source-bound report.

\subsection{Appendix J --- Expert derivations, exact interfaces, and native ownership}\label{appendix-j-expert-derivations-exact-interfaces-and-native-ownership}

This supplement answers ten of the eleven then-open finite targets by
constructive theorems, exact obstruction criteria or explicit countermodels.
CH10-O01 is completed in Appendix K after exact original source recovery.
All 154 original identifiers remain unique. No completion selects physical
laws or authority.
Full proofs and implementation evidence are in the linked expert notes and
\texttt{docs/derivation\_review.md}; new Lean milestones have separate source-bound
axiom verification records, without altering the historical 202-declaration census.

\subsubsection{J.1 Retention components and succession}\label{j.1-retention-components-and-succession}

Fix finite typed sets of sizes \(n=|M|>0\), \(k=|I|\), \(m=|O|>0\). Retention
extensions are precisely independent choices of \(u:I\times M\to M\),
\(r:M\to O\) and a subset of \(O\times O\), hence their count is
\(n^{nk}m^n2^{m^2}\). Projection and assembly are inverse: typing supplies no
cross-component law. Restricting updates to permutations gives \((n!)^k\);
restricting readouts to injections gives \(m!/(m-n)!\) when \(n\le m\), and none
otherwise. Reflexive symmetric comparisons number \(2^{m(m-1)/2}\);
equivalence comparisons number the Bell number \(B_m\). A congruence condition
mentioning both updates and reports is an additional coupling law. Identity
updates and diagonal comparison are canonical choices; a readout cannot be
selected invariantly under a fixed-point-free independent output symmetry.

For arbitrary comparison \(C\), define \(W(x,y)\) by
\(C(r(u_wx),r(u_wy))\) for every finite input word \(w\). The empty word gives
report respect; prepending a label gives forward compatibility. Induction
shows any compatible relation lies in \(W\). Deleting immediately compared
pairs whose successors fail computes \(W\) on a finite carrier. Reflexivity and
symmetry are inherited from \(C\), but transitivity need not be. A state quotient requires \(W\) to be an equivalence; for reflexive symmetric comparison this is exactly the additional transitivity check. A literal descended readout additionally requires
\(W(x,y)\Rightarrow r(x)=r(y)\). Neither condition is supplied by tolerance.

Under total deterministic updates and validated report equivalence, the complete six-predicate retention classification has 19 signatures:
nontrivial carrier, reversible updates, raw injectivity, realized comparison
identification, behavioral minimality and active update. The only constraints
are the singleton signature, immediate recognized injectivity implying
minimality, and inactive updates being reversible with minimality exactly
when the recognized report is injective. Cycles and delayed-report chains
construct every permitted signature on at most four states. These invariants
supplement the structural-data theorem, rather than replace its vocabulary.

For finite labelled relational updates with exact reports or validated report equivalence, unions of strong bisimulations give greatest
bisimilarity; inverse and composite bisimulations prove equivalence. The
bounded quotient satisfies both forward and back clauses and is minimal.
Its finite refinement and quotient-isomorphism classification need no global
clock. A raw directed relation embeds in a strict transitive relation iff it
has no directed cycle: transitive closure establishes sufficiency, and a cycle
contradicts irreflexivity. SCC condensation always gives an order on the
changed carrier. Swapping a selected ordered pair in a fully symmetric
carrier establishes that a nonempty invariant strict order cannot be selected.

Full proofs: \href{docs/swarm_retention_choices.md}{typed choices},
\href{docs/swarm_functional_classification.md}{predicate classification}, and
\href{docs/swarm_relational_retention.md}{relational retention}.

\subsubsection{J.2 Routing, reflection and report descent}\label{j.2-routing-reflection-and-report-descent}

Fix a context alphabet. In the category of profile equivalences with total context actions, arrows
preserve equivalence and intertwine every action. The least action-stable
equivalence containing the source profile exists by intersection. The identity
function is its unit. Any equivariant profile-preserving map into a congruent
target respects that closure by induction on its generators and factors
uniquely through the unit. The construction is functorial and the unit is
natural. The unit is initial in the comma category \(S\downarrow\mathrm{Cong}\), with an initial same-carrier \textbf{coarsening}. If the unit must also
reflect profile equality, equivariance and target congruence force source
congruence; a noncongruent three-state source therefore has no such completion.
This refutes the stronger initiality proposal without losing the paid reflector.

The typed native routing adapters first replay the actual quotient and closure
certificates, then apply a separately declared partial three-class terminal
table, coarsening/value-edit permissions and requests. Missing policy is refused.
A cochain lift supplies checked maps \(L_0,L_1,L_2\) with
\(L_1d_0=d'_0L_0\) and \(L_2d_1=d'_1L_1\), actual coordinate inclusion and field
transport. It induces a quotient map and preserves closedness; it need not
preserve defect norms. Guarded edits keep the window and weights fixed,
strictly decrease replayed \(\Phi_1\), and bound the actual \(\Gamma_2\) damage by
the declared nonnegative budget. A stabilized terminal additionally requires
an actual fixed point of a complete declared deterministic update table;
equal diagnostic values alone do not suffice. Finite checked sequences obey
the summed damage bound, without proving an infinite limit or physical stability.
Every verification record is bound to its actual problem. Neither a defect norm nor a
reflective unit chooses a routing policy.

For sheaf semantics use the finite inclusion category of an intersection-closed
context family, with families covering by union. Report sets and coherent
total restriction maps define a FinSet-valued presheaf. The category of elements
gives actual observer/report presentations. For a cover, the map from global
reports to compatible local families is bijective iff every fiber has one
element; fiber sizes zero and greater than one are respectively obstruction
and ambiguity. For constrained assignment charts, matching gives a unique
pointwise union; it is a permitted global report iff the global constraint
admits that union. On \(\{x,y\}\) with the equality constraint, singleton reports
0 and 1 match on the empty overlap but do not glue. Of four local families,
two glue and two obstruct. Adding the two missing global assignments restores
descent. No additive cohomology conclusion follows from these set-valued data.

Full proofs: \href{docs/swarm_higher_category.md}{reflection and routing} and
\href{docs/swarm_observer_descent.md}{observer descent}.

\subsubsection{J.3 Functional self-reference and archival observers}\label{j.3-functional-self-reference-and-archival-observers}

Declare independent maps \(q:X\to P\), \(s:X\to R\), \(t:X\to X\) and \(b:X\to R\).
Let \(F\) mean \(s\) factors through \(q\), \(P\) presentation stability, \(S\) report
stability, \(C\) equality of report and target, and \(B\) a presentation collision.
Exactly 20 signatures occur: \(\neg B\Rightarrow F\) and \(F\land P\Rightarrow S\)
are the only Boolean constraints. Necessity is fiber constancy and composition;
the explicit constructors on at most three states prove sufficiency. Choosing
\(b=s\) or changing one target independently toggles calibration success.

The original R1/R2/R3 vocabulary additionally needs accessible state,
reflexive provenance, later-state replay and a possibly partial comparison
protocol. These are now explicit checked maps. Linked-report generation is
replayed against captured and re-accessed self-state and observer identity.
Report stability compares actual transported reports. Comparison availability
requires both reports/targets and protocol admission; success additionally
requires their declared comparison to pass. A rejected comparison is available.
Later targets factor through a map iff constant on its fibers. Successful
linked-report plans therefore require actual generation, comparison-domain
and later-target compatibility, without silently completing partial tables.
Presentation blind spots constrain presentation-only functions; richer
accessible state can distinguish them. No phenomenology conclusion is used.

The legacy inventory now reproduces all 44 original source texts and every
referenced blob in15 absorbed snapshots (529 occurrences). Appendix K adds the
complete typed adapter criterion and original-source constructions. The
original six-component observer tuple, ten-component operational observer, five-component world-observer tuple, five-component policy tuple,
nine-component normalized observer tuple and general report/sensing/operator forms have distinct
semantics and explicit adapter domains.

The finite layers below retain actual maps and verified verification records. Linear
report/repair semantics belong to the functional owner using the matrix owner;
quantum operators remain full payloads rather than lossy diagonal reports.
Memory, policy, optional closure and source authority metadata stay separate.
See \href{docs/observer_adapter_completion.md}{original-source completion} and
\href{docs/swarm_legacy_observers.md}{the recovered-family adapters}.

\subsubsection{J.4 Reconstruction obstruction, composition and phases}\label{j.4-reconstruction-obstruction-composition-and-phases}

A signed normalized weight vector represents a classical distribution on a
declared report partition iff every block sum is nonnegative. Representing
every individual indicator is possible iff every individual weight is
nonnegative. An outcome-resolved signed update has a stochastic quotient iff
its block aggregation factors through the report map and the resulting
outcome matrices are nonnegative with stochastic total. Fiber constancy establishes
necessity and constructs the unique quotient on realized report classes.

For four binary measurements let
\(w(\lambda)=(1+\operatorname{CHSH}(\lambda))/16\). Eight entries are \(3/16\) and
eight are \(-1/16\); all four pair marginals are nonnegative PR distributions.
Its negativity is \(1/2\). For negativity \(N\), the total absolute mass is
\(1+2N\), so \(|S|\le2(1+2N)\); PR saturates the resulting \(N\ge1/2\) lower bound.
Accepted measurement tables can realize these observable marginals, while no
positive joint classical model can. They also obstruct standard quantum reconstruction preserving fixed bipartite local measurements, since the additional quantum package below excludes PR. A centralized classical machine depending on the measurement settings can reproduce PR; context-indexed global POVMs \(E_{ab}^{xy}=p_{\rm PR}(a,b|x,y)I\) can do so even on a one-dimensional Hilbert space. These encodings do not preserve the required joint-run or local-factor structure.

Declare finite complex Hilbert spaces, density positivity, local Hermitian
involutions, tensor composition and trace pairing. The CHSH operator obeys
\(\mathsf S^2=4I-[A_0,A_1]\otimes[B_0,B_1]\preceq8I\); positivity and trace
Cauchy-Schwarz give \(S^2\le8\). Native exact Hermitian and positive-semidefinite audits verify the stated
matrix bound. Commuting complete matrix units identify the composite as
\(\mathbb C^m\otimes\mathbb C^n\otimes\mathbb C^r\): simultaneous diagonal
blocks have a common multiplicity, and off-diagonal units identify them.
The dimension condition \(\dim H=mn\) forces \(r=1\). This paid package supplies
the intended composition while excluding PR and hidden multiplicity under
the additional dimension clause; CP and primitive difference do not select it.

For a real cochain declare \(\chi_\theta(a)=e^{i\theta a}\). A path amplitude
multiplies by the character of its edge sum. Gauge change \(a\mapsto a+d f\)
telescopes to \(e^{i\theta(f(v)-f(u))}\), so common-endpoint phase-blind readouts
descend. For \(\theta\ne0\), cycle periods are retained modulo
\((2\pi/\theta)\mathbb Z\), not faithfully as real cochain classes. A chosen
logarithm anchors \(\chi_\theta(1)=u\); continuous real characters are periodic.
The exact rational-turn/cyclotomic branch is a distinct scale, not an extension
of the fixed infinite-order Gaussian-rational phase. Declared effects and Kraus
maps provide a covariant operational interpretation. For amplitudes \((1,2)\),
normalized quadratic readout gives \((1/5,4/5)\) and quartic gives \((1/17,16/17)\);
both are phase invariant. Neither normalization nor gauge invariance selects Born.

Full proofs: \href{docs/swarm_qm_reconstruction.md}{reconstruction, tensor factors and real/rational phases}.

\subsubsection{J.5 Arbitrary-context repair and conditional sampling}\label{j.5-arbitrary-context-repair-and-conditional-sampling}

For context frequencies \(f_i\), minimize \(t\) over compatible probabilities
\(p_i\ge0\), coordinate envelopes \(d\ge|p-f|\) and context sums
\(\sum_s d_{i,s}\le t\). Equal overlap marginals and normalization are linear
equalities. Nonnegative slacks give a rational standard-form LP. Every feasible
LP bounds the desired maximum context distance; every compatible family
extends with \(d=|p-f|\), proving equality of the optima. Uniform global
marginals ensure feasibility. A primal/dual certificate \(Ax=b\), \(x\ge0\),
\(A^Ty\le c\), \(c\cdot x=b\cdot y\) establishes global optimality by weak duality
against every feasible candidate. Replay reconstructs the actual problem.

Three point-mass distributions on a three-outcome context require the uniform
repair with radius \(4/3\): distances are \(2(1-p_i)\) and
\(\min_i p_i\le1/3\). All pair defects are 2, so half the maximum is insufficient.
PR data are already compatible and require radius zero, without a positive
global joint distribution.

For \(n_i\) iid categorical trials on \(m_i\) sections,
\(\mathbb E\|f_i-p_i\|_2^2=(1-\sum_s p_i(s)^2)/n_i\) and Cauchy-Schwarz give
\(\mathbb E\|f_i-p_i\|_1^2\le(m_i-1)/n_i\). Markov gives failure at most
\(\min(1,(m_i-1)/(n_i r_i^2))\) for a declared \(r_i>0\); for \(r_i\ge2\) the strict
tail event is empty. Union bounds give simultaneous confidence without
cross-context independence. If the population marginals are compatible,
marginal L1 contraction and triangle inequality bound pair defects by
\(r_i+r_j\) on the good event. The data do not certify the population or iid
premises. Full proof: \href{docs/swarm_empirical_repair.md}{repair and trial protocol}.

\subsubsection{J.6 Repeatability, full verification records and geometric soundness}\label{j.6-repeatability-full-verification-records-and-geometric-soundness}

For two uses of the same complete instrument and distinct outcomes \(j\ne o\), the outcome-pair switching probability is
\(\sum_{a,b}\operatorname{tr}(K_{jb}K_{oa}\rho K_{oa}^*K_{jb}^*)\) (and the total switching probability sums over all unequal outcome pairs \(j\ne o\)).
All cross-outcome products (\(j\ne o\)) vanishing establishes repeatability. Conversely test
the full-rank state \(I/d\): the switching probability is a sum of nonnegative
squared matrix entries divided by \(d\), forcing every product to vanish.
Repeatability need not force projective effects or idempotent state update.

The native adapter checks actual instrument/state/event, external retained-mark
re-access, interaction/evaluator acceptance, indication, composed calibration
and enclosure, units and the declared traceability DAG. Rehashing false
records cannot satisfy the arithmetic replay. Documentary hierarchy checks
do not authenticate physical standards. Full proof:
\href{docs/swarm_quantum_measurement.md}{measurement and repeatability}.

Legacy geometry bridges are also corrected. A zero period for one cochain
does not prove path homotopy; a nonzero period obstructs it only after closedness
on all declared faces is checked. Horn diagnostics check existing faces and
their typed boundaries and cannot fabricate Kan fillers. Simplicial conversion
refuses oriented-cell collapse or missing face boundaries. Path kernels check
composability, closed cycles and nonnegative zero-absorbing mass traces.
No Kan fibrancy or HoTT equivalence is inferred from these finite diagnostics.

\subsection{Appendix K --- Complete source-bound observer adaptation}\label{appendix-k-complete-source-bound-observer-adaptation}

This appendix completes CH10-O01 at its typed partial-adapter scope. All
nineteen registered targets now have finite or conditional resolutions. The
original 44 files are recovered byte for byte from the recorded source commit
and retained in \texttt{docs/source\_archive/coda\_original\_source\_snapshot/}. The 529 absorbed occurrences
are independently checked at their recorded Git commits. The inventory indexes
all 573 occurrences and reviews all 79 definition-cue records, including every
one of the 53 semantic candidates. Custody identifies the actual statements;
the following proofs establish the adapter semantics.

\subsubsection{K.1 Necessary and sufficient adapter criterion}\label{k.1-necessary-and-sufficient-adapter-criterion}

Fix actual source carriers and a typed signature of partial operations and
relations. Declare a surjection on each sort. A unique target interpretation
preserving these mapped semantics exists exactly when each mapped operation
result is constant on its input fibers and each relation's truth value is
constant on its input fibers. Treat undefined as a distinct result: both
definedness and undefinedness are retained. Numerical or terminal values that
must remain literal use identity output maps.

Necessity evaluates the same target input twice. For sufficiency choose a
source preimage of each target input and use its mapped result. Fiber constancy
makes the choice independent; surjectivity supplies every preimage and makes
the target unique. Apply this to every symbol. Products of surjections are
surjective, including the one empty tuple of a nullary operation. Empty input
sorts have no rows. Identity sort maps give a faithful adapter for every
well-typed declared structure.

A required missing semantic field is recoverable from an onto visible
projection exactly when it is constant on the visible fibers. Two admitted
completions with identical visible data and different field values refute
recovery. The actual joint image of visible and required fields is the
coarsest faithful completion retaining both; any other onto faithful
completion has a unique comparison to it. Unused target inputs destroy
whole-target uniqueness, so an image restriction or explicit extensions are
required. These are complete criteria, not a default completion policy.

The category owner decides the finite table case, returning the first conflicting
representatives. The WDP wrapper supplies 13 typed operations and three
predicates. Independent countermodels cover every operation, every predicate,
definedness loss and an omitted layer. The standalone Lean result establishes the
same algebraic iff statements on arbitrary actual carriers; a finite table does
not certify an arbitrary program or a real vector space.

\subsubsection{K.2 The original six-component observer tuple, ten-component operational observer and five-component world-observer tuple}\label{k.2-the-original-six-component-observer-tuple-ten-component-operational-observer-and-five-component-world-observer-tuple}

The original source packet is \texttt{(M,U,p,rho,G,T)}, not the later education
six-tuple. Its context acts on all finite histories. Put
\texttt{q(w)=U*(m0,w)}. A unique context on reached memory states exists iff
\texttt{q(w)=q(v)} implies \texttt{q(rho(w))=q(rho(v))}. Necessity substitutes into the
context; sufficiency chooses a representative history, with uniqueness on the
actual reachable image. Conversely a state context lifts to histories iff
its images of reached states remain reachable; choose a representative of
each such image.

For a checked finite transducer, a product vertex records original memory,
transducer state and memory after its emitted word. Induction on word length
shows the initial vertex and each input edge represent exactly the source
history. Applying the final emission gives \texttt{q(rho(w))}. Every reachable vertex
has a representative word, and the graph is closed under all input symbols.
Checking its fibers therefore decides the criterion for every finite history,
with at most the transducer-state count times the square of the memory-state
count vertices. A failing fiber gives two actual words. This is all-word
verification, not a length cutoff.

The original O1 criterion is nonconstant induced profile. O2 requires two
equal-profile histories separated by a context. The product graph also decides
these criteria, grouping by source profile for O2. A context changing one
profile is insufficient: parity flip changes every profile while preserving
profile equality. Grammar and terminal routing do not occur in O1/O2, so
changing only G/T cannot change them; the source's six-field independence
proposal is refuted. An input-label-independent toggle retains length parity,
also refuting its proposed O0 corollary. A profile-to-grammar encoder remains
an additional explicit map.

The ten-component operational observer has an input-dependent Mealy report and an
independent profile. Both are retained. The five-component world-observer tuple
\texttt{(image\ p,C,H,W,G)} instead supplies a history presentation with actual
contexts, witnesses and grammar. Contexts need not descend to profile equality.
Its facts require an explicit grammar interpretation and forward stability:
membership in a subset must survive every declared context. The finite world
adapter checks the actual image, history closure, witnesses and this condition;
it adds no update or evaluator.

\subsubsection{K.3 Compressed policy, normalized evaluator and report families}\label{k.3-compressed-policy-normalized-evaluator-and-report-families}

The five-component policy tuple \texttt{(C,D,G,R,T)} uses G for transformations and R for
repairs. It is distinct from the source grammar G and profile codomain R. Actual
comparison defects are computed from the carrier coordinates and restriction
maps. A scalar norm is separately declared. In the checked rational
specialization, simultaneous finite signed permutations act on states,
comparisons, defects and repairs; actual intertwining and norm invariance are
checked. Reindexing repairs then preserves the quotient infimum. A verified
primal/dual verification record must match the actual defect, repair matrix and weights;
a valid certificate for another problem is refused. The declared terminalizer
is replayed separately.

The nine-component normalized observer tuple retains its carrier/stalk data, comparison, actual defect,
repair, optional closure, gauge, transformation class and terminalizer. A finite
interaction carrier does not make its coefficient spaces finite. Linear and
operator instances therefore use their actual symbolic matrices/certificates;
finite verification record domains use checked tables. Optional closure descends iff it
annihilates the repair image. An undefined closure ratio or omitted control
budget is not supplied by a default.

Aggregation/transformation/gauge triples need independent carrier, defect,
repair and terminal declarations. General report maps, token predicates,
invariant classifiers and finite-capacity reports retain their actual codomains
without acquiring other layers. General linear sensing may expose repairs:
it becomes a quotient report iff it annihilates the repair image. A supplied
blindspot witness is checked by applying both actual matrices. Memory uses its
exact directed comparison, rather than a symmetric or reflexive closure.
Measurement and self-target extensions retain their separately checked bridges.

\subsubsection{K.4 Operator source and necessary numerical hypotheses}\label{k.4-operator-source-and-necessary-numerical-hypotheses}

For a finite-dimensional Hermitian repair space W and the declared weighted
half-trace norm, quotient optimization has an attained minimum. A Jordan
decomposition of the residual and a polar-bounded linear functional annihilating
W give matching primal and dual costs. The polar is taken for the declared
ambient norm: on trace-zero Hermitian inputs a functional is represented modulo
scalar identities, so the dual bound cannot be copied unchanged from the full
Hermitian space while also retaining the old divergence constraint.
Weak duality establishes sufficiency; finite
norm attainment and a supporting quotient functional prove necessity. The
native rational-complex certificate checker replays every equation, positivity
condition, bound and actual-input binding. Some exact spectral witnesses are
irrational; they remain covered by the mathematical theorem and outside that
rational certificate representation.

A comparison transport descends to quotients iff it sends the actual source
repair image into the target repair image. Positive trace-preserving maps
contract trace norm on Hermitian inputs by transporting their positive and
negative parts. Together with calibrated weights and repair inclusion this
bounds the transported quotient obstruction. Shape or CPTP contraction alone does not
supply the repair-inclusion hypothesis.

For state-generated source defects, a compatible trace-one reference section
sigma yields the traceless repair \texttt{rho-sigma} and
\texttt{delta(rho)=d0(rho-sigma)}. If this repair is admitted, the obstruction vanishes.
An arbitrary free-edge circulation cannot replace those data. Incompatible
reset comparisons instead give a genuine nonzero state-generated fixture.
The source's dg Lie enrichment is also preserved separately: degree-zero
nilpotence does not imply nilpotent action on degree one, so its finite
exponential requires the representation-level condition.

\subsubsection{K.5 Source and implementation evidence}\label{k.5-source-and-implementation-evidence}

\href{docs/observer_adapter_completion.md}{The completion proof and family crosswalk}
bind every family to its actual source and declared domain. Detailed proofs:
\href{docs/swarm_observer_adapter_completeness.md}{unrestricted adapters and sensing},
\href{docs/swarm_original_genesis_observer.md}{original histories},
\href{docs/source_observer_world_adapter.md}{world facts},
\href{docs/swarm_policy_observer_adapter.md}{policy evaluator}, and
\href{docs/source_operator_observer_adapters.md}{operator/vector geometry}.
The \href{docs/source_recovery_report.md}{recovery report} records byte custody.

Reusable partial-structure and history algorithms belong to the category owner;
linear sensing belongs to the functional owner using the matrix owner; trace
certificates belong to the quantum owner. WDP assembles their source packets
and independently replays verification records. Lean establishes the complete stated algebraic
and all-history criteria, with source-bound axiom reports. The canonical row
retains partial Lean coverage because its full source reconciliation and all
matrix/analysis arguments are not one mechanized corpus theorem.

\subsection{Appendix L. Separating the archived source derivation from WDP, and replay bridges}\label{appendix-l.-separating-the-archived-source-derivation-from-wdp-and-replay-bridges}

The programme is oriented from the archived source derivation and its ongoing
continuation through WDP to
mathematical realizations and checked reasoning protocols. Its source origins
and remaining obligations are recorded in
\href{docs/genesis_to_wdp_derivation_programme.md}{the dependency programme} and
\href{docs/genesis_bridge_register.md}{its separate bridge register}. These
additional results do not enlarge the scope of the earlier 154 canonical rows.

\subsubsection{L.1 Primitive provenance}\label{l.1-primitive-provenance}

The source seed treats \texttt{\textless{}\textgreater{}} as one opaque datum. Its source establishes inertia
and the distinction between proposal, supply and admission. The relation
interpretation of \texttt{\textless{}arb\textgreater{}}, the protected prohibition represented by \texttt{\textless{}ara\textgreater{}},
and the resources for witnessing them require explicit admission and law
provenance. Once R(a,b) and irreflexivity are given, a=b would imply R(a,a),
so the distinction and identification obstruction follow by substitution.
The source's provenance proof also shows why those inputs do not follow
from a single opaque inscription alone. The continuation's distinct seeds
have distinct commitments and are kept as named source contexts.

\subsubsection{L.2 Evaluation, permission and exact information}\label{l.2-evaluation-permission-and-exact-information}

For a report E:X-\textgreater Reports and a protected decision A:X-\textgreater Decisions, a decoder
on the actual report image exists iff E(x)=E(y) implies A(x)=A(y). Necessity
follows by composing with the decoder. Sufficiency assigns each occupied
report its representative's decision; the fiber condition makes that choice
well-defined. Equal successful evaluations with opposite permissions supply
a finite countermodel to recovering independent permission from evidence.
No deterministic postprocessing of that same report removes the collision.

The joint signature (E,A) identifies exactly agreement of these two meanings.
Its projections recover both, and every report sufficient for both maps to
it. This is the coarsest protected-question signature; it need not identify
full source states. A declared policy can make permission a function of
evidence. Such a policy and its recognized standing are additional inputs.

One program may implement proposal, evaluation and admission as distinct
typed functions. Three programs may share a false grant flag. The semantic
separation criterion therefore does not prove a required physical actor count.
\texttt{ProposalAdmissionBridge.lean} establishes the abstract image criterion, its
countermodel, joint projections and equality characterization. Native and
Python adapters reuse the category owner's finite descent checks.

\subsubsection{L.3 Replay and certification}\label{l.3-replay-and-certification}

For a declared deterministic step, replay is an ordered fold and obeys
replay(s,first++second)=replay(replay(s,first),second). Induction on the first
log establishes this law. Induction on a declared run establishes equality to replay;
conversely replay supplies a declared run. In the guarded admission model,
invalid or unpermitted proposals leave admitted state unchanged. These laws
are mechanized in the same standalone Lean module.

A record of histories reconstructs protected semantics iff those semantics
are constant on its fibers. Multisets can lose event order, and final states
can lose occurrence count. Newly hashed false records remain false. The
executable verification record check binds independently supplied initial state, ordered
events, declaration, scope and actual transitions; it does not mint authority.
Faithful future execution additionally needs extension congruence and
preservation of protected future observations.

One retained \href{code/pfp_calculus/README.md}{proposal-protocol calculus episode} now connects
an actual model proposal to host audit, extraction and Stage 5. Its unchanged
seed asks for the real integral of \texttt{f(x)=x} from \texttt{0} to \texttt{10}. A fresh Space
Bunny response supplied \texttt{F(x)=x\^{}2/2} and \texttt{50}; the host checked the derivative
and endpoint difference using the original calculus checker. The control
\texttt{F(x)=5x} has the same endpoint difference but is refused because \texttt{F\textquotesingle{}=5}
is not \texttt{f(x)=x}. The saved packet, raw response, candidate identity, source
hashes and \href{code/pfp_calculus/live_20261004/host_receipt.json}{host verification record}
support this one execution. Audit and extraction here are host operations,
not additional model calls. Provider-reported cost is zero; this is not an
independent account-balance reconciliation or authorization grant.

\subsubsection{L.4 Directed source quotients and scalar realizations}\label{l.4-directed-source-quotients-and-scalar-realizations}

The fixed-carrier continuation bridge preserves all n\^{}2 integer assertion
coordinates, including loops and independent reverse arcs. Raising supplies
d0g(x,y)=g(x)-g(y). Marked cyclic diagnostics cancel these repairs exactly,
so D\_F d0=0. Rational repair equality reflects integer equality: fixing a
root o in a nonempty carrier yields integer h(x)=e(x,o) and d0h=e whenever
e=d0g for rational g. The empty carrier has the trivial equality. These are
repair-equality results, not equality of discrete change and norm costs.

A filled mark need not balance, and balancing marked triples need not detect
loops or all other defects. Antisymmetric simplicial interpretation requires
its extra hypotheses. Source term-addition countermodels prevent inferring
a global dynamic quotient update. The adapter uses the existing matrix and
functional owners, with no loop-dropping graph surrogate.

The separate finite quantum realification theorem starts with declared
complex Hilbert states, effects and instruments. The block map R preserves
multiplication and adjoint, has real trace 2Re of the complex trace, and
represents states as R(rho)/2 while effects and Kraus operators remain
unscaled. It preserves outcome probabilities, post-states and sequential
composition in the real submodel carrying the chosen complex structure.
Local composition now has an explicit balancing map T from the ordinary
4mn-dimensional real tensor to the 2mn-dimensional real view of the
declared complex tensor. Its kernel is exactly the complex balancing
relation, and P=T\^{}t T/2 is the balanced projector. Operators compress with
a factor 1/2; products of individually normalized realified densities do
not. Arbitrary complex composites have the normalized lift T\^{}t R(rho)T/4,
preserving local effects and declared CP instruments. General written
finite-matrix proofs, selected integer coefficient proofs in Lean and
exact Q(i) native checks are separated in
\href{docs/complex_tensor_balancing.md}{the tensor supplement}.
This is representation equivalence. It does not select complex scalars,
Hilbert geometry, Born readout or unrestricted real/tensor composition from
an inscription. The proof and exact owner checks are recorded in
\href{docs/genesis_qm_realification_bridge.md}{the quantum bridge}.

\subsubsection{L.5 Complete raising and marked results}\label{l.5-complete-raising-and-marked-results}

The \href{docs/gauge_result_refinement.md}{two-relation model supplement} defines
the carrier and source correspondence used here. For that model,
fixed-root normalization N retains the complete directed
integer table, including loops. At a fixed writing, N(e)=N(f) iff
f-e is an integer gradient. Minimizing N after each common rewriting gives
the complete source result key. A sorted collection of loop, pair and cycle
values can lose incidence and therefore need not determine this key.
The source's nothing result is characterized by zero loops, two-way sums
and ascending triangle sums. These criteria have carrier-wide written proofs
and selected standalone Lean proofs; their concrete programs have separate
finite original/native comparisons.

For the five-edit model, the complete mark is the jointly minimized pair (N(e),F),
where the same writing transports both the table and the oriented fillings F.
Separate minima can lose their joint incidence. At a fixed writing, mark
equality is exactly equality of fillings together with a gradient difference.
Marked cycle sums are invariant under that gradient. Raw canonicalization
and alignment retain the original assertion multiplicities needed to replay
the declared costs.

The retained original-function verification records compare 779 two-relation model observations,
16,292 bounded bulk raisings, 3,643 rewritings and 779 alignments, followed by
2,314 five-edit model observations, 17,489 raisings, 12,907 rewritings and
2,314 alignments. These are finite differential evidence. The sixteen and
nineteen selected Lean declarations for these two supplements prove their
stated algebraic laws and countermodels; they do not formally verify the
concrete Python or Rust implementations. AddTerm and Identify can separate
equal result marks, and identification can destroy marked balance. Ordered
raw intermediates remain necessary for the actual source execution.

Full statements, source pins and evidence:
\href{docs/gauge_result_refinement.md}{raising results} and
\href{docs/marked_result_refinement.md}{marked results}.

\subsubsection{L.6 Larger frontier}\label{l.6-larger-frontier}

The remaining foundational questions concern the primitive admission floor,
end-to-end source/run refinements, dynamic carrier/context coherence,
operational probability/scalar reconstruction, and recognized governance
and execution maps. Jurisdiction, appeal, facts and agreement each require
their own declared meaning and soundness. Checked repair, the admission,
adapter-soundness and continuation protocol, and the discrete deformation
steps must bind to actual carriers, detector adequacy, priced changes,
scoped terminal policies and separately recognized permission. Existing
finite witnesses and proofs support these constructions at stated scopes;
none silently licenses the entire enlarged programme.

The L.3 calculus episode closes one concrete model-to-host path. It does not
establish the entire staged demo campaign, every recursive re-entry, or the
general source-refinement obligation. The \href{docs/current_claim_evidence.md}{current evidence
index} separates this execution from older
reported runs and from written or mechanized mathematical results.

\subsection{Appendix M. Meron characters, finite shadows and conditional limits}\label{appendix-m.-meron-characters-finite-shadows-and-conditional-limits}

This appendix integrates the finite-to-continuum source merged in
PR \#401, at revision
\texttt{3bec3e2dd5659b53c33889d967fcd2b09689043c}. The
\href{docs/meron_coda_continuum_bridge.md}{source review and replay boundaries}
identify the actual functions, assumptions and unfinished targets. These
additional bridge results use the separate source-bridge register; the canonical
154 rows retain their original scopes. Operational authority remains null.

\subsubsection{M.1 Meron cancellation and the observables it preserves}\label{m.1-meron-cancellation-and-the-observables-it-preserves}

Declare a complete orientation orbit U=(Z/2Z)\^{}k, a constant nonnegative
weight A, a baseline sign s in \{+1,-1\}, and a sign character
chi\_m(u)=(-1)\^{}(m dot u), for m in U. Its signed partition contribution is

\[
Z_U=A s\sum_{u\in U}\chi_m(u)
=\begin{cases}A s\,2^k,&m=0,\\0,&m\ne0.\end{cases}
\]

For m nonzero choose a coordinate j with m\_j=1. The involution
u -\textgreater{} u+e\_j pairs equal weights with opposite signs, proving cancellation.
For m=0 every character value is one. A positive surviving sector further
requires s=+1. Completeness of the orbit, invariant weights and the character
law are explicit hypotheses to be established for the selected model.

For a declared real or complex observable O, put

\[
\widehat O(m)=2^{-k}\sum_{u\in U}\chi_m(u)O(u).
\qquad N_U=A s\,2^k\widehat O(m).
\]

For A\textgreater0, deleting a nonzero-meron orbit preserves its observable numerator
iff its coefficient at m vanishes. Orbit-constant O is sufficient, since its
coefficient is then proportional to the cancelled character sum; it is not
necessary. With k=1, A=1, s=1 and O(u)=(-1)\^{}u, the partition contribution is
zero but the numerator is two. For m=0 the retained quantity is the orbit
average of O. One orientation represents it only under an additional
constancy condition. Normalized expectations also require a nonzero total Z.

This supplies an interface between WDP and the shadow-to-completion bridge:
a compressed report determines a
protected query on its actual image exactly when that query is constant on
report fibers. If partition-only reports lose an observable coefficient,
retain that coefficient or its separately verified improved estimator.
Cancellation alone does not preserve arbitrary queries, paid histories or
future source transitions. The image-factor criterion of Appendix K applies
to the chosen report; the character calculation supplies its model-specific
query, rather than supplying an institutional terminal verdict.

\subsubsection{M.2 Finite characters and the declared circle carrier}\label{m.2-finite-characters-and-the-declared-circle-carrier}

For a positive integer N and a declared complex root zeta of exact order N,
the characters of Z/NZ
are chi\_j(n)=zeta\^{}(jn). For a nontrivial character, its sum S obeys
(1-zeta\textsuperscript{j)S=1-(zeta}j)\^{}N=0, so S=0. The trivial sum is N. Products give the
corresponding finite-torus statement. The finite dual has N characters,
or the product of the declared finite orders, rather than a free circle.
Meron orientation characters use only the subgroup \{+1,-1\}.

After admitting the unbounded discrete group Z and the topological group
U(1), every character is uniquely n -\textgreater{} z\^{}n for its value z at one. Evaluation
at one is a bijection onto U(1). For the compact-open topology it is a
homeomorphism: compact subsets of discrete Z are finite, and finite families
of powers vary continuously with z. The same argument identifies the dual
of Z\^{}d with U(1)\^{}d.~These are conditional representation theorems once the
unbounded group, circle and topologies have been supplied.

They do not construct a lattice limit from a single finite period. A
refinement programme still needs compatible groups/maps and a chosen limit
topology. Continuous translations R\^{}d require another declared carrier and
the appropriate continuous-character theorem. Radian coordinates and Haar
measure normalization are separately specified; rational interval
enclosures for pi do not make pi an exact rational coefficient or select the
circle from an L1 obstruction. The source's staged declarations record these
changes of scope. Its orthogonality status functions encode the theoretical
expected delta; they do not independently evaluate root-of-unity matrices.
This describes the PR 401 source reviewed on 2026-10-02. The later PR 407
finite presentation also sums rational character intervals and checks
containment of its assigned expected delta. The new exact native evaluation
in \href{docs/dyadic_refinement_and_duality.md}{the dyadic supplement} instead
actually sums cyclotomic roots; these are different evidence scopes.

\subsubsection{M.3 What a continuum enclosure assumes}\label{m.3-what-a-continuum-enclosure-assumes}

Let a\textgreater0, C\textgreater=0 and p\textgreater0. Let I\_a={[}l\_a,u\_a{]} contain the actual discrete value
d\_a, and suppose a
model-specific proof supplies a common real target c with
\textbar d\_a-c\textbar{} \textless= C a\^{}p.~Then

\[
c\in[l_a-Ca^p,u_a+Ca^p].
\]

Indeed d\_a-Ca\^{}p \textless= c \textless= d\_a+Ca\^{}p, and the input enclosure extends both
endpoints. For several mesh directions replace C a\^{}p by the sum of their
proved nonnegative bounds. Intersecting such enclosures remains sound.
If their widths tend to zero they contain at most one common target, since
any two common targets have distance bounded by every width. Existence from
nested nonempty bounded closed intervals additionally uses completeness of
the declared real carrier.

The merged defect carriers compute these expanded intervals and record
overlap and a contraction flag. They accept the constants and orders as
inputs; overlap does not prove a derivative bound, a PDE estimate or a
common physical target. A smaller declared error term need not shrink the
total width if the discrete input uncertainty grows. Compatible field,
action, state and observable maps require their own preservation and
convergence arguments.

\subsubsection{M.4 Finite-size extrapolation and its missing remainder}\label{m.4-finite-size-extrapolation-and-its-missing-remainder}

The actual occupation-basis checkerboard transfer replay gives
Z(4)=912, Z(6)=33012 and Z(8)=973824 for repetitions=2 and q=3. This function
traces over all 2\^{}L occupation states; it performs no fixed-particle-number
projection. Define the dimensionless log-partition density
g(L)=ln(Z(L))/L. Interpreting it as a thermodynamic free energy requires the
separately specified temperature, sign and normalization conventions.

For x\_i=1/L\_i\^{}2, the source's Neville tableau encloses the value at zero of
the interpolating quadratic. Its exact weights for L=(4,6,8) are

\[
(w_4,w_6,w_8)=(4/15,-81/35,64/21),\qquad
\widehat f=\sum_i w_i g(L_i).
\]

The weights obey sum w\_i=1, sum w\_i x\_i=0 and sum w\_i x\_i\^{}2=0. Thus under
the exact quadratic law g(L)=f+c\_1/L\textsuperscript{2+c\_2/L}4 the result encloses f.~More
generally, if a proof binds the same f and coefficients to all selected
values with remainders \textbar r\_i\textbar{} \textless= b\_i, then

\[
f\in[\widehat f_{lo}-B,\widehat f_{hi}+B],
\qquad B=\sum_i|w_i|b_i.
\]

This follows by cancelling the first two correction terms and bounding
sum w\_i r\_i by triangle inequality. The required b\_i are model-error bounds,
distinct from the arithmetic intervals for the logarithms.

Three exact values alone cannot supply them. Set
H(x)=(x-1/16)(x-1/36)(x-1/64). The zero function and H agree at all three
sample nodes; their limits at zero are respectively zero and -1/36864.
Both are polynomials with expansions in powers of 1/L\^{}2. Adding any multiple
of H to a fitted quadratic preserves the samples while changing the limit.
The original extrapolator returns {[}0,0{]} on those three zero inputs. This
countermodel refutes the claim that exact arithmetic or a finite polynomial
fit alone bounds an unrestricted thermodynamic limit.

For the actual three Meron values, the computed extrapolation is near
1.6941183355 with arithmetic width below 10\^{}-19. Its current verification record carries
no independently justified scaling remainder. That width therefore measures
the enclosed interpolation calculation. Appendix O separately establishes
time/volume and selected spatial-source limits using the actual Hamiltonian
and explicit bounds. Those proofs do not validate this three-point quadratic
fit or supply its missing model remainder.

\subsubsection{M.5 The field-theory frontier}\label{m.5-the-field-theory-frontier}

At the PR 401 source pin reviewed on 2026-10-02, the five higher-dimensional
targets provide a useful order of work:
an actual 2D Meron ensemble, metric/deficit-angle data, 3D winding/linking,
4D simplicial transport and action, and a field-theoretic topological sector.
Their verification records at that pin have different scopes. T1 reuses the 1D value 912
without constructing its labelled 16-site torus model. T2 encloses the
declared angle pi/3 for a five-triangle toy geometry. T3 records proposed
winding/linking data; T4 supplies zero simplicity and Bianchi defects as
constants; T5 supplies charge and sign fields rather than deriving them
from a field configuration and a fermion determinant.

In particular, the declared sign (-1)\^{}(2Q) equals -1 at both Q=1/2 and
Q=-1/2. Charge reversal alone therefore cannot justify the opposite sign
assigned in T5. An actual sign-changing involution, weight invariance and
observable transport need independent proofs. Wire syntax and verification record hashes
certify their respective artifacts, without establishing these physical laws.

Later merges add exact Regge code and further 2D/3D/4D Meron implementations.
The preceding classifications describe the pinned source audit; those later
implementations require their own source correspondence and model-error proofs.

A completed route to the selected field theory must bind actual carrier
refinements, fields and action, model-derived error bounds, a limit
construction, and preserved observables and dynamics. Meron cancellation
can reduce a sign obstruction within such a route. Finite shadow and refinement
interfaces can expose its missing assumptions. The current additions advance
these interfaces while retaining the enlarged O03/O04 obligations.

\subsubsection{M.6 The actual dyadic refinement and its topology choice}\label{m.6-the-actual-dyadic-refinement-and-its-topology-choice}

The later PR 407 presentation declares C\_n=Z/(2\^{}n) with k -\textgreater{} 2k.
The coordinates k/2\^{}n are unchanged under inclusion; equality at a common
later level is exactly equality modulo integers. Thus the algebraic direct
limit is D=Z{[}1/2{]}/Z. It is not already the full circle.

With the discrete topology on D, compatible finite character indices
a\_(n+1) mod 2\^{}n=a\_n identify its continuous dual with Z\_2. With the
circle-subspace metric topology, D is dense in the supplied circle and
its metric completion is that circle. Continuous characters in this
second branch extend to the circle and are indexed by ordinary integers.
The \href{docs/dyadic_refinement_and_duality.md}{refinement supplement} establishes
the common-level initial property, inverse-family classification and
completion/character arguments. Algebraic refinement selects neither topology.

Each finite character has exactly two adjacent extensions, a and a+2\^{}n.
Fifteen selected Lean residue laws prove refinement, restriction and
this two-element fiber at every natural level. They do not formalize the
infinite groups or topology. Eight math cases and one WDP case actually
evaluate finite roots, retain checked prefixes and compare 441 original
source observations under explicit resource limits.

After supplying circle length measure and a continuous periodic F,
its dyadic average differs from its integral by at most omega\_F(2\^{}-n).
For an L-Lipschitz function the bound is L/2\^{}(n+1), by integrating the
cellwise linear error. This conditional limit argument requires the
declared carrier, measure and regularity; it supplies no determinant,
action or thermodynamic remainder for a physical model.

The original enclosure flag checks adjacent overlap and nonincreasing
widths. It accepts {[}0,1{]}, {[}1/2,3/2{]}, {[}1,2{]}, {[}3/2,5/2{]} at declared rate 1/2;
the original mark lift consequently sets its witness flag true. Yet the
first and last intervals are disjoint and all widths stay one. The retained
original-function output therefore refutes using that flag as a common-target
or contraction certificate. Even a sound finite nested prefix still needs
an infinite refinement law and model-specific common-target bounds.
GEN-WDP-20 through GEN-WDP-22 record these results without changing the
canonical 154 rows or closing the five enlarged obligations.

\subsection{Appendix N. Source form specifications and native mathematical ownership}\label{appendix-n.-source-form-specifications-and-native-mathematical-ownership}

The source-derivation course already contributes the opaque admission boundary,
the separation of source/object/apparatus declarations, ordered replay,
protected reports and actual source raising/marked results to Appendix L.
Its 19 family labels and level 00--100 correspondence are educational indexes,
not evidence that all these structures are forced by the initial inscription.
The \href{../../../../../omega_engine_education/documentation/020_genesis/wdp_integration.md}{source-graded education crosswalk} makes
the relationship to Chapters 1--4, 6, 8--17, 20--21 and Appendices K--M explicit.

\subsubsection{N.1 What the 14 form specifications reconstruct}\label{n.1-what-the-14-form-specifications-reconstruct}

The existing source YAML specifications declare their relational-profile
qualification gap. Apartness constructs an extensional quotient under its
supplied laws; isomorphism canonicalization recovers labeled data only with a
transport; primitive integer rays recover original weights only with their
scale; a final-state digest does not reconstruct an omitted history. These
different reconstruction targets must not be conflated with a relational
lossless project-join theorem or a declaration-level decision procedure.

This distinction matters for continuation. The five-edit model counterexample of
Appendix L.5 already shows why a current canonical mark cannot replace the
complete raw representative needed for future edits. Protected current reports
and protected future operations impose different preservation obligations.
An index entry or an \texttt{executable} metadata flag does not discharge either one.

\subsubsection{N.2 One implementation per native owner}\label{n.2-one-implementation-per-native-owner}

The \href{docs/native_owner_decomposition.md}{native decomposition} moves exact matrix path
kernels into \texttt{omega\_epk}, replayed diagnostic/proposal routing into
\texttt{omega\_coda}, and finite lift/closure/terminal policies into
\texttt{omega\_prime\_algos}. Existing WDP paths remain compatibility surfaces. The
owners use the existing graph, matrix, certificate and cochain owners and do
not depend back on WDP. Source custody and paper-specific presentation remain
in WDP.

These are explicitly scoped mathematical profiles. The checked-repair profile supplies
no operational authority grant, and the admission-protocol profile does not replace the
nine-conjunct operational continuation gate. Agreement on a supplied overlap
list is not a general sheaf gluing theorem; a declared contracting sequence is
not proof that a source algorithm produced it. Those boundaries travel with
the extracted implementations.

\subsubsection{N.3 The shared language as a declared interface}\label{n.3-the-shared-language-as-a-declared-interface}

The repository's schema grammar and typed intermediate representation are software interfaces for
versioned carriers, maps, normal forms, protected questions and evidence.
The \href{../../../../../omega_engine_education/documentation/020_genesis/normal_forms/shared_language_boundary.md}{language boundary} proposes one SQL profile
as an initial executable target, with explicit set/bag, NULL, numeric,
dependency and reconstruction semantics. This source increment does not claim
that SQL lowering, a HoTT kernel or an all-encompassing mathematical language has
already been implemented.

Finite restriction/composition laws and extension fibers can be checked by the
current report-descent owners; exact cochains and primal/dual certificates by
their current algebraic owners. Higher identity and transport laws require
their separately declared formal backend. A hash or supplied proof flag refers
to evidence rather than checking it. A vanishing class in an incomplete
obstruction detector does not establish a global extension. The canonical
claim ledger and the five open source-derivation obligations retain their existing
statuses; this appendix supplies integration and ownership, not an unearned
foundational completion.

\subsection{Appendix O. Presentation closure from finite shadows to completion}\label{appendix-o.-presentation-closure-from-finite-shadows-to-completion}

\subsubsection{O.1 Relative continuity closes the realized image}\label{o.1-relative-continuity-closes-the-realized-image}

The \href{docs/cs_q_presentation_completion.md}{presentation-completion supplement}
combines factorization on the realized image with uniform extension. For
p:X→P into a metric space and q:X→Y into a complete metric space, require

\[
\forall\varepsilon>0\ \exists\delta>0:\quad
d_P(p(x),p(x'))<\delta\ \Longrightarrow\ d_Y(q(x),q(x'))<\varepsilon.
\]

This condition is equivalent to a uniformly continuous image factor
g:p(X)→Y. Its unique extension to the completion of p(X) gives the completed
protected query. Metric separation supplies exact fibre constancy at equal
profiles. The extension applies to the actual image, not unused labels.
If p is also uniformly continuous on a metric source, its completed transport
commutes with the completed query by uniqueness.

The separate \href{lean4/WitnessedDifference/PresentationCompletion.lean}{metric Lean module}
now establishes this equivalence and commuting square using actual Mathlib metric
completion, with fifteen selected declarations audited by its own axiom report.
The dated 2026-10-08 record in the
\href{docs/lane_e_character_duality_completion.md\#8-local-verification-record}{completion/duality supplement}
reports the scoped Lake build on Lean 4.33.1 and the pinned Mathlib revision.
A separate 2026-10-13 compatibility replay compiled the unchanged completion
module and its fifteen-declaration report with Lean 4.34.0 and Mathlib
\texttt{5ed2965256430c3649e86755f9576b54eca72435}; its reported axiom union was
\texttt{propext}, \texttt{Quot.sound} and \texttt{Classical.choice}, without \texttt{sorryAx} or a custom
axiom. This is not a fresh build of the full pinned 4.33.1 package: the current
4.33.1 cache lacked compiled Mathlib completion modules. This classical
existence result supplies no executable modulus.
The character, Haar and physical-model statements below have separate scopes;
the historical formal censuses retain their original meanings.

Source Lipschitz control alone is insufficient. On discrete X=\{0,1,2,\ldots\},
p(0)=0, p(n)=1/n, q(0)=0 and q(n)=1 give injective p and a 1-Lipschitz q,
yet the image query is discontinuous at zero. Every finite factorization
can pass while the all-level relative continuity condition fails.

\subsubsection{O.2 Character closure selects the circle branch}\label{o.2-character-closure-selects-the-circle-branch}

For dyadic torsion D with the circle-subspace topology, every continuous
character extends to T. The character group is Z, and its compact-open
topology is discrete: K=\{0\}∪\{1/2\^{}n:n≥1\} is compact and detects every
nonzero integer character at value −1. Taking the character dual again
gives T, with evaluation D→T equal to its dense completion embedding.
The proof uses uniform extension and explicit dual topologies; D in this
topology is not locally compact, so general LCA reflexivity is not applied
to it. Discrete D instead has dual Z\_2 and bidual discrete D.

Thus a coherent character tower requires the selected continuity condition.
The alternating-bit 2-adic family is coherent but sends inputs 1/4\^{}m→0 to
character turns tending to 1/3. It is excluded from the circle branch.
Neither topology is forced by the primitive mark.

The same proof works for each supplied finite product D\^{}d: coordinate
characters give discrete Z\^{}d and character bidual T\^{}d, with dense evaluation
embedding. Coordinatewise Lipschitz bounds give product-grid sampling error
at most Σ\_i L\_i/(2N\_i). This supplies the selected 2-, 3- and 4-torus branch;
it does not identify the physical dynamics on those spaces.

\subsubsection{O.3 Observable closure and its remaining model premises}\label{o.3-observable-closure-and-its-remaining-model-premises}

Band-limited trigonometric integrals descend exactly to sufficiently fine
finite grids. General L-Lipschitz integrals do not: zero and distance to
the grid have identical samples and different integrals. Averages of a
single admitted observable recover its circle integral with error
L/(2N). Arbitrary measurable observables do not share this sampling theorem.

The incoming Haar source returns a recovery flag on a prefix below its
threshold, and a direct shrinking-overlap flag need not bind a common
target. The supplement retains actual calls and distinguishes the wrapper
that detects the illustrated empty intersection. Its positive example
computes actual samples of one piecewise-linear circle observable.

Gaussian-rational scalars with the Archimedean metric also have the standard
coordinate completion R², identified with C; extending multiplication
requires bounded-set control. Fixed-size matrix sums, products and traces
then extend entrywise. Infinite-dimensional or unbounded-operator limits,
state selection and physical field models need further premises.

These results provide a specific conditional presentation closure and an
observable transport criterion. A Meron route must still bind its ensemble,
sign/weight identity, protected numerator, topology and model-specific time
and volume errors. Numeric convergence does not extend a discontinuous
Boolean threshold or supply institutional authorization. The five enlarged
source-derivation obligations remain open; the canonical 154-row ledger is unchanged.

\subsubsection{O.4 A common target for the actual finite Meron chain}\label{o.4-a-common-target-for-the-actual-finite-meron-chain}

The \href{docs/meron_time_and_volume_refinement.md}{time/volume supplement} binds
the original full-Fock transfer to the admitted spinless U=t Hamiltonian.
Its local block is I+(q-1)P\_x for an actual CAR projection. At M repetitions
the physical temperature parameter is theta=M log(q); fixed q with growing
M changes temperature. The removed scalar subtracts M log(q)/4 from the
unscaled log partition density.

The positive similar product and matrix-logarithm bounds give time error
at most M(log q)\^{}2 q\^{}L. For rational q=1+r/M and a common parameter in
{[}r-rho,r+rho{]}, the corrected density has error at most
r\^{}2 q\textsuperscript{L/M+(3/4)(r}2/(2M)+rho) against the exact finite Hamiltonian.
Open-chain splitting changes one bond of norm 3/4, so its log partitions
are near-additive. The resulting thermodynamic target f(theta) satisfies
\textbar ell\_L(theta)-f(theta)\textbar\textless=3 theta/(2L) for periodic chains.

Thus the corrected actual source interval, expanded by
r\^{}2 q\textsuperscript{L/M+(3/4)(r}2/(2M)+rho)+3(r+rho)/(2L), encloses one target.
For bounded r and M\textgreater=L, this radius tends to zero as even L and M grow
and rho tends to zero. This establishes a selected joint time/volume limit
without assuming a quadratic finite-size scaling polynomial. The target
is 3/4-Lipschitz in theta, so its rational-parameter presentation closes
to the admitted nonnegative real parameter line by PC-1.

The analytic model, near-additivity and limit proofs are written. The new
native quantum owner reuses the existing CAR/matrix owners; WDP retains
the declared plan and actual trace. Nineteen Python and five native checks
cover their separate finite scopes. A spatial-mesh field theory, a complete
infinite operator algebra, all observable numerators and general sign-problem
complexity are not established by this partition-density closure.

\subsubsection{O.5 Spin-sector and actual observable transport}\label{o.5-spin-sector-and-actual-observable-transport}

The \href{docs/meron_spin_and_observable_transport.md}{spin/observable supplement}
derives the periodic fermion boundary in the actual occupation carrier. A
staggered diagonal unitary turns the bulk into the antiferromagnetic
Heisenberg chain. Odd parity uses periodic spin coupling; even parity uses
a pi twist. Both matching restricted sectors are needed for full-Fock
readouts. The source and unrestricted periodic spin chain have the same
thermodynamic log density at fixed theta by a one-bond norm bound, while
their finite spectral readouts can differ.

Expanding the actual local breakups gives a graph/orientation-to-history
bijection and binds the SQL staggered observable to two insertion
derivatives of the transfer trace. In the canonical fixture, the raw squared
numerator is 36864 with partition 912. Its time-average square is 48/19;
the equal-time square is 149/57. Retaining the query therefore matters
even after the Hamiltonian has been identified.

A three-block matrix-jet lift establishes the normalized time-average query
converges at fixed volume to the finite Hamiltonian's imaginary-time Kubo
correlation. Its explicit bound preserves parameter uncertainty and gives
derivative convergence directly. It does not establish the growing-volume
correlation or spatial scaling limit. The general analytic proofs remain
written; generic exact derivative multiplication belongs to linear algebra,
the model to quantum physics, and WDP retains the actual plan/readout.

\subsubsection{O.6 A checked finite language path}\label{o.6-a-checked-finite-language-path}

The \href{docs/sql_normal_form_presentation.md}{normal-form presentation} closes
one existing repository schema grammar specification profile through actual typed intermediate representation, native FD checking,
SQL storage and independent reconstruction. On legal set relations, the
binary projection has natural join as a left inverse. The written closure
and two-tuple countermodel proofs establish the schema criterion independently
of the current instance. BCNF components and lossless join remain distinct
from dependency preservation; the existing S,J,T loss counterexample is replayed.

The category module supplies generic finite functional-dependency mathematics. The schema grammar supplies
profile admission, real typed intermediate representation construction, SQL emission and whole-artifact
replay. WDP retains the declared instance and recovery target. Eight new native
and twenty-two actual compiler/SQLite checks pass. The profile binds original
inputs and the executing native binary and refuses substituted certificates,
illegal rows and incompatible recovery targets. It is an immutable text-set
snapshot path; later component edits, richer dependency theories, a second
domain and an all-encompassing mathematical language require further evidence.

\textbf{A worked relational instance: point coincidence.} In general relativity the
hole argument and the point-coincidence reply {[}REF-REL-001; REF-REL-002{]} lead to
labelling events by coincidences of invariant fields rather than by coordinates;
Komar's four scalar invariants are the standard example {[}REF-REL-003;
REF-REL-004{]}. Read as a finite relation, a table of events with attributes
\(\varphi_0,\dots,\varphi_3\) (and further fields) satisfies point coincidence
exactly when equal values of \(\Lambda=\{\varphi_0,\dots,\varphi_3\}\) identify the
row, that is, when \(\Lambda\) is a \textbf{superkey} of the instance. This is the
factorization criterion of CH16-P02 (and R2) with the row identity as the
target: the identity factors through the projection onto \(\Lambda\) exactly when
rows with equal projections are equal, so a superkey is injectivity of the
projection on the instance (compare CH14-P02). Calling \(\Lambda\) a \emph{candidate} key adds
minimality, which is a separate claim about the instance and its declared
dependencies, not a consequence of four dimensions. The archived source paper
made that conflation; on its own three sample rows the minimal keys are
\(\{\text{kretschmann}\}\), \(\{\text{matter density}\}\), \(\{\varphi_0\}\),
\(\{\varphi_1,\varphi_2\}\) and \(\{\varphi_1,\varphi_3\}\), so \(\Lambda\) is a
superkey that is not minimal. A constructed five-row instance has \(\Lambda\) as
its only candidate key, so minimality is instance-level. Heath's decomposition then applies as in
the \href{docs/sql_normal_form_presentation.md}{normal-form presentation}: if
\(X\to Y\) holds, \(R=R[X\cup Y]\bowtie R[X\cup Z]\). That is a statement about one
snapshot. It does not establish that updates are anomaly-free, and the S,J,T
countermodel above shows locally legal component edits whose rejoin violates a
dependency. The \href{code/relational_coincidence/README.md}{relational coincidence checks}
replay these statements exactly (28 tests, including lossy-join negative
controls). No physical claim about spacetime follows; the instance shows only
that point coincidence has a precise relational reading.

\subsubsection{O.7 Actual source pressure and response closure}\label{o.7-actual-source-pressure-and-response-closure}

The \href{docs/meron_source_field_response.md}{source-field supplement} establishes
a common thermodynamic pressure for the declared full-Fock U=t chain with
staggered source lambda. Even-block near-additivity gives field-uniform
volume error 3theta/(2L). A positive similar transfer product gives an
explicit joint time/volume radius; for bounded target parameters it vanishes
when M/L grows, L grows and supplied parameter radii shrink. The original
SQL histories, weighted by v\^{}Q, equal the actual nonzero-source transfer.

Convexity and uniform pressure convergence give a quantitative weak limit
of source-response measures. Three pressure enclosures also close the
triangularly averaged response at every fixed field resolution epsilon,
with error amplified by epsilon\^{}-2. This is a protected response with its
resolution retained. A global uniform-convergence countermodel shows why
pointwise susceptibility cannot be inferred by twice differentiating the
pressure bound. The later O.8 supplies uniform derivative control in an explicit
high-temperature region; outside that region and for spatial field scaling,
the remaining derivative, sector and temperature-scaling obligations stay open.

The written analytic results are separate from twenty-one new Python and
seven new native checks. Thirteen affected Python and eleven affected native
cases are replayed. Quantum physics owns the model/field plan, measurement
owns generic response interval propagation, and WDP retains the declared
plan and trace. These results add no Lean coverage or operational authority.

\subsubsection{O.8 Pointwise response in a proved analytic region}\label{o.8-pointwise-response-in-a-proved-analytic-region}

The \href{docs/meron_analytic_response.md}{analytic-source supplement} uses even
CAR support factorization and two exact polymer expansions to prove a
volume- and time-uniform complex-source bound for 0\textless=theta\textless=1/192,
\textbar z\textbar\textless=1/64. It applies the Kotecky--Preiss criterion with all model premises
derived explicitly, including the periodic wrapping bond and both parity
sectors. Cauchy estimates then supply the derivative control missing from
pressure convergence alone.

For the actual source curve the intensive query \textless Q\textsuperscript{2\textgreater/((2M)}2L) converges
to the common thermodynamic Kubo response. Its exact error majorant is
4B/epsilon\^{}2 + 64epsilon\textsuperscript{2/(19R}4), with B the joint pressure error and
R=1/64. A polynomial schedule L=4k\^{}4, M=L\^{}2, epsilon=R/(4k) gives O(k\^{}-2).
The written proof is separate from twelve new Python and five new native
checks, with thirty-four/fifteen affected earlier cases replayed.

The quantum owner retains the admitted domain and bound; WDP retains its
actual intensive readout. The q=3,M=2 fixture lies outside this region and
receives no analytic certificate. All-temperature susceptibility and interacting spacetime/critical
mesh/coupling scaling remain open. O.9 closes a specific static spatial
source-functional profile under declared fluctuation scaling. No new Lean
coverage, general sign-problem algorithm or authority is claimed.

\subsubsection{O.9 Spatial Gaussian closure of the actual source functional}\label{o.9-spatial-gaussian-closure-of-the-actual-source-functional}

The \href{docs/meron_spatial_gaussian_source.md}{spatial-source supplement}
keeps the same full periodic Meron chain and original signed time histories.
A weighted polymer expansion bounds the local source-response kernel by
C\_K exp(-distance/4), uniformly in volume and time count. For one fixed
Lipschitz circle profile F, this identifies the smeared quadratic query
\textless Q\_F\textsuperscript{2\textgreater/((2M)}2L) with chi(theta) integral F\^{}2, with explicit response and
spatial-variation errors.

At mesh a=1/L and source amplitude sqrt(a), Cauchy estimates and
particle-hole symmetry make the higher source terms vanish. The actual
normalized source functional converges to exp{[}u\^{}2 chi(theta) integral F\^{}2/2{]}
along the existing polynomial common-temperature schedule. Integer circle
characters construct its static Gaussian distribution-valued target. Its
variance is nonzero at theta=0 and on the explicit sufficient range
theta\textless=2\^{}-20. The general source-functional theorem retains theta\textless=1/192.

The generic functional owner supplies exact periodic test profiles and
derived bounds; quantum physics supplies model bounds and weighted moments;
WDP retains those objects and existing bounded readings. Twenty new Python
and six new native checks pass, with forty-six/twenty-two affected cases
replayed. Four new comparisons use original SQL histories. These finite
checks do not formalize the written infinite-family proof in Lean.

The finite signed generating functions are not assumed to be probability
characteristic functions. The limiting Gaussian field is constructed
explicitly; no finite-measure tightness, real-time CCR dynamics, physical
coupling renormalization or interacting quantum field theory is claimed.
The selected circle geometry, mesh and fluctuation scaling remain premises.

\subsubsection{O.10 Time-labelled sources and higher connected responses}\label{o.10-time-labelled-sources-and-higher-connected-responses}

The \href{docs/meron_time_dependent_analyticity.md}{time-labelled source derivation}
keeps the full periodic CAR transfer and inserts an independent source at
each site and half-step. Its occupation-basis expansion retains the periodic
hop signs and the factor \(q^{-LM/4}\). The original SQL rows do not carry
these extra time labels, so the identity is an exact transfer extension,
not a claim of new SQL execution.
At zero coupling it reduces exactly to
\(\prod_x2\cosh((4M)^{-1}\sum_j z_{j,x})\), fixing the time-source
normalization and showing that nonzero temporal modes require interaction.

At \(q=1+r/M\), \(0\le r\le1/192\), the \(LM\) bond and \(2ML\) source
factors have exponent norm at most \(1/(256M)\). Their physical even-CAR
support graph has degree at most \(7M-1\), independently of the source
values. The rooted polymer majorant is at most \(1/(57M)\); its
incompatibility sum is at most \(7|C|/57<|C|\). Consequently the imported
polymer criterion gives, for every \(|z_{j,x}|\le R=1/64\), a zero-free
polydisk and the uniform bound

\[
 \bigl|\log X_{L,M}(\mathbf z)-L\log2\bigr|\le L/19.
 \tag{O-10}
\]

For a bounded spacetime array \(\Psi\), let
\(\Phi_{L,M}(\Psi)=(2M\sqrt L)^{-1}\sum_{j,x}\Psi_{j,x}a_x(s_j)\).
Its connected responses are derivatives of the logarithm of the
\emph{signed} normalized source sum. Cauchy's estimate on O-10 gives

\[
 |\kappa_n(\Phi_{L,M}(\Psi))|
 \le {2n!\over19}\left({\|\Psi\|_\infty\over R}\right)^n
 L^{1-n/2}. \tag{O-11}
\]

The bound is uniform in the time count. Particle-hole symmetry kills odd
orders; every connected order \(n\ge3\) vanishes under the displayed
fluctuation scaling. For fixed \(u\), the nonquadratic source-logarithm
remainder is \(O(L^{-1})\), uniformly in \(M\). This closes the higher-cumulant
subproblem for bounded time-labelled profiles. It does \textbf{not} identify the
non-equal-time covariance, prove a joint spacetime limit or positivity of a
finite signed-history measure, or select a spacetime field topology. Those
are separate open obligations.

\subsubsection{O.11 Positive finite-volume Euclidean covariance and temporal modes}\label{o.11-positive-finite-volume-euclidean-covariance-and-temporal-modes}

The \href{docs/meron_euclidean_covariance_positivity.md}{Euclidean covariance derivation}
uses the exact Gibbs trace of the same full periodic CAR Hamiltonian. Its
occupation-basis matrix is real symmetric, including the periodic sign,
and the staggered site operators are real diagonal. In a real energy
eigenbasis, pair the two orientations of every unequal-energy transition.
For integer time frequency \(m\), each such pair contributes a nonnegative
multiple of the spatial outer product
\(v_{\alpha\beta}(x)v_{\alpha\beta}(y)\), with coefficient

\[
 {2(p_\alpha-p_\beta)\theta(E_\beta-E_\alpha)
 \over [\theta(E_\beta-E_\alpha)]^2+(2\pi m)^2}
 \quad(E_\alpha<E_\beta). \tag{O-12}
\]

Equal-energy terms contribute only at zero frequency. Particle-hole
symmetry makes every site mean zero. Thus every finite-volume Fourier
coefficient matrix of the exact time-ordered covariance is positive
semidefinite. The staggered spatial row sum is the \(A_L=\sum_xa_x\)
covariance divided by \(L\). Because \([H_L/t,A_L]\ne0\) already on the
one-particle sector, its Fourier coefficient is strictly positive at
every nonzero time mode for every even \(L\ge4\) and \(\theta>0\). At
\(\theta=0\) the row-sum kernel is exactly \(1/4\) and these modes vanish.
The zero mode is exactly \(p_L''(\theta,0)\) by the Duhamel formula, so the
earlier static susceptibility limit is the selected zero-frequency limit.

The local double commutator obeys
\(0\le D_{L,\theta}:=\operatorname{Tr}\rho[A_L,[H_L/t,A_L]]\le2L\):
its operator is twice the sum of nearest-neighbor hopping terms, each of
norm one. The spectral coefficients therefore also give, uniformly in
even \(L\ge4\),

\[
 \operatorname{Lip}_{\mathbb T}(k_{L,\theta})\le\theta,\qquad
 0<\widehat k_{L,\theta}(m)\le{\theta\over2\pi^2m^2}
 \quad(m\ne0),\qquad
 |k_{L,\theta}(u)-p_L''(\theta,0)|\le\theta/6.
 \tag{O-13}
\]

In the declared analytic region, the source-disk bound controls
\(p_L''\) uniformly. Hence the row-sum time kernels are uniformly bounded
and equicontinuous. They have uniformly convergent subsequences with
nonnegative Fourier coefficients and the proved static zero mode.

This closes finite-volume Hamiltonian positivity and temporal
nontriviality plus a uniform temporal upper bound. It does not prove
uniqueness of that subsequential limit, a uniform positive lower bound
for nonzero modes as \(L\to\infty\), spatial summability of the
non-equal-time kernel, or comparison with the time-labelled discrete
transfer. Consequently the spacetime field topology and joint continuum
schedule remain open. The finite signed SQL histories are not declared
probability measures.

\subsubsection{O.12 The unique Euclidean thermodynamic limit of the row-sum covariance}\label{o.12-the-unique-euclidean-thermodynamic-limit-of-the-row-sum-covariance}

The \href{docs/meron_euclidean_thermodynamic_limit.md}{thermodynamic-limit derivation}
uses the same object as O.11. Take even \(L\ge4\), the full \(2^L\) CAR
space with the physical periodic hopping sign, \(P_x\), \(K_L\),
\(a_x=(-1)^x(n_x-1/2)\) and \(\rho_{L,\theta}\). Let
\(k_{L,\theta}(u)=\sum_yC^L_{0y}(u)\). It uses the polymer representation of
the analytic response together with the Kotecký--Preiss criterion that
representation imports, together with finite-dimensional spectral theory,
matrix trace inequalities, Cauchy estimates and standard Fourier/measure
representation results. The model-specific estimates are proved in the linked
derivation; the named standard results remain mathematical inputs.

At fixed \(u\), the two Euclidean segments of length \(1-u\) and \(u\) give a
gas of bond and source labels. Its label norm is \(1/128\) and its degree
is \(7\), so \(S_1=3/43\) and \((\Delta+1)S_1=24/43<1\). Every cluster that
contains both sources contains a chain of bonds joining them. Cauchy and
Schwarz estimates then give, uniformly in even \(L\) and \(u\) on
\(|\theta|\le1/96\),

\[
 |C^L_{0y}(u)|\le M_1(96|\theta|)^{d_L(0,y)},\qquad M_1={24576\over43},
 \qquad
 \sup_L\sup_u\sum_ye^{\mu d_L(0,y)}|C^L_{0y}(u)|\le M_1{1+q\over1-q}
 \tag{O-14}
\]

for \(\mu\ge0\) and \(q=96|\theta|e^\mu<1\). The absolute value is essential:
at \(\theta=-1/96\), \(\mu=0\), a signed ratio would give a zero upper bound,
although \(C^L_{00}(0)=1/4\) because \(a_0^2=I/4\). The corrected ratio is
\(q=1\), outside the weighted bound's domain. This does not affect the separate
positive-temperature theorem. The open-chain functions on any two intervals
containing \([-R,R]\) differ by at most \(2M_1(96|\theta|)^{R+1}\). So the
line two-point function \(C^\infty\) exists, does not depend on the
exhaustion, is translation invariant, and inherits (O-14). The wrapping
term is the CAR hopping between modes \(L-1\) and \(0\). A cluster with fewer
than \(L\) bonds therefore lives on an arc, which is the same CAR system as
a segment of the line. The same holds with the wrap sign erased, which is the
spin-periodic chain. The bounds above use only five premises: each bond term is four times a projector; every word on pairwise disjoint components has a factorizing normalized trace; bond terms and sources commute with every local operator off their supports; the staggered source has vanishing mean; and, on the physical ring, bond terms and sources are even. They are CAR statements for the physical ring and spin-language statements for the erased one. The word-level factorization follows from the support and evenness premises (the locality lemma of TL-1). They are checked exactly on \(L=4,6\) for every set of disjoint bond terms and sources. The check also confirms that the Hamiltonian is built from those bond terms, and it rejects four mutations that pass weaker versions of it: a nonlocal bond term, a source dressed by the total parity, an odd bond term, and an extra nonlocal term added to the Hamiltonian. An exact per-site comparison on \(L=4,6\) shows the two
kernels agree below order \(L\) and first differ at \(L\). The ring kernel minus its wrapped line images
vanishes below order \(L\), and Schwarz on \(|\theta|\le1/128\) gives, for
\(0\le\theta<1/128\),

\[
\begin{gathered}
 \sup_u|k_{L,\theta}(u)-k_\theta(u)|\le4M_1L(128\theta)^L\to0,\\
 \|E_L-C^\infty\|_{\ell^1_\mu(\mathbb Z;C[0,1])}\le M_1q^{L/2}\Bigl(4L+{4\over1-q}\Bigr).
\end{gathered}
 \tag{O-15}
\]

Here \(k_\theta=\sum_nC^\infty_{0n}\). The second bound is for
\(0\le\theta\le1/192\), \(\mu\ge0\), and \(q=96|\theta|e^\mu<1\), where \(E_L\) is the ring kernel placed on the window
\(-L/2<n\le L/2\). The limit is unique. All of O.11's subsequential limits
equal \(k_\theta\), and the rate is uniform on \(0\le\theta\le1/192\).

The limit has the following properties.

\begin{itemize}
\tightlist
\item
  Its zero mode is the static susceptibility \(\chi(\theta)\).
\item
  At \(\theta=0\), \(k_0=1/4\).
\item
  \(k_\theta(u)=k_\theta(1-u)\).
\item
  \(k_\theta(u)>0\) for \(\theta<43/56631360\).
\item
  Every Fourier mode is nonnegative, as a uniform limit of O.11.
\end{itemize}

Strict positivity is a separate statement. Integrating by parts twice over
the time circle gives

\[
 (2\pi m)^2\widehat k_\theta(m)=2\theta F(\theta)-\theta^2\widehat g(m).
\]

Here \(F\) is the bond energy and \(g\) is a sum of local Euclidean
correlations. The constants are
\(\|B_0\|^2\le3\), certified exactly,
\(|F-\theta/4|\le M_F(96\theta)^2\) and \(|g-1/2|\le288G\theta\). This yields,
for every \(m\ne0\),

\[
 \widehat k_\theta(m)\ge{\theta^2\over16\pi^2m^2}>0
 \qquad(0<\theta\le\theta_*={2537\over100281465600}).
 \tag{O-16}
\]

The intermediate bound \(\theta^2(1/2-c\theta)\) is strictly positive for
\(\theta<1/(2c)=2\theta_*\), so strict positivity holds on
\(0<\theta<2\theta_*\).

Evidence grade:

\begin{itemize}
\tightlist
\item
  The analytic estimates are written proofs, with the polymer criterion imported.
\item
  Exact rational CAR series on rings \(L=4,6,8\) and open chains check the vanishing orders, the lift, the zero mode, the integration-by-parts identity and the \(B_0\) certificate.
\item
  Mutations of the normalization, the staggering, the \(P\) scale and the wrap sign each move an exact coefficient.
\item
  \texttt{code/tests/test\_meron\_thermodynamic\_limit.py} runs these checks.
\item
  The Lean module \texttt{MeronThermoLimitConstants.lean} states only the rational constants. It compiles standalone with Lean 4.33.1, and \texttt{\#print\ axioms} reports no axiom for its 18 theorems; the receipt \texttt{\_intelligence/research/briefs/meron\_thermo\_lean\_receipt\_20261007.json} records the file hash and exit codes. It is not imported by a build root; the full package build was not run.
\item
  Floating-point values are diagnostics only.
\end{itemize}

The later window argument in
\href{docs/meron_euclidean_thermodynamic_limit.md\#window-positivity-on-0theta1192-added-2026-10-07}{the thermodynamic-limit supplement}
shows in this model \(\widehat k_\theta(m)\ge\theta^2/(36\pi^2m^2)>0\) for
\(0<\theta\le1/192\) and nonzero integer \(m\), under its imported polymer
criterion. The earlier small-temperature estimate is no longer the full
positivity range proved in this package.

Not proved:

\begin{itemize}
\tightlist
\item
  anything for \(\theta>1/128\);
\item
  real-time or spectral-density statements;
\item
  the \(X_{LM}\) comparison;
\item
  the joint spacetime schedule and the distribution topology.
\end{itemize}

The finite signed SQL histories are still not declared probability
measures.

\subsection{Appendix P --- The earlier component calculi decomposed into reusable WDP mathematics}\label{appendix-p-the-earlier-component-calculi-decomposed-into-reusable-wdp-mathematics}

The earlier nine-layer stack, its eight component papers and its
layered-authority monograph are consolidated
here by mathematical dependency. Their nine runtime stages apply six
constructions; they do not supply nine independent foundational calculi.
The \href{docs/coda_mathematical_decomposition.md}{source and implementation correspondence}
retains every primary text, its extra assumptions and its native owner.
This appendix gives written proofs. It adds no claim of a fresh all-estate
Lean build, general Master Normal Form or operational authority.

\subsubsection{P.1 Reports, queries and one factorization criterion}\label{p.1-reports-queries-and-one-factorization-criterion}

Let \texttt{p:X→A} be a report and \texttt{q:X→B} a protected query. Put \texttt{P=p(X)}.
There is a unique \texttt{q̄:P→B} with \texttt{q=q̄∘p} exactly when

\[p(x)=p(y)\ \Longrightarrow\ q(x)=q(y).\]

Necessity follows by applying \texttt{q̄} to the equal reports. For sufficiency,
define \texttt{q̄(p(x))=q(x)}; fiber constancy makes this independent of the
representative. Each realized report has a representative, giving uniqueness.
Extending a decoder to all of \texttt{A} requires additional values outside \texttt{P};
the theorem does not supply them. The finite-shadow and shadow-to-completion
constructions use this same criterion.

If \texttt{p=h∘p\_fine}, equality of fine reports implies equality of coarse
reports. Thus fiber constancy for the coarse report implies it for the
fine report. The reverse fails for \texttt{X=\{0,1\}²}, \texttt{p(a,b)=a} and \texttt{q(a,b)=b}:
the complete pair report decodes \texttt{q}, while the first-coordinate report does
not. Post-processing \texttt{p} cannot recover the lost second coordinate.

In a pseudometric output space, a decoder satisfying
\texttt{d(q(x),D(p(x)))≤ε} implies, on each report fiber,

\[d(q(x),q(y))\leq 2\varepsilon.\]

This is the triangle inequality through the common decoded value. It is a
necessary bound, not a converse in an arbitrary metric space. Three points
at mutual distance 2, with outputs taking all three values and epsilon 1,
satisfy that bound but have no radius-1 center in that three-point space.
At epsilon zero, output equality additionally needs metric separation:
the indiscrete pseudometric has zero distances between distinct values.
An input Lipschitz estimate is another premise, not this decoder theorem.

Standing is itself a protected query when it is to be decoded from a report.
Two states with equal defect numbers and different grant outcomes violate
its fiber criterion. This explains the collision between descent standing and
layered standing without making
authorization a consequence of observation.

\paragraph{P.1.a Decoders and collisions (DP-1, audited 2026-10-06)}\label{p.1.a-decoders-and-collisions-dp-1-audited-2026-10-06}

\textbf{Statement.} Let \texttt{X}, \texttt{A} and \texttt{B} be finite sets. Let \texttt{p:X→A} and \texttt{q:X→B}
be total tables, and put \texttt{P=p(X)}. That each value of \texttt{q} lies in \texttt{B} is a
typing premise of the table; it is not a further hypothesis to be tested.
Exactly one of the following holds.

\begin{itemize}
\tightlist
\item
  \textbf{(F)} There is a decoder \texttt{D:P→B} with \texttt{D(p(x))=q(x)} for every \texttt{x∈X}.
\item
  \textbf{(K)} There is a collision: a pair \texttt{x,y∈X} with \texttt{p(x)=p(y)} and
  \texttt{q(x)≠q(y)}.
\end{itemize}

In case (F) the decoder is unique as a function on \texttt{P}. The statement
includes \texttt{X=∅}. Then \texttt{P=∅}, and the empty function is the unique decoder,
also when \texttt{B=∅}.

\textbf{Proof.} Suppose (K) holds and \texttt{D} is a decoder. Then
\texttt{q(x)=D(p(x))=D(p(y))=q(y)}, which contradicts \texttt{q(x)≠q(y)}. So (F) and (K)
exclude each other.

Suppose (K) fails. Then \texttt{q} is constant on every fiber \texttt{p⁻¹(a)} with \texttt{a∈P}.
The construction above sets \texttt{D(a)=q(x)} for any \texttt{x} with \texttt{p(x)=a}. Every
\texttt{a∈P} has such an \texttt{x}, so two decoders agree at every point of \texttt{P}.

For \texttt{X=∅}, the decoder equation is vacuous. A function \texttt{∅→B} exists and is
unique for every \texttt{B}.

A finite scan decides which case holds. For each realized report it records
the first query value seen, and it stops at the first disagreement. It returns
either the table of \texttt{D} or a colliding pair, and each step compares two
entries of the supplied tables.

\textbf{DP-1′ (scope of uniqueness).} DP-1′ supersedes the candidate statement
``\texttt{Factorable}: a unique table \texttt{D:P→B} satisfies \texttt{D(p(x))=q(x)} for every
\texttt{x∈X}'' in one respect. Uniqueness is asserted for functions on \texttt{P=p(X)}.
Nothing is asserted, neither existence nor uniqueness, for functions on all
of \texttt{A}. Two countermodels show why.

\begin{enumerate}
\def\labelenumi{\arabic{enumi}.}
\tightlist
\item
  Let \texttt{X=\{x\}}, \texttt{A=B=\{0,1\}} and \texttt{p(x)=q(x)=0}. The decoder on \texttt{P=\{0\}} is
  unique. Both total extensions, with value \texttt{0} or \texttt{1} at the unrealized
  report \texttt{1}, satisfy \texttt{D(p(x))=q(x)}.
\item
  Let \texttt{X=B=∅} and \texttt{A=\{a\}}. The empty decoder on \texttt{P=∅} exists, but no total
  function \texttt{A→B} exists.
\end{enumerate}

\textbf{Refinement.} Suppose \texttt{p=h∘p\_f} with \texttt{h:A\_f→A}. Then (F) for \texttt{(p,q)}
implies (F) for \texttt{(p\_f,q)}, with decoder \texttt{D\_f=D∘h} restricted to \texttt{p\_f(X)},
because \texttt{D(h(p\_f(x)))=D(p(x))=q(x)}. The converse fails in the example
above: \texttt{X=\{0,1\}²}, \texttt{p\_f} the identity, \texttt{h(a,b)=a} and \texttt{q(a,b)=b}. There
\texttt{(0,0)} and \texttt{(0,1)} collide for the coarse report.

\textbf{Worked cases.} Take \texttt{X=\{x0,x1,x2\}}, \texttt{p=(0,0,1)} and \texttt{q=(1/3,1/3,2/3)}. The
decoder is \texttt{D(0)=1/3}, \texttt{D(1)=2/3}. Changing \texttt{q(x1)} to \texttt{2/3} produces the
collision \texttt{(x0,x1)}.

The test reads only the tables \texttt{p} and \texttt{q}. Profile metadata attached to a
report is neither sufficient nor necessary for factorization. A constant report
with \texttt{q} the identity on \texttt{\{0,1\}} collides at \texttt{(0,1)}, whatever metadata
accompanies the report.

\textbf{Evidence grade.}

\begin{itemize}
\tightlist
\item
  \emph{Written proof:} the argument above.
\item
  \emph{Exact finite replay} (2026-10-06 audit):

  \begin{itemize}
  \tightlist
  \item
    an exhaustive enumeration over \texttt{\textbar{}X\textbar{}≤3} and \texttt{\textbar{}A\textbar{},\textbar{}B\textbar{}∈\{2,3\}}, 1423 table
    pairs, each in exactly one of (F) and (K) and never with two decoders
    (\texttt{dp1\_exhaustive\_factorable\_xor\_collision\_sizes\_le3},
    \texttt{dp1\_exhaustive\_case\_count});
  \item
    \texttt{dp1\_positive\_fixture\_decoder}, \texttt{dp1\_mutated\_decoder\_rejected};
  \item
    \texttt{dp1\_reverse\_refinement\_fine\_factorable},
    \texttt{dp1\_reverse\_refinement\_coarse\_collides},
    \texttt{dp1\_refinement\_transfer\_exhaustive\_violations};
  \item
    \texttt{dp1\_metadata\_not\_sufficient\_collision};
  \item
    \texttt{dp1\_decoder\_on\_all\_A\_not\_unique},
    \texttt{dp1\_empty\_image\_but\_no\_function\_A\_to\_B},
    \texttt{dp1\_empty\_X\_empty\_B\_unique\_empty\_decoder}.
  \end{itemize}
\item
  \emph{Lean kernel:} the following, stated for arbitrary types, which contain the
  finite case. Codomain membership is carried by the Lean typing, matching the
  typing premise above.

  \begin{itemize}
  \tightlist
  \item
    \texttt{WitnessedDifference.ComponentReports.decoder\_exists\_iff\_fiber\_constant},
    \texttt{WitnessedDifference.ComponentReports.decoder\_unique},
    \texttt{WitnessedDifference.ComponentReports.not\_fiber\_constant\_iff\_collision},
    \texttt{WitnessedDifference.ComponentReports.decoder\_xor\_collision};
  \item
    \texttt{WitnessedDifference.ComponentReports.empty\_source\_has\_decoder};
  \item
    \texttt{WitnessedDifference.ComponentReports.refinement\_transfers},
    \texttt{WitnessedDifference.ComponentReports.refinement\_converse\_fails};
  \item
    \texttt{WitnessedDifference.ComponentReports.fixture\_factorable},
    \texttt{WitnessedDifference.ComponentReports.constant\_report\_collides};
  \item
    \texttt{WitnessedDifference.ComponentReports.decoder\_off\_image\_not\_unique}.
  \end{itemize}

  The axiom report is \texttt{lean4/component\_reports\_axioms.txt}. The module was
  checked standalone with Lean 4.33.1; the rooted \texttt{lake\ build} was not run.
\item
  \emph{Native exact implementation:} exact rational tables with typed refusals.

  \begin{itemize}
  \tightlist
  \item
    \texttt{omega\_category\_theory::finite\_report\_query} decides on total exact
    tables, with decisions \texttt{Factorable} and \texttt{Collision}. It refuses duplicate
    identities, duplicate, missing or extraneous assignments, and values
    outside the codomain.
  \item
    \texttt{omega\_coda} \texttt{Layer3Cs::evaluate\_report\_query} and \texttt{omega\_wdp}
    \texttt{PresentationMap::decide\_query} route through it. The legacy
    \texttt{Layer3Cs::evaluate} used to skip unmapped tokens and accept an
    extraneous row; it now refuses both. \texttt{PresentationMap::new} likewise used
    to accept keys outside the carrier and now refuses them. Each defect
    matched a typing gap named by the DP-1′ correction.
  \end{itemize}
\item
  \emph{Unresolved residual:} countermodel 2 (no total function \texttt{A→B}) is replayed
  but not stated in Lean. The termination and cost of the scan rest on the
  written argument.

  \begin{itemize}
  \tightlist
  \item
    In the native decision, membership in the codomain is a typing premise.
    The decision applies only to total exact tables.
  \end{itemize}
\end{itemize}

\paragraph{P.1.b Separation of declared output classes (DP-2, audited 2026-10-06)}\label{p.1.b-separation-of-declared-output-classes-dp-2-audited-2026-10-06}

\textbf{Statement.} Let \texttt{S} be a finite set of section records and \texttt{d} an exact
rational pseudometric on \texttt{S}. Put \texttt{Z=\{(s,t):d(s,t)=0\}}, and let \texttt{E} be an
equivalence relation on \texttt{S}.

\begin{enumerate}
\def\labelenumi{\arabic{enumi}.}
\tightlist
\item
  \texttt{d\_E({[}s{]},{[}t{]}):=d(s,t)} is a well-defined pseudometric on \texttt{S/E} iff
  \texttt{E⊆Z}.
\item
  It is a metric iff \texttt{E=Z}.
\item
  Suppose \texttt{E=Z} and \texttt{S/E} has at least two classes. Then
  \texttt{Δ=min\{d(s,t):\ not\ sEt\}} is attained by some pair, and \texttt{Δ\textgreater{}0}.
\item
  Let \texttt{p:X→A} be a report and \texttt{q:X→S} a query into the declared family, so
  that \texttt{q(X)⊆S}. Let \texttt{D:p(X)→T} be a decoder into an admitted pseudometric
  space \texttt{T⊇S} that extends \texttt{d}. Let \texttt{ε≥0} with \texttt{d(q(x),D(p(x)))≤ε} for every
  \texttt{x}. If \texttt{E=Z}, at least two classes exist and \texttt{2ε\textless{}Δ}, then the
  class-valued query \texttt{x↦{[}q(x){]}} factors through \texttt{p}.
\end{enumerate}

In the weighted coordinate instance, every record is a total table \texttt{v\_s:Y→Q}
on one declared common coordinate set \texttt{Y}. The distance is
\texttt{d(s,t)=Σ\_\{y∈Y\}w\_y\textbar{}v\_s(y)−v\_t(y)\textbar{}}, with rational \texttt{w\_y\textgreater{}0}. Then \texttt{Z} is
equality of the represented functions.

\textbf{Proof.}

\emph{Weighted coordinate instance.} Each summand is nonnegative and symmetric.
The scalar triangle inequality, multiplied by \texttt{w\_y\textgreater{}0} and summed, gives the
triangle inequality. With all weights positive, a zero sum forces every
coordinate difference to vanish.

\emph{Item 1.} If \texttt{s\ E\ s′}, \texttt{t\ E\ t′} and \texttt{E⊆Z}, then
\texttt{d(s,t)≤d(s,s′)+d(s′,t′)+d(t′,t)=d(s′,t′)}. Interchanging the primes gives
the reverse inequality. The laws of a pseudometric then descend to \texttt{S/E}.
Conversely, if \texttt{d\_E} is well defined and \texttt{s\ E\ s′}, then
\texttt{d(s,s′)=d\_E({[}s{]},{[}s{]})=0}.

\emph{Item 2.} Separation of \texttt{d\_E} says exactly \texttt{d(s,t)=0\ ⇒\ sEt}, that is \texttt{Z⊆E}.

\emph{Item 3.} By item 2, every interclass distance is positive. There are
finitely many interclass pairs and at least one, so the minimum is one of
them.

\emph{Item 4.} If \texttt{p(x)=p(y)=a}, then
\texttt{d(q(x),q(y))≤d(q(x),D(a))+d(D(a),q(y))≤2ε\textless{}Δ}. Since \texttt{q(x),q(y)∈S}, item 3
forces \texttt{q(x)\ E\ q(y)}, and P.1 gives the factorization. At \texttt{ε=0} the same
chain gives \texttt{d(q(x),q(y))=0}, which yields class equality whenever \texttt{E=Z},
with no positive gap needed. Literal equality of outputs needs \texttt{E} to be
literal equality.

\textbf{Worked example.} Take \texttt{Y=\{u,v\}}, weights \texttt{(1,2)} and records
\texttt{a=(0,0)}, \texttt{a′=(0,0)}, \texttt{b=(1/3,1/2)} and \texttt{c=(1,0)}, with classes \texttt{\{a,a′\}},
\texttt{\{b\}} and \texttt{\{c\}}. Then:

\begin{itemize}
\tightlist
\item
  \texttt{d(a,a′)=0};
\item
  \texttt{d(a,b)=1/3+2·1/2=4/3};
\item
  \texttt{d(a,c)=1};
\item
  \texttt{d(b,c)=2/3+1=5/3}.
\end{itemize}

So \texttt{E=Z} and \texttt{Δ=1}, attained at \texttt{(a,c)}. A decoder with \texttt{ε=1/3} has fiber
diameter at most \texttt{2/3\textless{}1}, so the queried class is constant on report fibers.
If instead \texttt{E} were identity on record names, \texttt{(a,a′)} would be a separation
collision, even though the least nonzero distance is still \texttt{1}. A minimum
taken after skipping zero pairs therefore certifies nothing.

\textbf{Countermodels to dropped hypotheses.}

\begin{itemize}
\tightlist
\item
  \emph{Indiscrete pseudometric.} Every distance is zero. A constant report,
  \texttt{q(0)=0}, \texttt{q(1)=1}, decoder \texttt{0} and \texttt{ε=0} pass every error check, yet the
  outputs differ.
\item
  \emph{Strict threshold.} In \texttt{Q} with \texttt{\textbar{}·\textbar{}}, the outputs \texttt{\{0,1\}} have \texttt{Δ=1}. The
  decoder \texttt{1/2} gives \texttt{ε=1/2}, so \texttt{2ε=Δ}, and the outputs still differ.
\item
  \emph{Declared domains.} Records \texttt{s:\{x\}→Q} and \texttt{t:\{y\}→Q}, both with value \texttt{0},
  have equal length and entries but different domains. They are not a
  zero-distance pair of one comparison problem.
\item
  \emph{Partial pair checks.} Take values \texttt{0,0,1} with literal \texttt{E}. Checking only
  the pairs that involve the third record misses the zero collision of the
  first two.
\item
  \emph{Non-equivalence.} Declaring \texttt{aEb} and \texttt{bEc} without \texttt{aEc} gives no
  quotient. The declaration is refused rather than silently closed.
\item
  \emph{Zero weight.} \texttt{w=0} makes \texttt{0} and \texttt{1} zero-distance on one coordinate.
\item
  \emph{Triangle law.} Symmetric distances \texttt{3,1,1} on \texttt{\{a,b,c\}} fail the triangle
  law. The decoder \texttt{c}, with errors \texttt{1} and \texttt{1}, does not bound \texttt{d(a,b)=3}
  by \texttt{2}.
\item
  \emph{Symmetry.} A directed table with \texttt{d(a,c)=d(b,c)=0} and every other
  off-diagonal entry \texttt{1} passes the directed triangles. It has zero errors
  toward \texttt{c}, yet \texttt{d(a,b)=1}.
\item
  \emph{Fiber bound without a center.} The three-point space of P.1 shows that a
  fiber bound alone is not sufficient.
\end{itemize}

\textbf{Evidence grade.}

\begin{itemize}
\tightlist
\item
  \emph{Written proof:} the argument above.
\item
  \emph{Exact finite replay:} \texttt{dp2\_distance\_table},
  \texttt{dp2\_weighted\_is\_pseudometric}, \texttt{dp2\_E\_equals\_Z\_and\_gap},
  \texttt{dp2\_two\_eps\_below\_gap}, \texttt{dp2\_identity\_E\_has\_zero\_collision},
  \texttt{dp2\_indiscrete\_pseudometric\_errors\_zero\_values\_differ},
  \texttt{dp2\_strict\_threshold\_control}, \texttt{dp2\_fiber\_bound\_not\_sufficient},
  \texttt{dp2\_triangle\_control\_refused}, \texttt{dp2\_symmetry\_control}.
\item
  \emph{Lean kernel:} axiom report \texttt{lean4/component\_reports\_axioms.txt}.

  \begin{itemize}
  \tightlist
  \item
    \texttt{WitnessedDifference.ComponentReports.Pseudometric.zero\_equivalence},
    \texttt{WitnessedDifference.ComponentReports.Pseudometric.descends\_of\_zero};
  \item
    \texttt{WitnessedDifference.ComponentReports.Pseudometric.well\_defined\_iff}
    (item 1), \texttt{WitnessedDifference.ComponentReports.Pseudometric.quotDist\_mk};
  \item
    \texttt{WitnessedDifference.ComponentReports.Pseudometric.separated\_iff}
    (item 2);
  \item
    \texttt{WitnessedDifference.ComponentReports.fiber\_diameter\_le},
    \texttt{WitnessedDifference.ComponentReports.strict\_gap\_forces\_equivalence}
    (item 4);
  \item
    \texttt{WitnessedDifference.ComponentReports.equality\_threshold\_countermodel},
    \texttt{WitnessedDifference.ComponentReports.indiscrete\_quotient\_singleton},
    \texttt{WitnessedDifference.ComponentReports.worked\_example\_distances}.
  \end{itemize}
\item
  \emph{Native exact implementation:}

  \begin{itemize}
  \tightlist
  \item
    \texttt{omega\_category\_theory::finite\_section\_separation}
    (\texttt{decide\_separation}, \texttt{verify\_separation},
    \texttt{epsilon\_decoder\_class\_factorization}) returns either a zero-distance
    collision or a separation with an attained gap and its attaining pair. It
    accepts at most 512 points.
  \item
    The certified path in \texttt{omega\_coda} is
    \texttt{Layer4CsQ::evaluate\_declared\_separation}.
  \end{itemize}
\item
  \emph{Unresolved residual:} the attained minimum of item 3 is a written proof
  with finite replay. The declared-domain discipline is a specification on inputs,
  and it is not formalised.

  \begin{itemize}
  \tightlist
  \item
    The legacy \texttt{Layer4CsQ::evaluate} remains a screen labelled ``not a
    separation certificate'', and it is not a certificate. It still accepts a
    pair flagged equivalent at positive distance, and it succeeds when nothing
    is compared.
  \end{itemize}
\end{itemize}

\subsubsection{P.2 Repair quotients, closure and exact certificates}\label{p.2-repair-quotients-closure-and-exact-certificates}

Admit finite-dimensional real spaces, linear maps
\texttt{d0:C0→C1}, \texttt{d1:C1→C2}, \texttt{a2:U2→C2}, the law \texttt{d1\ d0=0}, and strictly
positive coordinate weights \texttt{w1,w2}. Write \texttt{R1=im\ d0}, \texttt{R2=im\ a2} and

\[\Phi(x)=\min_f\|x+d_0f\|_{1,w_1},\qquad
\Gamma(x)=\min_u\|d_1x+a_2u\|_{1,w_2}.\]

These are quotient norms on \texttt{Q1=C1/R1} and \texttt{Q2=C2/R2}, respectively,
with \texttt{Γ} evaluated after the induced map \texttt{D̄:Q1→Q2}. This map is
well-defined because \texttt{d1} kills \texttt{R1}. A minimizing residual exists: the
affine repair coset is closed, and its intersection with a sufficiently
large norm ball is compact. Hence \texttt{Φ=0} iff \texttt{x∈R1}, and \texttt{Γ=0} iff
\texttt{d1x∈R2}. \textbf{The latter is closure modulo the second repair space.} Taking
\texttt{R2=\{0\}} recovers raw closure \texttt{d1x=0}. More generally, the two zero tests
agree for every \texttt{x∈C1} if and only if \texttt{im\ d1\ ∩\ R2=\{0\}}: only vectors in
that intersection can be both nonzero raw defects and zero quotient defects.
The re-presentation papers' raw closed/open square labels must retain this
distinction from the general trichotomy of the finite obstruction calculus.

Since \texttt{Φ=0} forces \texttt{d1x=0}, there are exactly three numerical regimes:
\texttt{Φ=0,Γ=0}; \texttt{Φ\textgreater{}0,Γ=0}; \texttt{Φ\textgreater{}0,Γ\textgreater{}0}. This classification does not prove
that a coherent residual survives every later allowed correction. Such
survival quantifies over a declared move system. Strictly positive weights
change prices but not either nullity condition; zero weights can destroy
that statement. Block-diagonal independent systems have additive \texttt{Φ} and
\texttt{Γ}, because their variables and coordinate objectives separate.

For a general repair matrix \texttt{A}, finite weighted L1 duality is

\[\min_f\|x+Af\|_{1,w}
=\max\{z^Tx:A^Tz=0,\ |z_i|\leq w_i\}.\]

To derive the dual, minimize \texttt{Σ\_i\ w\_i\ t\_i} subject to
\texttt{−t≤x+Af≤t}. The two nonnegative multiplier vectors \texttt{α,β} satisfy
\texttt{α+β=w} and \texttt{Aᵀ(α−β)=0}; setting \texttt{z=α−β} gives precisely the stated
capacities and annihilator. Conversely \texttt{α=(w+z)/2}, \texttt{β=(w−z)/2}
recover feasible multipliers. Finite linear-programming strong duality
applies to this feasible bounded-below primal. For rational input and
weights, rational LP gives rational optimal certificates; compactness of
a rational ball is not used.

A replayable optimum needs primal feasibility, \textbf{all} annihilator
equations, all capacities and equality of the two objectives, bound to
the actual input. Weak duality establishes the equality certifies optimality:
\texttt{zᵀx=zᵀ(x+Af)≤Σ\_i\ w\_i\textbar{}x\_i+(Af)\_i\textbar{}}. One selected annihilator vector
is only a lower bound. For \texttt{A=(1,0,0)ᵀ}, \texttt{x=(0,0,1)ᵀ}, \texttt{z=(0,1,0)ᵀ}
is feasible but pairs to zero while the optimum is one.

Compatible transport maps satisfy
\texttt{T1\ d0=d0′\ T0}, \texttt{T2\ d1=d1′\ T1}, \texttt{T2\ a2=a2′\ S2}. Transport a minimizing
repair to obtain

\[\Phi'(T_1x)\leq\|T_1\|\Phi(x),\qquad
\Gamma'(T_1x)\leq\|T_2\|\Gamma(x).\]

For weighted L1, \texttt{‖T‖=max\_j\ (Σ\_i\ w′\_i\textbar{}T\_ij\textbar{})/w\_j}. Equality or
reversibility requires stronger hypotheses. The finite Python and primary
Rust checked-repair pipelines implement an unweighted first quotient and a raw next
defect norm; they do not implement the whole weighted two-repair theorem.

\paragraph{P.2.a Raw and quotient second defects (DP-3, audited 2026-10-06)}\label{p.2.a-raw-and-quotient-second-defects-dp-3-audited-2026-10-06}

\textbf{Statement.}

\emph{Setting.} Write \texttt{D} for the rational matrix of \texttt{d1:C1→C2} and \texttt{A} for that
of \texttt{a2:U2→C2}. Let \texttt{w=w2} be strictly positive, with \texttt{m} coordinates. Define

\[\Gamma_{\rm raw}(x)=\|Dx\|_{1,w},\qquad
\Gamma_{\rm quot}(x)=\Gamma(x)=\min_u\|Dx+Au\|_{1,w}.\]

Let \texttt{B\_w=\{z:\textbar{}z\_i\textbar{}≤w\_i\}} be the capacity box, and let
\texttt{K\_w=\{z∈B\_w:Aᵀz=0\}}.

\emph{Claims.}

\begin{enumerate}
\def\labelenumi{\arabic{enumi}.}
\tightlist
\item
  \texttt{Γquot(x)≤Γraw(x)} for every \texttt{x}.
\item
  The following are equivalent.

  \begin{itemize}
  \tightlist
  \item
    \textbf{(G)} \texttt{Γquot(x)=Γraw(x)} for every \texttt{x}.
  \item
    \textbf{(P)} \texttt{DᵀK\_w=DᵀB\_w}.
  \item
    \textbf{(GV)} For every sign vector \texttt{σ∈\{−1,1\}\^{}m} there is \texttt{z∈K\_w} with
    \texttt{Dᵀz=Dᵀ(w∘σ)}. Here \texttt{w∘σ} is the coordinatewise product.
  \end{itemize}
\item
  The two zero tests agree for every \texttt{x} iff \texttt{im\ D∩R2=\{0\}}, as in P.2.
  Claim (G) implies this agreement. The converse fails.
\item
  For every vector \texttt{y},
  \texttt{(min\_i\ w\_i)‖y‖\_1≤‖y‖\_\{1,w\}≤(max\_i\ w\_i)‖y‖\_1}.
  This bounds \(\Gamma_{\rm raw}\) and \(\Gamma_{\rm quot}\) by their
  unweighted values times the minimum and maximum entries of \(w_2\).
  For \(\Phi\) use its own weight vector \(w_1\) and the corresponding extrema.
  On an empty coordinate space all these values are zero; no minimum or
  maximum of an empty weight vector is used.
\end{enumerate}

\textbf{Proof.}

\emph{Claim 1.} Take \texttt{u=0}.

\emph{Claim 2.} By the P.2 duality,
\texttt{Γquot(x)=max\{(Dᵀz)ᵀx\ :\ z∈K\_w\}}. This is the support function of the
polytope \texttt{DᵀK\_w}. Putting \texttt{A=0} gives \texttt{Γraw(x)=max\{(Dᵀz)ᵀx\ :\ z∈B\_w\}}, the
support function of \texttt{DᵀB\_w}. Both sets are nonempty, compact and convex, and
\texttt{DᵀK\_w⊆DᵀB\_w}.

\begin{itemize}
\tightlist
\item
  \begin{enumerate}
  \def\labelenumi{(\Alph{enumi})}
  \setcounter{enumi}{6}
  \tightlist
  \item
    ⇒ (P). Suppose \texttt{y∈DᵀB\_w∖DᵀK\_w}. Exact separation of a point from a
    rational polytope gives a rational \texttt{x} with \texttt{xᵀy\textgreater{}max\_\{y′∈DᵀK\_w\}xᵀy′}.
    Hence \texttt{Γraw(x)≥xᵀy\textgreater{}Γquot(x)}, and (G) fails.
  \end{enumerate}
\item
  \begin{enumerate}
  \def\labelenumi{(\Alph{enumi})}
  \setcounter{enumi}{15}
  \tightlist
  \item
    ⇒ (G). Equal sets have equal support functions.
  \end{enumerate}
\item
  \begin{enumerate}
  \def\labelenumi{(\Alph{enumi})}
  \setcounter{enumi}{15}
  \tightlist
  \item
    ⇔ (GV). Positive weights make \texttt{B\_w} a full box whose vertices are
    the points \texttt{w∘σ}. Every \texttt{b∈B\_w} is the convex combination
    \texttt{b=Σ\_σ\ c\_σ(w∘σ)}, with
    \texttt{c\_σ=∏\_i(1+σ\_ib\_i/w\_i)/2≥0}.
    These are the product weights of independent signs with means \texttt{b\_i/w\_i},
    so \texttt{Σ\_σc\_σ=1} and \texttt{Σ\_σc\_σσ\_iw\_i=b\_i}. Therefore
    \texttt{DᵀB\_w=conv\{Dᵀ(w∘σ)\}}. The convex set \texttt{DᵀK\_w} equals it exactly when it
    contains each \texttt{Dᵀ(w∘σ)}, which is (GV).
  \end{enumerate}
\end{itemize}

A vertex that fails (GV) has a replayable negative certificate. The rational
system \texttt{\textbar{}z\_i\textbar{}≤w\_i}, \texttt{Aᵀz=0}, \texttt{Dᵀz=Dᵀ(w∘σ)} is infeasible exactly when
Fourier--Motzkin elimination returns a multiplier vector \texttt{η≥0}. Writing the
system as \texttt{Mz≤b}, the vector satisfies \texttt{ηᵀM=0} and \texttt{ηᵀb\textless{}0}. A checker
replays these equations without trusting the solver.

\emph{Claim 3.} \texttt{Γraw(x)=0} iff \texttt{Dx=0}, and \texttt{Γquot(x)=0} iff \texttt{Dx∈R2}. (G)
equates the two values for every \texttt{x}, so it equates their zero sets.

\emph{Claim 4.} Compare coordinatewise.

\textbf{Hypotheses are needed.}

\begin{itemize}
\tightlist
\item
  \emph{Zero weight.} One coordinate with \texttt{D=1}, \texttt{A=0} and \texttt{w=0} gives \texttt{Γ=0}
  at \texttt{x=1} although \texttt{Dx≠0}.
\item
  \emph{Negative weight.} This can make the quotient objective unbounded below:
  with one coordinate, \texttt{D=A=1} and \texttt{w=-1}, it is \texttt{-\textbar{}x+u\textbar{}}.
  It need not do so: \texttt{Dx=(x,x)}, \texttt{Au=(u,u)} and \texttt{w=(-1,2)} give
  \texttt{-\textbar{}x+u\textbar{}+2\textbar{}x+u\textbar{}=\textbar{}x+u\textbar{}}, with minimum zero at \texttt{u=-x}.
\item
  \emph{Missing annihilator.} With \texttt{A=1}, \texttt{D=1} and \texttt{w=1}, the vector \texttt{z=1} meets
  the capacity bound and pairs to \texttt{1} with \texttt{Dx} at \texttt{x=1}. Yet
  \texttt{Γquot(1)=min\_u\textbar{}1+u\textbar{}=0}. Every annihilator equation must be checked.
\end{itemize}

\textbf{Example 1 (zero tests agree, values differ).} Take \texttt{Dx=(x,x)},
\texttt{Au=(u,0)} and \texttt{w=(1,1)}.

\begin{itemize}
\tightlist
\item
  \texttt{Γraw(x)=2\textbar{}x\textbar{}}, while \texttt{Γquot(x)=min\_u(\textbar{}x+u\textbar{}+\textbar{}x\textbar{})=\textbar{}x\textbar{}}, attained at \texttt{u=−x}.
\item
  \texttt{K\_w=\{(0,z\_2):\textbar{}z\_2\textbar{}≤1\}}, so \texttt{DᵀK\_w={[}−1,1{]}}, whereas \texttt{DᵀB\_w={[}−2,2{]}}.
\item
  The vertices \texttt{(1,1)} and \texttt{(−1,−1)} fail (GV): any \texttt{z∈K\_w} has
  \texttt{Dᵀz=z\_2≤1\textless{}2}. The separator \texttt{x=1} gives raw \texttt{2} and quotient \texttt{1}, with
  dual \texttt{z=(0,1)}.
\item
  The zero tests agree, since \texttt{im\ D∩R2=\{(x,x)\}∩\{(u,0)\}=\{0\}}.
\item
  With weights \texttt{(2,3)}, \texttt{Γraw(1)=5} and \texttt{Γquot(1)=3}. The quotient value is
  attained at \texttt{u=−1} and certified by \texttt{z=(0,3)}. The ratio is \texttt{5/3}.
\end{itemize}

\textbf{Example 2 and DP-3′.} Keep \texttt{Dx=(x,x)} and set \texttt{Au=(u,−u)}.

\begin{itemize}
\tightlist
\item
  With unit weights, \texttt{Γquot(x)=min\_u(\textbar{}x+u\textbar{}+\textbar{}x−u\textbar{})=2\textbar{}x\textbar{}=Γraw(x)}.
\item
  For each \texttt{σ}, the vector \texttt{z\_σ=((σ\_1+σ\_2)/2,(σ\_1+σ\_2)/2)} lies in \texttt{K\_w} and
  has \texttt{Dᵀz\_σ=σ\_1+σ\_2}.
\item
  So (G) holds although \texttt{R2=span(1,−1)≠\{0\}}.
\end{itemize}

DP-3′ supersedes the candidate wording ``demonstrating that \texttt{R2=\{0\}} is
sufficient but not necessary for value agreement''. The corrected wording
reads: \texttt{R2=\{0\}} implies (G), since then \texttt{K\_w=B\_w}, and with unit weights
\texttt{w=(1,1)} this example has (G) with \texttt{R2≠\{0\}}. The example is weight-specific.
With weights \texttt{(1,3)} and the same \texttt{D} and \texttt{A}, (GV) fails at \texttt{σ=(1,1)}. There
\texttt{Dᵀ(w∘σ)=4}, while every \texttt{z∈K\_w} has \texttt{z\_1=z\_2} and \texttt{\textbar{}z\_1\textbar{}≤1}, so \texttt{Dᵀz≤2}.
Directly, \texttt{Γraw(1)=4} and \texttt{Γquot(1)=min\_u(\textbar{}1+u\textbar{}+3\textbar{}1−u\textbar{})=2}, at \texttt{u=1}.

\textbf{Evidence grade.}

\begin{itemize}
\tightlist
\item
  \emph{Written proof:} claims 1--4 and the certificate replay.
\item
  \emph{Exact finite replay:}

  \begin{itemize}
  \tightlist
  \item
    \texttt{dp3\_ex1\_raw}, \texttt{dp3\_ex1\_quot\_lp}, \texttt{dp3\_ex1\_dual\_certificate},
    \texttt{dp3\_ex1\_zero\_tests\_agree\_on\_samples}, \texttt{dp3\_ex1\_weights\_23\_values},
    \texttt{dp3\_ex1\_weight\_bounds}, \texttt{dp3\_ex1\_gv\_vertex\_status},
    \texttt{dp3\_ex1\_farkas\_and\_separator};
  \item
    \texttt{dp3\_farkas\_replay\_rejects\_tampered\_certificates};
  \item
    \texttt{dp3\_ex2\_four\_vertex\_certificate}, \texttt{dp3\_ex2\_values\_equal};
  \item
    \texttt{dp3\_new\_ex2\_weights\_13\_vertex\_fails}, \texttt{dp3\_new\_ex2\_weights\_13\_values};
  \item
    \texttt{dp3\_convex\_coefficients};
  \item
    \texttt{dp3\_control\_mixed\_weights}, \texttt{dp3\_control\_zero\_weight\_nonzero\_defect},
    \texttt{dp3\_control\_missing\_annihilator\_false\_certificate},
    \texttt{dp3\_control\_negative\_weight\_unbounded}.
  \end{itemize}
\item
  \emph{Lean kernel:} these cover claim 1 and Example 1 only. The axiom report is
  \texttt{lean4/component\_repair\_axioms.txt}.

  \begin{itemize}
  \tightlist
  \item
    \texttt{WitnessedDifference.ComponentRepair.lower\_bound\_le\_raw},
    \texttt{WitnessedDifference.ComponentRepair.quot\_le\_raw};
  \item
    \texttt{WitnessedDifference.ComponentRepair.countermodel\_raw},
    \texttt{WitnessedDifference.ComponentRepair.countermodel\_quot},
    \texttt{WitnessedDifference.ComponentRepair.countermodel\_dual},
    \texttt{WitnessedDifference.ComponentRepair.countermodel\_zero\_tests},
    \texttt{WitnessedDifference.ComponentRepair.countermodel\_values\_differ};
  \item
    \texttt{WitnessedDifference.ComponentRepair.weighted\_countermodel}.
  \end{itemize}
\item
  \emph{Native exact implementation:}

  \begin{itemize}
  \tightlist
  \item
    \texttt{omega\_functional::gamma\_reconciliation} provides the raw defect and the
    certificate verifiers for the quotient value, the vertex table, the
    vertex Farkas certificate, the value gap and the zero-test gap.
  \item
    \texttt{omega\_solvers::gamma\_reconciliation\_search} provides the quotient
    defect and the agreement decisions.
  \item
    Replays: Example 1 has raw \texttt{2} and quotient \texttt{1} at unit weights, and raw
    \texttt{5} and quotient \texttt{3} at weights \texttt{(2,3)}. Example 2 at unit weights extends
    all four vertices, and so agrees.
  \end{itemize}
\item
  \emph{Unresolved residual:}

  \begin{itemize}
  \tightlist
  \item
    Native vertex enumeration is capped at \texttt{16} rows. Above that it returns a
    typed refusal, not a verdict.
  \item
    The CODA first layer computes the raw \texttt{Γ} only and takes no repair map, so
    the quotient \texttt{Γ} is not yet wired into it.
  \item
    The general equivalence (G)⇔(P)⇔(GV) is a written proof, exercised on the
    examples above; it is not stated in Lean.
  \item
    Zero-test agreement in Example 1 is replayed on samples. Its general form
    is the P.2 argument.
  \item
    Example 2 and the weights-\texttt{(1,3)} failure are finite replay only.
  \end{itemize}
\end{itemize}

\subsubsection{P.3 Lawful transport and an exact carrier-change criterion}\label{p.3-lawful-transport-and-an-exact-carrier-change-criterion}

For linear repair spaces \texttt{R⊆V}, \texttt{R′⊆V′}, a linear raw operation \texttt{T:V→V′}
induces \texttt{V/R→V′/R′} iff \texttt{T(R)⊆R′}. Sufficiency follows by comparing
representatives; necessity follows by applying the induced map to
\texttt{{[}r{]}={[}0{]}}. This is P.1 specialized to quotient reports.

For general reports \texttt{p:X→A}, \texttt{p′:X′→A′}, a lawful transport is the square
\texttt{p′F=F̄p}. Two such squares paste by substitution:
\texttt{p″GF=Ḡp′F=ḠF̄p}. Equal source reports remain equal after transport.
A noninjective collapse has no global left inverse: any two distinct
source values in one fiber would both equal the same recovered value.
Pasting provides neither target standing nor bicategorical coherence.

Here is a complete criterion for the \textbf{actual} fiber-sum operation on
directed integer pair tables. For a total finite map \texttt{f:X→Y}, retain all
declared target vertices, including empty fibers. Set

\[m_a=|f^{-1}(a)|,\quad s_a=\sum_{f(x)=a}g_x,\quad
(T_fe)_{ab}=\sum_{f(x)=a,\ f(y)=b}e_{xy}.\]

For the gradient \texttt{(dg)xy=gx−gy}, direct expansion gives

\[T_f(dg)_{ab}=m_b s_a-m_a s_b.\]

Its triangle sum is

\[(m_b-m_c)s_a+(m_c-m_a)s_b+(m_a-m_b)s_c.\]

\textbf{Finite carrier theorem.} \texttt{T\_f} carries every integer gradient to a
gradient iff at least one of these conditions holds: \texttt{\textbar{}Y\textbar{}≤2}; at most
one target fiber is occupied; all target fiber sizes are equal.

For at most two vertices, every antisymmetric zero-loop table is a
gradient. With at most one occupied fiber the displayed transported
gradient is zero. With common fiber size \texttt{k}, it is \texttt{d(ks)}.
For necessity suppose there are at least three target vertices, at least
two occupied fibers and unequal sizes. Some occupied \texttt{a} has two other
vertices \texttt{b,c} with \texttt{m\_b≠m\_c}. Otherwise all sizes away from each occupied
vertex agree; using two occupied vertices and a third vertex forces all
sizes to agree, a contradiction. Put a unit potential at one source
vertex over \texttt{a}. Then \texttt{s\_a=1} and all other sums vanish, so the triangle
is the nonzero integer \texttt{m\_b−m\_c}. No gradient has that cycle. This also
constructs a counterwitness for independent raw transport replay.

On a surjective map to at least three vertices, the criterion reduces
to equal fiber sizes. Empty fibers cannot be discarded to manufacture a
pass. The term-addition operation has fiber sizes \texttt{(1,…,1,0)} and descends exactly
when the old carrier has at most one vertex. Identify has one size-2
fiber and the remaining size-1 fibers and descends exactly when the old
carrier has at most three vertices. With map \texttt{(0,0,1,2)} and potential
\texttt{(0,0,1,0)}, the target cycle on \texttt{(0,1,2)} is \texttt{−1}.

Raw pair transports compose by regrouping finite sums. If both preserve
gradients, their composite does too. A root-normalized target potential
is \texttt{h\_a=T\_f(dg)\_\{a0\}}; it is well-defined on the potential modulo constants,
with the empty potential used for an empty target, and successive
normalization agrees with the composite. In particular, finite carriers
with these lawful total maps form a subcategory of finite sets: identities
have size-1 fibers and composition preserves gradients. The assignment
\texttt{X\ ↦\ Z\^{}(X×X)/\{dg\}} with arrows induced by \texttt{T\_f} is a functor to abelian
groups, since identity and composition hold already on raw pair tables.
This is a category for the selected quotient profile, rather than a claim
that every actual source edit is an arrow in it. The native owner
replays these identities exactly. This result covers fixed-writing gradient
quotients; raw counts, oriented marks, edit prices, future enabled moves
and all source-specific result equivalences require their own laws.

\subsubsection{P.4 A residual section is a chosen splitting}\label{p.4-a-residual-section-is-a-chosen-splitting}

Let \texttt{C:V→V} be linear with \texttt{C²=C}, and put \texttt{L=I−C}. Then
\texttt{L²=L}, \texttt{CL=LC=0}, \texttt{ker\ L=im\ C}, and \texttt{im\ L=ker\ C}. Direct expansion
establishes the first identities; \texttt{Cx=x} characterizes membership in \texttt{im\ C},
and \texttt{Ly=y} characterizes \texttt{ker\ C}. Consequently

\[V/\operatorname{im}C\ \simeq\ \ker C,\qquad [x]\longmapsto Lx.\]

Well-definedness uses \texttt{LC=0}; injectivity uses \texttt{ker\ L=im\ C}; surjectivity
uses \texttt{Ly=y} on \texttt{ker\ C}. If \texttt{im\ C=R1}, this represents the obstruction quotient,
and \texttt{Φ(x)\textgreater{}0} iff \texttt{Lx≠0}. Without that image equality, the link fails:
take \texttt{V=Q}, \texttt{C=I}, \texttt{d0=0}, \texttt{x=1}. Then \texttt{Lx=0} but \texttt{Φ(x)=1}.

If \texttt{TC=C′T}, expansion gives \texttt{TL=L′T}, and residual sections respect
transport. An arbitrary repair or optimizer need not satisfy that
commutation law. Choosing \texttt{C} is extra structure; an L1 quotient norm
does not choose a canonical projection. Nor does this vector-space
splitting establish A's broader ontology or persistence policy.

The dressing-field method {[}REF-REL-005{]} is the group-action analogue of a
chosen splitting. Let a finite group \texttt{G} act on configurations by \texttt{φ↦g·φ}. A
dressing field \texttt{u} is \textbf{equivariant} when \texttt{u(g·φ)=g·u(φ)}, and the dressed
configuration is \texttt{φ\^{}u=u(φ)\^{}\{-1\}·φ}. Then \texttt{(g·φ)\^{}u=u(φ)\^{}\{-1\}g\^{}\{-1\}g·φ=φ\^{}u}, so
the dressed configuration is invariant, and \texttt{u(φ\^{}u)=u(φ)\^{}\{-1\}u(φ)=e} by
equivariance, so dressing is idempotent. Both hold for nonabelian \texttt{G}. Two
printed conventions in the archived general-relativity source fail. Dressing
by \texttt{ρ(u)φ} without the inverse is not invariant even for the abelian group
\texttt{Z3}. A proposition requiring \texttt{u(g·φ)=g\^{}\{-1\}u(φ)} together with that dressing
has no solution in \texttt{S3} (it forces \texttt{(g\_1g\_2)\^{}\{-1\}=g\_1\^{}\{-1\}g\_2\^{}\{-1\}}). Its proof
step moves \texttt{ρ(g\^{}\{-1\})} past \texttt{ρ(u)}, which needs commutativity. The \href{code/relational_coincidence/README.md}{exact
checks} enumerate all \texttt{6\^{}6} maps from the
free \texttt{S3} orbit to \texttt{S3}: exactly six are equivariant, each is invariant and
idempotent, and the printed conventions fail at stated configurations. As with
\texttt{C} above, a dressing field is chosen extra structure. On a finite set an
equivariant field exists exactly when the action is free: \texttt{g·φ=φ} would give
\texttt{u(φ)=g·u(φ)} and so \texttt{g=e}, while for a free action one representative per
orbit defines \texttt{u(g·r)=g}. Nothing here concerns continuous gauge groups, where a
global choice can also fail for topological reasons.

\paragraph{P.4.a Contractive splittings in weighted L1 (DP-4, audited 2026-10-06)}\label{p.4.a-contractive-splittings-in-weighted-l1-dp-4-audited-2026-10-06}

\textbf{Setting.} Let \texttt{V=Qⁿ} carry the norm \texttt{‖x‖=Σ\_iw\_i\textbar{}x\_i\textbar{}}, with rational
\texttt{w\_i\textgreater{}0}. Let \texttt{R⊆V} be a subspace with basis matrix \texttt{B} (\texttt{n×r}). Choose \texttt{N}
(\texttt{n×(n−r)}) so that \texttt{{[}B\ N{]}} is invertible, and write its inverse in block
rows \texttt{{[}U;Q{]}}. Then \texttt{BU+NQ=I}, \texttt{UB=I}, \texttt{QN=I}, \texttt{UN=0} and \texttt{QB=0}. Put
\texttt{Φ\_R(x)=min\_\{y∈R\}‖x−y‖}, the quotient norm of P.2 with \texttt{R1=R}. For a matrix
\texttt{M}, the induced norm is

\[\|M\|=\max_j\frac{\sum_iw_i|M_{ij}|}{w_j}.\]

\textbf{Statement.}

\begin{enumerate}
\def\labelenumi{\arabic{enumi}.}
\tightlist
\item
  The idempotents \texttt{C} with \texttt{im\ C=R} are exactly \texttt{C\_H=B(U+HQ)}, for \texttt{H}
  ranging over the rational \texttt{r×(n−r)} matrices. For each,
  \texttt{L\_H=I−C\_H=(N−BH)Q} and \texttt{ker\ C\_H=im(N−BH)}.
\item
  For each such \texttt{C}, with \texttt{L=I−C}, every \texttt{x} satisfies
  \texttt{Φ\_R(x)≤‖Lx‖≤‖L‖Φ\_R(x)}.
\item
  If \texttt{R≠V}, then \texttt{‖L‖≥1}. The identity \texttt{‖Lx‖=Φ\_R(x)} for every \texttt{x} holds iff
  \texttt{‖L‖=1}; call \texttt{C} \emph{contractive} in that case. If \texttt{R=V}, then \texttt{L=0} and
  \texttt{Φ\_R=0}; this is a separate branch.
\item
  A contractive \texttt{C} exists iff the following rational linear program in
  \texttt{(H,t)} is feasible:
  \texttt{t\_ij≥(L\_H)\_ij}, \texttt{t\_ij≥−(L\_H)\_ij} and \texttt{Σ\_iw\_it\_ij≤w\_j}.
\item
  In \texttt{Q³} with unit weights and \texttt{R=span(e)}, \texttt{e=(1,1,1)}, every projection
  onto \texttt{R} has \texttt{‖L‖≥4/3}. Equality holds for \texttt{a=(1/3,1/3,1/3)}, so no
  contractive complement exists.
\item
  If \texttt{R} has codimension one, a contractive \texttt{C} exists, for every positive
  weight vector. In particular this holds for every line in weighted \texttt{ℓ1²}.
\item
  Let \texttt{T} be rational linear, and let \texttt{C} and \texttt{C′} be projections. Then
  \texttt{TC=C′T} iff \texttt{T(im\ C)⊆im\ C′} and \texttt{T(ker\ C)⊆ker\ C′}. When these hold,
  \texttt{TL=L′T}.
\end{enumerate}

\textbf{Proof.}

\emph{Item 1.} Expanding with \texttt{UB=I} and \texttt{QB=0} gives \texttt{C\_H²=C\_H}, and
\texttt{im\ C\_H=R}. Conversely, let \texttt{C} be idempotent with \texttt{im\ C=R}. Then \texttt{C=BM} for
some \texttt{M}, and \texttt{C²=C} with \texttt{B} injective gives \texttt{MB=I}. Hence
\texttt{M=M(BU+NQ)=U+(MN)Q}, which is \texttt{C\_H} with \texttt{H=MN}. The other two formulas
follow by expansion.

\emph{Item 2.} \texttt{Cx∈R} gives \texttt{Φ\_R(x)≤‖x−Cx‖}. For \texttt{y∈R}, \texttt{Ly=0} and so
\texttt{‖Lx‖=‖L(x−y)‖≤‖L‖‖x−y‖}; minimize over \texttt{y}. The column formula is the
induced norm: \texttt{‖Mx‖≤Σ\_j\textbar{}x\_j\textbar{}Σ\_iw\_i\textbar{}M\_ij\textbar{}}, with equality at the
maximizing \texttt{e\_j}.

\emph{Item 3.} A nonzero idempotent satisfies \texttt{‖L‖=‖L²‖≤‖L‖²}. If \texttt{‖L‖=1},
item 2 gives equality. Conversely, if \texttt{‖Lx‖=Φ\_R(x)≤‖x‖} for every \texttt{x}, then
\texttt{‖L‖≤1}.

\emph{Item 4.} \texttt{L\_H} is affine in \texttt{H}.

\emph{Item 5.} Every projection onto \texttt{R} has the form \texttt{C=eaᵀ} with \texttt{a\_1+a\_2+a\_3=1}.
Column \texttt{j} of \texttt{L} has norm \texttt{c\_j=\textbar{}1−a\_j\textbar{}+2\textbar{}a\_j\textbar{}}. Since \texttt{\textbar{}1−a\_j\textbar{}≥1−a\_j} and
\texttt{\textbar{}a\_j\textbar{}≥a\_j}, we get \texttt{c\_j≥1+a\_j}. So \texttt{c\_1+c\_2+c\_3≥4} and \texttt{max\_jc\_j≥4/3}.

A replayable infeasibility certificate for the LP of item 4 has thirteen
rows:

\begin{itemize}
\tightlist
\item
  \texttt{−a\_j−t\_jj≤−1} for each \texttt{j} (3 rows);
\item
  \texttt{a\_j−t\_ij≤0} for \texttt{i≠j} (6 rows);
\item
  \texttt{Σ\_it\_ij≤1} for each \texttt{j} (3 rows);
\item
  \texttt{−Σ\_ja\_j≤−1} (1 row).
\end{itemize}

Adding all thirteen gives \texttt{0≤−1}. At \texttt{a=(1/3,1/3,1/3)} every \texttt{c\_j=4/3}.
For \texttt{x=e1}, \texttt{Lx=(2/3,−1/3,−1/3)} has norm \texttt{4/3}. But \texttt{Φ\_R(e1)=1}, attained
at \texttt{y=0} and certified by the dual \texttt{z=(1,−1,0)}, which has \texttt{zᵀe=0} and
\texttt{\textbar{}z\_i\textbar{}≤1}.

\emph{Item 6.} Write \texttt{R=ker\ ψ} for a nonzero rational functional. Rescale \texttt{ψ} so
that \texttt{max\_i\textbar{}ψ\_i\textbar{}/w\_i=1}, and let \texttt{i} attain the maximum. Put
\texttt{v=sign(ψ\_i)e\_i/w\_i}. Then \texttt{‖v‖=1} and \texttt{ψ(v)=1}. Set \texttt{L=vψᵀ} and \texttt{C=I−L}.

\begin{itemize}
\tightlist
\item
  \texttt{L²=L}, \texttt{ker\ L=R}, \texttt{ψ(Cx)=0}, and \texttt{Cx=x} on \texttt{R}, so \texttt{im\ C=R}.
\item
  \texttt{‖Lx‖=\textbar{}ψ(x)\textbar{}‖v‖≤Σ\_i\textbar{}ψ\_i\textbar{}\textbar{}x\_i\textbar{}≤‖x‖}, so \texttt{‖L‖=1}.
\end{itemize}

The \texttt{4/3} obstruction of item 5 therefore needs codimension at least two, and
so \texttt{n≥3}.

\emph{Item 7.} (⇒) For \texttt{x∈im\ C}, \texttt{Tx=TCx=C′Tx∈im\ C′}. For \texttt{x∈ker\ C},
\texttt{C′Tx=TCx=0}. (⇐) Write \texttt{x=Cx+Lx}. Then \texttt{C′TCx=TCx} and \texttt{C′TLx=0}, so
\texttt{C′Tx=TCx}. In the control below, \texttt{C={[}{[}1,0{]},{[}0,0{]}{]}} and \texttt{C′={[}{[}1,1{]},{[}0,0{]}{]}}
share the image \texttt{span(e1)}. Their kernels \texttt{span(e2)} and \texttt{span(1,−1)} differ,
\texttt{T=I} fails the condition, and \texttt{Le2=(0,1)} while \texttt{L′e2=(−1,1)}.

\textbf{Worked examples.} In \texttt{Q²} with unit weights and \texttt{R=span(1,1)}:

\begin{itemize}
\tightlist
\item
  \texttt{C=(1/2)\ {[}{[}1,1{]},{[}1,1{]}{]}} gives \texttt{L=(1/2)\ {[}{[}1,−1{]},{[}−1,1{]}{]}}, with \texttt{‖L‖=1}.
\item
  \texttt{C0={[}{[}0,1{]},{[}0,1{]}{]}} is also idempotent with image \texttt{R}. It gives
  \texttt{L0={[}{[}1,−1{]},{[}0,0{]}{]}}, with \texttt{‖L0‖=1}.
\item
  For \texttt{x=(3,1)}, \texttt{Φ\_R(x)=2}, certified by \texttt{z=(1,−1)}. Then \texttt{Lx=(1,−1)} and
  \texttt{L0x=(2,0)}, both of norm \texttt{2}.
\end{itemize}

So two contractive sections represent one quotient class by different
residuals, and the choice of \texttt{C} is extra structure. The scalar \texttt{C=2} on \texttt{Q} is
not idempotent and is refused.

\textbf{Evidence grade.}

\begin{itemize}
\tightlist
\item
  \emph{Written proof:} items 1--7. Item 6 was added in this audit pass, and it
  resolves the audit's open question about lines in weighted \texttt{ℓ1²}.
\item
  \emph{Exact finite replay:}

  \begin{itemize}
  \tightlist
  \item
    \texttt{dp4\_diag\_line\_lp4\_infeasible\_H\_param}, \texttt{dp4\_thirteen\_row\_certificate},
    \texttt{dp4\_diag\_line\_attaining\_norm}, \texttt{dp4\_diag\_line\_Lx\_and\_phi};
  \item
    \texttt{dp4\_2d\_C\_idempotent\_norm\_Lx}, \texttt{dp4\_2d\_C0\_alternative}, \texttt{dp4\_2d\_phi\_x};
  \item
    \texttt{dp4\_new\_l1\_squared\_lines\_contractive} (four sampled lines in weighted
    \texttt{ℓ1²}, each feasible);
  \item
    \texttt{dp4\_zero\_quotient\_norm\_zero}, \texttt{dp4\_R\_zero\_norm\_one};
  \item
    \texttt{dp4\_transport\_control}, \texttt{dp4\_control\_C2\_not\_idempotent}.
  \end{itemize}
\item
  \emph{Lean kernel:} these cover item 5 and the planar example only. The axiom
  report is \texttt{lean4/component\_repair\_axioms.txt}.

  \begin{itemize}
  \tightlist
  \item
    \texttt{WitnessedDifference.ComponentRepair.colNorm\_ge},
    \texttt{WitnessedDifference.ComponentRepair.some\_column\_ge\_four\_thirds},
    \texttt{WitnessedDifference.ComponentRepair.uniform\_column\_norm},
    \texttt{WitnessedDifference.ComponentRepair.repaired\_example\_norm},
    \texttt{WitnessedDifference.ComponentRepair.example\_distance\_to\_R},
    \texttt{WitnessedDifference.ComponentRepair.example\_ratio};
  \item
    \texttt{WitnessedDifference.ComponentRepair.planar\_control}.
  \end{itemize}
\item
  \emph{Native exact implementation:}

  \begin{itemize}
  \tightlist
  \item
    \texttt{omega\_functional::residual\_splitting} provides the parametrized
    splitting, the weighted operator norm, and the splitting and transport
    verifiers.
  \item
    \texttt{omega\_solvers::rational\_feasibility} decides the LP of item 4 and replays
    Farkas certificates.
  \item
    \texttt{omega\_solvers::contractive\_complement} decides the distortion bound and
    whether a contractive complement exists. Pivot-budget exhaustion is a
    refusal, not a verdict.
  \item
    Replays:

    \begin{itemize}
    \tightlist
    \item
      the diagonal line in \texttt{ℓ1³} has minimum distortion \texttt{4/3}, attained;
    \item
      the bound \texttt{1333/1000} is infeasible;
    \item
      the thirteen-row certificate replays, and a copy with its last
      multiplier zeroed is rejected;
    \item
      the planar diagonal has a contractive complement.
    \end{itemize}
  \end{itemize}
\item
  \emph{Unresolved residual:}

  \begin{itemize}
  \tightlist
  \item
    The \texttt{omega\_wdp} projection bounds do not yet call the contractive
    complement decision.
  \item
    The Lean module takes the induced norm to be the largest column norm
    without formalising that identity.
  \item
    It treats only rank-one projections onto the diagonal line.
  \item
    Items 1--4, 6 and 7 are not in Lean.
  \item
    Item 6 has no exact replay beyond the four samples.
  \end{itemize}
\end{itemize}

\subsubsection{P.5 Overlaps, comparison spans and retained histories}\label{p.5-overlaps-comparison-spans-and-retained-histories}

Pairs \texttt{(a,b)} and \texttt{(b′,c)} glue to a triple with both projections iff
\texttt{b=b′}; when they do, \texttt{(a,b,c)} is the unique triple. This is the smallest
fiber-product argument. For spans \texttt{X←M→Y} and \texttt{Y←N→Z}, composition uses

\[M\times_YN=\{(m,n):\operatorname{right}(m)=\operatorname{left}(n)\}.\]

Matching pairs have the evident outer endpoints. An incompatible overlap
supplies no member. Iterated fiber products have the canonical tuple
rebracketing bijection; a general enriched comparison category still needs
its declared morphisms, identities and coherence data. The inherited
faithful point-span groupoid embedding claim is not reinstated: a
one-point span loses distinct arrows with the same endpoints.

The layered-authority source defines each composite comparison price to be the sum of its two
declared prices. Associativity of that sum follows from associativity of
addition. It is not a theorem about an optimized minimum price. Erasing
a trace can preserve endpoints and span composition while losing steps,
work and the intermediate raw states. Distinct histories can have the
same terminal state and different future permissions or accumulated cost.
Thus those protected queries must be tested against P.1 before erasure.

An empty support sieve, a missing terminalizer, an invalid leg or a
mismatched target blocks the standing predicate that requires those fields.
This is a consequence of a declared predicate. Finite overlap checks and
these non-creation results do not construct a general sheaf, global section,
omniversal comparison object or higher cohomological optimizer.

\paragraph{P.5.a Finite function gluing and the presheaf boundary (DP-6, audited 2026-10-06)}\label{p.5.a-finite-function-gluing-and-the-presheaf-boundary-dp-6-audited-2026-10-06}

\textbf{Statement.} Let \texttt{X} be a finite set and \texttt{(U\_i)\_\{i∈I\}} a finite family of
declared subsets of \texttt{X}. Let \texttt{B} be a finite output alphabet, and let
\texttt{s\_i:U\_i→B} be total function tables. Put \texttt{U=⋃\_iU\_i}. The following are
equivalent:

\begin{enumerate}
\def\labelenumi{\arabic{enumi}.}
\tightlist
\item
  \emph{(compatibility)} \texttt{s\_i(x)=s\_j(x)} for every \texttt{i,j} and every
  \texttt{x∈U\_i∩U\_j};
\item
  \emph{(unique glue)} there is exactly one \texttt{s:U→B} with \texttt{s\textbar{}U\_i=s\_i} for every
  \texttt{i}.
\end{enumerate}

A global report on a declared target \texttt{X} additionally requires \texttt{U=X}. When
\texttt{U≠X}, the missing points are reported as \texttt{IncompleteCoverage}, and gluing on
the smaller union is a separate problem. The case \texttt{I=∅} is included: then
\texttt{U=∅}, and the unique glue is the empty function.

For a general finite presheaf, the question is different. Let \texttt{F} be a
presheaf on a finite site, validated for totality, identities and
composition. Fix a cover of \texttt{X}, let \texttt{M} be its set of matching families, and
let \texttt{r:F(X)→M} be the restriction map. The extensions of a family \texttt{m∈M} are
exactly the fiber \texttt{r⁻¹(m)}. That fiber can have zero, one or more elements,
and pairwise compatibility does not decide which.

\textbf{Proof.}

\emph{(1)⇒(2).} Assume compatibility, and let
\texttt{G=⋃\_i\{(x,s\_i(x)):x∈U\_i\}}. The domain of \texttt{G} is \texttt{U}.

\begin{itemize}
\tightlist
\item
  \texttt{G} is a graph of a function. Suppose \texttt{(x,b)∈G} and \texttt{(x,c)∈G}. Then there
  are \texttt{i,j} with \texttt{x∈U\_i∩U\_j}, \texttt{b=s\_i(x)} and \texttt{c=s\_j(x)}, and compatibility
  gives \texttt{b=c}. So \texttt{G} is the graph of some \texttt{s:U→B}.
\item
  Restriction recovers every table. For \texttt{x∈U\_i}, \texttt{G} contains
  \texttt{(x,s\_i(x))}, so \texttt{s(x)=s\_i(x)}.
\item
  The glue is unique. Let \texttt{t:U→B} restrict to every \texttt{s\_i}. Each \texttt{x∈U} lies
  in some \texttt{U\_i}, and then \texttt{t(x)=s\_i(x)=s(x)}.
\end{itemize}

\emph{(2)⇒(1).} At \texttt{x∈U\_i∩U\_j}, \texttt{s\_i(x)=s(x)=s\_j(x)}. This direction uses
existence only. No step selects an element, so the empty union is covered.

\emph{Presheaf fibers.} By functoriality, every global restricts to a matching
family, so \texttt{r} is well defined. From the definitions of fiber, onto and
one-to-one:

\begin{itemize}
\tightlist
\item
  every matching family extends iff \texttt{r} is onto;
\item
  extensions are unique whenever they exist iff \texttt{r} is one-to-one;
\item
  the sheaf condition for this cover is both.
\end{itemize}

\textbf{Worked example.} Take \texttt{X=\{a,b,c\}}, \texttt{U\_0=\{a,b\}}, \texttt{U\_1=\{b,c\}} and
\texttt{B=\{0,1/2,1\}}. Set \texttt{s\_0(a)=0}, \texttt{s\_0(b)=1/2}, \texttt{s\_1(b)=1/2} and \texttt{s\_1(c)=1}.
The only overlap comparison is \texttt{1/2=1/2}. The glue is
\texttt{\{(a,0),(b,1/2),(c,1)\}}, and every global extending both tables takes these
values.

Mutating \texttt{s\_1(b)} to \texttt{1} gives the incompatibility witness \texttt{(b,1/2,1)}. Its
discrepancy is \texttt{1−1/2=1/2}.

\textbf{Countermodels for the dropped hypotheses.}

\begin{itemize}
\tightlist
\item
  \emph{Coverage omitted.} Let \texttt{X=\{a,b\}} with only \texttt{U\_0=\{a\}}, \texttt{s\_0(a)=0} and
  \texttt{B=\{0,1\}}. The glue on \texttt{\{a\}} is unique, but \texttt{(a↦0,b↦0)} and
  \texttt{(a↦0,b↦1)} both extend it on \texttt{X}.
\item
  \emph{Local totality omitted.} Let \texttt{X=U\_0=\{a\}} with an empty table. Both \texttt{a↦0}
  and \texttt{a↦1} complete it, so no local function on \texttt{U\_0} is specified.
\item
  \emph{Existence fails for a presheaf.} Use the category of the four subsets of
  \texttt{\{a,b\}} under inclusion.

  \begin{itemize}
  \tightlist
  \item
    Put \texttt{F(∅)=\{*\}} and \texttt{F(V)=\{0,1\}} for every nonempty \texttt{V}.
  \item
    Restrictions between nonempty sets are identities, and restrictions to
    \texttt{∅} are constant.
  \item
    Every functor law holds by direct evaluation.
  \item
    Cover \texttt{X} by \texttt{\{a\}} and \texttt{\{b\}}, with locals \texttt{0} and \texttt{1}. They match at \texttt{∅}.
  \item
    The globals restrict to \texttt{(0,0)} and \texttt{(1,1)}, neither of which is \texttt{(0,1)}.
    So the fiber is empty.
  \end{itemize}
\item
  \emph{Uniqueness fails for a presheaf.} On the same category, put
  \texttt{F(X)=\{g0,g1\}} and \texttt{F(V)=\{*\}} for every proper \texttt{V}. The family \texttt{(*,*)} is
  matching and has the two extensions \texttt{g0} and \texttt{g1}.
\item
  \emph{The empty intersection is skipped.} Put \texttt{F(V)=\{0,1\}} for all four subsets,
  \texttt{∅} included, with identity restrictions. The disjoint locals \texttt{0} and \texttt{1}
  restrict to unequal values at \texttt{∅}, so they do not match. A check that skips
  empty overlaps would accept them. For function tables this cannot happen,
  since there is exactly one function on \texttt{∅}.
\item
  \emph{Erased histories.} Take histories \texttt{h0,h1} with the same endpoint
  \texttt{p(h0)=p(h1)}.

  \begin{itemize}
  \tightlist
  \item
    Their costs differ: \texttt{c(h0)=0} and \texttt{c(h1)=1}.
  \item
    Their future permissions differ: \texttt{u(h0)=false} and \texttt{u(h1)=true}.
  \item
    By P.1 each query has a collision, so neither factors through the
    endpoint.
  \end{itemize}

  Erasure is valid for a specified query exactly when that query is constant
  on endpoint fibers.
\end{itemize}

\textbf{Evidence grade.}

\begin{itemize}
\tightlist
\item
  \emph{Written proof:} the argument above.
\item
  \emph{Exact finite replay:}

  \begin{itemize}
  \tightlist
  \item
    \texttt{dp6\_function\_glue}, \texttt{dp6\_function\_incompatible};
  \item
    \texttt{dp6\_empty\_cover\_empty\_function};
  \item
    \texttt{dp6\_existence\_presheaf\_functor\_laws}, \texttt{dp6\_existence\_presheaf\_fibers};
  \item
    \texttt{dp6\_uniqueness\_presheaf\_functor\_laws}, \texttt{dp6\_uniqueness\_presheaf\_fibers};
  \item
    \texttt{dp6\_empty\_intersection\_family\_01};
  \item
    \texttt{dp6\_history\_query\_collisions}.
  \end{itemize}
\item
  \emph{Lean kernel:} the axiom report is \texttt{lean4/component\_reports\_axioms.txt}.

  \begin{itemize}
  \tightlist
  \item
    \texttt{WitnessedDifference.ComponentReports.glue\_exists\_unique},
    \texttt{WitnessedDifference.ComponentReports.compatible\_of\_restricts},
    \texttt{WitnessedDifference.ComponentReports.not\_compatible\_iff\_witness},
    \texttt{WitnessedDifference.ComponentReports.empty\_cover\_glues};
  \item
    \texttt{WitnessedDifference.ComponentReports.CoverPresheaf.restriction\_matches},
    \texttt{WitnessedDifference.ComponentReports.CoverPresheaf.every\_matching\_extends\_iff};
  \item
    \texttt{WitnessedDifference.ComponentReports.existence\_failure\_zero\_extensions},
    \texttt{WitnessedDifference.ComponentReports.uniqueness\_failure\_two\_extensions}.
  \end{itemize}
\item
  \emph{Native exact implementation:}

  \begin{itemize}
  \tightlist
  \item
    \texttt{omega\_sheaf\_theory::finite\_report\_gluing} glues on the union and on the
    declared target, verifies function glues, and classifies presheaf
    gluing as a unique glue, no extension, or ambiguous. It also provides the
    history-erasure decision.
  \item
    The \texttt{omega\_wdp} presheaf descent wraps the existing fiber enumeration.
    It replays fibers of size \texttt{0}, \texttt{1} and more than one.
  \end{itemize}
\item
  \emph{Unresolved residual:}

  \begin{itemize}
  \tightlist
  \item
    The existing general sheaf equalizer routine still skips empty
    intersections. A test documents this, and the routine is unchanged.
  \item
    The native presheaf carrier is capped at \texttt{64} points. The theorem is not a
    claim that every finite size is admitted.
  \item
    The history countermodel is finite replay and written proof only.
  \end{itemize}
\end{itemize}

\subsubsection{P.6 Completion, common targets and independent standing}\label{p.6-completion-common-targets-and-independent-standing}

Let \texttt{P=p(X)} carry a declared metric, and let \texttt{B} be complete. After P.1,
the image query \texttt{q̄:P→B} has a unique uniformly continuous extension to
the metric completion exactly when it is uniformly continuous on \texttt{P}.
Equivalently, for every positive epsilon there is a positive delta such
that \texttt{d(p(x),p(y))\textless{}delta} implies \texttt{d(q(x),q(y))\textless{}epsilon}, for all source
states. Extend by limits of Cauchy representatives; uniform continuity
makes the image sequence Cauchy and the limit independent of the
representative. Restriction establishes necessity and density establishes uniqueness.

This is the actual CS-to-Q interface. Finite factorization alone is
insufficient, and source continuity for an unrelated metric does not
replace relative continuity. Appendix O.1 supplies the formal selected
metric criterion; O.2 retains the separate circle and discrete dyadic
dual topologies. O.4--O.9 bind model-specific common targets and error
moduli. No primitive mark chooses those metrics, physical models or limits.

If every interval \texttt{{[}l\_n,u\_n{]}} contains the same real \texttt{b} and its width
tends to zero, its midpoint error is at most \texttt{(u\_n−l\_n)/2}. Interval
order and shrinking widths without common-target membership do not give
that conclusion. Q's rational interval stage is not a construction of an
unknown enclosed point. Completing a field of scalars likewise supplies
neither an infinite operator domain nor interacting field dynamics.

Finally, operational descent standing consists of an actual terminal \texttt{t}, a
pipeline equality to that terminal and its seven evidence-derived guards.
Regime standing has the different form

\[\operatorname{Standing}_R(g,p) \Longleftrightarrow\
(\exists e.\operatorname{verifies}_R(p,e))\ \land\
\operatorname{recordedGrant}_R(g,p).\]

Assume verification implies truth and a recorded grant implies entitlement.
Applying those two assumptions to the standing witnesses establishes truth and
entitlement. This is a sufficient soundness theorem, not a proof of those
assumptions or an equivalence recovering them. Constant-success verification
and grant predicates with false truth give an immediate countermodel
without the assumptions. An interpretation relating operational and regime
standing is another obligation.

Refusal propagation and terminal-history reconstruction follow by unfolding
the partial pipeline; disjoint terminal constructors follow by case analysis.
These elementary laws are useful but do not prove general Kleisli
factorization, conservation of a tower or institutional authorization.
Python's finite application verifies a configured one-use HMAC grant;
the primary Rust descent stage currently refuses supplied grant strings because
its external verifier is unavailable. Neither result selects the grantor's
institutional standing. Proposal, mathematical evaluation, evidence and
authorization therefore retain distinct interfaces throughout the compression.

\paragraph{P.6.a Finite enclosure verification records and an exact square-root enclosure (DP-5, audited 2026-10-06)}\label{p.6.a-finite-enclosure-verification-records-and-an-exact-square-root-enclosure-dp-5-audited-2026-10-06}

\textbf{Correction.} The candidate asked for ``a nested-without-common-target
countermodel''. The audit read it as: a nested family of nonempty, closed,
bounded real intervals whose intersection is empty. No such family exists, so
that reading of the request is refuted. \textbf{DP-5′ supersedes the candidate
statement.} It keeps the theorem below and replaces the requested example
with controls that each change exactly one premise or the ambient field.

\textbf{Statement.}

\begin{enumerate}
\def\labelenumi{\arabic{enumi}.}
\tightlist
\item
  \emph{(Finite midpoint bound.)} Let \texttt{{[}l\_j,u\_j{]}}, \texttt{0≤j≤N}, be a nonempty prefix
  with rational endpoints. Let \texttt{b} be a declared target with replayable
  evidence \texttt{l\_j≤b≤u\_j} for every submitted \texttt{j}. Then
  \texttt{\textbar{}b−(l\_j+u\_j)/2\textbar{}≤(u\_j−l\_j)/2} for every \texttt{j≤N}. If \texttt{u\_j−l\_j≤2ε}, the
  midpoint has error at most \texttt{ε}. Nesting is not needed here.
\item
  \emph{(Common point.)} Every nested sequence of nonempty, closed, bounded real
  intervals has a common point. If the widths tend to zero, that point is
  unique.
\item
  \emph{(Finite intersection.)} A nonempty finite family of rational intervals has a
  common point iff \texttt{max\_jl\_j≤min\_ju\_j}. This finds an unnamed common point.
  It does not identify a separately designated target.
\item
  \emph{(Exact square-root enclosure.)} Let \texttt{a=p/q≥0} be reduced with \texttt{q\textgreater{}0}, and
  let \texttt{δ\textgreater{}0} be rational.

  \begin{itemize}
  \tightlist
  \item
    If \texttt{a=0}, return \texttt{{[}0,0{]}}.
  \item
    Otherwise put \texttt{l\_0=0} and \texttt{u\_0=H=max(1,a)}. Bisect at \texttt{m=(l+u)/2}:
    return \texttt{{[}m,m{]}} if \texttt{m²=a}; replace \texttt{l} by \texttt{m} if \texttt{m²\textless{}a}; replace \texttt{u} by
    \texttt{m} if \texttt{m²\textgreater{}a}.
  \item
    Each comparison is an integer comparison: for \texttt{m=s/t}, compare \texttt{s²q}
    with \texttt{pt²}.
  \item
    Every stage satisfies \texttt{0≤l\_n≤u\_n}, \texttt{l\_n²≤a≤u\_n²},
    \texttt{{[}l\_\{n+1\},u\_\{n+1\}{]}⊆{[}l\_n,u\_n{]}} and \texttt{u\_n−l\_n≤H/2\^{}n}.
  \item
    The procedure reaches width \texttt{≤δ} within the least \texttt{n} with \texttt{2\^{}nδ≥H}.
  \item
    Taking \texttt{δ=2ε} gives a midpoint with error at most \texttt{ε} from the
    nonnegative root.
  \end{itemize}
\item
  \emph{(Prefix limit.)} A checked finite prefix does not certify a modulus
  \texttt{n≥M(ε)⇒u\_n−l\_n≤2ε} for later indices. It also does not certify membership
  of a target that was never witnessed.
\end{enumerate}

\textbf{Proof.}

\emph{Item 1.} \texttt{l≤b≤u} and \texttt{m=(l+u)/2} give \texttt{−(u−l)/2≤b−m≤(u−l)/2}.

\emph{Item 2.} Let \texttt{b=sup\_nl\_n}. This exists, since \texttt{u\_0} bounds every \texttt{l\_n}.
Fix \texttt{j}. For \texttt{n≥j}, nesting gives \texttt{l\_n≤u\_n≤u\_j}. For \texttt{n\textless{}j},
\texttt{l\_n≤l\_j≤u\_j}. So \texttt{b≤u\_j}, and \texttt{l\_j≤b} by the definition of \texttt{b}. Two common
points are at most every width apart, so if the widths tend to zero they
coincide.

\emph{Item 3.} A common point lies in \texttt{{[}max\ l,min\ u{]}}. Conversely, if
\texttt{max\ l≤min\ u}, then \texttt{max\ l} is a common point.

\emph{Item 4.}
- \emph{Base case.} If \texttt{a≤1}, then \texttt{u\_0=1} and \texttt{1≥a}. If \texttt{a≥1}, then \texttt{a²≥a}.
Either way \texttt{l\_0²=0≤a≤u\_0²}.
- \emph{Step.} The step keeps whichever half preserves \texttt{l²≤a≤u²}, and it halves
the width. Order is preserved because \texttt{l≤m≤u}.
- \emph{Common root.} Squaring is increasing on nonnegative numbers. So each
interval contains the nonnegative root, and only that root.
- \emph{Termination.} Repeated integer doubling of \texttt{δ} finds the least \texttt{n} with
\texttt{2\^{}nδ≥H}.

Nonnegativity is needed. With \texttt{a=1} and \texttt{l=u=−1}, \texttt{l²≤a≤u²} holds, but the
interval excludes the positive root.

\emph{Item 5.} After any prefix \texttt{{[}0,2\^{}\{−j\}{]}}, \texttt{j≤N}, extend in two ways:

\begin{itemize}
\tightlist
\item
  continue the same rule, which gives widths \texttt{2\^{}\{−n\}} and target \texttt{0};
\item
  set every later interval to \texttt{{[}2,3{]}}.
\end{itemize}

The two extensions agree on the inspected prefix. The second fails both the
target obligation and the modulus obligation.

\textbf{Worked table (\texttt{a=2}).}

\begin{longtable}[]{@{}lrrrl@{}}
\toprule\noalign{}
\texttt{n} & \texttt{l\_n} & \texttt{u\_n} & width & decisive comparison \\
\midrule\noalign{}
\endhead
\bottomrule\noalign{}
\endlastfoot
0 & 0 & 2 & 2 & \texttt{0≤2≤4} \\
1 & 1 & 2 & 1 & \texttt{1\textless{}2} \\
2 & 1 & 3/2 & 1/2 & \texttt{9/4\textgreater{}2} \\
3 & 5/4 & 3/2 & 1/4 & \texttt{25/16\textless{}2} \\
4 & 11/8 & 3/2 & 1/8 & \texttt{121/64\textless{}2} \\
5 & 11/8 & 23/16 & 1/16 & \texttt{529/256\textgreater{}2} \\
6 & 45/32 & 23/16 & 1/32 & \texttt{2025/1024\textless{}2048/1024} \\
\end{longtable}

The last midpoint is \texttt{91/64}, with error at most \texttt{1/64}.

\begin{itemize}
\tightlist
\item
  \emph{Tamper.} Replacing the final lower endpoint by \texttt{23/16} is rejected,
  because \texttt{529/256\textgreater{}2}.
\item
  \emph{Exact root.} For \texttt{a=9/16}, the second midpoint is \texttt{3/4}, and the procedure
  returns \texttt{{[}3/4,3/4{]}}.
\item
  \emph{Zero.} \texttt{a=0} returns \texttt{{[}0,0{]}}.
\end{itemize}

\textbf{Controls for DP-5′.} Each changes one premise of item 2 or item 1.

\begin{itemize}
\tightlist
\item
  \emph{Incomplete field.} The table above, read in \texttt{Q}, consists of nonempty
  closed sets \texttt{{[}l\_n,u\_n{]}∩Q} that are nested and shrinking. Their intersection
  is empty.

  \begin{itemize}
  \tightlist
  \item
    A common rational \texttt{q} would satisfy \texttt{l\_n²≤q²≤u\_n²} for every \texttt{n}. Also
    \texttt{u\_n²−l\_n²=(u\_n+l\_n)(u\_n−l\_n)≤4H/2\^{}n} tends to zero (the Archimedean step).
    So \texttt{q²=2}.
  \item
    For reduced \texttt{q=m/n}, \texttt{m²=2n²} makes \texttt{m} even, then \texttt{n} even, a
    contradiction.
  \end{itemize}
\item
  \emph{Closedness dropped.} \texttt{(0,1/(n+1))} has empty intersection.
\item
  \emph{Target membership dropped.} \texttt{{[}1,1+2\^{}\{−n\}{]}} is closed, nested and
  shrinking, but it misses the designated \texttt{b=0}, and every midpoint error is
  above \texttt{1}.
\item
  \emph{Nesting dropped.} \texttt{{[}n,n+1/(n+1){]}} has no common point.
\item
  \emph{Shrinking dropped.} \texttt{{[}0,2{]}} at every index keeps midpoint error \texttt{1}.
\item
  \emph{Boundedness dropped.} \texttt{{[}n,∞)} is nested, closed and nonempty, with empty
  intersection.
\end{itemize}

\textbf{Evidence grade.}

\begin{itemize}
\tightlist
\item
  \emph{Written proof:} items 1--5 and every control.
\item
  \emph{Exact finite replay:}

  \begin{itemize}
  \tightlist
  \item
    \texttt{dp5\_sqrt2\_table}, \texttt{dp5\_midpoint\_and\_error},
    \texttt{dp5\_invariants\_width\_nesting},
    \texttt{dp5\_long\_run\_invariant\_a2\_delta\_2\_pow\_minus40};
  \item
    \texttt{dp5\_termination\_bound\_a1/2\_d1/100}, \texttt{dp5\_termination\_bound\_a2\_d1/32},
    \texttt{dp5\_termination\_bound\_a7\_d1/10};
  \item
    \texttt{dp5\_exact\_root\_9\_16}, \texttt{dp5\_zero\_radicand};
  \item
    \texttt{dp5\_tamper\_lower\_endpoint\_rejected}, \texttt{dp5\_negative\_endpoint\_control};
  \item
    \texttt{dp5\_no\_rational\_sqrt2\_small\_denominators};
  \item
    \texttt{dp5\_finite\_intersection\_max\_lo\_le\_min\_hi};
  \item
    \texttt{dp5\_prefix\_cannot\_separate\_tails}.
  \end{itemize}
\item
  \emph{Lean kernel:} the axiom report is \texttt{lean4/component\_completion\_axioms.txt}.

  \begin{itemize}
  \tightlist
  \item
    \texttt{WitnessedDifference.ComponentCompletion.inv\_start},
    \texttt{WitnessedDifference.ComponentCompletion.inv\_step},
    \texttt{WitnessedDifference.ComponentCompletion.inv\_iter};
  \item
    \texttt{WitnessedDifference.ComponentCompletion.width\_step},
    \texttt{WitnessedDifference.ComponentCompletion.step\_nested},
    \texttt{WitnessedDifference.ComponentCompletion.width\_iter};
  \item
    \texttt{WitnessedDifference.ComponentCompletion.root\_in\_bracket},
    \texttt{WitnessedDifference.ComponentCompletion.mid\_close},
    \texttt{WitnessedDifference.ComponentCompletion.negative\_endpoint\_control};
  \item
    \texttt{WitnessedDifference.ComponentCompletion.start\_two},
    \texttt{WitnessedDifference.ComponentCompletion.table\_row1} through
    \texttt{WitnessedDifference.ComponentCompletion.table\_row6}, and
    \texttt{WitnessedDifference.ComponentCompletion.table\_iter6};
  \item
    \texttt{WitnessedDifference.ComponentCompletion.sq\_mod\_two},
    \texttt{WitnessedDifference.ComponentCompletion.nat\_sq\_two\_descent},
    \texttt{WitnessedDifference.ComponentCompletion.no\_rational\_sqrt\_two}.
  \end{itemize}
\item
  \emph{Native exact implementation:}

  \begin{itemize}
  \tightlist
  \item
    \texttt{omega\_kernel::square\_root\_enclosure} provides the bisection, trace
    replay, the step count for a width, finite-prefix verification and the
    finite common point, with a designated target.
  \item
    At its resource limit it returns the last certified enclosure, never the
    requested precision.
  \end{itemize}
\item
  \emph{Unresolved residual:}

  \begin{itemize}
  \tightlist
  \item
    The Archimedean step of the incomplete-field control is a written proof
    and is not formalised.
  \item
    Item 2 is the real-analysis argument above. It is not stated in Lean.
  \item
    A modulus for an arbitrary generator must come with its own proof. The
    square-root generator is the only one supplied here.
  \end{itemize}
\end{itemize}

\subsubsection{P.7 Complete balance-preservation rules for the five permitted edits}\label{p.7-complete-balance-preservation-rules-for-the-five-permitted-edits}

The model in P.7--P.21 has a finite term set, two nonnegative integer
relation-count tables and oriented filled triples. Its five permitted edits
identify two terms, remove one assertion, fill a triple, add a term, or add
one assertion. This is the model specified in the
\href{https://github.com/dbccpage/H3V7N5QF/blob/5d9488324ce976817987bd83a28d901371cfc073/paper/twentyfifth/paper.md}{twenty-fifth source derivation}
and implemented in
\href{https://github.com/dbccpage/H3V7N5QF/blob/5d9488324ce976817987bd83a28d901371cfc073/code/twentyfifth/state.py}{\texttt{twentyfifth/state.py}}.
The source number is only an archive index. The orientation convention,
identification of filled triples and ordering of offered edits are explicitly
chosen readings in that implementation, not consequences of the initial
inscription or external standards. The statements below use those stated
choices as hypotheses.

This is a fixed-context bounded statement about the actual directed five-edit model
transition system. Fix its ordered carrier \texttt{0,…,n−1} and its sorted set \texttt{F} of
canonical cyclic triples. Both directions and loops remain independent entries.
For integer excess \texttt{e=r−s}, put \texttt{C\_(a,b,c)(e)=e\_ab+e\_bc+e\_ca}. Write \texttt{D\_F} for
the matrix with these cycle rows; balance means \texttt{D\_F\ e=0}. Every integer table
is realizable by nonnegative raw counts, using its positive and negative parts.
Adding the same positive offset to both tables leaves excess unchanged and
enables any selected removal. A marking records an equation to check; the source
kernel permits an unbalanced marking and does not solve it during \texttt{Fill}.

Call an enabled event \textbf{universally balance-preserving in this context} if
every balanced raw state with this carrier and face set, on which the same event
is enabled, has a balanced actual output. This quantifier differs from balance
of the particular supplied input/output and from descent of all gradient classes.

\begin{longtable}[]{@{}
  >{\raggedright\arraybackslash}p{(\linewidth - 2\tabcolsep) * \real{0.5000}}
  >{\raggedright\arraybackslash}p{(\linewidth - 2\tabcolsep) * \real{0.5000}}@{}}
\toprule\noalign{}
\begin{minipage}[b]{\linewidth}\raggedright
Actual edit
\end{minipage} & \begin{minipage}[b]{\linewidth}\raggedright
Necessary and sufficient condition
\end{minipage} \\
\midrule\noalign{}
\endhead
\bottomrule\noalign{}
\endlastfoot
\texttt{AddTerm} & Always, for the existing carried face set. \\
\texttt{AddAssertion} or \texttt{Remove} at directed \texttt{(x,y)}, in either relation & No face of \texttt{F} uses that directed arc. \\
\texttt{Identify(x,y)} with \texttt{x\textless{}y} & No surviving output face contains the merged vertex \texttt{x}. \\
\texttt{Fill(t)} with \texttt{t} absent from \texttt{F} & The row \texttt{C\_t} belongs to the rational row span of \texttt{D\_F}. \\
\end{longtable}

\textbf{Proof for AddTerm.} The actual operation preserves the old table entries and
old faces, pads incident entries with zero and adds no face. Thus every old cycle
sum is unchanged. This holds even when P.3 says that the same operation fails to
descend on the full gradient quotient: the new carrier has additional coordinates
whose quotient semantics differ from the old face equations.

\textbf{Proof for assertion changes.} The excess changes by \texttt{δ=±E\_xy}: plus for adding
R or removing S, minus for adding S or removing R. Each old cycle changes by
\texttt{C\_t(δ)}. All such changes vanish exactly when the directed arc occurs in no
filled face. Necessity is witnessed by zero excess, realized with common count
one on every coordinate so every removal is enabled. A face using the arc then
has output value \texttt{+1} or \texttt{−1}. The reverse arc is a different coordinate and loops
are absent from triples, so neither is silently identified with the tested arc.

\textbf{Proof for Identify.} The actual endpoint map is \texttt{f(y)=x} and
\texttt{f(v)=v−{[}v\textgreater{}y{]}} otherwise. It aggregates both count tables over endpoint fibers,
carries faces by cyclic normalization, drops degeneracies and merges duplicates.
If no surviving face meets \texttt{x}, each endpoint of every surviving face has a
singleton preimage. Its cycle therefore equals that of an old filled face and
remains zero. Conversely choose a surviving cyclic face \texttt{(a,b,c)} meeting \texttt{x}
and rotate so \texttt{a=x}. Its three fiber sizes are \texttt{(2,1,1)}. Let the old potential
be one at the unique preimage of \texttt{b} and zero elsewhere, and let \texttt{e=dg}.
Every old cycle vanishes. P.3's actual aggregate formula gives cycle
\texttt{(m\_b\ s\_a−m\_a\ s\_b)+(m\_c\ s\_b−m\_b\ s\_c)+(m\_a\ s\_c−m\_c\ s\_a)=−1}.
Positive/negative parts with common count one give a legal balanced raw source
and an actual unbalanced identified output. This argument does not require any
antisymmetry law for arbitrary source tables.

\textbf{Proof for Fill.} The operation keeps both tables and adds just the new
diagnostic row. If \texttt{C\_t=cᵀD\_F}, every old balanced table satisfies the new equation.
If no such rational coefficient vector exists, exact linear separation of
\texttt{D\_Fᵀc=C\_tᵀ} gives a rational vector \texttt{w} with \texttt{D\_F\ w=0} and \texttt{C\_t\ w≠0}.
Clear all denominators by their positive product to obtain an integer table
with the same properties. Realize it as raw counts. The selected filling is
enabled, its old faces are balanced, and its newly filled face is unbalanced.
Thus rational row-span membership is equivalent even when all source excess
tables are integral. Clearing denominators is exact and need not be minimal.

\textbf{A directed tetrahedron falsifier.} With old faces \texttt{013,023,123}, set only
\texttt{e\_20=1}. Every old cycle is zero, but filling \texttt{012} gives value one. The familiar
three-face implication for antisymmetric cochains does not apply to this raw
table. With both orientations retained, the actual identity is

\[C_{012}-C_{021}-C_{013}+C_{031}+C_{023}-C_{032}-C_{123}+C_{132}=0.\]

Hence seven directed tetrahedron face equations imply the eighth. This identity
is checked coordinate by coordinate; it does not add an antisymmetry assumption.

The generic native owner \texttt{linear\_algebra::row\_constraint\_implication} delegates
elimination to the existing exact-field solver and independently replays either
\texttt{cᵀD\_F=C\_t} or \texttt{D\_Fw=0,\ C\_tw≠0}. The WDP adapter
\texttt{genesis\_balance\_transition::Source25BalanceProfile} retains the actual source,
event and target, reports their separate balance values, and supplies a complete
checked raw counterstate for a negative carrier-wide verdict. It never substitutes
that counterstate for the supplied input. Dense Fill allocation/solver budgets
refuse explicitly; arbitrary integer bit lengths and total source execution
cost are outside those budgets.

The written classification above covers all five cases and arbitrary finite
carriers. The new Lean module checks selected cycle, update and directed
tetrahedron laws; it does not formalize the complete all-case necessity theorem
or verify the Rust compiler. Exhaustive native kernel replay covers all 263
marked contexts through four terms and their 19,388 enabled edits. That finite
evidence is separate from the written carrier-wide proof. These rules close one
dynamic balance question. General higher fillers, global site/HoTT semantics,
physical limits and institutional soundness in GEN-WDP-O01--O05 remain open.

\subsubsection{P.8 Affine balance paths and their exact coherence criterion}\label{p.8-affine-balance-paths-and-their-exact-coherence-criterion}

P.7 classifies one enabled edit at a time. A path needs a second question:
whether its \textbf{endpoint} is balanced for every balanced input, even when an
intermediate state is not. The following result answers that question at the
integer-excess layer without identifying a raw history with its endpoint.

For a declared five-edit model context \texttt{C=(X,F)}, let
\texttt{E\_C=Z\^{}(X×X)} retain every directed coordinate and loop, let \texttt{D\_C} be its
integer marked-cycle matrix, and put \texttt{K\_C=ker\_Z\ D\_C}. A context change has an
integer affine excess expression \texttt{F(e)=Ae+b}, where \texttt{A:E\_C→E\_C\textquotesingle{}} is an integer
matrix. The actual five edit forms have exactly these expressions: AddTerm is
zero padding, Identify is the pair-fiber sum of P.3, Fill leaves \texttt{e} fixed,
and AddAssertion/Remove have \texttt{A=I} and \texttt{b=±E\_xy} with the signs in P.7.
The expression is a statement about excess; the original raw counts, edit
enablement, occurrences, price and ordered history remain separately typed.

\textbf{Affine balance theorem.} The expression sends every balanced integer
excess table to a balanced target table,
\texttt{F(K\_C)⊆K\_C\textquotesingle{}}, if and only if

\[D_{C'}b=0\quad\text{and}\quad
 D_{C'}A=\Lambda D_C\quad\text{for some rational matrix }\Lambda.\]

The displayed pair is a finite positive certificate. The constant condition
is necessary by applying \texttt{F} to zero. After it is imposed, preservation says
that each row of \texttt{D\_C\textquotesingle{}A} annihilates \texttt{ker\_Z\ D\_C}. The integer kernel spans
\texttt{ker\_Q\ D\_C}: every rational kernel vector becomes integral after clearing
denominators. Over \texttt{Q}, the row space lies in the annihilator of the kernel;
both have dimension \texttt{rank\ D\_C}, so they coincide and give \texttt{Λ}. The converse
follows by substitution. If the constant condition fails, \texttt{e=0} is a
negative witness. If the row relation
fails, exact rational elimination gives \texttt{w∈ker\_Q\ D\_C} with
\texttt{D\_C\textquotesingle{}Aw≠0}; clearing denominators gives a replayable integer witness. Its
positive and negative parts are legal nonnegative assertion tables. For an
individual actual edit, adding a common sufficiently large count to both
relations also enables a selected Remove without changing \texttt{e}; therefore
P.7's all-states quantifier over enabled one-step raw states really ranges
over all of \texttt{K\_C}. This row-space test certifies a full-carrier \emph{homogeneous
implication}. It is not a certificate that an unrelated integer equation
\texttt{Mx=b} has a solution; that question needs an integer solution or Smith
divisibility evidence.

The affine expressions compose exactly:

\[ (A_2,b_2)\circ(A_1,b_1)
   =(A_2A_1,\ A_2b_1+b_2). \]

Identity has \texttt{(I,0)}. If the two steps have positive certificates
\texttt{D\_1A\_1=Λ\_1D\_0} and \texttt{D\_2A\_2=Λ\_2D\_1}, then the composite has
\texttt{D\_2A\_2A\_1=(Λ\_2Λ\_1)D\_0} and
\texttt{D\_2(A\_2b\_1+b\_2)=0}. Thus finite contexts and their balance-preserving affine
maps form a category at this selected excess level; certificates admit maps
but are not part of arrow equality. A composite can pass the endpoint
criterion when an intermediate step fails it: on a context
with marked face \texttt{012}, AddAssertion in R at \texttt{(0,1)} followed by Remove of
that new assertion returns the complete raw tables to their starting state.
The first step can make the face sum one, but the composite is identity with
\texttt{b=0} and row certificate \texttt{Λ=I}.
Consequently checking each step is sufficient, not necessary, for endpoint
balance. The converse failure is concrete rather than a claim about policy.

There is also an exact \emph{path comparison} test. For two affine paths with the
same source and target coordinate lists, \texttt{(A,b)} and \texttt{(B,c)} agree on every
balanced integer excess table if and only if

\[b=c\quad\text{and}\quad A-B=H D_C
  \quad\text{for some rational matrix }H.\]

Apply equality at zero, then use the same integer-kernel/row-space argument
for every row of \texttt{A-B}; the displayed equations also prove sufficiency.
Unequal constants or an integer kernel vector on which \texttt{A-B} is nonzero are
exact negative witnesses. In particular, a proposed restriction/edit square
commutes on balanced excess precisely when its two route expressions pass
this comparison test. Equalities of these extensional maps paste under
composition when the adjacent maps preserve their stated kernels, giving
coherent finite diamonds \textbf{for this declared layer}; they do not provide
fillers between arbitrary source histories or a Kan
structure.

\textbf{Path-saturation lemma for the five permitted edits.} Fix an initial context and a finite word
of the five actual edit descriptors. Assume the word has at least one legal
raw execution. Its intermediate contexts are determined by that context and
the word; only Remove's positive-count test depends on the tables. Every raw
prefix has a nonnegative integer linear part: AddTerm pads zeros, Identify
aggregates fibers, Fill is the identity, and assertion edits add constants.
Adding any common nonnegative table to both initial relations therefore
cannot disable a previously legal execution. For an arbitrary integer excess
\texttt{e}, start with \texttt{(e\_+,e\_−)} and add a common scalar \texttt{t} to \textbf{every} initial
pair coordinate. At each Remove the tested count is \texttt{αt+β(e)} with integer
\texttt{α≥0}. If \texttt{α=0}, the count has no initial-table dependence and equals its
positive value in the known legal execution. If \texttt{α\textgreater{}0}, a sufficiently large
\texttt{t} enables that Remove. One finite \texttt{t} handles the whole word. Consequently
the excess projection of its legal initial states is the entire \texttt{E\_C}, and
the projection of legal \textbf{balanced} states is exactly \texttt{K\_C}. If two words
have one common legal initial state, the same argument handles both words.
Thus the all-kernel endpoint criterion above is necessary as well as
sufficient for an actual nonempty five-edit model path, not just for its affine
extension. This lemma assumes unbounded nonnegative integer multiplicities;
a native allocation or bit-budget refusal is a separate execution outcome.

For a different guarded move system, let \texttt{S⊆K\_C} be the excess projection of
its legal whole-path inputs. When nonempty, choose \texttt{e\_0∈S} and put
\texttt{L\_S=span\_Q\{e-e\_0\ :\ e∈S\}}. The exact guarded criterion is
\texttt{D\_C\textquotesingle{}F(e\_0)=0} and \texttt{D\_C\textquotesingle{}A(L\_S)=0}; for two paths replace \texttt{D\_C\textquotesingle{}F} by their
difference on their \textbf{common} legal domain. A supplied finite rational basis
of \texttt{L\_S} makes this a finite check. Without that basis, a tested sample cannot
certify the full-domain guarded statement. An empty common domain is reported
as no enabled comparison, rather than as successful agreement.

Endpoint \textbf{raw-count} equality has a stronger exact test. Fix the same
ordered target carrier and face set for two raw affine path expressions.
Write the difference between their complete output count pairs as
\texttt{U\ r+V\ s+c}, with \texttt{r,s∈N\^{}(X×X)} and source balance \texttt{D\_C(r−s)=0}.
It vanishes on every balanced raw input if and only if

\[c=0,\qquad U+V=0,\qquad U=H D_C
  \quad\text{for some rational matrix }H.\]

Zero raw counts force \texttt{c=0}. Taking \texttt{r=s=t} for each nonnegative unit table
forces \texttt{U+V=0}. The difference is then \texttt{U(r−s)}. Every integer
\texttt{e∈K\_C} occurs as \texttt{r−s} by its positive/negative parts, so the earlier
row-space argument forces \texttt{U=H\ D\_C}; these equations conversely make the
difference zero. A failed equation has, respectively, the zero state, a
diagonal unit-count state, or an integer-kernel counterstate as its witness.
For two five-edit model words with a nonempty common legal domain, the displayed
test is also necessary and sufficient \textbf{on that domain}. The saturation
lemma supplies a legal representative of every balanced \texttt{e}. Moreover,
adding a diagonal unit count to both tables at one initial coordinate cannot
disable either word, so equality on the legal domain forces \texttt{U+V=0} one
coordinate at a time. A legal representative with \texttt{e=0} then forces \texttt{c=0},
and the representatives of all \texttt{e∈K\_C} force \texttt{U=H\ D\_C}. Failed conditions
have legal witnesses after common-offset padding; no arbitrary invalid raw
candidate is counted as a counterexample. For other guarded move systems,
their actual common domain must be checked separately. If the five-edit model raw
test passes and both paths end in the same declared context, all
deterministic future \emph{state} questions agree, including source-enabled moves
and plain-state price. It does not equate the paths' occurrences or accumulated
costs. Identity and the path that adds one R and one S assertion at the same
arc have identical endpoint excess on every input but different raw counts;
at the zero state, only the latter enables their removal. P.3's AddTerm
example likewise shows that balance preservation alone need not descend to
the gradient-result quotient. These meanings require the matching raw,
quotient, history and permission interfaces rather than an excess-only label.

\subsubsection{P.9 Whole-word descent to fixed-writing result classes}\label{p.9-whole-word-descent-to-fixed-writing-result-classes}

P.8 concerns balance and equality of raw or excess endpoints. The actual
five-edit model result mark asks a different question. At one fixed ordered writing
and fixed oriented face set \texttt{F}, the source result relation is

\[e\sim_X e'\quad\Longleftrightarrow\quad e-e'\in G_X,
  \qquad G_X=\operatorname{im}_{\mathbb Z}d_X,\quad
  (d_Xg)_{xy}=g_x-g_y.\]

The source-specific completeness proof in the \href{docs/marked_result_refinement.md}{marked-result
supplement} identifies this relation with
equality of the normalized table component when the same \texttt{F} is retained.
It does not identify raw counts or ordered histories. Write
\texttt{Q\_(X,F)=(E\_X/G\_X,F)} for this \textbf{fixed-writing} class carrier; the fully
permutation-minimized \texttt{result\_mark} has an additional orbit relation.

Fix an initial context \texttt{(X,F)} and a finite word \texttt{w} of the five actual
five-edit model edit descriptors with at least one legal execution. Follow each
original vertex through AddTerm inclusions and Identify mergers to obtain a
total map \texttt{f\_w:X→Y} to the final ordered carrier. The final face set \texttt{F\_w}
depends on \texttt{(X,F)} and \texttt{w}, but not on the initial count tables. Pair-table
aggregation satisfies \texttt{T\_gT\_f=T\_(g∘f)} by regrouping finite sums. AddTerm
is aggregation along an inclusion with an empty new fiber; Identify is
aggregation along its merger; Fill and assertion edits have the identity
linear part. Induction over the word therefore gives the \textbf{exact} expression

\[e_w=T_{f_w}e+b_w,\qquad
  b_w=\sum_{j\text{ assertion edit}}T_{j\rightsquigarrow w}\,\delta_j,\]

where \texttt{δ\_j=±E\_xy} is the excess change at an AddAssertion or Remove edit
and \texttt{T\_(j↝w)} is aggregation through the remaining carrier changes. Remove
has the corresponding opposite sign from AddAssertion in the same relation.
The formula does not reconstruct raw counts, costs or occurrences; those
fields must remain available through their separate interfaces.

\textbf{Fixed-writing whole-word descent theorem.} The word defines a total,
representative-independent map \texttt{{[}e{]}↦{[}T\_(f\_w)e+b\_w{]}} from \texttt{Q\_(X,F)} to
\texttt{Q\_(Y,F\_w)} by choosing any \emph{legal} raw representative of each source class
and applying the word if and only if

\[T_{f_w}(G_X)\subseteq G_Y.\]

Sufficiency is immediate: two equivalent source excesses differ by \texttt{d\_Xg},
and their output difference is \texttt{T\_(f\_w)d\_Xg}; the translation \texttt{b\_w} cancels.
For necessity, if one integer gradient fails the inclusion, P.8's path
saturation supplies legal raw representatives with excesses \texttt{0} and \texttt{d\_Xg}
and the same initial \texttt{F}. They have one source class and two target classes.
The legal domain therefore cannot hide a failed gradient. The source and
target face sets in this theorem are fixed; a word whose descriptors are
invalid in the chosen context has no induced map by vacuous truth.

P.3 supplies the complete carrier test for this actual \texttt{T\_f}. The inclusion
holds \textbf{exactly} when \texttt{\textbar{}Y\textbar{}≤2}, at most one fiber of \texttt{f\_w} is occupied, or all
target fiber sizes are equal. Empty target fibers, including new surviving
terms, count. A positive integral certificate is a matrix \texttt{H} of integer
potentials satisfying \texttt{T\_(f\_w)d\_X=d\_YH}, one column per source unit
potential. For \texttt{Y≠∅}, each proposed target column can be checked directly:
all loops vanish, reverse directed entries are negatives, and every
distinct-vertex triangle sum vanishes. Set the root potential to zero and
\texttt{h\_a=t\_(a,0)}; the triangle through the root then establishes
\texttt{t\_(a,b)=h\_a-h\_b} with \textbf{integer} \texttt{h}. For \texttt{Y=∅}, the pair table is zero.
If the fiber-size test fails, P.3 constructs a source unit potential whose
transported gradient has a nonzero target triangle; this is an exact negative
witness using two legal raw representatives after common-offset padding.
Rational span membership alone would not certify an integer lift for a
general lattice; this particular full directed-pair gradient has the explicit
integer root construction.

The result-class and balance requirements are independent. A word induces a
map on \textbf{balanced} fixed-writing classes
\texttt{(K\_(X,F)/G\_X)→(K\_(Y,F\_w)/G\_Y)} precisely when its fiber-size test passes
and the P.8 endpoint equations also hold:

\[D_{F_w}b_w=0,\qquad D_{F_w}T_{f_w}=\Lambda D_F
  \quad\text{for some rational }\Lambda.\]

Every gradient lies in the appropriate balance kernel, so both quotients are
well-defined. AddTerm from two old terms to three with no faces preserves
balance vacuously but has fiber sizes \texttt{(1,1,0)} and fails result descent.
Conversely, adding one R assertion on a directed arc of a marked triangle
has \texttt{f\_w=id} and passes result descent, but \texttt{D\_(F\_w)b\_w=1} and fails balance.
These are countermodels to replacing the two tests with one status flag.

Two words with the same initial context and the \textbf{same} final ordered
carrier and face set, each passing descent and having a nonempty common
legal domain, induce the same map on every fixed-writing result class iff

\[b_w-b_v\in G_Y\quad\text{and}\quad
 (T_{f_w}-T_{f_v})(E_X)\subseteq G_Y.\]

Evaluate at zero and at every source unit table for necessity; those excesses
have common legal raw representatives by P.8. The same equations prove
sufficiency by additivity. Exact integer potential witnesses for the constant
and each matrix column certify the condition; a failed column plus a unit
source table is a counterwitness. If both words also pass the P.8 balance
test, comparison restricted to balanced classes replaces the unit-column
test by an \textbf{integer} generating set for \texttt{K\_(X,F)}, with the constant still
tested separately. A rational kernel basis alone does not certify integral
membership in a target gradient lattice. Different final face sets fail
fixed-writing equality before any table comparison.

For example, on \texttt{X=\{0,1\}}, the potential \texttt{g=(1,0)} gives \texttt{e\_01=1} and
\texttt{e\_10=-1}. AddTerm produces \texttt{Y=\{0,1,2\}} with zero new incident entries,
so the target triangle \texttt{(0,1,2)} has sum one: descent fails even for this
legal representative. Follow AddTerm instead by Identify of the new term
with old term \texttt{0}. The complete word has final fiber sizes \texttt{(1,1)} and
acts as the identity on every initial pair coordinate, including loops.
It descends although its first prefix does not. It still incurs two
occurrences and a term-addition price. The theorem classifies fixed
canonical writings and complete word endpoints; it does not make a word
equivariant under arbitrary source renamings, equate raw histories, recover
permission, or construct the missing global site/higher fillers in O02/O03.

\subsubsection{P.10 Full result-mark descent on one source face orbit}\label{p.10-full-result-mark-descent-on-one-source-face-orbit}

P.9 fixes a writing. The five-edit model's complete \texttt{result\_mark} also identifies
the simultaneous renamings of the excess table and oriented fillings.
Here is an exact test for \textbf{one fixed literal edit word} on the orbit of
one initial face context. It does not silently rename the edit descriptors.

Write \texttt{E\_X=Z\^{}(X×X)}, \texttt{G\_X=im\ d\_X}, \texttt{Q\_X=E\_X/G\_X}, and let \texttt{P\_p} transport
pair coordinates under a permutation \texttt{p∈S\_X}. The same \texttt{p} carries each
oriented filling through cyclic canonicalization; write its result as
\texttt{pF}. Let \texttt{w} have one legal execution on \texttt{(X,F)}, with the P.9 excess
expression \texttt{Ae+b}, where \texttt{A=T\_(f\_w)}, and let \texttt{F\_1=w(F)} be its final
face set. For every \texttt{p∈S\_X}, run only the context skeleton of the same
word on \texttt{(X,pF)}. Let \texttt{Adm\_w(F)} contain the \texttt{p} for which all descriptors
are structurally valid, and write \texttt{F\_p=w(pF)} for their final face sets.
The identity belongs to \texttt{Adm\_w(F)}. The carrier maps and excess
translations of the five edits do not depend on fillings, so every such
writing has the \textbf{same} \texttt{A,b}, applied to its own input table. The known
legal base count run is also a legal count run in each structurally valid
face context; P.8 saturation then supplies a legal raw representative of
every integer excess there. Invalid contexts are excluded, not counted as
successful executions.

\textbf{Full-mark word theorem.} Applying \texttt{w} to any legal raw representative
defines a total map from the full five-edit model result classes with face set
in the orbit of \texttt{F} to full target result classes if and only if:

\begin{enumerate}
\def\labelenumi{\arabic{enumi}.}
\item
  \texttt{A(G\_X)⊆G\_Y};
\item
  for every \texttt{p∈Adm\_w(F)}, some \emph{single} \texttt{q∈S\_Y} satisfies
  \texttt{qF\_1=F\_p} and both integer-gradient conditions

  \[b-P_qb\in G_Y,\qquad
    (AP_p-P_qA)(E_X)\subseteq G_Y.\]
\end{enumerate}

The word \emph{single} means that \texttt{q} works for every source excess table for
that \texttt{p}; the witness \texttt{q} may differ between source permutations. These
conditions give a finite, exact certificate: check the final oriented
face lists, then supply integer target potentials for the constant and
each source-unit column. Condition 1 is the P.3 fiber-size criterion.
The certificate asserts a map on result classes, not equality of raw
histories or a group homomorphism \texttt{p↦q}. Enumerating source and target
permutations and the negative grid can be factorial or exponential; no
polynomial-time or implemented decision procedure is asserted.

For sufficiency, condition 1 makes the output independent of a source
gradient within any admitted writing. Given another source
representative with table \texttt{P\_p\ e+g} and faces \texttt{pF}, where \texttt{g∈G\_X},
condition 2 gives
\texttt{{[}A(P\_p\ e+g)+b{]}={[}P\_q(Ae+b){]}} in \texttt{Q\_Y} and carries \texttt{F\_1} to \texttt{F\_p}.
Thus the two outputs have one full joint mark. Every source class has a
legal representative in the base writing by saturation, so the induced
map is total on the stated face orbit.

For necessity of condition 1, take \texttt{g∈G\_X}. All tables \texttt{kg}, for integer
\texttt{k≥0}, have the same source class as zero and have legal representatives.
Their target classes are \texttt{{[}b{]}+k{[}Ag{]}}. The target permutation orbit of
\texttt{({[}b{]},F\_1)} is finite. The lattice \texttt{Q\_Y} is torsion-free: when \texttt{Y≠∅},
subtract the integer gradient with potential \texttt{h\_a=e\_(a,0)} for \texttt{a≠0}
and \texttt{h\_0=0}; this splits \texttt{E\_Y} into gradients and tables whose nonroot
entries in the root column vanish. Hence \texttt{{[}Ag{]}≠0} would produce
infinitely many distinct target classes in a finite orbit. In fact, one
of \texttt{k=1,…,\textbar{}S\_Y\textbar{}} is already a legal counterwitness. Therefore every
source gradient must map into \texttt{G\_Y}.

For necessity of condition 2, fix an admitted \texttt{p}. Full-mark equality
for every \texttt{e∈E\_X} says that for each \texttt{e}, at least one target
permutation \texttt{q} carrying \texttt{F\_1} to \texttt{F\_p} makes the affine difference

\[[AP_p e+b-P_q(Ae+b)]\in Q_Y\]

zero. There are finitely many candidate \texttt{q}. If none made this difference
identically zero, choose a nonzero integer affine coordinate of the
difference for each candidate, using the preceding lattice splitting.
Their product is a nonzero polynomial over \texttt{Q}, which cannot vanish on
all of \texttt{Z\^{}(X×X)}: induction on the number of variables establishes that the
integer lattice is Zariski dense. Thus one \texttt{e} would defeat every \texttt{q},
contradicting full-mark equality. An identically zero candidate gives
exactly the two displayed gradient conditions, by evaluating at zero
and each source-unit table. If there are \texttt{m} face-compatible candidate
permutations but none passes, the same polynomial argument finds a
countertable in the explicit grid \texttt{\{0,…,m\}\^{}(X×X)}; its degree in each
variable is at most \texttt{m}. If \texttt{m=0}, zero already witnesses the face
mismatch. P.8 provides legal raw representatives for the base and
renamed contexts. This finite-union step uses the torsion-free target
gradient quotient; it would not hold for an arbitrary finite torsion
report, where proper congruence classes can cover the input lattice.

The distinction is concrete. On three unmarked terms, \texttt{Identify(0,1)}
passes P.9 because the target has two terms. Take only diagonal excesses
\texttt{(e\_00,e\_11,e\_22)=(0,1,2)}, with all other entries zero. The identified
output has diagonal \texttt{(1,2)}. Rename the source by swapping terms \texttt{0}
and \texttt{2}, then apply the \textbf{same} \texttt{Identify(0,1)}: its output diagonal is
\texttt{(3,0)}. Gradients preserve loops, and no target permutation changes
the unordered loop multiset \texttt{\{1,2\}} into \texttt{\{0,3\}}. Both raw inputs are
legal and have the same source result mark, so the fixed literal word
does \textbf{not} descend to that full mark. Conversely, AddTerm followed by
Identify of the new last term with old term \texttt{0} returns every original
raw pair coordinate and filling on any nonempty carrier. That word acts
as identity on full result marks despite its non-descending AddTerm
prefix when at least two old terms are present. The complete criterion
remains finite and source-specific. P.11 constructs the separate raw-layer
transport of \textbf{renamed descriptors}; neither result supplies cross-system
refinements, a global site or higher fillers.

\subsubsection{P.11 Transported five-edit model words commute with renaming}\label{p.11-transported-five-edit-model-words-commute-with-renaming}

P.10 holds the edit descriptors literally fixed while renaming the input.
For a different, useful square, rename the descriptors \emph{with} the input.
Let \texttt{C=(X,r,s,F)} be an actual five-edit model count/filling state on the
ordered carrier \texttt{X=\{0,…,n−1\}}. A permutation \texttt{p∈S\_X} acts by simultaneous
transport of both directed count tables and every oriented filling;
it regenerates the source's canonical labels \texttt{t0,…,t(n−1)} rather than
permuting those literal strings. Write the result \texttt{pC}. The following
table defines the transported edit
\texttt{a\^{}p} and a target-carrier permutation \texttt{p\_a} for a well-formed edit
descriptor \texttt{a}, including one whose count or face guard later refuses.

\begin{longtable}[]{@{}
  >{\raggedright\arraybackslash}p{(\linewidth - 4\tabcolsep) * \real{0.3333}}
  >{\raggedright\arraybackslash}p{(\linewidth - 4\tabcolsep) * \real{0.3333}}
  >{\raggedright\arraybackslash}p{(\linewidth - 4\tabcolsep) * \real{0.3333}}@{}}
\toprule\noalign{}
\begin{minipage}[b]{\linewidth}\raggedright
Edit \texttt{a}
\end{minipage} & \begin{minipage}[b]{\linewidth}\raggedright
Transported descriptor \texttt{a\^{}p}
\end{minipage} & \begin{minipage}[b]{\linewidth}\raggedright
Target permutation \texttt{p\_a}
\end{minipage} \\
\midrule\noalign{}
\endhead
\bottomrule\noalign{}
\endlastfoot
\texttt{AddTerm} & \texttt{AddTerm} with the same new last term & Extend \texttt{p} by fixing that new term. \\
\texttt{AddAssertion(R/S,x,y)} & Same relation at \texttt{(p\ x,p\ y)} & \texttt{p} \\
\texttt{Remove(R/S,x,y)} & Same relation at \texttt{(p\ x,p\ y)} & \texttt{p} \\
\texttt{Fill(t)} & \texttt{Fill(cyc(p\ t))}, rotating the least vertex first without reversing orientation & \texttt{p} \\
\texttt{Identify(x,y)}, \texttt{x\textless{}y} & \texttt{Identify(u,v)}, where \texttt{u=min(p\ x,p\ y)} and \texttt{v=max(p\ x,p\ y)} & The unique \texttt{p\_a} satisfying \texttt{p\_a\ f\_(x,y)=f\_(u,v)p}. \\
\end{longtable}

Here \texttt{f\_(x,y)} is the actual canonical quotient map that identifies
\texttt{y} with \texttt{x} and shifts labels above \texttt{y} down by one. For Identify, the
displayed commuting equation defines a \textbf{bijection} on target carriers:
\texttt{p} carries exactly the merged fiber \texttt{\{x,y\}} to \texttt{\{u,v\}} and singleton
fibers to singleton fibers, so \texttt{f\_(u,v)p} is constant on each fiber of
the surjective \texttt{f\_(x,y)}. The induced map is well-defined and bijective,
and surjectivity makes it unique. The order of the two renamed endpoints
only chooses the source's canonical surviving label; it does not erase
the target permutation.

\textbf{One-step naturality.} On the raw pair/filling layer, admission and
the output commute exactly:

\[a\text{ enabled on }C\ \Longleftrightarrow\
  a^p\text{ enabled on }pC,\qquad
  a^p(pC)=p_a(aC).\]

AddTerm zero-pads both tables and leaves old faces, and its new vertex
is fixed by the extended permutation. Assertion edits transport their
unit coordinate because \texttt{P\_p\ E\_(x,y)=E\_(p\ x,p\ y)}; a Remove's positive
count guard is the same integer at the transported coordinate. Fill's
absent-face guard and oriented cyclic canonicalization commute with
renaming. For Identify, pair aggregation obeys
\texttt{T\_(f\_(u,v))\ P\_p=P\_(p\_a)\ T\_(f\_(x,y))} by reindexing each fiber sum.
The same commuting square carries filled triples: distinctness,
degeneracy after merging, cyclic orientation and duplicate removal are
preserved by the bijective \texttt{p\_a}. These checks use both count tables;
an excess-only argument would not prove Remove admission.

For a legal word \texttt{w=a\_k⋯a\_1}, transport each descriptor using the
permutation at its \textbf{current} carrier, starting with \texttt{p\_0=p}. The table
produces \texttt{a\_i\^{}(p\_(i−1))} and \texttt{p\_i}; call the resulting word \texttt{w\^{}p}.
Induction over the one-step square establishes

\[w^p(pC)=p_k(wC).\]

The transported word is legal on \texttt{pC}; transporting it back by \texttt{p⁻¹}
recovers \texttt{w}, so the same statement holds in both directions. A refused
one-step guard has the corresponding refusal under renaming, but an
invalid descriptor does not acquire an output by this theorem.

The rule is coherent under both compositions. For consecutive words
\texttt{w} then \texttt{v}, transport of the concatenation is \texttt{w\^{}p} followed by
\texttt{v\^{}(p\_k)}. For initial renamings \texttt{p} then \texttt{q}, transporting first by
\texttt{p} and then by \texttt{q} gives the same literal transported descriptors as
transporting once by \texttt{qp}, and the terminal permutations compose.
At Identify this follows from the uniqueness of the induced target
permutation; at the other edits it is coordinate transport or the
extension that fixes the new vertex. Thus identities and finite word
composition give strict, replayable renaming squares on this \textbf{selected
raw source layer}.

The raw equality implies corresponding affine excess conjugacy,
\texttt{A\_(w\^{}p)P\_p=P\_(p\_k)A\_w} and \texttt{b\_(w\^{}p)=P\_(p\_k)b\_w}, and transports marked
cycle values with their oriented faces. It therefore carries a passing
P.8 balance test or P.9 gradient-descent test to the transported word
in the renamed context. It does not make a \textbf{fixed literal} word act
on every full result mark: the three-term Identify failure in P.10
keeps \texttt{Identify(0,1)} fixed, whereas its transported descriptor after
swapping source terms \texttt{0} and \texttt{2} is \texttt{Identify(1,2)}. Neither exact
renaming square supplies a gradient quotient map when P.3 fails,
nor does it supply institutional permission or standing. P.12 checks
the selected source prices and complete ordered runs separately.
Cross-system source refinements, covering restrictions and higher
fillers remain O02/O03 obligations.

\subsubsection{P.12 Five-edit model prices and complete runs under descriptor transport}\label{p.12-five-edit-model-prices-and-complete-runs-under-descriptor-transport}

The declared five-edit model native execution
adapter
retains two count tables, canonical term labels, all intermediate raw
states and the ordered edit occurrences. Its two state quantities and
its run price have the exact definitions

\[S(C)=n+\sum_{x,y}(r_{xy}+s_{xy}),\qquad
  P(C)=n+\sum_{x,y}|r_{xy}-s_{xy}|,\]

\[A(R)=(\#\operatorname{AddTerm},\#\operatorname{AddAssertion}).\]

These are respectively \texttt{source\_size}, \texttt{plain\_state\_cost} and the
two-component \texttt{addition\_price}. \texttt{Fill}, \texttt{Identify} and \texttt{Remove} each
contribute \texttt{(0,0)} to the latter, although each remains an occurrence.
For a legal run
\texttt{R=(C\_0,a\_1,C\_1,…,a\_k,C\_k)}, use P.11 to transport every state and
its selected descriptor with successive permutations \texttt{p\_i}. The
result \texttt{R\^{}p} is a legal run on \texttt{pC\_0}, and

\[|R^p|=|R|=k,\qquad A(R^p)=A(R),\qquad
  S(p_iC_i)=S(C_i),\quad P(p_iC_i)=P(C_i)\ (0\le i\le k).\]

The proof is coordinatewise: each \texttt{p\_i} bijects the directed pair
coordinates and preserves the carrier size; hence both sums are
unchanged. The transported descriptor has the same event kind and
relation, so its addition-price increment is unchanged. P.11 gives
the exact equality of every transported intermediate raw state.
Transport by the inverse initial permutation recovers the original
run, with all occurrences in the original order. Thus transport is a
bijection between \textbf{complete legal runs} starting at \texttt{C\_0} and those
starting at \texttt{pC\_0}; it does not collapse different runs that happen
to have equal endpoints or result marks.

The native run splice \texttt{R.then(T)} requires complete equality of the
joining raw state and adds the two price pairs. Since renaming is
bijective on complete states, a defined splice transports exactly as
\texttt{(R.then(T))\^{}p=R\^{}p.then(T\^{}(p\_k))}; a mismatched join remains a
mismatch. At one state, the enabled \textbf{set} also transports exactly:

\[\operatorname{Enabled}(pC)=
  \{a^p:a\in\operatorname{Enabled}(C)\}.\]

This is P.11's admission equivalence in both directions. It does not
preserve the source's lexicographic \textbf{list order} of enabled events:
swapping terms \texttt{1} and \texttt{2} on a three-term state exchanges the order
of the distinct assertion descriptors \texttt{(0,1)} and \texttt{(0,2)}. Selected
run order is preserved because its occurrences are transported
positionwise, not regenerated from that enumeration list.

The limits are witnessed on the same carrier. On two terms, zero
excess and the gradient with \texttt{e\_01=1}, \texttt{e\_10=-1} have the same
fixed-writing result class and no fillings, but their plain-state
costs are respectively \texttt{2} and \texttt{4}; positive/negative parts give
legal raw count tables. Thus \texttt{P} does not descend to that result
class. On any nonempty carrier, identity replay and the legal word
AddTerm followed by Identify of the new term with old term \texttt{0}
have the same complete raw endpoint, but addition prices \texttt{(0,0)}
and \texttt{(1,0)} and different occurrence counts. Price and history are
not functions of the endpoint even when its raw tables agree.

These equalities close price and history transport for this declared
renaming profile. They do not make \texttt{P(C)} constant on gradient-result
classes, make a fixed literal word equivariant, prove equivalence to
every original executable source run, or transport institutional
authorization. Those are separate O02/O03/O05 obligations.

\subsubsection{P.13 Permutation cubes for initially addressed permitted edits}\label{p.13-permutation-cubes-for-initially-addressed-permitted-edits}

P.11 transports one word under a renaming; P.12 retains that word's ordered
history. Here the question is different: when may a fixed collection of
source edits be scheduled in different orders? The result concerns the
declared native five-edit model transition
semantics
on constructible finite carriers. It does not treat an edit addressed to a
term created \emph{during} the schedule.

Let \texttt{C=(X,r,s,F)} be one valid raw state, with \texttt{X=\{0,...,n-1\}}, two
nonnegative directed count tables, and a set of oriented cyclic faces. Let
\texttt{M} be a finite \textbf{tagged multiset} of well-formed descriptors, each addressed
to this initial \texttt{X}: \texttt{AddTerm}, \texttt{AddAssertion}, \texttt{Remove}, \texttt{Fill}, or \texttt{Identify}.
Tags distinguish planned occurrences of an identical descriptor; they do not
change the native event type. On any prefix, residualize each still-pending
descriptor by sending its original vertices through the prefix's current
quotient map. An Identify uses the smaller of the two current labels as its
survivor; a nondegenerate Fill is cyclically rotated to put its least vertex
first, without reversing orientation. \texttt{AddTerm} appends a fresh zero row and
column in each table. No residual may silently become a no-op: a collapsed
Identify, degenerate Fill, already-filled face, or empty Remove is refused.

Write \texttt{I} for the multigraph on \texttt{X} whose tagged edges are the Identify
occurrences. Let \texttt{\textasciitilde{}} be the equivalence relation generated by those edges,
and \texttt{q:X\ -\textgreater{}\ X/\textasciitilde{}} the quotient with classes labelled in increasing order of
their least original vertex. For a face \texttt{t}, write \texttt{q\#t} for its cyclically
canonical image if its three images are distinct, and for \texttt{none} otherwise;
\texttt{q\#F} discards \texttt{none} and deduplicates. If \texttt{a\_R(x,y)} and \texttt{m\_R(x,y)} count
AddAssertion and Remove occurrences at an original coordinate in relation
\texttt{R}, define analogous counts for \texttt{S}.

\textbf{Raw endpoint theorem.} Any two \emph{legal complete} orderings of \texttt{M} have the
same final five-edit model state and the same addition-price pair, although their
ordered runs remain distinct. The final carrier has \texttt{\textbar{}X/\textasciitilde{}\textbar{}\ +\ \#AddTerm}
canonical labels. On its original-class block, its tables are

\[
 \begin{aligned}
 r_{\mathrm{fin}}&=q_*r+\sum_{x,y}(a_R(x,y)-m_R(x,y))E_{q(x),q(y)},\\
 s_{\mathrm{fin}}&=q_*s+\sum_{x,y}(a_S(x,y)-m_S(x,y))E_{q(x),q(y)}.
 \end{aligned}
\]

All coordinates incident to newly added terms are zero. Its face set is
\texttt{q\#(F\ union\ \{t\ :\ Fill(t)\ occurs\ in\ M\})}, and its addition price is
\texttt{(\#AddTerm,\ \#AddAssertion)}. This statement does not assert that any complete
ordering exists. It also does not identify the two run objects: their edit
occurrences and intermediate raw states retain their actual orders.

\textbf{All-order legality and cube criterion.} Every permutation of the tagged
occurrences is a legal complete run if and only if all three conditions hold:

\begin{enumerate}
\def\labelenumi{\arabic{enumi}.}
\tightlist
\item
  \texttt{I} is a forest as a multigraph; a repeated Identify edge is a two-edge
  cycle.
\item
  At every original directed coordinate and separately in each relation,
  the initial count is at least its total number of planned Remove
  occurrences: \texttt{r\_xy\ \textgreater{}=\ m\_R(x,y)} and \texttt{s\_xy\ \textgreater{}=\ m\_S(x,y)}.
\item
  Every planned Fill has a nondegenerate \texttt{q\#t}; these final oriented faces
  are pairwise distinct and none belongs to \texttt{q\#F}.
\end{enumerate}

When these conditions hold, every subset \texttt{S} of the tagged occurrences has a
well-defined raw endpoint \texttt{C\_S} independent of its order. For each pending
tag \texttt{a}, its residual gives an enabled edge \texttt{C\_S\ -\textgreater{}\ C\_(S\ union\ \{a\})}.
Consequently the subset lattice carries an \textbf{endpoint-coherent Boolean-cube
skeleton} of actual raw states and edits: each square and every
higher-dimensional boundary reaches the same terminal raw state. The diagram
records the boundary routes as distinct runs. Endpoint equality supplies
replayable boundary data; it does not by itself define a 2-cell between runs,
a 3-cell between squares, or a higher-category filler. Those structures
still require a declared comparison relation and composition laws.

\textbf{Proof.} For any finite vertex map \texttt{f}, pair-table pushforward is linear
and \texttt{(g\ f)\_*=g\_*\ f\_*}. Each Identify step aggregates both tables by this
pushforward; each AddTerm uses the inclusion pushforward; an assertion edit
adds or subtracts one basis table at its residual coordinate. Sending that
unit through every later quotient yields \texttt{E\_(q(x),q(y))}. Thus the displayed
table formulas depend only on \texttt{M}, not its order. The source's Identify
operation removes the higher current label, so induction shows that current
classes are labelled by increasing minimum original vertex. Fresh AddTerm
vertices stay after the original classes because no descriptor in \texttt{M} names
one. Oriented-face transport satisfies \texttt{(g\ f)\#=g\#\ f\#}: once a face becomes
degenerate it cannot become distinct again, and cyclic rotation commutes
with taking the nondegenerate image. Fill adds a singleton face, while
Identify transports and deduplicates the set. This establishes the face formula
even when a Fill later collapses or merges with another face. The native
addition price counts kinds, so it is also order-independent.

For necessity of the three conditions, schedule all Identifies first. A
cycle's last edge has endpoints already joined and is refused. Schedule all
Removes before any Identify or AddAssertion: repeated subtraction at each
original coordinate succeeds only with the stated initial supply. Schedule
all Identifies before all Fills: each Fill must survive \texttt{q}, avoid \texttt{q\#F},
and avoid every earlier Fill image, in any ordering. For sufficiency, a
forest edge always joins distinct components of every preceding edge
subset. At any intermediate quotient coordinate, the remaining count is
the sum of original supplies in that fiber, plus preceding additions, minus
preceding removals; the per-coordinate supply inequalities make the next
Remove positive. A face that is nondegenerate and distinct after the \emph{full}
quotient was nondegenerate and distinct after every partial quotient;
otherwise later identifications could not separate it. Hence every Fill
guard succeeds. AddTerm and AddAssertion are unguarded on their well-formed
residual addresses. The same argument applies to each subset and establishes
the cubical assertion. This is an arbitrary-finite proof; bounded replay is
only a check of its implementation-facing cases.

The obstructions are actual, not inferred from an endpoint quotient. Three
Identifies on the edges of a triangle are individually enabled but no
complete ordering can execute the last edge. On one term with \texttt{r\_00=1},
two tagged \texttt{Remove(R,0,0)} occurrences exhaust their supply after the first.
On three terms, \texttt{Fill(0,1,2)} followed by \texttt{Identify(0,1)} is legal, but
reversing them makes the Fill degenerate. On four terms with
\texttt{F=\{cyc(0,2,3)\}}, \texttt{Fill(1,2,3)} is initially enabled; after
\texttt{Identify(0,1)}, its image collides with the existing face and is refused.
The two cyclic orientations remain distinct unless an actual quotient
identifies their oriented images.

For a two-occurrence fork, the criterion specializes to exact strict-swap
guards: an Identify/Fill pair swaps precisely when the filled face survives
the identification and avoids the transported old faces; two Removes at one
relation/coordinate require count at least two; two identical Fills or
Identifies do not swap. Other pairs of individually enabled initial-carrier
descriptors satisfy the criterion. This includes overlapping but distinct
Identify edges: their two-edge graph is still a forest. The theorem itself
does not cover a later edit that first names an AddTerm-created vertex;
P.14 supplies an additional named-token presentation for that case. Selected
gradient-result quotient descent, arbitrary restrictions/covers, permission
and institutional standing remain separate. Resource and allocation refusals
are outside these mathematical guards.

\subsubsection{P.14 Causal names and exact guards for dynamically created terms}\label{p.14-causal-names-and-exact-guards-for-dynamically-created-terms}

P.13 compares edits addressed only to the starting carrier. A literal
numeric address in this model is not a stable name for a newly created term when two
\texttt{AddTerm} occurrences change order. Supply that missing comparison datum
explicitly, without modifying the native five-edit
kernel.

Fix an initial valid raw state \texttt{C=(X,r,s,F)}. Let \texttt{E} be a finite set of
tagged edit occurrences with an acyclic precedence relation. Each \texttt{AddTerm}
occurrence births one unique token; put \texttt{U=X\ disjoint-union\ B} for the initial
and birth tokens. Every other descriptor names tokens in \texttt{U}, not current
numeric labels. \texttt{Identify} names an unordered token pair; \texttt{Fill} names an
oriented cyclic token triple. Require the birth of each named token to
precede its use.
Additional declared precedence edges are allowed. A linear extension is
executed natively by maintaining a map from every born token to its current
five-edit model label: \texttt{AddTerm} appends its token, \texttt{Identify} sends the lesser and
greater current labels as native \texttt{(x,y)} with \texttt{x\textless{}y} and canonically renumbers,
and \texttt{Fill} cyclically rotates its three current labels so the least is first.
\texttt{Remove} and \texttt{AddAssertion} use their named tokens' current labels in order.
A token remains a valid name after its class
is merged. A reference before birth is an invalid presentation, not a
negative mathematical instance.

For an order ideal \texttt{J} of the event poset, let \texttt{U\_J} be \texttt{X} plus the tokens
born in \texttt{J}. Let \texttt{I\_J} be its Identify-edge multigraph and \texttt{q\_J} the quotient
of \texttt{U\_J} by those edges. Use a fixed total order on token names to label its
classes by their least token; this is a \textbf{symbolic} writing, which may differ
from the native writing reached by a chosen linear extension. Define two
formal count tables and one oriented face set on these classes:

\[
\begin{aligned}
 R_J &= (q_J)_*r+
   \sum_{\substack{a\in J\\a=\mathrm{AddAssertion}(R,u,v)}}E_{q_Ju,q_Jv}
  -\sum_{\substack{a\in J\\a=\mathrm{Remove}(R,u,v)}}E_{q_Ju,q_Jv}, \\
 S_J &= (q_J)_*s+
   \sum_{\substack{a\in J\\a=\mathrm{AddAssertion}(S,u,v)}}E_{q_Ju,q_Jv}
  -\sum_{\substack{a\in J\\a=\mathrm{Remove}(S,u,v)}}E_{q_Ju,q_Jv},
\end{aligned}
\]

\[
 F_J=(q_J)_\#\left(F\cup
          \{t:\mathrm{Fill}(t)\in J\}\right).
\]

The initial tables are extended by zero on birth-token coordinates before
pushforward. \texttt{q\#} cyclically normalizes nondegenerate oriented triples and
discards degenerate ones. These formulas can have negative entries for an
unreachable ideal; calling them a formal profile does not admit an invalid
five-edit model state.

\textbf{Causal endpoint theorem.} Every \emph{legal} linear extension of \texttt{E} ends in
the symbolic profile \texttt{(R\_E,S\_E,F\_E)} up to the unique permutation that maps
each final native label to its token-equivalence class. Thus any two legal
linear extensions have isomorphic complete raw endpoints. Their occurrence
counts and AddTerm/AddAssertion addition-price pairs agree; P.12 then gives
equal endpoint source size and plain-state cost and transports their future
enabled \textbf{sets}. Their ordered histories, intermediate prices, and
lexicographic enabled-event lists need not agree. The theorem does not assert
that a legal linear extension exists.

\textbf{All-linear-extension criterion.} Assume the Identify-edge multigraph on
all of \texttt{U} is a forest. Then \emph{every} linear extension is legal if and only if,
for every order ideal \texttt{J} and every occurrence \texttt{a} whose predecessors lie in
\texttt{J} but which is not in \texttt{J}, the following applicable guard holds in the
formal profile of \texttt{J}:

\begin{longtable}[]{@{}
  >{\raggedright\arraybackslash}p{(\linewidth - 2\tabcolsep) * \real{0.5000}}
  >{\raggedright\arraybackslash}p{(\linewidth - 2\tabcolsep) * \real{0.5000}}@{}}
\toprule\noalign{}
\begin{minipage}[b]{\linewidth}\raggedright
Pending kind
\end{minipage} & \begin{minipage}[b]{\linewidth}\raggedright
Exact ideal-local guard
\end{minipage} \\
\midrule\noalign{}
\endhead
\bottomrule\noalign{}
\endlastfoot
\texttt{Remove(R/S,u,v)} & \texttt{R\_J(q\_Ju,q\_Jv)\textgreater{}0} or \texttt{S\_J(q\_Ju,q\_Jv)\textgreater{}0}, in the named relation. \\
\texttt{Fill(u,v,w)} & The three \texttt{q\_J} images are distinct and their oriented cyclic face is absent from \texttt{F\_J}. \\
\texttt{Identify(u,v)} & The two images differ; this follows from the forest assumption. \\
\texttt{AddTerm}, \texttt{AddAssertion} & Enabled once the declared birth precedences make their names available. \\
\end{longtable}

The test ranges over the finite lattice of order ideals, rather than a sample
of linear extensions. It is a finite exact decision procedure: enumerate
ideals, construct their quotient and the displayed integer tables/faces,
and replay the applicable guard. It may be exponential in \texttt{\textbar{}E\textbar{}}; no claim of
an efficient generic solver is made. If the forest condition fails, no
complete native execution can perform all its Identify occurrences. If a
guard fails, choose a smallest failing ideal. Its proper ideal prefixes
pass, so a linear extension reaching that ideal gives a concrete native
refusal at the named pending occurrence. Refusal is never counted as a
successful no-op.

\textbf{Proof.} In any legal prefix, the token-to-native-label map has fibers
exactly the connected components of the performed Identify edges. Indeed,
an \texttt{AddTerm} creates a singleton fiber, and an enabled Identify unions two
different fibers; no other edit changes them. It therefore induces a unique
bijection from the symbolic \texttt{q\_J} classes to current native labels. By P.11's
raw transport, the two writings have the same guards and results under that
bijection. Pair-table pushforward is functorial, and an assertion update at
tokens \texttt{(u,v)} contributes exactly one signed unit at their current class
pair, then at their final class pair under later identifications. Fill and
Identify act on oriented faces by functorial nondegenerate image and set
union. Induction on the prefix establishes the three symbolic formulas for every
legal prefix, hence the endpoint and price statements.

Every order ideal is the prefix set of some linear extension of a finite
poset. If all complete extensions are legal, its actual prefix realizes the
symbolic profile, and every currently available next occurrence must pass
the displayed guard. A cycle of Identify edges has a last edge in every
linear extension; its endpoints are already joined, so a legal complete
extension is unrealizable. Conversely assume the forest and all ideal-local
guards. Starting at the empty ideal, induction makes every available next
event legal and identifies its native output with the symbolic profile of
the enlarged ideal. Thus every linear extension succeeds. A smallest bad
ideal in a failed test has all proper prefixes admitted, giving the claimed
replayable negative witness. Adjacent swaps of incomparable events connect
any two linear extensions; the same ideal endpoints make their squares
commute on raw states up to the uniquely induced token permutation. This is
endpoint coherence, not an unprovided 2-cell between the run objects.

The added names matter. Start with one term and zero tables, birth tokens
\texttt{b} and \texttt{c} in either order, then add one \texttt{R} assertion at \texttt{(b,c)}. Both
orders are legal and their final tables agree after aligning the token
classes; the correct native descriptor is \texttt{(1,2)} in one order and \texttt{(2,1)}
in the other. Holding one literal numeric descriptor fixed would ask a
different question. A guard obstruction is equally concrete: from one
initial term and zero tables, birth \texttt{b}, then plan AddAssertion and Remove
in \texttt{R} at \texttt{(0,b)} with no precedence between those two. The order birth,
Remove, AddAssertion refuses. Adding AddAssertion-before-Remove precedence
admits its linear extension, and the ideal-local table test sees both cases.
P.13 is recovered when no later descriptor names a birth token and the
event poset is discrete: its forest, initial Remove-supply and final
Fill-survival conditions are the simplified all-ideal guards.

Stable token names and event precedence are \textbf{additional presentation
data}. Any one native run can be annotated noncanonically by tagging its
births and choosing token representatives of the current classes, but the
native enum does not itself assert that two runs share one such presentation.
P.14 therefore advances the selected O03 dynamic carrier/guard question
without proving a global source quotient functor, compatible restrictions
over covers, higher cells on histories, or institutional authorization.

The retained \href{scripts/check_causal_execution.py}{combined execution verifier}
now checks a bounded supplied P.14 presentation against the actual
\texttt{omega\_wdp} five-edit model transitions through a thin Rust JSON adapter. It binds
the initial raw tables and faces, occurrence tags, birth tokens, precedence,
all replayed native states, token-class alignments and protected prices.
Admission is capped at eight events and 10,000 linear extensions. A family
is accepted only after every extension actually executes. The \href{code/causal_certificates/README.md}{retained
fixtures} include all five edit kinds
across 24 legal schedules; a valid poset with an obstructed \texttt{Remove} is
refused. Binary and source hashes are checked before and after replay.
This is native execution evidence for the supplied finite family, not a
formal verification of the native implementation or an exhaustive test proof.

\subsubsection{P.15 Exact price descent on the causal ideal lattice}\label{p.15-exact-price-descent-on-the-causal-ideal-lattice}

P.14 fixes the raw state of each order ideal, up to its token-class writing.
It does \textbf{not} make every price accumulated along a schedule independent of
that schedule. Here is an exact test for an additive price. Fix one P.14
family whose \emph{every} linear extension is legal, and write \texttt{L=Idl(E)} for its
finite lattice of event order ideals. For each cover \texttt{J\ -\textgreater{}\ J+e} choose a value
\texttt{w(J,e)} in an abelian group \texttt{A}. An actual source-derived price must first
be invariant under the token-to-native relabeling of P.14, so that it really
is a function of the ideal and event, not of an unrecorded prefix history.
The theorem below is algebraic for any \texttt{A}; the retained checker handles
integer-vector prices \texttt{A=Z\^{}d} only.

For two distinct events \texttt{e,f} both available after \texttt{J}, all four edges of
their diamond exist. Define its \textbf{price curl}

\[
 \kappa_w(J;e,f)=w(J,e)+w(J+e,f)-w(J,f)-w(J+f,e).
\]

\textbf{Exact price theorem.} The following conditions are equivalent:

\begin{enumerate}
\def\labelenumi{\arabic{enumi}.}
\tightlist
\item
  Every two complete linear extensions from the empty ideal to \texttt{E} have
  the same accumulated \texttt{w} price.
\item
  Every elementary diamond has \texttt{kappa\_w(J;e,f)=0}.
\item
  There is a unique \texttt{Phi:L\ -\textgreater{}\ A} with \texttt{Phi(empty)=0} and
  \texttt{w(J,e)=Phi(J+e)-Phi(J)} on every cover.
\end{enumerate}

When they hold, \emph{every} pair of saturated routes from any \texttt{J} to any
containing \texttt{K} have the same price \texttt{Phi(K)-Phi(J)}. Thus the degree-one
square-closed cochains on this declared ideal lattice are exact. This is a
finite schedule statement, not a cohomology calculation for a proposed
global source site.

\textbf{Proof.} A potential telescopes along every route, proving (3) implies
(1) and (2). To obtain (2) from (1), choose any legal prefix reaching \texttt{J}
and a common continuation from \texttt{J+e+f} to \texttt{E}. The two full runs differing
only by \texttt{e,f} versus \texttt{f,e} have price difference exactly \texttt{kappa\_w(J;e,f)};
by (1) it is zero. For (2) implying (3), any two linear extensions of an
ideal are connected by adjacent swaps of incomparable events. To see this,
move the first event of one extension leftward in the other. Each event it
passes is incomparable with it, because both orders are linear extensions;
then induct on the suffix. A swap changes the path price by precisely one
diamond curl. Zero curls therefore make the sum from the empty ideal to
\texttt{J} independent of the extension. Define that sum as \texttt{Phi(J)}. Appending one
cover gives its required difference. Every ideal is reachable, so the root
normalization makes \texttt{Phi} unique. The same swap argument works between any
two fixed ideals.

For the native five-edit model addition-price pair, \texttt{w(J,AddTerm)=(1,0)},
\texttt{w(J,AddAssertion)=(0,1)}, and the other three kinds have \texttt{(0,0)}.
The curls vanish; \texttt{Phi(J)} is the pair of event-kind counts in \texttt{J}. Likewise
the \emph{difference} \texttt{source\_size(C\_(J+e))-source\_size(C\_J)} is exact with
potential \texttt{source\_size(C\_J)-source\_size(C\_empty)}, since P.12 makes the state
quantity invariant under relabeling. These are different price definitions:
an arbitrary fee evaluated \emph{at} a state need not be a state difference.

For a concrete obstruction, start with one five-edit model term, zero counts and
no fillings. Let \texttt{t=AddTerm} and \texttt{a=AddAssertion(R,0,0)} be incomparable.
Both orders are legal and have the same raw endpoint. Charge zero for \texttt{t};
charge \texttt{source\_size(C\_J)} for \texttt{a} when executed after ideal \texttt{J}. The route
\texttt{t,a} charges \texttt{2}, while \texttt{a,t} charges \texttt{1}, so
\texttt{kappa\_w(empty;t,a)=1}. The two complete orders are an exact negative
certificate for schedule-independent price. Endpoint equality cannot
silently replace this protected accumulated-cost question.

The \href{scripts/check_ideal_prices.py}{integer ideal-price checker} enumerates
all ideals and their covers, checks every diamond in \texttt{Z\^{}d}, and returns
either the root-normalized potential or the first nonzero curl together with
two complete topological orders and their different exact totals. Missing
edges, dimension mismatches and cyclic precedence are malformed input, not
negative mathematical instances. An \href{code/unit_tests/test_ideal_price_cohomology.py}{independent schedule-enumeration
test} compares its decisions
with all complete orders on 256 deterministic four-event DAG/price cases,
plus the two displayed positive and negative controls. The checker does not
itself verify P.14's native admissibility premise or source-price binding.

The combined P.14 verifier supplies those premises for its admitted bounded
profile. It derives the two native addition prices, source-size difference
and plain-state-cost difference from actual replay, verifies their ideal
bindings, and invokes the price checker on those exact values. Just
submitting an independently valid potential is insufficient.

If a curl is nonzero, choosing one path to each ideal still defines a
candidate potential and its exact edge differences. Replacing \texttt{w} by those
differences is a \textbf{changed price specification}, not a proof that the original
price descends. Evaluation of the change and any authorization to use it
remain separate. Arbitrary history-sensitive functionals, guarded
subgraphs missing ideal diamonds, and higher cells between raw histories
are outside this theorem; GEN-WDP-O03 remains open at its full scope.

\subsubsection{P.16 Causal swap certificates and a chosen thin history 2-category}\label{p.16-causal-swap-certificates-and-a-chosen-thin-history-2-category}

P.14 gives endpoint-coherent diamonds; it does not assert that a five-edit model
history has a native 2-morphism. Here is one explicit, scoped construction
which retains the ordered 1-morphisms and shows exactly where a new higher
identification is chosen. Fix the same finite tagged event poset \texttt{E}, stable
birth tokens, and a P.14 certificate that \textbf{every} linear extension is
legal. For each ideal \texttt{J}, use its token-class normal form \texttt{C\_J} and keep the
ideal label even if two different ideals happen to have identical packets.

On a cover \texttt{J\ -\textgreater{}\ J+e}, translate \texttt{e} to the current native descriptor,
execute that five-edit model step, and record the unique \emph{token-induced} bijection
from its native output labels to the normal writing \texttt{C\_(J+e)}. The bijection
is unique because it must send each native label to its named token class;
it is not a claim that the underlying packet has no other automorphisms.
P.11 transports the next descriptor through this bijection. A \textbf{framed
history} from \texttt{J} to \texttt{K} is a saturated path of these actual one-step
executions with all intermediate packets, occurrence tags and alignment
maps retained. Concatenation splices at the common framed boundary and
concatenates lists. P.12's exact-state splice and P.11's renaming law make
this associative, with the empty list as identity. The resulting 1-category
has distinct arrows for distinct event orders, even when their endpoints
and addition prices agree. An alignment is comparison data, \textbf{not} a sixth
native edit: transporting later descriptors back through the accumulated
alignments recovers one literal native run, while the framed form keeps the
same boundaries for composition.

For two different events \texttt{e,f} available after \texttt{J}, define an \textbf{elementary
swap certificate} with the two native paths \texttt{e;f} and \texttt{f;e}, their four
ideal-labelled boundary packets, the translated descriptors and the two
token-induced maps to \texttt{C\_(J+e+f)}. P.14 establishes both paths legal and both
endpoint maps correct. Thus the verification record has two genuinely different ordered
1-boundaries and one checked common framed endpoint. Prefix and suffix
framed histories can be pasted onto it; a list of such pasted verification records is a
proof-relevant swap chain. Reversing the list reverses the comparison, and
concatenating two lists composes comparisons vertically. No equality of
distinct proof-relevant swap chains is inferred from their endpoint alone.

\textbf{Connectivity and optimal verification record theorem.} Let \texttt{h,h\textquotesingle{}} be any two framed
histories from \texttt{J} to \texttt{K}. Count unordered pairs of occurrences in \texttt{K\textbackslash{}J}
whose order is opposite in \texttt{h} and \texttt{h\textquotesingle{}}; call this integer \texttt{inv(h,h\textquotesingle{})}.
Comparable events cannot contribute. There is a swap chain from \texttt{h} to
\texttt{h\textquotesingle{}} of exactly \texttt{inv(h,h\textquotesingle{})} elementary verification records, and no shorter chain.
To construct it, scan \texttt{h\textquotesingle{}} from left to right. Move its next wanted event
leftward through the current word until it reaches the wanted position.
Each crossed event lies after it in \texttt{h\textquotesingle{}}, hence is incomparable with it;
the two are simultaneously available after their common prefix ideal.
Their P.14 diamond supplies the required actual native swap certificate. Each
swap fixes exactly one inverted pair and no other pair order. Conversely an
adjacent swap changes only one pair order, so every chain needs at least
\texttt{inv(h,h\textquotesingle{})} steps. This also establishes that the verification record graph of histories
between any fixed pair of ideals is connected.

The \href{scripts/check_ideal_swaps.py}{minimum-swap checker} validates two
linear extensions and emits precisely those prefix ideals, event pairs and
positions. An \href{code/unit_tests/test_ideal_swap_cells.py}{independent BFS test}
replays every emitted diamond and compares its length with shortest-path
search for two selected order pairs in each of 64 four-event forward DAGs;
it also checks a nonempty start ideal and invalid schedules. The checker
verifies the poset and swap layer only. It does not execute native five-edit model
packets; the actual source-step square is supplied by P.14's separate
admission and transport proof.

There is a \textbf{declared thin completion}: for any two framed histories with
the same ideal-labelled source and target, identify \emph{all} swap chains with
those boundaries into one invertible 2-cell. Such a cell exists by the
connectivity theorem. The hom-category is then the indiscrete groupoid on
its distinct history 1-arrows. Horizontal composition is induced by
concatenation and whiskering, and vertical composition is the unique cell
between its boundaries. Identity, associativity and middle-four
interchange hold strictly because both sides are parallel cells in a thin
hom-category. This constructs a strict 2-category for the selected causal
family. Thinness is an \textbf{extra quotient of verification records}, not a property of
native five-edit model or a proof that distinct histories are equal as 1-arrows.

P.15 gives the exact price obstruction to this completion. Send every
object to the one object of the discrete 2-category whose 1-arrows are an
abelian price group \texttt{A} under addition, and send a history to its accumulated
ideal-local price. This extends to a strict 2-functor precisely when every
diamond curl vanishes: the two boundaries of each elementary 2-cell must
have the same target price, and P.15 establishes that condition sufficient for
all histories. Native five-edit model addition prices satisfy it. The legal
AddTerm/assertion square charged by current source size in P.15 has curl
one, so no price-preserving 2-functor sends that square to an identity 2-cell.
The unpriced swap certificate still exists; an evaluator protecting that fee
must refuse the proposed price-preserving comparison.

For three independent events, the six orders form a swap hexagon. Its
native endpoint alignments agree by P.14, and its two parallel swap chains
become equal \textbf{in the declared thin completion}. They remain different
proof-relevant verification records before that quotient. No nontrivial 3-cell between
those verification records, Kan filler, compatible restriction over arbitrary source
covers, result-mark quotient functor, or institutionally authorized history
identification follows. GEN-WDP-O03 therefore remains open beyond this
selected source-specific 2-category.

Within the bounded combined profile, \href{scripts/check_causal_three_cells.py}{swap and local-relation
replay} now binds each swap to the two
actual native histories, their common prefix, framed two-step boundary,
token alignments and protected prices. The older poset-only checker remains
combinatorial; native evidence is credited only to this combined path.

\subsubsection{P.17 Local 3-witnesses between causal swap proofs}\label{p.17-local-3-witnesses-between-causal-swap-proofs}

P.16 made all parallel swap chains equal by a declared thin quotient. For
the same \textbf{selected} P.14 all-order-legal tagged event poset, that equality
has a finite local presentation. Retain each framed actual run as a distinct
1-arrow and each elementary swap certificate as a generating 2-comparison. A
swap proof from one history \texttt{h} to another is a word in adjacent-position
swaps \texttt{s\_i}, with every intermediate event order a linear extension. Its
2-boundary is the pair of framed histories, not just their shared packet.

Adjoin the following invertible \textbf{local 3-witnesses}, with arbitrary
admissible swap-proof prefixes and suffixes whiskered onto either side:

\begin{enumerate}
\def\labelenumi{\arabic{enumi}.}
\tightlist
\item
  A backtrack \texttt{s\_i\ s\_i\ \textless{}-\textgreater{}\ id}, where the two events at positions \texttt{i,i+1}
  are incomparable at the insertion site. Deletion from an already valid
  proof automatically satisfies this condition.
\item
  A distant square \texttt{s\_i\ s\_j\ \textless{}-\textgreater{}\ s\_j\ s\_i} for \texttt{\textbar{}i-j\textbar{}\textgreater{}1}. These swaps touch
  disjoint positions and retain the same two incomparable event pairs.
\item
  A braid hexagon \texttt{s\_i\ s\_(i+1)\ s\_i\ \textless{}-\textgreater{}\ s\_(i+1)\ s\_i\ s\_(i+1)}. A valid side
  swaps each of the three pairs among its three tagged events. All three
  pairs must therefore be incomparable, so the other side is valid too.
\end{enumerate}

Each boundary order of a local witness has its P.16 framed native run and
token-induced alignment. The square and hexagon are comparisons of \textbf{swap
proofs}, not new native permitted edits. In particular, a square here is
not the ideal-local price diamond of P.15, though its constituent swaps
have such diamonds as verification records.

\textbf{Local coherence theorem.} Any two valid swap proofs with the same
framed source history and target history are joined by a finite sequence
of these local 3-witnesses. Number the event positions in the common
source order. A swap proof gives a word in the adjacent generators of the
symmetric group and acts on that order to produce its target; two proofs
with the same boundary therefore represent the same permutation. The
\href{https://archive.ymsc.tsinghua.edu.cn/pacm_download/116/7904-11512_2006_Article_28.pdf}{Coxeter word property},
Theorem 2.2, reduces each word using braid moves and deletion of \texttt{s\_i\ s\_i},
and connects any two reduced expressions by braid moves. In this symmetric
group the braid moves are exactly the distant square and the three-term
hexagon above. The restriction to legal event orders does not break the
argument: deleting a valid backtrack leaves a valid path; exchanging
disjoint swaps preserves both incomparable pairs; and the three-term move
is valid on its alternate route because its original route already swapped
all three pairs. Thus both reductions, their reduced-word connection, and
the reverse of the second reduction remain entirely inside the graph of
linear extensions. P.14 supplies actual native execution and endpoint
alignment at each intermediate order. The result includes the empty and
single-event intervals, where the swap proof is necessarily empty.

A \textbf{3-certificate} is the finite list of such local replacements, naming
the swap-word offset and, for an inserted backtrack, its adjacent position.
Replay checks every intermediate linear extension and the fixed two
boundary histories. This makes the P.16 thin 2-cell the quotient of
proof-relevant swap chains by explicitly witnessed local relations, rather
than an unexplained identification. For a three-event antichain, the two
three-swap routes from one order to its reverse are distinct 2-proofs and
the braid hexagon is their first nontrivial 3-witness. P.15 zero diamond
curls remain exactly the condition for an ideal-local additive price to
respect the generating 2-comparisons; the 3-witnesses introduce no new
price equality between different boundaries.

The \href{scripts/check_causal_three_cells.py}{executable 3-certificate checker}
checks backtrack, distant-square and braid replacements, including
whiskered replacements at a declared swap-word offset. Every intermediate
order must be legal in the bound native family and both framed boundary
histories must remain fixed. Retained tests refuse a comparable-event
braid, an invalid offset, changed boundaries and corrupted intermediate
states. Distinct replacement lists retain distinct certificate identities.
The replay closes the selected certificate problem; the general existence
argument above remains a written theorem, not a consequence of enumeration.

These 3-witnesses are \textbf{chosen presentation data}, not canonical native
five-edit model operations. The theorem says parallel 2-proofs are connected;
it does not identify different 3-certificates, furnish 4-cells or a
coherent tricategory, define restrictions over arbitrary covers, extend to
histories with dynamic carrier changes outside P.14, or grant authority.
GEN-WDP-O03 remains open at those broader scopes.

\subsubsection{P.18 Guarded histories and the componentwise price criterion}\label{p.18-guarded-histories-and-the-componentwise-price-criterion}

P.14--P.17 assumed that \textbf{every} linear extension of the declared causal
event family is legal. Drop that assumption, but retain stable event tokens,
birth-before-use precedences, the actual five-edit model guards, and at least one
legal complete execution. P.14's endpoint formula still identifies the
complete raw endpoint of \emph{each legal} execution up to token-class writing.
It does not say that those executions can be connected by legal swaps.

Let \(\mathcal H\) be the finite set of legal complete linear extensions.
Build a graph \(\Gamma\) on \(\mathcal H\): an edge exchanges two adjacent,
incomparable tagged events when both resulting complete executions are
legal. Such an edge has a common prefix ideal \(J\), the two orders
\(e,f\) and \(f,e\), and a common suffix. Both two-step prefixes end at
the same formal state of \(J+e+f\) by P.14, so the suffix has the same
guard outcomes after token alignment. A proposed swap for which one
two-step prefix refuses is not an edge.

Let \(A\) be an abelian group and let \(w(J,e)\in A\) be an ideal-local
price on each legal cover \(J\to J+e\), invariant under the
token-to-native relabeling as in P.15. The total \(W(h)\) is its sum
along a legal complete history. For a realized edge of \(\Gamma\), put

\[
 \kappa_w(J;e,f)=w(J,e)+w(J+e,f)-w(J,f)-w(J+f,e).
\]

\textbf{Guarded componentwise price theorem.} The totals \(W(h)\) are constant
on each connected component of \(\Gamma\) if and only if the price curl
vanishes on every realized edge. They are constant on \emph{all} legal complete
histories if and only if those curls vanish \textbf{and} one representative
total from every connected component has the same value. In particular,
zero local curls suffice without the second check when \(\Gamma\) is
connected. P.15 is the all-order-legal special case, where adjacent swaps
connect every pair of linear extensions. This theorem does not claim a
potential on the entire ideal lattice: some ideals may have no legal
prefix, and the legal-prefix graph need not be the full ideal lattice.

\textbf{Proof.} For one edge, the common prefix and suffix cancel in the
difference of the two totals, leaving exactly \(\kappa_w(J;e,f)\).
Therefore constancy on a component forces every edge curl to vanish.
Conversely, zero edge curls make the total unchanged along every edge
path. Comparing one total per component is then both necessary and
sufficient for global constancy. All statements use actual legal full
histories, so a diamond that cannot be continued to a complete legal
execution is not silently promoted to a complete-run price condition.

\textbf{Actual five-edit-kernel counterexample to connectivity.} Start with
one term, zero \texttt{R} and \texttt{S} count tables, and no filled faces. Use only
the native \texttt{AddAssertion(R,0,0)} and \texttt{Remove(R,0,0)} kinds, with five
distinct tagged occurrences:

\begin{longtable}[]{@{}lll@{}}
\toprule\noalign{}
Tag & Descriptor & Declared predecessor \\
\midrule\noalign{}
\endhead
\bottomrule\noalign{}
\endlastfoot
\texttt{p} & \texttt{AddAssertion(R,0,0)} & none \\
\texttt{x} & \texttt{Remove(R,0,0)} & none \\
\texttt{a} & \texttt{AddAssertion(R,0,0)} & \texttt{x} \\
\texttt{y} & \texttt{Remove(R,0,0)} & none \\
\texttt{b} & \texttt{AddAssertion(R,0,0)} & \texttt{y} \\
\end{longtable}

The edges \texttt{x\textless{}a} and \texttt{y\textless{}b} are permitted additional causal precedences
in P.14. Initially \texttt{x} and \texttt{y} fail the source's positive-count guard;
\texttt{a} and \texttt{b} are unavailable by precedence. Hence \texttt{p} must occur first,
leaving count one. If \texttt{x} occurs next, the count becomes zero, so \texttt{y}
cannot occur before \texttt{a}; precedence forces the unique continuation
\texttt{a,y,b}. The symmetric choice \texttt{y} forces \texttt{b,x,a}. Thus the \textbf{only} legal
complete histories are

\[
 h_x=(p,x,a,y,b),\qquad h_y=(p,y,b,x,a).
\]

Both reach the same raw one-term state with \texttt{R(0,0)=1}, zero \texttt{S}, no
faces, and native addition-price pair \texttt{(0,3)}. Yet neither has a legal
adjacent swap: swapping \texttt{p} with its following Remove fails the initial
count guard; swapping a Remove with its paired Add violates precedence;
swapping the middle Add with the other Remove attempts a Remove at count
zero. Therefore \(\Gamma\) consists of two isolated vertices. Equal
raw endpoints do \textbf{not} imply legal-swap connectivity.

To see why the price theorem needs its component comparison, set
\texttt{w(J,e)=1} only for \texttt{J=\{p\},\ e=x}, and zero on every other legal cover.
This is an ideal-local integer price. There is no realized legal swap
diamond, so every local curl condition holds vacuously, but
\texttt{W(h\_x)=1} and \texttt{W(h\_y)=0}. This does not contradict the native
addition price \texttt{(0,3)}, which is equal on both histories.

The counterexample settles guarded \textbf{connectivity} negatively and
classifies price descent componentwise. It does not decide \textbf{coherence}
of parallel swap proofs within a connected component. P.17's Coxeter
reductions used paths through all linear extensions; with illegal
schedules removed, each backtrack, distant square and braid hexagon must
be checked for legality before it is a 3-witness. Connectivity alone
does not supply those checks or identify distinct 3-certificates.
P.20 gives a sufficient guard condition under which every component is
coherent, and states what the bounded native search found when the condition
fails. GEN-WDP-O03 remains open at its broader source/site scope.

\subsubsection{P.19 Admission, transport, and finite continuation}\label{p.19-admission-transport-and-finite-continuation}

Checked repair and the admission, adapter-soundness and continuation protocol
have different jobs in the same declared execution. Checked repair tests the
proposed mathematical witness against its carrier and task. The protocol
separates admission from evaluation, requires the supplied
representation change to meet its transport obligations, and accounts for
continuation. An authorization requirement remains independently supplied.
The \href{docs/coda_prime_integration.md}{consolidation map} locates the component
sources and their original proof and executable evidence inside this package.

\textbf{Admission precedence.} Let \(A_P(w)\) be the admission predicate for a
candidate \(w\) on the original task \(P\). Define an evaluation operation
only on pairs equipped with evidence of \(A_P(w)\). Then every successful
evaluation carries an admitted candidate. This is immediate by projection
from that dependent input. In the retained Lean model of the protocol, the input is
instead a packet with a Boolean \texttt{admitted} field and a partial \texttt{measure}
operation. Its theorem \texttt{Prime.admission\_precedes\_measurement} establishes the
corresponding Boolean gating law. Neither formulation derives the admission
predicate, its checker soundness, or authority from packet syntax.

\textbf{Two independent transport obligations.} Suppose an admitted change sends
\(P\) to \(P'\), and let \(\operatorname{Sol}(P)\) retain the original
solution relation. A backward map
\(b:\operatorname{Sol}(P')\to\operatorname{Sol}(P)\) supplies sound
decoding. A forward map
\(f:\operatorname{Sol}(P)\to\operatorname{Sol}(P')\) supplies solution
coverage. Backward maps compose contravariantly and forward maps compose
covariantly along a finite history, by ordinary function composition. The
laws needed for original-witness recovery must be stated separately; these
two maps alone need not be inverses. Soundness can hold vacuously when the
target has no solutions, so it cannot replace coverage. This is the same
distinction used in §§20.6--20.8 and the finite constraint representation.

\textbf{Continuation conjunction.} Fix the nine declared continuation conditions
\(c_1,\ldots,c_9\): an obligation, its witness, a declared next carrier,
adapter soundness, detector adequacy, minimal lift, skeleton preservation,
remaining budget, and absence of the declared terminal refusal. Set

\[
 S=\bigwedge_{i=1}^{7}c_i,\qquad B=c_8,\qquad N=c_9.
\]

Then \(\bigwedge_{i=1}^{9}c_i\iff S\land B\land N\), by associativity
of conjunction. The retained theorem \texttt{Prime.cc9\_implies\_cc3} establishes one
direction for this explicit grouping. An implementation whose three fields
omit an original conjunct has only a weaker condition; matching field counts
does not establish equivalence. Verifying the supplied predicates is also
different from just receiving nine \texttt{true} values.

\textbf{Finite operational stopping.} Assume a shared initial budget
\(B_0\ge0\), positive cost floor \(\delta>0\), and executed operation
costs \(c_i\ge\delta\). Require every admitted prefix to satisfy
\(\sum_{i=1}^n c_i\le B_0\), including operations executed after
recursive re-entry. Then

\[
 n\delta\le\sum_{i=1}^n c_i\le B_0,
 \qquad n\le\lfloor B_0/\delta\rfloor.
\]

For positive rational costs, multiply by a common denominator to obtain the
integer inequality. The retained \texttt{Prime.ops\_le\_floor\_div} mechanizes its
natural-number arithmetic form under the supplied cumulative bound. It does
not establish that the concrete runtime charges every operation. A re-entry
that resets the budget violates the prefix premise and escapes the theorem.
Uncharged proposal rounds likewise require their own finite bound or charge.

Stopping at this bound yields resource exhaustion when the task has not been
resolved. It supplies no impossibility theorem and no zero-residual witness.
For example, a controller with budget two and unit-cost operations stops
before a declared three-operation solution even though the task is solvable.
Finite impossibility requires the exhaustive coverage certificate of §20.8;
guaranteed resolution and its optimal cost require the additional arena and
certificate hypotheses of Appendix S.

\textbf{The common finite specification.} A checked resolution therefore binds the
original question, its admitted representation history, the submitted
witness, the independent checker result, and the shared cost ledger. Local
satisfaction may still leave an overlap disagreement; a quotient-zero
diagnostic may still lack a realizable solving operation; finite protocol
stopping may still leave an undetermined question. These are different
obligations in one account, and none can discharge the others by renaming a
terminal status. The finite Python application retains its narrower specification:
its local authorization variant requires a caller-configured, task-bound
grant and is not a proof of institutional standing or native Rust execution.

\paragraph{P.19.a Rational instruments, no-signalling and an algebraic CHSH certificate (DP-9, audited 2026-10-06)}\label{p.19.a-rational-instruments-no-signalling-and-an-algebraic-chsh-certificate-dp-9-audited-2026-10-06}

\textbf{Statement.} Let \texttt{X,Y,Z} be finite nonempty carriers, and let \texttt{A,B} be
finite nonempty outcome sets.

\begin{itemize}
\tightlist
\item
  A \emph{classical rational instrument} \texttt{K:X→Y} is a family of matrices
  \texttt{K\_a∈Q\_\{≥0\}\^{}\{Y×X\}} with \texttt{Σ\_\{a,y\}K\_a(y,x)=1} for every \texttt{x}. Write
  \texttt{U(K)=Σ\_aK\_a} for its outcome-forgetting map.
\item
  For a second instrument \texttt{L:Y→Z}, put \texttt{(L∘K)\_\{(a,b)\}=L\_bK\_a}.
\end{itemize}

Then:

\begin{enumerate}
\def\labelenumi{\arabic{enumi}.}
\tightlist
\item
  \texttt{L∘K} is an instrument.
\item
  Composition is associative, up to reassociation of outcome tuples.
\item
  Forgetting is a homomorphism: \texttt{U(L∘K)=U(L)U(K)}.
\item
  For total maps \texttt{f:A→A′} and \texttt{g:B→B′}, coarse-graining the composite gives
  \texttt{(Σ\_\{g(b)=b′\}L\_b)(Σ\_\{f(a)=a′\}K\_a)} at \texttt{(a′,b′)}. This needs \texttt{L} to be fixed
  independently of the first outcome. An adaptive continuation \texttt{L\^{}a} forgets
  to \texttt{Σ\_\{a,b\}L\^{}a\_bK\_a} instead.
\item
  For a joint table \texttt{p(a,b)} and a complete local instrument \texttt{K} on \texttt{B}, put
  \texttt{q(o,a,b′)=Σ\_bK\_o(b′,b)p(a,b)}. Summing over all \texttt{o} and \texttt{b′} returns the
  \texttt{A} marginal of \texttt{p}.
\item
  The PR box \texttt{p(a,b\textbar{}x,y)=1/2} when \texttt{a⊕b=xy}, and \texttt{0} otherwise, satisfies
  every no-signalling equation. It has no global assignment table.
\item
  Let \texttt{r} be the algebraic number with \texttt{r²=2} and \texttt{1\textless{}r\textless{}3/2}. For every
  finite bipartite density and all Hermitian involutions \texttt{A\_x}, \texttt{B\_y} acting
  on separate tensor factors, the CHSH value
  \texttt{S=E\_00+E\_01+E\_10−E\_11} satisfies \texttt{−2r≤S≤2r}. The bound is attained.
\item
  Every Hermitian expectation over \texttt{Q(i)} is rational. So no example over
  \texttt{Q(i)} attains \texttt{S²=8}.
\end{enumerate}

\textbf{Proof.}

\emph{Items 1--4.} Every entry of a product is a finite sum of nonnegative
rationals. For fixed \texttt{x},

\[\sum_{a,b,z,y}L_b(z,y)K_a(y,x)
 =\sum_{a,y}K_a(y,x)\Big[\sum_{b,z}L_b(z,y)\Big]=1.\]

For a third instrument \texttt{M}, both bracketings give the \texttt{(a,b,c)} entry
\texttt{Σ\_\{z,y\}M\_c(w,z)L\_b(z,y)K\_a(y,x)} at \texttt{(w,x)}. Forgetting and coarse-graining
interchange the same finite sums.

\emph{Item 5.} Interchange sums:
\texttt{Σ\_\{o,b′\}q=Σ\_bp(a,b)Σ\_\{o,b′\}K\_o(b′,b)=Σ\_bp(a,b)}.

\emph{Item 6.} Each local marginal is \texttt{1/2}, so the no-signalling equations hold.
A global bit assignment \texttt{(a0,a1,b0,b1)\ ∈\ \{0,1\}⁴} would need \texttt{a0⊕b0=0}, \texttt{a0⊕b1=0},
\texttt{a1⊕b0=0} and \texttt{a1⊕b1=1}. Their xor gives \texttt{0=1}.

A nonnegative global distribution that reproduces the zero cells is supported
on such assignments. So no such distribution exists. Converting bits \texttt{a\_x,b\_y\ ∈\ \{0,1\}} to signs \texttt{A\_x=(-1)\^{}\{a\_x\}} and \texttt{B\_y=(-1)\^{}\{b\_y\}} in \texttt{\{-1,+1\}}, a deterministic
assignment has CHSH value \texttt{A0(B0+B1)+A1(B0−B1)=±2}, and averaging keeps
\texttt{\textbar{}S\textbar{}≤2}. The PR box has \texttt{S=4}.

\emph{Item 7.} Let \texttt{B=A0⊗(B0+B1)+A1⊗(B0−B1)}. Put
\texttt{F0=A0⊗I−I⊗(B0+B1)/r} and \texttt{F1=A1⊗I−I⊗(B0−B1)/r}. Expand, using:

\begin{itemize}
\tightlist
\item
  \texttt{A\_x²=B\_y²=I};
\item
  the fact that operators on different factors commute;
\item
  \texttt{r²=2}.
\end{itemize}

The two mixed terms \texttt{±I⊗(B0B1+B1B0)/2} cancel. This yields

\[F_0^*F_0+F_1^*F_1=4I-rB,\qquad 2rI-B=\tfrac1r(F_0^*F_0+F_1^*F_1)\succeq0.\]

Reversing both inner signs gives \texttt{2rI+B⪰0}. For \texttt{ρ=Σ\_jd\_jv\_jv\_j\^{}*} with
\texttt{d\_j≥0}, \texttt{Tr(ρF\^{}*F)=Σ\_jd\_j‖Fv\_j‖²≥0}. So \texttt{\textbar{}Tr(ρB)\textbar{}≤2r}.

The cancellation does not use \texttt{B0B1=B1B0}. The identity therefore holds for
any Hermitian involutions on the second factor. The audit checked it with
\texttt{B0=Z} and \texttt{B1={[}{[}3/5,4/5{]},{[}4/5,−3/5{]}{]}}.

\emph{Attainment.} Take \texttt{Z=diag(1,−1)} and \texttt{X={[}{[}0,1{]},{[}1,0{]}{]}}. Put \texttt{A0=Z},
\texttt{A1=X}, \texttt{B0=(Z+X)/r} and \texttt{B1=(Z−X)/r}. Let \texttt{ρ=vvᵀ/2} with \texttt{v=(1,0,0,1)ᵀ}.
- \texttt{ZX+XZ=0}, so each \texttt{B\_y²=I}.
- \texttt{Tr(ρZ⊗Z)=Tr(ρX⊗X)=1} and \texttt{Tr(ρZ⊗X)=Tr(ρX⊗Z)=0}.
- So the four correlators are \texttt{(1/r,1/r,1/r,−1/r)}.
- Hence \texttt{S=4/r=2r} and \texttt{S²=8}. The bracket \texttt{2\textless{}S\textless{}3} selects the positive
root.

In this instance \texttt{F0=Z⊗I−I⊗Z} and \texttt{F1=X⊗I−I⊗X}, and
\texttt{F0²+F1²=4I−2(Z⊗Z+X⊗X)}.

Arithmetic is exact in pairs \texttt{(u,v)} representing \texttt{u+vr}, with
\texttt{(u,v)(s,t)=(us+2vt,ut+vs)}.

\emph{Item 8.} The trace of a matrix over \texttt{Q(i)} lies in \texttt{Q(i)}. Hermiticity makes
the imaginary part zero, so the expectation is rational. A rational \texttt{S} with
\texttt{S²=8} would give a rational square root of \texttt{2}, and the parity argument of
P.6.a rules that out.

\textbf{Worked instrument.} Take \texttt{K0=diag(1/2,1/3)}, \texttt{K1=diag(1/2,2/3)},
\texttt{L0={[}{[}1/3,1/2{]},{[}0,0{]}{]}} and \texttt{L1={[}{[}0,0{]},{[}2/3,1/2{]}{]}}. Then

\begin{itemize}
\tightlist
\item
  \texttt{L0K0={[}{[}1/6,1/6{]},{[}0,0{]}{]}} and \texttt{L1K0={[}{[}0,0{]},{[}1/3,1/6{]}{]}};
\item
  \texttt{L0K1={[}{[}1/6,1/3{]},{[}0,0{]}{]}} and \texttt{L1K1={[}{[}0,0{]},{[}1/3,1/3{]}{]}}.
\end{itemize}

Their sum is \texttt{U(L)U(K)={[}{[}1/3,1/2{]},{[}2/3,1/2{]}{]}}. On \texttt{p=(2/5,3/5)ᵀ} the four
masses are \texttt{1/6}, \texttt{7/30}, \texttt{4/15} and \texttt{1/3}, summing to \texttt{1}. The forgotten
state is \texttt{(13/30,17/30)ᵀ}.

\textbf{Correction to the dropped-nonnegativity control.} The audit wrote: ``If
nonnegativity is dropped, one-outcome \texttt{{[}{[}2,0{]},{[}-1,1{]}{]}} has each column sum
one but sends \texttt{(1,0)\^{}T} to a vector with negative mass.'' The image is
\texttt{(2,−1)ᵀ}, whose total mass is \texttt{1}. \textbf{DP-9′ supersedes that sentence:} the
image has a \emph{negative entry}, \texttt{−1}, with total mass \texttt{1}. The control still
refutes nonnegativity, but by a negative cell, not by negative mass.

\textbf{Further controls.}

\begin{itemize}
\tightlist
\item
  \emph{Adaptive continuation.} \texttt{K0=diag(1,0)}, \texttt{K1=diag(0,1)}, \texttt{L\^{}0=I} and
  \texttt{L\^{}1={[}{[}0,1{]},{[}1,0{]}{]}} forget to \texttt{{[}{[}1,1{]},{[}0,0{]}{]}}. The unconditioned average
  gives \texttt{{[}{[}1/2,1/2{]},{[}1/2,1/2{]}{]}} instead.
\item
  \emph{Postselection.} For \texttt{p={[}{[}1/3,1/6{]},{[}1/6,1/3{]}{]}}, the \texttt{A} marginal is
  \texttt{(1/2,1/2)}. Keeping only \texttt{B=0} and renormalizing gives \texttt{(2/3,1/3)}.
\item
  \emph{Normalization dropped.} The scalar \texttt{{[}2{]}} doubles mass.
\item
  \emph{Unnormalized state.} \texttt{2ρ} gives \texttt{4r}.
\item
  \emph{Non-involutions.} \texttt{A0=A1=2I} with \texttt{B0=B1=I} gives \texttt{S=4}.
\item
  \emph{Non-Hermitian involutions.} \texttt{A0=A1={[}{[}1,4{]},{[}0,−1{]}{]}}, with \texttt{B0=B1=I} on the
  plus state, gives \texttt{S=4}.
\item
  \emph{No tensor separation.} On one qubit, take \texttt{A0=Z}, \texttt{A1=−X}, \texttt{B0=Z} and
  \texttt{B1=X}. Then \texttt{B=2I+2ZX}, which is not Hermitian. Its expectation on
  \texttt{(I+Y)/2} is \texttt{2+2i}. With \texttt{A1=X} instead the operator is \texttt{0}.
\end{itemize}

\textbf{Evidence grade.}

\begin{itemize}
\tightlist
\item
  \emph{Written proof:} items 1--8, the correction and the controls.
\item
  \emph{Exact finite replay:}

  \begin{itemize}
  \tightlist
  \item
    \texttt{dp9\_instrument\_products}, \texttt{dp9\_masses}, \texttt{dp9\_forgotten\_state},
    \texttt{dp9\_associativity}, \texttt{dp9\_forgetting\_homomorphism};
  \item
    \texttt{dp9\_adaptive\_vs\_average}, \texttt{dp9\_postselection},
    \texttt{dp9\_local\_marginals\_zero}, \texttt{dp9\_pr\_box};
  \item
    \texttt{dp9\_chsh\_correlators\_and\_value}, \texttt{dp9\_root\_selector},
    \texttt{dp9\_sos\_identity\_upper\_and\_lower}, \texttt{dp9\_new\_sos\_generic\_involutions};
  \item
    \texttt{dp9\_hermiticity\_control}, \texttt{dp9\_nonlocal\_control\_c1\_is\_zero},
    \texttt{dp9\_nonlocal\_control\_c2};
  \item
    \texttt{dp9\_line779\_negative\_entry\_not\_negative\_mass}.
  \end{itemize}
\item
  \emph{Lean kernel:} the axiom report is \texttt{lean4/component\_quantum\_axioms.txt}.

  \begin{itemize}
  \tightlist
  \item
    \texttt{WitnessedDifference.ComponentQuantum.pauli\_relations};
  \item
    \texttt{WitnessedDifference.ComponentQuantum.sos\_identity},
    \texttt{WitnessedDifference.ComponentQuantum.sos\_factors\_symmetric};
  \item
    \texttt{WitnessedDifference.ComponentQuantum.chshBs\_table},
    \texttt{WitnessedDifference.ComponentQuantum.bs\_quad\_le},
    \texttt{WitnessedDifference.ComponentQuantum.bs\_quad\_bell};
  \item
    \texttt{WitnessedDifference.ComponentQuantum.deterministic\_chsh},
    \texttt{WitnessedDifference.ComponentQuantum.pr\_box\_not\_deterministic};
  \item
    \texttt{WitnessedDifference.ComponentQuantum.single\_qubit\_minus\_x},
    \texttt{WitnessedDifference.ComponentQuantum.single\_qubit\_plus\_x},
    \texttt{WitnessedDifference.ComponentQuantum.single\_qubit\_minus\_x\_not\_symmetric}.
  \end{itemize}
\item
  \emph{Native exact implementation:}

  \begin{itemize}
  \tightlist
  \item
    \texttt{omega\_finite\_probability::classical\_instrument} provides fixed and
    adaptive composition, forgetting, the PR box, signalling witnesses,
    local-model verification and CHSH. It finds that the PR box has no local
    table by exhaustive search over the \texttt{16} deterministic assignments.
  \item
    \texttt{omega\_quantum\_physics::finite\_chsh\_correlations} admits the extreme
    value only for parties embedded in separate tensor factors. It takes the
    sign from an exact ordered test, not from \texttt{S²}.
  \item
    Both single-carrier controls are refused. With \texttt{A1=X}, the Bell operator
    is zero. With \texttt{A1=−X}, it is \texttt{2I+2ZX}.
  \end{itemize}
\item
  \emph{Unresolved residual:}

  \begin{itemize}
  \tightlist
  \item
    The Lean module states the bound in the rescaled form for
    \texttt{Z⊗Z+X⊗X}, which carries no \texttt{r}. It records non-symmetry of the
    one-qubit operator rather than the value \texttt{2+2i}.
  \item
    Items 1--5 are written proof with finite replay. They are not stated in
    Lean.
  \end{itemize}
\end{itemize}

\paragraph{P.19.b Exact density and channel representations (DP-8, audited 2026-10-06)}\label{p.19.b-exact-density-and-channel-representations-dp-8-audited-2026-10-06}

\textbf{Setting.} Fix positive finite dimensions and an ordered basis. Take
\texttt{F=Q}, with transpose as the adjoint, or \texttt{F=Q(i)}, with conjugate transpose.

\begin{itemize}
\tightlist
\item
  A \emph{density} is a Hermitian \texttt{ρ} over \texttt{F} with trace one and nonnegative
  quadratic form.
\item
  A linear map \texttt{Φ} is given on the matrix units \texttt{E\_ij}. Its input-first Choi
  matrix is \texttt{J\_Φ=Σ\_\{ij\}E\_ij⊗Φ(E\_ij)}.
\end{itemize}

\textbf{Statement.}

\begin{enumerate}
\def\labelenumi{\arabic{enumi}.}
\tightlist
\item
  \emph{(PSD decision.)} For a Hermitian \texttt{H} over \texttt{F}, the following are
  equivalent:

  \begin{itemize}
  \tightlist
  \item
    \texttt{H⪰0};
  \item
    every principal minor is nonnegative;
  \item
    there is an exact factorization \(PHP^*=LDL^*\), where \(P\) is a
    permutation matrix, \(L\) is unit lower triangular, and \(D\) is diagonal
    with nonnegative entries.
  \end{itemize}

  When \texttt{H} is not PSD, the procedure returns an exact vector \texttt{v} with
  \texttt{v\^{}*Hv\textless{}0}. Leading principal minors alone do not suffice.
\item
  \emph{(Choi.)} \texttt{Φ} is completely positive iff \texttt{J\_Φ⪰0}. It is trace preserving
  iff \texttt{Tr\_outJ\_Φ=I\_in}.
\item
  \emph{(Weighted Kraus.)} When \texttt{J\_Φ⪰0}, there are rational weights \texttt{d\_s≥0} and
  operators \texttt{K\_s} over \texttt{F} with \texttt{Φ(X)=Σ\_sd\_sK\_sXK\_s\^{}*}. When \texttt{Φ} is also trace
  preserving, \texttt{Σ\_sd\_sK\_s\^{}*K\_s=I}. A unit-weight family over \texttt{F} always
  exists: a weight \texttt{p/q} becomes \texttt{pq} copies of \texttt{K/q}.
\item
  \emph{(Realification.)} For \texttt{M=A+iB}, put \texttt{R(M)={[}{[}A,−B{]},{[}B,A{]}{]}}. Then:

  \begin{itemize}
  \tightlist
  \item
    \texttt{R(MN)=R(M)R(N)}, \texttt{R(M\^{}*)=R(M)ᵀ} and \texttt{Tr\ R(M)=2\ Re\ Tr\ M};
  \item
    \texttt{ρ\_R=R(ρ)/2} is a real density;
  \item
    \texttt{Tr(ρ\_RR(E))=Tr(ρE)} for every Hermitian \texttt{E}.
  \end{itemize}

  Realification of a tensor product is not the product of the
  realifications.
\end{enumerate}

\textbf{Proof.}

\emph{Item 1.} Write \texttt{H={[}{[}a,b\^{}*{]},{[}b,B{]}{]}} with \texttt{a\textgreater{}0}, and put \texttt{S=B−bb\^{}*/a}.
Completing the square gives

\[(x,y)^*H(x,y)=a\,|x+b^*y/a|^2+y^*Sy,\]

so \texttt{H⪰0} iff \texttt{S⪰0}, and this is one LDL step. Each case either reduces the
dimension or gives a witness:

\begin{itemize}
\tightlist
\item
  A negative diagonal entry gives a basis-vector witness.
\item
  A zero diagonal entry with a nonzero entry \texttt{c} beside it gives a witness.
  In \texttt{{[}{[}0,c{]},{[}c\^{}*,d{]}{]}} with \texttt{d≥0}, take \texttt{v=(−tc,1)} with
  \texttt{t=(d+1)/(2\textbar{}c\textbar{}²)}. Then \texttt{v\^{}*Hv=−1}.
\item
  A zero row and column is deleted.
\end{itemize}

A witness \texttt{y} found for \texttt{S} lifts to \texttt{(−b\^{}*y/a,y)} for \texttt{H}. Each step lowers
the dimension, so the procedure terminates.

For the minors:

\begin{itemize}
\tightlist
\item
  \emph{Necessity.} A principal submatrix of a PSD matrix is PSD, since it extends
  vectors by zero coordinates.
\item
  \emph{Sufficiency.} For every index set \texttt{T} of the remaining block,
  \texttt{det(S\_T)=det(H\_\{\{1\}∪T\})/a≥0}. Induct on the dimension. If no diagonal
  entry is positive, the two-by-two minors force \texttt{H=0}.
\end{itemize}

\emph{Item 2.} Put \texttt{Ω=Σ\_ie\_i⊗e\_i}.

\begin{itemize}
\tightlist
\item
  (⇒) Complete positivity gives \texttt{(id⊗Φ)(ΩΩ\^{}*)=J\_Φ⪰0}.
\item
  (⇐) Item 1 gives \texttt{J\_Φ=Σ\_sd\_sv\_sv\_s\^{}*}. Unvectorize with
  \texttt{(K\_s)\_\{ri\}=(v\_s)\_\{i·d\_out+r\}}. Comparing entries gives
  \texttt{Φ(X)=Σ\_sd\_sK\_sXK\_s\^{}*} on every matrix unit. Each term keeps every ancilla
  extension PSD.
\item
  The \texttt{(i,j)} entry of \texttt{Tr\_outJ\_Φ} is \texttt{Tr\ Φ(E\_ij)}. So the partial-trace
  equation is trace preservation on a basis. By cyclicity of the trace it
  equals \texttt{Σ\_sd\_sK\_s\^{}*K\_s=I}.
\end{itemize}

\emph{Item 3.} This is the (⇐) step of item 2. For the unit-weight family,
\texttt{pq·(1/q²)=p/q}.

\emph{Item 4.} The identities follow by block multiplication. For \texttt{z=x+iy},
\texttt{z\^{}*Hz=(x,y)ᵀR(H)(x,y)}, so \texttt{R} preserves Hermiticity and positivity. Also
\texttt{Tr(ρ\_RR(E))=Re\ Tr(ρE)=Tr(ρE)}.

Dimensions show the tensor statement: \texttt{R(A⊗B)} has size \texttt{2mn}, while
\texttt{R(A)⊗R(B)} has size \texttt{4mn}. Already \texttt{R(1)/2=I2/2}, while the independently
doubled product is \texttt{I4/4}.

\textbf{Worked examples and controls.}

\begin{itemize}
\tightlist
\item
  \emph{Transpose.} On dimension two,
  \texttt{J\_T={[}{[}1,0,0,0{]},{[}0,0,1,0{]},{[}0,1,0,0{]},{[}0,0,0,1{]}{]}}. For \texttt{v=(0,1,−1,0)ᵀ},
  \texttt{v\^{}*J\_Tv=−2}. So the transpose is positive and trace preserving but not
  completely positive.
\item
  \emph{Doubled identity.} \texttt{X↦2X} is completely positive but not trace preserving.
\item
  \emph{Leading-minor trap.} \texttt{diag(0,2,−1)} has every leading principal minor
  equal to \texttt{0} and trace one, but \texttt{e3} gives value \texttt{−1}.
\item
  \emph{Non-Hermitian.} \texttt{(1/2)\ {[}{[}1,1{]},{[}−1,1{]}{]}} has positive principal minors and
  trace one, but it is not Hermitian, so it is not a density.
\item
  \emph{Field boundary.} \texttt{ρ\_y=(1/2)\ {[}{[}1,−i{]},{[}i,1{]}{]}} is a Gaussian-rational density.
  It has no real-rational form in this basis.
\item
  \emph{Realification controls.}

  \begin{itemize}
  \tightlist
  \item
    Complex rank one maps to real rank two.
  \item
    In complex dimension one, \texttt{diag(1,0)} is a real density outside the image
    of \texttt{R}.
  \end{itemize}
\item
  \emph{Hadamard-type channel.} Take
  \texttt{H={[}{[}1,−1,−1,0{]},{[}1,1,0,1{]},{[}1,0,1,−1{]},{[}0,−1,1,1{]}{]}}. Then \texttt{HᵀH=3I4}, and
  \texttt{Φ(X)=HXHᵀ/3} is trace preserving. Its Choi matrix is
  \texttt{vec(H)vec(H)ᵀ/3}, of rank one.
\end{itemize}

\textbf{Correction on minimal Kraus counts.} The audit wrote: ``The coefficient
data thus agree as channels, while their minimal unweighted presentations
differ by field.'' \textbf{DP-8′ supersedes that sentence with the exact minima.}
For the channel above, the least number of unit-weight Kraus operators is:

\begin{longtable}[]{@{}lrl@{}}
\toprule\noalign{}
field & minimum & witness \\
\midrule\noalign{}
\endhead
\bottomrule\noalign{}
\endlastfoot
\texttt{R} & \texttt{1} & \texttt{H/√3} \\
\texttt{Q(i)} & \texttt{2} & \texttt{c\_1H}, \texttt{c\_2H} with \texttt{c\_1=(1+i)/3} and \texttt{c\_2=1/3} \\
\texttt{Q} & \texttt{3} & three copies of \texttt{H/3} \\
\end{longtable}

\emph{Proof.}

\begin{itemize}
\tightlist
\item
  \emph{Every family is a multiple of \texttt{H}.} A rank-one Choi matrix forces every
  Kraus operator to be \texttt{K\_s=c\_sH}, with \texttt{Σ\_s\textbar{}c\_s\textbar{}²=1/3}.
\item
  \emph{One operator.} This needs \texttt{\textbar{}c\textbar{}²=1/3}. Write \texttt{c=(a+ib)/d} with integers of
  joint gcd one. Then \texttt{3(a²+b²)=d²}. The squares modulo three are \texttt{0} and \texttt{1},
  so \texttt{3∣d}. Then \texttt{3∣a²+b²}, so \texttt{3∣a} and \texttt{3∣b}, which contradicts the gcd. So
  no single operator exists over \texttt{Q(i)}, and hence none over \texttt{Q}.
\item
  \emph{Two operators over \texttt{Q(i)}.} \texttt{\textbar{}c\_1\textbar{}²+\textbar{}c\_2\textbar{}²=2/9+1/9=1/3}.
\item
  \emph{Two operators over \texttt{Q}.} This needs \texttt{a²+b²=1/3} with rational \texttt{a,b}.
  Clearing denominators gives the same descent, so no such pair exists.
\item
  \emph{Three operators over \texttt{Q}.} \texttt{3·(1/9)=1/3}.
\end{itemize}

So the field changes the minimal unit-weight count, \texttt{1}, \texttt{2} or \texttt{3}. It does
not change the channel or its Choi matrix.

\textbf{Evidence grade.}

\begin{itemize}
\tightlist
\item
  \emph{Written proof:} items 1--4, DP-8′ and the controls.
\item
  \emph{Exact finite replay:}

  \begin{itemize}
  \tightlist
  \item
    \texttt{dp8\_identity\_choi\_psd}, \texttt{dp8\_transpose\_choi\_matrix},
    \texttt{dp8\_transpose\_choi\_negative}, \texttt{dp8\_transpose\_refused\_by\_ldl},
    \texttt{dp8\_doubled\_identity\_cp\_not\_tp};
  \item
    \texttt{dp8\_zero\_diagonal\_witness\_minus\_one}, \texttt{dp8\_singular\_leading\_minor\_trap},
    \texttt{dp8\_non\_hermitian\_positive\_minors}, \texttt{dp8\_rho\_y\_density};
  \item
    \texttt{dp8\_HtH\_3I}, \texttt{dp8\_choi\_rank\_one\_vecH}, \texttt{dp8\_three\_copy\_completeness},
    \texttt{dp8\_two\_copy\_mutation\_not\_tp}, \texttt{dp8\_mod3\_descent\_residues};
  \item
    \texttt{dp8\_new\_rational\_two\_squares\_absent\_d\_le\_60},
    \texttt{dp8\_new\_rational\_three\_copies\_minimum\_witness},
    \texttt{dp8\_new\_gaussian\_two\_kraus\_completeness},
    \texttt{dp8\_new\_gaussian\_two\_kraus\_same\_choi};
  \item
    \texttt{dp8\_realification\_hom\_trace\_adjoint},
    \texttt{dp8\_realification\_rank\_one\_to\_rank\_two},
    \texttt{dp8\_realification\_tensor\_dims}, \texttt{dp8\_realification\_unit\_scaling}.
  \end{itemize}
\item
  \emph{Lean kernel:} the axiom report is \texttt{lean4/component\_quantum\_axioms.txt}.

  \begin{itemize}
  \tightlist
  \item
    \texttt{WitnessedDifference.ComponentQuantum.choi\_transpose\_negative},
    \texttt{WitnessedDifference.ComponentQuantum.choi\_transpose\_tp},
    \texttt{WitnessedDifference.ComponentQuantum.leading\_minor\_trap};
  \item
    \texttt{WitnessedDifference.ComponentQuantum.kraus\_h\_orthogonal},
    \texttt{WitnessedDifference.ComponentQuantum.three\_copies\_complete};
  \item
    \texttt{WitnessedDifference.ComponentQuantum.sq\_mod\_three},
    \texttt{WitnessedDifference.ComponentQuantum.three\_dvd\_of\_sq},
    \texttt{WitnessedDifference.ComponentQuantum.mod\_three\_descent}.
  \end{itemize}

  \texttt{mod\_three\_descent} states, over natural numbers, that \texttt{3(a²+b²)=d²}
  forces \texttt{a=b=d=0}. That supports both the one-operator absence over \texttt{Q(i)}
  and the two-operator absence over \texttt{Q}.
\item
  \emph{Native exact implementation:}

  \begin{itemize}
  \tightlist
  \item
    \texttt{omega\_quantum\_physics::exact\_hermitian} decides PSD. When the matrix is
    not PSD it returns an exact negative vector with its value.
  \item
    The channel controls replay the transpose Choi matrix with value \texttt{−2}.
    They also replay the absence of a single Gaussian-rational operator for
    the rank-one channel, and the two-operator family. Realification sizes
    replay as \texttt{2ab} against \texttt{4ab}.
  \item
    The \texttt{omega\_wdp} density and channel adapters re-check their premises. A
    shape mismatch is refused.
  \end{itemize}
\item
  \emph{Unresolved residual:}

  \begin{itemize}
  \tightlist
  \item
    The \texttt{omega\_wdp} carriers remain real symmetric over \texttt{Q}.
  \item
    A budget refusal in unit-weight splitting is not a proof of nonexistence.
  \item
    The equivalences of items 1 and 2 are not stated in Lean.
  \item
    The rank-one step and realification are written proof with finite
    replay only.
  \item
    The \texttt{Q} two-square absence is replayed exhaustively only for denominators
    up to \texttt{60}. The general case rests on the descent.
  \end{itemize}
\end{itemize}

\paragraph{P.19.c A derived short-time product-defect bound (DP-7, audited 2026-10-06)}\label{p.19.c-a-derived-short-time-product-defect-bound-dp-7-audited-2026-10-06}

\textbf{Hypotheses.} Import the following finite-dimensional complex Hilbert-space
operations: trace, partial trace, the trace norm \texttt{‖·‖₁}, the operator norm
\texttt{‖·‖\_∞}, and constant-Hamiltonian evolution with \texttt{ħ=1}. They are premises; no
quantum dynamics is constructed from finite scalar verification records.

\begin{itemize}
\tightlist
\item
  Let \texttt{V} be a finite set of sites, with positive local dimensions and a
  declared tensor order.
\item
  Let \texttt{H=Σ\_αh\_α} be a finite sum of Hermitian terms. Each \texttt{h\_α} is
  time-independent, has support \texttt{X\_α} with \texttt{1≤\textbar{}X\_α\textbar{}≤2}, and acts as the
  identity off it.
\item
  Put \texttt{g=max\_v\ Σ\_\{α:v∈X\_α\}‖h\_α‖\_∞}. One-site fields are included.
\item
  Fix \texttt{σ⊆V} with \texttt{k=\textbar{}σ\textbar{}}, the number of sites.
\item
  Let \texttt{ρ(t)=e\^{}\{−itH\}ρ(0)e\^{}\{itH\}}, and write \texttt{ρ\_S(t)} for its reduction to a
  set of sites \texttt{S}.
\item
  Put \texttt{τ\_σ(t)=⊗\_\{i∈σ\}ρ\_i(t)} and \texttt{F\_σ(t)=ρ\_σ(t)−τ\_σ(t)}.
\item
  Put \texttt{D\_σ(t)=‖F\_σ(t)‖₁}, \texttt{n\_σ=F\_σ′(0)} and \texttt{ν\_act=‖n\_σ‖₁}.
\item
  Assume \texttt{F\_σ(0)=0}; a product initial state is sufficient.
\end{itemize}

\textbf{Statement (DP-7′).} For every real \texttt{t≥0},

\[D_\sigma(t)\ \ge\ \nu_{\mathrm{act}}t-4k(k+1)g^2t^2,\]

and for rational bounds \texttt{0≤ν≤ν\_act} and \texttt{G≥g} the same holds with \texttt{ν} and
\texttt{G} in place of \texttt{ν\_act} and \texttt{g}.

When \texttt{k≥2},

\[D_\sigma(t)\ \ge\ \nu t-6k^2G^2t^2 .\]

If \texttt{ν\textgreater{}0} and \texttt{G\textgreater{}0}, the weaker envelope has vertex \texttt{t*=ν/(12k²G²)}, and
\texttt{D\_σ(t)≥νt/2} for \texttt{0≤t≤t*}.

For the trace distance \texttt{T\_σ=D\_σ/2}, the valid weak form is
\texttt{T\_σ(t)≥(ν/2)t−3k²G²t²}.

\textbf{Correction.} The audit summary stated: ``D is at least nu\emph{t minus
4}k\emph{(k+1)}g\textsuperscript{2\emph{t\^{}2, hence the stated 6}k}2 bound.'' \textbf{DP-7′ supersedes that
sentence:} the step from \texttt{4k(k+1)} to \texttt{6k²} needs \texttt{k≥2}, because
\texttt{6k²−4k(k+1)=2k(k−2)}.

At \texttt{k=1} the difference is \texttt{−2}. There the relaxation runs the wrong way, and
the \texttt{6k²} form is stronger than the inequality it claims to follow from. The
bound itself is not violated at \texttt{k=1}: the residual of a single site is
\texttt{F≡0}, so \texttt{ν=0} and both sides are compatible. But ``hence'' is not a valid
step there.

\textbf{Proof.}

\emph{Step 1 (support counting).} For any finite matrices,
\texttt{‖{[}A,B{]}‖\_∞≤2‖A‖\_∞‖B‖\_∞}. Let \texttt{O\_S} be supported on \texttt{s} sites. Only terms
meeting \texttt{S} fail to commute with it, so

\[\|[H,O_S]\|_\infty\le2sg\|O_S\|_\infty .\]

A contributing \texttt{{[}h\_α,O\_S{]}} is supported on \texttt{S∪X\_α}, which has at most \texttt{s+1}
sites, because \texttt{\textbar{}X\_α\textbar{}≤2} and \texttt{X\_α} meets \texttt{S}. Hence

\[\|[H,[H,O_S]]\|_\infty\le4s(s+1)g^2\|O_S\|_\infty .\]

\emph{Step 2 (reduced derivatives).} \texttt{ρ″=−{[}H,{[}H,ρ{]}{]}}. For Hermitian \texttt{A},
\texttt{‖A‖₁=max\{Tr(WA):W=W\^{}*,−I⪯W⪯I\}}. Also, for a density \texttt{ρ},
\texttt{\textbar{}Tr(ρB)\textbar{}≤‖B‖\_∞}. Then partial trace and cyclicity give, for every \texttt{t}:

\begin{itemize}
\tightlist
\item
  \texttt{‖ρ\_S″(t)‖₁≤4s(s+1)g²};
\item
  \texttt{‖ρ\_i′(t)‖₁≤2g};
\item
  \texttt{‖ρ\_i″(t)‖₁≤8g²}.
\end{itemize}

None of these assumes the evolved state stays a product.

\emph{Step 3 (product of marginals).} Apply the product rule:

\[\tau_\sigma''=\sum_i\rho_i''\otimes\bigotimes_{j\ne i}\rho_j
+2\sum_{i<j}\rho_i'\otimes\rho_j'\otimes\bigotimes_{l\ne i,j}\rho_l .\]

The trace norm is multiplicative on tensor products of Hermitian matrices, and
\texttt{‖ρ\_i‖₁=1}. So

\[\|\tau_\sigma''\|_1\le8kg^2+4k(k-1)g^2=4k(k+1)g^2 .\]

With step 2, this gives \texttt{‖F\_σ″(t)‖₁≤8k(k+1)g²}.

\emph{Step 4 (normal slope and remainder).} For trace-zero Hermitian \texttt{Z}, put
\texttt{P\_σZ=Σ\_i(Tr\_\{σ∖i\}Z)⊗(⊗\_\{j≠i\}ρ\_j(0))}.

\begin{itemize}
\tightlist
\item
  \texttt{P\_σZ} has the same one-site marginals as \texttt{Z}. So \texttt{P\_σ} is idempotent, and
  \texttt{n\_σ=(I−P\_σ)ρ\_σ′(0)}.
\item
  Taylor's formula with integral remainder gives
  \texttt{F\_σ(t)=tn\_σ+∫\_0\^{}t(t−u)F\_σ″(u)du}.
\item
  The remainder has norm at most \texttt{8k(k+1)g²·t²/2=4k(k+1)g²t²}.
\item
  The reverse triangle inequality gives the first bound.
\item
  When \texttt{k≥2}, \texttt{2k(k−2)≥0} gives the second.
\end{itemize}

\emph{Step 5 (window).} Put \texttt{c=6k²G²}. Then
\texttt{f(t)=νt−ct²=ν²/(4c)−c(t−ν/(2c))²}, so the vertex is \texttt{t*=ν/(2c)}. Also
\texttt{f(t)−νt/2=t(ν/2−ct)≥0} on \texttt{{[}0,t*{]}}. The remainder estimate holds at every
\texttt{t≥0}; \texttt{t*} only bounds the half-slope claim.

\textbf{Exact models.}

\begin{itemize}
\tightlist
\item
  \emph{Two sites.} Take \texttt{H=X⊗X} on two qubits, starting at \texttt{\textbar{}00⟩⟨00\textbar{}}. Then
  \texttt{g=G=1}.

  \begin{itemize}
  \tightlist
  \item
    \texttt{n=i\textbar{}00⟩⟨11\textbar{}−i\textbar{}11⟩⟨00\textbar{}}. The dual witness \texttt{W=n} has
    \texttt{Tr(Wn)=2}, so \texttt{ν=2}.
  \item
    At \texttt{k=2}, both coefficients equal \texttt{24}, and \texttt{t*=1/24}.
  \item
    At \texttt{t=1/100}, the bound is \texttt{2/100−24/10000=11/625}.
  \item
    The exact defect is \(D(t)=|\sin(2t)|+\tfrac12\sin^2(2t)\). From
    \texttt{x−x³/6≤sin\ x}, \texttt{sin(1/50)≥14999/750000}, which exceeds
    \texttt{11/625=13200/750000}.
  \end{itemize}
\item
  \emph{Three sites.} Add a spectator qubit in \texttt{\textbar{}0⟩⟨0\textbar{}}, with \texttt{H=(X⊗X)⊗I} and
  \texttt{k=3}.

  \begin{itemize}
  \tightlist
  \item
    \texttt{t*=2/108=1/54}.
  \item
    At \texttt{t=1/200} the weak bound is \texttt{1/100−54/40000=173/20000}.
  \item
    The sharper bound is \texttt{1/100−48/40000=11/1250}.
  \end{itemize}
\end{itemize}

\textbf{Controls (each drops one hypothesis).}

\begin{enumerate}
\def\labelenumi{\arabic{enumi}.}
\tightlist
\item
  \emph{Fields omitted from \texttt{g}.} Take \texttt{H=X⊗X+100(Z⊗I)}, so \texttt{H²=10001I} and
  \texttt{g=101}. At \texttt{t0=π/√10001}, which satisfies \texttt{0\textless{}t0\textless{}1/25\textless{}1/24}, the defect is
  \texttt{D(t0)=0}. But the pair-only value \texttt{g=1} would claim \texttt{2t0−24t0²\textgreater{}0}.
\item
  \emph{\texttt{F\_σ(0)≠0}.} Take \texttt{ψ=c\textbar{}00⟩+is\textbar{}11⟩} with \texttt{c=9999/10001} and
  \texttt{s=200/10001}; then \texttt{9999²+200²=10001²}. At
  \texttt{θ=2arctan(1/100)\textless{}1/50\textless{}1/48}, the state is \texttt{\textbar{}00⟩}, so \texttt{D(θ)=0}. Yet
  \texttt{ν=1} would claim \texttt{θ−24θ²\textgreater{}0}.
\item
  \emph{Time dependence.} Take \texttt{H(t)=(1−2t/ε)(X⊗X)} with \texttt{ε=1/100}. It returns
  to \texttt{\textbar{}00⟩} at \texttt{t=ε}, so \texttt{D(ε)=0}. Yet the bound gives
  \texttt{2ε−24ε²=11/625\textgreater{}0}.
\item
  \emph{Higher-body terms.} With \texttt{N=27} sites,
  \texttt{H=X\^{}\{⊗27\}+Σ\_\{j=2\}\^{}\{27\}Z\_j} and \texttt{O=Z\_1}, we have \texttt{g=2}. The double
  commutator has expectation \texttt{104} on the GHZ-minus ray, which exceeds the
  two-local bound \texttt{96}.
\item
  \emph{Half-trace mixing.} In the two-site model, \texttt{T(1/100)≤101/10000}. That is
  less than \texttt{176/10000}, the raw bound misapplied to \texttt{T}.
\item
  \emph{Malformed hosts.} These are refused:

  \begin{itemize}
  \tightlist
  \item
    a non-Hermitian term \texttt{{[}{[}0,1{]},{[}0,0{]}{]}};
  \item
    \texttt{diag(2,−1)}, which is not positive;
  \item
    \texttt{2\textbar{}0⟩⟨0\textbar{}}, which has trace two.
  \end{itemize}
\item
  \emph{Degenerate parameters.}

  \begin{itemize}
  \tightlist
  \item
    At \texttt{k=1}, \texttt{F≡0}.
  \item
    \texttt{G=0} forces every term to be zero.
  \item
    \texttt{ν=0} gives no positive window.
  \end{itemize}
\item
  \emph{Raw commutator.} For \texttt{H=X⊗I} from \texttt{\textbar{}00⟩⟨00\textbar{}}, the raw derivative is
  nonzero. But \texttt{F\_σ(t)=0} for all \texttt{t}, and the normal slope is \texttt{0}.
\end{enumerate}

\textbf{Evidence grade.}

\begin{itemize}
\tightlist
\item
  \emph{Written proof:} steps 1--5, DP-7′ and the controls.
\item
  \emph{Exact finite replay:}

  \begin{itemize}
  \tightlist
  \item
    \texttt{dp7\_support\_growth\_constant}, \texttt{dp7\_P\_sigma\_idempotent},
    \texttt{dp7\_normal\_derivative\_n}, \texttt{dp7\_n\_survives\_P\_sigma},
    \texttt{dp7\_local\_tangent\_projected\_to\_zero};
  \item
    \texttt{dp7\_relaxation\_gap\_is\_2k\_k\_minus\_2},
    \texttt{dp7\_relaxation\_fails\_exactly\_at\_k1};
  \item
    \texttt{dp7\_tstar\_half\_slope\_k2}, \texttt{dp7\_tstar\_half\_slope\_k3},
    \texttt{dp7\_k2\_value\_t\_1\_100}, \texttt{dp7\_k3\_values\_t\_1\_200}, \texttt{dp7\_sin\_enclosure},
    \texttt{dp7\_inert\_spectator\_keeps\_nu\_2}, \texttt{dp7\_I\_pm\_H\_square\_identity};
  \item
    \texttt{dp7\_field\_anticommute\_H2}, \texttt{dp7\_understated\_field\_enclosure},
    \texttt{dp7\_correlated\_pythagorean}, \texttt{dp7\_time\_dependent\_return};
  \item
    \texttt{dp7\_higher\_body\_double\_commutator},
    \texttt{dp7\_higher\_body\_ghz\_expectation};
  \item
    \texttt{dp7\_half\_trace\_control}, \texttt{dp7\_half\_trace\_upper\_value}.
  \end{itemize}
\item
  \emph{Lean kernel:} scalar algebra only. The axiom report is
  \texttt{lean4/component\_quantum\_axioms.txt}.

  \begin{itemize}
  \tightlist
  \item
    \texttt{WitnessedDifference.ComponentQuantum.coefficient\_relax},
    \texttt{WitnessedDifference.ComponentQuantum.coefficient\_relax\_iff},
    \texttt{WitnessedDifference.ComponentQuantum.coefficient\_relax\_fails\_at\_one};
  \item
    \texttt{WitnessedDifference.ComponentQuantum.half\_slope\_window},
    \texttt{WitnessedDifference.ComponentQuantum.envelope\_monotone},
    \texttt{WitnessedDifference.ComponentQuantum.vertex\_value},
    \texttt{WitnessedDifference.ComponentQuantum.envelope\_relax};
  \item
    \texttt{WitnessedDifference.ComponentQuantum.replay\_k2},
    \texttt{WitnessedDifference.ComponentQuantum.replay\_k3},
    \texttt{WitnessedDifference.ComponentQuantum.sine\_enclosure\_comparison};
  \item
    \texttt{WitnessedDifference.ComponentQuantum.field\_control\_square},
    \texttt{WitnessedDifference.ComponentQuantum.correlated\_control\_normalization},
    \texttt{WitnessedDifference.ComponentQuantum.time\_dependent\_control}.
  \end{itemize}
\item
  \emph{Native exact implementation:}

  \begin{itemize}
  \tightlist
  \item
    \texttt{omega\_quantum\_physics::product\_defect\_bound} checks premises only. It
    certifies the premises, the lower envelope and the half-slope endpoint.
  \item
    It accepts the weak \texttt{6k²} coefficient only for \texttt{k≥2}.
  \item
    It refuses an understated field norm and a correlated initial state.
  \item
    It names the analytic remainder as an explicit residual and does not
    discharge it.
  \end{itemize}
\item
  \emph{Unresolved residual:}

  \begin{itemize}
  \tightlist
  \item
    The CODA k-cell evaluator is not yet bound to this premise checker.
  \item
    Steps 1--4 use imported analytic facts and are a written proof:
    trace-norm duality, multiplicativity of the trace norm, the matrix
    exponential and Taylor's formula. They are not stated in Lean.
  \item
    The existing scalar evaluator does not bind its inputs to a state, a
    Hamiltonian, norm witnesses or the host hypotheses.
  \item
    Bounds on \texttt{π} and on \texttt{arctan} in controls 1 and 2 are written rational
    enclosures.
  \end{itemize}
\end{itemize}

\subsubsection{P.20 Coherence inside guarded components}\label{p.20-coherence-inside-guarded-components}

P.18 left one question open: inside a connected component of the legal-swap
graph \(\Gamma\), do legal backtracks, distant squares and braid hexagons
connect every pair of parallel legal swap proofs? This section gives a
sufficient guard condition with a proof, shows that every native five-edit
family with compatibly ordered precedences satisfies it, and records a bounded
native search for the case where it fails.

\textbf{The complex.} Let \(\mathcal L\) be the legal complete linear extensions of
P.18. Let \(X_{\mathcal L}\) be the 2-complex with vertex set \(\mathcal L\),
the edges of \(\Gamma\), a square for each pair of disjoint adjacent swaps
whose four orders are all legal, and a hexagon for each window of three
positions whose six orders are all legal. A square or hexagon in
\(X_{\mathcal L}\) is exactly a P.17 local relation whose intermediate orders
are all legal. A combinatorial homotopy of edge paths inserts or deletes a
backtrack, or replaces one arc of a face boundary by the complementary arc.
Such a replacement decomposes into backtrack insertions followed by a P.17
square or braid move, with every intermediate order a vertex of the same
face. Two parallel legal swap proofs are therefore joined by legal local
3-witnesses exactly when they are homotopic relative to their endpoints in
\(X_{\mathcal L}\). Coherence inside a component is simple connectivity of
that component.

\textbf{Definition (star-shaped legality).} Fix \(e_0\in\mathcal L\). For an order
\(h\), let \(\operatorname{Inv}(h)\) be the set of event pairs that \(h\)
orders opposite to \(e_0\). A \emph{lower cover} of \(h\) is \(h\cdot s_i\) where
the adjacent pair at positions \(i,i+1\) lies in \(\operatorname{Inv}(h)\).
The legal set is \emph{star-shaped from} \(e_0\) when every lower cover of every
legal order is legal. Descending by lower covers shows that this is
equivalent to the down-set property: every order whose inversion set is
contained in that of a legal order is legal.

\textbf{GEN-WDP-69 (guarded coherence under star-shaped legality).} If
\(\mathcal L\) is star-shaped from some \(e_0\), then \(\Gamma\) is connected
and \(X_{\mathcal L}\) is 1-connected (\(\pi_1=0\)). Any two legal swap proofs with the
same framed source and target histories are joined by a finite sequence of
legal backtracks, distant squares and braid hexagons, and every intermediate
order is legal.

\emph{Proof.} (i) \emph{Faces below a legal order.} Let \(v\in\mathcal L\) have inverted
positions \(i<j\). If \(j>i+1\), the square spanned by \(s_i\) and \(s_j\) at
\(v\) has its other three orders below \(v\): their inversion sets are
contained in \(\operatorname{Inv}(v)\), so they are legal. If \(j=i+1\), the
window holds events \(A,B,C\) with \((A,B)\) and \((B,C)\) inverted. Then
\(e_0\) orders \(C\) before \(B\) before \(A\), so \((A,C)\) is inverted as
well, and the six orders of the window are below \(v\) and legal. In both
cases each swapped pair is ordered differently by \(e_0\) and by \(v\), both of
which satisfy the declared precedences, so the pair is incomparable and the
cell is a legal P.17 relation. (ii) \emph{Connected lower links.} The lower covers
of \(v\) are pairwise joined by the faces of (i), whose remaining vertices
have smaller inversion count. The lower link of \(v\) is therefore a complete
graph. It is nonempty unless \(v=e_0\), because a nonidentity order relative
to \(e_0\) has an inverted adjacent pair. (iii) \emph{Attachment.} An adjacent swap
changes \(|\operatorname{Inv}|\) by exactly one, so every edge joins different
ranks, and every face is a weak-order interval with a unique top vertex. Build
\(X_{\mathcal L}\) from the point \(e_0\) by attaching the legal orders in
increasing rank. Each attachment adds the vertex, its edges to its lower covers
and the faces with that vertex as top. This is a cone on its lower link. A
cone on a nonempty connected graph, glued along that graph, preserves
connectedness, and by van Kampen it preserves simple connectivity: the new
loops through \(v\) are generated by pairs of lower covers, and each pair is
killed by a face. (iv) The first paragraph converts homotopy into local
moves. \(\square\)

The proof is constructive. Let \(D(h)\) be the canonical descent that
repeatedly swaps the leftmost inverted pair until it reaches \(e_0\). For an
inverted position \(i\) of \(x\), induction on rank turns the word
\([i]\,D(x s_i)\) into \(D(x)\). If \(i\) is the leftmost inversion this is
the definition. Otherwise both descents pass through the bottom of the face
of (i), one square or braid move exchanges them, and the inductive hypothesis
handles the two lower paths. Rewriting a word \(w\) from \(h\) to \(h'\) edge
by edge then gives \(w\simeq D(h)D(h')^{-1}\). Applying this to two parallel
words yields an explicit list of local replacements. The native owner emits
that list, and an independent checker replays every intermediate order
against the native guards.

P.17 is the special case \(\mathcal L=\) all linear extensions of the declared
poset. That set is a down-set from each of its elements, so the theorem
recovers P.17's coherence without its all-order premise. P.18's family is not
star-shaped from either legal history: the condition fails exactly where
connectivity fails.

\textbf{Native corollary (compatibly ordered precedence).} Rank the five native
edit kinds as AddTerm 0, AddAssertion and Fill 1, Identify 2, Remove 3. Suppose
the transitive closure of the declared precedences never places a
higher-ranked event before a lower-ranked one. Then any kind-sorted linear
extension \(e_0\) is a valid base, and every five-edit model family of that shape
satisfies GEN-WDP-69. (The bounded search below tests only the
lexicographically least kind-sorted extension.) It suffices to check that an
adjacent pair that a legal history orders against \(e_0\) can be swapped toward
\(e_0\) with the run staying legal and ending in the same framed state. Such a
pair has the lower-ranked event second, or has equal ranks.

\begin{itemize}
\tightlist
\item
  A Remove followed by an AddAssertion on the same count key: adding first
  only raises the count at the intermediate prefix, and both orders end with
  the same count.
\item
  A Remove followed by an Identify: identification pushes counts forward by
  summation, so the pooled entry seen by the later Remove is at least the
  entry it saw before. Both orders end in the same pooled table.
\item
  An Identify followed by a Fill: if the fill is legal after the
  identification it is nondegenerate and new there. Its preimage is then
  nondegenerate and new before, and both orders give the same filled set.
\item
  An Identify followed by an AddAssertion: both are always enabled, and the
  pushed-forward counts agree.
\item
  Events on different current keys, Fills against counts, and AddTerm
  against any event not using the term it creates commute, up to the
  token-to-native alignment of P.14.
\item
  Two events of equal rank (two Removes, AddAssertions, Fills, Identifies or
  AddTerms) are legal in both orders whenever they are legal in one. Two
  identifications that would merge the same pair of classes fail in either
  order, so no legal history contains both.
\end{itemize}

The \href{lean4/WitnessedDifference/GuardedHistoryCoherence.lean}{Lean module}
checks the abstract swap step (\texttt{helper\_swap\_preserves\_run}), its count-key
instance (\texttt{count\_kernel\_swap}), and the fact that the converse swap fails at
count zero. It also checks that a cochain closed on the declared faces is
invariant under backtrack and face moves (\texttt{moves\_preserve\_total}). The ten
declarations in its axiom report use only \texttt{propext} and \texttt{Quot.sound}. The case analysis for
Identify and Fill is a written proof checked on bounded families by the
native search below.

\textbf{When the condition fails (bounded native search).} The native module
\texttt{omega\_wdp::genesis\_guarded\_history} enumerated 132,278 declared families in
five tiers. The tiers use one to three initial terms; count, identification,
fill and term-creation events; up to six events in three tiers and five in
the other two; and up to three extra precedence edges. The menu is fixed in the search binary. Every one of the
38,158 kind-compatible families with a legal history had a star-shaped legal
set, one component and zero rational first homology. The 65,253 constructive
witness lists it generated all replayed. Outside the condition, 126 families
have disconnected legal sets. The smallest has four events: count one, and two
removals, each followed by its own addition. Exactly one family has a
connected component with positive rational first homology.

\textbf{Native counterexample to unconditional coherence.} Take one term with
\(R(0,0)=2\), three removals \(x_1,x_2,x_3\) of that assertion and three
additions \(a_1,a_2,a_3\), with the declared precedences \(x_i<a_i\). Let
coordinate \(i\) count the events of pair \(i\) already performed. Each
schedule is then a monotone lattice path from \((0,0,0)\) to \((2,2,2)\) in
\(\{0,1,2\}^3\), and an adjacent swap of events from different pairs flips one
corner of the path. The count after a prefix is \(2\) minus the number of
pairs at coordinate \(1\). A schedule is therefore legal exactly when its path
avoids the centre \((1,1,1)\). Of the \(90\) paths, \(36\) pass through the
centre, so \(54\) are legal. They form one swap component with \(84\) legal
swaps and \(30\) legal squares. There is no legal hexagon: a hexagon fills a
unit cube of the grid, every unit cube has the centre as a corner, and some of
its six paths pass through it. Exact elimination gives rational first
homology of rank one. The retained certificate supplies two legal swap proofs
from \(x_1a_1x_2a_2x_3a_3\) to \(x_3x_2a_3x_1a_1a_2\) with the same framed
boundaries. It also supplies a 1-cochain supported on four edges, closed on
all \(30\) legal squares, whose values on the two proofs are \(-1\) and \(0\).
By \texttt{moves\_preserve\_total}, no sequence of legal backtracks, squares and braids
joins them. Coherence inside a guarded component therefore fails without the
guard condition, already for the count kernel. Dropping any one of the three
precedences removes the hole. With at most two extra precedences, the search
found no component with positive rational first homology. The counterexample is not star-shaped from any legal
history, in agreement with GEN-WDP-69.

The bounded search supports, and does not prove, the native corollary for
Identify and Fill. Rational first homology zero would not by itself give
simple connectivity; the coherence statements rest on the written proof.

\textbf{Scope.} GEN-WDP-69 is a written theorem with a constructive
reduction. Its Lean counterpart covers only the local swap and invariance
steps. Native evidence covers declared bounded families. The theorem does not
give 4-cells, identify distinct 3-certificates, define restrictions over
arbitrary source covers, or grant authority. GEN-WDP-O03 remains open beyond
the components it covers.

\subsubsection{P.21 Minimum retained information for two native edits}\label{p.21-minimum-retained-information-for-two-native-edits}

\textbf{Definition.} Fix the literal raw states with at most three terms,
nonnegative integer count tables \(R,S\) of combined total at most four,
and every subset of the available oriented nondegenerate filled triples.
Use the 42 literal five-edit descriptors with addresses in \(0..2\), and
at most two successful edits. Protect the joint result mark, the full
literal balance record, native guard outcomes, carrier-bound admission,
native addition-price pairs, and successful occurrences. A native-enabled
edit excluded by the carrier bound retains its native price but is not an
executed edge. A false native guard is a distinct outcome. Neither refusal
consumes a successful-edit slot.

At a terminal prefix, one profile observes only the present mark and balance
record; the second additionally permits nonexecuting status probes. At
positive layers both observe every offer and all admitted continuations.
These two profiles are declared separately, before computing their quotients.

\textbf{Finite proposition (GEN-WDP-73).}
For either profile at horizon two, the coarsest adequate report partition is
raw-state equality: 29,771 classes in total and \(|M^{-1}(m)|\) classes in each
result-mark fiber. An attaining report is
\[
 H(s)=\bigl(M(s),F(s),\min(R,2),\min(S,2)\bigr).
\]
This is a minimum number of information classes, not a minimum number of
fields or bytes.

\emph{Written proof.} Equal protected futures determine whether each count is
positive and whether it remains positive after one removal, hence every
count capped at two. Removals stay inside the bound. Summing the root-normal
table cancels its potentials, so the joint mark determines
\(\sigma=\sum R-\sum S\). Under total bound four, either no cap is two
(every count is known), two caps are two (both counts are exactly two and
all others zero), or exactly one cap is two. In the last case its fixed
relation sign and \(\sigma\) recover its count after subtracting the other
known signed counts. The mark supplies the term count and the literal
balance keys supply the filled set. Thus equal futures force raw equality;
raw equality also preserves all five native operations. The displayed report
decodes by the same argument. Counting the carriers gives
\(1+15+495+4\binom{22}{4}=29771\). ∎

The \href{docs/derivations/native_future_retention.md}{complete derivation}
constructs the layer signatures, forward projections and backward path lifts
from every supplied raw starting state. It gives identity and composition
laws with the remaining horizon retained, whole-field deletion witnesses,
and explicit AddTerm and Identify controls. The reciprocal-count pair
\((R_{01},R_{10})=(1,2)\) versus \((2,1)\) has equal one-step reports only
under the present-only terminal profile: after one removal, a terminal guard
probe separates it under the second profile. There is no
convention-independent least-horizon claim.

Seven native tests, independent replay of all 29,771 states and 1,250,382
offers against the native and pinned original implementations, and 21 selected
formal declarations support different parts of this result. The selected
formal model is not a verification of either concrete program. The
source-specific finite result
does not replace the broader source-refinement and higher-coherence
obligations GEN-WDP-O02 and GEN-WDP-O03. Larger count bounds, transported
action names, unbounded histories, observation instruments and authorization
remain outside this statement.

\subsubsection{P.22 Exact topology and coherence thresholds for paired count guards}\label{p.22-exact-topology-and-coherence-thresholds-for-paired-count-guards}

P.20's count-kernel counterexample extends to an unbounded family with an exact
classification. Fix \(n\ge2\) and \(1\le c<n\). On one valid Source-25 count
key, start with count \(c\) and take distinct event tags \(r_i,a_i\),
\(1\le i\le n\). Each \(r_i\) is a \texttt{Remove}, each \(a_i\) an \texttt{AddAssertion},
and every tag occurs once. Declare the task precedences \(r_i<a_i\).
These tags and precedences are additional presentation data; the native count
update does not supply them. Other counts, terms and filled faces stay fixed.

For any pair-respecting prefix \(u\), its active tasks are those whose removal
has occurred and whose addition has not. Native replay succeeds exactly when
all prefixes have at most \(c\) active tasks, and its count then equals
\(c-|\operatorname{Act}(u)|\). This follows by induction: removal opens one
task and requires positive count; addition closes its already opened task.
Every complete legal history returns to count \(c\).

Let \(K(n,c)\) be the full subcomplex of the event permutohedron consisting
of ordered partitions \(B=(B_1|\cdots|B_k)\) whose every total refinement is
legal. Its vertices are complete tagged histories, and a face has dimension
\(2n-k\). Its two-dimensional faces are exactly P.20's legal squares and braid
hexagons. Higher faces are explicitly included by the same refinement rule.
A finite face test is
\[
 b(r_i)<b(a_i)\quad\text{for all }i,\qquad
 \max_t\#\{i:b(r_i)\le t\le b(a_i)\}\le c,
\]
where \(b(e)\) is the block containing \(e\). To prove the test, put every
removal in a block before every addition there: this attains its largest
possible occupancy. Both interval endpoints are included because an opening
task can overlap a closing task within that refinement.

Define \(T(n,c)\) as the task permutohedron subcomplex consisting of ordered
partitions of \([n]\) whose blocks have size at most \(c\). Also put
\[
 M(n,c+1)=\mathbb R^n\setminus
 \bigcup_{\substack{J\subseteq[n]\\|J|=c+1}}
 \{x:x_i=x_j\text{ for all }i,j\in J\}.
\]
This is the classical real no-\((c+1)\)-equal arrangement complement.

\textbf{Theorem (guarded schedule classification).} For every such \(n,c\),
\[
 |K(n,c)|\simeq |T(n,c)|\simeq M(n,c+1).
\]
For \(c=1\), there are exactly \(n!\) contractible components. For
\(2\le c<n\), the complex is \((c-2)\)-connected and is not
\((c-1)\)-connected; its first nonzero reduced integral homology occurs
in degree \(c-1\). In particular, legal swap proofs have square-and-braid
coherence exactly when \(c\ge3\), within the connected range. Moreover,
\[
 K(n,n-1)\simeq S^{n-2}.
\]

\emph{Proof of the classification.} Index three finite covers by all strict task
posets \(P\) of width at most \(c\). Let \(L(Q)\) denote the full
linear-extension complex of a finite poset \(Q\). It is contractible in
every dimension: the \href{docs/guarded_schedule_topology.md\#3-contractibility-of-a-full-linear-extension-complex}{proved contractibility lemma}
attaches the decreasing-run faces in increasing inversion rank; each face
meets the preceding complex in a contractible polyhedral antistar.
For each \(P\), take
\[
 U_P=L(Q_P),\qquad T_P=L(P),\qquad
 V_P=\{x:x_i<x_j\text{ whenever }i<_Pj\},
\]
where \(Q_P\) is generated by \(r_i<a_i\) and \(a_i<r_j\) for \(i<_Pj\).
Active tasks in \(U_P\) form a \(P\)-antichain, so \(U_P\subseteq K(n,c)\).
Conversely, a legal face defines intervals \([b(r_i),b(a_i)]\) and the
interval order \(i<_Pj\) iff \(b(a_i)<b(r_j)\). Pairwise intersecting
intervals on a line share a point; therefore its width is the maximum
occupancy in the face test, at most \(c\). This shows that the \(U_P\) cover.

Every block of a face in \(L(P)\) is a \(P\)-antichain, so
\(T_P\subseteq T(n,c)\). Conversely, the strict order between the blocks
of any face of \(T(n,c)\) is a poset of width at most \(c\) covering that
face. Similarly, equal-coordinate blocks of \(x\in M(n,c+1)\), ordered by
their values, give a width-at-most-\(c\) poset with \(x\in V_P\).
A point of \(V_P\) cannot have \(c+1\) equal coordinates, since they would
form an antichain. Thus all three are covers of their stated spaces.

A finite intersection in any cover is nonempty exactly when the union of
its task orders is acyclic. In the event cover, a task cycle gives an event
cycle through \(r_i<a_i\), while an event cycle contracts to a task cycle.
When acyclic, the two polyhedral intersections are linear-extension complexes
of the respective transitive closures and hence contractible. The Euclidean
intersection is a nonempty open convex cone. The covers therefore have the
same nerve and contractible nonempty intersections. The finite-subcomplex and
open-good-cover forms of the nerve theorem give both equivalences.

\emph{Proof of the consequences.} At capacity one, \(T(n,1)\) consists of the
\(n!\) task orders. At capacity \(c\ge2\), it contains the full
\((c-1)\)-skeleton of the task permutohedron: a \(d\)-face has no block
larger than \(d+1\). Adjoining the omitted cells, all of dimension at least
\(c\), cannot change \(\pi_j\) for \(j\le c-2\), and produces a
contractible polytope. This establishes \((c-2)\)-connectivity.
Projection from \(M(n,c+1)\) to its first \(c+1\) coordinates has a
continuous section: place the remaining coordinates successively at one,
two, and so on above the maximum retained coordinate. Its retract
\(M(c+1,c+1)\) is \(\mathbb R^{c+1}\) minus its diagonal line and has
homotopy type \(S^{c-1}\). Hence \(\pi_{c-1}\) is nonzero, and
\(H_{c-1}\) has \(\mathbb Z\) as a direct summand. Finally,
\(T(n,n-1)\) is exactly the boundary of the \((n-1)\)-dimensional task
permutohedron, proving the sphere statement. Fundamental groups depend only
on the two-skeleton; its path homotopies are precisely legal backtracks,
squares and braid replacements. This demonstrates the coherence assertion. \(\square\)

For \((n,c)=(3,2)\), this recovers P.20's circle obstruction. At \((4,2)\),
\(T\) has 24 vertices, 36 edges and six square faces. Each square has a free
edge belonging to no other retained face. Collapsing the six squares along
such edges leaves a connected graph with 24 vertices and 30 edges, hence
\[
 H_1(K(4,2);\mathbb Z)\cong\mathbb Z^7.
\]
At \((4,3)\), \(K\)'s cell counts by dimension are
\((1944,5400,5340,2184,302)\), totaling 15,170. The smaller model \(T\) has
\((24,36,14)\), totaling 74; it is a two-sphere. This size reduction preserves
the homotopy type and its obstructions. It is not an execution quotient:
\(T\)'s vertices are task orders, not complete event histories. No preservation
of every native answer, cost, permission or future continuation follows.

The \href{docs/guarded_schedule_topology.md}{full derivation} supplies the written
proof and relates it to established arrangement-complement and semaphore
execution-space theory. The contribution here is the exact binding to the
declared native guard and its legal-cell convention. The
\href{lean4/WitnessedDifference/GuardedScheduleTopology.lean}{Lean module} addresses
prefix replay, the all-refinements capacity criterion, tagged refinements and
selected order/cycle laws. These are separate from the written topological
proof; they do not formalize the nerve theorem or sphere classification.
The \href{scripts/verify_guarded_schedule_topology.py}{finite replay} checks declared
instances, native guard correspondence and cell/homology certificates; finite
checks do not prove the theorem's unbounded quantifiers. Fresh run results and
formal axiom dependencies belong in the verification receipt.

The full complex is essential for higher homology: discarding its higher cells
can introduce additional classes. The theorem covers one count key and the
specified paired events. Additional precedences, multiple keys, repeated
loops, identification and filling require new analysis. It advances O03 for
this family without asserting a classification of every guarded presentation.

\subsubsection{P.23 Sparse integral transport and polynomial bounds at fixed degree}\label{p.23-sparse-integral-transport-and-polynomial-bounds-at-fixed-degree}

The \href{docs/sparse_guarded_chain_transport.md}{complete signed construction and bounds}
make P.22's reduction effective on supplied chains. Fix the same tagged events,
legal-cell convention and integral cellular differential. There are explicit
augmentation-preserving chain maps and a degree-one operator
\[
\begin{gathered}
p:C_*(K(n,c))\longrightarrow C_*(T(n,c)),\\
i:C_*(T(n,c))\longrightarrow C_*(K(n,c)),\\
pi=1,\quad ip-1=\partial H+H\partial.
\end{gathered}
\]
The image of a requested cell is computed from its own carrier and recursive
boundaries. No list of all execution cells, all task cells, or all linear
extensions is needed.

\textbf{Theorem (minimum-sweep sparse transport).} Use the minimum-sweep contraction
specified below, and fix a largest queried degree \(d\). Put
\[
R=\binom{2n}{2},\qquad B=2^{d+2}-2,\qquad C=\max(2,RB).
\]
For \(0\le q\le d\), a degree-\(q\) basis cell has projection and inclusion
images with total absolute coefficient sum at most \(C^q\), and homotopy
image with sum at most \(2R(q+1)C^{2q}\). For an input with \(s\) nonzero
coefficients, these bounds multiplied by \(s\) bound the output supports.
If the coefficients have at most \(b\) magnitude bits and \(M\) is the
applicable unit-column bound, the output uses
\[
O\bigl(sM[n\log(n+1)+b+\log(sM+1)]\bigr)
\]
bits. Evaluation has polynomial bit complexity in \(n,s,b\) for each fixed
\(d\). Checking a supplied certificate is polynomial in \(n\) and that
certificate's own sparse encoding, with explicit bounds in the supplement.
The exponent depends on \(d\); no polynomial bound uniform in degree is asserted.

\emph{Construction and proof.} A native cell \(A\) determines the interval order
\(P_A\), in which a task precedes another when its finish block is strictly
before the other's start block. P.22 establishes its width is at most \(c\).
A task cell \(B\) determines the order \(P_B\) between its blocks. Write
\(Q_P\) for the event order generated by pair precedences and these task
precedences. These carriers shrink under faces.

For a finite order \(Q\), choose its lexicographically first extension.
Move its first, minimal label left across the blocks before it, assigning
each merged cell the incidence sign of the face that puts the label after
that block. The signed sum \(s_x\) satisfies
\(\partial s_x+s_x\partial=1-I_xD_x\), where \(D_x\) deletes a singleton
\(x\) and is zero on a nonsingleton containing it, and \(I_x\) prepends
the singleton. The supplement verifies all cancellation signs. Recursing on
the remaining labels constructs \(h_Q\) with
\(\partial h_Q+h_Q\partial=1-r\epsilon\). It produces at most
\(\binom m2\) distinct unit-coefficient cells on \(m\) labels, using
\(O(m^3)\) label operations. Every crossed label is incomparable with the
minimal one, so the construction remains inside \(L(Q)\).

Map vertices to the corresponding roots and sequential paired words. For
positive degrees, define
\[
p(A)=h_{P_A}p(\partial A),\qquad
i(B)=h_{Q_{P_B}}i(\partial B),\qquad
H(A)=h_{Q_{P_A}}\bigl(ip(A)-A-H(\partial A)\bigr).
\]
The homotopy formula also applies to vertices with \(H\partial=0\).
Nested carriers and the contraction identity prove both chain laws and the
homotopy identity by degree induction. The composite \(pi(B)\) is carried
by the closed face \(B\), whose top chain group has only that generator.
It fixes vertices; its boundary then forces \(pi(B)=B\) in every degree.
A degree-\(q\) cell has at most \(2^{q+1}-2\) facets. Combining this
bound with the quadratic contraction bound gives the stated recurrences
and estimates. All choices are capacity-independent, so \(p,i,H\) commute
strictly with increasing-capacity inclusions. \(\square\)

For a cycle \(z\), the protected question is exactly whether it bounds over
the integers. If \(\partial b=p(z)\), then
\[
\partial\bigl(i(b)-H(z)\bigr)=z.
\]
The reconstructed filling is checked directly against \(z\), using legality
tests and boundaries of the supplied cells. The checker needs no transport
producer or ambient matrix. Closed cochains pull back along \(p\), and
their pairings with included cycles are preserved by \(pi=1\). A pullback
cochain may be represented by its evaluation rule; sparse input does not
guarantee sparse support for its complete coordinate vector. If a finite
cochain support is supplied, checking all legal adjacent-block merges of its
support establishes global closure, rather than checking only visited cells.

The minimum-sweep maps are an explicit second version. The earlier
first-descent contraction has different coordinates and only an exponential
general bound from its inversion-descent proof. Both versions are retained
and named in the native interface. Generic contractions and chain checking
belong to \texttt{omega\_homology}; WDP supplies the event carriers. The written
all-size proof, selected Lean algebraic/recurrence proofs, and bounded Rust
execution evidence remain distinct. These results concern the integral
boundary question; they confer no authorization and do not identify arbitrary
native histories or establish an empirical search advantage.

\subsection{Appendix Q --- Sound transfer of learned rejection rules and its accounted cost}\label{appendix-q-sound-transfer-of-learned-rejection-rules-and-its-accounted-cost}

The \href{docs/antibody_transfer_and_amortization.md}{full derivation and correction
record} connect checked rejection
to reuse without promoting a single failed candidate into an unproved
general rule. The set-theoretic ingredients and cost arithmetic are standard.
The WDP contribution here is their binding to declared tasks, adapters,
candidate identity and independently replayed evidence. Supplemental local
\texttt{ANT-*} results and GEN-WDP-48 leave the 154 canonical rows unchanged.

\subsubsection{Q.1 One rejection versus every rejection}\label{q.1-one-rejection-versus-every-rejection}

Declare candidate carriers \texttt{X,Y}, valid sets \texttt{G\_S\ subset\ X,G\_T\ subset\ Y},
and total adapter \texttt{alpha:Y-\textgreater{}X}. A sound source rejection rule rejects
\texttt{D\ subset\ X\textbackslash{}G\_S}. Then

\[\bigl[\forall D\subseteq X\setminus G_S:\
  \alpha^{-1}(D)\cap G_T=\varnothing\bigr]
  \quad\Longleftrightarrow\quad \alpha(G_T)\subseteq G_S.\]

\textbf{Proof.} Inclusion makes a valid target's image disjoint from every such
\texttt{D}. Conversely, if \texttt{y} is valid but \texttt{alpha(y)} is invalid, the singleton
\texttt{D=\{alpha(y)\}} is a sound source rejection that rejects \texttt{y}. For one fixed
rejection rule only, the exact condition is disjointness of its preimage
from \texttt{G\_T}; this can hold when the full adapter audit fails. Image
inclusion expresses one-way validity preservation, not general predicate
factorization. A soundness certificate need not be a linear dual witness.

\subsubsection{Q.2 Task parameters and partial adapters}\label{q.2-task-parameters-and-partial-adapters}

For task families \texttt{G\_S(d),G\_T(e)} and task-data map \texttt{beta}, the overall
criterion is \texttt{alpha(G\_T(e))\ subset\ G\_S(beta(e))} for every admitted \texttt{e}.
For derivative validation, the concrete condition on valid candidates is
\texttt{partial(alpha(F))=beta(e)} whenever \texttt{partial(F)=e}. Identity maps satisfy
this for the parameterized predicate \texttt{partial(F)!=e}. The source-bound
predicate \texttt{partial(F)!=x} does not become safe for integrand \texttt{5} just
because its code is reused: \texttt{F=5x} is then valid.

For a partial adapter with domain \texttt{A\ subset\ Y}, declare that out-of-domain
inputs pass through unfiltered. The overall condition becomes
\texttt{alpha(G\_T\ intersect\ A)\ subset\ G\_S}. If out-of-domain instead means reject,
one additionally needs \texttt{G\_T\ subset\ A}. A missing-input error is neither
policy by default. Presence or equality of metadata fields does not prove
the semantic inclusion; the actual predicates and maps must be checked.

\subsubsection{Q.3 Conservative accumulation}\label{q.3-conservative-accumulation}

If every transferred rejection \texttt{R\_j} is disjoint from \texttt{G\_T}, then their
union is disjoint from \texttt{G\_T}, so every valid candidate survives. This
follows elementwise and requires either a checked transfer for each rule or
a complete adapter audit plus each source soundness certificate. Source
soundness alone is insufficient. Per-rule checks against the full declared
valid set compose; a general adapter audit is a reusable alternative.
Checking only candidates left by a potentially unsound earlier filter
cannot establish preservation of all valid candidates.
Incremental checking is sound when the invariant that all valid candidates
remain is itself certified and each new step preserves it.

For the linear repair task, an exact annihilator \texttt{y\ D=0} gives a reusable
rejection predicate \texttt{y\ b!=0}, sound for \texttt{b\ in\ D(A)} for any declared repair
set \texttt{A}. Capacity bounds additionally give quantitative weighted-norm lower
bounds; they are unnecessary for qualitative nonmembership. If acceptance
allows \texttt{Phi\textless{}=epsilon}, rejection requires a certified lower bound strictly
greater than \texttt{epsilon}, not just a nonzero pairing. Changing \texttt{D}, the
coefficient domain or the task requires a new justification. Retaining
all valid candidates establishes neither completeness of the rejection
family nor efficient search among the survivors.

\subsubsection{Q.4 Amortization on a declared workload}\label{q.4-amortization-on-a-declared-workload}

For \texttt{N} candidate occurrences, setup cost \texttt{C\_0\textgreater{}=0}, guard cost \texttt{g\textgreater{}=0},
full evaluation cost \texttt{v\textgreater{}0} and \texttt{h} sound early rejections, the declared
constant-cost model gives

\[T_{bare}=Nv,\qquad T_{guard}=C_0+Ng+(N-h)v,\qquad
\operatorname{saving}=hv-Ng-C_0.\]

Strict benefit occurs exactly when \texttt{h\ v\ \textgreater{}\ N\ g\ +\ C\_0}. With \texttt{N\textgreater{}0} this is
\texttt{h/N\ \textgreater{}\ g/v+C\_0/(Nv)}. If \texttt{g\textgreater{}=v}, strict benefit is unrealizable; equality of
costs can occur. If all costs vanish there is no strict benefit. For
\texttt{N=10000,v=1,g=1/1000,C\_0=2}, at least \texttt{13} hits are required. This is a
conditional calculation, not a measured hit rate or performance claim.
A singleton source rejection can have many target preimages even on a pool
of distinct candidates. Repeated occurrences also count repeatedly, so
a singleton guard need not have at most one hit.

Comparing \texttt{C\_alpha+k\ c} with \texttt{k\ c\textquotesingle{}} for reusable versus per-rule audit is
also conditional arithmetic. Division by \texttt{c\textquotesingle{}-c} is valid only when that
difference is positive; no theorem forces one implementation's check cost
to exceed another's. Target hit counts require a named workload or a proved
bound. There is no canonical uniform frequency on an infinite rational
carrier, and source hit rates do not supply target hit rates.

\subsubsection{Q.5 Evidence, reuse and remaining obligations}\label{q.5-evidence-reuse-and-remaining-obligations}

The \href{code/antibody_transfer/source_provenance.json}{finite checker and original
submission} retain separate
source soundness, adapter, per-rule transfer, accumulation and cost scopes.
Replaying a candidate's certificate never grants authorization. A finite
table audit establishes its table instance; a symbolic law is needed to cover
an infinite polynomial or cochain family. Porting a guard to a new
implementation requires preservation of the actual predicate and task
binding, with a new cost assessment.

Restricting the valid target set or enlarging the valid source set preserves
an already proved adapter inclusion. Enlarging the valid source set can
invalidate the source rejection rule, however, so the complete transfer bundle
must retain both obligations. These results justify sound reuse under
explicit premises. Demonstrating that accumulated knowledge makes new
tasks cheaper requires measured or proved target-workload costs in addition
to soundness; general recursive re-entry remains a separate obligation.

\subsubsection{Q.6 A checked finite example across two representation changes}\label{q.6-a-checked-finite-example-across-two-representation-changes}

The \href{docs/derivations/rejection_rule_reuse.md}{full finite derivation} separates
sound answers, coverage of all original solutions and preservation of a cost
bound. Let \texttt{D=\{-1,0,1\}}, \texttt{X0=D²}, and let a valid original answer satisfy
\texttt{x+y=d}, where \texttt{d∈\{-1,0,1,2\}}. Define exact image carriers and maps

\[f_{01}(x,y)=(x+y,x-y),\quad X_1=f_{01}(X_0),\quad
g_{10}(u,v)=((u+v)/2,(u-v)/2),\]
\[f_{12}(u,v)=(u,v+u),\quad X_2=f_{12}(X_1),\quad
g_{21}(r,t)=(r,t-r).\]

The maps are inverse bijections on these carriers. The composed original-answer
decoder is \texttt{H(r,t)=(t/2,r-t/2)}, inverse to \texttt{F=f12∘f01}. The valid sets are
\texttt{G0(d)=\{x+y=d\}} within \texttt{X0}, \texttt{G1(d)=\{u=d\}} within \texttt{X1} and \texttt{G2(d)=\{r=d\}}
within \texttt{X2}; thus \texttt{F(G0(d))=G2(d)}. For \texttt{d=1}, the complete solution paths are
\texttt{(0,1)→(1,-1)→(1,0)} and \texttt{(1,0)→(1,1)→(1,2)}, decoded to the respective
original pairs. Reporting the first coordinate alone does not decode the pair.

For the \emph{actual} emittable set \texttt{E⊆X2}, including menu and fuel restrictions,

\[H(E)\subseteq G_0(d)\ \Longleftrightarrow\ E\subseteq G_2(d),\qquad
G_0(d)\subseteq H(E)\ \Longleftrightarrow\ G_2(d)\subseteq E.\]

The first equivalence follows from pointwise validity transport. For the second,
every valid original answer has a unique preimage under the bijection \texttt{H}; hence
all are emittable exactly when those valid preimages belong to \texttt{E}. Soundness can
hold vacuously when no solution is reachable, so it does not imply coverage.

The parameterized rejection \texttt{R\_d(u,v)≡u≠d} transfers to level 2 through \texttt{g21},
giving \texttt{r≠d}; its level-0 pullback uses \texttt{f01}, giving \texttt{x+y≠d}. Fixed \texttt{R1} is
not a cross-task rule: at \texttt{d=2} it rejects the unique valid level-2 candidate
\texttt{(2,2)}. With a mandatory guard and no fallback the resulting finite arena has
an initial-containing, target-disjoint loss trap. Restoring the original route
restores solvability. Also \texttt{(1,0)∉X1} but \texttt{(1,0)∈X2}: admitting candidates by
field names would accept an out-of-domain level-1 decoder value \texttt{(1/2,1/2)}.

With proposal/check/emission prices \texttt{1,2,1}, the supplied menu permits naming
any candidate. Its optimal resolution cost is exactly \texttt{4}: select a valid pair,
check and emit; every successful run must pay these three charges. Making a safe
guard mandatory adds \texttt{1/2} per occurrence, giving operating cost \texttt{9/2} and cost
\texttt{15/2} including setup \texttt{3}. Coverage survives, but a protected total bound of
\texttt{4} does not. Installation is therefore not necessary to solve this finite
profile within budget \texttt{18}, including task \texttt{d=2}.

A \textbf{different, restricted menu} removes named selection, methods and alternate
routes and draws without replacement in adversarial order. With \texttt{v} valid
candidates and \texttt{s} soundly rejected invalid ones among nine, stopping at the
first valid emission gives

\[T_0=31-3v,\qquad T_R=36-\tfrac72v-2s.\]

These values count all invalid candidates first, an attainable worst case.
Operating nonincrease requires \texttt{4s≥10-v}; including setup on the guarded side
requires \texttt{4s≥16-v}. They are scan-family values, not global optima of the
named-proposal menu. Even maximal rejection at \texttt{d=2} costs \texttt{33/2} to operate
and \texttt{39/2\textgreater{}18} with setup. A check-free guarded extension would cost \texttt{5/2} per
successful candidate, not \texttt{2}, and is not the supplied check-retained menu.

One successful shared-ledger run at \texttt{d=1} pays \texttt{4} for the parent answer, \texttt{3}
for learning/checking/installation, \texttt{1} for transport to level 2, and \texttt{9/2}
for each of its two decoded child answers: total \texttt{17≤18}. Re-entry carries
spent fuel unchanged. A separate outer continuation checker validates all fifteen
steps, receiving completed solve episodes without treating an arena's terminal
emission as a free continuation. Ten hostile traces are refused for skipped
checks/setup/transport, changed bindings, incorrect emissions, reset/undercharge
or an over-budget repeated episode. No wrapper optimality claim is made.
Resetting after two transports and installation would
expose \texttt{5+18=23} against the original \texttt{18}, not create a valid new budget.

Reuse after a change follows the affected mathematical dependencies. A checked
finite verdict-equivalence table can preserve acceptance across a checker
refactoring, but cost additionally depends on the menu, successors, reports,
targets and prices. A changed digest is a custody mismatch, not itself a
mathematical counterexample. Adding an unused expensive action can leave the
optimum unchanged; optimality still needs its lower inequality checked. A checker
that admits sums in \texttt{\{d-1,d\}} instead admits invalid original \texttt{(0,0)} at \texttt{d=1}.

The \href{code/rejection_reuse/replay.py}{Python replay} checks all nine carrier rows,
6,144 emission subsets, 6,144 rejection subsets, 46 independently replayed finite
arena certificates and 22 controls. The \href{code/rejection_reuse/rejection_reuse.lean}{Lean 4 formalization}
checks the image/inverse/validity laws, generic soundness and coverage criteria,
safe union, guard coverage, scan arithmetic, fuel laws and finite witnesses.
It uses standard Lean logical axioms, with no custom axioms or \texttt{sorry}.
Its image-restricted integer decoder and the Python off-image rational control
have separate stated scopes. These checks complete this finite example, not a
general correspondence between recursive engine execution and external tasks.

\subsection{Appendix R --- A plain-language picture of observers and authority}\label{appendix-r-a-plain-language-picture-of-observers-and-authority}

This appendix is an interpretive companion for readers outside the
mathematics. It states no new result; every claim in it points to a section
where the corresponding construction is defined or proved.

\textbf{The rule card.} Three children each say their block tower is best. The
dispute has no answer until someone puts a card on the table: today ``best''
means tallest, measured with this stick from the floor; a loose block may be
fixed once; a tower that falls does not count; equal heights are a tie. The
card is the observer of Chapter 10. It is not a child, an eye or a tower. It
says what is compared, which fixes are allowed and which outcomes exist. A new
card (``best'' means the tower that survives a gentle push) can name a different
winner with the same towers and children. Neither card is wrong; they evaluate
different declared questions (§10.1).

\textbf{The clock.} Two clocks read 12:00 and 12:03. ``Three minutes'' is a defect only
once a time standard, a unit, a tolerance and the permitted repair
(recalibration or replacement) are declared. In a classroom the difference may
be acceptable; in a trading system three milliseconds may not be; in a
distributed database the relevant defect may be whether event order can still
be reconstructed. The same disagreement yields different defects under
different evaluators (§§7.5--7.6, CH10-N01).

\textbf{Five questions to ask of any reported result.} What was compared? What
counted as a defect? Which repairs were allowed? What residual remained? Which
declared rule assigned the final status? A report that cannot answer them is a
proposal, not an evaluation (§§10.1--10.2, §15.6).

\textbf{Gauge, without fog.} A difference that an allowed redescription removes
(units, offsets, naming, formatting) is not an obstruction. A difference that
survives every allowed redescription is one (§4.12). Seeing a difference does
not make it real, and calling everything a convention does not make it go
away.

\textbf{Evaluation is not permission.} A passed check is evidence for the question
the check declares. Whether anyone may act on it is decided by a separately
declared authority regime: who may propose, who may grant, within what scope,
with what appeal (Chapter 15, Appendix P.6, P.19). The useful division of
labour is that a model or person proposes, an evaluator evaluates, and a regime
may authorize. Collapsing any two of these hides the step at which a claim
acquired its standing.

\subsection{Appendix S --- Certified progress in a finite belief arena}\label{appendix-s-certified-progress-in-a-finite-belief-arena}

This is the dynamic counterpart of Appendix Q's static rejection-rule transfer laws.
A restrictive transferred guard can raise cost or destroy winning; a transfer
violation alone is not a losing certificate. The submitted appendix and supplied
adversarial review are retained byte for byte in
\href{docs/source_archive/appendix_s_review/source_manifest.json}{source archive manifest} in \texttt{docs/source\_archive/appendix\_s\_review/}.
The reviewed propositions below keep their proposed WDP-CHS identifiers, outside
the 154-row canonical ledger.

\subsubsection{S.0 Evidence and contribution}\label{s.0-evidence-and-contribution}

D01--D10 are definitions; P01--P19 are finite mathematical propositions;
C01--C04 are conditional interpretations; X01--X08 are retained counterexamples;
N01--N08 are nonclaims; O01--O03 remain open. Written proofs, Lean mathematical
models, and concrete Python replay are distinct evidence.

Finite reachability attractors, alternating bisimulation, partition refinement,
and minimax are standard mathematics. Their task-bound certificate specifications,
exact cost obligations, and separation of evaluation from authorization are the
construction studied here. No parity-game theorem or general Tarski theorem
is needed for these elementary finite proofs.

\href{code/certified_progress/README.md}{Implementation},
\href{lean4/WitnessedDifference/CertifiedProgress.lean}{certificate laws}, and
\href{lean4/WitnessedDifference/BeliefGameQuotient.lean}{quotient laws} have separate
scope. The \href{docs/certified_progress_proof_and_evidence.md}{proof/evidence supplement}
maps each obligation to its actual evidence.

\subsubsection{S.1 Declared inputs and interpretation}\label{s.1-declared-inputs-and-interpretation}

\textbf{WDP-CHS-D01.} Declare a nonempty finite world carrier S, finite actions A,
enablement En(s), nonempty successor sets δ(s,a) for enabled pairs, post-state
reports Λ(s',a), nonempty initial belief b₀, and a finite listed resolution
vocabulary R with total validation tables. Successful targets T are sets of
nonempty beliefs, supplied by the model, independently of certificate keys.
Targets stop execution. Reachable enabled belief/action pairs have positive
rational costs. Task, assumptions, and authorization requirements are inputs.

\textbf{WDP-CHS-D02.} A belief is goal-resolvable when one listed resolution passes
the supplied check in every possible world. Validity does not itself enable an
emission blocked by a guard. Every enabled valid emission must deterministically
reach a declared successful terminal. Therefore

\[V(b)\le\min_{r:\,\text{valid and emit}(r)\text{ enabled at }b}
c(b,\text{emit}(r)).\]

This is an upper bound. Another emission or a preliminary route can be cheaper.
Equality requires accounting for the alternative routes in the specific model.

\textbf{WDP-CHS-D03.} Check successful target membership before reporting a dead end.
A successful terminal with no actions is successful. Otherwise use:

\begin{longtable}[]{@{}
  >{\raggedright\arraybackslash}p{(\linewidth - 2\tabcolsep) * \real{0.5000}}
  >{\raggedright\arraybackslash}p{(\linewidth - 2\tabcolsep) * \real{0.5000}}@{}}
\toprule\noalign{}
\begin{minipage}[b]{\linewidth}\raggedright
Report
\end{minipage} & \begin{minipage}[b]{\linewidth}\raggedright
Required evidence
\end{minipage} \\
\midrule\noalign{}
\endhead
\bottomrule\noalign{}
\endlastfoot
Refused & Controller actually performed a declared refusal action \\
Stalled & Non-target belief has no admissible actions \\
Exhausted & Checked optimal finite value exceeds available rational budget \\
Impossible & Checked initial-containing loss trap and an explicit task policy admitting this report \\
LosingInDeclaredArena & Checked loss without that reporting policy \\
UndeterminedBudget & Only a feasible upper bound exceeds the budget \\
\end{longtable}

An invalid certificate yields RefusedCertificate; it does not mean the controller
executed a refusal. Chapter 20's Proved/Refuted/Refused/Undetermined meanings
remain unchanged. Report labels are not new institutional authorizations.

\textbf{WDP-CHS-N01.} Reporting precedence never removes dead ends from a trap or
turns a successful terminal into a loss.

\paragraph{S.1.1 Assumptions and conditional conclusions}\label{s.1.1-assumptions-and-conditional-conclusions}

\begin{longtable}[]{@{}
  >{\raggedright\arraybackslash}p{(\linewidth - 2\tabcolsep) * \real{0.5000}}
  >{\raggedright\arraybackslash}p{(\linewidth - 2\tabcolsep) * \real{0.5000}}@{}}
\toprule\noalign{}
\begin{minipage}[b]{\linewidth}\raggedright
Assumption
\end{minipage} & \begin{minipage}[b]{\linewidth}\raggedright
Exact scope
\end{minipage} \\
\midrule\noalign{}
\endhead
\bottomrule\noalign{}
\endlastfoot
A1 & Accepted resolutions decode to valid original answers \\
A2 & Finite listed resolution vocabulary and terminating table checks \\
A3 & Each supplied resolution is correctly classified in both directions \\
A4 & Every listed post-state observation is adversarially realizable \\
A5 & Infinite nonterminal play is losing; legal loops can persist indefinitely \\
A6 & All reachable enabled prices are at least some rational ε\textgreater0 \\
\end{longtable}

\textbf{WDP-CHS-C01.} A checked win supplies a genuine original answer if targets
carry sound original-answer decoders and A1 holds. Proof: compose P01's target
reachability with the decoder, as in §20.6's backward transport.

\textbf{WDP-CHS-C02.} A non-target with empty admissibility is losing. Proof: its
singleton is a target-disjoint trap with vacuous closure.

\textbf{WDP-CHS-C03.} Checked optimal value above budget rules out guaranteed success
within that budget. Proof: P16's lower bound applies to every successful plan.
A finite success above budget is not an unrealizable task.

\textbf{WDP-CHS-C04.} An arena loss implies source unsatisfiability only with an
explicit map taking every original solution to an admissible winning strategy,
preserving task and requirements. Proof: contrapositive of that coverage map.
A3 alone checks supplied resolutions and supplies no missing solution or enabled
operation. This is §20.6's constructive forward coverage obligation.

\textbf{WDP-CHS-N02.} Loss is relative to the declared arena, moves, and target.
\textbf{WDP-CHS-N03.} Effective checks do not prove A1 or A3.

\subsubsection{S.2 Beliefs and dynamic updates}\label{s.2-beliefs-and-dynamic-updates}

\textbf{WDP-CHS-D04.} A belief is a nonempty subset of S. Let B be the beliefs
reachable from b₀ before stopping at targets. Finite exploration gives
\textbar B\textbar≤2\^{}\textbar S\textbar−1.

\textbf{WDP-CHS-D05.} Adm(b) consists of actions enabled at every s∈b.

\textbf{WDP-CHS-D06.} For admissible a,

\[\operatorname{Post}_a(b)=\bigcup_{s\in b}\delta(s,a),\qquad
\operatorname{Obs}_a(b)=\{\Lambda(s',a):s'\in\operatorname{Post}_a(b)\},\]
\[b_{a,o}'=\operatorname{Post}_a(b)\cap\Lambda_a^{-1}(o),\qquad
\operatorname{Succ}(b,a)=\{b_{a,o}':o\in\operatorname{Obs}_a(b)\}.\]

Each successor is nonempty because its observation is realized in the post-set.
For b=\{0\}, δ(0,a)=\{1\}, and observation fibre \{1,2\}, the posterior is \{1\};
neither prior intersection nor the unrestricted fibre is correct.

\textbf{WDP-CHS-D07.} A resolution procedure returns a listed r with a bound supplied
check. The finite implementation represents validation tables. Computation,
checking, and emission remain separately declared operations and costs.

\subsubsection{S.3 Sharp finite reachability}\label{s.3-sharp-finite-reachability}

\textbf{WDP-CHS-P01.} Starting with W₀=T∩B, put

\[W_{n+1}=W_n\cup\{b\in B\setminus T:\exists a\in\operatorname{Adm}(b),
\operatorname{Succ}(b,a)\subseteq W_n\}.\]

There are at most \textbar B\textbar{} strict rounds. The stable W is exactly the winning region.
It admits a memoryless strategy on W∖T and a natural rank decreasing on every
selected successor, with rank zero at targets. Its complement L=B∖W is a
target-disjoint trap: every admissible action at an L-belief has an L-successor.

\emph{Proof.} Every strict round adds a new belief. At each admitted nonterminal
choose a witnessing action and its first positive layer as rank. All successors
have smaller layers, so every adversarial play reaches a target within that
rank. At stabilization, an action with every successor in W would have admitted
its source. Hence each enabled action at an L-belief has an L-successor. Choosing
these successors prevents target reachability against every controller, including
a history-dependent one; play either continues in L or ends at a non-target
dead end. This establishes both directions. ∎

\textbf{WDP-CHS-P02.} An avoidably reachable trap does not establish initial loss.
Proof: s has finish→₅t and trap→₁ℓ with ℓ looping. \{ℓ\} is a trap, but s wins.
A sound loss witness must contain b₀ itself. ∎

\textbf{WDP-CHS-P03.} Winning need not be upward closed under belief inclusion.
Proof: p enables only a→t and q only b→t; both singletons win but \{p,q\} stalls. ∎

\textbf{WDP-CHS-P04.} An initial-containing target-disjoint L with an actual
L-successor for every L/enabled pair establishes loss without computing W.
Proof: follow the supplied closure at every controller choice; targets remain
excluded. Dead ends remain in L, where closure is vacuous. ∎

\textbf{WDP-CHS-P05.} Finite-table rejection-rule transfer α(G\_T)⊆G\_S is decidable and a
failure has a target-good/source-bad element. It need not imply modified-arena
loss. Proof: remove a direct cost-one emission but retain a cost-one repair
followed by cost-one emission. Value rises from one to two and remains finite.
Loss still requires P04. ∎

\textbf{WDP-CHS-P06.} P01 quantifies exactly over the listed successors. Changing
adversarial realizability can change winning. Proof: a sole action to a target
wins; adjoining an unlisted looping non-target successor makes it losing.
Suppressing that bad successor reverses the effect. ∎

\textbf{WDP-CHS-P07.} W and L partition B with explicit controller and adversary
strategies on their respective sides. Proof: the constructions of P01 cover B;
rank descent and target-disjoint closure cannot hold at the same initial belief. ∎

\textbf{WDP-CHS-N04.} This reachability argument does not prove parity objectives.

\subsubsection{S.4 Submitted certificates}\label{s.4-submitted-certificates}

\textbf{WDP-CHS-D08.} WIN\_OPTIMAL supplies a complete disjoint W/L partition of
input-derived B, b₀∈W, T∩B⊆W, finite nonnegative rational values and exact natural
ranks exactly on W, strategy exactly on W∖T, and resolution witnesses exactly
where a common supplied resolution exists. Terminal values and ranks are zero.
WIN\_FEASIBLE supplies the same sound partition/strategy but does not claim lower
optimality inequalities.

\textbf{WDP-CHS-D09.} LOSS supplies L⊆B containing b₀ and excluding all model targets.
Closure rows cover exactly all L/enabled pairs and name actual successors in L.
The stalled list is exactly those L-members with no admissible action. Empty win
fields and a false optimality flag are required. A standalone L need not cover B.

Both formats bind schema, model digest, task, assumptions, authorization metadata,
and initial belief. The checker reconstructs targets and B from original inputs,
checks every stalled claim, every map domain, exact scalar/rank types, and each
supplied resolution on every possible world. It never replaces L by L∖stall.
It neither searches for a strategy nor substitutes a solver's answer.

\textbf{WDP-CHS-N05.} A checked packet records evidence, not operational authority.

\subsubsection{S.5 Belief-arena quotient and constructive refinement}\label{s.5-belief-arena-quotient-and-constructive-refinement}

\textbf{WDP-CHS-P08.} Let Π partition B. Within every block require equal target
truth, enabled action sets, action prices, and successor-block sets. The map
q:B→B/Π preserves winning and every achievable finite rational budget, hence
optimal values including losing infinity. These are sufficient conditions,
not necessary conditions for accidental agreement of one answer.

\emph{Proof.} Define quotient tables by a representative; the laws make them
well-defined. Every fine successor maps to a quotient successor, and every
quotient successor at q(b) has an actual fine successor at b. Induction on
horizon establishes target reachability equivalence: at horizon zero target tests
agree; at the next horizon actions and full successor tests agree by the
two directions and induction. Subtracting equal action prices gives the same
induction for rational remaining budgets. Quantify over horizons. ∎

This is a partition of beliefs. Targets and belief costs cannot be tested on
untyped concrete-world representatives.

\textbf{WDP-CHS-P09.} A new target separating a merged block cannot factor through q.
Proof: equal q-images would require equal target truth. ∎

\textbf{WDP-CHS-P10.} Accidental cost equality does not imply congruence. Proof:
p and q both have cost-one d→t; p additionally has a→₂t and q has b→₂t.
Both values are one but action sets differ. Retarget p as successful: V(p)=0,
whereas q's only successors t are non-target dead ends and V(q)=∞. ∎

\textbf{WDP-CHS-P11.} Information compression can destroy singleton wins, as P03
shows. That merge is not a validated P08 game quotient. ∎

\textbf{WDP-CHS-D10.} Starting at Π₀, split each block by the finite signature

\[\left(1_T(b),\
\{(a,c(b,a),\{q_k(y):y\in\operatorname{Succ}(b,a)\}):
a\in\operatorname{Adm}(b)\}\right).\]

Repeat until unchanged. Blocks are only split; new measurements are not invented.

\textbf{WDP-CHS-P12.} The procedure terminates in at most \textbar B\textbar−\textbar Π₀\textbar{} strict rounds
and yields the coarsest P08-congruence refining Π₀.

\emph{Proof.} Every strict round increases the number of nonempty blocks by at least
one, bounded by \textbar B\textbar. Stable signatures give all four congruence laws.
Let Ξ be any congruence refining Π₀. If it refines Π\_k, its common target/action/
price data agree inside every Ξ-block; equal successor Ξ-block sets imply equal
successor Π\_k-block sets. Thus the next signature split never separates a Ξ-block.
Induct over rounds. Ξ refines the stable output, proving coarseness and uniqueness. ∎

\textbf{WDP-CHS-P13.} A refinement verification record binds original model, input partition,
stable partition, and strict-round count. Reconstructing signatures checks the
result; a changed model needs fresh replay. Proof: all four signature components
are finite original tables. A verification record does not implement an external sensor. ∎

\textbf{WDP-CHS-N06.} Deploying q requires an actual observer realizing its
distinctions and preserving future enabled operations and prices. Section
14.7 constructs and optimizes such observation procedures for declared
finite transition/report tables; physical realization and a new engine's
source correspondence still require their own evidence.

\textbf{WDP-CHS-P14.} Quotient strategies lift by σ(b)=σ̄(q(b)); class-constant fine
strategies descend with the same worst-case budgets. An arbitrary fine strategy
need not literally descend. Optimal quotient strategies exist by P15 and lift
as optimal class-constant fine ones. Proof: P08 matches paths in both directions
with equal cost and target truth; class constancy makes σ̄ well-defined. ∎

\subsubsection{S.6 Costs, rank, and independent optimality}\label{s.6-costs-rank-and-independent-optimality}

\textbf{WDP-CHS-P15.} On a nonterminal winning arena, let ε be the minimum enabled
reachable price, c\_max its maximum, R the maximum attractor rank, M=c\_max R,
and H=⌊M/ε⌋. Use synchronous finite-horizon recursion, with V₀=0 on targets
and infinity elsewhere:

\[V_{n+1}(b)=\begin{cases}
0,&b\in T,\\
\min_{a\in\operatorname{Adm}(b)}
\left(c(b,a)+\max_{y\in\operatorname{Succ}(b,a)}V_n(y)\right),&b\notin T.
\end{cases}\]

The empty minimum is infinity. V\_H is the exact optimum on W; outside W it is
infinity. If only a successful initial target is reachable, H=0 and value zero.

\emph{Proof.} An attractor strategy has cost at most M at every winner. Every strategy
of guaranteed cost at most M uses at most H actions on every branch, because
n actions cost at least nε. There are finitely many contingent plans of depth H;
backward induction computes their minimum. Any better strategy would also have
cost at most M, hence depth at most H. This establishes attainment and optimality.
Decomposing an optimal plan at its first action establishes the Bellman equation;
conversely optimal successor plans can be composed after observing that action.
Every minimizing action has V(y)≤V(b)−ε for all successors. Its graph is acyclic.
In increasing value order set ρ(t)=0 and ρ(b)=1+max\_y ρ(y). This is the strategy's
rank, not necessarily its earlier attractor rank. ∎

\textbf{WDP-CHS-P16.} S.4's checked partition, rank, and terminal boundaries, together
with the selected-action equality and all-action lower inequalities,

\[V(b)=c(b,\sigma(b))+\max_{y\in\operatorname{Succ}(b,\sigma(b))}V(y),\]
\[V(b)\le c(b,a)+\max_{y\in\operatorname{Succ}(b,a)}V(y)
\quad(a\in\operatorname{Adm}(b)),\]

are sufficient for exact optimality. Infinity in the second line is permitted
only for a successor in the checked losing trap. P15 supplies a certificate in
this normal form for every winning instance.

\emph{Proof.} Induct on submitted rank to obtain a terminating selected strategy
with cost at most V(b); terminal cost is zero. For a successful contingent plan,
induct on finite depth and choose a successor attaining the lower line's maximum.
An L-successor contradicts guaranteed success by trap closure; a W-successor
has subtree cost at least V(y) by induction. Thus worst-case total is at least
c(b,a)+V(y)≥V(b). Finite branching implies any strategy all of whose maximal
plays reach a target has bounded depth: an unbounded tree would have an infinite
nonterminal branch. Hence the lower argument also covers history-dependent
guaranteed strategies. Upper and lower bounds agree. ∎

Terminal zero fixes a constant otherwise unconstrained by local Bellman
equalities. These laws define this certificate normal form; differently
formatted proofs of one initial optimum need not contain all the same tables.

\textbf{WDP-CHS-P17.} A certified rank-decreasing strategy costs at most c\_max ρ(b).
Proof: at most ρ(b) paid actions occur and each costs at most c\_max. ∎

\textbf{WDP-CHS-P18.} D03's reporting procedure returns one label under its evidence
and precedence. Proof: successful targets are tested first, then non-target
stalls; W and L are disjoint, and budget comparisons use winning finite values.
Feasible upper cost above budget supplies no negative budget result. ∎

\textbf{WDP-CHS-N07.} Losing infinity, finite budget exhaustion, and bounded
computation left Undetermined are different claims.

\subsubsection{S.7 Optional actions and restrictive guards}\label{s.7-optional-actions-and-restrictive-guards}

\textbf{WDP-CHS-P19.} (1) Adding optional actions while preserving all old targets,
enablements, prices, and successor sets cannot remove a win or raise optimum.
Proof: every old strategy remains available unchanged. Removing adversarial
branches of an old action also cannot worsen its guaranteed budget.
(2) A restrictive guard can raise finite cost or destroy winning; actual loss
requires P04. P05 gives the finite-cost alternative; replacing a sole emission
with a non-target loop gives a singleton trap.
(3) Setup payback requires finite operating savings greater than
C\_learn+C\_check+C\_install. Positive savings alone do not suffice. For a declared
deployment distribution or count, sum finite savings and compare the specified
total setup cost. ∎

\textbf{WDP-CHS-N08.} Deployment probabilities, empirical hit rates, and setup costs
are external declarations, not consequences of the game theorem.

\subsubsection{S.8 Retained examples and changed-task controls}\label{s.8-retained-examples-and-changed-task-controls}

X01/X02: P03 refutes the singleton-world hope argument; D06 refutes prior/fibre
posterior substitutes. X03/X04: P10 separates accidental equality from structural
congruence, and P11 demonstrates a destructive information merge.

X05: For s→₁u→₁t and s→₅t, optimum is two. Budget 3/2 yields Exhausted, not
Impossible. A cost-five feasible upper bound does not refute budget three.
The selected cheap strategy has rank two, though the direct action admits s
in attractor layer one.

X06: On the filled triangle with edges (01,12,02),
d₀f=(f₁−f₀,f₂−f₁,f₂−f₀). For δ=(0,1,0),(1,1,0),(0,2,0), the unit L¹
quotient obstructions are 1,2,2. Primal f=0 attains these raw norms. Dual
y=(1,1,−1) annihilates d₀, has infinity norm one, and pairs to the same values.
Thus all three are nonexact. Blind checks reporting nonexact classify exactness
but do not identify the vector. An optional self-loop preserves the available
classification emission win; the identification task with only blind loops
is losing. R7 must name which task it tests.

X07: False stall u cannot mask X05's cheaper route. Terminal value 100 cannot
shift the cost-one optimum to 101. Empty winning witnesses omitting b₀,
certificate-defined target omissions, and successful-terminal loss claims
are refused. Retaining all of L repairs both genuine dead-end loss and optimal
winning games with alternative losing detours.

X08: F=5x is invalid for the antiderivative task f=x and valid for f=5. A source
rejection rule rejecting it cannot be transferred unchanged to that target task.
If the only valid emission is replaced by a non-target loop, its singleton is
a trap. Another solving route prevents that inference. In a sole-emission
cost-1/2 example the optimum is 1/2; this does not imply D02's old general equality.

\subsubsection{S.9 Executable replay}\label{s.9-executable-replay}

The authoritative implementation is the separate producer, model validator,
and checker under code/certified\_progress. No embedded obsolete sketch is an
alternative checker. The producer uses exact rational values and an internal
missing-value marker for infinity; no float infinity is a serialized witness.
Strict disk JSON admission refuses duplicate keys and nonfinite literals.

Replay retains original model, submitted certificate, verification record, mutation controls,
and source hashes, with no model calls. The tests include all seven supplied
review failures, avoidable traps, incorrect resolutions, coefficient/type
refusals, changed bindings, target and partition controls, both triangle tasks,
and independent policy enumeration over 4,096 small three-state/two-action games.
These are concrete finite checks, separately from the all-finite written proof.

\subsubsection{S.10 Acceptance and open boundaries}\label{s.10-acceptance-and-open-boundaries}

\begin{longtable}[]{@{}
  >{\raggedright\arraybackslash}p{(\linewidth - 4\tabcolsep) * \real{0.3333}}
  >{\raggedright\arraybackslash}p{(\linewidth - 4\tabcolsep) * \real{0.3333}}
  >{\raggedright\arraybackslash}p{(\linewidth - 4\tabcolsep) * \real{0.3333}}@{}}
\toprule\noalign{}
\begin{minipage}[b]{\linewidth}\raggedright
Obligation
\end{minipage} & \begin{minipage}[b]{\linewidth}\raggedright
Mathematical evidence
\end{minipage} & \begin{minipage}[b]{\linewidth}\raggedright
Concrete evidence
\end{minipage} \\
\midrule\noalign{}
\endhead
\bottomrule\noalign{}
\endlastfoot
Initial membership and complete win partition & P01/P04/P16; Lean rank/trap laws & Forged initial/partition controls \\
Input targets and terminal zero & P04/P16 & Target-as-loss and shifted-value refusals \\
Dead ends remain in L & P01/P04 & Genuine dead-end and losing-detour replay \\
Exact optimum and selected rank & P15/P16; Lean upper/lower budget laws & Independent policies and cost-two witness \\
Dynamic posterior & D06 & Changed-world and ambiguous-report controls \\
Congruence suffices & P08/P14; Lean two-sided transport & Quotient and retargeting tests \\
Constructive coarsest refinement & P12; Lean finite relation construction & Signature splitting and stability/bound checks \\
Guard failure is not automatically loss & P05/P19 & Finite-cost and singleton-trap examples \\
Budget/report separation & C03/P18 & Feasible versus optimal comparisons \\
\end{longtable}

\textbf{WDP-CHS-O01.} Full parity/alternating objectives need their own certificate
laws and proofs.

\textbf{WDP-CHS-O02.} Original-world validation soundness, completeness, and
constructive winning-strategy coverage remain profile-specific adapter duties.
Declaring A1--A3 does not discharge them.

\textbf{WDP-CHS-O03.} Native Rust binding and verification of Python source remain
open. The written theorem, Lean model, and concrete replay have separate scopes.
Physical observer realization, general guarded-history coherence, and
GEN-WDP-O01--O05 retain their existing open meanings.

\subsection{Theorem map --- local acceptance, projection obstruction, checked repair, and WDP}\label{theorem-map-local-acceptance-projection-obstruction-checked-repair-and-wdp}

Four names meet on one finite question, and they do not do the same work.

\begin{itemize}
\tightlist
\item
  Local acceptance supplies local satisfaction and the actual compatibility defect. A vector is locally invisible when it lies in \(K=\ker Q\). It is an actual source defect when it also lies in the realizable image \(R=\operatorname{im} D_s\).
\item
  Projection Obstruction identifies what a report hides. For an observer \(A\), the hidden set is \(\ker A\). A hidden realizable defect is a nonzero vector of \(R\cap\ker A\).
\item
  Checked repair supplies checked transport and declared repair equivalence. Two vectors represent the same repair class when their difference lies in \(B=\operatorname{im} D_r\). A certificate is accepted only when an independent checker replays the submitted witness. That acceptance is not an authorization to perform a repair.
\item
  WDP retains the original question and the realizability obligations: which added observations distinguish realizable defects from one another and distinguish repair classes, and which of those observations are obtainable from the declared operations. An algebraic identity \(B=V\) makes the quotient trivial. It does not exhibit an admissible source operation.
\end{itemize}

The equal-image count in §10.6 is the special case \(R=B\). Theorem WDP-DIM-1 does not change that statement. Imported linear algebra is named in Appendix T before it is used. The WDP construction is the separation of \(R\) from \(B\), the joint row count, the restricted menus, and the separate operation witness.

Targets 2--5 of \texttt{docs/five\_outstanding\_derivations.md}, general guarded-history coherence, unrestricted source refinement, physical reconstruction, and joint continuum foundations stay open.

\subsection{Appendix T --- Sharp completion with distinct images (WDP-DIM-1)}\label{appendix-t-sharp-completion-with-distinct-images-wdp-dim-1}

\textbf{Coefficient field.} \(\mathbb Q\). All spaces below are finite-dimensional rational vector spaces, given by explicit matrices. No completeness, topology, or estimate is used.

\textbf{Data.} Linear maps \(D_s:U\to V\), \(D_r:F\to V\), and \(Q:V\to Y\). Write \(R=\operatorname{im} D_s\) for the realizable defect space and \(B=\operatorname{im} D_r\) for the algebraic repair space. Do not assume \(R=B\). Let \(K=\ker Q\), and let \(E:V\to\mathbb Q^m\) be an additional row map. Set \(A=(Q,E):V\to Y\oplus\mathbb Q^m\).

\textbf{Imported linear algebra, used as algebra and not as a WDP theorem.} Rank-nullity. For subspaces \(X,Y\), \(\dim(X+Y)+\dim(X\cap Y)=\dim X+\dim Y\). A linear map is injective exactly when its kernel is \(\{0\}\). If \(\pi:Y\oplus\mathbb Q^m\to (Y\oplus\mathbb Q^m)/A(B)\) is the quotient projection, then \(\ker(\pi\circ A)=A^{-1}(A(B))=\ker A+B\), so the induced map \(V/B\to (Y\oplus\mathbb Q^m)/A(B)\) has kernel \((\ker A+B)/B\cong\ker A/(\ker A\cap B)\), and that induced map is injective exactly when \(\ker A\subseteq B\).

\textbf{WDP-DIM-1.} Let \(T=R\cap K\cap B\). The three sharp row counts are
\[
\begin{aligned}
m_{\mathrm{raw}}&=\dim(R\cap K),\\
m_{\mathrm{quot}}&=\dim K-\dim(K\cap B),\\
m_{\mathrm{both}}&=\dim K-\dim(K\cap B)+\dim T.
\end{aligned}
\]
The joint minimum row count is \(m_{\mathrm{both}}=\dim K-\dim(K\cap B)+\dim T\). The sum of the separate minima exceeds the joint minimum by \(\dim((R\cap K)/T)\). The two counts agree exactly when \(R\cap K\subseteq B\), which includes the equal-image case \(R=B\).

\subsubsection{T.1 Raw sufficiency}\label{t.1-raw-sufficiency}

Raw sufficiency means that \(A\) distinguishes realizable defects: \(R\cap\ker A=\{0\}\). Since \(\ker A=K\cap\ker E\), this is exactly \((R\cap K)\cap\ker E=\{0\}\). Thus \(E\) must be injective on \(R\cap K\), and rank-nullity gives the lower bound \(m\ge\dim(R\cap K)\).

The bound is attained. Choose a basis of \(R\cap K\), extend it to a basis of \(V\), and take the dual coordinate rows of that first block. Those \(m_{\mathrm{raw}}\) rows are injective on \(R\cap K\). A lower certificate is any nonzero vector of \(R\cap K\cap\ker E\), or the dimension count when \(m<m_{\mathrm{raw}}\). An upper certificate is the constructed row matrix, replayed by checking that its rank on a basis of \(R\cap K\) equals \(\dim(R\cap K)\).

\subsubsection{T.2 Quotient sufficiency}\label{t.2-quotient-sufficiency}

The induced map \(V/B\to(Y\oplus\mathbb Q^m)/A(B)\) is injective if and only if \(\ker A\subseteq B\). For a block observer \(A=(Q,E)\) one has \(\ker A=K\cap\ker E\), so the criterion is \(K\cap\ker E\subseteq B\). The largest subspace of \(K\) that may remain in the kernel is \(K\cap B\), of codimension \(m_{\mathrm{quot}}\) in \(K\). This count does not depend on \(R\).

Attaining rows: extend a basis of \(K\cap B\) to a basis of \(K\) and then to a basis of \(V\), and take the dual coordinates of the complementary block inside \(K\). Those rows vanish on \(K\cap B\) and are injective on a complement, so \(K\cap\ker E=K\cap B\subseteq B\).

The kernel of the induced map is \((\ker A+B)/B\). It is zero exactly under the criterion above. A lower certificate is a vector of \(\ker A\) outside \(B\). An upper certificate is the constructed row matrix, replayed by checking \(K\cap\ker E\subseteq B\).

\subsubsection{T.3 The joint optimum is not the sum}\label{t.3-the-joint-optimum-is-not-the-sum}

Both properties hold together exactly when
\[
K\cap\ker E\subseteq B\qquad\text{and}\qquad (K\cap\ker E)\cap R=\{0\}.
\]
The second condition then says \(K\cap\ker E\subseteq K\cap B\) meets \(T\) only at \(0\). Inside \(K\cap B\), the largest subspace meeting \(T\) only at \(0\) has dimension \(\dim(K\cap B)-\dim T\). Hence
\[
\dim(K\cap\ker E)\le\dim(K\cap B)-\dim T
\]
whenever both properties hold. Rank-nullity on \(K\) gives \(\dim(K\cap\ker E)\ge\dim K-m\), so
\[
m\ge\dim K-\dim(K\cap B)+\dim T=m_{\mathrm{both}}.
\]
If \(m<m_{\mathrm{both}}\) and \(K\cap\ker E\not\subseteq B\), the quotient property fails. If \(m<m_{\mathrm{both}}\) and \(K\cap\ker E\subseteq B\), the dimension count forces a nonzero intersection with \(T\), so the raw property fails. This case split is the proof that no smaller joint row count exists. It does not add \(m_{\mathrm{raw}}\) to \(m_{\mathrm{quot}}\).

The excess of that sum is
\[
m_{\mathrm{raw}}+m_{\mathrm{quot}}-m_{\mathrm{both}}=\dim(R\cap K)-\dim T=\dim((R\cap K)/T).
\]
It is zero exactly when \(R\cap K\subseteq B\).

The count is attained. Split \(K\cap B=T\oplus S_0\) and \(K=(K\cap B)\oplus C\), and set \(W=T\oplus C\). Then \(\dim W=m_{\mathrm{both}}\) and \(S_0\cap R=\{0\}\). Extend a basis that begins with \(S_0\) and continues with \(W\) to a basis of \(V\). The dual rows of the \(W\) block satisfy \(K\cap\ker E=S_0\subseteq B\) and \(S_0\cap R=\{0\}\).

When \(R=B\), one has \(T=R\cap K=B\cap K\), so \(m_{\mathrm{raw}}=\dim(B\cap K)\), \(m_{\mathrm{quot}}=\dim K-\dim(B\cap K)\), and \(m_{\mathrm{both}}=\dim K\). The sum equals the joint count. That is the §10.6 specialization, recovered here and not redefined.

A submitted witness carries a separate field, \texttt{lower\_columns}. The checker accepts that field only when its columns are an obstruction basis for the claimed target. Width and rank equal to the sharp count are required, and they are not sufficient.

\begin{itemize}
\tightlist
\item
  Raw: the columns are a basis of \(R\cap K\).
\item
  Quotient: the columns lie in \(K\), meet \(B\) only at \(0\), and number \(m_{\mathrm{quot}}\).
\item
  Joint: the columns lie in \(K\), number \(m_{\mathrm{both}}\), and meet \(K\cap B\) only inside \(T\).
\end{itemize}

The empty matrix is that basis when the sharp count is \(0\). On the \(B\subset R\) example of T.4, the produced joint basis is a basis of \(K=\operatorname{span}\{e_1,e_2\}\). The two columns e\_2 and e\_3 have width \(2\) and rank \(2\), and \(e_3\notin\ker Q\), so they are not a joint obstruction basis. The checker rejects them. This acceptance rule does not replace the kernel-vector lower bounds proved in T.1 and T.2. Those vectors show that a too-small row map fails. The field \texttt{lower\_columns} shows that the sharp count is forced, independently of which attaining rows were submitted.

\subsubsection{T.4 Three inclusions}\label{t.4-three-inclusions}

All three examples use the last-coordinate observer, so \(K\) is the span of the earlier coordinates. The shipped checker recomputes the dimensions from the matrices.

\textbf{\(R\subset B\).} \(V=\mathbb Q^3\), \(D_s\) has the single column \((1,0,0)\), \(D_r\) has columns \((1,0,0)\) and \((0,1,0)\), and \(Q=(0,0,1)\). Then \(B=K\), \(R=\operatorname{span}\{e_1\}\), \(T=R\), \(m_{\mathrm{raw}}=1\), \(m_{\mathrm{quot}}=0\), \(m_{\mathrm{both}}=1\), and the separate sum is \(1\).

\textbf{\(B\subset R\).} The same \(V\) and \(Q\), with \(D_s\) now both of \((1,0,0)\) and \((0,1,0)\), and \(D_r\) the single column \((1,0,0)\). Then \(R=K\), \(B=\operatorname{span}\{e_1\}\), \(T=B\), \(m_{\mathrm{raw}}=2\), \(m_{\mathrm{quot}}=1\), \(m_{\mathrm{both}}=2\), and the separate sum is \(3\). The joint optimum is strictly smaller.

\textbf{Neither inclusion.} \(V=\mathbb Q^4\), \(Q=(0,0,0,1)\), \(D_s\) has columns \((0,1,0,0)\) and \((0,0,1,0)\), and \(D_r\) has columns \((1,0,0,0)\) and \((0,1,0,0)\). Then \(K=\operatorname{span}\{e_1,e_2,e_3\}\), \(B=\operatorname{span}\{e_1,e_2\}\), \(R=\operatorname{span}\{e_2,e_3\}\), and \(T=\operatorname{span}\{e_2\}\). Neither image contains the other. Here \(m_{\mathrm{raw}}=2\), \(m_{\mathrm{quot}}=1\), \(m_{\mathrm{both}}=2\), and the separate sum is \(3\).

A raw completion need not be a quotient completion. On \(V=\mathbb Q^3\) with \(R=\operatorname{span}\{e_1\}\), \(B=\operatorname{span}\{e_3\}\), and \(Q=(0,0,1)\), the single raw row \((1,0,0)\) kills every realizable invisible defect, while \(e_2\) remains in \(\ker A\) and outside \(B\). The converse gap is the quotient row on the \(B\subset R\) example: \((0,1,0)\) forces \(K\cap\ker E=\operatorname{span}\{e_1\}\subseteq B\), and \(e_1\) is a nonzero realizable defect still in the kernel.

\subsubsection{T.5 Restricted menus}\label{t.5-restricted-menus}

Unrestricted algebraic rows are arbitrary elements of the dual of \(V\). A restricted menu is a smaller declared set.

\textbf{Rows required to annihilate \(B\).} Then \(B\subseteq\ker E\), so \(T\subseteq R\cap\ker A\). If \(T\neq\{0\}\), every such row system fails raw sufficiency and fails the joint property. The obstruction certificate is any nonzero vector of \(T\). The \(R\subset B\) example is this impossibility: \(T=\operatorname{span}\{e_1\}\). If \(T=\{0\}\), put a basis of \(B\) first, then a complement of \(K\cap B\) inside \(K\), then extend to \(V\). The dual rows of that complement annihilate \(B\), and they attain \(m_{\mathrm{both}}\), which equals \(m_{\mathrm{quot}}\) in this case. Quotient sufficiency alone remains feasible at \(m_{\mathrm{quot}}\) whether or not \(T\) vanishes: the same annihilator rows may kill \(T\), and the quotient property allows that.

\textbf{A finite instrument shelf.} The rows must be chosen from a supplied finite list, not from their linear span. Positive rational costs are part of the shelf. Adding a row preserves sufficiency, so a minimum-cost sufficient subset exists whenever any subset is sufficient, and it is found by exhaustive enumeration of the finitely many subsets in increasing cost. The producer returns one such subset. A separate checker enumerates again from the submitted index set and cost. It rejects an insufficient subset, a cost that does not match the submitted indices, and any sufficient subset of strictly smaller cost. If no subset works, the checker accepts impossibility only after the full measurement family itself fails the same sufficiency test.

\subsubsection{T.6 A repair claim needs an operation}\label{t.6-a-repair-claim-needs-an-operation}

Choosing \(B=V\) makes \(K\cap B=K\), so \(m_{\mathrm{quot}}=0\) and the induced quotient map is automatically injective. That fact does not solve the original problem. A claim that a repair solves the original defect is accepted only with a separate witness \((D_{\mathrm{adm}},u,r_{\star},\delta)\) satisfying \(D_{\mathrm{adm}}u=r_{\star}\), \(r_{\star}\in B\), and \(r_{\star}=\delta\), where \(\delta\) is the declared realizable defect. The checker replays those matrix identities. It does not invent an operation. Certificate acceptance remains separate from an authorization to use the operation.

\subsubsection{T.7 What this theorem does not say}\label{t.7-what-this-theorem-does-not-say}

WDP-DIM-1 is a finite rational row-count theorem. It does not assert a continuum reconstruction, a quantum code, a physical instrument, or a repair that has not been given as an admissible map. The Chapter 23 corrections are a different carrier and are proved in Appendix U.

\subsubsection{T.8 Repair-class separation restricted to realizable candidates}\label{t.8-repair-class-separation-restricted-to-realizable-candidates}

\textbf{WDP-DIM-R (realizable-domain corollary).} Retain the rational linear data
of WDP-DIM-1, but restrict candidates to \(r\in R\). The protected answer is
\(q(r)=r+B\), and the observed report is the actual value \(A(r)=(Q(r),E(r))\).
A decoder on \(A(R)\) exists exactly when

\[
R\cap\ker A\subseteq B.
\]

Writing \(L=R\cap K\) and \(T=L\cap B\), the minimum number of additional rational
linear observation rows is

\[
m_{\mathrm{quot},R}=\dim L-\dim T
=\dim(R\cap K)-\dim(R\cap K\cap B).
\]

\emph{Proof.} Equal reports for \(r,r'\in R\) mean \(r-r'\in R\cap\ker A\);
equal protected answers mean \(r-r'\in B\). Conversely every vector in
\(R\cap\ker A\) occurs as the difference of itself and zero, so the fiber
criterion is exactly the displayed inclusion. Equivalently,
\(\ker(E|_L)\subseteq T\). Rank-nullity therefore gives
\(m\ge\operatorname{rank}(E|_L)\ge\dim L-\dim T\).
For attainment, extend a rational basis of \(T\) to one of \(L\), and then to one
of \(V\). Take the dual coordinate rows of the complementary block inside \(L\).
Their kernel on \(L\) is exactly \(T\), and there are \(\dim L-\dim T\) rows.
This includes the zero-row case. \(\square\)

For example, let \(V=\mathbb Q^2\), \(R=\operatorname{span}\{e_1\}\), \(B=0\)
and \(Q=0\). The row \(E=(1,0)\) suffices on \(R\), whereas the global repair-class
question needs two rows: \(e_2\) and zero still collide under this \(E\).
If both raw and repair-class separation are required only on \(R\), raw
injectivity already supplies both, so the minimum is \(\dim(R\cap K)\).
The joint count in T.3 instead includes the global quotient obligation.

The restricted quotient presentation is
\(R/(R\cap B)\to A(R)/A(R\cap B)\). It is injective under exactly the same
criterion: if \(A(r)=A(b)\) for \(b\in R\cap B\), then
\(r-b\in R\cap\ker A\), and the converse follows by taking \(b=0\).
Do not replace its report quotient by the ambient quotient by \(A(B)\).
For \(R=\operatorname{span}\{e_1\}\), \(B=\operatorname{span}\{e_2\}\) and
\(A=(1,1)\), actual reports distinguish every candidate in \(R\), but
\(A(B)=\mathbb Q\) erases the entire report. Repairs outside \(R\) change that
observation contract.

\subsubsection{T.9 Finite encodings and instrument-menu boundaries}\label{t.9-finite-encodings-and-instrument-menu-boundaries}

The \href{docs/observation_scope_and_literature.md}{observation-scope supplement}
gives two direct finite consequences of the report-fiber criterion. An
unrestricted auxiliary binary encoding needs
\(\lceil\log_2\max_{a\in p(X)}|q(p^{-1}(a))|\rceil\) bits for a nonempty finite
candidate set; this standard counting result does not price an instrument
shelf. A four-state, five-query shelf has pairwise non-derivable binary queries
and inclusion-minimal separating sets of sizes two and three. Those sets
cannot be the bases of a matroid. The exact table, proof and canonical native
replay are supplied there; no theorem is attributed to an unread comparison
paper. Section 14.6's hitting-set criterion remains the general finite-menu
result. The supplement also states the calibrated-noise limitation and the
remaining primary-source comparison obligations.

\subsection{Appendix U --- Corrected finite toric instrument mathematics}\label{appendix-u-corrected-finite-toric-instrument-mathematics}

The corrected finite toric instrument mathematics has been promoted to canonical Chapter 23 (§23.1--§23.9). The sections below preserve the detailed proof calculations, the \(L=3, t=2\) collision regression analysis, and the mechanized boundary for reference. Lattice assumption throughout: the periodic carrier \(\mathbb Z_L\times\mathbb Z_L\) with \(L\ge 3\), one qubit on each edge of the grid. There are \(n=2L^2\) edges. Horizontal edge \(h_{i,j}\) has index \(2(iL+j)\), and vertical edge \(v_{i,j}\) has index one greater. Face \(f_{i,j}\) is incident to \(h_{i,j}\), \(h_{i+1,j}\), \(v_{i,j}\), and \(v_{i,j+1}\). Star \(s_{i,j}\) is incident to \(h_{i,j}\), \(h_{i,j-1}\), \(v_{i,j}\), and \(v_{i-1,j}\). Indices are periodic. The Pauli space is \(\mathbb F_2^n\oplus\mathbb F_2^n\). The weight convention is
\[
\operatorname{wt}(e)=|\operatorname{supp} e_X|+|\operatorname{supp} e_Z|,
\]
so a single \(Y\) contributes \(2\). Detectors in this appendix are passive queries: a row returns a bit. No post-state is declared, and a classical bit table is not a proof that a physical measurement leaves a state fixed.

\subsubsection{\texorpdfstring{U.1 The commutation kernel is \(S^\perp\), not \(S\)}{U.1 The commutation kernel is S\^{}\textbackslash perp, not S}}\label{u.1-the-commutation-kernel-is-sperp-not-s}

Let \(A\) be the \(L^2\times n\) face-incidence matrix and \(A_v\) the star-incidence matrix, both over \(\mathbb F_2\).

Each edge lies in two faces, so the sum of all rows of \(A\) is zero. The same holds for \(A_v\). The dual graph, with a vertex per face and an edge per qubit, is the connected torus grid. A vector in the left kernel of \(A\) is constant on the two faces of every edge, hence constant on a connected graph. The left kernel is the constants and \(\operatorname{rank} A=L^2-1\). The same connectedness argument gives \(\operatorname{rank} A_v=L^2-1\).

The full stabilizer commutation syndrome is the block map
\[
Q(x\mid z)=(Ax,\,A_v z).
\]
Thus \(\operatorname{rank} Q=2L^2-2\) and, on the Pauli space of dimension \(4L^2\),
\[
\dim\ker Q=4L^2-(2L^2-2)=2L^2+2.
\]

A star meets every face evenly. The four edges of \(s_{i,j}\) are shared two at a time with the four faces \(f_{i,j}\), \(f_{i,j-1}\), \(f_{i-1,j}\), and \(f_{i-1,j-1}\), and every other face is disjoint from that star because those four edges have only eight face incidences. Therefore \(A A_v^{\mathsf T}=0\), and transposition gives \(A_v A^{\mathsf T}=0\).

The stabilizer space \(S\) is spanned by the vectors \((\text{row of }A_v\mid 0)\) and \((0\mid\text{row of }A)\). The commutation relations put \(S\subseteq\ker Q\). The two blocks are independent, so \(\dim S=2L^2-2\). Hence
\[
\dim(\ker Q/S)=(2L^2+2)-(2L^2-2)=4.
\]
The kernel of \(Q\) is the set of Pauli vectors that commute with every generator of \(S\). That set is the symplectic orthogonal \(S^\perp\). It properly contains \(S\), because the quotient dimension is \(4\), not \(0\).

\textbf{WDP-CH23-P01C.} On this carrier, with this weight convention, the full commutation syndrome has kernel \(S^\perp\neq S\), and \(\dim(S^\perp/S)=4\). The imported sentence that a full independent stabilizer measurement has kernel \(S\) is withdrawn. The arithmetic that subtracts \(\operatorname{rank} Q\) from the edge count \(n\), rather than from the Pauli dimension \(2n\), is the source of that mistake. The imported identifier WDP-CH23-P01 keeps its imported wording, including the quotient display of dimension \(4\) already present there; P01C is the statement that is used from here on.

An exhibiting vector is the weight-\(L\) support \(Z_1\), the horizontal edges \(h_{i,0}\) for \(i=0,\ldots,L-1\), placed in the \(X\) block with the \(Z\) block zero. It lies in \(\ker A\). The row-zero horizontal flux, the sum of \(h_{0,0},\ldots,h_{0,L-1}\), is zero on every star, because each star meets that row in zero or two horizontal edges, and it equals \(1\) on \(Z_1\) for every \(L\ge3\): the row intersects that support in exactly \(h_{0,0}\). Thus \(Z_1\) is outside the star span, the Pauli vector is outside \(S\), and its weight is \(L\). At \(L=3\) the shipped binary reconstruction and the Lean encoding both check kernel membership, exclusion from \(S\), and weight \(3\).

\subsubsection{U.2 Distance under the declared weight}\label{u.2-distance-under-the-declared-weight}

Let \(x\in\ker A\). Write \(\alpha_i(x)\) for the sum of the horizontal edges in lattice row \(i\), and \(\beta_j(x)\) for the sum of the vertical edges in lattice column \(j\). The sum of the face rows in lattice row \(i\) equals the horizontal edges of row \(i\) plus those of row \(i+1\), because each vertical edge of that strip lies in two of those faces. Therefore \(Ax=0\) forces \(\alpha_i=\alpha_{i+1}\). Summing the face rows in lattice column \(j\) cancels the horizontal edges and forces the vertical fluxes in columns \(j\) and \(j+1\) to agree. Each star has two horizontal edges in one row and two vertical edges in one column, so both fluxes vanish on the row span of \(A_v\). The pair \((\alpha,\beta)\) therefore descends to \(\ker A/\operatorname{rowspan}(A_v)\) and separates the four classes: it is an isomorphism onto \(\mathbb F_2^2\), since the quotient is two-dimensional and both nontrivial values occur on the exhibited loops of weight \(L\).

If \(\alpha=1\), every lattice row contributes at least one horizontal edge, so the weight is at least \(L\). If \(\beta=1\), the weight is at least \(L\). The same lower bound holds in the other sector. A class in \(\ker Q\setminus S\) is nontrivial in at least one sector, so its weight is at least \(L\). The exhibited loops have weight \(L\), and that is the minimum.

A class that is nontrivial in both sectors has weight at least \(2L\) under the declared weight, because the two supports are counted separately and a shared qubit contributes to both. The imported value \(2L-1\) for the diagonal class counts non-identity Paulis, with a shared qubit read as one \(Y\). That is a different convention.

\textbf{WDP-CH23-P07C.} Under \(\operatorname{wt}(e)=|\operatorname{supp}e_X|+|\operatorname{supp}e_Z|\), the diagonal class has distance \(2L\), not \(2L-1\). However, this distance does not establish the fault-tolerance threshold \(t\le L-1\) for the coarser question \(B=S+\langle u+v\rangle\). The relevant separation distance is
\[
d_B=\min\{\operatorname{wt}(z):z\in\ker Q\setminus B\},
\]
measuring forbidden differences. Because the diagonal class \(u+v\) belongs to the allowed subspace \(B\), its weight of \(2L\) does not determine that distance. The forbidden difference set \(\ker Q\setminus B\) still contains the pure single-sector logical classes \(u\) and \(v\), each of weight \(L\). Thus \(d_B=L\), retaining the single-sector threshold \(2t<L\). The inference that \(2t<2L\) tolerates \(t\le L-1\) fails: at \(L=3, t=2\), the errors \(e_1=h_{0,0}\) and \(e_2=\{h_{1,0},h_{2,0}\}\) have weights one and two, identical full syndromes, and difference \(u\notin B\), producing an undetectable collision. \textbf{WDP-CH23-P08C} withdraws the imported slogans that this halves the required lattice and that the coarser question is factorially cheaper. No cost identity and no lattice-size identity of that form is proved. The coarser question retains the threshold requirement \(2t<L\).

\subsubsection{U.3 The restricted fiber criterion}\label{u.3-the-restricted-fiber-criterion}

Let \(E_{\mathrm{adm}}\) be a declared finite set of Pauli vectors, let \(Q_D\) be a declared detector map, and let \(q\) be a protected question with \(B=\ker q\), here a subspace.

\textbf{WDP-CH23-P03C.} For every admitted pair \(e,e'\),
\[
Q_D(e)=Q_D(e')\quad\text{implies}\quad q(e)=q(e')
\]
if and only if a recovery witness exists: one admitted representative \(e_s\) may be chosen in each nonempty fiber, and the decoder returns that representative. Both selection and verification terminate by finite enumeration. Equivalently, when \(q\) is linear,
\[
(E_{\mathrm{adm}}-E_{\mathrm{adm}})\cap\ker Q_D\subseteq B.
\]
The forward direction subtracts the two recovery representatives. The reverse direction defines the answer on a fiber by the common value of \(q\). The imported wording of WDP-CH23-P03 that begins the forward direction by assuming the difference already lies in \(B\) is not this proof.

A global rank defect is not an admitted collision. The quotient \(S^\perp/S\) has dimension \(4\) and its nontrivial classes have weight at least \(L\). At \(t=0\) the admitted set is \(\{0\}\), so its difference set is
\(\{0\}\subseteq B\). The forbidden-difference set outside \(B\) is empty,
and the inclusion holds. \textbf{WDP-CH23-P10C.} Deleting one independent row from a full-rank syndrome drops the rank by \(1\) and raises \(\dim(\ker Q_D/S)\) from \(4\) to \(5\). The imported count \(\dim\ker Q_D=n-r_D\), using the edge count in place of the Pauli dimension, is withdrawn, and so is the imported quotient dimension \(1\). Whether the enlarged kernel contains a low-weight admitted pair is settled by searching \(E_{\mathrm{adm}}\), not by quoting the dimension.

\subsubsection{\texorpdfstring{U.4 Full rank, the empty purchase at \(t=0\), and the matroid boundary}{U.4 Full rank, the empty purchase at t=0, and the matroid boundary}}\label{u.4-full-rank-the-empty-purchase-at-t0-and-the-matroid-boundary}

Let \(\mathcal E_t\) be the set of Pauli vectors of weight at most \(t\), and let \(B=S\).

\textbf{WDP-CH23-P06C.} Let \(1\le t\) and \(2t<L\), and let \(B=S\). A passive stabilizer-check set of rank \(2L^2-2\) separates \(\mathcal E_t\) up to \(S\). Sufficiency uses the distance proved in U.2: the kernel of the full syndrome is \(S^\perp\), and every vector of weight at most \(2t\) in \(S^\perp\) already lies in \(S\). The converse is withdrawn. Full rank is not necessary for that separation. The sentence that every proper row subspace kills a weight-one column is withdrawn with it.

The column sentence already fails on the face-incidence matrix at \(L=3\). That matrix has rank \(L^2-1=8\). Its rows indexed by \(\{0,1,3,5,7,8\}\) have rank \(6\). No column of this submatrix is zero, so no weight-one edge lies in its kernel. The edge pairs \((0,1)\), \((4,13)\), and \((5,8)\) have columns summing to zero, so the kernel contains those weight-two vectors. The pair \((0,1)\) lies outside the star span. The same six rows, placed in the \(X\) block and accompanied by no star row, do not separate \(\mathcal E_1\).

Full rank also fails to be necessary for the full commutation syndrome. At \(L=4\) that syndrome has shape \((32,64)\) and rank \(30=2\cdot 4^2-2\). Delete rows 1, 10, and 16. The remaining \(29\) rows are independent and have rank \(29\). On \(\mathcal E_1\), the admitted-collision search returns no pair, so this suite separates \(\mathcal E_1\) up to \(S\). Here \(t=1\) and \(2t=2<4\). A dropped row can enlarge the kernel by a vector heavier than \(2t\), and the fiber criterion can still hold.

\textbf{WDP-CH23-P11C and WDP-CH23-P13C.} At \(t=0\) the empty detector set is sufficient and has cost \(0\). A positive purchase is not required. For \(1\le t\) and \(2t<L\) with \(B=S\), rank \(2L^2-2\) is a sufficient independent count. It need not be the minimum independent count: the \(L=4\) suite above separates \(\mathcal E_1\) with \(29\) independent checks. The uniform-cost minimum on a declared finite shelf is the cost returned by the enumeration of P12C.

\textbf{WDP-CH23-P12C.} Greedy selection in the stabilizer-check matroid returns a minimum-cost basis of the check space only after sufficiency has been proved equal to full check-space rank for the obligation in hand. That hypothesis is withdrawn for \(\mathcal E_t\) in the regime \(1\le t\), \(2t<L\), \(B=S\), by the rank-\(29\) suite of P06C. The matroid shortcut is not used for this obligation. For a general finite family \(E_{\mathrm{adm}}\) the obligation is the hitting condition of P03C. The terminating exact procedure is exhaustive enumeration of the shelf subsets under positive rational costs. The optimality checker is a second enumeration. It does not call the producer.

\subsubsection{U.5 Adaptivity, schedules, and noise}\label{u.5-adaptivity-schedules-and-noise}

An adaptive policy is a function from finite transcripts to the next passive query. A transcript is a finite list of selected shelf indices and the bits they returned. The information state is the set of admitted errors consistent with that transcript.

\textbf{WDP-CH23-P17C.} The imported unrestricted claim that adaptivity buys nothing is false. Take the four hypotheses \(\{w,x,y,z\}\) and three passive tests: Split of cost \(2\), partitioning \(\{w,x\}\) from \(\{y,z\}\); Wx of cost \(2\), separating \(w\) from \(x\) and not \(y\) from \(z\); Yz of cost \(2\), separating \(y\) from \(z\). Any two tests leave a colliding pair, so the nonadaptive minimum uses all three and costs \(6\). The adaptive policy buys Split and then the one test named by that bit. Every branch costs \(4\). A claimed adaptive lower bound of \(6\) is false, and the indistinguishable nonadaptive pairs are the certificate. The shipped checker rejects that false bound.

The further sentence that, for \(t\ge 1\), \(2t<L\), and \(B=S\), every proper stabilizer-check subspace hides a weight-one error, and that the all-zero branch must accumulate rank \(2L^2-2\), is withdrawn. The \(L=4\) rank-\(29\) suite of P06C separates \(\mathcal E_1\), so a passive policy may stop on that suite. \textbf{WDP-CH23-P18C} withdraws every unrestricted detector-purchase bound, including that full-rank branch. The \(t=0\) family remains a counterexample to any rule that buys a detector on every admitted family.

\textbf{WDP-CH23-P15C.} The imported cone radius is the definition \(\lceil v_{\mathrm{LR}}\Delta t\rceil\). Comparing it with \(\lfloor L/2\rfloor\) is a geometric comparison. It is not an admitted collision and not a proof that a schedule fails. A failure certificate is still a pair in \((E^R-E^R)\cap\ker Q_{D,R}\) outside \(B^R\), at the decision time, inside the declared family.

\textbf{WDP-CH23-P16C.} If a finite joint family of data and readout errors is actually supplied, together with a map to the observed record and a protected subspace, the fiber criterion of P03C is necessary and sufficient for deterministic factorization on that family. The imported spacetime chain, the Bernoulli readout model, and the repetition threshold \(R\ge L\) are not reconstructed from the bare toric incidence matrices. They remain conditional on those missing matrices. Repeated measurements at different times are not the same row unless a proved state-evolution map identifies them. The detectors used here are passive queries. Nothing in the binary table is a theorem of physical non-disturbance.

\textbf{Current answers and later operations.} A passive query can match the current protected bit while leaving a later declared operation disabled. A witness that sets the current report matched, the future operation disabled, and nevertheless claims that the future operation remains available, is rejected. Preserving the current answer is not that later operation.

\subsubsection{U.6 Executable and mechanized boundary}\label{u.6-executable-and-mechanized-boundary}

The Python reconstruction in \texttt{code/toric\_code\_observer\_completion.py} builds the incidence matrices, the commutation syndrome, the stabilizer rows, the weight-\(L\) logical, the fiber search, the recovery representative, and the separate minimum-cost checker. \texttt{ToricSyndromeCorrespondence.lean} establishes the concrete \(L=3\) encoding: face rank \(8\), star rank \(8\), the logical vector in the face kernel and outside the star span, its weight \(3\), the row flux, and the identity \((2L^2+2)-(2L^2-2)=4\) for every integer \(L\ge 1\). Those Lean facts are not a formalization of the flux-distance argument for every \(L\), of the rank-\(29\) separation counterexample, of adaptivity, or of any physical schedule. The general-\(L\) distance and fiber proofs above are written proofs. The \(L=3\) face-row counterexample and the \(L=4\) rank-\(29\) suite are executable witnesses in \texttt{code/toric\_code\_observer\_completion.py}, replayed by \texttt{code/tests/test\_toric\_code\_observer\_completion.py}.

\subsection{Appendix V --- Accepted locality proofs L1--L4}\label{appendix-v-accepted-locality-proofs-l1l4}

The text below is the complete accepted locality argument, copied into this
manuscript. Exact finite observer completion, conditional Lieb--Robinson
readout bounds, reconstruction conditioning, noise refusal, and CAR locality
with gauge constraints stay separate inside that argument. A Lieb--Robinson
upper bound is not a reconstruction-error floor. The Bell, gauge, star-exchange,
and upper-envelope counterexamples are part of the proof, not a citation.
The live file \texttt{unified\_lieb\_robinson\_observer\_completion\_derivation.md} remains
in place because locality replay pins its path. This copy, not a link to that
file, is the canonical proof.

\section{Locality, sufficient observation, and faithful recovery}\label{locality-sufficient-observation-and-faithful-recovery}

Corrected derivation, 2026-10-05.
The submitted draft is preserved byte for byte in
\href{docs/source_archive/locality_and_smearing/source_manifest.json}{\texttt{docs/source\_archive/locality\_and\_smearing/source\_manifest.json}} in \texttt{docs/source\_archive/locality\_and\_smearing/}.
Its statements are not inherited as proved propositions.

\subsection{1. The precise bridge}\label{the-precise-bridge}

There are two different restrictions: information discarded by an observer,
and operations excluded by physical locality. They meet at a declared readout
map; neither restriction supplies the missing physical adapter for the other.
The theorem below gives exact finite recovery conditions, then bounds the
conditioning of a physical readout under explicit locality hypotheses.
It does not identify a positive error upper bound with an error floor.

Fix a nonempty finite set E of admitted errors in V = F₂ⁿ, a repair subspace
B ≤ V, and a linear syndrome map Q:V→F₂ᵐ with B ≤ ker Q. Repair success
means e+r∈B. An error in B has zero quotient obstruction; an error outside B
does not become zero by selecting another representative. These are
algebraic declarations. An actual decoder operation realizing r must still
be supplied to use this model for a device or a proposal-protocol repair.

\textbf{Theorem L1 (sharp admitted-error recovery).} A syndrome-only decoder
r:Q(E)→V succeeds for every e∈E if and only if

\[
 Qe=Qf\quad\Longrightarrow\quad e+f\in B \qquad(e,f\in E).
\]

\textbf{Proof.} If r succeeds, add e+r(Qe)∈B and f+r(Qf)∈B; the common repair
cancels. Conversely select one admitted representative e\_s from each nonempty
fiber Q⁻¹(s)∩E and set r(s)=e\_s. The hypothesis gives e+e\_s∈B.
Both selection and verification terminate by finite enumeration. A positive
certificate lists every used syndrome and its repair. A negative certificate
is a pair e,f in the same fiber with e+f∉B. The checker replays the supplied
repair or pair; it does not replace an incorrect decoder with a new one. ∎

This is not a left inverse of Q, RQ=id\_V. That equation is unrealizable for a
noninjective Q. Recovery is only required on E, modulo the declared B.
For a full Hamming ball E=\{e:wt(e)≤t\}, define
d=min\{wt(v):v∈ker Q∖B\}, with d=∞ when this set is empty. Then L1 is
equivalent to 2t\textless d: a forbidden vector of weight at most 2t splits into two
vectors of weight at most t, and the sum of any two admitted errors has
weight at most 2t. For a restricted admitted family, L1 is the criterion;
the distance inequality is sufficient but need not be necessary.

\textbf{Theorem L2 (minimum linear observer completion).} Additional binary rows
must annihilate B. A row suite C permits recovery precisely when every
forbidden difference

\[
 F=\{e+f:e,f\in E,\ Qe=Qf,\ e+f\notin B\}
\]

has a nonzero value under at least one row of C. Thus its minimum size is
the minimum finite hitting suite among the rows in Bᵖᵉʳᵖ. There is an
exhaustive terminating procedure: enumerate row subsets by increasing size,
discard rows that do not annihilate B, and test all forbidden differences.
A lower certificate must cover every smaller admitted suite with a retained
collision. Just displaying a successful suite is only an upper bound.

If E is itself a linear subspace containing B, the minimum is
dim(ker(Q\textbar E)/B). Necessity follows from injectivity of the added map on that
quotient. Sufficiency follows by choosing dual coordinate functionals there,
extending them to E/B and then V/B. This dimension formula is not asserted
for an arbitrary restricted E.

The retained example has n=3, Q(x,y,z)=x+y, B=\{000,110\}, E=F₂³.
The original syndrome collides on 000 and 001. One extra row z produces
(x+y,z), with repair (0,x+y,z). Zero extra rows fail and one succeeds, so
the exact minimum is one. \texttt{LocalObservationCompletion.lean} establishes this
concrete decoder separately from its abstract fiber criterion.

\subsection{2. Conditional physical locality theorem}\label{conditional-physical-locality-theorem}

Let a finite tensor-product quantum system have a declared region R. Let
ρ,σ be density matrices. Let O be a Hermitian evolved readout and O\_R a
Hermitian operator supported on the \textbf{joint} region R. Assume

\[
 \operatorname{Tr}_{R^c}\rho=\operatorname{Tr}_{R^c}\sigma,
 \qquad \|O-O_R\|_\infty\le\beta.
\]

\textbf{Theorem L3 (local readout bound).}

\[
 |\operatorname{Tr}((\rho-\sigma)O)|
 \le\|\rho-\sigma\|_1\,\beta\le2\beta.
\]

\textbf{Proof.} Equal joint marginals give
Tr((ρ−σ)O\_R)=0. Subtract this identity and apply trace/operator norm
duality. The trace distance bound follows from positivity and unit trace:
∥ρ−σ∥₁≤∥ρ∥₁+∥σ∥₁=2. ∎

With approximate marginals at trace-norm distance d\_R, the bound is instead
d\_R∥O\_R∥∞+∥ρ−σ∥₁β. No correlation-decay assumption is needed for equal
joint marginals. Equal \textbf{single-site} marginals do not give equal joint
marginals on a multi-site region.

Here is an explicit way to obtain O\_R from finite commutator estimates.
For a qubit y outside R let E\_y be the average of conjugation by I,X,Y,Z
on y. It is a norm contraction and removes dependence on that qubit.
Composing these maps gives E\_R, the conditional expectation onto operators
supported on R. If ∥{[}O,P\_y{]}∥∞≤η\_y for each of the three nonidentity
Paulis, then

\[
 \|O-E_R O\|_\infty\le\tfrac34\sum_{y\notin R}\eta_y.
\]

Indeed O−E\_yO is the average of the three differences O−P\_yOP\_y;
each has the commutator norm. A telescoping sum over y, with preceding
conditional expectations contractive, establishes the claim. Higher local
dimensions use a declared finite unitary twirl and its corresponding average.

An imported Lieb--Robinson bound can therefore provide a β of the form
C\_R exp{[}−μ(r−v\textbar t\textbar){]} after the spatial tail has actually been summed.
C\_R includes the geometry, operator norm, and interaction hypotheses.
This derivation does not infer such a bound from the name of a Hamiltonian.
The standard locality results are imported from
\href{https://arxiv.org/abs/quant-ph/0603121}{Bravyi--Hastings--Verstraete} and
\href{https://arxiv.org/abs/1810.02428}{Nachtergaele--Sims--Young}.
The finite twirl and the readout deduction above are self-contained.

The equality Tr({[}ρ,O{]})=0 holds for every pair of finite square matrices.
It is \textbf{not} an expression for Tr((ρ−σ)O). Confusing these two expressions
would wrongly erase every readout difference.

\subsection{3. Conditioning and noisy refusal}\label{conditioning-and-noisy-refusal}

Let z=R(ρ) and z'=R(σ) be readout vectors, in a declared norm. Suppose
∥z−z'∥≤η, a decoder F satisfies ∥F(z)−F(z')∥≤κ∥z−z'∥, and its endpoint
errors against the desired targets a,b are at most α.

\textbf{Theorem L4.} If ∥a−b∥=δ, then

\[
 \delta\le2\alpha+\kappa\eta.
\]

For exact endpoint decoding and η≤2β, κ must satisfy κ≥δ/(2β) when β\textgreater0.
If β=0 and δ\textgreater2α, no such decoder exists. This is an obstruction to stable
reconstruction, not a lower bound on the physical error of the device.
For multiple channels their common β and vector norm must be declared;
one cannot silently use a scalar bound for an arbitrary-dimensional norm.

\textbf{Proof.} Insert F(z),F(z') between a and b and apply the triangle
inequality. Combining the locality bound uses κ≥0. ∎

If two readouts have distance at most 2ε in a normed real space, their closed
noise balls of radius ε intersect at the midpoint. At the common report,
any reconstruction has error at least δ/2 at one endpoint. Thus β≤ε is
a sufficient indistinguishability condition in the scalar bound of L3.
It is not a necessary physical horizon, nor evidence that the system has
intrinsic error ε. The Lean proof establishes the common-report contradiction
for exact rational targets and errors.

\subsection{4. Retained carrier counterexamples}\label{retained-carrier-counterexamples}

\textbf{Bell marginals.} The rational density matrices of
(\textbar00⟩±\textbar11⟩)/√2 have identical individual marginals I/2. The observable
X⊗X has expectations +1 and −1. Keeping current single-site outputs loses
the joint readout. Setting O\_R=Z⊗I and O=O\_R+εX⊗X gives readout difference
2ε and tail norm ε: the bound in L3 is sharp on these finite data.
The executable checker checks density positivity, exact marginal equality,
the submitted norm certificates, and the actual readout separately.

\textbf{Gauge cycle.} For oriented edges a→b→c→a and α=(0,1,0), the gradient
is (1,−1,0) if dα(edge)=α(head)−α(tail), or (−1,1,0) for the opposite
convention. Its L¹ norm is two, and its circulation is zero. Every gauge
gradient telescopes around a cycle. Translation of the admitted repair set
therefore preserves the quotient obstruction and the infimal repair cost.
It cannot extinguish a nonzero cycle class. Changing the carrier, weights,
or admissible repairs is a different operation requiring a preservation map.

\textbf{Star exchange.} Four physical sites 0,1,2,3 form a star centered at 0.
The two-particle configuration loop is
\{1,2\}→\{0,2\}→\{2,3\}→\{0,3\}→\{1,3\}→\{0,1\}→\{1,2\}.
Its legal directed hops are 1→0,0→3,2→0,0→1,3→0,0→2.
In ascending CAR occupation order their signs are 1,−1,1,1,−1,−1,
with product −1. Six signs all equal to −1 would instead multiply to +1.
The exchange loop lives in the configuration graph, not the physical site
tree. The CAR/exterior-power carrier construction must be declared; it is
not an inclusion of physical site cycles into configuration cycles.

\textbf{Upper envelopes.} An inequality 0≤error≤positive\_bound does not imply
error\textgreater0. Infimizing an upper envelope does not change this fact. A lower
obstruction requires a collision, a separated pair, or another independently
checked witness of the missing capability.

\subsection{5. Scope relative to WDP, checked repair and local acceptance}\label{scope-relative-to-wdp-checked-repair-and-local-acceptance}

L1 separates local syndrome satisfaction from global admitted-error recovery,
the local-acceptance distinction. L2 specifies what an adequate observer must retain.
The quotient ker Q/B connects the raw defect to the declared repair equivalence.
L3--L4 give a physical readout restriction \textbf{only when a real carrier/readout
adapter and locality hypotheses are supplied}. WDP binds those source,
observer, repair and evaluation maps, instead of letting a name substitute
for them. Proposal, independent checking and authorization remain separate.

There is no claim here that MUB6 is resolved, that a Majorana device realizes
these declarations, that every noise model is correctable, or that a physical
error floor follows from a Lieb--Robinson upper bound. The standard quantum
error-correction setting is imported separately from
\href{https://arxiv.org/abs/quant-ph/9604034}{Knill--Laflamme}; L1 is a finite
binary admitted-error theorem and is not a proof of that full operator theorem.

\subsection{6. Evidence boundary}\label{evidence-boundary}

The full written theorem uses ordinary finite linear algebra, trace norms,
and conditional imported analytic locality bounds. Python replays finite
binary recovery, exact rational matrices, the star exchange, and rejection
controls. \texttt{LocalObservationCompletion.lean} establishes the fiber criterion,
the concrete binary decoder, rational locality/conditioning consequences,
and cycle/sign identities. It does not formalize the imported analytic
Lieb--Robinson theorem or verify the Python program. The reference package
and source-bound verification records are described in
\href{docs/derivation_evidence.md}{\texttt{derivation\_evidence.md}}.

\subsection{Appendix W --- Accepted Trotter proofs and the Trotter-response boundary}\label{appendix-w-accepted-trotter-proofs-and-the-trotter-response-boundary}

The text below is the complete accepted Trotter argument, copied into this
manuscript: finite pairing, fixed-carrier time consistency, positivity, the
strict margin, the rational schedule, and cochain descent. Exact finite
rational certificates, the all-\(M\) statement for one compressed channel,
full fixed-four-site time convergence, finite spectral observations at
selected \(M\), general all-\(M\) inertia, and infinite-volume or joint
spacetime convergence stay separate. The live file
\texttt{trotter\_cochain\_extinction\_derivation.md} remains in place because locality
replay pins its path. This copy, not a link to that file, is the canonical proof.

\section{Trotter response, smooth observations, and cochain descent}\label{trotter-response-smooth-observations-and-cochain-descent}

Corrected derivation, 2026-10-05.
Original submitted bytes are retained under
\href{docs/source_archive/locality_and_smearing/source_manifest.json}{\texttt{docs/source\_archive/locality\_and\_smearing/source\_manifest.json}} in \texttt{docs/source\_archive/locality\_and\_smearing/}.
The theorem closes a fixed finite spatial carrier's time limit. It does not
assert a joint infinite-volume/continuum limit or Gaussian field law.

\subsection{1. The exact supplied finite carrier}\label{the-exact-supplied-finite-carrier}

Take L=4 sites and the full 16-dimensional CAR Fock space. In the ascending
occupation basis, c\_x removes site x with sign (−1) raised to the number of
occupied sites preceding x. Define

\[
 n_x=c_x^*c_x,\quad a_x=(-1)^x(n_x-\tfrac12I),\qquad
 P_x=\tfrac12(n_x+n_{x+1}-2n_xn_{x+1}+c_x^*c_{x+1}+c_{x+1}^*c_x).
\]

Indices are periodic. Each P\_x is an orthogonal projection. Disjoint even
bonds commute, as do disjoint odd bonds. Let 0\textless θ≤1, M≥1, r=θ/M, N=2M,
T₀=(I+rP₀)(I+rP₂), T₁=(I+rP₁)(I+rP₃), and K\_j=T\_(j mod 2).
The input partition is

\[
 Z_M(z)=(1+r)^{-M}\operatorname{Tr}
 \bigl(K_{N-1}e^{\sum_xz_{N-1,x}a_x/2}\cdots K_0e^{\sum_xz_{0,x}a_x/2}\bigr),
 \qquad G_M=4\,D_z^2\log Z_M(0).
\]

The scalar prefactor cancels from G. Its matrix dimension is LN=8M;
this is a spacetime response matrix, not a one-particle Hilbert-space matrix.
For sampled real profiles ψ\_(j,x), the declared physical quadratic is

\[
 q_M(\psi)=\frac{1}{L N^2}\psi^TG_M\psi.
\]

This normalization is used throughout. The original Trotter-response implementation
is reconstructed exactly at M=1; the new reference extends the same insertion
convention, using full traces rather than guessed spectral formulas.

\subsection{2. Reconstructed finite spectra}\label{reconstructed-finite-spectra}

Independent congruence witnesses establish these inertias at θ=1/10.
Counts below are (positive, negative, zero).

\begin{center}
\small
\begin{longtable}{cc>{\centering\arraybackslash}p{0.32\textwidth}>{\centering\arraybackslash}p{0.34\textwidth}c}
\toprule
$M$ & step $r$ & normalized partition & staggered eigenvalue & inertia \\
\midrule
1 & $1/10$ & $\frac{17661}{1100}$ & $-\frac{1}{706440}$ & $(4,1,3)$ \\[0.4em]
2 & $1/20$ & $\frac{755374987}{47040000}$ & $-\frac{35281}{90644998440}$ & $(8,1,7)$ \\[0.4em]
3 & $1/30$ & $\frac{1395062306051761}{86870556000000}$ & $-\frac{12188547841}{55802492242070440}$ & $(12,1,11)$ \\[0.4em]
4 & $1/40$ & $\frac{2379246859363301568641}{148151258316800000000}$ & $-\frac{14588472130734881}{95169874374532062745640}$ & $(16,1,15)$ \\
\bottomrule
\end{longtable}
\end{center}

For M=1 the full spectrum is 10/21,21/40,420/841 (twice),
0 (three times), and −1/706440. At a symbolic one-period r,

\[
 \lambda_1(r)=\frac{-r^4}{4(r+2)(r^2+4r+8)}.
\]

Thus λ₁(1/192)=−1/87426958080. The denominator is
64+64r+24r²+4r³, giving
λ₁=−r⁴/64+r⁵/64−5r⁶/512+O(r⁷).
At r=1/10 the quadratic factor is 841/100, not 281/100.
The table verifies four finite matrices; it is not an all-M inertia theorem.
The time-limit theorem below does not depend on that unproved extrapolation.

\subsection{3. A sharp finite pairing theorem}\label{a-sharp-finite-pairing-theorem}

Let v\_(j,x)=(−1)\^{}j, so ∥v∥²=LN, and put
bar A=(LN)⁻¹Σ\_(j,x)(−1)\^{}jψ\_(j,x). If
\textbar ψ\_(2k,x)−ψ\_(2k+1,x)\textbar≤D/N, then

\[
 |\bar A|\le\frac{D}{2N}=\frac{D}{4M}.
\]

\textbf{Proof.} Pair the N terms into M adjacent differences. The total absolute
sum is at most LM D/N; divide by LN and use N=2M. A linear sampled
profile ψ\_j=j/N attains equality, so the finite pair bound is sharp. A
C¹ profile with sup\textbar∂\_τψ\textbar≤D supplies this premise by the mean-value theorem.
Time-independent profiles have bar A=0 exactly. ∎

For the rank-one compressed channel P\_vG\_MP\_v, set
λ\_v=vᵀG\_Mv/(LN). Its contribution is exactly

\[
 q_{v,M}=\frac{\lambda_v}{N}\bar A^2,\qquad
 |q_{v,M}|\le\frac{|\lambda_v|D^2}{4N^3}
             =\frac{|\lambda_v|D^2}{32M^3}.
\]

No eigenvector hypothesis is needed for this compression identity. When v
is the sole negative eigenvector, as in the four checked instances, it is
the entire negative spectral contribution. Outside those instances this
identification needs its own proof.

\textbf{All-M extinction of this channel.} Set E=e\^{}(4θ). Products of the K\_j
have norm at most E. Tr((T₁T₀)\^{}M)≥16 because T₀,T₁≥I and T₁T₀ is
similar to T₀\textsuperscript{(1/2)T₁T₀}(1/2)≥I. First and second normalized insertion
moments have magnitudes at most E/2 and E/4. Hence each response entry
has magnitude at most (E+E²)/4, and
\textbar λ\_v\textbar≤2M(E+E²). Consequently

\[
 |q_{v,M}|\le\frac{(E+E^2)D^2}{16M^2}\longrightarrow0.
\]

This establishes extinction for the declared smooth observed channel without
postulating bounded eigenvalues or extrapolating the finite inertia table.
It does not by itself prove positivity of the full q\_M.

\subsection{4. The fixed-L time-continuum theorem}\label{the-fixed-l-time-continuum-theorem}

Let H=Σ\_xP\_x, and let ψ\_x(τ) be real, periodic C¹ functions with
\textbar ψ\_x\textbar≤B and \textbar ψ'\_x\textbar≤D. Set A(τ)=Σ\_xψ\_x(τ)a\_x, Z=Tr e\^{}(θH),
and ρ=e\^{}(θH)/Z. Define the continuous time-ordered covariance kernel

\[
 K_{xy}(s)=Z^{-1}\operatorname{Tr}(e^{(1-s)\theta H}a_y e^{s\theta H}a_x)
              -\operatorname{Tr}(\rho a_x)\operatorname{Tr}(\rho a_y),\quad0\le s\le1.
\]

For τ\_i≤τ\_j use s=τ\_j−τ\_i; reverse order transposes x,y. This produces
a continuous periodic symmetric kernel on pairs: at the endpoints the
trace is cyclic, and at coinciding times the a\_x commute. Define

\[
 Q_\theta(\psi)=\frac1L\sum_{x,y}\int_0^1\!\int_0^1
 \psi_x(\tau)K_{xy}(\sigma-\tau)\psi_y(\sigma)\,d\tau d\sigma.
\]

\textbf{Theorem T1 (full consistency, explicit envelope).} For the supplied
L=4 carrier and samples ψ\_(j,x)=ψ\_x(j/N),

\[
 |q_M-Q_\theta|\le\frac{A_\theta(B,D)}{M},\qquad
 A_\theta=L\{30(\theta+\theta^2)e^{16\theta}B^2
                   +2\theta e^{4\theta}B^2+\tfrac12e^{16\theta}BD\}.
\]

\textbf{Proof, with conservative constants.} Write H₀=P₀+P₂,H₁=P₁+P₃.
Their norms are at most two and ∥H∥≤4. Taylor's integral remainder gives
∥T\_i−e\textsuperscript{(rH\_i)∥≤r²e}(2r). Duhamel's double integral gives
∥e\textsuperscript{(rH₁)e}(rH₀)−e\textsuperscript{(rH)∥≤(r²/2)∥{[}H₁,H₀{]}∥e}(4r)
≤4r²e\^{}(4r). Thus a complete period has error at most 6r²e\^{}(4r).
The same bound holds in the other layer order. Telescoping k≤M periods
gives 6Mr²e\^{}(4θ). For an unmatched half-period, the identity expansion
and ∥H\_i∥,∥H/2∥≤2 give
∥T\_i−e\textsuperscript{(rH/2)∥≤2r+r²+(e}(2r)−1)≤6re\^{}(2r), since r≤1.
Multiplying the remaining segment and adding the complete-period error
therefore bounds \textbf{every} ordered segment by
ε\_M=6(θ+θ²)e\^{}(8θ)/M against its ideal e\^{}(θ length H/N).
This covers both starting parities and cyclic arcs.

For each trace there are 16 dimensions and ∥a\_x∥=1/2. Both discrete and
ideal normalizing traces are at least 16; their difference is at most
16ε\_M. The quotient identity u/d−u'/d'=(u−u')/d+u'(d'−d)/(dd')
then bounds first-moment errors by ε\_M(1+E), where E=e\^{}(4θ).
For the two-arc second moment the product telescopes into two errors,
giving ε\_M(2E+E²). Products of first moments differ by at most
E ε\_M(1+E). Thus each covariance-entry error is at most
5ε\_ME²=30(θ+θ²)e\^{}(16θ)/M.
Multiplying by (LN²)⁻¹Σ\textbar ψ\_iψ\_j\textbar{} yields the first term of A\_θ/M.

For temporal quadrature, \textbar K\_xy\textbar≤(E+E²)/4≤e\^{}(16θ)/2.
Differentiating a time segment introduces θ{[}H,a\_x{]}; trace duality gives
\textbar K'\_xy\textbar≤2θE. This bound also controls the kernel across its continuous
coinciding-time seam. The integrand has Lipschitz bound
2θEB²+(e\^{}(16θ)/2)BD in each time coordinate. Each grid cell has total
coordinate error at most 2/N. Summing all x,y and dividing by L yields
L(2θEB²+(e\^{}(16θ)/2)BD)/M. Add the two bounds. ∎

No spatial quadrature is taken in T1: the x indices are the original four
sites. A joint L,M limit needs uniform spatial estimates and an actual
spatial reconstruction map. Sampling a spatially oscillating function on
a fixed grid cannot supply those hypotheses.

\subsection{5. Positivity, strict margins and exact rational schedules}\label{positivity-strict-margins-and-exact-rational-schedules}

\textbf{Theorem T2.} Q\_θ(ψ)≥0. Strict positivity requires a nonzero channel
in the spectral formula; a nonzero profile alone need not suffice.

\textbf{Proof.} The supplied H and a\_x are real symmetric. Choose a real
orthogonal eigenbasis for H, writing its eigenvalues h\_a
and thermal weights w\_a=e\^{}(θh\_a)/Z. Expand A(τ) in temporal Fourier modes.
For each unequal-energy ordered pair a,b and frequency ω\_p=2πp, the
real covariance coefficient is

\[
 \frac{(w_a-w_b)\,\theta(h_a-h_b)}
 {\theta^2(h_a-h_b)^2+\omega_p^2}\ge0.
\]

Its contribution is this coefficient times \textbar Â\_ab(p)\textbar²/L. The unsymmetrized
coefficient is (w\_a−w\_b)/(θ(h\_a−h\_b)+iω\_p). Since A(τ) is real symmetric
in this eigenbasis, Â\_ab(p)=Â\_ba(p); the a,b and b,a coefficients are
complex conjugates at the same frequency, so their imaginary parts cancel.
This real-source hypothesis is part of this theorem; the argument is not
asserted for arbitrary time-dependent complex matrix sources.
At equal energies only p=0 remains. Their combined term is the thermal
variance of the energy-block-diagonal part of the temporal average of A,
which is nonnegative by the weighted square inequality. Subtracting the
mean occurs in this variance term, not separately at every frequency.
These terms exhaust the finite spectral trace and its Fourier expansion.
C¹ profiles justify the summation by square summability. ∎

This is the usual finite thermal spectral argument, derived here from the
declared H. It contains neither an unexplained \textbar k\textbar{} Fourier weight nor a
replacement of time-dependent Fourier coefficients with only their p=0
averages. Positivity of a covariance also does not establish Gaussianity.

\textbf{Corollary T3 (eventual finite positivity).} If an independently justified
margin Q\_θ(ψ)≥m\textgreater0 is supplied, q\_M\textgreater0 for every

\[
 M\ge M_0=\max(1,\lfloor A_\theta/m\rfloor+1).
\]

The inequality is strict. A ceiling alone fails when A\_θ/m is an integer.
The result uses a uniform error envelope for all later M, not a single
positive sample or an assumption of monotonicity.

There is a concrete source-derived positive-margin example: ψ\_x(τ)=1.
Then A=Σa\_x has trace zero, Tr A²=16, and thermal mean zero. The last
statement follows from the unitary cyclic site translation, which preserves
H and sends A to −A. In the constant-frequency Kubo trace each spectral
weight before division by Z is ∫e\^{}((1-s)θh\_a+sθh\_b)ds≥1, since H≥0.
Therefore Q\_θ≥Tr A²/(LZ)≥1/(4e\^{}(4θ)).

For exact rational schedules replace e\^{}(kθ) by
E\_k=(1−kθ/32)\^{}(-32), k∈\{4,8,16\}, 0≤θ≤1. The inequality
−log(1−u)≥u establishes E\_k≥e\^{}(kθ). The reference computes a rational A
and m=1/(4E₄), then checks M₀m\textgreater A independently. The margin is bound
to the constant source profile; it is not silently assigned to other fields.

An optional general kernel reconstruction has cell values L G\_ij, with each
cell measure 1/(LN), because its double integral must equal q\_M. If its
L¹ kernel error is δ and profile interpolation error η, then a separate
declared absolute kernel bound κ yields a quadratic error at most
δB²+κ(2Bη+η²). A convolution of K with the uniform grid converges to
the \textbf{average} of K; it does not reproduce K. This alternative cannot
substitute for T1 without its own kernel-binding evidence.

\subsection{6. Exact cochain descent and its limits}\label{exact-cochain-descent-and-its-limits}

On a declared periodic N-cycle let (Df)\emph{j=f}(j+1)−f\_j. Then
im D=\{x:Σx\_j=0\}; a primitive is constructed by cumulative sums.
The Euclidean quotient norm on C¹/im D is
√N times \textbar mean(x)\textbar. Constants are harmonic, not exact cochains.
The staggered vector is exact: D(−v/2)=v.

\textbf{Theorem T4 (adapter descent).} For a declared target vector space W (with scalar R as the one-dimensional case), a linear adapter A:C¹→W factors
through C¹/B if and only if A(B)=0. An orthogonal low-pass projector Π\_W:C¹→W
alone need not satisfy this. The valid contractive adapter onto W is
Π\_W(I−P\_B), with P\_B the orthogonal projector onto the \textbf{declared} B.
It can discard further quotient information; descent is not faithfulness.

\textbf{Proof.} Descent implies A(x+b)=A(x), hence Ab=0. Conversely this
identity makes the class formula {[}x{]}↦Ax well-defined. Orthogonal projections
are contractions, so their composition is contractive. ∎

For the full cycle B=im D, the adapter keeps only the constant mode if
0∈W, and is zero otherwise. A nonzero low-frequency exact mode is a
counterexample to the claim that low-pass filtering alone kills B. To retain
low-frequency fluctuations while rejecting only the staggered mode, one
must instead declare the smaller B=span(v), and must justify that repair
space for the represented task. Neither choice is automatically the CAR
physics' repair relation. In a two-dimensional cochain complex, the general
Hodge decomposition also includes im d₁*, the coexact component; only on
cocycles does quotienting im d₀ leave precisely the harmonic component.

For integrated edge cochains, subdivision by k maps each edge x\_j to k
copies of x\_j/k. Potentials interpolate linearly, so D'F₀=F₁D.
Summing the k fragments is a left inverse, and periods are preserved.
If fine edge weight is k, Σ\_fine k(x\_j/k)²=Σ\_coarse x\_j², giving the
declared weighted L² isometry and quotient-norm preservation. Unweighted
L² is not invariant. This convention is separate from WDP's L¹ price norm.

\subsection{7. Evidence and remaining scope}\label{evidence-and-remaining-scope}

The written T1--T3 prove the fixed four-site time limit and its positive-margin
case. The general rank-one smearing theorem does not require an all-M
negative-eigenspace classification. Python retains four original-operator
matrices and independent congruence certificates, pairing witnesses,
cycle primitives, weighted subdivisions and mutation controls. Lean establishes
the finite pairing inequality, normalization, square/energy bounds,
strict-margin consequences, and quotient-adapter/cycle identities.
The analytic exponential, spectral and C¹ arguments are written proofs;
they are not advertised as Lean formalization of the physical CAR evolution.
Python is also not formally verified by the Lean mathematical model.

General all-M inertia, uniform infinite-volume estimates, joint spatial
reconstruction, cumulant control, Gaussian limits and native physical repair
bindings remain separate obligations. This scope does not close the paper's
full continuum or guarded-history registers.

\subsubsection{W.8 Trotter-response, and the closures that stay open}\label{w.8-trotter-response-and-the-closures-that-stay-open}

The active Trotter-response checker and its historical verification record stay where the locality
replay already hashes them:
\texttt{docs/meron\_field\_limit\_derivations/code/gate\_e\_exact\_counterexample.py} and
\texttt{docs/meron\_field\_limit\_derivations/code/gate\_e\_exact\_counterexample\_receipt.json}.
This section copies the verification record's mathematical content. It does not replace
the verification record and it does not rerun the writer that would overwrite it.

The checked point is \(L=4\), \(M=1\), \(r=1/10\). The verification record gives
eigenvalue \(-1/706440\), test vector \((-1,-1,-1,-1,+1,+1,+1,+1)\), and
quadratic value \(v^{\mathsf T}Gv=-1/88305\), with
\texttt{is\_positive\_semidefinite} false. The point \(r=1/10\) lies outside the
declared region \(r\le 1/192\). One finite point does not establish a
unrestricted finite-\(M\) claim, and it does not establish an asymptotic formula.
The four inertias in the Trotter section above are four finite observations.
No finite list of inertias is an all-\(M\) inertia formula.

The compressed-channel extinction proved above is an all-\(M\) statement for
that declared rank-one channel only. Theorem T1 is full fixed-four-site time
consistency. Neither statement is a spatial continuum theorem. A cochain
factorization, including Theorem T4, does not supply a realizable physical
operation.

Independently supported static Gaussian and bounded time-labelled cumulant
results remain under the hypotheses already stated for those results. They
are not a joint spacetime field limit. A finite signed generating function is
not a probability characteristic function. Euclidean covariance positivity,
including Theorem T2, is not a physical quantum reconstruction theorem.

Earlier model responses proposed closing the uniform-volume, joint-limit,
and extinction/positivity obligations. Those responses remain proposals.
Their missing obligations are the
estimates those responses name and do not finish: a uniform infinite-volume
bound, a joint spacetime limit under the source normalization and sign
conventions, and an extinction or positivity claim that is not the one-point
Trotter-response verification record. Those obligations stay open. This appendix does not enlarge
the Meron programme and does not give those claims a new name in place of the
missing estimate.

\subsection{Appendix X --- Calibrated reports, finite recovery, and observation cost}\label{appendix-x-calibrated-reports-finite-recovery-and-observation-cost}

This appendix connects the calibration bounds of Chapter 13 and Appendix I.4
to the report-fiber criterion and the finite cost results of Sections
14.6--14.7 and Appendix S. The new step is an exact finite representation of
continuous interval reports. Calibration, units, and possible report
combinations are explicit premises. A checked answer does not authorize an
operation or establish that a physical instrument satisfies those premises.

\subsubsection{X.1 Reports and their error bounds}\label{x.1-reports-and-their-error-bounds}

Let \(X\) be a nonempty finite candidate set, \(p:X\to P\) a free observation,
and \(q:X\to Q\) the protected answer. There are finitely many instruments
\(i\in I\), each with price \(c_i>0\) in one declared cost unit. A query leaves
\(x\in X\) unchanged. A complete report vector is fixed for the experiment;
a query reveals one coordinate. Strategies are deterministic functions of
the reports already received. Each instrument may be queried at most once. Repeated
fresh measurements, changing states, and statistical noise reduction need a
different model; Section 14.7 treats supplied state changes.

Instrument \(i\) has a declared raw coordinate and a calibrated coordinate.
For each admitted realization of candidate \(x\), let \(t_i\) be its ideal raw
indication and \(r_i\) its recorded raw report. Assume

\[|r_i-t_i|\le\eta_i(x),\qquad
  |a_it_i+b_i-f_i(x)|\le\kappa_i(x),\]

where \(a_i\ne0\), \(\eta_i(x),\kappa_i(x)\ge0\). The parser applies the fixed
nominal conversion \(z_i=a_i r_i+b_i\). The ideal indication \(t_i\) may vary
across admitted calibration realizations; it is not another observed input.
The gain has calibrated-unit/raw-unit
dimension; \(b_i\), \(f_i\), and \(\kappa_i\) have calibrated units. The first
bound includes any rounding or quantization error in the raw indication.
The second includes the declared calibration-model error. These are uniform
bounds over the admitted realizations for \(x\), not confidence levels.

\textbf{Proposition X.1 (transported enclosure).} Every such report belongs to

\[J_i(x)=[f_i(x)-\epsilon_i(x),f_i(x)+\epsilon_i(x)],\qquad
  \epsilon_i(x)=|a_i|\eta_i(x)+\kappa_i(x).\]

\textbf{Proof.} Subtract \(f_i(x)\) from \(a_i r_i+b_i\), insert
\(a_it_i\), and apply the triangle inequality. Both terms have calibrated
units. No independence assumption is used. \(\square\)

The calibration chain in Appendix I.4 can provide an error budget: its
radius update is \((|a|+\alpha)r+\alpha|u|+\beta+\eta\) at nominal input
center \(u\). To use such a budget here, its direction and units must match
the reported coordinate and it must hold uniformly over the candidate's
admitted raw indications. A bound on a different quantity or only one
calibration sample does not supply this premise. Tighter verified enclosures
may replace \(J_i(x)\) throughout the argument.

Let \(L_x\) be the set of physically or symbolically admitted complete report
vectors for \(x\). \textbf{Outer-bound validity} means

\[\varnothing\ne L_x\subseteq\prod_{i\in I}J_i(x).\]

This condition suffices for answer soundness and conservative upper cost
bounds. \textbf{Full report realizability} is the stronger condition

\[L_x=\prod_{i\in I}J_i(x).\]

It is an adversarial product model, not a statement about stochastic
independence. It permits every combination of the individual interval
values for each fixed candidate. Error propagation alone does not establish
it. Correlated calibration errors can violate it even when each individual
interval is exact.

\subsubsection{X.2 Certified answers and abstention}\label{x.2-certified-answers-and-abstention}

A history consists of the free report \(p_0\) and calibrated readings \(z_i\)
for a queried subset \(H\subseteq I\). Define

\[C(h)=\{x\in X:p(x)=p_0,\ z_i\in J_i(x)\text{ for every }i\in H\}.\]

\textbf{Proposition X.2 (answer certificate).} If outer-bound validity holds and
the true candidate supplied the history, it belongs to \(C(h)\). A proposed
answer \(v\) is therefore correct if the checker establishes

\[C(h)\ne\varnothing\quad\text{and}\quad
  (\forall x\in C(h))\ q(x)=v.\]

If \(C(h)\) contains two different answers, this checker abstains or requests
another instrument. If \(C(h)=\varnothing\), it reports inconsistency with
the model; it cannot accept by vacuous truth.

\textbf{Proof.} Each genuine reading satisfies its enclosure, and its free
report equals \(p(x)\). Thus the true candidate satisfies every defining
constraint. Apply the checked answer equality to that candidate. \(\square\)

Under full report realizability this criterion is also necessary for a
uniformly correct answer at a compatible history. Indeed each candidate
in \(C(h)\) admits a complete report vector extending the history: choose any
point in each unqueried nonempty interval. Two such candidates with different
answers give the same observed history, so one answer cannot be correct for
both. This is the existing report-fiber obstruction. With only outer bounds,
some candidates in \(C(h)\) may be spurious; abstention can then be conservative.

An explicit parser is part of this statement. The finite reference code
accepts a dictionary with a model digest, a string-valued free observation,
and a list of distinct named instrument readings. Each reading contains the
expected raw-unit label and a canonical rational string, such as \texttt{-1/2}.
Unknown fields, wrong units, repeated instruments, floating-point inputs,
and a digest for another model are rejected. Calibration coefficients,
enclosures, prices, and the answer table come from the caller's separately
supplied model. The model digest binds those data; it does not authenticate
the source of a reading or a calibration standard. Soundness remains
conditional on that source and on the bounds above.

\subsubsection{X.3 Exact finite representation of continuous reports}\label{x.3-exact-finite-representation-of-continuous-reports}

For each instrument define its membership signature

\[S_i(z)=\{x\in X:z\in J_i(x)\},\qquad
  \mathcal S_i=\{S_i(z):z\in\mathbb R,\ S_i(z)\ne\varnothing\}.\]

A signature retains exactly which candidates are compatible with one
reading. This does not round the reading to a nearby center.

\textbf{Theorem X.3 (finite report reduction).} The interval product model has
the same possible answer histories and the same optimal fixed and adaptive
worst-case observation costs as the finite deterministic model

\[\widehat X=\{(x,s):x\in X,\ s\in\prod_i\mathcal S_i,
                         \ x\in s_i\text{ for every }i\},\]

\[\widehat p(x,s)=p(x),\qquad \widehat q(x,s)=q(x),\qquad
  m_i(x,s)=s_i,\]

with the same instrument prices. A compatible original candidate is exactly
the projection of a compatible lifted state. There are at most \(4|X|-1\)
nonempty signatures for each instrument.

\textbf{Proof.} List the distinct interval endpoints in increasing order, say
\(e_1<\cdots<e_k\), with \(k\le2|X|\). Membership is constant on each singleton
\(\{e_j\}\) and each open interval \((e_j,e_{j+1})\). Both exterior intervals
have empty membership. Keep the nonempty signatures and merge equal ones.
There are at most \(2k-1\le4|X|-1\). In particular the singleton endpoints
must be retained: two closed intervals may meet at just one point. Rational
endpoints give rational representatives, using endpoints and midpoints.
Consequently restriction to exact rational reports retains every signature
and both costs for rational data, including a nonzero rational raw conversion.

A continuous vector admitted for \(x\) yields a lifted state by applying
\(S_i\) to each coordinate. Conversely choose one representative \(z_i(s_i)\)
for each nonempty signature. If \(x\in s_i\), that representative lies in
\(J_i(x)\). Full report realizability makes the vector of representatives an
admitted vector for \(x\).

For a partial history, \(x\in C(h)\) precisely when it belongs to all the
observed signatures and has the right free report. Each unqueried interval
is nonempty, so select one of its signatures to complete a lifted state.
The reverse implication follows from the observed coordinates. Hence the
projection of every compatible lifted fiber is exactly \(C(h)\).

A strategy on signatures works on continuous readings by applying \(S_i\)
after each query. Every resulting history is admitted in the finite model,
so correctness and worst-case cost bounds transfer. In the other direction,
run any continuous-reading strategy on the selected representatives
\(z_i(s_i)\) when a signature is observed. Every such path is admissible for
its underlying candidate, including after adaptive choices, because every
partial product assignment extends to a full one. It follows that the
finite strategy has no greater worst-case cost. The two translations equate
the infima. There are finitely many deterministic trees on the finite
alphabet with depth at most \(|I|\), so every feasible finite optimum is
attained; its continuous translation attains the same value. Fixed suites
are unchanged in both translations. This establishes equality of both optima.
\(\square\)

The theorem compares information and cost, not the literal numeric contents
of histories. Without full report realizability, every actual history still
maps into this finite model, so its certified strategies remain sound upper
bounds. The converse translation, and consequently physical optimality or
impossibility, need not hold.

\subsubsection{X.4 Separation and robust cost bounds}\label{x.4-separation-and-robust-cost-bounds}

Write \(F_{\rm box}\) and \(A_{\rm box}\) for the optimal fixed and adaptive
costs of the interval product model. The fixed suite is chosen before the
free report; an adaptive strategy receives that report before its first
paid query. An infeasible optimum is \(+\infty\).

\textbf{Corollary X.4 (fixed robust cost).} A fixed suite \(J\subseteq I\)
ensures recovery under declared assumptions exactly when every pair \(x,y\) with \(p(x)=p(y)\) and
\(q(x)\ne q(y)\) has some \(i\in J\) with

\[J_i(x)\cap J_i(y)=\varnothing,
  \quad\text{equivalently}\quad
  |f_i(x)-f_i(y)|>\epsilon_i(x)+\epsilon_i(y).\]

Thus \(F_{\rm box}\) is the existing minimum weighted hitting-set cost on
these separating-instrument sets.

\textbf{Proof.} A disjoint coordinate prevents the pair from sharing a report.
If no queried coordinate is disjoint, choose a common value in each of
its intersections. Complete the unqueried coordinates separately for
\(x\) and \(y\). The two resulting lifted states have the same observed report
and different answers. The equivalence with the strict inequality is the
closed-interval intersection criterion. Section 14.6 now applies to the
finite lift. \(\square\)

The strict inequality matters: equality leaves a common endpoint. Under
outer-bound validity alone the separation condition is sufficient, but its
failure is not proof of actual indistinguishability.

For completeness, the existing adaptive recursion becomes particularly
simple on a pair \((C,R)\) of a nonempty compatible candidate set and the
remaining instruments. It is zero if \(q\) is constant on \(C\). Otherwise,

\[W(C,R)=\min_{i\in R}\left[c_i+
       \max_{s\in\mathcal S_i:\ C\cap s\ne\varnothing}
                W(C\cap s,R\setminus\{i\})\right],\]

where the minimum over an empty set is \(+\infty\). The initial adaptive cost
is the maximum of \(W(p^{-1}(p_0),I)\) over nonempty free-report fibers.
The recursion terminates because \(R\) decreases, including when a report
leaves \(C\) unchanged. Tracking only \(C\) and assuming every branch strictly
shrinks it would be incorrect. By Theorem X.3 these are the same minimax
costs as Section 14.6 on the lift; its complete-table checker can be reused.
The maximum lists all realizable adversarial outcomes, not selected examples.

\textbf{Corollary X.5 (comparison and stability).} Let \(F_0,A_0\) be the costs
for the nominal deterministic reports \(f_i(x)\). Then

\[A_{\rm box}\le F_{\rm box},\qquad
  F_0\le F_{\rm box},\qquad A_0\le A_{\rm box}.\]

Enlarging any intervals while retaining their centers cannot decrease either
product-model optimum. Moreover, suppose each instrument separates all of
its distinct nominal report values despite the errors:

\[f_i(x)\ne f_i(y)\ \Longrightarrow\
 |f_i(x)-f_i(y)|>\epsilon_i(x)+\epsilon_i(y)
 \qquad(i\in I,\ x,y\in X).\]

Then \(F_{\rm box}=F_0\) and \(A_{\rm box}=A_0\).

\textbf{Proof.} Querying a sufficient fixed suite is an adaptive strategy with
no greater worst-case cost. A nominal vector is an admitted product-model
vector, so restriction of any robust strategy to nominal reports establishes
the other two lower bounds. The same restriction argument establishes
monotonicity for nested supports. Under the separation hypothesis, any
compatible reading determines its unique nominal-value group. Decode that
group and follow an optimal nominal strategy or fixed suite. The true
candidate remains in the decoded group, so this strategy is correct at
its nominal cost. Combine this upper bound with the lower bounds. \(\square\)

For an actual correlated model \(L_x\subseteq\prod_iJ_i(x)\), only
\(F_{\rm actual}\le F_{\rm box}\) and
\(A_{\rm actual}\le A_{\rm box}\) follow without further assumptions.
The product-model lower bounds must not be reported as physical lower bounds.
A finite upper certificate ensures recovery under declared assumptions within its budget; under
full realizability, a budget below the exact optimum cannot guarantee an
answer on every admitted history. A sound procedure must then sometimes
abstain, request a larger budget, or leave the stated model.

\subsubsection{X.5 Changes of units}\label{x.5-changes-of-units}

\textbf{Proposition X.6 (coordinate invariance).} For any invertible affine
change of each calibrated coordinate,

\[z'_i=\lambda_i z_i+\tau_i,\quad
 f'_i(x)=\lambda_i f_i(x)+\tau_i,\quad
 \epsilon'_i(x)=|\lambda_i|\epsilon_i(x),\qquad\lambda_i\ne0,\]

candidate compatibility, separation, and both observation costs are unchanged,
provided the instrument prices retain their original cost unit. For ordinary
changes of physical units \(\lambda_i>0\); a negative value also covers a
reversal of coordinate orientation.

\textbf{Proof.} The image interval has the stated center and radius.
Invertibility gives \(z_i\in J_i(x)\) if and only if
\(z'_i\in J'_i(x)\). Thus all signatures and their allowed candidate memberships
are unchanged. Both sides of the separation inequality are multiplied by
\(|\lambda_i|\). Apply Theorem X.3 to the identical finite observation table.
Scaling the centers but leaving radii or unit tags unchanged is not this
coordinate change. \(\square\)

\subsubsection{X.6 Exact examples and a necessary counterexample}\label{x.6-exact-examples-and-a-necessary-counterexample}

Take four candidates with distinct protected answers and one constant free
report. Their nominal calibrated reports and prices are

\begin{longtable}[]{@{}lrrrrr@{}}
\toprule\noalign{}
Instrument & 00 & 01 & 10 & 11 & Price \\
\midrule\noalign{}
\endhead
\bottomrule\noalign{}
\endlastfoot
A & 0 & 0 & 10 & 10 & 1 \\
B & 0 & 10 & 0 & 0 & 3 \\
C & 0 & 0 & 0 & 10 & 3 \\
\end{longtable}

Give every report radius \(\epsilon\) in the same calibrated unit. For
\(0\le\epsilon<5\), Corollary X.5 applies. The fixed cost is exactly \(7\):
the pairs \((00,01)\), \((00,10)\), and \((10,11)\) force B, A, and C,
respectively. All three suffice. The adaptive cost is exactly \(4\): query
A, then B on its low branch or C on its high branch. A budget below \(4\)
cannot query A and a second distinct instrument; querying B or C alone
also leaves different answers unresolved. At \(\epsilon=5\), the vector
\((5,5,5)\) is compatible with all four candidates, so both optima are infinite
in the product model. Larger radii retain this obstruction. The reference
uses the raw conversion \(z=2r+1\); changing coordinate units preserves the
finite costs when the whole error interval is transported.

Full report realizability cannot be dropped from the lower-bound argument.
Let \(x\in\{0,1\}\), with constant free observation and protected answer \(x\),
and let the same unknown
\(e\in[-1,1]\) affect two unit-price instruments:

\[z_1=x+e,\qquad z_2=-x+e.\]

Their individual intervals are \([-1,1]\) for \(x=0\), and \([0,2]\), \([-2,0]\)
for \(x=1\). Each instrument alone has overlapping supports. The interval
product model admits the report \((0,0)\) for both candidates and therefore
has infinite robust recovery cost. But the actual correlated model obeys
\((z_1-z_2)/2=x\) for every admitted \(e\). Two queries recover exactly, and
one cannot guarantee recovery because of the individual overlaps. Its
actual fixed and adaptive optima are both \(2\). This discrepancy persists
even though every single-instrument interval is exact.

\subsubsection{X.7 What is checked}\label{x.7-what-is-checked}

\href{code/calibrated_recovery/bridge.py}{The reference checker} constructs the finite lift and checks
answer proposals against a separately supplied model. It reuses
\texttt{finite\_observation.solver} for cost proposals and for complete-table
verification; the verification path never calls the solver to replace a
bad certificate. The reference has exact rational endpoints and refuses a
lift above the existing ten-state limit. A resource refusal is not mathematical
infeasibility. The four-candidate endpoint collision is checked directly,
since its complete lift exceeds that limit. No new native evaluator or
production observation-cost solver is claimed.

The cost certificate explicitly names \texttt{ideal\_calibrated\_interval\_product}.
It certifies the full interval product before raw-encoding restrictions. The
raw-report API separately requires canonical rational strings of at most 160
characters and numerators/denominators of at most 256 bits, before calibration.
Model admission does not guarantee closure under inverse calibration within
those bounds. For example, gain \(2^{-255}\), offset \(0\), and point supports
\(\{0\}\) and \(\{2^{255}\}\) admit a two-state model with ideal fixed and adaptive
cost \(1\). The high state's only raw reading is \(2^{510}\) (511 bits), refused
as \texttt{INVALID\_REPORT}; the low state's raw reading \(0\) verifies. This is a
bounded API refusal, not ideal-model infeasibility or acceptance of a wrong
answer. Applying the ideal costs to this API additionally requires proof that
the needed reports have admitted raw encodings; this certificate does not
provide that proof.

\href{code/calibrated_recovery/TransferLaws.lean}{The six transfer lemmas} compile
using Lean core alone. Its six declarations
establish partial-product extension, the exact projection of compatible
lifted states, answer transfer, preservation under outer supports, answer
soundness, and the shared-report obstruction. The first three use classical
choice. These declarations do not formalize real interval arithmetic,
physical calibration, or the cost comparison. Those proofs are written
above and use the finite cost results already stated in this paper.
\href{code/unit_tests/test_calibrated_recovery.py}{The Python controls} test
endpoint collisions, exact cost certificates,
unit changes, invalid reports, false answers, substituted models, and
correlated-noise limits. They are finite deterministic tests, not a device
experiment or statistical validation.

\subsection{Appendix Y --- Exact recovery costs from joint report supports}\label{appendix-y-exact-recovery-costs-from-joint-report-supports}

Appendix X makes exact cost claims for products of intervals. This appendix
retains the correlations between reports. It first characterizes the costs
for arbitrary declared joint supports. It then gives a finite construction
for a specified class of shared affine errors, including the example in
which the product enclosure has infinite cost but two actual queries suffice.
The costs concern deterministic, exact, worst-case recovery of a supplied
answer. They do not price a statistical confidence level or authorize a query.

\subsubsection{Y.1 Joint supports and compatible histories}\label{y.1-joint-supports-and-compatible-histories}

Let \(X\) be nonempty and finite, with free report \(p:X\to P\) and protected
answer \(q:X\to Q\). Let \(I\) be a finite instrument shelf with positive
rational prices \(c_i\), all in one declared cost unit. For each \(x\), supply a
nonempty set

\[L_x\subseteq\mathbb R^I.\]

A world is a pair \((x,z)\) with \(z\in L_x\); write \(\Omega\) for their set.
The protected answer of this world is \(q(x)\).
One complete vector is fixed throughout the experiment. Querying \(i\)
reveals \(z_i\) and leaves the world unchanged. Each instrument is queried
at most once. As in Appendix X, strategies are deterministic functions of
the exact reports received. The general statements below impose no
computability, measurability, or probability model on those functions.

For a history \(h=(p_0,H,r_H)\), define

\[\begin{aligned}
B(h)&=\{(x,z)\in\Omega:p(x)=p_0,\ z_i=r_i\text{ for every }i\in H\},\\
C(h)&=\{x:\exists z,\ (x,z)\in B(h)\}.
\end{aligned}\]

Equivalently, \(x\in C(h)\) precisely when \(p(x)=p_0\) and
\(r_H\in\pi_H(L_x)\), where \(\pi_H\) is coordinate projection.
The same complete vector must extend every reading in the history.
Checking each coordinate against its marginal support does not establish
this condition.

A proposed answer \(v\) is certified when \(C(h)\ne\varnothing\) and
\(q(x)=v\) for every \(x\in C(h)\). If the true world belongs to the declared
support, this answer is correct. If two compatible candidates have
different answers, their admitted complete vectors realize the same
history; a decoder of that history cannot be correct for both. Empty
compatibility means inconsistency with the model, not acceptance.
This is the report-fiber criterion applied to actual joint histories.

These distinctions also govern physical use. If the actual support is
only known to be a nonempty subset of \(L_x\), certification remains sound,
but abstention can be conservative. Exact physical lower costs require
that the declared joint supports are realizable. A bound on each
coordinate does not establish that premise.

\subsubsection{Y.2 Fixed suites and projected intersections}\label{y.2-fixed-suites-and-projected-intersections}

\textbf{Proposition Y.1 (exact fixed cost).} A fixed suite \(J\subseteq I\),
chosen before receiving the free report, is sufficient exactly when

\[\pi_J(L_x)\cap\pi_J(L_y)=\varnothing
\quad\text{whenever}\quad p(x)=p(y),\ q(x)\ne q(y).\]

Consequently its optimum is

\[F_\Omega=\min\left\{\sum_{i\in J}c_i:
J\subseteq I\text{ satisfies this condition}\right\},\]

with value \(+\infty\) when no suite is sufficient.

\textbf{Proof.} A point in a projected intersection is supplied by complete
worlds with the same free and queried reports but different answers.
Conversely, failure of answer constancy on a report fiber supplies just
such an intersection. There are finitely many suites, so a finite minimum
is attained. \(\square\)

This remains a hitting-set problem, but the obligations are now pairs
of complete worlds. For \(p(x)=p(y)\) and \(q(x)\ne q(y)\), put

\[E_{(x,z),(y,w)}=\{i\in I:z_i\ne w_i\},
\qquad z\in L_x,\ w\in L_y.\]

A suite is sufficient precisely when it meets every such set. Although
there can be infinitely many world pairs, there are at most \(2^{|I|}\)
distinct sets of this form. Identifying which sets are realized can
require a support-membership or existence oracle; the finite count is
not an algorithm for an arbitrary supplied set.

In particular, recovery using the whole shelf is feasible exactly when
the full supports \(L_x,L_y\) are disjoint for every cross-answer pair
with the same free report. A shared complete vector obstructs every
adaptive strategy too. When there is no such vector, querying the whole
shelf and checking its fiber supplies a correct mathematical decoder.
Marginal overlap by itself is not this obstruction.

\subsubsection{Y.3 Adaptive cost depends on the residual joint support}\label{y.3-adaptive-cost-depends-on-the-residual-joint-support}

Let \(B\) be a nonempty set of worlds contained in one already observed
free-report fiber, and let \(R\) be the remaining instruments. No further
free observation is available. Write \(V(B,R)\) for the least worst-case additional cost of a strategy
that returns the correct answer on every world in \(B\). If no such strategy
exists, set \(V(B,R)=+\infty\). A strategy that sometimes abstains does not
establish a finite value of this guaranteed-recovery objective.
For convenience only, set \(V(\varnothing,R)=0\) on unreachable branches.

For \(i\in R\), let

\[B_{i,t}=\{(x,z)\in B:z_i=t\},
\qquad T_i(B)=\{z_i:(x,z)\in B\}.\]

\textbf{Theorem Y.2 (joint-history recursion).} If \(q\) is constant on \(B\),
then \(V(B,R)=0\). Otherwise,

\[V(B,R)=\min_{i\in R}\left[c_i+
\max_{t\in T_i(B)}V(B_{i,t},R\setminus\{i\})\right].\]

An empty minimum is \(+\infty\). Every displayed maximum is attained,
even if \(T_i(B)\) is infinite. A finite value is attained by a strategy.
The initial adaptive optimum is

\[A_\Omega=\max_{p_0\in p(X)}
V(\{(x,z)\in\Omega:p(x)=p_0\},I).\]

\textbf{Proof.} Induct on \(|R|\). With no remaining query, a nonconstant answer
cannot be recovered. For the induction step, every finite child value
belongs to the finite set of subset sums

\[\mathcal K_{R\setminus\{i\}}
=\left\{\sum_{j\in S}c_j:S\subseteq R\setminus\{i\}\right\}.\]

Thus the nonempty set of child values lies in
\(\mathcal K_{R\setminus\{i\}}\cup\{+\infty\}\) and has an attained maximum.
Any correct strategy with first query \(i\) has worst-case cost at least
the corresponding displayed option. For a finite minimum, choose a
minimizing \(i\) and use an optimal continuation on each realized branch.
One can define these choices at all histories by fixing an order on \(I\)
and always taking the first minimizing instrument; recursion decreases
\(|R|\). This supplies a strategy attaining the value. Its value is again
a subset sum, since the maximum selects one child value and adds a
previously unqueried price. The same argument handles infeasible
branches with the extended value \(+\infty\). The finite free-report
image supplies the last maximum. \(\square\)

Restriction to a smaller set of worlds cannot increase either cost.
Also \(A_\Omega\le F_\Omega\), by asking a sufficient fixed suite in order.
Neither fact requires independent errors.

It is essential to retain \(B(h)\), or an exact representation of its
continuations. The pair consisting only of \(C(h)\) and the remaining
instruments is generally insufficient. For example, let the free report
be constant, the answer be \(x\in\{0,1\}\), and

\[L_0=\{(e,e):0\le e\le1\},\qquad
L_1=\{(e,2e-\tfrac13):0\le e\le1\}.\]

After a first reading of \(0\) or \(1/3\), the same candidates \(\{0,1\}\)
remain. At \(0\), a second unit-price reading distinguishes them. At
\(1/3\), both coordinates agree and no remaining query can distinguish
them. Their remaining costs are \(1\) and \(+\infty\), respectively.

\subsubsection{Y.4 Finite cost witnesses exist, but are not report encoders}\label{y.4-finite-cost-witnesses-exist-but-are-not-report-encoders}

The preceding recursion gives an existence result for arbitrary joint
supports, without claiming an effective way to inspect them.

\textbf{Theorem Y.3 (finite cost witness).} There is a finite
\(W\subseteq\Omega\), meeting each candidate state, whose fixed and adaptive
optima equal \(F_\Omega\) and \(A_\Omega\). If \(m=|I|\), one bound is

\[|W|\le |X|+2^{m+1}+\tau_m,\qquad
\tau_0=2,\quad \tau_r=2+r\tau_{r-1}.\]

\textbf{Proof.} First consider any nonempty \(B\) contained in one fixed
free-report fiber, as in Y.3, and an instrument set \(R\) of size \(r\).
We show that a nonempty finite \(S\subseteq B\) preserves \(V(B,R)\) and
has at most \(\tau_r\) worlds. If the answer is constant, one world suffices.
For \(r=0\) and a nonconstant answer, choose two worlds with different
answers. For \(r>0\) and a nonconstant answer, choose two such worlds and,
for every \(i\in R\), choose a reading \(t_i\) attaining the maximum in
Theorem Y.2. By induction take \(S_i\subseteq B_{i,t_i}\) preserving that
branch's remaining value, with \(|S_i|\le\tau_{r-1}\). Let \(S\) contain
the chosen pair and every \(S_i\).

Restriction gives \(V(S,R)\le V(B,R)\). Since the chosen pair has different
answers, any correct strategy on \(S\) must start with a query, say \(i\).
Its \(t_i\) branch contains \(S_i\). Monotonicity therefore gives a remaining
cost at least \(V(S_i,R\setminus\{i\})\), which is the worst child value for
query \(i\) on \(B\). Minimizing over possible first queries gives the reverse
inequality. This argument also applies when a value is infinite. The
cardinality is at most \(2+r\tau_{r-1}\).

Choose a free-report fiber attaining \(A_\Omega\) and retain the preceding
finite witness for that fiber. For every insufficient fixed suite, choose
two admitted worlds with the same reports on that suite and different
answers. There are at most \(2^m\) suites. Finally choose one admitted world
per candidate state. Let \(W\) be the union of these selections. Every
insufficient fixed suite still fails on its retained pair, while every
sufficient suite stays sufficient after restriction. The retained worst
free-report fiber gives \(A_W\ge A_\Omega\), and restriction gives the
opposite inequality. Counting the selections yields the stated bound.
\(\square\)

The selections use full support information and exact branch optima.
This is an existence proof, not an algorithm that constructs a witness
from arbitrary set descriptions. A guessed finite sample alone supplies
lower costs, which can be strict. Its finite certificate becomes an exact
certificate for the declared model only with a valid completeness argument.
Nor does \(W\) represent every literal report or supply a decoder for every
unsampled history.

There need not even be a sufficient finite encoding of each coordinate
separately. For two unit-price instruments, take

\[L_0=\{(u,u):u\in[0,1]\},\qquad
L_1=\{(u,v):u,v\in[0,1],\ u\ne v\},\]

with constant free report and answer \(0\) or \(1\). The exact fixed and
adaptive costs are both \(2\). For any finite-valued first-coordinate
encoding, some distinct \(a,b\) have the same code. Then \((b,b)\in L_0\)
and \((a,b)\in L_1\) have identical encoded reports, regardless of the
second-coordinate encoding. The encoded problem cannot recover the
answer. This example uses arbitrary, not necessarily closed supports;
it makes no assertion about finite-precision instrument performance.
It separates a finite witness of a cost from a finite encoding of all
the reports needed for recovery.

\subsubsection{Y.5 A constructive reduction for shared affine error}\label{y.5-a-constructive-reduction-for-shared-affine-error}

A useful explicit class does admit a finite cost construction. Index a
finite nonempty family of segments by \(k\in K\). Each has an owner
\(x(k)\in X\), with every candidate owning at least one segment, and

\[z_i=a_{ki}e+b_{ki},\qquad
e\in[\ell_k,u_k],\quad \ell_k\le u_k,\quad a_{ki}\ne0.\]

All endpoints and coefficients are rational. The joint support \(L_x\) is
exactly the union of the segments owned by \(x\). Different components may
overlap. A component label and its parameter are hidden variables, not
free observations. Degenerate parameter intervals are permitted. The
nonzero-slope condition is imposed for every component and instrument.
The case \(I=\varnothing\) is handled by the free report alone.

Fix a free report \(p_0\) and a possible first query \(i\). On component \(k\)
with \(p(x(k))=p_0\), its value \(t\) lies in the closed interval \(D_{ki}\),
the image of \([\ell_k,u_k]\) under \(e\mapsto a_{ki}e+b_{ki}\).
That reading fixes the component's parameter and predicts every later
coordinate as

\[g^{\,i}_{kj}(t)
=\frac{a_{kj}}{a_{ki}}\,t+
b_{kj}-\frac{a_{kj}}{a_{ki}}\,b_{ki}.\]

Make a finite set of critical values from:

\begin{enumerate}
\def\labelenumi{\arabic{enumi}.}
\tightlist
\item
  every endpoint of every \(D_{ki}\) in this free-report fiber;
\item
  every isolated equality point
  \(g^{\,i}_{kj}(t)=g^{\,i}_{lj}(t)\), for every component pair and every
  remaining coordinate \(j\), that lies in \(D_{ki}\cap D_{li}\).
\end{enumerate}

Identical affine functions agree throughout their common domain; parallel
unequal functions never agree. Neither contributes an isolated point.
Partition the range at these critical values, retaining singleton points
as well as open intervals between consecutive points. Use the point itself
or the interval midpoint as a representative. Exterior intervals contain
no admitted first reading.

On each cell, the active component indices are constant, as are all
equalities between their predicted values for every remaining instrument.
Hence the finite remaining observation problems at any two readings in
the same cell differ only by a relabeling of outcomes. Their protected
answers and costs agree. Overlapping components that describe the same
candidate and full vector may be identified; duplicating such a world
does not create an observable distinction.

Construct \(W\) by retaining one endpoint world from each segment and, for
every \(p_0,i\) and cell whose compatible answers are not all equal, every
world in the complete first-reading fiber at its representative. A first
reading determines at most one world per component, so every such fiber
is finite. Cells with one answer need no further observation.

\textbf{Theorem Y.4 (affine critical-value reduction).} This construction is
finite and has

\[F_W=F_\Omega,\qquad A_W=A_\Omega.\]

It therefore supplies an instance for the existing finite solver and
certificate checker. For rational input, every retained world is rational.

\textbf{Proof.} Activity can change only at an image endpoint. An equality of
two predicted affine values can change only at an isolated root, unless
the functions are identical. This establishes the cell statement above,
including equality at a single critical point.

For fixed suites, \(W\subseteq\Omega\) gives one direction. Suppose a
nonempty suite \(J\) fails on \(\Omega\). Choose its conflicting admitted
worlds and some \(i\in J\). Their common \(i\) value belongs to a cell.
The same two components are active at its representative and continue
to agree on every coordinate in \(J\), while their answers differ. This
representative fiber was included in \(W\), so \(J\) fails there too. For
\(J=\varnothing\), the endpoint worlds preserve the candidate answers in
every free-report fiber. Thus precisely the same fixed suites suffice.

For the adaptive result, first consider a nonconstant free-report fiber.
For any proposed first query \(i\), every nonconstant continuation problem
is equivalent to a representative problem included in \(W\). Moreover the
fiber of \(W\) at that representative is the complete fiber of \(\Omega\):
it contains all its worlds, and every sampled world is admitted in
\(\Omega\). Constant-answer continuations have value zero. The remaining
values have only finitely many cell types, so the largest continuation
value is represented whenever it is positive or infinite. If all
continuations have value zero, any nonempty restricted fiber has that
value too. Thus no strategy on \(W\) with first query \(i\) can have a lower
worst-case cost than the corresponding option on \(\Omega\). Theorem Y.2
and minimization over first queries give \(A_W\ge A_\Omega\). Restriction
gives the reverse inequality. Constant-answer free-report fibers have
zero cost in both models. \(\square\)

There are finitely many endpoints and pairwise affine roots; enumerating
them requires only exact scalar rational arithmetic. With \(n=|K|\) and
\(m=|I|\ge1\), a coarse bound on the number of sampled worlds is

\[n+mn\left(4n+2(m-1)\binom n2-1\right).\]

Indeed each free-report fiber has at most
\(2n+(m-1)\binom n2\) critical values. It contributes at most twice this
number minus one cells, each with at most its number of components.
Summing the component counts over free-report fibers gives \(n\).
This bounds the finite construction, not the subsequent exponential
cost-certificate computation.

A complete-vector collision with the same free report and different
protected answers gives an earlier exit: retain the two admitted
conflicting worlds. Both costs are then
infinite by Proposition Y.1. For affine segments this collision can be
checked by requiring all predicted-coordinate equalities at one common
first-coordinate value, within both parameter-image intervals.
This shortcut is separate from truncating an oversized sample.

The reduction also explains how a finite optimal policy can be used on
continuous reports. If the free report already determines the answer,
return it without querying an instrument. Otherwise take the finite
policy's first query for that free report.
If the returned value leaves one answer, stop. Otherwise retain that
exact value \(t\), find its cell representative \(s\), and use the complete
finite fiber at \(s\). For a later coordinate, map the actual observed
value to \(g^{\,i}_{kj}(s)\) using any initially active component whose
prediction at \(t\) equals that observation. The map is well-defined
because predicted-value equalities are constant on the cell. Filter by
the accumulated history and follow the finite policy. This preserves
answers and never exceeds the certified worst-case cost. It is an
encoding that uses the first reading and the history; it is not a fixed
finite code for each coordinate in isolation.

\subsubsection{Y.6 Exact examples and a misleading sampling baseline}\label{y.6-exact-examples-and-a-misleading-sampling-baseline}

For the shared-error supports

\[L_x=\{(x+e,-x+e):-1\le e\le1\},\qquad x\in\{0,1\},\]

with constant free report, protected answer \(x\), and unit prices, neither
coordinate alone suffices, but \((z_1-z_2)/2=x\). Both exact costs are \(2\).
The reference construction retains ten worlds and its finite certificate
has the same costs. The product of the marginal intervals instead admits
\((0,0)\) for both answers and has infinite cost, as shown in Appendix X.
At the actual joint report \((0,0)\), only \(x=0\) is compatible.

Multiple components can encode correlation too. Let

\[L_0=\{(0,0),(1,1)\},\qquad
L_1=\{(0,1),(1,0)\}.\]

These are degenerate affine segments. Their marginal supports are
identical, but their joint supports are disjoint. Both unit-price
queries are necessary and sufficient. Matching the first reading to
one component and the second reading to another component of the same
candidate would give a false compatibility result.

The crossing example in Y.3 has the shared complete report
\((1/3,1/3)\). Both costs are infinite. A grid that happens to omit this
point does not establish feasibility.

An isolated critical value can also change a finite optimum. Take
three instruments A, B, C, with prices \(1,1,3\), and

\[L_0=\{(e,e,e):-1\le e\le1\},\qquad
L_1=\{(e,-e,e+10):-1\le e\le1\}.\]

The free report is constant and the protected answer is the support
index. Instrument C alone separates its disjoint ranges, at cost \(3\).
Any budget below \(3\) can use only A and B. Their joint report \((0,0)\)
is compatible with both answers, so no such budget suffices. The fixed
and adaptive optima are therefore both \(3\).

Sampling only the parameter endpoints yields four worlds on which
A and B suffice at cost \(2\). This finite answer is correct for that
sample and false for the declared continuous model. The critical-value
construction retains the interior collision of A and B, as well as
the differing C values, and returns \(3\). This comparison tests the
completeness of the reduction, not a statistical estimate of noise.

\subsubsection{Y.7 Units, calibration and additional error}\label{y.7-units-calibration-and-additional-error}

The calibrated coordinate \(z_i\) and its intercept \(b_{ki}\) have the
instrument's declared output unit. The parameter \(e,\ell_k,u_k\) have
the component's declared parameter unit; \(a_{ki}\) has their unit ratio.
Raw reports use the declared nonzero conversion
\(z_i=\gamma_i r_i+\delta_i\), as in Appendix X. All admitted effects of
calibration must already be covered by the supplied joint support.
The unit labels and a model digest cannot establish that physical premise.

For a coordinate change \(z'_i=\lambda_i z_i+\tau_i\), with
\(\lambda_i\ne0\), transport every component by

\[a'_{ki}=\lambda_i a_{ki},\qquad
b'_{ki}=\lambda_i b_{ki}+\tau_i,\]

and transport the raw conversion in the same way. This is a bijection
of full and partial report spaces, preserving their compatible answers.
It conjugates any query strategy without changing the prices. Hence
both costs are unchanged. Negative \(\lambda_i\) reverses orientation.
Changing a unit label or only a nominal center is not this transformation.

A segment's marginal interval has center
\(a_{ki}(\ell_k+u_k)/2+b_{ki}\) and radius
\(|a_{ki}|(u_k-\ell_k)/2\), with the corresponding output unit. Their
product is generally larger than the segment. It remains an outer
bound, with the limitations already stated in Appendix X.

The affine-segment algorithm does not cover an unmodeled independent
rounding or calibration residual added to each coordinate. Such a
residual thickens the support. After one reading, the latent parameter
need not be uniquely determined, so the finite-fiber argument in Y.5
does not apply as written. The general joint-support criterion and
recursion still apply to a correctly enlarged support; an effective
reduction for that larger class is a separate obligation.

\subsubsection{Y.8 Reference implementation and proof scope}\label{y.8-reference-implementation-and-proof-scope}

\href{code/joint_recovery/reference.py}{The joint-support reference} validates
exact rational segment data, checks reports against a single common
parameter in one component, constructs the critical finite witness, and
reuses the existing finite observation solver and complete-table checker.
Cost verification reconstructs the witness from the separately supplied
model and checks the submitted certificate; it never invokes the proposal
solver to replace a rejected certificate.

The model digest binds all components, unit labels, raw conversions,
free reports, answer labels, prices and their common cost unit. It is not
source authentication.
The report parser reuses the bounded canonical-rational parser from
Appendix X and rejects wrong units, repeated instruments, unknown fields
and substituted models. A checked answer supplies no permission to act.

The reference admits at most ten segments and twelve instruments,
with 256-bit rational input numerators and denominators. It refuses a
constructed witness above the existing ten-world solver limit rather
than deleting worlds or reporting infeasibility. The collision shortcut
uses a genuine two-world obstruction. Duplicate descriptions of the
same world are removed before the cap is applied.

\href{code/joint_recovery/certificates.json}{The stored examples} and
\href{code/unit_tests/test_joint_recovery.py}{targeted controls} cover shared
error, unions of components, isolated collisions, a missed interior cost
peak, free observations, unit changes, certificate tampering and resource
refusals. Degenerate supports are compared with a direct finite-checker
baseline. These are exact deterministic examples, not instrument trials
or evidence that a physical error model is valid.

The results in this appendix have written proofs. They are not new
Lean formalizations, and no general joint-support decision procedure
is implemented. The earlier six core Lean transfer lemmas retain their
stated scope. The reference certifies costs for its admitted affine
models; it does not implement a deployed controller for the continuous
policy transformation described above.

\subsection{Limitations \& Explicit Nonclaims}\label{limitations-explicit-nonclaims}

This paper makes narrow, scoped claims under declared assumptions. Operational authority is null. Computational demonstrations and finite checks do not constitute unmechanized carrier-scoped results or institutional authorization.

\end{document}
